% This is pdfTeX, Version 3.14159265-2.6-1.40.19 (MiKTeX 2.9.6880)
% LaTeX2e <2018-12-01>

% Package: `showlabels' v1.7 <2015/12/08>
% Package: aliasctr 2014/04/21 v66
% Package: amsbsy 1999/11/29 v1.2d Bold Symbols
% Package: amsmath 2018/12/01 v2.17b AMS math features
% Package: amsopn 2016/03/08 v2.02 operator names
% Package: amstext 2000/06/29 v2.01 AMS text
% Package: amsthm 2017/10/31 v2.20.4
% Package: atbegshi 2016/06/09 v1.18 At begin shipout hook (HO)
% Package: atveryend 2016/05/16 v1.9 Hooks at the very end of document (HO)
% Package: auxhook 2016/05/16 v1.4 Hooks for auxiliary files (HO)
% Package: bigintcalc 2016/05/16 v1.4 Expandable calculations on big integers (HO
% Package: bitset 2016/05/16 v1.2 Handle bit-vector datatype (HO)
% Package: bm 2017/01/16 v1.2c Bold Symbol Support (DPC/FMi)
% Package: calc 2017/05/25 v4.3 Infix arithmetic (KKT,FJ)
% Package: cleveref 2018/03/27 v0.21.4 Intelligent cross-referencing
% Package: enumitem 2018/11/30 v3.6 Customized lists
% Package: epstopdf-base 2016/05/15 v2.6 Base part for package epstopdf
% Package: etexcmds 2016/05/16 v1.6 Avoid name clashes with e-TeX commands (HO)
% Package: etoolbox 2018/08/19 v2.5f e-TeX tools for LaTeX (JAW)
% Package: everyshi 2001/05/15 v3.00 EveryShipout Package (MS)
% Package: expl3 2018-12-12 L3 programming layer (code)
% Package: expl3 2018-12-12 L3 programming layer (loader)
% Package: gettitlestring 2016/05/16 v1.5 Cleanup title references (HO)
% Package: graphics 2017/06/25 v1.2c Standard LaTeX Graphics (DPC,SPQR)
% Package: graphicx 2017/06/01 v1.1a Enhanced LaTeX Graphics (DPC,SPQR)
% Package: grfext 2016/05/16 v1.2 Manage graphics extensions (HO)
% Package: hobsub 2016/05/16 v1.14 Construct package bundles (HO)
% Package: hobsub-generic 2016/05/16 v1.14 Bundle oberdiek, subset generic (HO)
% Package: hobsub-hyperref 2016/05/16 v1.14 Bundle oberdiek, subset hyperref (HO)
% Package: hopatch 2016/05/16 v1.3 Wrapper for package hooks (HO)
% Package: hycolor 2016/05/16 v1.8 Color options for hyperref/bookmark (HO)
% Package: hyperref 2018/11/30 v6.88e Hypertext links for LaTeX
% Package: ifluatex 2016/05/16 v1.4 Provides the ifluatex switch (HO)
% Package: ifpdf 2018/09/07 v3.3 Provides the ifpdf switch
% Package: ifthen 2014/09/29 v1.1c Standard LaTeX ifthen package (DPC)
% Package: ifvtex 2016/05/16 v1.6 Detect VTeX and its facilities (HO)
% Package: ifxetex 2010/09/12 v0.6 Provides ifxetex conditional
% Package: infwarerr 2016/05/16 v1.4 Providing info/warning/error messages (HO)
% Package: intcalc 2016/05/16 v1.2 Expandable calculations with integers (HO)
% Package: keyval 2014/10/28 v1.15 key=value parser (DPC)
% Package: kvdefinekeys 2016/05/16 v1.4 Define keys (HO)
% Package: kvoptions 2016/05/16 v3.12 Key value format for package options (HO)
% Package: kvsetkeys 2016/05/16 v1.17 Key value parser (HO)
% Package: letltxmacro 2016/05/16 v1.5 Let assignment for LaTeX macros (HO)
% Package: ltxcmds 2016/05/16 v1.23 LaTeX kernel commands for general use (HO)
% Package: mathtools 2018/01/08 v1.21 mathematical typesetting tools
% Package: mhsetup 2017/03/31 v1.3 programming setup (MH)
% Package: nameref 2016/05/21 v2.44 Cross-referencing by name of section
% Package: parseargs 2014/04/21 v66
% Package: pdfescape 2016/05/16 v1.14 Implements pdfTeX's escape features (HO)
% Package: pdftexcmds 2018/09/10 v0.29 Utility functions of pdfTeX for LuaTeX (HO
% Package: pgf 2015/08/07 v3.0.1a (rcs-revision 1.15)
% Package: pgfcomp-version-0-65 2007/07/03 v3.0.1a (rcs-revision 1.7)
% Package: pgfcomp-version-1-18 2007/07/23 v3.0.1a (rcs-revision 1.1)
% Package: pgfcore 2010/04/11 v3.0.1a (rcs-revision 1.7)
% Package: pgffor 2013/12/13 v3.0.1a (rcs-revision 1.25)
% Package: pgfrcs 2015/08/07 v3.0.1a (rcs-revision 1.31)
% Package: pgfsys 2014/07/09 v3.0.1a (rcs-revision 1.48)
% Package: refcount 2016/05/16 v3.5 Data extraction from label references (HO)
% Package: rerunfilecheck 2016/05/16 v1.8 Rerun checks for auxiliary files (HO)
% Package: scalerel 2016/12/29 v1.8 Routines for constrained scaling and
%          stretching of objects, relative to a reference object or in
%          absolute terms
% Package: showlabels 2015/12/08 v1.7
% Package: stix2 2018/04/02 v2.0.0-latex STIX Two fonts support package
% Package: textcomp 2018/08/11 v2.0j Standard LaTeX package
% Package: thm-amsthm 2014/04/21 v66
% Package: thm-autoref 2014/04/21 v66
% Package: thm-kv 2014/04/21 v66
% Package: thm-listof 2014/04/21 v66
% Package: thm-patch 2014/04/21 v66
% Package: thm-restate 2014/04/21 v66
% Package: thmtools 2014/04/21 v66
% Package: tikz 2015/08/07 v3.0.1a (rcs-revision 1.151)
% Package: tikz-cd 2018/11/19 v0.9f Commutative diagrams with TikZ
% Package: trig 2016/01/03 v1.10 sin cos tan (DPC)
% Package: uniquecounter 2016/05/16 v1.3 Provide unlimited unique counter (HO)
% Package: url 2013/09/16  ver 3.4  Verb mode for urls, etc.
% Package: xcolor 2016/05/11 v2.12 LaTeX color extensions (UK)
% Package: xcolor-patch 2016/05/16 xcolor patch
% Package: xparse 2018-10-17 L3 Experimental document command parser
% Package: xstring 2018/12/09 v1.81 String manipulations (CT)
% Package: xy 2013/10/06 Xy-pic version 3.8.9

% Version 2
% 1.  Define \into
% 2.  Define m-chart at a point and subchart at a point
% 3,  Define (local) m-morphism
% 4.  Define local M-atlas morphism
% 5.  Change definition of M-atlas morphism
% 6.  Define classic M-atlas morphism
% 7.  Change definition of Ck-atlas morphism
% 8.  Define classic Ck-atlas morphism

\documentclass{article}
\usepackage{amsmath}
\usepackage{amsthm}
\usepackage{bm}
\usepackage{enumitem}
\usepackage{ifthen}
\usepackage{mathtools}
\usepackage{scalerel}
\usepackage{stix2}[notext,not1] %does not reqire XeTeX or luaTeX
\usepackage{thmtools}
\usepackage{tikz-cd}
\usepackage{xparse} % loads expl3
%See interface3.pdf
\usepackage{xstring}

\usepackage[cmtip,all,barr]{xy}
%Remove when xybarr.tex bug fixed
\newdir_{ (}{{ }*!/-.5em/@_{(}}

% Group action
\newcommand \action [1] []
   { \ifthenelse {\equal{#1}{}}
       {\mathbin \star}
       {\mathbin{\overset {#1} {\star}}}
   }

%Morphisms of category
\newcommand \Ar {\mathrm{Ar}}

\DeclareMathOperator \arin {\stackrel{\Ar}{\in}}

%Category of atlases from E to C
\newcommand \Atl [1] []
   { \ifthenelse {\equal{#1}{}}
     {\mathscr{A\!t\!l}}
     {\mathscr{A\!t\!l}^{\mathrm{#1}}}
   }

\DeclareMathOperator \Atlas {\mathrm{Atlas}}

\newcommand \Bun {\mathrm{Bun}}

\newcommand \BunProd {\mathrm{Bun}-\mathrm{prod}}

\newcommand \C {\mathrm{C}}

% Names of standard categories
\newcommand \cat [1] {\mathbf{#1}}

%Category of model space
\newcommand \Cat {\mathrm{Cat}}

% Select font for categories and sequences of categories
\newcommand \catname [1] {\mathscr{#1}}
%\newcommand \catseqname [1] {\mathbfscr{#1}}
\newcommand \catseqname [1] {\bm{\mathscr{#1}}}
% Too many fonts for \mathbfscr without XeTeX

\newcommand \Ck {\mathrm{C^k}}

\newcommand \Classic {\mathrm{Classic}}

\DeclareMathOperator \codomain {\mathrm{codomain}}

%\newcommand \compat [2] [] {\underset{\mathrm{{#1}compat}{#2}}

% Extended composition operators for function sequences
\newcommand {\compose} [1] [def]
   { \ifthenelse {\equal{#1}{def}}
       {\mathbin \circ}
       {\mathbin{\overset{#1} {\circ}}}
   }

% Compose head
\newcommand {\composeh} [1] [def]
   { \ifthenelse {\equal{#1}{def}}
       {\mathbin \odot}
       {\mathbin{\overset{#1} {\odot}}}
   }

% Compose tail
\newcommand {\composet} [1] [def]
   { \ifthenelse {\equal{#1}{def}}
       {\mathbin \cdot}
       {\mathbin{\overset{#1} {\cdot}}}
   }

\DeclareMathOperator \defeq {\stackrel{\mathrm{def}}{=}}

\DeclareMathOperator \domain {\mathrm{domain}}

\DeclareMathOperator \Domain {\seqname{domain}}

\newcommand \false {\mathrm{False}}

\newcommand \Fib   {\mathrm{Fib}}

%\newcommand \full [2] [] {\underset{\mathrm {{#1}full}}{#2}}

\newcommand \fullcref [1]
   { \ifthenelse {\equal{\nameref{#1}}{}}
       {\cref{#1}}
       {\cref{#1} (\nameref{#1})}
   }

% Select font for functions and sequences of functions
% \newcommand \funcname [1] {\mathit{#1}}
% \newcommand \funcseqname [1] {\bm{#1}}

\DeclareMathOperator \Functor {\mathop{\mathscr{F}}}

\newcommand \group {\mathrm{group}}

\DeclareMathOperator \head {\mathrm{head}}

\DeclareMathOperator \Hom {\mathrm{Hom}}

%Identity function
%\DeclareMathOperator* \Id {\mathrm{Id}}

%Sequence of identity functions
%\DeclareMathOperator* \ID {\mathbf{Id}}

\newcommand \into {\negthickspace\mon}

%Propositional function isCk_{f,A}(x,y,...)
\DeclareMathOperator \isCk {\mathrm{isCk}}

%Propositional function isAtl(A,E,C)
\newcommand \isAtl [1] []
   { \ifthenelse {\equal{#1}{}}
     {\mathrm{isAtl}}
     {\mathrm{isAtl}^{\mathrm{#1}}}
   }

%Propositional function isLCS(L, M, A, F, Sigma)
\DeclareMathOperator \isLCS {\mathrm{isLCS}}

\DeclareMathOperator \iso {\stackrel{\sim}{=}}

% Join two tuples
\DeclareMathOperator \join {\mathrm{join}}

%Category of M-Sigma local coordinate spaces
\newcommand \LCS {\mathrm{LCS}}

\DeclareMathOperator \length {\mathrm{length}}

\DeclareMathOperator \lengtho {\mathrm length0}

\newcommand \M {\mathrm{M}}

\newcommand \Man {\mathrm{Man}}

% Adjust : spacing for f maps a to b
\newcommand \maps {\!\colon}

%Set long names upright, short names italic
\newcommand{\mathvarname}[1]{%
  \begingroup\noexpandarg
  \StrLen{#1}[\temp]%
  \ifnum\temp>1
    \mathrm{#1}%
  \else
    #1%
  \fi
  \endgroup
}

\newcommand \minimal [2] [] {\underset{\mathrm{{#1}min}}{#2}}

%Model category
\newcommand \Mod {\mathrm{Mod}}

%Morphisms between two objects
\newcommand \Mor {\mathrm{Mor}}

% Morphism of category
\DeclareMathOperator \morphin {\stackrel{\Ar}{\in}}

%Object of category
\newcommand \Ob {\mathrm{Ob}}
\DeclareMathOperator \objin {\stackrel{\Ob}{\in}}

\newcommand \onto {\epi}

% All non-null open sets
\newcommand \op [2] [] {\underset{\mathrm {{#1}op}}{#2}}

\newcommand \optriv [2] [] {\underset{\mathrm {{#1}op-\mathscr{T\!r\!i\!v}}}{#2}}

\newcommand \pagecref [1]
   { \ifthenelse {\equal{\nameref{#1}}{}}
       {\cref{#1} on \cpageref{#1}}
       {\cref{#1} (\nameref{#1}) on \cpageref{#1}\negmedspace}
   }

\newcommand \Pagecref [1]
   { \ifthenelse {\equal{\nameref{#1}}{}}
       {\Cref{#1} on \cpageref{#1}}
       {\Cref{#1} (\nameref{#1}) on \cpageref{#1}}
   }

\DeclareMathOperator \range {\mathrm{range}}

\DeclareMathOperator \Range {\seqname{range}}

\newcommand \restrictto {\negthickspace\restriction}

\DeclareMathOperator \seqeq {\stackrel{()}{=}}

\DeclareMathOperator \seqin {\stackrel{()}{\in}}

%Font for sequence
%\newcommand \seqname [1] {\bm{\mathit{\mathsf{#1}}}}

%Underlying set of object in concrete category
\DeclareMathOperator* \Set {\mathbf{Set}}
%Use \cat{Set} for the standard category Set.
%
%\newcommand \sing [2] [] {\underset{\mathrm{{#1}Sing}{#2}}
%\newcommand \Sing [2] [] {\underset{\bm{\mathrm{{#1}Sing}}{#2}}

\newcommand \singcat [2] [] {\underset{{#1}\mathscr{S\!i\!n\!g}}{#2}}
\newcommand \Singcat [2] [] {\underset{\bm{{#1}\mathscr{S\!i\!n\!g}}}{#2}}
% Too many fonts to use \mathbfscr without XeTeX

%newcommand \strict [2] [] {\underset{\mathrm {{#1}strict}}{#2}}

%Subcategory
\newcommand \subcat [1] [] {\overset{\mathrm{{#1}cat}}{\subseteq}}

%Sequence of subcategories
\newcommand \SUBCAT [1] [] {\overset{\mathbf{{#1}cat}}{\subseteq}}

%Model subspace
\newcommand \submod [1] [] {\overset{\mathrm{{#1}mod}}{\subseteq}}

\DeclareMathOperator \SUBSETEQ {\stackrel{()}{\subseteq}}

%subspace of topological space
\newcommand \subsp [1] [] {\overset{\mathrm{{#1}subsp}}{\subseteq}}

\DeclareMathOperator \tail {\mathrm{tail}}

\newcommand \toiso {\,\to/{>}->>/^{\iso}}

%Topological space of model space
%Underlying topological space of object in concrete category
\newcommand \Top {\mathrm{Top}}
%Use \cat{Top} for the standard category Top.

%Topological category of model space
\newcommand \Topcat {\mathscr{T\!o\!p}}

% Font for topology
\newcommand \topname [1] {\mathfrak{#1}}

%Topology of topological space and relative topology of subset
\newcommand \Topology {\mathfrak{T\!o\!p}}

% \newcommand \triv [2] [] {\underset{{#1}\mathrm {triv}}{#2}}
% \newcommand \Triv [2] [] {\underset{{#1}\mathbf {triv}}{#2}}

\newcommand \trivcat [2] [] {\underset{\mathrm{#1}\mathscr{T\!r\!i\!v}}{#2}}
%\newcommand \Trivcat [2] [] {\underset{\mathrm{#1}\mathbfscr{T\!r\!i\!v}}{#2}}
\newcommand \Trivcat [2] [] {\underset{\mathrm{#1}\bm{\mathscr{T\!r\!i\!v}}}{#2}}
% Too many fonts to use \mathbfscr without XeTeX

\newcommand \true {\mathrm{True}}
\newcommand \truthcat {\mathscr{Truth}}
\newcommand \truthset {\mathbb{T}}
\newcommand \truthspace {\mathrm{Truthspace}}
\newcommand \truthtop {\mathfrak{Truthtop}}

\newcommand \unioncat [1] [] {\overset{\mathrm{{#1}cat}}{\cup}}
\newcommand \UNIONCAT [1] [] {\overset{\mathbf{{#1}cat}}{\cup}}

% Existential and universal quamtifiers
% Set former
% Union and intersection

% --------------------------------------------------------------------
% | From here to closing --- belongs in package                      |
% |                                                                  |

% Copyright 2016 Shmuel (Seymour J.) Metz
% I grant permission to the AMS, arXiv.org, the LATEX3 project and the
% TeX users group to incorporate these commands in any LaTeX package
% that may be freely redistributed, provided that they attribute the
% source.

% These commands are intended to allow semantic markup for some
% common mathermtical constructs:
%   \equant{variables}{proposition}     Existential qunatifier
%   \uquant{variables}{proposition}     Universal   quantifier
%   \union[indices]{set}                Union
%   \intersection[indices]{set}         Intersection
%   \set {elements}[propositions]       Set of
%   \set{elements}[propositions]*       Set of, split
%   \setupquant{optionstring}           Style of \equant, \uquant
%   \setupset{optionstring}             Style of \set, \union, \intersection
%   \full[qualifiers]{name}             Render name underset with some variation of full
%   \maxfull[qualifiers]{name}          Render name underset with some variation of max or full
%   \maximal[qualifiers]{name}          Render name underset with some variation of max
%   \strict[qualifiers]{name}           Render name underset with some variation of strict
%   \seqname{name}                      Render sequence/set/tuple name
%   \Triv[qualifiers]{name}             Render name underset with some variation of bold triv

% Because LaTeX has problems with parameters containg \\, the \set command
% has code to split the line between the elements and the proposition if
% the invocation is \set* and the environment is multline.

% Surround individual indices in \intersection, individual variables in
% quantifiers, individual propositions in \set and individual indices in
% \union with braces, e.g., \equant {{x \in X},{y \in Y}}{P(x,y)},
% \set{x}[{P(x)}, {Q(x)}], \union[{i \in I},{j \in J}]{O(i.j)}.

% Setup keywords for \equant, \uquant
% subscript  =none            Default
%                             \exists var1 ... \exists varn prop
%             stacked         \exists_\substack{var1 \\ ... varn} prop
%             multiple        \exists_{var1,...,varn} prop
% parentheses=none            Default
%             single          (\exists var1 ... \exists varn) prop
%             multiple        (\exists var1) ... (\exists varn) prop
% separator=                  Default {}

% Setup keywords for \intersection, \set, \union
% subscript  =none            \bigcap ix1, ..., ixn set
%             stacked         Default
%                             \bigcap__\substack{ix1 \\ ... ixn} set
%             multiple        \bigcap_{ix1,...,ixn} set
% separator=                  Default {\mid}
%                             {elements separator prop1 ^ ... propn}

\ExplSyntaxOn

\int_gzero_new:N \g_style_quant_parens_int
\int_gzero_new:N \g_style_quant_subscr_int
\int_gzero_new:N \g_style_set_subscr_int

\NewDocumentCommand{\equant}{mm}
  {
    \quant:nnn {\exists} {#1} {#2}
  }

\NewDocumentCommand{\uquant}{mm}
  {
    \quant:nnn {\forall} {#1} {#2}
  }

\NewDocumentCommand \setupquant {m}
  {
    \keys_set:nn {shmuel / quant} {#1}
  }

\keys_define:nn {shmuel / quant}
 {
  subscript            .choices:nn =
    {
      {
        none,
        stacked,
        multiple
      }
      {
        \int_gset:Nn \g_style_quant_subscr_int {\l_keys_choice_int-1}
      }
    },
  subscript            .default:n = multiple,
  subscript            .initial:n = none,
  parentheses          .choices:nn =
    {
      {
        none,
        single,
        multiple
      }
      {
        \int_gset:Nn \g_style_quant_parens_int {\l_keys_choice_int - 1}
      }
    },
  parentheses          .default:n = multiple,
  parentheses          .initial:n = none,
  separater            .tl_set:N = \g_style_quant_sep_tl,
  separater            .default:n = {.},
  separater            .initial:n = {}
 }

\cs_new:Npn \quant:nnn #1 #2 #3
  {
    %\int_show:N \g_style_quant_parens_int
    %\int_show:N \g_style_quant_subscr_int
    % g_style_quant_parens_int \ \int_use:N \g_style_quant_parens_int \
    % g_style_quant_subscr_int \ \int_use:N \g_style_quant_subscr_int \
    \clist_set:Nn \l_tmpa_clist {#2}
    \int_case:nn
      {\g_style_quant_subscr_int}
      {
        {0}
        {
          % No subscript
          % Set separater to ) ( quantifier or just quantifier
          \int_compare:nTF {\g_style_quant_parens_int = 2}
            {\tl_set:Nn \l_tmpa_tl {\right ) \left ( #1}}
            {\tl_set:Nn \l_tmpa_tl {#1}}
          \int_compare:nT {\g_style_quant_parens_int > 0} {\left (}
          #1
          \clist_use:Nn \l_tmpa_clist {\l_tmpa_tl}
          \int_compare:nT {\g_style_quant_parens_int > 0} {\right )}
          \g_style_quant_sep_tl #3
        }
        {1}
        {
          % Stacked subscript on single quantifier
          \fp_set:Nn
            \l_tmpa_fp
            {ceil{\clist_count:N{\l_tmpa_clist} - 1} * .1 + 1}
          \int_compare:nT {\g_style_quant_parens_int > 0} {\left (}
          \scaleobj{\fp_to_decimal:N \l_tmpa_fp}{#1} \sb
             { \substack { \clist_use:Nn \l_tmpa_clist { \\ } } }
          \int_compare:nT {\g_style_quant_parens_int > 0} {\right )}
          #3
        }
        {2}
        {
          % Subscripts on separate quantifiers
          \clist_map_inline:Nn
            \l_tmpa_clist
            {
              % (quantifier \sb predicate) or quantifier \sb predicate
              {
                \int_compare:nT
                  {\g_style_quant_parens_int > 0}
                  {\left (}
                \scaleobj{1.2}{#1} \sb
                {##1}
                \int_compare:nT
                  {\g_style_quant_parens_int > 0}
                  {\right )}
              }
            }
          #3
        }
      }
  }

\NewDocumentCommand{\set}{mos}
  {
    \IfBooleanTF {#3}
    {
%      \msg_term:n {set with star}
%      \msg_term:n{{P1 #1}}
%      \msg_term:n{{P2 #2}}
      \bool_set_true:N \l_tmpa_bool
%      \bool_show:N \l_tmpa_bool
    }
    {
      \bool_set_false:N \l_tmpa_bool
    }
    \set_of:nnn {#1} {#2} {\l_tmpa_bool}
  }

%\tl_new:N \g_style_set_sep_tl
%\tl_gset:Nn \g_style_set_sep_tl {\mid}

\NewDocumentCommand \setupset {m}
  {
    \keys_set:nn {shmuel / set} {#1}
  }

\keys_define:nn {shmuel / set}
  {
    separater .tl_set:N = \g_style_set_sep_tl,
    subscript            .choices:nn =
      {
        {
          stacked,
          multiple
        }
        {
          \int_gset:Nn \g_style_set_subscr_int {\l_keys_choice_int-1}
        }
      },
    separater .initial:n = {\mid},
    subscript .initial:n = stacked
  }

\cs_new:Npn \set_of:nnn #1 #2 #3
  {
    \IfValueTF {#2}
    {
%      \msg_term:n {set_of:nn \ has \ predicates \ #2}
    }
    {
%      \msg_term:n {set_of:nn \ has \ no \ predicates}
    }
    \clist_set:Nn \l_tmpa_clist {#2}
%    \msg_term:n {l_tmpa_clist \ set}
    \tl_gset:Nn \g_tmpa_tl {\clist_use:Nn \l_tmpa_clist {\land}}
%    \msg_term:n {g_tmpa_tl \ set \ to \ \g_tmpa_tl}
    \IfValueTF {#2}
    {
      \bool_if:nTF {#3}
      {
        \scalerel*{\{}{#1\g_tmpa_tl}
        #1
%        \msg_term:n  {scalerel returns \scalerel{\g_style_set_sep_tl}{\g_tmpa_tl}}
        \clist_set:Nn \l_tmpa_clist {#2}
        \scalerel*{\mathbin{\g_style_set_sep_tl}}{#1\g_tmpa_tl}
        \\
        \clist_set:Nn \l_tmpa_clist {#2}
        \g_tmpa_tl
%         \tl_show:N \g_tmpa_tl
        \clist_set:Nn \l_tmpa_clist {#2}
        \scalerel*{\}}{#1\g_tmpa_tl}
      }
      {
        \left \{
        #1
        \clist_set:Nn \l_tmpa_clist {#2}
        \scalerel*{\mathbin{\g_style_set_sep_tl}}{#1\g_tmpa_tl}
        \g_tmpa_tl
        \right \}
      }
    }
    {
      \left \{ #1 \right \}
    }
  }

% Render underset with some variation of full
\NewDocumentCommand{\full}{O{full} m}
  {
    \show_full:nn {#1} {#2}
  }

\cs_new:Npn \show_full:nn #1 #2
  {
    \tl_set:Nn \l_tmpa_tl {#1}
    \str_if_eq:eeTF {\str_item:nn {#1} {-1}} {-}
      {
        \tl_set:Nx \l_tmpa_tl {{#1}max}
      }
      {
        \str_case:nn {#1}
        {
          {*}  {\tl_set:Nn \l_tmpa_tl {(full,S-max-full,max-full)}}
          {**} {\tl_set:Nn \l_tmpa_tl {full~(S-max-full,max-full)}}
        }
      }
    \underset {\textup{\l_tmpa_tl}} {#2}
  }

\NewDocumentCommand{\funcname}{m}
  {
    \funcname:n {#1}
  }

\cs_new:Npn \funcname:n #1
  {
    %code here \tl_count:n
    \int_compare:nTF {\tl_count:n{#1} > 1}
      {
        {\mathrm{#1}}
      }
      {
        {\mathit{#1}}
      }
  }

\NewDocumentCommand{\funcseqname}{m}
  {
    \funcseqname:n {#1}
  }

\cs_new:Npn \funcseqname:n #1
  {
    %code here \tl_count:n
    \int_compare:nTF {\tl_count:n{#1} > 1}
      {
        {\bm{\mathrm{#1}}}
      }
      {
        {\bm {\mathit{#1}}}
      }
  }

% Identity or inclusion function
\NewDocumentCommand \Id {O{}}
  {
     \Id:nnn {Id} {\operatorname{Id}} {#1}
  }

% Sequence of Identity or Inclusion functions
\NewDocumentCommand \ID {O{}}
  {
     \Id:nnn {ID} {\operatorname{\mathbf{Id}}} {#1}
  }

\cs_new:Npn \Id:nnn #1 #2 #3
  {
    \clist_set:Nn \l_tmpa_clist {#3}
    \int_case:nnF
      {\clist_count:N \l_tmpa_clist}
      {
        {0}
        {
         #2
        }
        {1}
        {
         #2 \sb {#3}
        }
        {2}
        {
         #2
           \sp {\clist_item:Nn \l_tmpa_clist 1}
           \sb {\clist_item:Nn \l_tmpa_clist 2}
        }
      }
      {
         \msg_error:nnnn {shmuel} {toomany} {two} {#1: [#3]}
      }
  }

    \msg_new:nnn
      {shmuel}
      {toomany}
      {More \ than \ #1  \ items \ for \  \c_backslash_str #2}

% Render underset with some variation of max-full
\NewDocumentCommand{\maxfull}{O{max-full} m}
  {
    \maximal:nnnnnnn {max-full} {#1} {#2}
                     {full,S-max,max,S-max-full,max-full}
                     {full,S-max,max,S-max-full}
                     {full~(S-max,max,S-max-full,max-full)}
                     {full~(S-max,max,S-max-full)}
  }

% Render underset with some variation of max
\NewDocumentCommand{\maximal}{O{max} m}
  {
    \maximal:nnnnnnn {max} {#1} {#2}
                     {S-max,max}
                     {S-max}
                     {S-max~(max)}
                     {S-max)}
  }

\cs_new:Npn \maximal:nnnnnnn #1 #2 #3 #4 #5 #6 #7
  {
    \tl_set:Nn \l_tmpa_tl {#2}
    \str_if_eq:eeTF {\str_item:nn {#2} {-1}} {-}
      {
        \tl_set:Nx \l_tmpa_tl {{#2}#1}
      }
      {
        \str_case:nn {#2}
        {
          {*}   {\tl_set:Nn \l_tmpa_tl {(#4)}}
          {*x}  {\tl_set:Nn \l_tmpa_tl {(#5)}}
          {**}  {\tl_set:Nn \l_tmpa_tl {#6}}
          {**x} {\tl_set:Nn \l_tmpa_tl {#7}}
        }
      }
    \underset {\textup{\l_tmpa_tl}} {#3}
  }

\NewDocumentCommand{\donotbreak} {m}
    {
      \mbox{#1}\discretionary{}{}{}
    }

\NewDocumentCommand{\nobreaks} {>{\SplitList{,}}m}
    {
      \ProcessList{#1}{\donotbreak}
    }

% Render cref (nameref) on cpageref
\NewDocumentCommand \pagecrefn {m}
    {
      \pagecref:nn {\cref} {#1}
    }

\cs_new:Npn \pagecref:nn #1 #2
    {
      \bool_if:nTF
        {
          \str_if_eq_p:nn {\str_range:nnn {#2} {1} {3}} {eq:} ||
          \str_if_eq_p:nn {\nameref{#2}}                {}
        }
        {
          #1{#2} on \cpageref{#2}
        }
        {
          #1{#2} (\nameref{#2}) on \cpageref{#2}
        }
    }

% Render sequence, set or tuple name
\NewDocumentCommand{\seqname}{m}
  {
    \seqname:n {#1}
  }

\cs_new:Npn \seqname:n #1
  {
    %code here \tl_count:n
    \int_compare:nTF {\tl_count:n{#1} > 1}
      {
        {\mathbf{#1}}
      }
      {
        {\mathbfit{#1}}
      }
  }

% Render underset with some variation of strict
\NewDocumentCommand{\strict}{O{strict} m}
  {
    \show_strict:nn {#1} {#2}
  }

\cs_new:Npn \show_strict:nn #1 #2
  {
    \tl_set:Nn \l_tmpa_tl {#1}
    \str_if_eq:eeTF {\str_item:nn {#1} {-1}} {-}
      {
        \tl_set:Nx \l_tmpa_tl {{#1}strict}
      }
      {
        \str_case:nn {#1}
        {
          {*}  {\tl_set:Nn \l_tmpa_tl {(semi-strict,strict)}}
          {**} {\tl_set:Nn \l_tmpa_tl {semi-strict~(strict)}}
        }
      }
    \underset {\textup{\l_tmpa_tl}} {#2}
  }

% Render underset with some variation of bold triv
\NewDocumentCommand{\Triv}{o m}
  {
    \Triv:nn {#1} {#2}
  }

\cs_new:Npn \Triv:nn #1 #2
  {
    \IfValueTF
      {#1}
      {
        \str_case:nnF {#1}
        {
           {full-submod-} {\underset{\mathbf{{#1}triv}}{#2}}
           {full-}        {\underset{\mathbf{{#1}triv}}{#2}}
           {submod-}      {\underset{\mathbf{{#1}triv}}{#2}}
        }
        {
          {\underset{#1\mathbf{triv}}{#2}}
        }
      }
      {
         {\underset{\mathbf{triv}}{#2}}
      }
  }

\NewDocumentCommand{\intersection}{om}
  {
    \unint_of:nnn \bigcap {#1} {#2}
  }

\NewDocumentCommand{\union}{om}
  {
    \unint_of:nnn \bigcup {#1} {#2}
  }

\cs_new:Npn \unint_of:nnn #1 #2 #3
  {
    \IfValueTF {#2}
    {
%      \int_show:N \g_style_set_subscr_int
      \clist_set:Nn \l_tmpa_clist {#2}
      \int_case:nn
        {\g_style_set_subscr_int}
        {
          {0}
          {
            % Stacked subscript
%            \msg_term:n {\clist_count:N{\l_tmpa_clist} \ tokens \ stacked}
%            \clist_show:N \l_tmpa_clist
            \fp_set:Nn
              \l_tmpa_fp
              {ceil{\clist_count:N{\l_tmpa_clist} - 1} * .1 + 1}
            \scaleobj{\fp_to_decimal:N \l_tmpa_fp}{#1} \sb
            {\substack { \clist_use:Nn \l_tmpa_clist { \\ } }}
            #3
          }
          {1}
          {
            % Subscripts comma separated
%            \msg_term:n {\clist_count:N{\l_tmpa_clist} \ tokens \ comma \ separated}
%            \clist_show:N \l_tmpa_clist
            #1 \sb
            {\clist_use:Nn \l_tmpa_clist {,}}
            #3
          }
        }
    }
    {
      % No subscript
%      \msg_term:n {No~subscript}
%      \clist_show:N \l_tmpa_clist
      #1 #3
    }
  }

\ExplSyntaxOff

% |                                                                  |
% | From opening to here --- belongs in package                      |
% --------------------------------------------------------------------



\theoremstyle{definition}

\theoremstyle{remark}

\numberwithin{equation}{section}

\newcommand \Alpha   A
\newcommand \Beta    B
\newcommand \Epsilon E
\newcommand \Zeta    Z
\newcommand \Eta     H
\newcommand \Iota    I
\newcommand \Kappa   K
\newcommand \Mu      M
\newcommand \Nu      N
\newcommand \Omicron O
\newcommand \Rho     P
\newcommand \Tau     T
\newcommand \Chi     S

\usepackage[colorlinks,hidelinks,draft=false]{hyperref}

\usepackage{cleveref}

% \newtheorem must follow cleveref
\newtheorem{theorem}{Theorem}[section]
\newtheorem{corollary}[theorem]{Corollary}
\def\corollaryautorefname{Corollary}         % Needed for \autoref
\newtheorem{definition}[theorem]{Definition}
\def\definitionautorefname{Definition}       % Needed for \autoref
\newtheorem{example}[theorem]{Example}
\newtheorem{lemma}[theorem]{Lemma}
\def\lemmaautorefname{Lemma}                 % Needed for \autoref
\newtheorem{remark}[theorem]{Remark}
\newtheorem{xca}[theorem]{Exercise}

\hypersetup {
   bookmarksnumbered=true,
   colorlinks,
   pdfinfo={
      Author={Shmuel (Seymour J.) Metz},
      Keywords={fiber bundles,manifolds},
      Subject={Topology},
      Title={Towards a taxonomy of atlases and of morphisms between them}
   }
}
\usepackage[draft]{showlabels}
% \usepackage[final]{showlabels}
\showlabels{cite}
\showlabels{cref}
\showlabels{crefrange}

%\DeclareMathAlphabet{\mathbdit}{T1}{\itdefault}{\mddefault}{\sldefault}
%\SetMathAlphabet{\mathbdit}{bold}{T1}{\itdefault}{\bfdefault}{\sldefault}
%
%\DeclareMathAlphabet{\mathsfbd}{T1}{\sfdefault}{\mddefault}{\sldefault}
%\SetMathAlphabet{\mathsfbd}{bold}{T1}{\sfdefault}{\bfdefault}{\sldefault}
%
%\DeclareMathAlphabet{\mathsfit}{T1}{\sfdefault}{}{\sldefault}
%\SetMathAlphabet{\mathsfit}{normal}{T1}{\sfdefault}{}{\sldefault}
%
%\DeclareMathAlphabet{\mathsfbdit}{T1}{\sfdefault}{\mddefault}{\sldefault}
%\SetMathAlphabet{\mathsfbdit}{bold}{T1}{\sfdefault}{\bfdefault}{\sldefault}

\begin{document}
\hyphenation{near pres-en-ta-tion}

\title%
{%
  Towards a taxonomy of atlases and of morphisms between them\thanks%
  {
  I wish to gratefully thank Walter Hoffman (z"l),
  Milton Parnes, Dr. Stanley H. Levy (z"l), the Mathematics department of
  Wayne State University, the Mathematics department of the State
  University of New York at Buffalo and others who guided my education.
  }
}
\author{Shmuel (Seymour J.) Metz
}
\maketitle

% \address{4963 Oriskany Drive\\Annandale, VA  22003-5141}
% \email{smetz3@gmu.edu}
% \urladdr{http://mason.gmu.edu/~smetz3}

% \subjclass[2010]{Primary 18F15; Secondary 55R65,57N99,58A05}
% 18F15  Abstract manifolds and fiber bundles
% 32Qxx  Complex manifolds
% 53-    Differential geometry
% 53Axx  Classical differential geometry
% 53A99  None of the above, but in this section
% 55Rxx  Fiber spaces and bundles
% 55R10  Fiber bundles
% 55R65  Generalizations of fiber spaces and bundles
% 57Nxx  Topological manifolds
% 57N99  None of the above, but in this section
% 57Rxx  Differential topology
% 58Axx  General theory of differentiable manifolds
% 58A05  Differentiable manifolds, foundations
% 58Bxx  Infinite-dimensional manifolds

% Primary 18F15  Abstract manifolds and fiber bundles
% Secondary 55R65,57N99,58A05
% \keywords{fiber bundles,manifolds}

\begin{abstract}
Manifolds and fiber bundles, while superficially different, have strong
parallels; in particular, they are both defined in terms of equivalence
classes of atlases or in terms of maximal atlases, with the atlases
treated as mere adjuncts.  This paper presents a unified view of atlases
for manifolds and fiber bundles as mathematical entities in their own
right. It defines some convenient notation, defines categories of, e.g.,
atlases, and constructs functors among them.

\cite{LCS} introduced the ideas presented here, but many of the details
are not needed there.  This paper fleshes out the concepts in more
detail than would be relevant there.
\end{abstract}

\setupquant {parentheses=multiple,subscript=stacked}
\setupset   {}

\tikzset{
    rot90/.style={anchor=south, rotate=90, inner sep=.2mm}
}

\part {Introduction}
\label{part:intro}

Historically, manifolds were formalized in terms of atlases based on
Euclidean spaces and fiber bundles were formalized in terms of atlases
based on product spaces, using either equivalence classes of atlases or
maximal atlases.  The concept of pseudo-groups allowed unifying
manifolds and manifolds with boundary.  The definitions of fiber bundles
and manifolds have strong parallels, and can be unified in a similar
fashion.  The standard presentations treat these atlases as secondary
adjuncts.

\cite{LCS} introduced the concept of an m-atlas and of morphisms among
them only as preliminaries to unifying manfolds and fiber bundles as
special cases of local coordiinate spaces.  This paper treats m-atlases
as objects of interest in their own right, and includes concepts not
relevant to local coordinate spaces.

This paper defines an approach using categories and
commutative diagrams that is designed for easy exposition at the possible
expense of abstractness and generality. In particular, I have chosen to
simplify proofs by stating stricter assumptions than necessary and have
assumed the Axiom of Choice (AOC).

Although this paper incidentally defines partial equivalents to
manifolds and fiber bundles using model spaces and model atlases, it
does not explicitly reflect the role of the group in fiber bundles.

This paper defines functors among categories of atlases, categories of
model spaces, categories of manifolds and categories of fiber bundles;
it constructs more machinery than is customary in order to facilitate
the presentation of those categories and functors.

\begin{remark}
The definitions of morphisms given here are nonstandard, but there are
also definitions of classical morphisms that are much closer to the
standard ones.
\end{remark}

\Crefrange{part:conv}{part:modeldef} present nomenclature and give basic
results.

\Cref{part:M-charts} defines m-atlases, m-atlas morphisms and categories
of them;
{
  \showlabelsinline
  \cref{lem:M-ATLiscat}
}
proves that the defined categories are indeed categories.

\pagecref{part:man} and \pagecref{part:bun} show the equivalence of
manifolds and fiber bundles with m-atlases by explicitly exhibiting
functors to and from the m-atlases.

\begin{remark}
The unconventional definitions of manifold and fiber bundle are intended
to make their relationship to local coordinate spaces more natural.
See \cite{LCS} for details.
\end{remark}

Most of the lemmata, theorems and corollaries in this paper should be
substantially identical to results that are familiar to the reader. What
is novel is the perspective and the material directly related to local
coordinate spaces. The presentation assumes only a basic knowledge of
Category Theory, such as may be found in the first chapter of
\cite{CftWM} or
{
  \showlabelsinline
  \cite{JoyCat}.
}

Many of the proofs are obvious and may be omitted from other versions of
this paper.

\section{New concepts and notation}
\label{sec:new}
This paper introduces a significant number of new concepts and some
modifications of the definitions for some conventional concepts. It also
introduces some notation of lesser importance.  The following are the
most important.

\begin{enumerate}
\item \hyperref[def:NCD]{Nearly commutative diagram (NCD)}, NCD at a point,
locally NCD and special cases with related nomenclature.
These are cases where a diagram can be modified to make it commutative.

\item\hyperref[def:model]{Model space} and related concepts.
A \hyperref[def:model]{Model space}\footnote{
  The phrase has been used before, but with a different meaning.
}
is a topological space with a category of permissible open sets and
mappings satisfying specified conditions.  It is similar to a
pseudo-group, with some important differences.

\item \hyperref[def:ModTop]{Model topology} and
\hyperref[def:M-para]{M-paracompactness}

\item
An \hyperref[def:m-atlas]{m-atlas} is a generalization of atlases in
manifolds.

\item
An \hyperref[def:M-ATLmorphnear]{m-atlas near morphism}
(\hyperref[def:M-ATLmorph]{m-atlas morphism}) is a generalization of a
$\Ck$ map between differentiable manifolds. It differs in using both a
function between the manifolds and a function between the coordinate
spaces.

\item
\hyperref[def:lin]{Linear space and related concepts}
A linear space is a subspace of a Banach space or of a Fr\'echet space
that is suitable for defining coordinate patches in a manifold.

\item
\hyperref[def:trivck]{Trivial $\Ck$ linear model space} and related
concepts.
A trivial $\Ck$ linear model space is a maximal model space
with morphisms $\Ck$ maps between open subsets of a linear space.

\item \hyperref[def:BunAtl]{$G$-$\rho$ bundle atlases}%
\footnote{Similar to coordinate bundles}
and related concepts
\end{enumerate}

\part {Conventions}
\label{part:conv}
An arrow with an Equal-Tilde ($A \toiso_\phi B$) represents an
isomorphism. One with a hook ($A \underset{i}{\hookrightarrow} B$)
represents an inclusion map. One with a tail ($A \into_i B$) represents
a monomorphism. One with a double head ($A \onto_\pi B$) represents a
surjection.

When a superscript ends in $-1$, e.g., $\funcname{f}^{i-1}$, it is to be
taken as function inverse rather than subtraction of 1.

All diagrams shown are commutative; none are exact.

A definition of a base term and several related terms may specify
modifiers in parenthese in the base definition and then give the
definitions for each modifier; usually the modifiers will be
restrictions and the base definition will not be repeated.  E.g.,
{\textquotedblleft}$\funcseqname{f}$ is a (semi-strict, strict)
prestructure morphism of $\seqname{P}^1$ to $\seqname{P}^2$
iff{\textquotedblright} is followed by the definition of prestructure
morphism, then by the restrictions for strict and semi-strict
prestructure morphisms.

When a definition defines a base propositional function and variant
propositional functions with a qualifier given as a superscript or
underset, then the form $\mathrm{base}^{(qualifiers)}$ will refer to
either $\mathrm{base}^{qualifier}$ or $\mathrm{base}$ and the form
$\underset{(qualifiers)}{base}$ will refer to either
$\underset{qualifier}{base}$ or $base$, e.g.,
$\strict[*]{\isAtl[(classic)]_\Ar}$ refers to either $\isAtl_\Ar$,
$\isAtl[classic]_\Ar$, $\strict[semi-]{\isAtl_\Ar}$,
$\strict[semi-]{\isAtl[classic]_\Ar}$, $\strict{\isAtl_\Ar}$ or
$\strict{\isAtl[classic]_\Ar}$.

The form $\mathrm{base}^{qualifier(,qualifiers)}$ will refer to either
$\mathrm{base}^{\mathrm{qualifier}}$ or to the variant with one or more
of the additional qualifiers in the superscript, e.g.,
$\strict{\isAtl[classic(,near)]_\Ar}$ refers to either
$\strict{\isAtl[classic]_\Ar}$ or
$\strict{\isAtl[classic,near]_\Ar}$.

Alternatively, a definition may specify a numbered list of alternatives,
and subsequently specify additional numbered lists with items
corresponding to those in the first list, e.g.,
$\funcseqname{f}$ is also a
\begin{enumerate}
\item full
\item semi-maximal
\item maximal
\item full semi-maximal
\item full maximal
\end{enumerate}
$E^1$-$E^2$ $\Ck$ morphism of
$\seqname{A}^1$ to $\seqname{A}^2$ in the coordinate spaces $C^1$,
$C^2$, abbreviated as
\begin{enumerate}
\item abbreviation for full
\item abbreviation for semi-maximal
\item abbreviation for maximal
\item abbreviation for full semi-maximal
\item abbreviation for full maximal
\end{enumerate}
iff
\begin{enumerate}
\item definition for full
\item definition for semi-maximal
\item definition for maximal
\item definition for full semi-maximal
\item definition for full maximal
\end{enumerate}

A Corollary, Lemma or Theorem that applies to related terms defined
with the above convention will specify the modifiers in parentheses
to indicate that it applies to the base term and to the variant
terms, e.g., {\textquotedblleft}a (semi-strict, strict)
$\seqname{E}^i$-$\seqname{E}^{i+1}$ m-atlas morphism{\textquotedblright}
If it applies only to more restrictive terms then it will specify the
first relevant modifier followed by the remaining relevant modifiers in
parentheses, e.g., {\textquotedblleft}If
$\catname{S}^i \SUBCAT[full-] \catname{S}'^i$ and $\funcseqname{f}^i$ is
a semi-strict (strict) prestructure morphism{\textquotedblright}.

Blackboard bold upper case will denote specific sets, e.g., the
Naturals.

Bold lower case italic letters will refer to sets, sequences and tuples
of functions, e.g.,
$\funcseqname{f} \defeq (\funcname{f}_1, \funcname{f}_2)$.

Bold lower case Latin letters will refer to sequence valued functions of
sequences and tuple valued functions of tuples, e.g., $\Range$ yields
the sequence of ranges of a sequence of functions.

Bold upper case italic letters will refer to sequences or tuples, e.g.,
$\seqname{A}=(x,y,z)$, to sets of them, to sets of topological spaces or
to sets of open sets.

Bold upper case script letters will refer to
sequences of categories, e.g., \\*
\nobreaks
{{
  $\catseqname{A} \defeq (\catname{A}_\alpha, \alpha \in \Alpha)$.
}}

Fraktur will refer to topologies and to topology-valued functions, e.g.,
$\Topology$.

Functions have a range, domain and relation, not just a relation. Unless
otherwise stated, they are assumed or proven to be continuous.

Groups are assumed to be topological groups. The ambiguous notation
$x^{-1}$ will be used when it is obvious from context what the group
operation $\star$ and the group identity $\mathbf{1}_G$ are.

Lower case Greek letters other than $\pi$, $\rho$, $\sigma$, $\phi$ and
$\psi$ will refer to ordinals, possibly transfinite, and to formal
labels.  A letter with a Greek superscript and a letter with a Latin or
numeric superscript always refer to distinct variables.

Lower case $\pi$ will refer to a projection operator

Lower case $\rho$ will refer to a continuous effective group action,
i.e., a continuous representation of a group in a homeomorphism group.

Lower case $\sigma$ will refer to a sequence of ordinals, referred to as
a signature.

Lower case $\phi$ will refer to a coordinate function.

Lower case italic and Latin letters will refer to
\begin{enumerate}
\item elements of a set or sequence
\item formal labels \\*
When the letters $u$ and $v$, with or without a superscript, refer to an
element of a set then they may also be used as formal labels for sets
associated with that element.
\item functions
\item morphisms of a category
\item natural numbers
\item objects of a category
\item ordinal numbers
\end{enumerate}

Upper case Greek letters other than $\Sigma$ may refer to
\begin{enumerate}
\item ordinal used as the limit of a sequence of consecutive ordinals,
e.g., $x_\alpha, \alpha \preceq \Alpha$
\item ordinal used as the order type of a sequence of consecutive
ordinals, e.g., $x_\alpha, \alpha \prec \Alpha$
\end{enumerate}

Upper case $\Sigma$ will refer to a sequence of signatures

Upper case Latin letters will refer to
\begin{enumerate}
\item Natural numbers
\item Topological spaces
\item Open sets
\item Formal labels for elements of a sequence or tuple of functions,
e.g., $\funcname{f}_E$ might be $\funcname{f}_0 \maps E_1 \to E_2$.
\end{enumerate}

Upper case Script Latin letters will refer to categories and functors.

Upright Latin letters will be used for long names.
In particular,
\begin{enumerate}
\item
$\Ar$, $\Hom$, $\Mor$ and $\Ob$ are the conventional set-valued
functions on small categories:

\begin{description}
\item[$\Ar(\catname{C})$]
All morphisms of $\catname{C}$
\item[$\Hom_\catname{C}(a,b)$]
All morphisms of $\catname{C}$ from $a$ to $b$
\item[$\Mor(\catname{C})$]
All morphisms of $\catname{C}$.
Synonymous with $\Ar(\catname{C})$.
\item[$\Ob(\catname{C})$]
All objects of $\catname{C}$
\end{description}

They will never be applied to large categories in this paper.

\item
$\cat{Set}$ and $\cat{Top}$ are the conventional categories of sets and
topological spaces.

\item
The letters E, G, X and Y will be used for formal labels.
\end{enumerate}

The term $\Ck$ includes $\mathrm{C}^\infty$ (smooth) and
$\mathrm{C}^\omega$ (analytic).

This paper uses the term morphism in preference to arrow, but uses
the conventional $\Ar$.

This paper uses the terms empty, null and void interchangeably.

The term sequence without an explicit reference to $\mathbb{N}$
will refer to a general ordinal sequence, possibly transfinite.

Sequence numbering, unlike tuple numbering, starts at 0 and the
exposition assumes a von Neumann definition of ordinals, so that
$\alpha \in \beta \equiv \alpha \prec \beta$.

Except where explicitly stated otherwise, all categories mentioned are
small categories with underlying sets, but the morphisms will often not
be set functions between the objects and there will not always be a
forgetful functor to $\cat{Set}$ or $\cat{Top}$. By abuse of language no
distinction will consistently be made between a category $\catname{A}$
of topological spaces and the concrete category
$(\catname{A},\catname{U})$ over $\cat{Top}$.  Similarly, no distinction
will consistently be made among the object $U \in \Ob(\catname{A})$, the
topological space $\catname{U}(U)$ and the underlying set.

The notation $G^V$ will refer only to the set of continuous functions
from $V$ to $G$, never to the set of all functions from $V$ to $G$.

When defining a category, the ordered pair $(\seqname{O}, \seqname{M})$
refers to the smallest concrete category over $\cat{Set}$ or $\cat{Top}$
whose objects are the elements in $\seqname{O}$, whose morphisms include
all of the elements of $\seqname{M}$ and whose morphisms from
$o^1 \in \seqname{O}$ to $o^2 \in \seqname{O}$ are functions
$\funcname{f} \maps o^1 \to o^2$ and whose composition is function
composition.

When defining a category, the ordered triple
$(\seqname{O}, \seqname{M}, C)$ refers to the small category whose
objects are in $\seqname{O}$, whose morphisms are in $\seqname{M}$,
whose $\Hom$ is

\begin{equation}
\Hom_{(\seqname{O}, \seqname{M}, C)}
  (o_1 \in \seqname{O}, o_2 \in \seqname{O})
\defeq
\set
  {
    (
      \funcseqname{f},
      o_1,
      o_2
    )
    \in \seqname{M}
  }
\end{equation}
and whose composition is C.

By abuse of language I may write
``$\catname{S}$'' for $\Ob(\catname{S})$,
``$A \in \catname{A}$'' for  $A \in\ \Ob(\catname{A})$,
``$A \subset \catname{A}$'' for $A \subset \Ob(\catname{A})$,
``$A \in \catname{A} \subset B \in \catname{B}$'' for
``the underlying set of $A$ is contained in the underlying set of $B$
and the inclusion $\funcname{i} \maps x \in A \hookrightarrow x \in B$
is a morphism'' and
``$\funcname{f} \maps A \to B$'' for
$\funcname{f} \in \Hom_{\catname{C}}(A,B)$, where $\catname{C}$ is
understood by context.

By abuse of language I shall use the same nomenclature for sequences and
tuples, and will use parentheses around a single expression both for
grouping and for a tuple with a single element; the intent should be
clear from context.

By abuse of language I will omit parenthese around the operands of
Functors when they can be assumed by context.

By abuse of language I shall use the $\times$ and $\bigtimes$ symbols
for both Cartesian products of sets and Cartesian products of
functions on those sets.

By abuse of language, and assuming AOC,  I shall refer to some sets as
ordinal sequences, e.g., ``$(C_\alpha, \alpha \in \Alpha)$'' for
``$\{C_\alpha \mid \alpha \in \Alpha\}$'', in cases where the order is
irrelevant.

Where a proof requires proving intermediate propositions, a bold leadin
to a paragraph will indicate such a proposition, e.g.

\begin{description}
\item[Intermediate proposition 1:] Intermediate proof 1.

\item[Intermediate proposition 2:]
\leavevmode
Intermediate proof 2.
\end{description}

By abuse of language, I may omit universal quantifiers in cases where
the intent is clear.

In some cases I define a notion similar to a conventional notion and
also need to refer to the conventional notion. In those cases I prefix
a letter or phrase to the term, e.g., m-paracompact versus paracompact.

\part {General notions}
\label{part:notions}
This part of the paper describes nomenclature used throughout the paper.
In some cases this reflects new nomenclature or notions, in others it
simply makes a choice from among the various conventions in the
literature.

\begin{definition}[Operations on categories]
\label{def:catprop}
If $\catname{C}$ is a category then
\begin{enumerate}
\item
$x \objin \catname{C} \equiv x \in \Ob(\catname{C})$

\item
$x \arin \catname{C} \equiv x \in \Ar(\catname{C})$

\item
If $\catname{C}$ is a concrete category over $\cat{Set}$ or $\cat{Top}$
and $o \objin \catname{C}$, then $\Set_\catname{C}(o)$ is the underlying
set of $o$. By abuse of language we will write $\Set(o)$ when the
category is understood from context.

\item
If $\catname{C}$ is a concrete category over $\cat{Top}$ and
$o \objin \catname{C}$, then $\Top_\catname{C}(o)$ is the underlying
topological space of $o$. By abuse of language we will write $\Top(o)$
when the category is understood from context.
\end{enumerate}

If $\catname{S}$ and $\catname{T}$ are categories then
$\catname{S} \subcat \catname{T}$ iff $S$ is a subcategory of
$\catname{T}$ and $\catname{S} \subcat[full-] \catname{T}$ iff $S$ is a
full subcategory of $\catname{T}$.

If $\catname{S}$ and $\catname{T}$ are categories then the category
union of $\catname{S}$ and $\catname{T}$, abbreviated \\*
$\catname{S} \unioncat \catname{T}$, is the category whose objects are
in $\catname{S}$ or in $\catname{T}$ and whose morphisms are in
$\catname{S}$ or in $\catname{T}$.
\end{definition}

\begin{definition}[Identity]
\label{def:Id}
$\Id[S]$ is the identity function on the space $S$,

$\Id[o]$ is the
identity morphism for the object $o$\footnote{
The object is often expressed as a tuple, e.g.,
$\Id[{{(\seqname{A}, \seqname{B})}}]$ is the identity morphism for the
object $(\seqname{A}, \seqname{B})$
},

$\Id[U,V]$, for $U \subseteq V$, is the inclusion map\footnote{
$U$ need not have the subspace topology.
}.
$\Id[{{U,V}}]$ is synonymous with $\Id[U,V]$.

$\Id[\catname{C}]$ is the identity functor on the category $\catname{C}$.

$\ID[\seqname{S}^i]$, $i=1,2$, is the sequence of identity functions
for the elements of the sequence
$\seqname{S}^i \defeq (\seqname{S}^i_\alpha,~ \alpha \prec \Alpha)$. Let
$\seqname{S}^1 \SUBSETEQ \seqname{S}^2$. Then
$\ID[\seqname{S}^1,\seqname{S}^2]$ is the sequence of inclusion maps
$
  \biggl (
    \Id[\seqname{S}^1_\alpha,\seqname{S}^2_\alpha],~ \alpha \prec \Alpha
  \biggr )
$
for the elements of the sequences $\seqname{S}^i$.
$\ID[{{\seqname{S}^1,\seqname{S}^2}}]$ is synonymous with
$\ID[\seqname{S}^1,\seqname{S}^2]$.  A similar definition applies for
tuples, except in a context where they are treated as objects of a
category.

$\ID[{{(A \dots Z)}}]$ is the tuple of identity maps
$\bigl (\Id[A], \dots, \Id[Z] \bigr)$.
$\ID[{(A^1, \dots, Z^1)},{(A^2, \dots, Z^2)}]$ is the tuple of inclusion maps
$\biggl ( \Id[A^1, A^2], \dots, \Id[Z^1, Z^2] \biggr )$.
\end{definition}

\begin{definition}[Images]
$\funcname{f} [U] \defeq \set { {\funcname{f}(x)} }[x \in U]$ is the
image of $U$ under $\funcname{f}$ and \\*
$\funcname{f}^{-1} [V] \defeq \set {x}[\funcname{f}(x) \in V]$ is the
inverse image of $V$ under $\funcname{f}$.

\begin{remark}
  This notation, adopted from \cite{GenTop}, avoids the ambiguity in
  the traditional $\funcname{f}(U)$ and $\funcname{f}^{-1}(V)$.
\end{remark}
\end{definition}

\begin{definition}[Projections]
\label{projections}
$\pi_\alpha$ is the projection function that maps a sequence into
element $\alpha$ of the sequence.  $\pi_i$ is also the projection
function that maps a tuple into element $i$ of the tuple.
\end{definition}

\begin{definition}[Topological category]
\label{def:topcat}
A topological category is a small subcategory of $\cat{Top}$ or its
concrete category over $\cat{Set}$.

$\catname{T}$ is a full topological category iff it is a topological
category and whenever $U^i,V^i \objin \catname{T}$, $i=1,2$,
$V^i \subsp U^i$,
$\funcname{f} \maps U^1 \to U^2 \arin \catname{T}$ and
$\funcname{f}[V^1] \subseteq V^2$ then \\
$\funcname{f} \maps V^1 \to V^2 \arin \catname{T}$.
\end{definition}

\begin{lemma}[Inclusions in topological categories are morphisms]
\label{lem:topInc}
Let $\catname{T}$ be a full topological category,
$S^i \objin \catname{T}$, $i=1,2$, and $S^1 \subsp S^2$. Then
$\Id[S^1,S^2]$ is a morphism of $\catname{T}$

\begin{proof}
\begin{enumerate}
\item
$U^i \defeq S^2 \objin \catname{T}$, $i=1,2$, by hypothesis.

\item
$V^1 \defeq S^1 \objin \catname{T}$ and
$V^2 \defeq S^2 \objin \catname{T}$ by hypothesis.

\item
$V^i \subsp U^i$,

\item
$\funcname{f} \defeq \Id[S^2] \maps U^1 \to U^2 \arin \catname{T}$.

\item
$\funcname{f}[V^1] = V^1 \subsp V^2$.
\end{enumerate}

Thus $\Id[S^1,S^2] = \funcname{f} \maps V^1 \to V^2 \arin \catname{T}$
by \cref{def:topcat}.
\end{proof}
\end{lemma}

\begin{definition}[Local morphisms]
\label{def:topLocal}
Let $\catname{S}^i$, $i=1,2$, be a full topological category and
$S^i \objin \catname{S}^i$, A continuous function
$\funcname{f} \maps S^1 \to S^2$ is a local
$\catname{S}^1$-$\catname{S}^2$ morphism of $S^1$ to $S^2$ iff
$\catname{S}^1 \subcat[full-] \catname{S}^2$ and for every
$u \in S^1$ there is an open neighborhood $U_u$ for $u$ and an open
neighborhood $V_u$ for $v \defeq \funcname{f}(u)$ such that
$\funcname{f}[U_u] \subseteq V_u$ and $\funcname{f} \maps U_u \to V_u$
is a morphism of $\catname{S}^2$.

\begin{remark}
The condition $\funcname{f} \maps U_u \to V_u \arin \catname{S}^2$
ensures that $U_u \objin \catname{S}^1$ and $V_u \objin \catname{S}^2$
\end{remark}

Let $\catname{T}$ be a full topological category and
$S^i \objin \catname{T}$, $i=1,2$. A continuous function
$\funcname{f} \maps S^1 \to S^2$ is a local $\catname{T}$ morphism of
$S^1$ to $S^2$ iff it is a local $\catname{T}$-$\catname{T}$ morphism
of $S^1$ to $S^2$.
\end{definition}

\begin{lemma}[Local morphisms]
\label{lem:topLocal}
Let $\catname{T}^i$, $i \in [1,3]$, be a full topological category,
$S^i \objin \catname{T}^i$ and
$\catname{T}^i \subcat[full-] \catname{T}^{i+1}$.

If $\funcname{f}^i \maps S^i \to S^{i+1} \arin \catname{T}^{i+1}$ then
$\funcname{f}^i$ is a local $\catname{T}^i$-$\catname{T}^{i+1}$
morphism of $S^i$ to $S^{i+1}$.

\begin{proof}
Let $u \in S^i$ and $v \defeq \funcname{f}^i(u) \in S^{i+1}$. $S^i$ is
an open for $u$, $S^{i+1}$ is an open neighborhood of $v$ and
$\funcname{f}^i \maps S^i \to S^{i+1} \arin \catname{T}^{i+1}$ by
hypothesis.
\end{proof}

If each $\funcname{f}^i \maps S^i \to S^{i+1}$, is a local
$\catname{T}^i$-$\catname{T}^{i+1}$ morphism of $S^i$ to
$S^{i+1}$ then
$\funcname{f}^2 \compose \funcname{f}^1 \maps S^1 \to S^3$ is a local
$\catname{T}^1$-$\catname{T}^3$ morphism of $S^1$ to $S^3$.

\begin{proof}
Since $\catname{T}^1 \subcat[full-] \catname{T}^2$ and
$\catname{T}^2 \subcat[full-] \catname{T}^3$,
$\catname{T}^1 \subcat[full-] \catname{T}^3$.
Let $u \in S^1$, $v \defeq \funcname{f}^1(u)$ and
$w \defeq \funcname{f}^2(v)$. There exist an open neighborhood $U_u$
for $u$, open neighborhoods $V_u$, $V'_v$ for $v$ and an open
neighborhood $W_v$ of $w$ such that
$\funcname{f}^1[U_u] \subseteq V_u$,
$\funcname{f}^1 \maps U_u \to V_u$ is a morphism of $\catname{T}^2$,
$\funcname{f}^2[V'_v] \subseteq W_v$ and
$\funcname{f}^2 \maps V'_v \to W_v$ is a morphism of $\catname{T}^3$.
Then $\hat{V_u} \defeq V_u \cap V'_v \neq \emptyset$, $\hat{V_u}$ is an
open neighborhood of $v$ and
$\hat{U_u} \defeq \funcname{f}^{i-1}_1[\hat{V_u}]$ is an open
neighborhood of $u$.  $\funcname{f}^1 \maps \hat{U_u} \to \hat{V_u}$
and $\funcname{f}^2 \maps \hat{V_u} \to W_v$ are morphisms of
$\catname{T}^3$ by \pagecref{def:topcat} and thus
$\funcname{f}^2 \compose \funcname{f}^1 \maps \hat{U_u} \to W_v$ is a
morphism of $\catname{T}^3$.
\end{proof}
\end{lemma}

\begin{corollary}[Local morphisms]
\label{cor:topLocal}
Let $\catname{T}^i$, $i=1,2$, be a full topological category,
$\catname{T}^i \subcat[full-] \catname{T}^{i+1}$,
$S^i \objin \catname{T}^i$ and $S^1 \subseteq S^2$. Then $\Id[S^1,S^2]$
is a local $\catname{T}^1$-$\catname{T}^2$ morphism of $S^1$ to $S^2$
and $\Id[{S^i}]$ is a local $\catname{T}$ morphism
of $S^i$ to $S^i$.

\begin{proof}
$S^1 \objin \catname{T}^2$ because $S^1 \objin \catname{T}^1$ and
$\catname{T}^1 \subcat \catname{T}^2$,
$S^2 \objin \catname{T}^2$ by hypothesis and \\*
$S^1 \subseteq S^2$ by hypothesis, so
$\Id[S^1,S^2] \arin \catname{T}^2$ by \cref{lem:topInc}.

$\Id[S^i] \defeq \Id[S^i,S^i]$.
\end{proof}
\end{corollary}

\begin{definition}[Sequence functions]
Let $\seqname{S} \defeq (s_\alpha, \alpha \prec \Alpha)$ be a sequence
of functions. Then
\begin{equation}
\Domain(\seqname{S}) \defeq \bigl ( \domain(s_\alpha), \alpha \prec \Alpha \bigr )
\end{equation}
\begin{equation}
\Range(\seqname{S}) \defeq \bigl ( \range(s_\alpha), \alpha \prec \Alpha \bigr )
\end{equation}

Let $\seqname{T} \defeq (t_\alpha, \alpha \prec \Alpha)$ be a sequence
of functions with $\Range(\seqname{S}) = \Domain(\seqname{T})$. Then
their composition is the sequence
$
  \seqname{T} \compose[()] \seqname{S} \defeq
  (t_\alpha \compose s_\alpha, \alpha \prec \Alpha)
$,

Let $\seqname{S} \defeq (s_\gamma, \gamma \preceq \Gamma)$, then these
functions extract information about the sequence:
\begin{equation}
\head(\seqname{S},\Omega) \defeq (s_\gamma, \gamma \prec \Omega)
\end{equation}
\begin{equation}
\head(\seqname{S}) \defeq \head(\seqname{S},\Gamma)
\end{equation}
\begin{equation}
\lengtho(\seqname{S}) \defeq \Gamma
\end{equation}
\begin{equation}
\tail(\seqname{S}) \defeq S_\Gamma
\end{equation}

Let $\seqname{S} \defeq (s_\gamma, \gamma \prec \Gamma)$, then
\begin{equation}
\length(\seqname{S}) \defeq \Gamma
\end{equation}

\begin{remark}
If $\lengtho(\seqname{S})$ is defined then
$\length(\seqname{S}) = \lengtho(\seqname{S}) + 1$.
$\lengtho(\seqname{S})$ is the ordinal type of
$\head(\seqname{S})$, not the ordinal type of $\seqname{S}$.
\end{remark}

Let $\catseqname{S} \defeq (\catname{S}_\alpha, \alpha \prec \Alpha)$
and $\catseqname{T} \defeq (\catname{T}_\alpha, \alpha \prec \Alpha)$ be
sequences of categories.  Then $\catseqname{S}$ is a subcategory
sequence of $\catseqname{T}$, abbreviated
$\catseqname{S} \SUBCAT \catseqname{T}$, iff every category in
$\catseqname{S}$ is a subcategory of the corresponding category in
$\catseqname{T}$, i.e.,
$
  \uquant%
    {\alpha \prec \Alpha}
    {\catname{S}_\alpha \subcat \catname{T}_\alpha}
$,
and $\catseqname{S}$ is a full subcategory sequence of $\catseqname{T}$,
abbreviated $\catseqname{S} \SUBCAT[full-] \catseqname{T}$, iff every
category in $\catseqname{S}$ is a full subcategory of the corresponding
category in $\catseqname{T}$, i.e.,
$
  \uquant%
    {\alpha \prec \Alpha}
    {\catname{S}_\alpha \subcat[full-] \catname{T}_\alpha}
$.

The category sequence union of $\catseqname{S}$ and $\catseqname{T}$,
abbreviated $\catseqname{S} \UNIONCAT \catseqname{T}$, is the sequence
of category unions of corresponding categories in $\catseqname{S}$ and
$\catseqname{T}$, i.e.,
$(\catseqname{S}_\alpha \unioncat \catseqname{T}_\alpha)$.

\end{definition}

\begin{lemma}[Sequence functions]
\label{lem:seqfunc}
Let $\funcseqname{f}^i \defeq (\funcname{f}^i_\alpha, \alpha \prec \Alpha)$,
$i \in [1,3]$, be sequences of functions with
\donotbreak
  {
    $
      \Domain(\funcseqname{f}^2) = \Range(\funcseqname{f}^1)
    $
  }
and
\donotbreak
  {
    $
      \Domain(\funcseqname{f}^3) = \Range(\funcseqname{f}^2)
    $.
  }
Then
$
(\funcseqname{f}^3 \compose[()] \funcseqname{f}^2) \compose[()] \funcseqname{f}^1 =
\funcseqname{f}^3 \compose[()] (\funcseqname{f}^2 \compose[()] \funcseqname{f}^1)
$.
\begin{proof}
\begin{equation*}
\begin{split}
  (\funcseqname{f}^3 \compose[()] \funcseqname{f}^2) \compose[()]
  \funcseqname{f}^1
  & =
  \bigl (
    (\funcname{f}^3_\alpha \compose \funcname{f}^2_\alpha) \compose
    \funcname{f}^1_\alpha,
    \alpha \prec \Alpha
  \bigr )
\\*
  & =
  \bigl (
    \funcname{f}^3_\alpha \compose
    (\funcname{f}^2_\alpha \compose \funcname{f}^1_\alpha),
    \alpha \prec \Alpha
  \bigr )
\\*
  & =
  \funcseqname{f}^3 \compose[()] (\funcseqname{f}^2 \compose[()] \funcseqname{f}^1)
\end{split}
\end{equation*}
\end{proof}

Let $\funcseqname{f} \defeq (\funcname{f}_\alpha, \alpha \prec \Alpha)$
be a sequence of functions, $\seqname{D} = \Domain(\funcseqname{f})$ and \\*
$\seqname{R} = \Range(\funcseqname{f})$. Then $\ID[\seqname{R}]$ is a
left $\compose[()]$ identity for $\funcseqname{f}$ and
$\ID[\seqname{D}]$ is a right $\compose[()]$ identity for
$\funcseqname{f}$.

\begin{proof}
\begin{equation*}
\begin{split}
  \ID[\seqname{R}] \compose[()] \funcseqname{f}
  & =
  (\Id[\range(\funcname{f}_\alpha)] \compose \funcname{f}_\alpha,
\alpha \prec \Alpha)
\\*
  & =
  (\funcname{f}_\alpha, \alpha \prec \Alpha)
\\*
  & =
  \funcseqname{f}
\end{split}
\end{equation*}
\begin{equation*}
\begin{split}
  \funcseqname{f} \compose[()] \ID[\seqname{D}]
  & =
  (\funcname{f}_\alpha \compose \Id[\domain(\funcname{f}_\alpha)],
\alpha \prec \Alpha)
\\*
  & =
  (\funcname{f}_\alpha, \alpha \prec \Alpha)
\\*
  & =
  \funcseqname{f}
\end{split}
\end{equation*}
\end{proof}
\end{lemma}

\begin{definition}[Tuple functions]
Let $\seqname{S} \defeq (s_n, n \in [1,N])$ be a tuple of functions. Then
\begin{equation}
\Domain(\seqname{S}) \defeq \bigl ( \domain(s_n), n \in [1,N] \bigr )
\end{equation}
\begin{equation}
\Range(\seqname{S}) \defeq \bigl ( \range(s_n), n \in [1,N] \bigr )
\end{equation}

Let $\seqname{T} \defeq (t_n, n \in [1,N])$ be a tuple of functions with
$\Range(\seqname{S})=\Domain(\seqname{T})$,
Then their composition is the tuple
$
  \seqname{T} \compose[()] \seqname{S} \defeq
  (t_n \compose s_n, n \in [1,N])
$

Let $\seqname{S} \defeq (s_m, m \in [1,M])$ and
$\seqname{T} \defeq (t_n, n \in [1,N])$ be tuples.
Then the following are tuple functioms
\begin{equation}
\head(S,I) \defeq (s_m, m \in [1,I])
\end{equation}
\begin{equation}
\head(S) \defeq \head(S,M-1)
\end{equation}
\begin{equation}
\tail(T,I) \defeq (t_n, n \in [I,N])
\end{equation}
\begin{equation}
\tail(T) \defeq t_N
\end{equation}
\begin{equation}
\join(S,T) \defeq (s_1, \dots, s_M, t_1, \dots, t_N)
\end{equation}
\end{definition}

\begin{lemma}[Tuple functions]
\label{lem:tupfunc}
Let $\funcseqname{f} \defeq (\funcname{f}_n, n \in [1,N])$ be a tuple of
functions,
$\seqname{D} \defeq \Domain(\funcseqname{f})$ and
$\seqname{R} \defeq \Range(\funcseqname{f})$. Then $\ID[\seqname{R}]$ is a
left $\compose[()]$ identity for $\funcseqname{f}$ and $\ID[\seqname{D}]$
is a right $\compose[()]$ identity for $\funcseqname{f}$.

\begin{proof}
\begin{equation*}
\begin{split}
  \ID[\seqname{R}] \compose[()] \funcseqname{f}
  & =
  (\Id[\range(\funcname{f}_\alpha)] \compose \funcname{f}_n, n \in [1,N])
\\*
  & =
  (\funcname{f}_n, n \in [1,N])
\\*
  & =
  \funcseqname{f}
\end{split}
\end{equation*}
\begin{equation*}
\begin{split}
  \funcseqname{f} \compose[()] \ID[\seqname{D}]
  & =
  (\funcname{f}_n \compose \Id[\domain(\funcname{f}_n)], n \in [1,N])
\\*
  & =
  (\funcname{f}_n, n \in [1,N])
\\*
  & =
  \funcseqname{f}
\end{split}
\end{equation*}
\end{proof}
\end{lemma}

\begin{lemma}[Tuple composition for unlabeled morphisms]
\label{lem:tupcomp}
Let \\*
$
  \funcseqname{f}^i \defeq (\funcname{f}^i_n,~ n \in [1,N]),~
  i \in [1,3]
$,
be tuples of functions with
\donotbreak
  {
    $
      \Domain(\funcseqname{f}^2) = \Range(\funcseqname{f}^1)
    $
  }
and
\donotbreak
  {
    $
      \Domain(\funcseqname{f}^3) = \Range(\funcseqname{f}^2)
    $.
  }
Then
$
(\funcseqname{f}^3 \compose[()] \funcseqname{f}^2) \compose[()] \funcseqname{f}^1 =
\funcseqname{f}^3 \compose[()] (\funcseqname{f}^2 \compose[()] \funcseqname{f}^1)
$.

\begin{proof}
\begin{equation*}
\begin{split}
  (\funcseqname{f}^3 \compose[()] \funcseqname{f}^2) \compose[()] \funcseqname{f}^1
  & =
  \bigl ( (\funcname{f}^3_n \compose \funcname{f}^2_n) \compose \funcname{f}^1_n, n \in [1,N] \bigr )
\\*
  & =
  \bigl ( \funcname{f}^3_n \compose (\funcname{f}^2_n \compose \funcname{f}^1_n), n \in [1,N] \bigr )
\\*
  & =
  \funcseqname{f}^3 \compose[()] (\funcseqname{f}^2 \compose[()] \funcseqname{f}^1)
\end{split}
\end{equation*}
\end{proof}
\end{lemma}

\begin{definition}[Tuple composition for labeled morphisms]
\label{def:labcomp}
Let $\seqname{M}^i \defeq (\funcseqname{f}^i, o^i_1, o^i_2)$, $i=1,2$,
be tuples such that each $\funcseqname{f}^i$ is a sequence of functions
or each $\funcseqname{f}^i$ is a tuple of functions,
$\Range(\funcseqname{f}^1) = \Domain(\funcseqname{f}^2)$ and
$0^1_2=o^2_1$. Then

\begin{equation}
\seqname{M}^2 \compose[A] \seqname{M}^1 \defeq
\bigl (
  \funcseqname{f}^2 \compose[()] \funcseqname{f}^1,
  o^1_1,
  o^2_2
\bigr )
\end{equation}
\end{definition}

\begin{lemma}[Tuple composition for labeled morphisms]
\label{lem:labcomp}
Let $\seqname{M}^i \defeq (\funcseqname{f}^i, o^i_1, o^i_2)$,
\donotbreak{$i \in [1,3]$,}
be a tuple such that $\funcseqname{f}^i$ is a sequence or tuple of
functions,
\donotbreak
  {
    $
      \Range(\funcseqname{f}^i) =
      \Domain(\funcseqname{f}^{i+1})
    $
  }
and $o^i_2 = o^{i+1}_1$, $i=1,2$.
Then
\begin{equation}
\seqname{M}^3 \compose[A] \bigl ( \seqname{M}^2 \compose[A] \seqname{M}^1 \bigr )
=
\bigl ( \seqname{M}^3 \compose[A] \seqname{M}^2 \bigr ) \compose[A] \seqname{M}^1
\end{equation}

\begin{proof}
From \fullcref{def:labcomp}, \\
{
  \showlabelsinline
  \pagecref{lem:seqfunc}
}
and \fullcref{lem:tupfunc}, we have
\begin{equation*}
\begin{split}
  \seqname{M}^3 \compose[A] (\seqname{M}^2 \compose[A] \seqname{M}^1)
  & =
  \seqname{M}^3 \compose[A] (\funcseqname{f}^2 \compose[()] \funcseqname{f}^1, o^1_1, o^2_2)
\\*
  & =
  (\funcseqname{f}^3 \compose[()] \funcseqname{f}^2 \compose[()] \funcseqname{f}^1, o^1_1, o^3_2)
\\*
  & =
  (\funcseqname{f}^3 \compose[()] \funcseqname{f}^2, o^2_1, o^3_2) \compose[A] \seqname{M}^1
\\*
  & =
  (\seqname{M}^3 \compose[A] \seqname{M}^2) \compose[A] \seqname{M}^1
\end{split}
\end{equation*}
\end{proof}

Let $\seqname{D}^i \defeq \Domain(\funcseqname{f}^i)$ and
$\seqname{R}^i \defeq \Range(\funcseqname{f}^i)$. Then
\begin{equation}
(\ID[\seqname{R}^i], o^i_2, o^i_2) \compose[A] \seqname{M}^i =
\seqname{M}^i
\end{equation}
\begin{equation}
\seqname{M}^i \compose[A] (\ID[\seqname{D}^i], o^i_1, o^i_1) =
\seqname{M}^i
\end{equation}

\begin{proof}
\begin{equation*}
\begin{split}
  (\ID[\seqname{R}^i], o^i_2, o^i_2) \compose[A] \seqname{M}^i
  & =
  (\ID[\seqname{R}^i] \compose[()] \funcseqname{f}^i, o^i_1, o^i_2)
\\*
  & =
  (\funcseqname{f}^i, o^i_1, o^i_2)
\\*
  & =
 \seqname{M}^i
\end{split}
\end{equation*}
\begin{equation*}
\begin{split}
  \seqname{M}^i \compose[A] (\ID[\seqname{D}^i], o^i_1, o^i_1)
  & =
  (\funcseqname{f}^i \compose[()] \ID[\seqname{D}^i], o^i_1, o^i_2)
\\*
  & =
  (\funcseqname{f}^i, o^i_1, o^i_2)
\\*
  & =
 \seqname{M}^i
\end{split}
\end{equation*}
\end{proof}
\end{lemma}

\begin{definition}[Cartesian product of sequence]
\label{def:Cart}
Let $\seqname{S}^i \defeq (S^i_\alpha, \alpha \prec \Alpha)$, $i=1,2$,
be a sequence and
$\funcseqname{f} \defeq (\funcname{f}_\alpha \maps S^1_\alpha \to
S^2_\alpha, \alpha \prec \Alpha)$
be a sequence of functions, then
$
  \bigtimes \seqname{S}^i \defeq
  \bigtimes_{\alpha \prec \Alpha} S^i_\alpha
$
is the generalized Cartesian product of the sequence $\seqname{S}^i$ and \\*
\donotbreak
  {
    $
      \bigtimes \funcseqname{f} \maps \seqname{S}^1 \to \seqname{S}^2
      \defeq
      \bigtimes_{\alpha \prec \Alpha}\funcname{f}_\alpha
    $
  }
is the generalized Cartesian product of the function sequence
$\funcseqname{f}$.
\end{definition}


\begin{definition}[Topology functions]
Let $S$ and $T$ be a topological spaces and $Y$ a subset of $S$. Then

\begin{enumerate}
\item
$S \subsp T$ iff $S$ is a subspace of $T$.
\item
$S \subsp[open-] T$ iff $S$ is an open subspace of $T$.
\item
$\Set(S)$ is the underlying set of $S$.
\item
$\Topology(S)$ is the topology of $S$.
\item
$\Topology(Y,S) \defeq \set{U \cap Y}[{U \in \Topology(S)}]$
is the relative topology of $Y$.
\item
$\Top(Y,S) \defeq \bigl ( Y, \Topology(Y,S) \bigr )$ is $Y$ with the
relative topology.
\item
$
  \op{S} \defeq
  \set
    {(U, \Topology(U,S))}%
    [{{ U \in \Topology(S) \setminus \{ \emptyset \} }}]
$
is the set of all non-null open subspaces of $S$.
\end{enumerate}

Let $\seqname{S}$ be a set of topological spaces. Then
$\op{\seqname{S}} \defeq \union [{S \in \seqname{S}}] { {\op{S}}}$ is the set
of open subspaces in $\seqname{S}$.

Let $S$ and $T'$ be spaces, $T \subseteq T'$ be a subspace and
$\funcname{f} \maps S \to T$ a function. Then
$\funcname{f} \maps S \to T' \defeq \Id[{T},{T'}] \compose \funcname{f}$
is $\funcname{f}$ considered as a function from $S$ to $T'$.

Let $S'$ and $T'$ be spaces, $S \subseteq S'$, $T \subseteq T'$ be
subspaces and $\funcname{f}' \maps S' \to T'$ a function such that
$\funcname{f}'[S] \subseteq T$. Then $\funcname{f}' \maps S \to T$, also
written $\funcname{f}' \restrictto_{S,T}$, is
$\funcname{f}' \restrictto_S$ considered as a function from $S$ to
$T$.

Let $\seqname{S}^i \defeq (S^i_\alpha, \alpha \preceq \Alpha)$, $i-1,2$,
be a sequence of spaces, $\seqname{S}^1 \SUBSETEQ \seqname{S}^2$ and \linebreak
$\funcname{f}^2 \maps \head(\seqname{S}^2) \to \tail(\seqname{S}^2)$ a
function.
$
  \funcname{f}^2 \restrictto_{\head(\seqname{S}^1)} \defeq
  \funcname{f}^2 \restrictto_{\bigtimes \head(\seqname{S}^1)}
$.
If \linebreak
$
  \funcname{f}^2 \restrictto_{\head(\seqname{S}^1)}
    [\bigtimes \head(\seqname{S}^1)]
  \subseteq \tail(\seqname{S}^1)
$
then
$ \funcname{f}^2 \maps \head(\seqname{S}^1) \to \tail(\seqname{S}^1)$,
also written
$
  \funcname{f}^2 \restrictto_%
    {\head(\seqname{S}^1),\tail(\seqname{S}^1)}
$,
is $\funcname{f}^2 \restrictto_{\head(\seqname{S}^1)}$ considered as a
function from $\bigtimes \head(\seqname{S}^1)$ to $\tail(\seqname{S}^1)$.
\end{definition}

\begin{definition}[Sequence inclusion]
\label{def:seqin}
Let $\seqname{S} \defeq (S_\alpha, \alpha \prec \Alpha)$ and
$\seqname{T} \defeq (T_\alpha, \alpha \prec \Alpha)$ be sequences.
$\seqname{S} \seqin \seqname{T}$ iff
$\uquant{\alpha \prec \Alpha} {S_\alpha \in T_\alpha}$ or
$\uquant{\alpha \prec \Alpha} {S_\alpha \objin T_\alpha}$.
\end{definition}

\begin{lemma}[Sequence inclusion]
\label{lem:seqin}
Let $\seqname{S} \defeq (S_\alpha, \alpha \prec \Alpha)$
be a sequence, \\*
\nobreaks
{{
  $
    \catseqname{T}^i \defeq
    (\catname{T}^i_\alpha, \alpha \prec \Alpha)
  $,
}}
$i=1,2$, a sequence of
categories, $\catseqname{T}^1 \SUBCAT \catseqname{T}^2$ and
$\seqname{S} \seqin \catseqname{T}^1$ Then
$\seqname{S} \seqin \catseqname{T}^2$.

\begin{proof}
If
$\uquant{\alpha \prec \Alpha} {S_\alpha \objin \catname{T}^1_\alpha}$
then
$\uquant{\alpha \prec \Alpha} {S_\alpha \objin \catname{T}^2_\alpha}$.
\end{proof}
\end{lemma}

\part{Nearly commutative diagrams}
\label{part:ncd}
The notion of commutative diagrams is very useful in, e.g., Algebraic
Topology. Often one encounters commutative diagrams in which two
outgoing terminal nodes can be connected by a bridging function such
that the resulting diagram is still commutative. This paper uses the
term nearly commutative to describe a restricted class of such diagrams.

Let $\catname{C}$ be a full topological category and $D$ a tree with two
branches, whose nodes are topological spaces $U_i$ and $V^j$ and whose
links are continuous functions $\funcname{f}_i \maps U_i \to U_{i+1}$
and $\funcname{f}'_j \maps U_j \to U_{j+1}$ between the sets:

\begin{equation}
\begin{split}
\label{eq:NCD}
D \defeq \{
  &
    U_0 = V_0 \to^{\funcname{f}_0} U_1,
    \dotsc,
    U_{m-1} \to^{\funcname{f}_{m-1}} U_m,
\\
  &
    U_0 = V_0 \to^{\funcname{f}'_0} V_1,
    \dotsc,
    V_{n-1} \to^{\funcname{f}'_{n-1}} V_n
\}
\end{split}
\end{equation}
with $U_0 = V_0$, $U_m \objin \catname{C}$ and
$V_n \objin \catname{C}$, as shown in
{
  \showlabelsinline
  \cref{fig:NCDb}.
}

\begin{figure}
\[ \bfig
\node u0(0,0)[U_0=V_0]
\node u1(-500,-500)[U_1]
\node v1(500,-500)[V_1]
\node d1(-500,-1000)[\vdots]
\node d2(500,-1000)[\vdots]
\node um(-500,-1500)[U_m]
\node vn(500,-1500)[V_n]
\arrow |l|[u0`u1;\funcname{f}_0]
\arrow |r|[u0`v1;\funcname{f}'_0]
\arrow |l|[u1`d1;\funcname{f}_1]
\arrow |r|[v1`d2;\funcname{f}'_1]
\arrow |l|[d1`um;\funcname{f}_{m-1}]
\arrow |r|[d2`vn;\funcname{f}'_{n-1}]
\efig \]

\caption{Uncompleted nearly commutative diagram}
\label{fig:NCDb}
\end{figure}

\begin{definition}[Nearly commutative diagrams in category $\catname{C}$]
\label{def:NCD}
$D$ is left (right) nearly commutative in category $\catname{C}$ iff the
two final nodes are in $\catname{C}$ and there is a morphism
$\hat{\funcname{f}} \maps U_m \to V_n \arin \catname{C}$
($\hat{\funcname{f}'} \maps V_n \to U_m \arin \catname{C}$)
making the graph a commutative diagram, as shown in \cref{fig:NCDa}.
$D$ is nearly commutative in category $\catname{C}$ iff it is eiher
left nearly commutative or right nearly commutative.

$D$ is strongly nearly commutative in category $\catname{C}$ iff the
morphism is an isomorphism of $\catname{C}$,
\end{definition}

\begin{figure}
\[ \bfig
\node u0(0,0)[U_0]
\node u1(-500,-500)[U_1]
\node v1(500,-500)[V_1]
\node d1(-500,-1000)[\vdots]
\node d2(500,-1000)[\vdots]
\node um(-500,-1500)[U_m]
\node vn(500,-1500)[V_n]
\arrow |l|[u0`u1;\funcname{f}_0]
\arrow |r|[u0`v1;\funcname{f}'_0]
\arrow |l|[u1`d1;\funcname{f}_1]
\arrow |r|[v1`d2;\funcname{f}'_1]
\arrow |l|[d1`um;\funcname{f}_{m-1}]
\arrow |r|[d2`vn;\funcname{f}'_{n-1}]
\arrow |a|/{@{->}@/^5pt/}/[um`vn;\hat{\funcname{f}}]
\arrow |b|/{@{->}@/^5pt/}/[vn`um;\hat{\funcname{f}'}]
\efig \]
\caption{Completed nearly commutative diagram}
\label{fig:NCDa}
\end{figure}

\begin{definition}[Nearly commutative diagrams in category $\catname{C}$ at a point]
\label{def:NCDp}
Let $D$ be as above and $x$ be an element of the initial node. $D$ is
(left,right,strongly) nearly commutative in $\catname{C}$ at $x$ iff
there are subobjects of the nodes such that the tree $D'$ formed by
replacing the nodes is (left,right,strongly) nearly commutative in
$\catname{C}$, and $x$ is in the new initial node, as shown in
\pagecref{fig:NCDl} below:
\begin{equation}
x \in U'_0 = V'_0
\end{equation}
\begin{equation}
U'_i \subseteq U_i
\end{equation}
\begin{equation}
V'_j \subseteq V_j
\end{equation}
\begin{equation}
\funcname{f}_i \restrictto_{U'_i} \maps U'_i \to U'_{i+1}
\end{equation}
\begin{equation}
\funcname{f}'_j \restrictto_{V'_j} \maps V'_j \to V'_{j+1}
\end{equation}

\begin{equation}
\begin{split}
\label{eq:NCDp}
D' \defeq \{
  &
    U'_0 = V'_0 \to^{\funcname{f}_0} U'_1,
    \dotsc,
    U'_{m-1} \to^{\funcname{f}_{m-1}} U'_m,
\\
  &
    U'_0 = V'_0 \to^{\funcname{f}'_0} V'_1,
    \dotsc,
    V'_{n-1} \to^{\funcname{f}'_{n-1}} V'_n
\}
\end{split}
\end{equation}
\end{definition}

\begin{figure}
\[ \bfig
\node x(0,500)[x]
\node u0(-1000,0)[U_0]
\node up0(0,0)[U'_0]
\node v0(1000,0)[V_0]
\node u1(-1000,-500)[U_1]
\node up1(-500,-500)[U'_1]
\node vp1(500,-500)[V'_1]
\node v1(1000,-500)[V_1]
\node d1(-1000,-1000)[\vdots]
\node d2(-500,-1000)[\vdots]
\node d3(500,-1000)[\vdots]
\node d4(1000,-1000)[\vdots]
\node um(-1000,-1500)[U_m]
\node upm(-500,-1500)[U'_m]
\node vpn(500,-1500)[V'_n]
\node vn(1000,-1500)[V_n]
\arrow |r|/_{ (}->/[x`up0;i]
\arrow |a|/_{ (}->/[up0`u0;i]
\arrow |b|[up0`up1;\funcname{f}_0]
\arrow |b|[up0`vp1;\funcname{f}'_0]
\arrow |a|/^{ (}->/[up0`v0;i]
\arrow |l|[u0`u1;\funcname{f}_0]
\arrow |r|[v0`v1;\funcname{f}'_0]
\arrow |l|[u1`d1;\funcname{f}_1]
\arrow |b|/_{ (}->/[up1`u1;i]
\arrow |r|[up1`d2;\funcname{f}_1]
\arrow |b|/^{ (}->/[vp1`v1;i]
\arrow |l|[vp1`d3;\funcname{f}'_1]
\arrow |r|[v1`d4;\funcname{f}'_1]
\arrow |l|[d1`um;\funcname{f}_{m-1}]
\arrow |b|/_{ (}->/[d2`d1;i]
\arrow |r|[d2`upm;\funcname{f}_{m-1}]
\arrow |l|/^{ (}->/[d3`d4;i]
\arrow |l|[d3`vpn;\funcname{f}'_{n-1}]
\arrow |r|[d4`vn;\funcname{f}'_{n-1}]
\arrow |b|/_{ (}->/[upm`um;i]
\arrow |a|/{@{->}@/^5pt/}/[upm`vpn;\hat{\funcname{f}}]
\arrow |b|/{@{->}@/^5pt/}/[vpn`upm;\hat{\funcname{f}'}]
\arrow |b|/^{ (}->/[vpn`vn;i]
\efig \]
% \[\begin{tikzcd}
%   &
%   &
%   x
%     \ar[d, hook, "\funcname{i}"]
% \\
%   U_0
%     \ar[d, "\funcname{f}_0"]
%   &
%   &
%   U'_0
%     \ar[ll, hook', swap, "\funcname{i}"]
%     \ar[dl, "\funcname{f}_0"]]
%     \ar[dr, swap, "\funcname{f}'_0"]
%     \ar[rr, hook,"\funcname{i}"]
%   &
%   &
%   U_0
%     \ar[d, "\funcname{f}'_0"]
% \\
%   U_1
%     \ar[d, "\funcname{f}_1"']
%   &
%   U'_1
%     \ar[l, hook', "\funcname{i}"]
%     \ar[d, "\funcname{f}_1"]
%   &
%   &
%   V'_1
%     \ar[r, hook, "\funcname{i}"']
%     \ar[d, "\funcname{f}'_1"]
%   &
%   V_1
%     \ar[d, "\funcname{f}'_1"]
% \\
%   \vdots
%     \ar[d, "\funcname{f}_{m-1}"]
%   &
%   \vdots
%     \ar[l, hook', "\funcname{i}"]
%     \ar[d, "\funcname{f}_{m-1}"]
%   &
%   &
%   \vdots
%     \ar[r, hook, "\funcname{i}"']
%     \ar[d, "\funcname{f}'_{n-1}"]
%   &
%   \vdots \ar[d, "\funcname{f}'_{n-1}"]
% \\
%   U_m
%   &
%   U'_m
%     \ar[l, hook', "\funcname{i}"]
%     \ar[rr, dotted, "\iso", "\hat{\funcname{f}}"'] &
%   &
%   V'_n
%     \ar[r, hook, "\funcname{i}"']
%   &
%   V_n
% \end{tikzcd}\]
\caption{Local nearly commutative diagram}
\label{fig:NCDl}
\end{figure}

\begin{definition}[Locally nearly commutative diagrams in category $\catname{C}$]
\label{def:NCDl}
Let $D$ be as above. $D$ is (left,right,strongly) locally nearly
commutative in $\catname{C}$ iff it is (left,right,strongly) nearly
commutative in $\catname{C}$ at $x$ for every $x$ in the initial node.
\end{definition}

\begin{remark}
It will often be clear from context what the relevant categories are.
This paper may use the term ``nearly commutative'' without explicitly
identifying the categories in which the modes are found.
\end{remark}

\begin{lemma}[Locally nearly commutative diagrams in category $\catname{C}$]
\label{lem:NCD}
Let $\catname{C}$ and $D$ be as above. If $U_0 = \emptyset$ then $D$ is
(left, right, strongly) locally nearly commutative in $\catname{C}$.

\begin{proof}
$D$ is vacuously locally nearly commutative at every $x \in U_0$ since
there is no such $x$.
\end{proof}

Let $D$ be locally nearly commutative in $\catname{C}$ and
let $\hat{U_0}=\hat{V_0} \subseteq U_0$ and $\hat{D}$ be $D$ with
$U_0=V_0$ replaced by $\hat{U_0}=\hat{V_0}$. Then $\hat{D}$ is locally
nearly commutative in $\catname{C}$.

\begin{proof}
If $x \in \hat{U_0}$ then $x \in U_0$ and hence $D$ is locally nearly
commutative in $\catname{C}$ at $x$.
Replacing $U'_0$ with $U'_0 \cap \hat{U_0}$ in the definition shows that
$\hat{D}$ is locally nearly commutative in $\catname{C}$ at $x$.
\end{proof}
\end{lemma}

\begin{lemma}[Nearly commutative diagrams in category $\catname{C}'$]
\label{lem:NCDsubcat}
Let $\catname{C} \subcat \catname{C}'$ and $D$ be as in diagram
\eqref{eq:NCD} above.

If $D$ is (left, right, strongly) nearly commutative in $\catname{C}$
then $D$ is (left, right, strongly) commutative in $\catname{C}'$.

\begin{proof}
If $U_m \objin \catname{C}$ then $U_m \objin \catname{C}'$.
If $V_n \objin \catname{C}$ then $V_n \objin \catname{C}'$.
If $\hat{\funcname{f}} \maps U_m \toiso V_n \arin \catname{C}$ then
$\hat{\funcname{f}} \arin \catname{C}'$
\end{proof}
\end{lemma}

\begin{corollary}[Locally nearly commutative diagrams in category $\catname{C}'$ at a point]
\label{cor:NCDp}
Let $\catname{C} \subcat \catname{C}'$ and $D$ be as in diagram
\eqref{eq:NCD} above.

If $D$ is locally nearly commutative in $\catname{C}$ at $x$ then $D$ is
locally nearly commutative in $\catname{C}'$ at $x$.

\begin{proof}
Let $U'_i$, $V'_j$ and diagram \eqref{eq:NCDp} be as in
\pagecref{def:NCDp}\!, as shown in diagram \fullcref{fig:NCDl}\!.  $D'$
is nearly commutative in category $\catname{C}$, and hence is nearly
commutative in category $\catname{C}$.
\end{proof}
\end{corollary}

\begin{corollary}[Locally nearly commutative diagrams in category $\catname{C}'$]
\label{cor:NCDn}
Let $\catname{C} \subcat \catname{C}'$ and $D$ be as in diagram
\eqref{eq:NCD} above.

If $D$ is is locally nearly commutative in $\catname{C}$ then $D$
is locally nearly commutative in $\catname{C}'$.

\begin{proof}
Let $x \in U_0 = V_0$. $D$ is is locally nearly commutative in
$\catname{C}$ at $x$, hence locally nearly commutative in $\catname{C}'$
at $x$.
\end{proof}
\end{corollary}

\part{Model spaces and allied notions}
\label{part:modeldef}
Let $S$ be a topological space. We need to formalize the notions of an
open cover by sets that are ``well behaved'' in some sense, e.g.,
convex, sufficiently small, and of {\textquotedblleft}well
behaved{\textquotedblright} functions among those sets, e.g., preserving
fibers, smooth. We do this by associating a category of acceptable sets
and functions.

\begin{remark}
Using pseudo-groups, as in \cite[p.~1]{FoundDiffGeo1}, would not allow
restricting model neighborhoods to, e.g., convex sets.
\end{remark}

\section{Model spaces}
\begin{definition}[Model spaces]
\label{def:model}
Let $S$ be a topological space and $\catname{S}$ a small category whose
objects are open subsets of $S$ and whose morphisms are continuous
functions. $\seqname{S} \defeq (S, \catname{S})$ is a model space for
$S$ iff

\begin{enumerate}
\item $\Ob(\catname{S})$ is an open cover for $S$. Note that it need
not be a basis for $S$.
\label{mod:cover}
\item $\Ob(\catname{S})$ is closed under finite intersections.
\label{mod:closed}
\item The morphisms of $\catname{S}$ are continuous functions in $S$.
\label{mod:continuous}
\item If $\funcname{f} \maps A \to B$ is a morphism,
$A' \objin \catname{S} \subseteq A \objin \catname{S}$,
$B' \objin \catname{S} \subseteq B \objin \catname{S}$ and
\nobreaks
{{
  $\funcname{f}[A'] \subseteq B'$
}}
then
$\funcname{f} \restrictto_{A'} \maps A' \to B'$ is a morphism.
\label{mod:restriction}
\item If $A' \objin \catname{S} \subseteq A \objin \catname{S}$ then
the inclusion map $\Id[{A'},{A}] \maps A' \hookrightarrow A$ is a
morphism.
\label{mod:inclusion}
\begin{remark}
This is actually a consequence of \cref{mod:restriction},
but it is convenient to give it here.
\end{remark}
\item Restricted sheaf condition:
\label{mod:sheaf}
informally, consistent morphisms can be glued together. Whenever
\begin{enumerate}
\item $U_\alpha$ and $V_\alpha$, $\alpha \prec \Alpha$,
are objects of $\catname{S}$.
\item $\funcname{f}_\alpha \maps U_\alpha \to V_\alpha$ are morphisms of
$\catname{S}$.
\item
$U = \union[\alpha \prec \Alpha]{U_\alpha}$ is an object of $\catname{S}$.
\item
$V = \union[\alpha \prec \Alpha]{V_\alpha}$ is an object of $\catname{S}$.
\item
$\funcname{f} \maps U \to V$ is a continuous function and
for every $\alpha \prec \Alpha$,
$\funcname{f}$ agrees with $\funcname{f}_\alpha$ on $U_\alpha$
\end{enumerate}
then $\funcname{f}$ is a morphism of $\catname{S}$.
\end{enumerate}

By abuse of language we write
{\textquotedblleft}$U \objin \seqname{S}${\textquotedblright} for
{\textquotedblleft}$U \objin \Cat{\seqname{S}}${\textquotedblright} and
{\textquotedblleft}$\funcname{f} \arin \seqname{S}${\textquotedblright}
for
{\textquotedblleft}$\funcname{f} \arin \Cat{\seqname{S}}${\textquotedblright}.

\begin{remark}
A model space is similar to a pseudo-group, but has some important
differences:
\begin{enumerate}
\item
Open subsets of objects of $\catname{S}$ need not be objects of
$\catname{S}$.

\item
The morphisms of $\catname{S}$ need not be homeomorphisms or even 1-1,
thus they need not have inverses.

\item
The morphisms of $\catname{S}$ need not satisfy the sheaf condition,
but only the restricted sheaf condition in \cref{mod:sheaf}.
\end{enumerate}
\end{remark}

$\seqname{S}$ is fine grained iff every open subset of an object of
$\catname{S}$ is an object of $ \catname{S}$. Note that it is not
sufficient that $\Ob(\catname{S})$ be a basis for $S$.

$\Top(\seqname{S}) \defeq S = \pi_1(\seqname{S})$ is the topological
space of $\seqname{S}$.

Let $\seqname{S}^i$, $i=1,2$, be a model space.
$\Cat(\seqname{S}^i) \defeq \catname{S}^i = \pi_2(\seqname{S}^i)$ is the
category of model space $\seqname{S}^i$.
By abuse of language we write $\seqname{S}^1 \subsp \seqname{S}^2$
$\bigl ( \seqname{S}^1 \subsp[open-] \seqname{S}^2 \bigr )$ for
$\Top(\seqname{S}^1) \subsp \Top(\seqname{S}^2)$
$\bigl ( \Top(\seqname{S}^1) \subsp[open-] \Top(\seqname{S}^2) \bigr )$ .

Let $U \objin \catname{S}$. Then $\Top(U,\seqname{S}) \defeq \Top(U,S)$
is $U$ with the relative topology.

$
  \Topcat(\seqname{S}) \defeq
  \Biggl (
    \set
      {{\Top(U,S)}}%
      [{U \objin \catname{S}}],
    \Ar(\catname{S})
  \Biggr )
$
is the topological category of $\seqname{S}$.

By abuse of language we write $U \subseteq \seqname{S}$ for
$U \objin \catname{S}$, $\Topology(\seqname{S})$ for
$\Topology \bigl ( \Top(\seqname{S}) \bigr )$, $\funcname{f}$ is a
morphism of $\seqname{S}$ for $\funcname{f}$ is a morphism of
$\catname{S}$ and $\funcname{f}$ is an isomorphism of $\seqname{S}$ for
$\funcname{f}$ is an isomorphism of $\catname{S}$.

Let $\seqname{C}$ be a set of model spaces. $\seqname{C}$ is fine
grained iff every $(C,\catname{C}) \in \seqname{C}$ is fine grained.
\end{definition}

\begin{lemma}[The topological category of a model space is a full topological category]
\label{lem:modtopcat}
Let $\seqname{M} \defeq (S, \catname{S})$ be a model space for $S$. Then
$\Topcat(\seqname{M})$ is a full topological category.
\begin{proof}
$\Topcat(\seqname{M})$ is a small subcategory of $\cat{Top}$ by
construction. $\Topcat(\seqname{M})$ is a full topological category by
\cref{mod:restriction} of
{
  \showlabelsinline
  \cref{def:model}.
}
\end{proof}
\end{lemma}

\begin{definition}[Model neighborhoods]
Let $\seqname{S} \defeq (S, \catname{S})$ be a model space for $S$.
Then the objects of $\catname{S}$ are model neighborhoods of
$\seqname{S}$. If $u \in U \objin \catname{S}$ then $U$ is a model
neighborhood of $\seqname{S}$ for $u$. If $U_i \objin \catname{S}$,
$i=1,2$, and $U_1 \subseteq U_2$ then $U_1$ is a model subneighborhood
of $U_2$; if $u^1 \in U_1$ then $U_1$ is also a model subneighborhood
for $u^1$.

If $\emptyset \objin \catname{S}$ then $\emptyset$ is a degenerate model
neighborhood of $\seqname{S}$; any nonvoid model neighborhood of
$\seqname{S}$ is a nondegenerate model neighborhood of $\seqname{S}$.
\end{definition}

\begin{definition}[Model subspaces]
\label{def:modsub}
$\seqname{S} \defeq (S, \catname{S})$ is a model subspace of
$\seqname{T} \defeq (T, \catname{T})$, abbreviated
$\seqname{S} \submod \seqname{T}$, iff

\begin{enumerate}
\item $\seqname{S}$ and $\seqname{T}$ are model spaces
\item $S$ is a subspace of $T$
\label{modsub:sub}

\item $\catname{S} \subcat \catname{T}$
\label{modsub:subcat}

\item every intersection of $S$ with an object of $\catname{T}$ is an object of$\catname{S}$
\label{modsub:rest}
\begin{equation}
\uquant%
  {U \objin \catname{T}}
  {U \cap S \objin \catname{S}}
\end{equation}

\item
\label{modsub:full}
Every object of $\catname{T}$ contained in $S$ is an object of
$\catname{S}$
\begin{equation}
\uquant%
  {U \objin \catname{T}: U \subseteq S}
  {U \objin \catname{S}}
\end{equation}

\begin{remark}
This is actually a consequence of \cref{modsub:rest}, but it is
convenient to state it here.
\end{remark}
\end{enumerate}

$\seqname{S}$ is also a full model subspace of $\seqname{T}$,
abbreviated $\seqname{S} \submod[full-] \seqname{T}$, iff
$S = T$.

$\seqname{S}$ is also a strict model subspace of $\seqname{T}$,
abbreviated $\seqname{S} \submod[strict-] \seqname{T}$, iff
$S$ is a model neighborhood of $\seqname{T}$.
\begin{remark}
$S$ is also a model neighborhood of $\seqname{S}$ by \cref{modsub:full}.
\end{remark}

Let $\seqname{S} \defeq (S, \catname{S})$ be a model space and $T$ a
model neighborhood of $\seqname{S}$. Then $\Mod(T,\seqname{S})$, the
relative model space of $T$, is $(T, \catname{T})$, where $\catname{T}$
is the full subcategory of $\catname{S}$ containing all model
subneighborhoods of $T$.

By abuse of language we write $T$ for $\Mod(T,\seqname{S})$ when
$\seqname{S}$ is understood by context.
\end{definition}

\begin{lemma}[Model subspaces]
\label{lem:modsub}

\begin{enumerate}

\item
$\submod$, $\submod[full-]$ and $\submod[strict-]$ are transitive.

\begin{proof}
Let $\seqname{S}^i \defeq (S^i, \catname{S}^i)$ $i=1,2,3$, be a model
space. If $\seqname{S}^i \submod \seqname{S}^{i+1}$, $i=1,2$, then

\begin{enumerate}
\item
$\seqname{S}^1$ and $\seqname{S}^3$ are model spaces by hypothesis.

\item Since $S^i$ is a subspace of $S^{i+1}$, $i=1,2$, then $S^1$ is a
subspace of $S^3$.

\item
Each $\catname{S}^i \subcat \catname{S}^{i+1}$ by
{
  \showlabelsinline
  \cref{modsub:subcat}
}
of \fullcref{def:modsub} above.
Then $\catname{S}^1 \subcat \catname{S}^3$.

\item
Each intersection of $S^i$ with an object of $\catname{S}^{i+1}$
is an object of $\catname{S}^i$ by
{
  \showlabelsinline
  \cref{modsub:rest}
}
of \cref{def:modsub}. Then every intersection of $S^1$ with an object
of $\catname{S}^3$ is an object of $\catname{S}^1$.

\item
Let $U \objin \seqname{S}^3$ and $U \subseteq S^1$.
$U \cap S^1 = U \objin \seqname{S}^1$.
\end{enumerate}

If each $\seqname{S}^i \submod[full-] \seqname{S}^{i+1}$, then each
$S^i = S^{i+1}$ by \cref{def:modsub}; then $S^1 = S^3$ and thus
$\seqname{S}^1 \submod[full-] \seqname{S}^3$.

If each $\seqname{S}^i \submod[strict-] \seqname{S}^{i+1}$, then each
$S^i$ is a model neighborhood of $\seqname{S}^{i+1}$ by
\cref{def:modsub}, and $S^1$ is a model neighborhood of
$\seqname{S}^3$ by item
{
  \showlabelsinline
  \cref{modsub:full}
}
of \cref{def:modsub}.
\end{proof}

\item
Let $\seqname{S} \defeq (S.\catname{S})$ be a model space and
$T \subseteq \seqname{S}$. Then
$\Mod(T,\seqname{S}) \submod \seqname{S}$

\begin{proof}
Let
$
  \seqname{T} = \bigl ( (T.\topname{T}), \catname{T} \bigr )
  \defeq \Mod(T,\seqname{S})
$.

\begin{enumerate}
\item
$\seqname{S}$ and $\seqname{T}$ are model spaces.

\item
$T$ is a subspace of $S$.

\item
$\catname{T} \subcat \catname{S}$ by construction.

\item
Let $U \objin \seqname{S}$. $U \cap T \objin \seqname{S}$. Since
$U \cap T \subseteq T$ is a model neighborhood of $\seqname{S}$, it is a
model neighborhood of $\Mod(T,\seqname{S})$ by construction.

\item
Let $U \objin \seqname{S}$ be contained in $T$. Then $U$ is a model
neighborhood of $\Mod(T,\seqname{S})$ by construction.
\end{enumerate}
\end{proof}

\item
Let $\seqname{S}^i \defeq (S^i,\catname{S}^i)$, $i=1,2$, and
$\seqname{S}^1 \submod \seqname{S}^2$.
Then $S^1$ is an open subspace of $S^2$.

\begin{proof}
The model neighborhoods of $\seqname{S}^1$ are model neighborhoods of
$\seqname{S}^2$, hence open in $S^2$. Their union, which is open in
$S^2$, is $S^1$ by \cref{mod:cover} of
{
  \showlabelsinline
  \pagecref{def:modsub}
}.
\end{proof}

\item
Let $S^i$, $i=1,2$, be a topological space and $S^1 \subsp[open-] S^2$.
Then $\Id[S^1,S^2]$ is open and continuous.

\begin{proof}
Let $U$ be open in $S^1$.
Since $S^1 \subsp[open-] S^2$, there is a $V$ open in $S^2$ such that
$U = S^1 \cap V$.  Since $V$ and $S^1$ are open in $S^2$, so is
$U = S^1 \cap V$.

Let $V$ be open in $S^2$. Since $S^1 \subsp S^2$,
${\Id[S^1,S^2]}^{-1}[V] = V \cap S^1$ is open in $S^1$.
\end{proof}
\end{enumerate}
\end{lemma}

\begin{corollary}[Model subspaces]
\label{cor:modsub}
Let $\seqname{S}^i \defeq (S^i,\catname{S}^i)$, $i=1,2$,
$\seqname{S}'^2 \defeq (S'^2,\catname{S}'^2)$ be model spaces, and
$\seqname{S}^1 \submod \seqname{S}^2 \submod \seqname{S}'^2$.
\begin{enumerate}

\item
$S^1 \subsp[open-] S'^2$.

\begin{proof}
$S^1 \subsp[open-] S^2 \subsp[open-] S'^2$.
\end{proof}

\item
Let $\funcname{f} \maps S^1 \to S^2$ be continuous. Then
$\funcname{f} \maps S^1 \to S'^2$ is continuous.
\begin{proof}
$S^1 \subsp[open-] S'^2$, Let $V \subseteq S'^2$ be an arbitrary open
set.  Since $\seqname{S}^2 \submod \seqname{S}'^2$ by hypothesis,
$V \cap S^2$ is open in $S^2$.  Since $\funcname{f} \maps S^1 \to S^2$
is continuous by hypothesis,
$\funcname{f}^{-1}[V] = \funcname{f}^{-1}[V \cap S^2]$
is open in $S^1$.
\end{proof}
\end{enumerate}
\end{corollary}

\section{M-nearly commutative diagrams}
\label{sec:M-NCD}
Let $\seqname{M} \defeq (M,\catname{M})$ be a model space and $D$ a tree
with two branches, whose nodes are topological spaces $U_i$ and $V_j$
and whose links are continuous functions
$\funcname{f}_i \maps U_i \to U_{i+1}$ and
$\funcname{f}'_j \maps V_j \to V_{j+1}$ between the spaces:

\begin{align*}
D \defeq \{
  &
    U_0 = V_0 \to^{\funcname{f}_0} U_1,
    \dotsc,
    U_{m-1} \to^{\funcname{f}_{m-1}} U_m,
\\
  &
    U_0 = V_0 \to^{\funcname{f}'_0} V_1,
    \dotsc,
    V_{n-1} \to^{\funcname{f}'_{n-1}} V_n
\}
\end{align*}
with $U_0 = V_0$, $U_m \objin \catname{M}$ and
$V_n \objin \catname{M}$, as shown in \pagecref{fig:NCDb}.

\begin{definition}[M-nearly commutative diagrams]
\label{def:M-NCD}
$D$ is (left,right,strongly) M-nearly commutative in model space
$\seqname{M}$ iff $D$ is (left,right,strongly) nearly commutative in
category $\catname{M}$.
\end{definition}

\begin{definition}[M-nearly commutative diagrams at a point]
\label{def:M-NCDp}
Let $\seqname{M}$ and $D$ be as above and $x$ be an element of the
initial node. $D$ is (left,right,strongly) M-nearly commutative in
$\seqname{M}$ at $x$ iff $D$ is (left,right,strongly) nearly commutative
in category $\catname{M}$ at $x$.
\end{definition}

\begin{definition}[M-locally nearly commutative diagrams]
\label{def:M-NCDl}
Let $\catname{C}$ and $D$ be as above. $D$ is (left,right,strongly)
M-locally nearly commutative in $\seqname{M}$ iff $D$ is
(left,right,strongly) nearly commutative in category $\catname{M}$ at
every point.
\end{definition}

\section{Trivial model spaces}
\label{sec:triv}
Informally, a trivial model space of a specific type is one that does
not restrict the potential objects and morphisms of its type.

\begin{definition}[Trivial model spaces]
\label{def:trivmod}
Let $S$ be a topological space and $\catname{S}$ the category of all
continuous functions between open sets of $S$.

$\Triv{S} \defeq (S, \catname{S})$ is the trivial model space
of $S$ and $\Triv{S}$ is a trivial model space.

Let $\seqname{S}$ be a set of topological spaces.

$
  \Triv{\seqname{S}}
  \defeq
  \set
    {\Triv{S'}}%
    [S' \in \seqname{S}]
$
is the set of all trivial model spaces in $\seqname{S}$.

$\Trivcat{\seqname{S}}$, the category of all continuous functions
between elements of $\Triv{\seqname{S}}$, is the trivial model category
of $\seqname{S}$.

$\optriv{\seqname{S}}$, the category of open trivial model spaces
in $\seqname{S}$, is the category whose objects are
$\set {\Triv{U}}[U \in \op{\seqname{S}}]$, the trivial
model spaces of non-null open sets of spaces in $\seqname{S}$,
and whose morphisms are all the continuous functions among them.
\end{definition}

\begin{lemma}[Trivial model spaces]
\label{lem:trivmod}
Let $S^i$, $i=1,2$, be a topological space.

\begin{enumerate}
\item
$\Triv{S^i}$ is a model space.

\begin{proof}
Let $\catname{S}^i \defeq \Cat \biggl ( \Triv{S^i} \biggr )$,
$\funcname{f} \maps A \to B$ be a morphism of $\catname{S}^i$,
$A' \objin \catname{S}^i \subseteq A$ and
$B' \objin \catname{S}^i \subseteq B$. Then the definitions of
continuity and of $\Triv{S^i}$ imply each of the following.
\begin{enumerate}
\item $\Ob(\catname{S}^i)=\Topology(S^i)$ and thus is an open cover of
$S^i$.
\item $\Ob(\catname{S}^i)=\Topology(S^i)$ and thus is closed under
finite intersections.
\item All morphisms in $\Ar(\catname{S}^i)$ are continuous.
\item Since $\funcname{f} \maps A \to B$ is a morphism of
$\catname{S}^i$, $\funcname{f}$ is continuous.
\item
Since $A' \objin \catname{S}^i$ then $A'$ is open.
\item
Since $B' \objin \catname{S}^i$ then $B'$ is open.
\item
Since $\funcname{f} \maps A \to B$ is continuous then
$\funcname{f} \restrictto_{A'} \maps A' \to B$ is continuous.
\item
Since $\funcname{f}[A'] \subseteq B'$ then
$\funcname{f} \restrictto_{A',B'} \maps A' \to B'$ is well defined and
continuous, hence a morphism.
\item If $A' \objin \catname{S}^i \subseteq A \objin \catname{S}^i$ then
\begin{enumerate}
\item $A$ and $A'$ are open.
\item
The inclusion map $\Id[A',A]$ is continuous.
\item The inclusion map $\Id[A',A]$ is a
morphism of $\catname{S}^i$ by the definition of $\Triv{S^i}$.

\item $\Triv{S^i}$ trivially satisfies the restricted sheaf condition
because every continuous function between open sets of $S^i$ is a
morphism of $\catname{S}^i$.
\end{enumerate}
\end{enumerate}
\end{proof}

\item
Let $S^1$ be an open subspace of $S^2$. Then
$\Triv{S^1} \submod[strict-] \Triv{S^2}$.

\begin{proof}
Let $\catname{S}^i \defeq \Cat \bigl ( \Triv{S^i} \bigr )$, $i=1,2$.
By \pagecref{def:trivmod}, $\catname{S}^i$ is the category of all
continuous functions between open sets of $S^i$,

$\Triv{S^1}$ and $\Triv{S^2}$ satisfy the six conditions of
\pagecref{def:modsub}:

\begin{description}
\item[$\Triv{S^1}$ and $\Triv{S^2}$ are model spaces:]
Proved above.

\item[$S^1$ is a subspace of $S^2$:]
\leavevmode \\*
$\Top \biggl ( \Triv{S^i} \biggr )= S^i$ by \pagecref{def:trivmod}.
$S^1$ is an open subspace of $S^2$ by hypothesis.

\item[\ensuremath{\catname{S}^1 \subcat[full-] \catname{S}^2}:]
\leavevmode \\*
If $U \objin \catname{S}^1$ then $U$ is open in $S^1$. Since $S^1$ is
an open subspace of $S^2$, $U$ is open in $S^1$ and thus
$U \objin \catname{S}^2$.

If $\funcname{f} \maps U \to V \arin \catname{S}^1$ then $\funcname{f}$
is a continuous function between open sets of $S^1$. Since $S^1$ is an
open subspace of $S^2$, $\funcname{f}$ is a continuous function between
open sets of $S^2$ and thus a morphism of $\catname{S}^2$.

If $\funcname{f} \maps U \to V \arin \catname{S}^2$ then $\funcname{f}$
is a continuous function between open sets of $S^2$. If
$U \objin \catname{S}^1$ and $U \objin \catname{S}^1$, then $U$ and
$V$ are open sets of $S^1$. Since $S^1$ is an open subspace of $S^2$,
$\funcname{f}$ is a continuous function between open sets of $S^1$ and
thus a morphism of $\catname{S}^1$.

\item[Every intersection of $S^1$ with an object of $\catname{S}^2$ is an object of $\catname{S}^1$:]
\leavevmode \\*
Let $U \objin \catname{S}^2$. $U$ and $S^1$
are open in $S^2$. $U /cup \Set(S^1)$ is open in $S^2$ and contained in
$S^1$. Since $S^1$ is an open subspace of $S^2$, $U /cup \Set(S^1)$ is
open in $S^1$ and thus an object of $\catname{S}^1$.

\item[Every object of $\catname{S}^2$ contained in $S^1$ is an object of $\catname{S}^1$:]
\leavevmode \\*
Let $U \objin \catname{S}^2$. $U$ is open in $S^2$. If $U$ is
contained in $S^1$ then it is open in $S^1$ and thus an object of
$\catname{S}^1$.
\end{description}

Since $S^1$ is an open subspace of $S^2$, $\Set(S^1)$ is a model
neighborhood of $\Triv{S^2}$ by \cref{def:trivmod} above
\end{proof}

$\Triv{S^1} \submod \Triv{S^2}$.

\begin{proof}
$S^1 \subsp[open-] S^2$ by \pagecref{lem:modsub}.
\end{proof}

\item
Let $\funcname{f} \maps S^1 \to S^2$ be continuous. Then $\funcname{f}$
is a model function from $\Triv{S^1}$ to $\Triv{S^2}$.

\begin{proof}
Let $U^i$ be a model neighborhood of $\Triv{S^i}$, $i=1,2$.
\begin{enumerate}
\item
Since $U^2$ is a model neighborhood of $\Triv{S^2}$, $U^2$ is open,
$\funcname{f}^{-1}[U^2]$ is open and thus a model neighborhood of
$\Triv{S^1}$.
\item
$\funcname{f}[U^1] \subseteq S^2$. Since $S^2$ is open, it is a model
neighborhood of $\Triv{S^2}$.
\end{enumerate}
\end{proof}
\end{enumerate}
\end{lemma}

\section{Minimal model spaces}
\label{sec:min}
Even if a set of open sets fails one of \cref{mod:cover,mod:closed} or a
set of functions among them fails one of
\crefrange{mod:continuous}{mod:sheaf} in the definition of a model
space, there is a minimal model space containing them.

\begin{definition}[Minimal model spaces]
\label{def:minmod}
Let $S$ be a topological space, $\seqname{O}$ a set of open sets in $S$,
$\funcseqname{f}$ a set of continuous functions between elements of
$\seqname{O}$ and $\catname{S}$ the smallest concrete category  over
$\cat{Top}$ having all sets in $\seqname{O}$ as objects, having all
functions in $\funcseqname{f}$ as morphisms and satisfying
\crefrange{mod:continuous}{mod:sheaf} of
{
  \showlabelsinline
  \pagecref{def:model}
}.

Then
\begin{equation}
\minimal{\Mod}
  (
    S,
    \seqname{O},
    \funcseqname{f}
  )
\defeq
  \left (
    \Top \bigl ( \union{O}, S \bigr ),
    \catname{S}
  \right )
\end{equation}
is the minimal model space of $S$ with neighborhoods $\seqname{O}$ and
neighborhood mappings $\funcseqname{f}$.
\begin{remark}
The trivial model space $\Triv{S}$ is a special case.
\end{remark}

Let $\seqname{S} \defeq (S, \catname{S})$ be a model space,
$\seqname{O}$ a set of model neighborhoods of $\seqname{S}$ and
$\funcseqname{f}$ a set of morphisms between elements of $\seqname{O}$.

Then
\begin{equation}
\minimal{\Mod}
  (
    \seqname{S},
    \seqname{O},
    \funcseqname{f}
  )
\defeq
\minimal{\Mod}
  (
    S,
    \seqname{O},
    \funcseqname{f}
  )
\end{equation}
is the minimal model space of $\seqname{S}$ with neighborhoods
$\seqname{O}$ and neighborhood mappings $\funcseqname{f}$.

\begin{remark}
$
  \minimal{\Mod}
    (
      \seqname{S},
      \seqname{O},
      \funcseqname{f}
    )
$
is not, in general, a model subspace of $\seqname{S}$, since
{
  \showlabelsinline
  \cref{modsub:rest,modsub:full}
}
of \pagecref{def:modsub} may fail.

While every model neighborhood of
$\minimal{\Mod}(\seqname{S}, \seqname{O}, \funcseqname{f})$
is a model neighborhood of
$\minimal{\Mod}(S, \seqname{O}, \funcseqname{f})$, the converse need not
be true.
\end{remark}
\end{definition}

\begin{lemma}[Minimal model spaces are model spaces]
\label{lem:minmod}
Let $S$ be a topological space, $\seqname{O}$ a set of open sets in $S$
and $\funcseqname{f}$ a set of continuous functions between elements of
$\seqname{O}$. Then
$(C,\catname{C}) \defeq \minimal{\Mod}(S, \seqname{O}, \funcseqname{f})$
is a model space.

\begin{proof}
\leavevmode
\begin{enumerate}
\item Finite intersections of open sets are open and
$C = \union[U \in \seqname{O}]{U}$ by construction
\item $\Ob(\catname{C})$ is closed under finite intersections by
construction
\item Compositions of continuous functions, inclusion maps and
restrictions of continuous functions are continuous
\item Restrictions of morphisms are morphisms by construction
\item Inclusion maps are morphisms by construction
\item The restricted sheaf condition holds by construction
\end{enumerate}
\end{proof}

Let $\seqname{S} \defeq (S, \catname{S})$ be a model space,
$\seqname{O}$ a set of model neighborhoods of $\seqname{S}$ and
$\funcseqname{f}$ a set of morphisms between elements of $\seqname{O}$.
Then
$\minimal{\Mod}(\seqname{S}, \seqname{O}, \funcseqname{f})$
is a model space.

\begin{proof}
$
  \minimal{\Mod}(\seqname{S}, \seqname{O}, \funcseqname{f}) =
  \minimal{\Mod}(S, \seqname{O}, \funcseqname{f})
$
and $\minimal{\Mod}(S, \seqname{O}, \funcseqname{f})$ is a model space.
\end{proof}
\end{lemma}

\section{M-paracompact model spaces}
\label{sec:m-para}
Paracompactness is an important property for topological spaces because
of partitions of unity. There is an analogous property for model
spaces.

\begin{definition}[Model topology]
\label{def:ModTop}
Let $\seqname{M} \defeq (S, \catname{S})$ be a model space for $S$.
Then the model topology $\topname{M}^*$ for $M$ is the topology
generated by $\Ob(\catname{S})$.
\begin{remark}
$\topname{M}^*$ is not guarantied to be T0 even if $S$ is T4. However,
$\topname{M}^*$ may be normal or regular even if $S$ is not.
\end{remark}
\end{definition}

\begin{definition}[m-paracompactness]
\label{def:M-para}
A model space $\seqname{M} \defeq (S, \catname{S})$ is \\
m-paracompact iff $\topname{M}^*$ is regular and every cover of $S$ by
model neighborhoods has a locally finite refinement by model neighborhoods.
\begin{remark}
This is a stronger condition than merely requiring $\topname{M}^*$ to
be paracompact.
\end{remark}
\end{definition}

\begin{theorem}[m-paracompactness and paracompactness]
If $\seqname{M} = (S, \catname{S})$ is m-paracompact and the model
neighborhoods form a basis for $S$ then $S$ is paracompact.
\begin{proof}
Let $R$ be an open cover of $S$.  Since the model neighborhoods form a
basis, every set in $R$ is a union of model neighborhoods and thus there
is a refinement $R_1$ by model neighborhoods.  Since by hypothesis
$\seqname{M}$ is m-paracompact, $R_1$ has a locally finite refinement
$R_2$ by model neighborhoods.  Since model neighborhoods are open sets,
$R_2$ is a locally finite refinement of $R$ in the conventional sense.
\end{proof}
\end{theorem}

\section{Model functions and model categories}
\label{sec:modcat}
It is convenient to have a notion of mappings between model spaces that
are well behaved in some sense, e.g., fiber preserving; using that
notion it is then possible to group model spaces into categories.

\subsection{Model functions}
\begin{definition}[Model functions]
\label{def:modfunc}
Let $\seqname{S}^i \defeq (S^i, \catname{S}^i)$, $i=1,2$, be a model
space and $\funcname{f} \maps S^1 \to S^2$ be a continuous function.
$\funcname{f}$ is a model function of $\seqname{S}^1$ to $\seqname{S}^2$
iff the inverse images of model neighborhoods are model neighborhoods.

\begin{equation}
\uquant%
  {V \objin \catname{S}^2}
  {\funcname{f}^{-1}[V] \objin \catname{S}^1}
\end{equation}

$\funcname{f}$ is also a constrained model function of $\seqname{S}^1$
to $\seqname{S}^2$ iff $\funcname{f}[S^1]$ is contained in a model
neighborhood of $\seqname{S}^2$.

\begin{equation}
\equant%
  {V \objin \catname{S}^2}
  {\funcname{f}[S^1] \subseteq V}
\end{equation}

\begin{remark}
When $\funcname{f} \maps \seqname{S}^1 \to \seqname{S}^2$ is
constrained, $S^1$ is a model neighborhood of $\seqname{S}^1$.
\end{remark}

By abuse of language we write
$\funcname{f} \maps \seqname{S}^1 \to \seqname{S}^2$ both for
$\funcname{f}$ considered as a model function and for $\funcname{f}$
considered as a continuous function.

Let $\seqname{S}^1 \defeq (S^1, \catname{S}^1)$ and
$\seqname{S}'^2 \defeq (S'^2, \catname{S}'^2)$ be model spaces,
$\seqname{S}^2 \defeq (S^2, \catname{S}^2) \submod \seqname{S}'^2$ be a
model subspace and $\funcname{f} \maps S^1 \to S^2$ be a model function
of $\seqname{S}^1$ to $\seqname{S}^2$. Then
$
  \funcname{f} \maps \seqname{S}^1 \to \seqname{S}'^2 \defeq
  \Id[{\seqname{S}^2},{\seqname{S}'^2}] \compose \funcname{f}
$
is $\funcname{f}$ considered as a model function from $\seqname{S}^1$ to
$\seqname{S}'^2$.

Let $\seqname{S}'^i \defeq (S'^i, \catname{S}'^i)$, $i=1,2$, be a model
space,
$\seqname{S}^i \defeq (S^i, \catname{S}^i) \submod \seqname{S}'^i$ be a
model subspace and $\funcname{f}'$ be a model function of
$\seqname{S}'^1$ to $\seqname{S}'^2$ such that
$\funcname{f}'[S^1] \subseteq S^2$. Then
$\funcname{f} \maps \seqname{S}^1 \to \seqname{S}^2$, also written
$\funcname{f}' \restrictto_{\seqname{S}^1,\seqname{S}^2}$, is
$\funcname{f}' \restrictto_{\seqname{S}^1}$ considered as a model
function of $\seqname{S}^1$ to $\seqname{S}^2$.
\end{definition}

\begin{lemma}[Model functions]
\label{lem:modfunc}
Let $\seqname{S}^i \defeq (S^i, \catname{S}^i)$ and
$\seqname{S}'^i \defeq (S'^i, \catname{S}'^i)$, $i=1,2$, be model
spaces, $\seqname{S}^i \submod \seqname{S}'^i$ and
$\funcname{f}' \maps S'^1 \to S'^2$ be a model function of
$\seqname{S}'^1$ to $\seqname{S}'^2$.

$\funcname{f} \defeq \funcname{f}' \restrictto_{S^1,S^2}$ is a model
function of $\seqname{S}^1$ to $\seqname{S}^2$.

\begin{proof}
$\funcname{f}'$ is a model function by hypothesis.
$\funcname{f}'$ is continuous by \cref{def:modfunc} above, so
$\funcname{f} \defeq \funcname{f}' \restrictto_{S^1,S^2}$ is continuous.

Let $U$ be a model neighborhood of $\seqname{S}^2$.
$\seqname{S}^2 \submod \seqname{S}'^2$ by hypothesis, so
$U$ is a model neighborhood of $\seqname{S}'^2$.
$\funcname{f}'$ is a model function of $\seqname{S}'^1$ to
$\seqname{S}'^2$ by hypothesis, so
$\funcname{f}'^{-1}[U] \objin \seqname{S}'^1$ by \cref{def:modfunc}.
Then
$
  \funcname{f}^{-1}[U] = \funcname{f}'^{-1}[U] \cap S^1 \objin
  \seqname{S}^1
$
by \cref{def:modfunc}.
\end{proof}
\end{lemma}

\begin{lemma}[Model functions and trivial model spaces]
\label{lem:modfunctriv}
Let $S^i$, $i=1,2$, be a topological space and
$\funcname{f} \maps S^1 \to S^2$ be a continuous function.
Then $\funcname{f} \maps \Triv{S^1} \to \Triv{S^2}$ is a model function.

\begin{proof}
Let $U^i$ be a model neighborhood of $\Triv{S^i}$, $i=1,2$.

$\funcname{f}^{-1}[U^2]$ is open, hence a model neighborhood of
$\Triv{S^1}$.

$S^2$ is a model neighborhood of $\Triv{S^2}$ and
$\funcname{f}[U^1] \subseteq U^2$.
\end{proof}

Let $\seqname{S}^i \defeq (S^i,\catname{S}^i)$, $i=1,2$, be a model
space and $\funcname{f} \maps \seqname{S}^1 \to \seqname{S}^2$ be a
model function. Then $\funcname{f} \maps \Triv{S^1} \to \Triv{S^2}$ is a
model function.

\begin{proof}
$\funcname{f} \maps S^1 \to S^2$ is continuous.
\end{proof}

Let $\seqname{S}^i \defeq (S^i,\catname{S}^i)$, $i=1,2$, be a model
space and $\funcname{f} \maps \seqname{S}^1 \to \seqname{S}^2$ be a
model function. Then $\funcname{f} \maps \Triv{S^1} \to \seqname{S}^2$
is a model function iff
$\funcname{f} \maps \seqname{S}^1 \to \seqname{S}^2$  is a constrained
model function.

\begin{proof}
let $U^2$ be a model neighborhood of $\seqname{S}^2$.
$\funcname{f}^{-1}[U^2]$ is open, hence a model neighborhood of
$\Triv{S^1}$.

$S^1$ is a model neighborhood of $\Triv{S^1}$. If $\funcname{f} \maps
\Triv{S^1} \to \seqname{S}^2$ is a model function then
$\funcname{f}[S^1]$ is contained in a model neighbothood of
$\seqname{S}^2$ and thus $\funcname{f}$ is constrained.

If $\funcname{f}$ is constrained, let $U^1$ be a model neighborhood of
$\Triv{S^1}$. Then $\funcname{f}[U^1] \subseteq \funcname{f}[S^1]$ is
contained in a model neighbothood of $\seqname{S}^2$.
\end{proof}
\end{lemma}

\begin{lemma}[Composition of model functions]
\label{lem:modcomp}
Let $\seqname{S}_i \defeq (S_i,\catname{S}_i)$, $i \in [1,3]$, be a
model space and
$\funcname{f}_i \maps \seqname{S}_i \to \seqname{S}_{i+1}$, $i=1,2$, be
a model function.

$\funcname{f}_2 \compose \funcname{f}_1$ is a model function.

\begin{proof}
Let $U_i \objin \catname{S}_i,$ $i=1,2,3$.

Since $\funcname{f}_2$ is a model function,
\nobreaks
  {{
    $\funcname{f}^{-1}_2[U_3] \objin \catname{S}_2$.
  }}
Since $\funcname{f}_1$ is a model function,
$
  (\funcname{f}_2 \compose \funcname{f}_1)^{-1}[U_3]
  =
  \funcname{f}^{-1}_1[\funcname{f}^{-1}_2[U_3]]
  \objin
  \catname{S}_1
$.
Since $\funcname{f}_1$ is a model functions, $\funcname{f}_1[U_1]$
is contained in a model neighborhood $V_2$.
Since $\funcname{f}_2$ is a model functions, $\funcname{f}_2[V_2]$
is contained in a model neighborhood $V_3$.
Then $(\funcname{f}_2 \compose \funcname{f}_1)[U_1] \subseteq V_3$.
\end{proof}

If each $\funcname{f}_i$ is constrained then
$\funcname{f}_2 \compose \funcname{f}_1$ is constrained.

\begin{proof}
$\funcname{f}_1[S^1] \subseteq S^2$ and
$\funcname{f}_2[S^2] \subseteq S^3$, hence
$\funcname{f}_2 \compose \funcname{f}_1[S^1] \subseteq S^3$.
\end{proof}
\end{lemma}

\begin{definition}[Model homeomorphisms]
\label{def:modelhomeo}
Let $\seqname{S}^i \defeq (S^i, \catname{S}^i)$, $i=1,2$, be a model
space and $\funcname{f} \maps S^1 \to S^2$ be a model function.
$\funcname{f}$ is a model homeomorphism iff it is also invertible and
its inverse is a model function.
\end{definition}

\subsection{Model categories}
\begin{definition}[Model categories]
\label{def:modcat}
A category $\catname{M}$ is a model category iff

\begin{enumerate}
\item the objects of $\catname{M}$ are model spaces

\item the morphisms of $\catname{M}$ are model functions.

\item composition is functional composition.

\item
\label{modcat:inc}
If $\seqname{S} \defeq (S, \catname{S}) \objin \catname{M}$ and
$S^i \objin \catname{S}$, $i=1,2$, then
$\seqname{S}^i \defeq \Mod(S^i,\seqname{S}) \objin \catname{M}$.
If further $S^1 \subseteq S^2$ then the inclusion map
$\Id[S^1,S^2] \maps \seqname{S}^1 \into \seqname{S}^2$ is a morphism
of $\catname{S}$.

\item
\label{modcat:rest}
If $\seqname{S}^i \defeq (S^i, \catname{S}^i) \objin \catname{M}$,
$U^i \objin \catname{S}^i$,
$i=1,2$, $\funcname{f} \maps \seqname{S}^1 \to \seqname{S}^2$ is a
morphism of $\catname{M}$ and $\funcname{f}[U^1] \subseteq U^2$ then
$\funcname{f} \maps \Mod(U^1,\seqname{S}^1) \to \Mod(U^2,\seqname{S}^2)$
is a morphism of $\catname{M}$.

\item
\label{modcat:morph}
If $\seqname{S}^i \defeq ({S}^i, \catname{S}^i)$, $i=1,2$, is a
subspace of $\seqname{S} \defeq (S, \catname{S}) \objin \catname{M}$
and $\funcname{f} \maps S^1 \to S^2$ is a morphism of $\catname{S}$ then
$\funcname{f} \maps \seqname{S}^1 \to \seqname{S}^2$ is a morphism of
$\catname{M}$.
\end{enumerate}

A model category $\catname{M}$ satisfies the restricted
sheaf condition iff whenever
\begin{enumerate}
\item $U$ and $V$ are objects of $\catname{S}$.

\item $U_\alpha \submod U$ and $V_\alpha \submod V$,
$\alpha \prec \Alpha$, are objects of $\catname{M}$.

\item
$\Set(U) = \union[\alpha \prec \Alpha]{\Set(U_\alpha)}$

\item
$\Set(V) = \union[\alpha \prec \Alpha]{\Set(V_\alpha)}$

\item $\funcname{f}_\alpha \maps U_\alpha \to V_\alpha$ are morphisms of
$\catname{M}$.

\item
$\funcname{f} \maps U \to V$ is a model function and
for every $\alpha \prec \Alpha$,
$\funcname{f}$ agrees with $\funcname{f}_\alpha$ on $U_\alpha$
\end{enumerate}
then $\funcname{f}$ is a morphism of $\catname{M}$.

Let $\catname{M}^i$, $i=1,2$, be a model category. $\catname{M}^1$ is a
strict subcategory of $\catname{M}^2$, abbreviated
$\catname{M}^1 \subcat[strict-] \catname{M}^2$, iff
$\catname{M}^1 \subcat[fullt-] \catname{M}^2$ and $\catname{M}^2$
satisfies the restricted sheaf condition.
\end{definition}

\begin{lemma}[Model categories]
\label{lem:modcat}
Let $\catname{M}$ be a model category,
$\seqname{S}^i \defeq (S^i, \catname{S}^i) \objin \catname{M}$. \\*
$i=1,2$,
$U^i \objin \catname{S}^i$,
$\seqname{U}^i \defeq \Mod(U^i,\seqname{S}^i)$ and
$\funcname{f} \maps \seqname{U}^1 \to \seqname{U}^2 \arin \catname{M}$.

$\funcname{f} \maps \seqname{U}^1 \to \seqname{S}^2 \arin \catname{M}$.

\begin{proof}
Since $\Id[\seqname{U}^2,\seqname{S}^2] \arin \catname{M}$ by
\cref{modcat:inc} of
{
  \showlabelsinline
  \cref{def:modcat}
}
above and
$\funcname{f} \maps \seqname{U}^1 \to \seqname{U}^2 \arin \catname{M}$
by hypothesis,
$
  \funcname{f} \maps \seqname{U}^1 \to \seqname{S}^2 =
    \Id[\seqname{U}^2,\seqname{S}^2] \compose
    \funcname{f} \maps \seqname{U}^1 \to \seqname{U}^2
  \arin \catname{M}
$
\end{proof}
\end{lemma}

\begin{definition}[Trivial model categories]
\label{def:modtrivcat}
Let $\seqname{S} \defeq (S, \catname{S})$ be a model space.

$
  \Triv[submod-]{\seqname{S}}
  \defeq
  \set
    {\seqname{M}}%
    [\seqname{M} \submod \seqname{S}]
$
is the set of all subspaces of $\seqname{S}$.

$\Trivcat[submod-]{\seqname{S}}$, the category of all model functions
between elements of $\Triv[submod-]{\seqname{S}}$, is the trivial
subspace model category of $\seqname{S}$.

$
  \Triv[full-submod-]{\seqname{S}}
  \defeq
  \set
    {\seqname{M}}%
    [\seqname{M} \submod[full-] \seqname{S}]
$
is the set of all full subspaces of $\seqname{S}$.

$\Trivcat[full-submod-]{\seqname{S}}$, the category of all full model
functions between elements of $\Triv[full-submod-]{\seqname{S}}$, is the
trivial full subspace model category of $\seqname{S}$.

Let $\seqname{S}$ be a set of model spaces.

$
  \Triv[submod-]{\seqname{S}}
  \defeq
  \set
    {\seqname{M}}%
    [
      \equant
        {\seqname{M}' \in \seqname{S}}
        {\seqname{M} \submod \seqname{M}'}
    ]
$
is the set of all model subspaces of model spaces in $\seqname{S}$.

$\Trivcat[submod-]{\seqname{S}}$, the category of all model functions
between elements of $\Triv[submod-]{\seqname{S}}$, is the trivial
subspace model category of $\seqname{S}$.

$
  \Triv[full-submod-]{\seqname{S}}
  \defeq
  \set
    {\seqname{M}}%
    [
      \equant
        {\seqname{M}' \in \seqname{S}}
        {\seqname{M} \submod[full-] \seqname{M}'}
    ]
$
is the set of all full subspaces in $\seqname{S}$.

$\Trivcat[full-submod-]{\seqname{S}}$, the category of all full model
functions between elements of $\Triv[full-submod-]{\seqname{S}}$, is the
trivial full subspace model category of $\seqname{S}$.
\end{definition}

\begin{lemma}[Trivial model categories]
\label{lem:modtrivcat}

Let $\seqname{S} \defeq (S, \catname{S})$ be a model space. \\*
\nobreaks
  {{
    $
      \Trivcat[full-submod-]{\seqname{S}} \subcat[full-]
      \Trivcat[submod-]{\seqname{S}}
    $.
  }}

\begin{proof}
$
  \uquant
    {\seqname{M} \submod[full-] \seqname{S}}
    {\seqname{M} \submod \seqname{S}}
$,
so
$
  \Ob \bigl (  \Trivcat[full-submod-]{\seqname{S}} \bigr ) \subseteq
  \Ob \bigl ( \Trivcat[submod-]{\seqname{S}} \bigr )
$.

It remains to show that for a pair of objects in both categories, the $\Hom$ sets are the same.
Let $\seqname{M}^i \submod[full-] \seqname{S}$, $i=1.2$.
\begin{enumerate}
\item
Let
$
  \funcname{f} \maps \seqname{M}^1 \to \seqname{M}^2 \arin
  \Trivcat[full-submod-]{\seqname{S}}
$.
Then $\funcname{f}$ is a model function and thus
$
  \funcname{f} \maps \seqname{M}^1 \to \seqname{M}^2 \arin
  \Triv[submod-]{\seqname{S}}
$.
\item
Let
$
  \funcname{f} \maps \seqname{M}^1 \to \seqname{M}^2 \arin
  \Triv[submod-]{\seqname{S}}
$.
Then $\funcname{f}$ is a model function and thus \\*
$
  \funcname{f} \maps \seqname{M}^1 \to \seqname{M}^2 \arin
  \Trivcat[full-submod-]{\seqname{S}}
$.
\end{enumerate}

\end{proof}
\end{lemma}

\subsection{(Local) m-morphisms}

\begin{definition}[Local m-morphisms]
\label{def:modlocalMmorph}
Let $\catname{M}^i$ be a model category, $i=1,2$, and
$\seqname{S}^i \defeq (S^i, \catname{S}^i) \objin \catname{M}^i$.

Let $U^i \subseteq \seqname{S}^i$, $i=1,2$. be open. A
function $\funcname{f} \maps U^1 \to U^2$ is a local
$\seqname{S}^1$-$\seqname{S}^2$ m-morphism of $U^1$ to $U^2$ iff
$\seqname{S}^1 \submod \seqname{S}^2$ and for every $u \in U^1$ there is
a model neighborhood $U_u \subseteq U^1$ for $u$ and a model
neighborhood $V_u \subseteq U^2$ for $v \defeq \funcname{f}(u)$ such
that $\funcname{f}[U_u] \subseteq V_u$ and
$\funcname{f} \maps U_u \to V_u$ is a morphism of $\seqname{S}^2$.

Let $U^1 \subseteq \seqname{S}^i$ and $U^2 \subseteq \seqname{S}^i$
be open.  A function
$\funcname{f} \maps U^1 \to U^2$ is a local $\seqname{S}^i$
m-morphism of $U^1$ to $U^2$ iff it is a local
$\seqname{S}^i$-$\seqname{S}^i$ m-morphism of $U^1$ to $U^2$.

A model function $\funcname{f} \maps \seqname{S}^1 \to \seqname{S}^2$ is
a local m-morphism of $\seqname{S}^1$ to $\seqname{S}^2$ iff
it is a local $\seqname{S}^1$-$\seqname{S}^2$ m-morphism of
$S^1$ to $S^2$.\footnote{
  This definition does not require that $S^i$ be a model neighborhod of
  $S^i$.
}
\newcounter{fn:nonmodel}
\setcounter{fn:nonmodel} {\value{footnote}}

A model function $\funcname{f} \maps \seqname{S}^i \to \seqname{S}^i$ is
a local m-morphism of $\seqname{S}^i$ iff it is a
local m-morphism of $\seqname{S}^i$ to $\seqname{S}^i$.

Let $U^i \objin \catname{S}^i$, $i=1,2$. A function
$\funcname{f} \maps U^1 \to U^2$ is a local
$\seqname{S}^1$-$\seqname{S}^2$-$\catname{M}^1$-$\catname{M}^2$
m-morphism of $U^1$ to $U^2$ iff
$\catname{M}^1 \subcat[full-] \catname{M}^2$ and for every $u \in U^1$
there is a model neighborhood $U_u \subseteq U^1$ for $u$ and a model
neighborhood $V_u \subseteq U^2$ for $v \defeq \funcname{f}(u)$ such
that $\funcname{f}[U_u] \subseteq V_u$ and
$
  \funcname{f} \maps \Mod(U_u,\seqname{S}^1)
  \to \Mod(V_u,\seqname{S}^2)
$
is a morphism of $\catname{M}^2$.

Let $U^1 \subseteq \seqname{S}^i$ and $U^2 \subseteq \seqname{S}^i$, be
open. A model function $\funcname{f} \maps U^1 \to U^2$ is a local
$\seqname{S}^i$-$\catname{M}^i$ m-morphism of $U^1$ to $U^2$ iff it is a
local $\seqname{S}^i$-$\seqname{S}^i$-$\catname{M}^i$-$\catname{M}^i$
m-morphism of $U^1$ to $U^2$.

A model function $\funcname{f} \maps \seqname{S}^1 \to \seqname{S}^2$ is
a local $\catname{M}^1$-$\catname{M}^2$ m-morphism of $\seqname{S}^1$ to
$\seqname{S}^2$ iff it is a local
$\seqname{S}^1$-$\seqname{S}^2$-$\catname{M}^1$-$\catname{M}^2$
m-morphism of $S^1$ to $S^2$.\footnotemark[\value{fn:nonmodel}]

A model function $\funcname{f} \maps \seqname{S}^i \to \seqname{S}^i$ is
a local $\catname{M}^i$ m-morphism iff it is a local
$\catname{M}^i$-$\catname{M}^i$ m-morphism of $\seqname{S}^i$ to
$\seqname{S}^i$.

The phrases {\textquotedblleft}locally a(n) ...
m-morphism{\textquotedblright} and {\textquotedblleft}a local ...
m-morphism{\textquotedblright} are equivalent.
\end{definition}

\begin{definition}[M-morphisms]
\label{def:modMmorph}
Let $\catname{M}^i$ be a model category, $i=1,2$, and \\*
$\seqname{S}^i \defeq (S^i, \catname{S}^i) \objin \catname{M}^i$.

Let $U^i \subseteq \seqname{S}^i$, $i=1,2$. be open. A model
function $\funcname{f} \maps U^1 \to U^2$ is an
$\seqname{S}^1$-$\seqname{S}^2$ m-morphism of $U^1$ to $U^2$ iff
$\seqname{S}^1 \submod \seqname{S}^2$ and $\funcname{f}$ is a morphism
of $\seqname{S}^2$.

A model function $\funcname{f} \maps \seqname{S}^1 \to \seqname{S}^2$ is
an m-morphism of $\seqname{S}^1$ to $\seqname{S}^2$ iff it is an
$\seqname{S}^1$-$\seqname{S}^2$ m-morphism of $S^1$ to
$S^2$.\footnotemark[\value{fn:nonmodel}]

A model function $\funcname{f} \maps \seqname{S}^i \to \seqname{S}^i$ is
an m-morphism of $\seqname{S}^i$ iff it is an m-morphism
of $\seqname{S}^i$ to $\seqname{S}^i$.

Let $U^i \subseteq \seqname{S}^i$, $i=1,2$. be model neighborhoods. A
model function $\funcname{f} \maps U^1 \to U^2$ is an
$\seqname{S}^1$-$\seqname{S}^2$-$\catname{M}^1$-$\catname{M}^2$
m-morphism of $U^1$ to $U^2$ iff
$\catname{M}^1 \subcat[full-] \catname{M}^2$ and \\*
$
  \funcname{f} \maps \Mod(U^1,\seqname{S}^1)
  \to \Mod(U^2,\seqname{S}^2)
$
is a morphism of $\catname{M}^2$.

Let $U^1 \subseteq \seqname{S}^i$ and $U^2 \subseteq \seqname{S}^i$,
$i=1,2$, be model neighborhoods. A model function
$\funcname{f} \maps U^1 \to U^2$ is an $\seqname{S}^i$-$\catname{M}^i$
m-morphism of $U^1$ to $U^2$ iff it is an
$\seqname{S}^i$-$\seqname{S}^i$-$\catname{M}^i$-$\catname{M}^i$
m-morphism of $U^1$ to $U^2$.

A model function $\funcname{f} \maps \seqname{S}^1 \to \seqname{S}^2$ is
an $\catname{M}^1$-$\catname{M}^2$ m-morphism of $\seqname{S}^1$ to
$\seqname{S}^2$ iff
\nobreaks{{$\catname{M}^1 \subcat[full-] \catname{M}^2$}} and
$\funcname{f}$ is a morphism of
$\catname{M}^2$.\footnotemark[\value{fn:nonmodel}]

A model function $\funcname{f} \maps \seqname{S}^i \to \seqname{S}^i$ is
an $\catname{M}^i$ m-morphism of $\seqname{S}^i$ iff it is an
$\catname{M}^i$-$\catname{M}^i$ m-morphism of $\seqname{S}^i$ to
$\seqname{S}^i$.
\end{definition}

\begin{lemma}[(Local) m-morphisms]
\label{lem:modMmorph}
Let $\catname{M}'^i$, $\catname{M}^i \subcat[full-] \catname{M}'^i$,
$i=1,2$, be model categories,
$\seqname{S}^i \defeq (S^i, \catname{S}^i) \objin \catname{M}^i$,
$\seqname{S}'^i \defeq (S'^i, \catname{S}'^i) \objin \catname{M}'^i$
and $\seqname{S}^i \submod \seqname{S}'^i$.

\begin{enumerate}
\item
\label{M-morph:openmorphislocal}
Let $U^i \subseteq \seqname{S}^i$, $i=1,2$, be open.  If
$\funcname{f} \maps U^1 \to U^2$ is a
$\seqname{S}^1$-$\seqname{S}^2$ m-morphism of $U^1$ to $U^2$ then
$\funcname{f} \maps U^1 \to U^2$ is a local
$\seqname{S}^1$-$\seqname{S}^2$ m-morphism of $U^1$ to $U^2$.

\begin{proof}
\leavevmode
\begin{enumerate}
\item
Since $\funcname{f}$ is an $\seqname{S}^1$-$\seqname{S}^2$ m-morphism of
$U^1$ to $U^2$, $\seqname{S}^1 \submod \seqname{S}^2$ and
$\funcname{f}$ is a morphism of $\seqname{S}^2$,

\item
$U^1$ and $U^2$ are model neighborhoods of $\seqname{S}^2$. $U^1$ is a
model neighborhood of $\seqname{S}^1$ by \cref{modsub:full} of
{
  \showlabelsinline
  \pagecref{def:modsub}.
}

\item
Let $u^1 \in U^1$, $u^2 \defeq \funcname{f}(u^1) \in U^2$, Then $U^1$ is
a model neighborhood for $u^1$ and $U^2$ is a model neighborhood for
$u^2$.  $\funcname{f}[U^1] \subseteq U^2$ by hypothesis.
$\funcname{f} \maps U^1 \to U^2$ is a morphism of $\seqname{S}^2$ as
shown above.
\end{enumerate}
\end{proof}

\item
\label{M-morph:openlocalismorph}
Let $U^i \subseteq \seqname{S}^i$ be a model neighborhood of
$\seqname{S}^i$, $i=1,2$. If $\funcname{f} \maps U^1 \to U^2$ is a
$\seqname{S}^1$-$\seqname{S}^2$ local m-morphism of $U^1$ to
$U^2$ then $\funcname{f} \maps U^1 \to U^2$ is a
$\seqname{S}^1$-$\seqname{S}^2$ m-morphism of $U^1$ to $U^2$.

\begin{proof}
$\seqname{S}^1 \submod \seqname{S}^2$ by hypothesis.
$U^1$ is a model neighborhood of $\seqname{S}^1$ by hypothesis, thus a
model neighborhood of $\seqname{S}^2$
\item
For each $u \in U^1$, let $U_u$ be a model subneighborhood of $U^1$  for
$u$ and $V_u$ be a model subneighborhood of $U^2$ for
$v \defeq \funcname{f}(u)$ such that $\funcname{f}[U_u] \subseteq V_u$
and $\funcname{f} \maps U_u \to  V_u$ is a morphism of $\seqname{S}^2$.
Then $\funcname{f} \maps S^1 \to S^2$ is a morphism of $\seqname{S}^2$
by
{
  \showlabelsinline
  \cref{mod:sheaf}
}
of \pagecref{def:model}.
\end{proof}

\item
\label{M-morph:opencatmorphislocal}
Let $U^i \subseteq \seqname{S}^i$ be a model neighborhood, $i=1,2$, If
$\funcname{f} \maps U^1 \to U^2$ is a
$\seqname{S}^1$-$\seqname{S}^2$-$\catname{M}^1$-$\catname{M}^2$
m-morphism of $U^1$ to $U^2$ then $\funcname{f} \maps U^1 \to U^2$ is a
local
$\seqname{S}^1$-$\seqname{S}^2$-$\catname{M}^1$-$\catname{M}^2$
m-morphism of $U^1$ to $U^2$.

\begin{proof}
\leavevmode
\begin{enumerate}
\item
Each $U^i$ is a model neighborhood of $\seqname{S}^i$ by hypothesis.

\item
Since $\funcname{f} \maps U^1 \to U^2$ is an
$\seqname{S}^1$-$\seqname{S}^2$-$\catname{M}^1$-$\catname{M}^2$
m-morphism of $U^1$ to $U^2$, \\*
\nobreaks
  {
    $\catname{M}^1 \subcat[full-] \catname{M}^2$
  }
and
$
  \funcname{f} \maps \Mod(U^1,\seqname{S}^1)
  \to \Mod(U^2,\seqname{S}^2)
$
is a morphism of $\catname{M}^2$ by \pagecref{def:modMmorph}.

\item
Let $u^1 \in U^1$, $u^2 \defeq \funcname{f}(u^1) \in U^2$.

\begin{enumerate}
\item
Each $U^i$ is a model neighborhood of $\seqname{S}^i$ for $u^i$,

\item
Each $U^i \subseteq U^i$.

\item
$\funcname{f}[U^1] \subseteq U^2$.
\end{enumerate}
\end{enumerate}
\end{proof}

\item
\label{M-morph:opencatlocalismorph}
Let $\catname{M}^2$ satisfy the restricted sheaf condition and each
$U^i \subseteq \seqname{S}^i$ be a model neighborhood, $i=1,2$, Every
$\funcname{f} \maps U^1 \to U^2$ that is a
$\seqname{S}^1$-$\seqname{S}^2$-$\catname{M}^1$-$\catname{M}^2$ local
m-morphism of $U^1$ to $U^2$ is a
$\seqname{S}^1$-$\seqname{S}^2$-$\catname{M}^1$-$\catname{M}^2$
m-morphism of $U^1$ to $U^2$.

\begin{proof}
\leavevmode
\begin{enumerate}
\item
Since each $U^i$ is a model neighborhood of $\seqname{S}^i$ by hypothesis,
each \\*
$\Mod(U^i,\seqname{S}^i) \objin \catname{M}^i$ by \cref{modcat:inc}
of
{
  \showlabelsinline
  \cref{def:modcat}
}

\item
$
  \funcname{f} \maps
  \Mod(U^1,\seqname{S}^1) \to
  \Mod(U^2,\seqname{S}^2)
$
is a model function by \cref{def:modlocalMmorph}

\item
Each $\seqname{S}^i \objin \catname{M}^i$ by hypothesis

\item
$\catname{M}^1 \subcat \catname{M}^2$ by hypothesis

\item
$\catname{M}^2$ satisfies the restricted sheaf condition by hypothesis

\item
For every $u^1 \in U^1$ there is a model neighborhood $U^1_{u^1}$ for
$u^1$ and a model neighborhood $U^2_{u^1}$ for
$u^2 \defeq \funcname{f}(u^1)$ such that
$\funcname{f}[U^1_{u^1}] \subseteq U^2_{u^1}$ and
$
  \funcname{f} \maps
  \Mod(U^1_{u^1},\seqname{S}^i) \to
  \Mod(U^2_{u^1},\seqname{S}^i)
$
is a morphism of $\catname{M}^i$.

\begin{enumerate}
\item
$\Mod(U^1_{u^1},\seqname{S}^1) \submod \seqname{S}^1$.

\item
Each
$
  \funcname{f} \maps
  \Mod(U^1_{u^1},\seqname{S}^1) \to
  \Mod(U^2_{u^1},\seqname{S}^2)
$
agrees with $\funcname{f}$ on $U^1_{u^1}$.

\item
Since $\union[{u^1 \in U^1}]{U^1_{u^1}} = U^1$ and
$\union[{u^1 \in U^1}]{U^2} = U^2$, then
$\funcname{f} \arin \catname{M}^2$ by \pagecref{def:modcat}.
\end{enumerate}
\end{enumerate}
\end{proof}

\item
\label{M-morph:localislocaltosup}
If $\funcname{f} \maps S^1 \to S^2$ is a local m-morphism of
$\seqname{S}^1$ to $\seqname{S}^2$ then $\funcname{f}$ is a local
m-morphism of $\seqname{S}^1$ to $\seqname{S}'^2$.

\begin{proof}
$\funcname{f} \maps S^1 \to S'^2$ is continuous by
\pagecref{cor:modsub}.

Let $u \in S^1$, $U_u$ be a model neighborhood of $\seqname{S}^1$ for
$u$ and $V_u$ be a model neighborhood of $\seqname{S}^2$ for
$v \defeq \funcname{f}(u)$ such that $\funcname{f}[U_u] \subseteq V_u$
and $\funcname{f} \maps U_u \to  V_u$ is a morphism of $\seqname{S}^2$.

Since $\seqname{S}^2 \submod \seqname{S}'^2$ by hypothesis,
$\funcname{f} \maps U_u \to  V_u$ is a morphism of $\seqname{S}'^2$.
\end{proof}

\item
\label{M-morph:morphismorphtosup}
If $\funcname{f} \maps S^1 \to S^2$ is an m-morphism of
$\seqname{S}^1$ to $\seqname{S}^2$, then $\funcname{f}$ is an m-morphism
of $\seqname{S}^1$ to $\seqname{S}'^2$.

\begin{proof}
$\seqname{S}^1 \submod \seqname{S}^2$ by \pagecref{def:modMmorph} and
$\seqname{S}^2 \submod \seqname{S}'^2$ by hypothesis, so
$\seqname{S}^1 \submod \seqname{S}'^2$.

$\funcname{f}$ is a morphism of $\seqname{S}^2$ by
\pagecref{def:modMmorph}, hence a morphism of $\seqname{S}'^2$.
\end{proof}

\item
\label{M-morph:catlocalislocaltosup}
If $\funcname{f} \maps S^1 \to S^2$ is a local
$\catname{M}^1$-$\catname{M}^2$ m-morphism of $\seqname{S}^1$ to
$\seqname{S}^2$, then $\funcname{f}$ is a local
$\catname{M}'^1$-$\catname{M}'^2$ m-morphism of $\seqname{S}^1$ to
$\seqname{S}'^2$.

\begin{proof}
Let $u \in S^1$, $U_u$ be a model neighborhood of $\seqname{S}^1$ for
$u$ and $V_u$ be a model neighborhood of $\seqname{S}^2$ for
$v \defeq \funcname{f}(u)$ such that $\funcname{f}[U_u] \subseteq V_u$
and $\funcname{f} \maps U_u \to  V_u$ is a morphism of $\catname{S}^2$.

Since $\catname{S}^2 \subcat[full-] \catname{S}'^2$ by hypothesis,
$\funcname{f} \maps U_u \to  V_u$ is a morphism of $\catname{S}'^2$.
\end{proof}

\item
\label{M-morph:catmorphismorphtosup}
If $\funcname{f} \maps S^1 \to S^2$ is an
$\catname{M}^1$-$\catname{M}^2$ m-morphism of $\seqname{S}^1$ to
$\seqname{S}^2$, then $\funcname{f}$ is an
$\catname{M}'^1$-$\catname{M}'^2$ m-morphism of $\seqname{S}^1$ to
$\seqname{S}'^2$.

\begin{proof}
$\catname{S}^1 \subcat[full-] \catname{S}^2$ by \pagecref{def:modMmorph} and \\*
$\catname{S}^2 \subcat[full-] \catname{S}'^2$ by hypothesis, so
$\catname{S}^1 \subcat[full-] \catname{S}'^2$.

$\funcname{f}$ is a morphism of $\catname{S}^1$ by
\cref{def:modMmorph}, hence a morphism of $\catname{S}'^2$.
\end{proof}

\item
\label{M-morph:contisstrict}
If $\seqname{S}^1 \submod \seqname{S}^2$ and
$\funcname{f} \maps S^1 \to S^2$ is continuous, then $\funcname{f}$ is
an m-morphism of $\Triv{S^1}$ to $\Triv{S^2}$.

\begin{proof}
\leavevmode \\*
\begin{description}
\item[$\funcname{f}$ is continuous:]
\leavevmode \\*
$\funcname{f} \maps S^1 \to S^2$ is continuous.
$
  \Top \biggl ( \Triv{S^i} \biggr ) =
  \Top(\seqname{S}^i) =
  S^i
$,
$i=1,2$.

\item [$\funcname{f}$ maps model neighborhoods into model neighborhoods:]
\leavevmode \\*
$S^2$ is open in $S^2$, hence a model neighborhood of $\Triv{S^2}$.
Let $U$ be a model neighborhood of $\Triv{S^1}$. Then
$\funcname{f}[U] \subseteq S^2$.

\item [Inverse images of model neighborhoods are model neighborhoods:]
\leavevmode \\*
Let $V$ be a model neighborhood of $\Triv{S^2}$. $V$ is open in $S^2$
and $\funcname{f}$ is continuous, so $\funcname{f}^{-1}[V]$ is open in
$S^1$ and hence a model neighborhood of $\Triv{S^1}$.
\end{description}
\end{proof}

\item
\label{M-morph:Idislocal}
If $\seqname{S}^1 \submod \seqname{S}^2$ then $\Id[S^1,S^2]$ is a local
m-morphism of $\seqname{S}^1$ to $\seqname{S}^2$. If each ${S}^i$ is
a model neighborhood of $\seqname{S}^i$ then $\Id[S^1,S^2]$ is a
morphism of $\catname{S}^2$.

\begin{proof}
Let $u \in S^1$ and $U_u$ a model neighborhood of $\seqname{S}^1$ for
$u$. Since $\seqname{S}^1 \submod \seqname{S}^2$, $U_u$ is also a
model neighborhood of $\seqname{S}^2$ for $u$, and hence
$\Id[U_u] \arin \catname{S}^2$.

If each $S^i$ is a model neighborhood of $\seqname{S}^i$, then since
 $\seqname{S}^1 \submod \seqname{S}^2$ by hypothesis. $S^1$ is a model
neighborhood of $\seqname{S}^2$.  The result follows by
\cref{mod:inclusion} of
{
  \showlabelsinline
  \pagecref{def:model}.
}
\end{proof}

\item
\label{M-morph:Idcatislocal}
If $\catname{M}^1 \subcat \catname{M}^2$ and
$\seqname{S}^1 \submod \seqname{S}^2$ then $\Id[S^1,S^2]$ is a
morphism of $\catname{M}^2$.

\begin{proof}
$\seqname{S}^1 \objin \catname{M}^2$ because
$\catname{M}^1 \subcat \catname{M}^2$.
$\Id[S^2] \arin \catname{M}^2$. Then
$\Id[S^1,S^2] \arin \catname{M}^2$ by
{
  \showlabelsinline
  \cref{modcat:inc}
}
of \pagecref{def:modcat}.
\end{proof}

\item
\label{M-morph:openmorphistrivmorph}
Let $U^i \subseteq \seqname{S}^i$, $i=1,2$, be open.  If
$\funcname{f} \maps U^1 \to U^2$ is a (local)
$\seqname{S}^1$-$\seqname{S}^2$ m-morphism of $U^1$ to $U^2$ then
$\funcname{f} \maps U^1 \to U^2$ is a $\Triv{S^1}$-$\Triv{S^2}$
m-morphism of $U^1$ to $U^2$.

\begin{proof}
Since $\seqname{S}^1 \submod \seqname{S}^2$, $\Triv{S^1} \subsp \Triv{^2}$.
Since $\funcname{f} \maps U^1 \to U^2$ is a model function, it is
continuous and hence a morphism of $\Triv{S^2}$
\end{proof}
\end{enumerate}
\end{lemma}

\begin{corollary}[(Local) m-morphisms]
\label{cor:modMmorph}
Let $\catname{M}'^i$, $\catname{M}^i \subcat[full-] \catname{M}'^i$,
$i=1,2$, be model categories,
$\seqname{S}^i \defeq (S^i, \catname{S}^i) \objin \catname{M}^i$,
$\seqname{S}'^i \defeq (S'^i, \catname{S}'^i) \objin \catname{M}'^i$
and $\seqname{S}^i \submod \seqname{S}'^i$.

\begin{enumerate}
\item
\label{M-morph:morphislocal}
If $\funcname{f} \maps \seqname{S}^1$ to $\seqname{S}^2$ is an
m-morphism of $\seqname{S}^1$ to $\seqname{S}^2$ then
$\funcname{f} \maps \seqname{S}^1$ to $\seqname{S}^2$ is a local
m-morphism of $\seqname{S}^1$ to $\seqname{S}^2$.

\begin{proof}
Since $\funcname{f} \maps \seqname{S}^1$ to $\seqname{S}^2$ is a
$\seqname{S}^1$-$\seqname{S}^2$ m-morphism of $S^1$ to $S^2$ then
$\funcname{f} \maps \seqname{S}^1$ to $\seqname{S}^2$ is a local
$\seqname{S}^1$-$\seqname{S}^2$ m-morphism of $S^1$ to $S^2$.
\end{proof}

\item
\label{M-morph:sublocalismorph}
If $U^1$ and If $U^2$ are model neighborhoods of $\seqname{S}^i$ then
every function $\funcname{f} \maps U^1 \to U^2$ that is a local
$\seqname{S}^i$ m-morphism of $U^1$ to $U^2$ is a $\seqname{S}^i$
m-morphism of $U^1$ to $U^2$.

\begin{proof}
Since $\funcname{f} \maps U^1 \to U^2$ is a local
$\seqname{S}^i$-$\seqname{S}^i$ m-morphism of $U^1$ to $U^2$ then it is
an $\seqname{S}^i$-$\seqname{S}^i$ m-morphism of $U^1$ to $U^2$.
\end{proof}

\item
\label{M-morph:localismorph}
If $\funcname{f} \maps S^1 \to S^2$ is a local m-morphism of
$\seqname{S}^1$ to $\seqname{S}^2$ and each $S^i$ is a model
neighborhood of $\seqname{S}^i$, then $\funcname{f} \maps S^1 \to S^2$
is an m-morphism of $\seqname{S}^1$ to $\seqname{S}^2$.

\begin{proof}
Since $\funcname{f} \maps S^1 \to S^2$ is a
$\seqname{S}^1$-$\seqname{S}^2$ local m-morphism of $S^1$ to $S^2$ then
it is an $\seqname{S}^1$-$\seqname{S}^2$ m-morphism of $S^1$ to
$S^2$.
\end{proof}

\item
\label{M-morph:catmorphislocal}
If $\funcname{f} \maps S^1 \to S^2$ is a $\catname{M}^1$-$\catname{M}^2$
m-morphism of $\seqname{S}^1$ to $\seqname{S}^2$ then
$\funcname{f} \maps S^1 \to S^2$ is a local
$\catname{M}^1$-$\catname{M}^2$ m-morphism of $\seqname{S}^1$ to
$\seqname{S}^2$.

\begin{proof}
Since $\funcname{f}$ is a
$\seqname{S}^1$-$\seqname{S}^2$-$\catname{M}^1$-$\catname{M}^2$
m-morphism of $S^1$ to $S^2$ then $\funcname{f}$ is a local
$\seqname{S}^1$-$\seqname{S}^2$-$\catname{M}^1$-$\catname{M}^2$
m-morphism of $S^1$ to $S^2$.
\end{proof}

\item
\label{M-morph:subcatlocalismorph}
If $\catname{M}^i$ satisfies the restricted sheaf condition and $U^j$ is
a model neighborhood of $\seqname{S}^i$, $j=1,2$, then every function
$\funcname{f} \maps U^1 \to U^2$ that is a local
$\seqname{S}^i$-$\catname{M}^i$ m-morphism of $U^1$ to $U^2$ is a
$\seqname{S}^i$-$\catname{M}^i$ m-morphism of $U^1$ to $U^2$.

\begin{proof}
\leavevmode
\begin{enumerate}
\item
$
  \funcname{f} \maps
  \Mod(U^1,\seqname{S}^i) \to
  \Mod(U^2,\seqname{S}^i)
$
is a local
\nobreaks%
  {
    {$\seqname{S}^i$-$\seqname{S}^i$-},
    {$\catname{M}^i$-$\catname{M}^i$}
  }
m-morphism of $U^1$ to $U^2$ by \pagecref{def:modlocalMmorph}

\item
$
  \funcname{f} \maps
  \Mod(U^1,\seqname{S}^i) \to
  \Mod(U^2,\seqname{S}^i)
$
is a
\nobreaks%
  {
    {$\seqname{S}^i$-$\seqname{S}^i$-},
    {$\catname{M}^i$-$\catname{M}^i$}
  }
m-morphism of $U^1$ to $U^2$ by \cref{M-morph:opencatlocalismorph} of
{
  \showlabelsinline
  \cref{lem:modMmorph}.
}

\item
The result follows from \cref{def:modlocalMmorph}.
\end{enumerate}
\end{proof}

\item
\label{M-morph:catlocalismorph}
Let each $S^i$ be a model neighborhood of $\seqname{S}^i$, Then every
$\funcname{f} \maps S^1 \to S^2$ that is a
$\catname{M}^1$-$\catname{M}^2$ local m-morphism of $\seqname{S}^1$ to
$\seqname{S}^2$ is a $\catname{M}^1$-$\catname{M}^2$ m-morphism
of $\seqname{S}^1$ to $\seqname{S}^2$.

\end{enumerate}

\begin{enumerate}
\item
If $\funcname{f} \maps S^1 \to S^2$ is a local m-morphism of
$\seqname{S}^1$ to $\seqname{S}^2$ and
$\seqname{S}^2 \submod \seqname{S}'^2$ then
$\funcname{f}$ is a local m-morphism of $\seqname{S}^1$ to
$\seqname{S}'^2$.

\begin{proof}
$\funcname{f} \maps S^1 \to S'^2$ is continuous by
\pagecref{cor:modsub}.
\end{proof}

\item
If $\funcname{f} \maps S^1 \to S^2$ is a(n) (local) m-morphism of
$\seqname{S}^1$ to $\seqname{S}^2$ then $\funcname{f}$ is an m-morphism
of $\Triv{S^1}$ to $\Triv{S^2}$.

\begin{proof}
$\funcname{f}$ is continuous by \pagecref{def:modlocalMmorph}.
$
  \Top \biggl ( \Triv{S^i} \biggr ) =
  \Top(\seqname{S}^i) =
  S^i
$,
$i=1,2$.
\end{proof}

\item
If $S^i$ is a model neighborhood of $\seqname{S}^i$ then
$\Id[{\seqname{S}^i}]$ is a morphism of $\catname{S}^i$.

\begin{proof}
$\seqname{S}^i \submod \seqname{S}^i$.
\end{proof}

\item
If $\funcname{f} \maps S^1 \to S^2$ is a local m-morphism of
$\seqname{S}^1$ to $\seqname{S}^2$ and
$\funcname{f}^2 \maps S^2 \to S^3$ is an m-morphism of $\seqname{S}^2$
to $\seqname{S}^3$ then $\funcname{f}^2 \compose \funcname{f}$ is a
local m-morphism of $\seqname{S}^1$ to $\seqname{S}^3$.

\begin{proof}
Since $\funcname{f}^2 \maps S^2 \to S^3$ is an m-morphism of
$\seqname{S}^2$ to $\seqname{S}^3$ then $\funcname{f}^2$ is a local
m-morphism of $\seqname{S}^2$ to $\seqname{S}^3$.
\end{proof}

\item
If $\funcname{f} \maps S^1 \to S^2$ is an m-morphism of
$\seqname{S}^1$ to $\seqname{S}^2$ and
$\funcname{f}^2 \maps S^2 \to S^3$ is a local m-morphism of
$\seqname{S}^2$ to $\seqname{S}^3$ then
$\funcname{f}^2 \compose \funcname{f}$ is a local m-morphism of
$\seqname{S}^1$ to $\seqname{S}^3$.

\begin{proof}
Since $\funcname{f} \maps S^1 \to S^2$ is an m-morphism of
$\seqname{S}^1$ to $\seqname{S}^2$ then $\funcname{f}$ is a local
m-morphism of $\seqname{S}^2$ to $\seqname{S}^2$.
\end{proof}

\item
If $\funcname{f} \maps S^1 \to S^2$ is a local
$\catname{M}^1$-$\catname{M}^2$ m-morphism of $\seqname{S}^1$ to
$\seqname{S}^2$ and $\funcname{f}^2 \maps S^2 \to S^3$ is an
$\catname{M}^2$-$\catname{M}^3$ m-morphism of $\seqname{S}^2$ to
$\seqname{S}^3$ then $\funcname{f}^2 \compose \funcname{f}$ is a local
$\catname{M}^1$-$\catname{M}^3$ m-morphism of $\seqname{S}^1$ to
$\seqname{S}^3$.

\begin{proof}
Since $\funcname{f}^2 \maps S^2 \to S^3$ is an
$\catname{M}^2$-$\catname{M}^3$ m-morphism of $\seqname{S}^2$ to
$\seqname{S}^3$ then $\funcname{f}^2$ is a local
$\catname{M}^2$-$\catname{M}^3$ m-morphism of $\seqname{S}^2$ to
$\seqname{S}^3$.
\end{proof}

\item
If $\funcname{f} \maps S^1 \to S^2$ is an
$\catname{M}^1$-$\catname{M}^2$ m-morphism of $\seqname{S}^1$ to
$\seqname{S}^2$ and $\funcname{f}^2 \maps S^2 \to S^3$ is a local
$\catname{M}^2$-$\catname{M}^3$ m-morphism of $\seqname{S}^2$ to
$\seqname{S}^3$ then $\funcname{f}^2 \compose \funcname{f}$ is a local
$\catname{M}^1$-$\catname{M}^3$ m-morphism of $\seqname{S}^1$ to
$\seqname{S}^3$.

\begin{proof}
Since $\funcname{f} \maps S^1 \to S^2$ is an
$\catname{M}^1$-$\catname{M}^2$ m-morphism of $\seqname{S}^1$ to
$\seqname{S}^2$ then $\funcname{f}$ is a local
$\catname{M}^1$-$\catname{M}^2$ m-morphism of $\seqname{S}^2$ to
$\seqname{S}^2$.
\end{proof}

\item
\label{M-morph:morphistrivmorph}
If $\funcname{f} \maps \seqname{S}^1 \to \seqname{S}^2$ is a (local)
m-morphism of $\seqname{S}^1$ to $\seqname{S}^2$ then
$\funcname{f} \maps S^1 \to S^2$ is an m-morphism of $\Triv{S^1}$ to
$\Triv{S^2}$.

\begin{proof}
Each $S^i$ is open, so $\funcname{f} \maps S^1 \to S^2$ is a
$\Triv{S^1}$-$\Triv{S^2}$ m-morphism of $S^1$ to $S^2$.
\end{proof}
\end{enumerate}
\end{corollary}

\begin{lemma}[Composition of local m-morphisms]
\label{lem:modlocalMmorphcomp}
Let $\catname{M}^i$, $i=1,2,3$, be model categories,
$\seqname{S}^i \defeq (S^i, \catname{S}^i) \objin \catname{M}^i$ and
$U^i \subseteq S^i$ be open.

\begin{enumerate}
\item
If $\funcname{f}^i \maps U^i \to U^{i+1}$, $i=1,2$, is a local
$\seqname{S}^i$-$\seqname{S}^{i+1}$ m-morphism of $U^i$ to $U^{i+1}$ then
$\funcname{f}^2 \compose \funcname{f}^1$ a local
$\seqname{S}^1$-$\seqname{S}^3$ m-morphism of $U^i$ to $U^{i+1}$.

\begin{proof}
Since $\seqname{S}^1 \submod \seqname{S}^2$ and
$\seqname{S}^2 \submod \seqname{S}^3$ then
$\seqname{S}^1 \submod \seqname{S}^3$, Let $u^1 \in U^1$,
$u^2 \defeq \funcname{f}^1(u^1)$ and $u^3 \defeq \funcname{f}^2(u^2)$.
Let $U_{u^i}$ and $V_{u^i}$ be model neighborhoods of $\seqname{S}^i$
and $\seqname{S}^{i+1}$ such that that
$\funcname{f}^i[U_{u^i}] \subseteq V_{u^i}$ and
$
  \funcname{f}^i \maps \Mod(U_{u^i},\seqname{S}^i)
  \to \Mod(V_{u^i},\seqname{S}^{i+1})
$
is a morphism of $\seqname{S}^{i+1}$.

Let
$
  U'_{u^1} \defeq
  U_{u^1} \cap (\funcname{f}^2 \compose \funcname{f}^1)^{-1}[V_{u^3}]
$.
Then $\funcname{f}^2 \compose \funcname{f}^1 \maps U'_{u^1} \to V_{u^3}$
is a morphism of $\seqname{S}^3$.
\end{proof}

\item
If $\funcname{f}^i \maps S^i \to S^{i+1}$, $i=1,2$, is a local
$\catname{M}^i$-$\catname{M}^{i+1}$ m-morphism of $\seqname{S}^i$ to
$\seqname{S}^{i+1}$ then $\funcname{f}^2 \compose \funcname{f}^1$ a
local $\catname{M}^i$-$\catname{M}^{i+2}$ m-morphism of
$\seqname{S}^1$ to $\seqname{S}^3$.

\begin{proof}
Since $\catname{M}^1 \subcat[full-] \catname{M}^2$ and
$\seqname{M}^2 \subcat[full-] \catname{M}^3$ then
$\seqname{M}^1 \subcat[full-] \catname{M}^3$.
Let
\nobreaks
  {{
    $u^1 \in \seqname{S}^1$,
  }}
$u^2 \defeq \funcname{f}^1(u^)$ and $u^3 \defeq \funcname{f}^2(u^2)$.
Let $U_{u^i}$ and $V_{u^i}$ be model neighborhoods of $\seqname{S}^i$
and $\seqname{S}^{i+1}$ such that
$\funcname{f}^i \maps U_{u^i} \to V_{u^i}$ is a morphism of
$\catname{M}^{i+1}$. Let
\nobreaks
  {{
    $
      U'_{u^1} \defeq
      U_{u^1} \cap (\funcname{f}^3 \compose \funcname{f}^2)^{-1}[V_{u^3}]
    $.
  }}
Then $\funcname{f}^2 \compose \funcname{f}^2 \maps U'_{u^1} \to V_{u^3}$
is a morphism of $\catname{M}^3$ by
{
  \showlabelsinline
  \cref{modcat:rest}
}
of \pagecref{def:modcat}.
\end{proof}
\end{enumerate}

If each $\funcname{f}^i \maps \seqname{S}^i \to \seqname{S}^{i+1}$,
$i=1,2$, is a local m-morphism of $\seqname{S}^i$ to
$\seqname{S}^{i+1}$, then
\nobreaks
  {{
    $\funcname{f}^2 \compose \funcname{f}^1 \maps S^1 \to S^3$
  }}
is a local m-morphism of $\seqname{S}^1 \maps S^1 \to S^3$.

\begin{proof}
Since $\seqname{M}^1 \submod \seqname{M}^2$ and
$\seqname{M}^2 \submod \seqname{M}^3$,
$\seqname{M}^1 \submod \seqname{M}^3$.
Let $u \in S^i$, $v \defeq \funcname{f}^1(u)$ and
$w \defeq \funcname{f}^2(v)$. There exist a model neighborhood $U_u$
for $u$, model neighborhoods $V_u$, $V'_u$ for $v$ and a model
neighborhood $W_v$ of $w$ such that
\nobreaks
  {{
    $\funcname{f}_1[U_u] \subseteq V'_u$,
  }}
$\funcname{f}_2[V_v] \subseteq W_v$,
$\funcname{f}_1 \maps U_u \to V'_u$ is a morphism of $\seqname{S}^2$
and $\funcname{f}_2 \maps V_v \to W_v$ is a morphism of
$\catname{S}^3$. Then $\hat{V_u} \defeq V_v \cap V'_u \neq \emptyset$,
$\hat{V_u}$ is a model neighborhood of $v$ and
$\hat{U_u} \defeq \funcname{f}^{i-1}_1[\hat{V_u}]$ is a model
neighborhood for $u$.  $\funcname{f}_1 \maps \hat{U_u} \to \hat{V_u}$
and $\funcname{f}_2 \maps \hat{V_u} \to W_v$ are morphisms of
$\catname{S}^3$ by
{
  \showlabelsinline
  \cref{mod:restriction}
}
of \pagecref{def:model} and thus
$\funcname{f}_2 \compose \funcname{f}_1 \maps \hat{U_u} \to W_v$ is a
morphism of $\seqname{S}^3$.
\end{proof}

If each $\funcname{f}^i \maps \seqname{S}^i \to \seqname{S}^{i+1}$ is
a local $\catname{M}^i$-$\catname{M}^{i+1}$ m-morphism of
$\seqname{S}^i$ to $\seqname{S}^{i+1}$ then \\*
\nobreaks
  {{
    $
      \funcname{f}^2 \compose \funcname{f}^1 \maps
      \seqname{S}^1 \to \seqname{S}^3
    $
  }}
is a local $\catname{M}^1$-$\catname{M}^3$ m-morphism of
$\seqname{S}^1$ to $\seqname{S}^3$.

\begin{proof}
Since $\catname{M}^1 \subcat[full-] \catname{M}^2$ and
$\catname{M}^2 \subcat[full-] \catname{M}^3$,
$\catname{M}^1 \subcat[full-] \catname{M}^3$.
Let \nobreaks{{$u \in \seqname{S}^1$,}} $v \defeq \funcname{f}^1(u)$ and
$w \defeq \funcname{f}^2(v)$. There exist a model neighborhood $U_u$
for $u$, model neighborhoods $V_v$, $V'_u$ for $v$ and a model
neighborhood $W_v$ of $w$ such that
$\funcname{f}^1[U_u] \subseteq V'_u$,
$\funcname{f}^1 \maps U_u \to V'_u$ is a morphism of $\catname{M}^2$,
$\funcname{f}^2[V_v] \subseteq W_v$ and
$\funcname{f}^2 \maps V_v \to W_v$ is a morphism of $\catname{M}^3$.
Then $\hat{V_u} \defeq V_v \cap V'_u \neq \emptyset$, $\hat{V_u}$ is a
model neighborhood of $v$ and
$\hat{U_u} \defeq \funcname{f}^{i-1}_1[\hat{V_u}]$ is a model
neighborhood for $u$.  $\funcname{f}^1 \maps \hat{U_u} \to \hat{V_u}$
and $\funcname{f}^2 \maps \hat{V_u} \to W_v$ are morphisms of
$\catname{M}^3$ by
{
  \showlabelsinline
  \cref{modcat:inc}
}
of \pagecref{def:modcat} and thus
$\funcname{f}^2 \compose \funcname{f}^1 \maps \hat{U_u} \to W_v$ is a
morphism of $\catname{M}^3$.
\end{proof}
\end{lemma}

\begin{corollary}[[Composition of local m-morphisms]
\label{cor:modlocalMmorphcomp}
Let $\catname{M}^i$, $i=1,2,3$, be model categories,
$\seqname{S}^i \defeq (S^i, \catname{S}^i) \objin \catname{M}^i$.

\begin{enumerate}
\item
Let $U^j \subseteq S^i$, $j=1,2,3$, be open.
If $\funcname{f}^j \maps U^j \to U^{j+1}$, $j=1,2$, is a local
$\seqname{S}^i$ m-morphism of $U^j$ to $U^{j+1}$ then
$\funcname{f}^2 \compose \funcname{f}^1$ is a local
$\seqname{S}^i$ m-morphism of $U^1$ to $U^3$.

\begin{proof}
Since each  $\funcname{f}^j \maps U^j \to U^{j+1}$ is a local
$\seqname{S}^i$-$\seqname{S}^i$ m-morphism of $U^j$ to $U^{j+1}$ then
$\funcname{f}^2 \compose \funcname{f}^1$ is a local
$\seqname{S}^i$-$\seqname{S}^i$ m-morphism of $U^j$ to $U^{j+1}$.
\end{proof}

\item
If $\funcname{f}^i \maps S^i \to S^{i+1}$, $i=1,2$, is a local
m-morphism of $\seqname{S}^i$ to $\seqname{S}^{i+1}$ then
$\funcname{f}^2 \compose \funcname{f}^1$ is a local m-morphism of
$\seqname{S}^1$ to $\seqname{S}^3$.

\begin{proof}
Since $\seqname{S}^1 \submod \seqname{S}^2$ and
$\seqname{S}^2 \submod \seqname{S}^3$ then
$\seqname{S}^1 \submod \seqname{S}^3$.
Let $u^1 \in \seqname{S}^1$, $u^2 \defeq \funcname{f}^1(u^1)$ and
$u^3 \defeq \funcname{f}^2(u^2)$.
Let $U_{u^i}$ and $V_{u^i}$ be model neighborhoods of $\seqname{S}^i$ and $\seqname{S}^{i+1}$
such that$\funcname{f}^i \maps U_{u^i} \to V_{u^i}$ is a morphism of $\seqname{S}^{i+1}$.
Let
\nobreaks
  {{
    $
      U'_{u^1} \defeq
      U_{u^1} \cap (\funcname{f}^2 \compose \funcname{f}^2)^{-1}[V_{u^3}]
    $.
  }}
Then $\funcname{f}^2 \compose \funcname{f}^2 \maps U'_{u^1} \to V_{u^3}$
is a morphism of $\seqname{S}^3$.
\end{proof}
\end{enumerate}

\end{corollary}

\begin{lemma}[Composition of m-morphisms]
\label{lem:modMmorphcomp}
Let $\catname{M}^i$, $i=1,2,3$, be model categories and
$\seqname{S}^i \defeq (S^i, \catname{S}^i) \objin \catname{M}^i$.

\begin{enumerate}

\item
If $\funcname{f}^i \maps S^i \to S^{i+1}$, $i=1,2$, is an m-morphism of
$\seqname{S}^i$ to $\seqname{S}^{i+1}$ then
$\funcname{f}^2 \compose \funcname{f}^1 \maps S^1 \to S^3$ is an
m-morphism of $\seqname{S}^1 \maps S^1 \to S^3$.

\begin{proof}
Since each $\funcname{f}^i$ is a morphism of $\seqname{S}^{i+1}$ by
\pagecref{def:modMmorph} and
$\seqname{S}^1 \submod \seqname{S}^2 \submod \seqname{S}^3$, then
each $\funcname{f}^i$ is a morphism of $\seqname{S}^3$ by
\pagecref{def:modMmorph}, and thus
$\funcname{f}^2 \compose \funcname{f}^1 \maps S^1 \to S^3$ is a morphism
of $\seqname{S}^3$.
\end{proof}

\item
If $\funcname{f}^i \maps S^i \to S^{i+1}$, $i=1,2$, is an
$\catname{M}^i$-$\catname{M}^{i+1}$ m-morphism of $\seqname{S}^i$ to
$\seqname{S}^{i+1}$ then $\funcname{f}^2 \compose \funcname{f}^1$ an
$\catname{M}^i$-$\catname{M}^{i+2}$ m-morphism of $\seqname{S}^1$ to
$\seqname{S}^3$.

\begin{proof}
Since $\catname{M}^1 \subcat[full-] \catname{M}^2$ and
$\seqname{M}^2 \subcat[full-] \catname{M}^3$ then
$\seqname{M}^1 \subcat[full-] \catname{M}^3$.
Since each $\funcname{f}^i \maps S^i \to S^{i+1}$ is a morphism of
$\catname{S}^3$, then so is $\funcname{f}^2 \compose \funcname{f}^1$
\end{proof}

\item
If each $\funcname{f}^i \maps \seqname{S}^i \to \seqname{S}^{i+1}$ is
an $\catname{M}^i$-$\catname{M}^{i+1}$ m-morphism of
$\seqname{S}^i$ to $\seqname{S}^{i+1}$ then \\*
\nobreaks
  {{
    $
      \funcname{f}^2 \compose \funcname{f}^1 \maps
      \seqname{S}^1 \to \seqname{S}^3
    $
  }}
is an $\catname{M}^1$-$\catname{M}^3$ m-morphism of $\seqname{S}^1$ to
$\seqname{S}^3$.

\begin{proof}
Since $\catname{M}^1 \subcat[full-] \catname{M}^2$ and
$\catname{M}^2 \subcat[full-] \catname{M}^3$ then
$\catname{M}^1 \subcat[full-] \catname{M}^3$ and
$\funcname{f}^i$ is a morphism of $\catname{S}^{i+1}$ by
\pagecref{def:modMmorph}.
Since $\catname{M}^2 \subcat[full-] \catname{M}^3$, $\funcname{f}^1$ is a
morphism of $\catname{S}^3$. and thus
$\funcname{f}^2 \compose \funcname{f}^1 \maps S^1 \to S^3$ is a morphism
of $\catname{S}^3$.
\end{proof}
\end{enumerate}
\end{lemma}

\section{Spaces and proper functions}
\label{sec:spaces}
Several of the following definitions involve spaces of a restrictive
character and specific types of mappings among them. Except where
otherwise qualified, the word \textit{space} will have this restricted
meaning.

\begin{definition}[Spaces]
\label{def:space}
A space is a topological space, a model space or either with an
additional associated structure.
\end{definition}

\begin{definition}[Proper functions]
\label{def:proper}
Let $S^1$, $S^2$ be spaces and $\funcname{f} \maps S^1 \to S^2$ a
continuous function; $\funcname{f}$ need not preserve any associated
algebraic structure\footnote{
  However, in practice commutation relations will often enforce the
  preservation of algebraic structures.
}.
$\funcname{f}$ is a proper\footnote{
  In the spirit of \cite[footnote, p.~112] {GenTop}
  \begin{quote}
    This nomenclature is an excellent example of the time-honored
    custom of referring to a problem we cannot handle as abnormal,
    irregular, improper, degenerate, inadmissible, and otherwise
    undesirable.
  \end{quote}
}
function iff

\begin{enumerate}
\item $S^1$ and $S^2$ are both $\truthspace$ or both
$\seqname{Truthspace}$, and $\funcname{f}(\true)=\true$.

\item $S^1$, $S^2$ are topological spaces other than
$\truthspace$.

\item $S^1$ is a topological space other than
$\truthspace$, $S^2$ is a model space other than
$\seqname{Truthspace}$ and the images of open sets are contained in
model neighborhoods.

\item $S^1$ is a model space other than $\seqname{Truthspace}$, $S^2$ is
a topological space other than $\truthspace$ and the
inverse images of open sets are model neighborhoods.

\item $S^1$, $S^2$ are both model spaces other than
$\seqname{Truthspace}$ and $\funcname{f}$ is a model function.
\end{enumerate}
\end{definition}

\begin{definition}[Singleton categories and sequences of singleton categories]
\label {def:singcat}
Let $S$ be a topological space. Then the singleton category of $S$,
abbreviated $\singcat{S}$, is the category whose sole object is $S$ and
whose sole morphism is the identity morphism of $S$ to itself.

Let $\seqname{S}$ be a model space. Then the singleton category of
$\seqname{S}$, abbreviated $\singcat{\seqname{S}}$, is the category whose
sole object is $\seqname{S}$ and whose sole morphism is the identity
morphism of $\seqname{S}$ to itself.

Let $\seqname{S} \defeq (S^\alpha, \alpha \prec \Alpha)$ be a sequence
of spaces. Then the singleton category sequence of $\seqname{S}$,
abbreviated $\Singcat{\seqname{S}}$, is
$(\singcat{S^\alpha}, \alpha \prec \Alpha)$.

Let $\seqname{S}$ be a set of spaces. Then the singleton category of
$\seqname{S}$, abbreviated $\Singcat{\seqname{S}}$, is
$\union[S \in \seqname{S}]{\singcat{S}}$.

\begin{remark}
By abuse of language the notation $\Singcat{\catseqname{S}}$ will be used
to name sequences of categories constructed with this and similar
functions.
\end{remark}
\end{definition}

% Let $\seqname{M}=(S, \catname{S})$ be a model space for $S$ and $P\subset \catname{S}$. We say that $P$ is a
% patching for $\catname{S}$, the elements of $P$ are patches and that $\seqname P \defeq
% (S, \catname{S},P)$ is a patched space iff
%
% \begin{enumerate}
% \item $P$ is a cover for  $S$.
% \item $P$ is a basis for the finite intersections of elements of  $P$.
% \end{enumerate}

\part{M-charts and m-atlases}
\label{part:M-charts}
The literature defines fiber bundles and manifolds using the language of
charts, atlases and transition functions; it has multiple equivalent
definitions. Some authors start with topological spaces and define
atlases over them. Some start with abstract sets and define atlases over
them, deriving the topology from the atlas. Some start with an indexed
set of open patches in a coordinate space and transition functions among
them satisfying a cocycle (compatibility) constraint, and then derive
the total space as a quotient space of the disjoin union of the patches.
Some authors use maximal atlases while others use equivalence classes of
atlases.

This paper uses the first approach, explicitly making the relevant
spaces topological spaces, but modifies the definitions of charts in
order to make them fit more natuarally into the context of local
coordinate spaces. It adds a prefix, e.g., $\Ck$, M, in order to avoid
confusion with the conventional definitions. An m-atlas based on a
topological spaces is closer to the conventional definitions of an atlas
for a manifold while an m-atlas based on a model space is suitable for
defining both manifolds and fiber bundles. Although simple manifolds and
fiber bundles could both be defined directly in terms of maximal
m-atlases, this paper has a different perspective, and uses the
m-atlases as part of the more general Local Coordinate Space (LCS),
presented in \cite{LCS}.

\Crefrange{sec:M-charts}{sec:M-atlas} define m-charts and m-atlasses.
Maximal m-atlasses are essentially the same as manifolds.

\Cref{sec:M-ATLmorph} defines morphisms between m-atlasses, referred to
as m-atlas morphisms, in a fashion tailored to use in defining local
coordinate spaces. It constructs categories of M-atlases and functors,
and proves some basic results.  It also defines a related concept of
m-atlas near morphisms, although the restriction needed to use them as
morphisms in a category of m-atlases requires them to be morphisms.

\Cref{sec:M-ATLcmorph} defines morphisms between m-atlasses in a
fashion similar to that conventionally used for $\Ck$ maps between
differentiable manifolds, in order to provide context. It constructs
categories of M-atlases and functors, and proves some basic results.

\section{M-charts}
{
  \showlabelsinline
  \label{sec:M-charts}
}
\begin{definition}[M-charts]
\label{def:M-chart}
Let $\seqname{E} \defeq (E, \catname{E})$ and
$\seqname{C} \defeq (C,\catname{C})$ be model spaces. An m-chart
$(U, V, \phi)$ of $\seqname{E}$ in the coordinate space $\seqname{C}$
consists of
\begin{enumerate}
\item A nonvoid model neighborhood $U \objin \catname{E}$, known as a
coordinate patch
\item A model neighborhood $V\objin \catname{C}$
\item A model homeomorphism $\phi \maps U \toiso V$, known as a
coordinate function
\end{enumerate}
\begin{remark}
I consider it clearer to explicate the range, rather than the
conventional usage of specifying only the domain and function or the
minimalist usage of specifying only the function.
\end{remark}
For evey $u \in U$, $(U, V, \phi)$ is an m-chart of $\seqname{E}$ in the
coordinate space $\catname{C}$ at $u$ and covers $u$.
$\seqname{E}$ is the total model space or total space for the chart and
$\seqname{C}$ is the coordinate model space or coordinate space for the
chart.

Let $E$ be a topological space and $\seqname{C} \defeq (C,\catname{C})$
be a model space. $(U, V, \phi)$ is an m-chart of $E$ in the coordinate
space $\catname{C}$ iff it is an m-chart of $\Triv{E}$ in the coordinate
space $\catname{C}$.

For evey $x \in U$, $(U, V, \phi)$ is an m-chart of $E$ in the
coordinate space $\catname{C}$ at $x$.
$E$ is the total space for the chart and $\seqname{C}$ is the coordinate
model space or coordinate space for the chart.
\end{definition}

\begin{lemma}[M-charts]
\label{lem:M-chart}
\leavevmode
\begin{enumerate}
\item
\label{M-chart:openC}
Let $\seqname{E} \defeq (E, \catname{E})$ and
$\seqname{C} \defeq (C,\catname{C})$ be model spaces, and $(U, V, \phi)$
be an m-chart of $\seqname{E}$ in the coordinate space $\seqname{C}$
such that every open subset of $V$ is a model neighborhood of
$\seqname{C}$. Then $(U, V, \phi)$ is an m-chart of $E$ in the
coordinate space $\seqname{C}$.

\begin{proof}
$U$, $V$ and $\phi$ satisfy the conditions of the definition:
\begin{description}
\item%
  [
    $U$ is a model neighborhood of \ensuremath{\Triv{E}}:
  ]
$U$ is a model neighborhood of $\seqname{E}$, hence open.

\item[$V$ is a model neighborhood of $\seqname{C}$:]
By hypothesis.

\item%
  [
    $\phi$ is a model homeomorphism from
    \ensuremath{\Mod \biggl ( U,\Triv{E} \biggr )} to
    $\Mod(V,\seqname{C})$:
  ]
\leavevmode
\begin{enumerate}
\item
Since $\phi$ is a model homeomorphism from $\Mod(U,\seqname{E})$ to
$\Mod(V,\seqname{C})$, it is a homeomorphism.

\item
If $V'$ is a model subneighborhood of $V$, then $\phi^{-1}[V']$ is open
and hence a model neighborhood of $\Triv{E}$.

\item
If $U'$ is a model neighborhood of $\Triv{E}$ then it is open, hence
$\phi[U]$ is an open subset of $V$ and thus by hypothesis a model
neighborhood of $\seqname{C}$.
\end{enumerate}
\end{description}
\end{proof}

\item
\label{M-chart:openE}
Let $(U, V, \phi)$ be an m-chart of $E$ in the coordinate space
$\seqname{C}$. $(U, V, \phi)$ is an m-chart of $\seqname{E}$ in the
coordinate space $\seqname{C}$ iff every open subset of $U$ is a model
neighborhood of $\seqname{E}$.

\begin{proof}
If $(U, V, \phi)$ is an m-chart of $E$ in the coordinate space
$\seqname{C}$ and also an m-chart of $\seqname{E}$ in the coordinate
space $\seqname{C}$ then let $U' \subseteq U$ be open. Since
$(U, V, \phi)$ is an m-chart of $E$ in the coordinate space
$\seqname{C}$, $\phi^{-1}$ is a model function from
$\Mod(V,\seqname{C})$ to $\Triv{E}$ and $\phi[U']$ is a model
neighborhood of $\seqname{C}$.  Since $(U, V, \phi)$ is an m-chart of
$\seqname{E}$ in the coordinate space $\seqname{C}$, $\phi$ is a model
function from $\Mod(U,\seqname{E})$ to $\seqname{C}$ and thus
$\phi^{-1} \bigl [ \phi[U'] \bigr ]$ is a model neighborhood of $\seqname{E}$.

If $(U, V, \phi)$ is an m-chart of $E$ in the coordinate space
$\seqname{C}$ and every open
\nobreaks
  {{
    $U'\subseteq U$
  }}
is a model neighborhood of $\seqname{E}$ then
\begin{enumerate}
\item Since $U \subseteq U$ is open, $U$ is a model neighborhood of
$\seqname{E}$ by hypothesis.

\item $V$ is a model neighborhood of $\seqname{C}$ by hypothesis.

\item If $V' \subseteq V$ is a model neighborhood of $\seqname{C}$ then
$\phi^{-1}[V']$ is open and hence a model neighborhod of $\seqname{E}$
by hypothesis.

$\phi^{-1}$ is a model function from $\Mod \biggl ( U,\Triv{E} \biggr )$ to
$\Mod(V,\seqname{C})$. if $U' \subseteq U$ is a model neighborhood of
$\seqname{E}$ then $U'$ is open and hence a model neighborhood of
$\Triv{E}$. Then $\phi[U']$ is a model neighborhood of $\seqname{C}$.
\end{enumerate}
\end{proof}
\end{enumerate}
\end{lemma}

\begin{corollary}[M-charts]
\label{cor:M-chart}
Let $\seqname{E} \defeq (E, \catname{E})$ and
$\seqname{C} \defeq (C,\catname{C})$ be model spaces, $\seqname{C}$ be
fine grained and $(U, V, \phi)$ be an m-chart of $\seqname{E}$ in the
coordinate space $\seqname{C}$.  Then $(U, V, \phi)$ is an m-chart of
$E$ in the coordinate space $\seqname{C}$.

\begin{proof}
Since $V$ is open and $\seqname{C}$ is fine grained, every open subset of
$V$ is a model neighborhood of $\seqname{C}$. The result follows by
\cref{M-chart:openC} of
{
  \showlabelsinline
  \fullcref{lem:M-chart}
  on \cpageref{M-chart:openC}.
}
\end{proof}
\end{corollary}

\begin{definition}[Subcharts]
\label{def:M-subchart}
Let $(U, V, \phi)$ be an m-chart of
$\seqname{E} \defeq (E, \catname{E})$ in the coordinate space
$(C,\catname{C})$ and $U' \subseteq U$ a nonvoid model neighborhood of
$\seqname{E}$. Then
$(U', V', \phi') \defeq (U', \phi[U'], \phi \restrictto_{U',V'})$ is a
subchart of $(U, V, \phi)$.

Let $(U, V, \phi)$ be an m-chart of $E$ in the coordinate
space $(C,\catname{C})$ and $U' \subseteq U$ an open set of
$E$. Then
$(U', V', \phi') \defeq (U', \phi[U'], \phi \restrictto_{U',V'})$
is a subchart of $(U, V, \phi)$.

$(U', V', \phi')$ is a subchart of $(U, V, \phi)$ at  $x$ iff
$(U', V', \phi')$ is a subchart of $(U, V, \phi)$ and $x \in U'$.
\end{definition}

\begin{lemma}[Subcharts]
\label{lem:M-subchart}
Let $\seqname{E} \defeq (E, \catname{E})$ and
$\seqname{C} \defeq (C,\catname{C})$ be model spaces.

Let $(U, V, \phi)$ be an m-chart of $\seqname{E}$ in the coordinate
space $\seqname{C}$ and $(U', V', \phi')$ a subchart of $(U, V, \phi)$.
Then $(U', V', \phi')$ is an m-chart of $\seqname{E}$ in the coordinate
space $\seqname{C}$.

\begin{proof}
$U'$, $V'$ and $\phi'$ satisfy the conditions of the definition of an m-chart:
\begin{enumerate}
\item $U'$ is a model neighborhood of $\seqname{E}$ by the definition of
subchart.
\item Since $U'$ is a model neighborhood and $\phi$ is a model
homeomorphism,
\nobreaks
  {{
    $V'=\phi[U']$
  }}
is a model neighborhood.
\item Since $\phi$ and $\phi^{-1}$ are model functions, so are
$\phi' = \phi \restrictto_{U',V'}$ and
$\phi'^{-1} = \phi^{-1} \restrictto_{V',U'}$. Thus $\phi'$ is a model
homeomorphism.
\end{enumerate}
\end{proof}

Let $(U, V, \phi)$ be an m-chart of $E$ in the coordinate space
$\seqname{C}$ and $(U', V', \phi')$ a subchart of $(U, V, \phi)$. Then
$(U', V', \phi')$ is an m-chart of $E$ in the coordinate space
$\seqname{C}$.

\begin{proof}
$U'$, $V'$ and $\phi'$ satisfy the conditions of the definition of an m-chart:
\begin{enumerate}
\item $U'$ is open by the definition of subchart.
\item Since $\phi$ is a homeomorphism and $U'$ is open, $V'=\phi[U']$ is open.
\item Since $\phi$ is a homeomorphism that maps open sets into model
neighborhoods, $\phi \restrictto_{U',V'}$ is a homeomorphism that maps
open sets into model neighborhoods.
\end{enumerate}
\end{proof}
\end{lemma}

\begin{definition}[M-compatibility]
\label{def:M-compat}
Let $\seqname{E} \defeq (E, \catname{E})$ and
$\seqname{C} \defeq (C, \catname{C})$ be model spaces and
$(U_i, V_i, \phi_i)$, $i=1,2$, be an m-chart of $\seqname{E}$ in the
coordinate space $\seqname{C}$. Then $(U_1, V_1, \phi_1)$ is
m-compatible with $(U_2, V_2, \phi_2)$ in the coordinate space
$\seqname{C}$ iff either

\begin{enumerate}
\item $U_1$ and $U_2$ are disjoint

\item The transition function
$
  \funcname{t}=\phi_2 \compose \phi^{-1}_1
  \maps \phi_1[U_1 \cap U_2] \to \phi_2[U_1 \cap U_2]
$
is an isomorphism of $\catname{C}$.
\end{enumerate}

$(U_1, V_1, \phi_1)$ is also properly m-compatible with
\nobreaks
  {{
    $(U_2, V_2, \phi_2)$
  }}
in the space $\seqname{E}$ and the coordinate space $\seqname{C}$, and
properly m-compatible in the space $\seqname{E}$ and the coordinate
space $\seqname{C}$ with
\nobreaks
  {{
    $(U_2, V_2, \phi_2)$,
  }}
iff
for every model neighborhoods $U'_i \subseteq U_i$

\begin{enumerate}
\item
For every morphism $\funcname{f} \maps U'_1 \to U'_2$ of $\seqname{E}$,
the composition \\*
$
  \phi_2 \compose \funcname{f} \compose \phi^{-1}_1
  \maps \phi_1[U'_1] \to \phi_2[U'_2]
$
is a morphism of $\seqname{C}$.

\item
For every morphism $\funcname{g} \maps U'_2 \to U'_1$ of $\seqname{E}$,
the composition
$
  \phi_1 \compose \funcname{g} \compose \phi^{-1}_2
  \maps \phi_2[U'_2] \to \phi_1[U'_1]
$
is a morphism of $\seqname{C}$.
\end{enumerate}

\begin{remark}
This is a stronger condition than requiring that for every
\nobreaks{{$\funcname{f} \maps U_1 \to U_2$}} and
\nobreaks{{$\funcname{g} \maps U_2 \to U_1$}} the compositions be
morphisms of $\seqname{C}$.
\end{remark}

Let $E$ be a topological space, $\seqname{C} \defeq (C, \catname{C})$ be
a model space and $(U_i, V_i, \phi_i)$, $i=1,2$, be an m-chart of $E$ in
the coordinate space $\seqname{C}$. Then $(U_1, V_1, \phi_1)$ is
m-compatible with $(U_2, V_2, \phi_2)$ iff either

\begin{enumerate}
\item $U_1$ and $U_2$ are disjoint
\item The transition function
$
  \funcname{t}=\phi_2 \compose \phi^{-1}_1
  \maps \phi_1[U_1 \cap U_2] \to \phi_2[U_1 \cap U_2]
$
is an isomorphism of $\catname C$
\end{enumerate}

$(U_1, V_1, \phi_1)$ is also properly m-compatible with
\nobreaks
  {{
    $(U_2, V_2, \phi_2)$
  }}
in the space $E$ and the coordinate space $\seqname{C}$, and properly
m-compatible in the space $E$ and the coordinate space $\seqname{C}$ with
\nobreaks
  {{
    $(U_2, V_2, \phi_2)$,
  }}
iff
for every open $U'_i \subseteq U_i$

\begin{enumerate}
\item
For every continuous $\funcname{f} \maps U'_1 \to U'_2$,
the composition
$
  \phi_2 \compose \funcname{f} \compose \phi^{-1}_1
  \maps \phi_1[U'_1] \to \phi_2[U'_2]
$
is a morphism of $\seqname{C}$.

\item
For every continuous $\funcname{g} \maps U'_2 \to U'_1$,
the composition
$
  \phi_1 \compose \funcname{g} \compose \phi^{-1}_2
  \maps \phi_2[U'_2] \to \phi_1[U'_1]
$
is a morphism of $\seqname{C}$.
\end{enumerate}

In both cases, the space and coordinate space may be omitted when they
are clear from context.
\end{definition}

\begin{lemma}[M-compatibility]
\label{lem:M-compat}
\mbox{}
\begin{enumerate}
\item
Let $(U_i,V_i,\phi_i)$, $i=1,2$. be a chart of $E$ in the coordinate
space $\seqname{C}$. Then $(U_1, V_1, \phi_1)$ is (properly)
m-compatible with
\nobreaks
  {{
    $(U_2, V_2, \phi_2)$
  }}
in the space $E$ and the coordinate space $\seqname{C}$ iff
$(U_1, V_1, \phi_1)$ is (properly) m-compatible with
\nobreaks
  {{
    $(U_2, V_2, \phi_2)$
  }}
in the space $\Triv{E}$ and the coordinate space $\seqname{C}$.

\begin{proof}
$U_1 \cap U_2$ and the transition function
\nobreaks
  {{
    $
      \funcname{t}=\phi_2 \compose \phi^{-1}_1
      \maps \phi_1[U_1 \cap U_2] \to \phi_2[U_1 \cap U_2]
$
}}
is the same for the case of $E$ and the case of $\Triv{E}$.

Any
$U'_i \subseteq U_i$ is a model neighborhood of $\Triv{E}$ iff it is open.

Any $\funcname{f} \maps U'_1 \to U'_2$ is a morphism of $\Triv{E}$ iff
it is continuous.

Any $\funcname{g} \maps U'_2 \to U'_1$ is a morphism of $\Triv{E}$ iff
it is continuous.

The compositions
$
  \phi_2 \compose \funcname{f} \compose \phi^{-1}_1
  \maps \phi_1[U'_1] \to \phi_2[U'_2]
$
and
$
  \phi_1 \compose \funcname{g} \compose \phi^{-1}_2
  \maps \phi_2[U'_2] \to \phi_1[U'_1]
$
are the same for the case of $E$ and the case of $\Triv{E}$.
\end{proof}
\end{enumerate}
\end{lemma}

\begin{corollary}
\label{cor:M-compat}
Let $(U_i,V_i,\phi_i)$, $i=1,2$. be a chart of $E$ in the coordinate
space $\seqname{C}$ snd let $(U'_i,V'_i,\phi'_i)$ be a subchart. If
$(U_1,V_1,\phi_1)$ is (properly) m-compatible with $(U_2,V_2,\phi_2)$
in the space $E$ and the coordinate space $\seqname{C}$ then
$(U'_1,V'_1,\phi'_1)$ is (properly) m-compatible with
$(U'_2,V'_2,\phi'_2)$.

\begin{proof}
By \cref{lem:M-compat} above, $(U'_1,V'_1,\phi'_1)$ is (properly)
m-compatible with $(U'_2,V'_2,\phi'_2)$ in the space $E$ and the
coordinate space $\seqname{C}$ iff $(U'_1,V'_1,\phi'_1)$ is (properly)
m-compatible with $(U'_2,V'_2,\phi'_2)$ in the space $\Triv{E}$ and the
coordinate space $\seqname{C}$.

The result follows from \cref{lem:M-compat}.
\end{proof}
\end{corollary}

\begin{lemma}[Symmetry of m-compatibility]
\label{lem:M-compatsym}
Let $(U_i, V_i, \phi_i)$, $i=1,2$, be an m-chart of
$(E, \catname{E})$ in the coordinate space $(C,\catname{C})$. Then
\nobreaks
  {{
    $(U_1, V_1, \phi_1)$
  }}
is (properly) m-compatible with
\nobreaks
  {{
    $(U_2, V_2, \phi_2)$
  }}
in the space $\seqname{E}$ and the coordinate space $\seqname{C}$ iff
\nobreaks
  {{
    $(U_2, V_2, \phi_2)$
  }}
is (properly) m-compatible with
\nobreaks
  {{
    $(U_1, V_1, \phi_1)$.
  }}

\begin{proof}
It suffices to prove the implication in only one direction.
\begin{enumerate}
\item $U_1 \cap U_2 = U_2 \cap U_1$.
\item Since the transition function
\nobreaks
  {{
    $
      \funcname{t}=\phi_2 \compose \phi^{-1}_1
      \maps \phi_1[U_1 \cap U_2] \to \phi_2[U_1 \cap U_2]
    $
  }}
is an isomorphism of $\catname C$, so is
\nobreaks
  {{
    $
      \funcname{t}^{-1} = \phi_1 \compose \phi_2^{-1}
      \maps \phi_2[U_1 \cap U_2] \to \phi_1[U_1 \cap U_2]
    $.
  }}
\item The constraint for proper maximality is intrinsically symmetrical.
\end{enumerate}
\end{proof}

Let $(U_i, V_i, \phi_i)$, $i=1,2$, be an m-chart of $E$ in the
coordinate space $(C,\catname{C})$. Then $(U_1, V_1, \phi_1)$ is
(properly) m-compatible with $(U_2, V_2, \phi_2)$ in the space $E$ and
the coordinate space $\seqname{C}$ iff iff $(U_2, V_2, \phi_2)$ is
(properly) m-compatible with $(U_1, V_1, \phi_1)$.

\begin{proof}
The same proof applies as above.
\end{proof}
\end{lemma}

\begin{lemma}[M-compatibility of subcharts]
\label{lem:M-compatsub}
Let $(U_i, V_i, \phi_i)$, $i=1,2$, be an m-chart of $(E, \catname{E})$
in the coordinate space $(C,\catname{C})$,
\nobreaks
  {{
    $(U_1, V_1, \phi_1)$
  }}
be (properly) m-compatible with
\nobreaks
  {{
    $(U_2, V_2, \phi_2)$
  }}
in the space $\seqname{E}$ and the coordinate space $\seqname{C}$ and
\nobreaks
  {{
    $(U'_i, V'_i, \phi'_i)$
  }}
be a subchart. Then
\nobreaks
  {{
    $(U'_1, V'_1, \phi'_1)$
  }}
is (properly) m-compatible with
\nobreaks
  {{
    $(U'_2, V'_2, \phi'_2)$.
  }}

\begin{proof}
Since subcharts are charts, $U'_i$ and $V'_i$ are model neighborhoods.
If \nobreaks{{$U_1 \cap U_2 = \emptyset$}} then
\nobreaks{{$U'_1 \cap U'_2 = \emptyset$}} and thus
\nobreaks{{$(U'_1, V'_1, \phi'_1)$}} is m-compatible with
\nobreaks{{$(U'_2, V'_2, \phi'_2)$}}.  Otherwise, the transition
function
\nobreaks
  {{
    $
      t^1_2 \defeq
      \phi_2 \compose \phi^{-1}_1
      \maps \phi_1[U_1 \cap U_2] \to \phi_2[U_1 \cap U_2]
    $
  }}
is an isomorphism of $\catname{C}$ and hence
\nobreaks
  {{
    $
      t^1_2 \maps \phi_1[U'_1 \cap U'_2] \toiso \phi_2[U'_1 \cap U'_2]
    $
 }}
is an isomorphism of $\catname{C}$.

If $(U_1, V_1, \phi_1)$ is properly m-compatible with
$(U_2, V_2, \phi_2)$,
$U''_i \subseteq U'_i$, $i=1,2$, is a model neighbothood and
$\funcname{f} \maps U''_1 \to U''_2$ is a morphism of $\seqname{E}$,
then each $U''_i \subseteq U_i$ and thus by \cref{def:M-compat} the
composition
$
  \phi_2 \compose \funcname{f} \compose \phi^{-1}_1
  \maps \phi_1[U''_1] \to \phi_2[U''_2]
$
is a morphism of $\seqname{C}$.
\end{proof}

Let $(U_i, V_i, \phi_i)$, $i=1,2$, be an m-chart of $E$ in the
coordinate space $(C,\catname{C})$, $(U'_i, V'_i, \phi'_i)$ be a
subchart and $(U_1, V_1, \phi_1)$ be (properly) m-compatible with
\nobreaks
  {{
    $(U_2, V_2, \phi_2)$
  }}
in the space $E$ and the coordinate space $\seqname{C}$. Then
\nobreaks
  {{
    $(U'_1, V'_1, \phi'_1)$
  }}
is (properly) m-compatible with
\nobreaks
  {{
    $(U'_2, V'_2, \phi'_2)$.
  }}

\begin{proof}
By \pagecref{lem:M-compat}, $(U_1, V_1, \phi_1)$ is (properly)
m-compatible with $(U_2, V_2, \phi_2)$ in the space $\Triv{E}$ and the
coordinate space $\seqname{C}$. Then by \cref{lem:M-compatsub} above,
$(U'_1,V'_1,\phi'_1)$ is (properly) m-compatible with
$(U'_2,V'_2,\phi'_2)$ in the space $\Triv{E}$ and the coordinate space
$\seqname{C}$.

The result follows from \cref{lem:M-compat}.
\end{proof}
\end{lemma}

\begin{corollary}[M-compatibility of subcharts]
\label{cor:M-compatsub}
Let $(U, V, \phi)$ be an m-chart of $(E, \catname{E})$ in
the coordinate space $(C,\catname{C})$ and $(U', V', \phi')$ a
subchart. Then $(U', V', \phi')$ is properly m-compatible with
\nobreaks
  {{
    $(U, V, \phi)$
  }}
in the space $\seqname{E}$ and the coordinate space $\seqname{C}$.

\begin{proof}
$(U, V, \phi)$ is m-compatible with itself and is a subchart of
itself,
\end{proof}

Let $(U, V, \phi)$ be an m-chart of $E$ in the coordinate
space $(C,\catname{C})$ and $(U', V', \phi')$ a subchart. Then
$(U', V', \phi')$ is m-compatible with
\nobreaks
  {{
    $(U,V,\phi)$
  }}
in the space $E$ and the coordinate space $\seqname{C}$.

\begin{proof}
$(U, V, \phi)$ is m-compatible with itself and is a subchart of
itself,
\end{proof}
\end{corollary}

\begin{definition}[Covering by m-charts]
Let $\seqname{A}$ be a set of charts of the topological space $E$ in
the coordinate space
$\seqname{C}=(C,\catname{C})$.
$\seqname{A}$ covers $E$ iff $\pi_1[\seqname{A}]$ covers $E$.

Let $\seqname{A}$ be a set of charts of the model space
$\seqname{E}=(E, \catname{E})$ in the coordinate space
$\seqname{C}=(C,\catname{C})$.  $\seqname{A}$ covers $\seqname{E}$ iff
$\pi_1[\seqname{A}]$ covers $E$.
\end{definition}

\section{M-atlases}
\label{sec:M-atlas}
A set of charts can be atlases for different coordinate model spaces
even if it is for the same total model space. In order to aggregate
atlases into categories, there must be a way to distinguish them.
Including the two\footnote{
The spaces are redundant, but convenient.}
spaces in the definitions of the categories serves the purpose.

\begin{definition}[M-atlases]
\label{def:M-atlas}
Let $\seqname{A}$ be a set of mutually (properly) m-compatible m-charts
of $\seqname{E}=(E, \catname{E})$ in the coordinate space
$\seqname{C}=(C,\catname{C})$. Then

\begin{enumerate}
\item
$\seqname{A}$ is an (a proper) m-atlas of $\seqname{E}$ in the
coordinate space $\seqname{C}$, abbreviated \\*
$\isAtl[(proper)]_\Ob(\seqname{A}, \seqname{E}, \seqname{C})$, iff
$\seqname{A}$ covers $E$.

\item
$\seqname{A}$ is a full (proper) atlas of $\seqname{E}$ in the
coordinate space $\seqname{C}$, abbreviated \\*
$\full{\isAtl[(proper)]_\Ob}(\seqname{A}, E, \seqname{C})$, iff
$\seqname{A}$ covers $E$ and $\pi_2[\seqname{A}]$ covers $C$.

\item
The triple $(\seqname{A}, \seqname{E}, \seqname{C})$ refers to
$\seqname{A}$ considered as an m-atlas of $\seqname{E}$ in the
coordinate space $\seqname{C}$.

\item
$\seqname{E}$ is the total model space or total space for the atlas

\item
$\seqname{C}$ is the coordinate model space or coordinate space for the
atlas.
\end{enumerate}

Let $E$ be a topological space, $\seqname{C}=(C,\catname{C})$ a model
space and $\seqname{A}$ be a set of mutually m-compatible m-charts of
$E$ in the coordinate space $\seqname{C}$. Then

\begin{enumerate}[resume]
\item
$\seqname{A}$ is an m-atlas of $E$ in the coordinate space
$\seqname{C}$, abbreviated as
$\isAtl_\Ob(\seqname{A}, E, \seqname{C})$, iff $\seqname{A}$ is an
m-atlas of $\Triv{E}$ in the coordinate space $\seqname{C}$.

\item
$\seqname{A}$ is a full atlas of $E$ in the coordinate space
$\seqname{C}$, abbreviated
$\full{\isAtl_\Ob}(\seqname{A}, E, \seqname{C})$, iff $\seqname{A}$ is a
full atlas of $\Triv{E}$ in the coordinate space $\seqname{C}$.

\item
The triple $(\seqname{A}, E, \seqname{C})$ refers to $\seqname{A}$
considered as an m-atlas of $E$ in the coordinate space $\seqname{C}$.

\item
$E$ is the total space for the atlas.

\item
$\seqname{C}$ is the coordinate model space or coordinate space for the
atlas.
\end{enumerate}

Let $\seqname{A}$ be an (a proper) m-atlas of $\seqname{E}$ in the
coordinate space $\seqname{C}$.  $\seqname{A}$ is fine grained iff
whenever $(U,V,\phi) \in \seqname{A}$, every subchart of
$(U,V,\phi)$ is also an atlas of $\seqname{A}$.

\begin{remark}
$\seqname{A}$ can be a fine graind (proper) atlas of $\seqname{E}$ in
the coordinate space $\seqname{C}$ even if $\seqname{C}$ is not fine
grained.
\end{remark}

Let $\seqname{A}$ be an (a proper) m-atlas of $E$ in the coordinate
space $\seqname{C}$.  $\seqname{A}$ is fine grained iff whenever
$(U,V,\phi) \in \seqname{A}$, every subchart of $(U,V,\phi)$ is also an
atlas of $\seqname{A}$.

By abuse of language we write $U \in \seqname{A}$ for
$U \in \pi_1[\seqname{A}]$.
\end{definition}

\begin{lemma}[M-atlases]
\label{lem:M-atlas}
Let $\seqname{A}$ be an m-atlas of $\seqname{E} \defeq (E, \catname{E})$
in the coordinate space $\seqname{C} \defeq (C,\catname{C})$. If
$\seqname{C}$ is fine grained then $\seqname{A}$ is an m-atlas of $E$ in
the coordinate space $\seqname{C}$.

\begin{proof}
Let $(U,V,\phi) \in \seqname{A}$. Then $(U,V,\phi)$ is an m-chart of $E$
in the coordinate space $\seqname{C}$ by \pagecref{cor:M-chart}.
\end{proof}
\end{lemma}

\begin{definition}[M-compatibility with m-atlases]
\label{def:M-CompatATL}
An m-chart $(U, V, \phi)$ of $(E, \catname{E})$ is (properly)
m-compatible with an m-atlas $\seqname{A}$ in the space $\seqname{E}$
and the coordinate space $\seqname{C}$ iff it is (properly) m-compatible
in the space $\seqname{E}$ and the coordinate space $\seqname{C}$
with every chart in $\seqname{A}$.

An m-chart $(U, V, \phi)$ of $E$ is (properly) m-compatible with an
m-atlas $\seqname{A}$ in the space $\seqname{E}$ and the coordinate
space $\seqname{C}$ iff it is (properly) m-compatible in the space
$\seqname{E}$ and the coordinate space $\seqname{C}$ with every chart in
$\seqname{A}$.
\end{definition}

\begin{lemma}[M-compatibility with m-atlases]
\label{lem:M-CompatATL}
Let $\seqname{A}$ be an m-atlas of $\seqname{E} \defeq (E, \catname{E})$
in the coordinate space $\seqname{C} \defeq (C,\catname{C})$. If
$\seqname{C}$ is fine grained then $\seqname{A}$ is an m-atlas of $E$ in
the coordinate space $\seqname{C}$.

\begin{proof}
\end{proof}
\end{lemma}

\begin{lemma}[M-compatibility of subcharts with atlases]
\label{lem:M-CompatsubATL}
Let $\seqname{A}$ be an (a proper) m-atlas of $(E, \catname{E})$ in the
coordinate space $(C,\catname{C})$ and $(U_1, V_1, \phi_1)$ an m-chart
in $\seqname{A}$. Then any subchart of $(U_1, V_1, \phi_1)$ is
properly m-compatible with $\seqname{A}$ in the space $\seqname{E}$
and the coordinate space $\seqname{C}$.

\begin{proof}
Let $(U'^1, V'^1, \phi'^1)$ be a subchart of $(U^1, V^1, \phi^1)$ and
$(U_2, V_2, \phi_2)$ be another chart in $\seqname{A}$.  Then
$(U_1, V_1, \phi_1)$ is properly compatible with $(U_2, V_2, \phi_2)$
by \pagecref{def:M-atlas}, $(U_2, V_2, \phi_2)$ is properly
m-compatible with itself and $(U_2, V_2, \phi_2)$ is a subchart of
itself. $(U'^1, V'^1, \phi'^1)$ is properly m-compatible with
$(U_2, V_2, \phi_2)$ by \pagecref{lem:M-compatsub}. Since
$(U_2, V_2, \phi_2)$ is arbitrary, $(U'^1, V'^1, \phi'^1)$ is properly
m-compatible with $\seqname{A}$ by \pagecref{def:M-CompatATL}.
\end{proof}

Let $\seqname{A}$ be an m-atlas of $E$ in the coordinate space
$(C,\catname{C})$ and $(U_1, V_1, \phi_1)$ an m-chart in $\seqname{A}$.
Then any subchart of $(U_1, V_1, \phi_1)$ is properly m-compatible with
$\seqname{A}$ in the space $E$ and the coordinate space $\seqname{C}$.

\begin{proof}
The same proof applies in as above.
\end{proof}
\end{lemma}

\begin{lemma}[Extensions of m-atlases]
\label{lem:M-ATLextensions}
Let $\seqname{E}=(E,\catname{E})$ and $\seqname{C}=(C,\catname{C})$ be
model spaces, $\seqname{A}$ an m-atlas of $\seqname{E}$ in the
coordinate space $\seqname{C}$, and $(U_i,V_i,\phi_i)$, $i=1,2$, an
m-chart of $\seqname{E}$ in the coordinate space $\seqname{C}$
(properly) m-compatible with $\seqname{A}$ in the space $\seqname{E}$
and the coordinate space $\seqname{C}$.  Then $(U_1, V_1, \phi_1)$ is
(properly) m-compatible with $(U_2, V_2, \phi_2)$.

\begin{proof}
If $U_1 \cap U_2 = \emptyset$ then $(U_1, V_1, \phi_1)$ is m-compatible
with $(U_2 ,V_2, \phi_2)$. Otherwise, $\phi_1$ is a model homeomorphism,
$U_1 \cap U_2$ is a model neighborhood of $\seqname{E}$,
$\phi_1[U_1 \cap U_2]$ is a model neighborhood of $\seqname{C}$,
$\phi_2[U_1 \cap U_2]$ is a model neighborhood of $\seqname{C}$ and
$
  \phi_2 \compose \phi^{-1}_1
  \maps \phi_1[U_1 \cap U_2] \toiso \phi_2[U_1 \cap U_2]
$
is a homeomorphism.  It remains to show that
$\phi_2 \compose \phi^{-1}_1 \restrictto_{\phi_1[U_1 \cap U_2]}$ is a
model homeomorphism, Let
$(U_\alpha,V_\alpha,\phi_\alpha)$, $\alpha \prec \Alpha$, be charts in
$\seqname{A}$ such that
$U_1 \cap U_2 \subseteq \union[\alpha \prec \Alpha]{U_\alpha}$
and
\nobreaks
  {{
    $
      \uquant
        {
          \alpha \prec \Alpha
        }
        {
          U_1 \cap U_2 \cap U_\alpha \neq \emptyset
        }
    $.
  }}
Since the charts are m-compatible with
\nobreaks
  {{
    $(U_\alpha, V_\alpha, \phi_\alpha)$,
  }}
the compositions \\*
\nobreaks
  {{
    $
      \phi_2 \compose \phi^{-1}_\alpha
      \maps \phi_\alpha[U_1 \cap U_2 \cap U_\alpha]
      \to \phi_2[U_1 \cap U_2 \cap U_\alpha]
    $
  }}
and \\*
\nobreaks
  {{
    $
      \phi_\alpha \compose \phi^{-1}_1
      \maps \phi_1[U_1 \cap U_2 \cap U_\alpha]
      \to \phi_\alpha[U_1 \cap U_2 \cap U_\alpha]
    $
  }}
are isomorphisms of $\seqname{C}$ and thus the composition
\nobreaks
  {{
    $
      \phi_2 \compose \phi^{-1}_1
      \maps \phi_1[U_1 \cap U_2 \cap U_\alpha]
      \to \phi_2[U_1 \cap U_2 \cap U_\alpha]
      =
    $
  }} \\*
\nobreaks
  {{
    $
      \phi_2 \compose \phi^{-1}_\alpha \compose
      \phi_\alpha \compose \phi^{-1}_1
      \maps \phi_1[U_1 \cap U_2 \cap U_\alpha]
      \to \phi_2[U_1 \cap U_2 \cap U_\alpha]
    $
  }}
is an isomorphisms of $\seqname{C}$. Then by
{
  \showlabelsinline
  \cref{mod:sheaf}
},
the restricted sheaf property, of \pagecref{def:model},
$\phi_2 \compose \phi^{-1}_1$
is an isomorphism of $\seqname{C}$.

If each $(U_i,V_i,\phi_i)$ is properly m-compatible with $\seqname{A}$
then let
\nobreaks
  {{
    $\funcname{f} \maps U'_1 \to U'_2$
  }}
and
\nobreaks
  {{
    $\funcname{g} \maps U'_2 \to U'_1$
  }}
be morphisms of $\seqname{E}$ and $U'_i \subseteq U_i$, $i=1,2$.
Let $u$ be an arbitary point in $U_1$ and
\nobreaks
  {{
    $(U_u, V_u, \phi_u) \in \seqname{A}$
  }}
be an arbitray chart at $u$. Since
\nobreaks
  {{
    $(U_1,V_1,\phi_1)$
  }}
is properly m-compatible with $\seqname{A}$,
it is properly m-compatible with
\nobreaks
  {{
    $(U_u, V_u, \phi_u) \in \seqname{A}$,
  }}
thus the compositions
\nobreaks
  {{
    $
      \phi_2 \compose \funcname{f} \compose \phi^{-1}_u
      \maps \phi^{-1}_u[U_u] \to \phi^{-1}_2[U'_2]
    $
  }}
and
\nobreaks
  {{
    $
      \phi_u \compose \phi^{-1}_1
      \maps \phi^{-1}_1[U_u] \to \phi^{-1}_u[U_u]
    $
  }}
are morphisms of $\seqname{C}$, their composition
\nobreaks
  {{
    $
      \phi_2 \compose \funcname{f} \compose \phi^{-1}_u \compose
      \phi_u \compose \phi^{-1}_1
      \maps \phi^{-1}_1[U_u] \to \phi^{-1}_2[U'_2] =
    $
  }}
\nobreaks
  {{
    $
      \phi_2 \compose \funcname{f} \compose \phi^{-1}_1
      \maps \phi^{-1}_1[U_u] \to \phi^{-1}_2[U'_2]
    $
  }}
is a morphism of $\seqname{C}$ and by
{
  \showlabelsinline
  \cref{mod:sheaf}
},
the restricted sheaf property, of \pagecref{def:model}, \\*
\nobreaks
  {{
    $
      \phi_2 \compose \funcname{f} \compose \phi^{-1}_1
      \maps \phi^{-1}_1[U'_1] \to \phi^{-1}_2[U'_2]
    $
  }}
is a morphism of $\seqname{C}$.

The same proof applies in the opposite direction.
\end{proof}

Let $E$ be a topological space, $\seqname{C}=(C,\catname{C})$ a model
space, $\seqname{A}$ an (a proper) m-atlas of $E$ in the coordinate
space $\seqname{C}$ and $(U_i, V_i, \phi_i)$, $i=1,2$, a chart
(properly) m-compatible with $\seqname{A}$ in the space $E$ and the
coordinate space $\seqname{C}$.  Then $(U_1, V_1, \phi_1)$ is (properly)
m-compatible with $(U_2, V_2, \phi_2)$.

\begin{proof}
The same proof applies as above.
\end{proof}
\end{lemma}

\begin{definition}[Maximal (semi-maximal) m-atlases]
\label{def:M-ATLmax}
Let $\seqname{E}$ and $\seqname{C}$ be model spaces and $\seqname{A}$ an
m-atlas of $\seqname{E}$ in the coordinate space $\seqname{C}$.
$\seqname{A}$ is a (proper) maximal m-atlas of $\seqname{E}$ in the
coordinate space $\seqname{C}$, abbreviated
\nobreaks
  {{
    $
      \maximal{\isAtl[(proper)]_\Ob}
      (\seqname{A}, \seqname{E}, \seqname{C})
    $,
  }}
iff $\seqname{A}$ cannot be extended by adding an additional (properly)
m-compatible chart. $\seqname{A}$ is a (proper) semi-maximal m-atlas of
$\seqname{E}$ in the coordinate space $\seqname{C}$, abbreviated
\nobreaks
  {{
    $
      \maximal[S-]{\isAtl[(proper)]_\Ob}
      (\seqname{A}, \seqname{E}, \seqname{C})
    $.
  }}
iff whenever

\begin{enumerate}
\item
$(U,V,\phi) \in \seqname{A}$

\item
$U' \subseteq U$, $V' \subseteq V$ and $V'' \objin \catname{C}$ are
model neighborhoods

\item
$\phi[U'] = V'$

\item
$\phi' \maps V' \toiso V''$ is an isomorphism of $\catname{C}$
\end{enumerate}

then $(U', V'', \phi' \compose \phi \restrictto_{U'}) \in \seqname{A}$.

\begin{equation}
  \maxfull{\isAtl[(proper)]_\Ob}(\seqname{A}, \seqname{E}, \seqname{C}) \defeq
    \full{\isAtl[(proper)]_\Ob}(\seqname{A}, \seqname{E}, \seqname{C}) \land
    \maximal{\isAtl[(proper)]_\Ob}(\seqname{A}, \seqname{E}, \seqname{C})
\end{equation}
\begin{equation}
  \maxfull[S-]{\isAtl[(proper)]_\Ob}(\seqname{A}, \seqname{E}, \seqname{C}) \defeq
    \full{\isAtl[(proper)]_\Ob}(\seqname{A}, \seqname{E}, \seqname{C}) \land
    \maximal[S-]{\isAtl[(proper)]_\Ob}(\seqname{A}, \seqname{E}, \seqname{C})
\end{equation}

Let $E$ be a topological space, $\seqname{C}=(C,\catname{C})$ be a model
space and $\seqname{A}$ be an m-atlas of $E$ in the coordinate space
$\seqname{C}$. $\seqname{A}$ is a (properly) maximal m-atlas of $E$ in
the coordinate space $\seqname{C}$, abbreviated
$\maximal{\isAtl[(proper)]_\Ob}(\seqname{A}, E, \seqname{C})$, iff
$\seqname{A}$ is an m-atlas that cannot be extended by adding an
additional (properly) m-compatible chart.  $\seqname{A}$ is a (properly)
semi-maximal m-atlas of $E$ in the coordinate space $\seqname{C}$,
abbreviated
$\maximal[S-]{\isAtl[(proper)]_\Ob}(\seqname{A}, E, \seqname{C})$, iff
whenever

\begin{enumerate}
\item
$(U,V,\phi) \in \seqname{A}$

\item
$U' \subseteq U$, $V' \subseteq V$ and $V'' \objin \catname{C}$ are
model neighborhoods

\item
$\phi[U'] = V'$

\item
$\phi' \maps V' \toiso V''$ is an isomorphism of $\catname{C}$
\end{enumerate}

then $(U', V'', \phi' \compose \phi \restrictto_{U'}) \in \seqname{A}$.

\begin{remark}
There is no sheaf condition\footnote{
  However, note \cref{mod:sheaf} (restricted sheaf condition) of
  {
    \showlabelsinline
    \pagecref{def:model}.
  }
};
the union of a set of coordinate patches in the maximal atlas whose
coordinate functions match on the intersections need not be a coordinate
patch in the atlas.
\end{remark}

\begin{equation}
  \maxfull{\isAtl[(proper)]_\Ob}(\seqname{A}, E, \seqname{C}) \defeq
    \full{\isAtl[(proper)]_\Ob}(\seqname{A}, E, \seqname{C}) \land
    \maximal{\isAtl[(proper)]_\Ob}(\seqname{A}, E, \seqname{C})
\end{equation}
\begin{equation}
  \maxfull[S-]{\isAtl[(proper)]_\Ob}(\seqname{A}, E, \seqname{C}) \defeq
    \full{\isAtl[(proper)]_\Ob}(\seqname{A}, E, \seqname{C}) \land
    \maximal[S-]{\isAtl[(proper)]_\Ob}(\seqname{A}, E, \seqname{C})
\end{equation}
\end{definition}

\begin{lemma}[Maximal m-atlases are semi-maximal m-atlases]
\label{lem:M-ATLmax}
Let $\seqname{E}=(E, \catname{E})$ and $\seqname{C}=(C,\catname{C})$ be
model spaces and $\seqname{A}$ a maximal m-atlas of $\seqname{E}$ in the
coordinate space $\seqname{C}$. Then $\seqname{A}$ is a semi-maximal
m-atlas of $\seqname{E}$ in the coordinate space $\seqname{C}$.

\begin{proof}
Let $(U,V,\phi) \in \seqname{A}$, $U' \subseteq U, V' \subseteq V$
and $V'' \objin \catname{C}$ be model neighborhoods, $\phi[U'] = V'$ and
$\phi' \maps V' \toiso V''$ be an isomorphism of $\catname{C}$.
$(U', V', \phi)$ is a subchart of $(U,V,\phi)$ and by
\pagecref{lem:M-Compat} is m-compatible with the charts of
$\seqname{A}$.  Since $\phi'$ is a model homeorphism,
$(U', V'', \phi' \compose \phi)$ is m-compatible with the charts of
$\seqname{A}$.  Since $\seqname{A}$ is maximal,
$(U', V'', \phi' \compose \phi)$ is a chart of $\seqname{A}$.
\end{proof}

Let $E$ be a topological space, $\seqname{C}=(C,\catname{C})$ a model
space and $\seqname{A}$ a maximal m-atlas of $E$ in the coordinate space
$\seqname{C}$. Then $\seqname{A}$ is a semi-maximal m-atlas of $E$ in
the coordinate space $\seqname{C}$.

\begin{proof}
Let $(U,V,\phi) \in \seqname{A}$, $U' \subseteq U$ open, $V' \subseteq V$
and $V'' \objin \catname{C}$ be model neighborhoods, $\phi[U'] = V'$ and
$\phi' \maps V' \toiso V''$ be an isomorphism of $\catname{C}$.
$(U', V', \phi)$ is a subchart of $(U,V,\phi)$ and by
\pagecref{lem:M-Compat} is m-compatible with the charts of
$\seqname{A}$.  Since $\phi'$ is a model homeorphism,
$(U', V'', \phi' \compose \phi)$ is m-compatible with the charts of
$\seqname{A}$.  Since $\seqname{A}$ is maximal,
$(U', V'', \phi' \compose \phi)$ is a chart of $\seqname{A}$.
\end{proof}
\end{lemma}

\begin{theorem}[Existence and uniqueness of maximal m-atlases]
\label{the:M-ATLmax}
Let $\seqname{A}$ be an (a proper) m-atlas of
$\seqname{E} \defeq (E, \catname{E})$ in the coordinate space
$\seqname{C} \defeq (C,\catname{C})$. Then there exists a unique
(proper) maximal m-atlas
\nobreaks
  {{
    $\maximal{\Atlas[(proper)]}(\seqname{A}, \seqname{E}, \seqname{C})$
  }}
of $\seqname{E}$ in the coordinate space $\seqname{C}$ m-compatible with
$\seqname{A}$.

\begin{proof}
Let $\seqname{P}$ be the set of all m-atlases of $\seqname{E}$ in the
coordinate space $\seqname{C}$ containing $\seqname{A}$ and (properly)
m-compatible in the coordinate space $\seqname{C}$ with all of the
m-charts in $\seqname{A}$. Let $\maximal{\seqname{P}}$ be a maximal
chain of $\seqname{A}$. Then
\nobreaks
  {{
    $
      \maximal{\Atlas[(proper)]}(\seqname{A}, \seqname{E}, \seqname{C})
      \defeq \union{\maximal{\seqname{P}}}
    $
  }}
is a (proper) maximal m-atlas of $\seqname{E}$ in the coordinate space
$\seqname{C}$ m-compatible with $\seqname{A}$. Uniqueness follows from
\pagecref{lem:M-ATLextensions}.
\end{proof}
\end{theorem}

\begin{corollary}[Existence and uniqueness of maximal m-atlases]
\label{cor:M-ATLmax}
Let $\seqname{A}$ be an (a proper) m-atlas of $E$
in the coordinate space $\seqname{C} \defeq (C,\catname{C})$. Then there
exists a unique (proper) maximal m-atlas
$\maximal{\Atlas[(proper)]}(\seqname{A}, E, \seqname{C})$ of $E$
in the coordinate space $\seqname{C}$ m-compatible with $\seqname{A}$.

\begin{proof}
$
  \maximal{\Atlas[(proper)]}(\seqname{A}, E, \seqname{C}) \defeq
  \maximal{\Atlas[(proper)]}(\seqname{A}, \Triv{E}, \seqname{C})
$.
\end{proof}
\end{corollary}

\begin{definition}[Sets of m-atlases]
\label{def:M-ATLsets}

Let $\seqname{E}$ and $\seqname{C}$ be model spaces.

\begin{multline}
\maxfull[*]{\Atl[(proper)]_\Ob(\seqname{E}, \seqname{C})} \defeq
\set
{{
   (\seqname{A}, \seqname{E}, \seqname{C})
}}%
[
  {
    \maxfull[*]
    {
      \isAtl[(proper)]_\Ob
      (\seqname{A}, \seqname{E}, \seqname{C})
    }
  }
]*
\end{multline}

Let $\seqname{E}$ and $\seqname{C}$ be sets of model spaces.

\begin{multline}
\maxfull[*]{\Atl[(proper)]_\Ob(\seqname{E}, \seqname{C})} \defeq
\set
{{
   \bigl ( \seqname{A}, (E,\catname{E}), (C,\catname{C}) \bigr )
}}%
[
  {(E,\catname{E}) \in \seqname{E}},
  {(C,\catname{C}) \in \seqname{C}},
  {
    \maxfull[*]
    {
      \isAtl[(proper)]_\Ob
      \bigl ( \seqname{A}, (E,\catname{E}), (C,\catname{C}) \bigr )
    }
  }
]*
\end{multline}

Let $\seqname{E}$ be a set of topological spaces and $\seqname{C}$ a
set of model spaces. Then
\begin{multline}
\maxfull[*]{\Atl[(proper)]_\Ob (\seqname{E}, \seqname{C})} \defeq
\set
{{
   \bigl ( \seqname{A}, E, (C,\catname{C}) \bigr )
}}%
[
  {E \in \seqname{E}},
  {(C,\catname{C}) \in \seqname{C}},
  {
    \maxfull[*]
    {
      \isAtl[(proper)]_\Ob
      \bigl ( \seqname{A}, E, (C,\catname{C}) \bigr )
    }
  }
]*
\end{multline}
\end{definition}

\begin{lemma}[Sets of m-atlases]
\label{lem:M-ATLsets}

Let $\seqname{E}$ and $\seqname{C}$ be sets of model spaces.

\begin{equation}
  \maxfull{\Atl[(proper)]_\Ob(\seqname{E}, \seqname{C})} \subseteq
  \maxfull[*x]{\Atl[(proper)]_\Ob(\seqname{E}, \seqname{C})}
\end{equation}

\begin{equation}
  \maxfull[**]{\Atl[(proper)]_\Ob(\seqname{E}, \seqname{C})} \subseteq
  \Atl[(proper)]_\Ob(\seqname{E}, \seqname{C})
\end{equation}

\begin{equation}
  \maximal[S-max-full~(max-full)]{\Atl[(proper)]_\Ob(\seqname{E}, \seqname{C})}
  \subseteq
  \maximal[S-max~(max)]{\Atl[(proper)]_\Ob(\seqname{E}, \seqname{C})}
\end{equation}

\begin{proof}
The result follows from
{
  \showlabelsinline
  \pagecref{def:M-atlas},
}
\pagecref{def:M-ATLmax}
and \cref{def:M-ATLsets} above.
\end{proof}

\begin{equation}
  \maxfull{\Atl[(proper)]_\Ob(\seqname{E}, \seqname{C})} \subseteq
  \maxfull[S-]{\Atl[(proper)]_\Ob(\seqname{E}, \seqname{C})} \subseteq
  \full{\Atl[(proper)]_\Ob(\seqname{E}, \seqname{C})}
\end{equation}

\begin{proof}
The result follows from \pagecref{def:M-ATLmax} and \cref{def:M-ATLsets}
above.
\end{proof}
\end{lemma}

\begin{corollary}[Sets of m-atlases]
\label{cor:M-ATLsets}

Let $\seqname{E}$ be a set of topological spaces and $\seqname{C}$ a
set of model spaces. Then

\begin{equation}
  \maxfull{\Atl[(proper)]_\Ob(\seqname{E}, \seqname{C})} \subseteq
  \maxfull[*x]{\Atl[(proper)]_\Ob(\seqname{E}, \seqname{C})}
\end{equation}

\begin{equation}
  \maxfull[**]{\Atl[(proper)]_\Ob(\seqname{E}, \seqname{C})} \subseteq
  \Atl[(proper)]_\Ob(\seqname{E}, \seqname{C})
\end{equation}

\begin{equation}
  \maximal[S-max-full~(max-full)]{\Atl[(proper)]_\Ob(\seqname{E}, \seqname{C})}
  \subseteq
  \maximal[S-max~(max)]{\Atl[(proper)]_\Ob(\seqname{E}, \seqname{C})}
\end{equation}

\begin{equation}
  \maxfull{\Atl[(proper)]_\Ob(\seqname{E}, \seqname{C})} \subseteq
  \maxfull[S-]{\Atl[(proper)]_\Ob(\seqname{E}, \seqname{C})} \subseteq
  \full{\Atl[(proper)]_\Ob(\seqname{E}, \seqname{C})}
\end{equation}

\begin{proof}
The result follows from \pagecref{def:M-atlas},  \pagecref{def:M-ATLmax}
and \cref{def:M-ATLsets} above.
\end{proof}
\end{corollary}

\section{M-atlas (near) morphisms and functors}
\label{sec:M-ATLmorph}
This section introduces a taxonomy for morphisms between m-atlases,
defines categories of m-atlases, defines functors amomg them and proves
some basic reasults.

\subsection{M-atlas (near) morphisms}
\label{sub:M-ATLmorph}
This subsection introduces the notions of m-atlas morphisms, m-atlas
near morphisms, strict m-atlas (near) morphisms and semi-strict m-atlas
(near) morphisms. Only (strict, semi-strict) m-atlas morphisms are
needed for the exposition of local coordinate spaces, as m-atlas near
morphisms between maximal atlases are proven to be m-atlas morphisms.

\subsubsection{Definitions of m-atlas (near) morphisms}
\begin{definition}[M-atlas morphisms for model spaces]
\label{def:M-ATLmorph}
Let $\catname{E}^i$, $\catname{C}^i$, $i=1,2$, be model categories,
$\seqname{E}^i \objin \catname{E}^i$,
$\seqname{C}^i \objin \catname{C}^i$, $\seqname{A}^i$ be an m-atlas of
$\seqname{E}^i$ in the coordinate space $\seqname{C}^i$ and
$
  \funcseqname{f} \defeq
  (
    \funcname{f}_0 \maps \seqname{E}^1 \to \seqname{E}^2,
    \funcname{f}_1 \maps \seqname{C}^1 \to \seqname{C}^2
  )
$
a pair of model functions.

$\funcseqname{f}$ is an
$\seqname{E}^1$-$\seqname{E}^2$ m-atlas morphism of $\seqname{A}^1$ to
$\seqname{A}^2$ in the coordinate spaces $\seqname{C}^1$,
$\seqname{C}^2$, abbreviated as
$
  \isAtl_\Ar
    (
      \seqname{A}^1,
      \seqname{E}^1,
      \seqname{C}^1,
      \seqname{A}^2,
      \seqname{E}^2,
      \seqname{C}^2,
      \funcname{f}_0,
      \funcname{f}_1
    )
$
and an $\catname{E}^1$-$\catname{E}^2$ m-atlas morphism of
$\seqname{A}^1$ to $\seqname{A}^2$ in the coordinate model categories
$\catname{C}^1$, $\catname{C}^2$, abbreviated as \\*
$
  \isAtl_\Ar
  (
    \seqname{A}^1,
    \catname{E}^1,
    \catname{C}^1,
    \seqname{A}^2,
    \catname{E}^2,
    \catname{C}^2,
    \funcname{f}_0,
    \funcname{f}_1
  )
$,
iff for any charts
$(U^i, V^i, \phi^i \maps U^i \toiso V^i) \in \seqname{A}^i$, $i=1,2$,
with $I \defeq U^1 \cap \funcname{f}^{-1}_0[U^2] \neq \emptyset$ and any
$u^1 \in I$, there exists a subchart
$
  (U'^1, V'^1, \phi^1 \maps U'^1 \toiso V'^1)
  \in \seqname{A}^1
$
at $u^1$, a model neighborhood $U'^2 \subseteq U^2$ of
$u^2 \defeq \funcname{f}_0(u^1)$ and a chart
$
  (U'^2 ,\hat{V}'^2, \phi'^2 \maps U'^2 \toiso \hat{V}'^2)
  \in \seqname{A}^2
$
at $u^2$ such that
$\funcname{f}_0[U'^1] \subseteq U'^2$ and diagram \eqref{eq:AtlM1}
below is commutative, as shown in \pagecref{fig:AtlM1}.

\begin{equation}
\begin{array}{lll}
\label{eq:AtlM1}
D \defeq & \Biggl \{
\\
  &&
  \set{u^1} \to^{\funcname{i}} U'^1
  \to/{>}->>/^{\phi^1}_\iso V'^1
  \to^{\funcname{f}_1} \hat{V}'^2,
\\
  &&
  U'^1 \to^{\funcname{f}_0}  U'^2 \to/{>}->>/^{\phi'^2}_\iso \hat{V}'^2
\\
& \Biggr \}
\end{array}
\end{equation}

\begin{figure}[ht]
\[ \bfig
\node e1(500,0)[u^1]
\node e2(1500,0)[u^2]
\node u1(-300,-500)[U^1]
\node up1(500,-500)[U'^1]
\node up2(1500,-500)[U'^2]
\node u2(2000,-500)[U^2]
\node v1(-300,-1000)[V^1]
\node vp1(500,-1000)[V'^1]
\node vhp2(1500,-1000)[\hat{V}'^2]
\node v2(2000,-1000)[V^2]
\arrow |a|[e1`e2;\funcname{f}_0]
\arrow |r|/^{ (}->/[e1`up1;\funcname{i}]
\arrow |r|/^{ (}->/[e2`up2;\funcname{i}]
\arrow |r|/>->/[u1`v1;\phi^1]
\arrow |l|//[u1`v1;\iso]
\arrow |a|/_{ (}->/[up1`u1;\funcname{i}]
\arrow |a|[up1`up2;\funcname{f}_0]
\arrow |r|[up1`vp1;{\phi^1}]
\arrow |l|//[up1`vp1;\iso]
\arrow |a|/^{ (}->/[up2`u2;i]
\arrow |r|[up2`vhp2;\phi'^2]
\arrow |l|//[up2`vhp2;\iso]
\arrow |r|[u2`v2;\phi^2]
\arrow |l|//[u2`v2;\iso]
\arrow |a|/_{ (}->/[vp1`v1;\funcname{i}]
\arrow |a|[vp1`vhp2;\funcname{f}_1]
\efig \]
\caption{Completed m-atlas morphism}
\label{fig:AtlM1}
\end{figure}

The triple
$
  \Bigl (
    \funcseqname{f},
    (\seqname{A}^1, \seqname{E}^1, \seqname{C}^1),
    (\seqname{A}^2, \seqname{E}^2, \seqname{C}^2)
  \Bigr )
$
will refer to $\funcseqname{f}$ considered as an
$\seqname{E}^1$-$\seqname{E}^2$ m-atlas morphism of
$\seqname{A}^1$ to $\seqname{A}^2$ in the coordinate spaces
$\seqname{C}^1$, $\seqname{C}^2$.

$\funcseqname{f}$ is also a constrained
$\seqname{E}^1$-$\seqname{E}^2$ m-atlas morphism of $\seqname{A}^1$ to
$\seqname{A}^2$ in the coordinate spaces $\seqname{C}^1$,
$\seqname{C}^2$, abbreviated as
$
  \isAtl[const]_\Ar
    (
      \seqname{A}^1,
      \seqname{E}^1,
      \seqname{C}^1,
      \seqname{A}^2,
      \seqname{E}^2,
      \seqname{C}^2,
      \funcname{f}_0,
      \funcname{f}_1
    )
$
and a constrained $\catname{E}^1$-$\catname{E}^2$ m-atlas morphism of
$\seqname{A}^1$ to $\seqname{A}^2$ in the coordinate model categories
$\catname{C}^1$, $\catname{C}^2$, abbreviated as
$
  \isAtl[const]_\Ar
  (
    \seqname{A}^1,
    \catname{E}^1,
    \catname{C}^1,
    \seqname{A}^2,
    \catname{E}^2,
    \catname{C}^2,
    \funcname{f}_0,
    \funcname{f}_1
  )
$,
iff $\funcname{f}_0$ and $\funcname{f}_1$ are constrained, i.e.,
$\funcname{f}_0[\seqname{E}^1]$ is contained in a model neighborhood of
$\seqname{E}^2$ and $\funcname{f}_1[\seqname{C}^1]$ is contained in a
model neighborhood of $\seqname{C}^2$.

\begin{remark}
If $\funcseqname{f}$ is constrained then $\Top(\seqname{E}^1)$ is a
model neighborhood of $\seqname{E}^1$ and $\Top(\seqname{C}^1)$ is a
model neighborhood of $\seqname{C}^1$.
\end{remark}
\end{definition}

\begin{definition}[M-atlas morphisms for topological spaces]
\label{def:M-ATLmorphtop}
Let $\catname{C}^i$, $i=1,2$, be a model category,
$\seqname{C}^i \objin \catname{C}^i$, $E^i$ be a topological space and
$\seqname{A}^i$ be an m-atlas of $E^i$ in the coordinate space
$\seqname{C}^i$.

A pair of functions
$
  \funcseqname{f} \defeq
  (
    \funcname{f}_0 \maps E^1 \to E^2,
    \funcname{f}_1 \maps \seqname{C}^1 \to \seqname{C}^2
  )
$
is an $E^1$-$E^2$ m-atlas morphism of
$\seqname{A}^1$ to $\seqname{A}^2$ in the coordinate spaces
$\seqname{C}^i$, abbreviated as \\*
$
  \isAtl_\Ar
    (
      \seqname{A}^1,
      E^1,
      \seqname{C}^1,
      \seqname{A}^2,
      E^2,
      \seqname{C}^2,
      \funcname{f}_0,
      \funcname{f}_1
    )
$,
and an $E^1$-$E^2$ m-atlas morphism of $\seqname{A}^1$ to $\seqname{A}^2$
in the coordinate model categories $\catname{C}^i$, abbreviated as \\*
$
  \isAtl_\Ar
    (
      \seqname{A}^1,
      E^1,
      \catname{C}^1,
      \seqname{A}^2,
      E^2,
      \catname{C}^2,
      \funcname{f}_0,
      \funcname{f}_1
    )
$,
iff it is an $\Triv{E^1}$-$\Triv{E^2}$ m-atlas morphism of
$\seqname{A}^1$ to $\seqname{A}^2$ in the coordinate model categories
$\catname{C}^1$, $\catname{C}^2$.

The triple
$
  \bigl (
    \funcseqname{f},
    (\seqname{A}^1, E^1, \seqname{C}^1),
    (\seqname{A}^2, E^2, \seqname{C}^2)
  \bigr )
$
will refer to $\funcseqname{f}$ considered as an $E^1$-$E^2$ m-atlas
(near) morphism of $\seqname{A}^1$ to $\seqname{A}^2$ in the coordinate
spaces $\seqname{C}^1$, $\seqname{C}^2$.

$\funcseqname{f}$ is also a constrained $E^1$-$E^2$ m-atlas morphism of
$\seqname{A}^1$ to $\seqname{A}^2$ in the coordinate spaces
$\seqname{C}^1$, $\seqname{C}^2$, abbreviated as
$
  \isAtl[const]_\Ar
    (
      \seqname{A}^1,
      E^1,
      \seqname{C}^1,
      \seqname{A}^2,
      E^2,
      \seqname{C}^2,
      \funcname{f}_0,
      \funcname{f}_1
    )
$,
and a constrained $E^1$-$E^2$ m-atlas morphism of $\seqname{A}^1$ to $\seqname{A}^2$
in the coordinate model categories $\catname{C}^1$, $\catname{C}^2$, abbreviated as
$
  \isAtl[const]_\Ar
    (
      \seqname{A}^1,
      E^1,
      \catname{C}^1,
      \seqname{A}^2,
      E^2,
      \catname{C}^2,
      \funcname{f}_0,
      \funcname{f}_1
    )
$,
iff $\funcname{f}_1$ is constrained, i.e.,
$\funcname{f}_1[\seqname{C}^1]$ is contained in a model neighborhood of
$\seqname{C}^2$.
\end{definition}

\begin{definition}[M-atlas near morphisms for model spaces]
\label{def:M-ATLmorphnear}
Let $\catname{E}^i$, $\catname{C}^i$, $i=1,2$, be model categories,
$\seqname{E}^i \objin \catname{E}^i$,
$\seqname{C}^i \objin \catname{C}^i$, $\seqname{A}^i$ be an m-atlas of
$\seqname{E}^i$ in the coordinate space $\seqname{C}^i$ and
$
  \funcseqname{f} \defeq
  (
    \funcname{f}_0 \maps \seqname{E}^1 \to \seqname{E}^2,
    \funcname{f}_1 \maps \seqname{C}^1 \to \seqname{C}^2
  )
$
a pair of model functions.

$\funcseqname{f}$ is an $\seqname{E}^1$-$\seqname{E}^2$ m-atlas near
morphism of $\seqname{A}^1$ to $\seqname{A}^2$ in the coordinate spaces
$\seqname{C}^1$, $\seqname{C}^2$, abbreviated as
$
  \isAtl[near]_\Ar
    (
      \seqname{A}^1,
      \seqname{E}^1,
      \seqname{C}^1,
      \seqname{A}^2,
      \seqname{E}^2,
      \seqname{C}^2,
      \funcname{f}_0,
      \funcname{f}_1
    )
$
and an $\catname{E}^1$-$\catname{E}^2$ m-atlas near morphism of
$\seqname{A}^1$ to $\seqname{A}^2$ in the coordinate model categories
$\catname{C}^1$, $\catname{C}^2$, abbreviated as
$
  \isAtl[near]_\Ar
  (
    \seqname{A}^1,
    \catname{E}^1,
    \catname{C}^1,
    \seqname{A}^2,
    \catname{E}^2,
    \catname{C}^2,
    \funcname{f}_0,
    \funcname{f}_1
  )
$
iff for any charts
$(U^i, V^i, \phi^i \maps U^i \toiso V^i) \in \seqname{A}^i$, $i=1,2$,
with $I \defeq U^1 \cap \funcname{f}^{-1}_0[U^2] \neq \emptyset$,
diagram \eqref{eq:AtlN1} below is strongly M-locally nearly commutative
in $\seqname{C}^2$, i.e., for any
$u^1 \in I$ there are model neighborhoods $U'^1 \subseteq I$,
$V'^1 \subseteq V^1$, $U'^2 \subseteq U^2$, $V'^2 \subseteq V^2$,
$\hat{V}'^2 \subseteq C^2$ and an isomorphism
$\hat{\funcname{f}} \maps V'^2 \toiso \hat{V}'^2$%
\footnote
{
  This reverses the direction of the arrow
  $\hat{\funcname{f}} \maps U_m \toiso V_n$ from
  \pagecref{def:NCD}.
  This is permissible since $\hat{\funcname{f}}$ is an isomorphism
}
\newcounter{fn:reverses}
\setcounter{fn:reverses} {\value{footnote}}
such that \crefrange{eq:xin}{eq:fphigeqgphi} below hold, as shown in
{
  \showlabelsinline
  \crefrange{fig:AtlN1}{fig:AtlN2}.
}

\begin{equation}
\label{eq:AtlN1}
\begin{array} {lll}
D \defeq & \Biggl \{
\\
  &&
  \funcname{f}_0 \maps I \defeq U^1 \cap \funcname{f}^{-1}_0[U^2] \to U^2,
  \phi^2 \maps U^2 \toiso V^2,
\\
  &&
  \phi^1 \maps I \into V^1,
  \funcname{f}_1 \maps  V^1 \to \seqname{C}^2
\\
& \Biggr \}
\end{array}
\end{equation}

\begin{equation}
u^1 \in U'^1
\label{eq:xin}
\end{equation}

\begin{equation}
\label{eq:fup1in}
\funcname{f}_0[U'^1] \subseteq U'^2
\end{equation}

\begin{equation}
\phi^1[U'^1] = V'^1
\end{equation}

\begin{equation}
\funcname{f}_1[V'^1] \subseteq \hat{V}'^2
\label{eq:fvp1in}
\end{equation}

\begin{equation}
\phi^2[U'^2] = V'^2
\end{equation}

\begin{equation}
\hat{\funcname{f}} \compose \phi^2 \compose \funcname{f}_0 =
\funcname{f}_1 \compose \phi^1
\label{eq:fphigeqgphi}
\end{equation}

\begin{remark}
$(U'^i, V'^i, \phi^i \restrictto_{U'^i, V'^i})$ need not be a chart of
$\seqname{A}^i$.
\end{remark}

\begin{figure}[ht]
\[ \bfig
\node uint(0,0)[{I = U^1 \cap \funcname{f}^{-1}_0[U^2]}]
\node u2(2000,0)[U^2]
\node v1(0,-1000)[V^1]
\node c21(1000,-1000)[C^2]
\node v2(2000,-1000)[V^2]
\arrow |a|[uint`u2;\funcname{f}_0]
\arrow |r|/>->/[uint`v1;\phi^1]
\arrow |r|[u2`v2;\phi^2]
\arrow |l|//[u2`v2;\iso]
\arrow |a|[v1`c21;\funcname{f}_1]
\efig \]
\caption{Uncompleted m-atlas near morphism}
\label{fig:AtlN1}
\end{figure}

\begin{figure}[ht]
\[ \bfig
\node x(500,0)[u^1]
\node uint(-300,-500)[{I=U^1 \cap \funcname{f}^{-1}_0[U^2]}]
\node up1(500,-500)[U'^1]
\node up2(1500,-500)[U'^2]
\node u2(2000,-500)[U^2]
\node v1(-300,-1000)[V^1]
\node vp1(500,-1000)[V'^1]
\node vhp2(1000,-1000)[\hat{V}'^2]
\node vp2(1500,-1000)[V'^2]
\node v2(2000,-1000)[V^2]
\arrow |r|/^{ (}->/[x`up1;\funcname{i}]
\arrow |r|/>->/[uint`v1;\phi^1]
\arrow |a|/_{ (}->/[up1`uint;\funcname{i}]
\arrow |a|[up1`up2;\funcname{f}_0]
\arrow |r|[up1`vp1;{\phi^1}]
\arrow |l|//[up1`vp1;\iso]
\arrow |a|/^{ (}->/[up2`u2;i]
\arrow |r|[up2`vp2;\phi^2]
\arrow |l|//[up2`vp2;\iso]
\arrow |r|[u2`v2;\phi^2]
\arrow |l|//[u2`v2;\iso]
\arrow |a|/_{ (}->/[vp1`v1;\funcname{i}]
\arrow |a|[vp1`vhp2;\funcname{f}_1]
\arrow |a|/ >.>>/[vp2`vhp2;\hat{\funcname{f}}]
\arrow |b|//[vp2`vhp2;\iso]
\arrow |a|/^{ (}->/[vp2`v2;\funcname{i}]
%\arrow |a|/^{ (}->/[v2`c2;\funcname{i}]
\efig \]
\caption{Completed m-atlas near morphism}
\label{fig:AtlN2}
\end{figure}

The triple
$
  \Bigl (
    \funcseqname{f},
    (\seqname{A}^1, \seqname{E}^1, \seqname{C}^1),
    (\seqname{A}^2, \seqname{E}^2, \seqname{C}^2)
  \Bigr )
$
will refer to $\funcseqname{f}$ considered as an
$\seqname{E}^1$-$\seqname{E}^2$ m-atlas near morphism of
$\seqname{A}^1$ to $\seqname{A}^2$ in the coordinate spaces
$\seqname{C}^1$, $\seqname{C}^2$.
\end{definition}

\begin{definition}[M-atlas near morphisms for topological spaces]
\label{def:M-ATLmorphtopnear}
Let $\catname{C}^i$, $i=1,2$, be a model category,
$\seqname{C}^i \objin \catname{C}^i$, $E^i$ be a topological space,
$\seqname{A}^i$ be an m-atlas of $E^i$ in the coordinate space
$\seqname{C}^i$ and
$
  \funcseqname{f} \defeq
  (
    \funcname{f}_0 \maps E^1 \to E^2,
    \funcname{f}_1 \maps \seqname{C}^1 \to \seqname{C}^2
  )
$
be a pair of continuous functions.

$\funcseqname{f}$ is an $E^1$-$E^2$ m-atlas near morphism of
$\seqname{A}^1$ to $\seqname{A}^2$ in the coordinate spaces
$\seqname{C}^1$, $\seqname{C}^2$, abbreviated as
$
  \isAtl[near]_\Ar
    (
      \seqname{A}^1,
      E^1,
      \seqname{C}^1,
      \seqname{A}^2,
      E^2,
      \seqname{C}^2,
      \funcname{f}_0,
      \funcname{f}_1
    )
$,
and a
$E^1$-$E^2$ m-atlas near morphism of $\seqname{A}^1$ to
$\seqname{A}^2$ in the coordinate model categories $\catname{C}^1$, $\catname{C}^2$,
abbreviated as
$
  \isAtl[near]_\Ar
    (
      \seqname{A}^1,
      E^1,
      \catname{C}^1,
      \seqname{A}^2,
      E^2,
      \catname{C}^2,
      \funcname{f}_0,
      \funcname{f}_1
    )
$,
iff it is an $\Triv{E^1}$-$\Triv{E^2}$ m-atlas near morphism of
$\seqname{A}^1$ to $\seqname{A}^2$ in the coordinate model spaces
$\seqname{C}^1$, $\seqname{C}^2$.

The triple
$
  \Bigl (
    \funcseqname{f},
    (\seqname{A}^1, E^1, \seqname{C}^1),
    (\seqname{A}^2, E^2, \seqname{C}^2)
  \Bigr )
$
will refer to $\funcseqname{f}$ considered as an
$E^1$-$E^2$ m-atlas near morphism of
$\seqname{A}^1$ to $\seqname{A}^2$ in the coordinate spaces
$\seqname{C}^1$, $\seqname{C}^2$.
\end{definition}

\begin{definition}[M-atlas inclusion and identity morphisms]
\label{def:M-ATLident}
Let
\begin{enumerate}
\item
$\catname{E}^i$, $\catname{C}^i$, $i=1,2$, be model categories
\item
$\seqname{E}^i \objin \catname{E}^i$
\item
$\seqname{E}^1 \submod \seqname{E}^2$
\item
$\seqname{C}^i \objin \catname{C}^i$
\item
$\seqname{C}^1 \submod \seqname{C}^2$
\item
$\seqname{A}^i$ be an m-atlas of $\seqname{E}^i$ in the coordinate
space $\seqname{C}^i$
\item
$\seqname{A}^1 \subseteq \seqname{A}^2$
\end{enumerate}

The $\seqname{E}^1$-$\seqname{E}^2$ m-atlas inclusion morphism of
$\seqname{A}^1$ to $\seqname{A}^2$ in the coordinate spaces
$\seqname{C}^1$, $\seqname{C}^2$ is
$
 \ID%
   [
     {(\seqname{E}^1, \seqname{C}^1)},
     {(\seqname{E}^2, \seqname{C}^2)}
   ]
$
considered as an $\seqname{E}^1$-$\seqname{E}^2$ m-atlas morphism of
$\seqname{A}^1$ to $\seqname{A}^2$ in the coordinate spaces
$\seqname{C}^1$, $\seqname{C}^2$.

The m-atlas inclusion morphism of
$(\seqname{A}^1, \seqname{E}^1, \seqname{C}^1)$ to
$(\seqname{A}^2, \seqname{E}^2, \seqname{C}^2)$ is \\*
$
  \Id%
    [
        {(\seqname{A}^1, \seqname{E}^1, \seqname{C}^1)},
        {(\seqname{A}^2, \seqname{E}^2, \seqname{C}^2)}
    ]
  \defeq
  \biggl (
    \ID%
      [
        {(\seqname{E}^1,\seqname{C}^1)},
        {(\seqname{E}^2,\seqname{C}^2)}
      ],
    (\seqname{A}^1, \seqname{E}^1, \seqname{C}^1),
    (\seqname{A}^2, \seqname{E}^2, \seqname{C}^2)
  \biggr )
$,\hspace{1pt}
$
 \ID%
   [
     {(\seqname{E}^1, \seqname{C}^1)},
     {(\seqname{E}^2, \seqname{C}^2)}
   ]
$
considered as an $\seqname{E}^1$-$\seqname{E}^2$ m-atlas morphism of
$\seqname{A}^1$ to $\seqname{A}^2$ in the coordinate spaces
$\seqname{C}^1$, $\seqname{C}^2$.

The $\seqname{E}^i$ m-atlas identity morphism of $\seqname{A}^i$ in the
coordinate space $\seqname{C}^1$, $\seqname{C}^2$ is
$
 \ID%
   [
     {(\seqname{E}^i, \seqname{C}^i)}
   ]
$
considered as an $\seqname{E}^1$-$\seqname{E}^2$ m-atlas morphism of
$\seqname{A}^1$ to $\seqname{A}^2$ in the coordinate spaces
$\seqname{C}^1$, $\seqname{C}^2$.

The m-atlas identity morphism of
$(\seqname{A}^i, \seqname{E}^i, \seqname{C}^i)$ is \\*
$
  \Id[{{(\seqname{A}^i, \seqname{E}^i, \seqname{C}^i)}}] \defeq
  \bigl (
    \ID[{{(\seqname{E}^i,\seqname{C}^i)}}],
    (\seqname{A}^i, \seqname{E}^i, \seqname{C}^i),
    (\seqname{A}^i, \seqname{E}^i, \seqname{C}^i)
  \bigr )
$,\hspace{1pt}
$
 \ID%
   [
     {(\seqname{E}^i, \seqname{C}^i)}
   ]
$\hspace{1pt}
considered as an $\seqname{E}^i$ m-atlas morphism of $\seqname{A}^i$ in
the coordinate space $\seqname{C}^i$.

Let
\begin{enumerate}
\item
$E^i$ be a topological space, $i=1,2$
\item
$E^1 \subsp E^2$
\item
$\catname{C}^i$ be a model category
\item
$\seqname{C}^i \objin \catname{C}^i$
\item
$\seqname{C}^1 \submod \seqname{C}^2$
\item
$\seqname{A}^i$ be an m-atlas of $E^i$ in the coordinate space
$\seqname{C}^i$
\item
$\seqname{A}^1 \subseteq \seqname{A}^2$
\end{enumerate}

The $E^1$-$E^2$ m-atlas inclusion morphism of
$\seqname{A}^1$ to $\seqname{A}^2$ in the coordinate spaces
$\seqname{C}^1$, $\seqname{C}^2$ is
$
 \ID%
   [
     {(E^1, \seqname{C}^1)},
     {(E^2, \seqname{C}^2)}
   ]
$.

The m-atlas inclusion morphism of
$(\seqname{A}^1, E^1, \seqname{C}^1)$ to
$(\seqname{A}^2, E^2, \seqname{C}^2)$ is \\*
$
  \Id%
    [
        {(\seqname{A}^1, E^1, \seqname{C}^1)},
        {(\seqname{A}^2, E^2, \seqname{C}^2)}
    ]
  \defeq
  \biggl (
    \ID%
    [
      {(E^1,\seqname{C}^1)},
      {(E^2,\seqname{C}^2)}
    ],
    (\seqname{A}^1, E^1, \seqname{C}^1),
    (\seqname{A}^2, E^2, \seqname{C}^2)
  \biggr )
$,\hspace{1pt}
$
 \ID%
   [
     {(E^1, \seqname{C}^1)},
     {(E^2, \seqname{C}^2)}
   ]
$
considered as an $E^1$-$E^2$ m-atlas morphism of
$\seqname{A}^1$ to $\seqname{A}^2$ in the coordinate spaces
$\seqname{C}^1$, $\seqname{C}^2$.

The $E^i$ m-atlas identity morphism of $\seqname{A}^i$ in the
coordinate space $\seqname{C}^i$ is
$
 \ID%
   [
     {(E^i, \seqname{C}^i)}
   ]
$.

The m-atlas identity morphism of
$(\seqname{A}^i, E^i, \seqname{C}^i)$ is \\*
$
\Id[{{(\seqname{A}^i, E^i, \seqname{C}^i)}}] \defeq
\bigl (
  \ID[{{(E^i,\seqname{C}^i)}}],
  (\seqname{A}^i, E^i, \seqname{C}^i),
  (\seqname{A}^i, E^i, \seqname{C}^i)
\bigr )
$,\hspace{1pt}
$
 \ID%
   [
     {(E^i, \seqname{C}^i)}
   ]
$
considered as an $E^i$ m-atlas morphism of $\seqname{A}^i$ in
the coordinate space $\seqname{C}^i$.
\end{definition}

\begin{definition}[Equivalence of m-atlas (near) morphisms]
\label{def:M-ATLmorphequiv}
Let $\seqname{E}^i$, $\seqname{C}^i$, $i=1,2$, be model spaces,
$\seqname{A}^i$ be an m-atlas of $\seqname{E}^i$ in the coordinate space
$\seqname{C}^i$, and \\*
$
  \funcseqname{f} \defeq
  (
    \funcname{f}_0 \maps \seqname{E}^1 \to \seqname{E}^2,
    \funcname{f}_1 \maps \seqname{C}^1 \to \seqname{C}^2
  )
$ and
$
  \funcseqname{g} \defeq
  (
    \funcname{g}_0 \maps \seqname{E}^1 \to \seqname{E}^2,
    \funcname{g}_1 \maps \seqname{C}^1 \to \seqname{C}^2
  )
$
be m-atlas (near) morphisms of $\seqname{A}^1$ to $\seqname{A}^2$ in the
coordinate spaces $\seqname{C}^1$, $\seqname{C}^2$.  Then
$\funcseqname{f}$ is equivalent to $\funcseqname{g}$ as an
$\seqname{E}^1$-$\seqname{E}^2$ m-atlas (near) morphism of
$\seqname{A}^1$ to $\seqname{A}^2$ in the coordinate spaces
$\seqname{C}^1$, $\seqname{C}^2$, abbreviated
$
  \funcseqname{f} \maps
  (\seqname{A}^1, \seqname{E}^1, \seqname{C}^1) \to
  (\seqname{A}^2, \seqname{E}^2, \seqname{C}^2)
$
is equivalent to
$
  \funcseqname{g} \maps
  (\seqname{A}^1, \seqname{E}^1, \seqname{C}^1) \to
  (\seqname{A}^2, \seqname{E}^2, \seqname{C}^2)
$
and
$
  \bigl (
    \funcseqname{f},
    (\seqname{A}^1, \seqname{E}^1, \seqname{C}^1),
    (\seqname{A}^2, \seqname{E}^2, \seqname{C}^2)
  \bigr )
$
is equivalent to
$
  \bigl (
    \funcseqname{g},
    (\seqname{A}^1, \seqname{E}^1, \seqname{C}^1),
    (\seqname{A}^2, \seqname{E}^2, \seqname{C}^2)
  \bigr )
$
iff $\funcname{f}_0 = \funcname{g}_0$.

By abuse of language we write $\funcseqname{f}$ is equivalent to
$\funcseqname{g}$ when the spaces and atlases are understood by context.

Let $E^i$, $i=1,2$, be a topological space, $\seqname{C}^i$ be a model
space, $\seqname{A}^i$ be an m-atlas of $E^i$ in the coordinate space
$\seqname{C}^i$ and
$
  \funcseqname{f} \defeq
  (
    \funcname{f}_0 \maps E^1 \to E^2,
    \funcname{f}_1 \maps \seqname{C}^1 \to \seqname{C}^2
  )
$, and \\*
$
  \funcseqname{g} \defeq
  (
    \funcname{g}_0 \maps E^1 \to E^2,
    \funcname{g}_1 \maps \seqname{C}^1 \to \seqname{C}^2
  )
$
be m-atlas (near) morphisms of $\seqname{A}^1$ to $\seqname{A}^2$ in the
coordinate spaces $\seqname{C}^1$, $\seqname{C}^2$.
Then $\funcseqname{f}$ is equivalent to $\funcseqname{g}$ as an
$E^1$-$E^2$ m-atlas (near) morphism  of $\seqname{A}^1$ to
$\seqname{A}^2$ in the coordinate spaces $\seqname{C}^1$,
$\seqname{C}^2$, abbreviated
$
  \funcseqname{f} \maps
  (\seqname{A}^1, E^1, \seqname{C}^1) \to
  (\seqname{A}^2, E^2, \seqname{C}^2)
$
is equivalent to
$
  \funcseqname{g} \maps
  (\seqname{A}^1, E^1, \seqname{C}^1) \to
  (\seqname{A}^2, E^2, \seqname{C}^2)
$,
iff $\funcname{f}_0 = \funcname{g}_0$.

By abuse of language we write $\funcseqname{f}$ is equivalent to
$\funcseqname{g}$ when the spaces and atlases are understood by context.
\end{definition}

\subsubsection{Definitions of semi-strict and strict}
\begin{definition}[Semi-strict m-atlas (near) morphisms]
\label{def:M-ATLmorphsemi}
Let $\catname{E}^i$, $\catname{C}^i$, $i=1,2$, be model categories,
$\seqname{E}^i \objin \catname{E}^i$,
$\seqname{C}^i \objin \catname{C}^i$, $\seqname{A}^i$ be an m-atlas of
$\seqname{E}^i$ in the coordinate space $\seqname{C}^i$ and
$
  \funcseqname{f} \defeq
  (
    \funcname{f}_0 \maps \seqname{E}^1 \to \seqname{E}^2,
    \funcname{f}_1 \maps \seqname{C}^1 \to \seqname{C}^2
  )
$
an $\seqname{E}^1$-$\seqname{E}^2$ m-atlas (near) morphism of
$\seqname{A}^1$ to $\seqname{A}^2$ in the coordinate spaces
$\seqname{C}^1$, $\seqname{C}^2$.

$\funcseqname{f}$ is a semi-strict $\seqname{E}^1$-$\seqname{E}^2$
m-atlas (near) morphism of $\seqname{A}^1$ to $\seqname{A}^2$ in the
coordinate spaces $\seqname{C}^1$, $\seqname{C}^2$, abbreviated as
$
  \strict[semi-]{\isAtl[(near)]_\Ar}
  (
    \seqname{A}^1,
    \seqname{E}^1,
    \seqname{C}^1,
    \seqname{A}^2,
    \seqname{E}^2,
    \seqname{C}^2,
    \funcname{f}_0,
    \funcname{f}_1
  )
$,
iff $\funcname{f}_0$ is a local m-morphism of $\seqname{E}^1$
to $\seqname{E}^2$, $\funcname{f}_1$ is a local m-morphism of
$\seqname{C}^1$ to $\seqname{C}^2$ and for every \\*
$(U^i, V^i, \phi^i \maps U^i \toiso V^i) \in \seqname{A}^i$, $i=1,2$,
with $I \defeq U^1 \cap \funcname{f}^{-1}_0[U^2] \neq \emptyset$,
$
  \phi^2 \compose \funcname{f}_0 \compose \phi^{1-1} \maps
  \phi^1[I] \to V^2
$
is a local $\seqname{C}^1$-$\seqname{C}^2$ m-morphism of $\phi^1[I]$ to
$V^2$.

It is also a coherent semi-strict
\nobreaks{{$\seqname{E}^1$-$\seqname{E}^2$}}
m-atlas (near) morphism of $\seqname{A}^1$ to $\seqname{A}^2$ in the
coordinate spaces $\seqname{C}^1$, $\seqname{C}^2$ iff
each $E^i \objin \seqname{E}^i$,
each $C^i \objin \seqname{C}^i$,
$\seqname{E}^1 \submod[strict-] \seqname{E}^2$ and
$\seqname{C}^1 \submod[strict-] \seqname{C}^2$.

$\funcseqname{f}$ is a semi-strict
$\seqname{E}^1$-$\seqname{E}^2$-$\catname{C}^1$-$\catname{C}^2$ m-atlas
(near) morphism of $\seqname{A}^1$ to $\seqname{A}^2$ in the coordinate
spaces $\seqname{C}^1$, $\seqname{C}^2$, abbreviated as
$
  \strict[semi-]{\isAtl[(near)]_\Ar}
  (
    \seqname{A}^1,
    \seqname{E}^1,
    \seqname{C}^1,
    \seqname{A}^2,
    \seqname{E}^2,
    \seqname{C}^2,
    \funcname{f}_0,
    \funcname{f}_1,
    \catname{C}^1,
    \catname{C}^2
  )
$,
iff $\funcname{f}_0$ is a local
m-morphism of $\seqname{E}^1$ to $\seqname{E}^2$, $\funcname{f}_1$ is a
local $\catname{C}^1$-$\catname{C}^2$ m-morphism of $\seqname{C}^1$ to
$\seqname{C}^2$ and for every
$(U^i,V^i, \phi^i \maps U^i \toiso V^i) \in \seqname{A}^i$, $i=1,2$,
with $I \defeq U^1 \cap \funcname{f}^{-1}_0[U^2] \neq \emptyset$,
$
  \phi^2 \compose \funcname{f}_0 \compose \phi^{1-1} \maps
  \phi^1[I] \to V^2
$
is a local $\catname{C}^1$-$\catname{C}^2$ m-morphism of $\phi^1[I]$ to
$V^2$.

It is also a coherent semi-strict
\nobreaks
  {
    {$\seqname{E}^1$-$\seqname{E}^2$-},
    {$\catname{C}^1$-$\catname{C}^2$}
  }
m-atlas (near) morphism of $\seqname{A}^1$ to $\seqname{A}^2$ in the
coordinate spaces $\seqname{C}^1$, $\seqname{C}^2$ iff
each $E^i \objin \seqname{E}^i$,
$\seqname{E}^1 \submod[strict-] \seqname{E}^2$ and
$\catname{C}^1 \subcat[strict-] \catname{C}^2$.

$\funcseqname{f}$ is a semi-strict
$\catname{E}^1$-$\catname{E}^2$-$\seqname{E}^1$-$\seqname{E}^2$ m-atlas
(near) morphism of $\seqname{A}^1$ to $\seqname{A}^2$ in the coordinate
spaces $\seqname{C}^1$, $\seqname{C}^2$, abbreviated as
$
  \strict[semi-]{\isAtl[(near)]_\Ar}
  (
    \seqname{A}^1,
    \seqname{E}^1,
    \seqname{C}^1,
    \seqname{A}^2,
    \seqname{E}^2,
    \seqname{C}^2,
    \funcname{f}_0,
    \funcname{f}_1,
    \catname{E}^1,
    \catname{E}^2
  )
$,
iff $\funcname{f}_0$ is a local $\catname{E}^1$-$\catname{E}^2$
m-morphism of $\seqname{E}^1$ to $\seqname{E}^2$, $\funcname{f}_1$ is
a local m-morphism of $\seqname{C}^1$ to $\seqname{C}^2$ and for
every $(U^i, V^i, \phi^i \maps U^i \toiso V^i) \in \seqname{A}^i$,
$i=1,2$, with
$I \defeq U^1 \cap \funcname{f}^{-1}_0[U^2] \neq \emptyset$,
$
  \phi^2 \compose \funcname{f}_0 \compose \phi^{1-1} \maps
  \phi^1[I] \to U^2
$
is a local $\seqname{C}^1$-$\seqname{C}^2$ m-morphism of $\phi^1[I]$ to $V^2$.

It is also a coherent semi-strict
\nobreaks
  {
    {$\catname{E}^1$-$\catname{E}^2$-},
    {$\seqname{E}^1$-$\seqname{E}^2$}
  }
m-atlas (near) morphism of $\seqname{A}^1$ to $\seqname{A}^2$ in the
coordinate spaces $\seqname{C}^1$, $\seqname{C}^2$ iff
each $C^i \objin \seqname{C}^i$,
$\catname{E}^1 \subcat[strict-] \catname{E}^2$ and
$\seqname{C}^1 \submod[strict-] \seqname{C}^2$.

$\funcseqname{f}$ is a semi-strict
\nobreaks
  {
    {$\catname{E}^1$-$\catname{E}^2$-},
    {$\seqname{E}^1$-$\seqname{E}^2$-},
    {$\catname{C}^1$-$\catname{C}^2$}
  }
m-atlas (near) morphism of $\seqname{A}^1$ to $\seqname{A}^2$ in the
coordinate spaces $\seqname{C}^1$, $\seqname{C}^2$, abbreviated as \\*
$
  \strict[semi-]{\isAtl[(near)]_\Ar}
  (
    \seqname{A}^1,
    \seqname{E}^1,
    \seqname{C}^1,
    \seqname{A}^2,
    \seqname{E}^2,
    \seqname{C}^2,
    \funcname{f}_0,
    \funcname{f}_1,
    \catname{E}^1,
    \catname{E}^2,
    \catname{C}^1,
    \catname{C}^2
  )
$,
iff $\funcname{f}_0$ is a local $\catname{E}^1$-$\catname{E}^2$
m-morphism of $\seqname{E}^1$ to $\seqname{E}^2$ and $\funcname{f}_1$ is
a local $\catname{C}^1$-$\catname{C}^2$ m-morphism of $\seqname{C}^1$ to
$\seqname{C}^2$ and for every \\*
$(U^i,V^i, \phi^i \maps U^i \toiso V^i) \in \seqname{A}^i$, $i=1,2$,
with $I \defeq U^1 \cap \funcname{f}^{-1}_0[U^2] \neq \emptyset$,
$
  \phi^2 \compose \funcname{f}_0 \compose \phi^{1-1} \maps
  \phi^1[I] \to V^2
$
is a local \nobreaks{{$\catname{C}^1$-$\catname{C}^2$}} m-morphism of
$\phi^1[I]$ to $V^2$.

It is also a coherent semi-strict
\nobreaks
  {
    {$\catname{E}^1$-$\catname{E}^2$-},
    {$\seqname{E}^1$-$\seqname{E}^2$-},
    {$\catname{C}^1$-$\catname{C}^2$}
  }
m-atlas (near) morphism of $\seqname{A}^1$ to $\seqname{A}^2$ in the
coordinate spaces $\seqname{C}^1$, $\seqname{C}^2$ iff
$\catname{E}^1 \subcat[strict-] \catname{E}^2$ and
$\catname{C}^1 \subcat[strict-] \catname{C}^2$.

By abuse of language we will write $\funcseqname{f}$ is coherent when it
is a
\nobreaks
  {
    {($\catname{E}^1$-$\catname{E}^2$-)},
    {$\seqname{E}^1$-$\seqname{E}^2$},
    {(-$\catname{C}^1$-$\catname{C}^2$)}
  }
m-atlas (near) morphism of $\seqname{A}^1$ to $\seqname{A}^2$ in the
coordinate spaces $\seqname{C}^1$, $\seqname{C}^2$ where the model
spaces, model categories and atlases are understood by context.
\end{definition}

\begin{definition}[Strict m-atlas (near) morphisms]
\label{def:M-ATLmorphstrict}
Let
\begin{enumerate}
\item
$\catname{E}^i$, $\catname{C}^i$, $i=1,2$, be model categories
\item
$\seqname{E}^i \objin \catname{E}^i$
\item
$\seqname{C}^i \objin \catname{C}^i$
\item
$\seqname{A}^i$ be an m-atlas of
$\seqname{E}^i$ in the coordinate space $\seqname{C}^i$
\item
$
  \funcseqname{f} \defeq
  (
    \funcname{f}_0 \maps \seqname{E}^1 \to \seqname{E}^2,
    \funcname{f}_1 \maps \seqname{C}^1 \to \seqname{C}^2
  )
$
be an m-atlas (near) morphism
\end{enumerate}

$\funcseqname{f}$ is a strict \nobreaks{$\seqname{E}^1-\seqname{E}^2$}
m-atlas (near) morphism of $\seqname{A}^1$ to $\seqname{A}^2$ in the
coordinate spaces $\seqname{C}^1$, $\seqname{C}^2$, abbreviated as
$
   \strict{\isAtl[(near)]_\Ar}
   (
     \seqname{A}^1,
     \seqname{E}^1,
     \seqname{C}^1,
     \seqname{A}^2,
     \seqname{E}^2,
     \seqname{C}^2,
     \funcname{f}_0,
     \funcname{f}_1
   )
$,
iff $\funcname{f}_0$ is an m-morphism of $\seqname{E}^1$ to
$\seqname{E}^2$, $\funcname{f}_1$ is an m-morphism of
$\seqname{C}^1$ to $\seqname{C}^2$ and for every \\*
$(U^i, V^i, \phi^i \maps U^i \toiso V^i) \in \seqname{A}^i$, $i=1,2$,
with $I \defeq U^1 \cap \funcname{f}^{-1}_0[U^2] \neq \emptyset$,
$
   \phi^2 \compose \funcname{f}_0 \compose \phi^{1-1} \maps
   \phi^1[I] \to U^2
$
is a morphism of $\Mod(V^2,\seqname{C}^2)$.

It is also a coherent strict
\nobreaks{{$\seqname{E}^1$-$\seqname{E}^2$}}
m-atlas (near) morphism of $\seqname{A}^1$ to $\seqname{A}^2$ in the
coordinate spaces $\seqname{C}^1$, $\seqname{C}^2$ iff
each $E^i \objin \seqname{E}^i$,
each $C^i \objin \seqname{C}^i$,
$\seqname{E}^1 \submod[strict-] \seqname{E}^2$ and
$\seqname{C}^1 \submod[strict-] \seqname{C}^2$.

$\funcseqname{f}$ is a strict
\nobreaks
  {
    {$\seqname{E}^1-\seqname{E}^2$-},
    {$\catname{C}^1-\catname{C}^2$}
  }
m-atlas (near) morphism of $\seqname{A}^1$ to $\seqname{A}^2$ in the
coordinate spaces $\seqname{C}^1$, $\seqname{C}^2$, abbreviated as
$
  \strict{\isAtl[(near)]_\Ar}
  (
    \seqname{A}^1,
    \seqname{E}^1,
    \seqname{C}^1,
    \seqname{A}^2,
    \seqname{E}^2,
    \seqname{C}^2,
    \funcname{f}_0,
    \funcname{f}_1,
    \catname{C}^1,
    \catname{C}^2
  )
$,
iff $\funcname{f}_0$ is an m-morphism of $\seqname{E}^1$ to
$\seqname{E}^2$, $\funcname{f}_1$ is an m-morphism of
$\catname{C}^1$ to $\catname{C}^2$ and for every \\*
$(U^i,V^i, \phi^i \maps U^i \toiso V^i) \in \seqname{A}^i$, $i=1,2$,
with $I \defeq U^1 \cap \funcname{f}^{-1}_0[U^2] \neq \emptyset$,
$
  \phi^2 \compose \funcname{f}_0 \compose \phi^{1-1} \maps
  \phi^1[I] \to V^2
$
is a morphism of $\catname{C}^2$.

It is also a coherent strict
\nobreaks
  {
    {$\seqname{E}^1$-$\seqname{E}^2$-},
    {$\catname{C}^1$-$\catname{C}^2$}
  }
m-atlas (near) morphism of $\seqname{A}^1$ to $\seqname{A}^2$ in the
coordinate spaces $\seqname{C}^1$, $\seqname{C}^2$ iff
each $E^i \objin \seqname{E}^i$,
$\seqname{E}^1 \submod[strict-] \seqname{E}^2$ and
$\catname{C}^1 \subcat[strict-] \catname{C}^2$.

$\funcseqname{f}$ is a strict
\nobreaks
  {
    {$\catname{E}^1$-$\catname{E}^2$-},
    {$\seqname{E}^1$-$\seqname{E}^2$}
  }
m-atlas (near) morphism of $\seqname{A}^1$ to $\seqname{A}^2$ in the
coordinate spaces $\seqname{C}^1$, $\seqname{C}^2$, abbreviated as
$
  \strict{\isAtl[(near)]_\Ar}
  (
    \seqname{A}^1,
    \seqname{E}^1,
    \seqname{C}^1,
    \seqname{A}^2,
    \seqname{E}^2,
    \seqname{C}^2,
    \funcname{f}_0,
    \funcname{f}_1,
    \catname{E}^1,
    \catname{E}^2
  )
$,
iff $\funcname{f}_0$ is an m-morphism of $\catname{E}^1$ to
$\catname{E}^2$, $\funcname{f}_1$ is an m-morphism of $\seqname{C}^1$ to
$\seqname{C}^2$ and for every \\*
$(U^i, V^i, \phi^i \maps U^i \toiso V^i) \in \seqname{A}^i$,
$i=1,2$, with
$I \defeq U^1 \cap \funcname{f}^{-1}_0[U^2] \neq \emptyset$,
$
  \phi^2 \compose \funcname{f}_0 \compose \phi^{1-1} \maps
  \phi^1[I] \to U^2
$
is a morphism of $\Mod(\phi^1[I],\seqname{C}^2)$.

It is also a coherent strict
\nobreaks
  {
    {$\catname{E}^1$-$\catname{E}^2$-},
    {$\seqname{E}^1$-$\seqname{E}^2$}
  }
m-atlas (near) morphism of $\seqname{A}^1$ to $\seqname{A}^2$ in the
coordinate spaces $\seqname{C}^1$, $\seqname{C}^2$ iff
$\seqname{C}^1 \submod[strict-] \seqname{C}^2$.
$\catname{E}^1 \subcat[strict-] \catname{E}^2$ and
$\seqname{C}^1 \submod[strict-] \seqname{C}^2$.

$\funcseqname{f}$ is a strict
\nobreaks
  {
    {$\catname{E}^1$-$\catname{E}^2$-},
    {$\seqname{E}^1$-$\seqname{E}^2$-},
    {$\catname{C}^1$-$\catname{C}^2$}
  }
m-atlas (near) morphism of $\seqname{A}^1$ to $\seqname{A}^2$ in the
coordinate spaces $\seqname{C}^1$, $\seqname{C}^2$, abbreviated as \\*
$
  \strict{\isAtl[(near)]_\Ar}
  (
    \seqname{A}^1,
    \seqname{E}^1,
    \seqname{C}^1,
    \seqname{A}^2,
    \seqname{E}^2,
    \seqname{C}^2,
    \funcname{f}_0,
    \funcname{f}_1,
    \catname{E}^1,
    \catname{E}^2,
    \catname{C}^1,
    \catname{C}^2
  )
$,
iff $\funcname{f}_0$ is an m-morphism of $\catname{E}^1$ to
$\catname{E}^2$, $\funcname{f}_1$ is an m-morphism of $\catname{C}^1$ to
$\catname{C}^2$ and for every
$(U^i,V^i, \phi^i \maps U^i \toiso V^i) \in \seqname{A}^i$, $i=1,2$,
with $I \defeq U^1 \cap \funcname{f}^{-1}_0[U^2] \neq \emptyset$,
$
  \phi^2 \compose \funcname{f}_0 \compose \phi^{1-1} \maps
  \phi^1[I] \to V^2
$
is a morphism of $\catname{C}^2$.

It is also a coherent strict
\nobreaks
  {
    {$\catname{E}^1$-$\catname{E}^2$-},
    {$\seqname{E}^1$-$\seqname{E}^2$-},
    {$\catname{C}^1$-$\catname{C}^2$}
  }
m-atlas (near) morphism of $\seqname{A}^1$ to $\seqname{A}^2$ in the
coordinate spaces $\seqname{C}^1$, $\seqname{C}^2$ iff
$\catname{E}^1 \subcat[strict-] \catname{E}^2$ and
$\catname{C}^1 \subcat[strict-] \catname{C}^2$.

By abuse of language we will write $\funcseqname{f}$ is coherent when it
is a
\nobreaks
  {
    {($\catname{E}^1$-$\catname{E}^2$-)},
    {$\seqname{E}^1$-$\seqname{E}^2$},
    {(-$\catname{C}^1$-$\catname{C}^2$)}
  }
m-atlas (near) morphism of $\seqname{A}^1$ to $\seqname{A}^2$ in the
coordinate spaces $\seqname{C}^1$, $\seqname{C}^2$ where the model
spaces, model categories and atlases are understood by context.
\end{definition}

\subsubsection{Abbreviated nomenclature for m-atlas morphisms}
\begin{definition}[Abbreviated nomenclature for m-atlas (near) morphisms]
\label{def:M-ATLmorphabbrev}

Let $\catname{E}^i$, $\catname{C}^i$, $i=1,2$, be model categories,
$\seqname{E}^i \objin \catname{E}^i$,
$\seqname{C}^i \objin \catname{C}^i$, $\seqname{A}^i$ be an m-atlas of
$\seqname{E}^i$ in the coordinate space $\seqname{C}^i$ and
$
  \funcseqname{f} \defeq
  (
    \funcname{f}_0 \maps \seqname{E}^1 \to \seqname{E}^2,
    \funcname{f}_1 \maps \seqname{C}^1 \to \seqname{C}^2
  )
$
a pair of model functions.

$\funcseqname{f}$ is a constrained semi-strict~(strict)
$\seqname{E}^1$-$\seqname{E}^2$ m-atlas (near) morphism of
$\seqname{A}^1$ to $\seqname{A}^2$ in the coordinate spaces
$\seqname{C}^1$, $\seqname{C}^2$, abbreviated as \\*
$
  \strict[**]{\isAtl[(const,const-near)]_\Ar}
  (
    \seqname{A}^1,
    \seqname{E}^1,
    \seqname{C}^1,
    \seqname{A}^2,
    \seqname{E}^2,
    \seqname{C}^2,
    \funcname{f}_0,
    \funcname{f}_1
  )
$,
iff it is a constrained $\seqname{E}^1$-$\seqname{E}^2$ m-atlas (near)
morphism of $\seqname{A}^1$ to $\seqname{A}^2$ in the coordinate spaces
$\seqname{C}^1$, $\seqname{C}^1$ and it is a semi-strict~(strict)
$\seqname{E}^1$-$\seqname{E}^2$ m-atlas (near) morphism of
$\seqname{A}^1$ to $\seqname{A}^2$ in the coordinate spaces
$\seqname{C}^1$, $\seqname{C}^1$.

$\funcseqname{f}$ is a (constrained) (semi-strict, strict)
$\seqname{E}^1$-$\seqname{E}^2$ m-atlas (near) morphism of
$\seqname{A}^1$ to $\seqname{A}^2$ in the coordinate space
$\seqname{C}^1$, abbreviated as \\*
$
  \strict[*]{\isAtl[(const,const-near,near)]_\Ar}
  (
    \seqname{A}^1,
    \seqname{E}^1,
    \seqname{C}^1,
    \seqname{A}^2,
    \seqname{E}^2,
    \funcname{f}_0,
    \funcname{f}_1
  )
$,
iff it is a (constrained) (semi-strict, strict)
$\seqname{E}^1$-$\seqname{E}^2$ m-atlas (near) morphism of
$\seqname{A}^1$ to $\seqname{A}^2$ in the coordinate spaces
$\seqname{C}^1$, $\seqname{C}^1$.

$\funcseqname{f}$ is a (constrained) (semi-strict, strict)
\nobreaks
  {
    {$\catname{E}^1$-$\catname{E}^2$-},
    {$\seqname{E}^1$-$\seqname{E}^2$}
  }
m-atlas (near) morphism of $\seqname{A}^1$ to $\seqname{A}^2$ in the
coordinate space $\seqname{C}^1$, abbreviated as \\*
$
  \strict[*]{\isAtl[(const,const-near,near)]_\Ar}
  (
    \seqname{A}^1,
    \seqname{E}^1,
    \seqname{C}^1,
    \seqname{A}^2,
    \seqname{E}^2,
    \funcname{f}_0,
    \funcname{f}_1,
    \catname{E}^1,
    \catname{E}^2
  )
$,
iff it is a (constrained) (semi-strict, strict)
\nobreaks
  {
    {$\catname{E}^1$-$\catname{E}^2$-},
    {$\seqname{E}^1$-$\seqname{E}^2$}
  }
m-atlas (near) morphism of $\seqname{A}^1$ to $\seqname{A}^2$ in
the coordinate spaces $\seqname{C}^1$, $\seqname{C}^1$.

$\funcseqname{f}$ is a (constrained) (semi-strict, strict)
$\seqname{E}^1$ m-atlas (near) morphism of $\seqname{A}^1$ to
$\seqname{A}^2$ in the coordinate spaces $\seqname{C}^1$,
$\seqname{C}^2$, abbreviated as \\*
$
  \strict[*]{\isAtl[(const,const-near,near)]_\Ar}
  (
    \seqname{A}^1,
    \seqname{E}^1,
    \seqname{C}^1,
    \seqname{A}^2,
    \seqname{C}^2,
    \funcname{f}_0,
    \funcname{f}_1
    )
$,
iff it is a (constrained) (semi-strict, strict)
$\seqname{E}^1$-$\seqname{E}^1$ m-atlas (near) morphism of
$\seqname{A}^1$ to $\seqname{A}^2$ in the coordinate spaces
$\seqname{C}^1$, $\seqname{C}^2$.

$\funcseqname{f}$ is a (constrained) (semi-strict, strict)
$\seqname{E}^1$ m-atlas (near) morphism of $\seqname{A}^1$ to
$\seqname{A}^2$ in the coordinate space $\seqname{C}^1$, abbreviated as
\\*
$
  \strict[*]{\isAtl[(const,const-near,near)]_\Ar}
  (
    \seqname{A}^1,
    \seqname{E}^1,
    \seqname{C}^1,
    \seqname{A}^2,
    \funcname{f}_0,
    \funcname{f}_1
  )
$,
iff it is a (constrained) (semi-strict, strict) $\seqname{E}^1$-$\seqname{E}^1$
m-atlas (near) morphism of
$\seqname{A}^1$ to $\seqname{A}^2$ in the coordinate spaces
$\seqname{C}^1$, $\seqname{C}^1$.

$\funcseqname{f}$ is a (constrained) (semi-strict, strict)
$\seqname{E}^1$-$\catname{C}^1$-$\catname{C}^2$ m-atlas
(near) morphism of $\seqname{A}^1$ to $\seqname{A}^2$ in the coordinate
spaces $\seqname{C}^1$, $\seqname{C}^2$, abbreviated as \\*
$
  \strict[*]{\isAtl[(const,const-near,near)]_\Ar}
    (
      \seqname{A}^1,
      \seqname{E}^1,
      \seqname{C}^1,
      \seqname{A}^2,
      \seqname{C}^2,
      \funcname{f}_0,
      \funcname{f}_1,
      \catname{C}^1,
      \catname{C}^2
    )
$,
iff it is a (constrained) (semi-strict, strict)
$\seqname{E}^1$-$\seqname{E}^1$-$\catname{C}^1$-$\catname{C}^2$ m-atlas
(near) morphism of $\seqname{A}^1$ to $\seqname{A}^2$ in the coordinate
spaces $\seqname{C}^1$, $\seqname{C}^2$.

$\funcseqname{f}$ is a (constrained) (semi-strict, strict)
$\seqname{E}^1$-$\catname{C}^1$ m-atlas (near) morphism of
$\seqname{A}^1$ to $\seqname{A}^2$ in the coordinate space
$\seqname{C}^1$, abbreviated as \\*
$
  \strict[*]{\isAtl[(const,const-near,near)]_\Ar}
  (
    \seqname{A}^1,
    \seqname{E}^1,
    \seqname{C}^1,
    \seqname{A}^2,
    \funcname{f}_0,
    \funcname{f}_1,
    \catname{C}^1
  )
$,
iff it is a (constrained) (semi-strict, strict)
$\seqname{E}^1$-$\seqname{E}^1$-$\catname{C}^1$-$\catname{C}^1$ m-atlas
(near) morphism of $\seqname{A}^1$ to $\seqname{A}^2$ in the coordinate
spaces $\seqname{C}^1$, $\seqname{C}^1$.

Let $\catname{C}^i$, $i=1,2$, be a model category,
$\seqname{C}^i \objin \catname{C}^i$, $E^i$ be a topological space,
$\seqname{A}^i$ be an m-atlas of $E^i$ in the coordinate space
$\seqname{C}^i$, $\funcname{f}_0 \maps E^1 \to E^2$ be continuous,
$\funcname{f}_1 \maps \seqname{C}^1 \to \seqname{C}^2$ be a model
function and
$
  \funcseqname{f} \defeq
  (
    \funcname{f}_0 \maps E^1 \to E^2,
    \funcname{f}_1 \maps \seqname{C}^1 \to \seqname{C}^2
  )
$.

$\funcseqname{f}$ is a (constrained) (semi-strict, strict) $E^1$-$E^2$
m-atlas (near) morphism of $\seqname{A}^1$ to $\seqname{A}^2$ in the
coordinate spaces $\seqname{C}^1$, $\seqname{C}^2$, abbreviated as \\*
$
  \strict[*]{\isAtl[(const,const-near,near)]_\Ar}
    (
      \seqname{A}^1,
      E^1,
      \seqname{C}^1,
      \seqname{A}^2,
      E^2,
      \seqname{C}^2,
      \funcname{f}_0,
      \funcname{f}_1
    )
$,
and a (constrained) (semi-strict, strict) $E^1$-$E^2$ m-atlas (near)
morphism of $\seqname{A}^1$ to $\seqname{A}^2$ in the coordinate model
categories $\catname{C}^1$, $\catname{C}^2$, abbreviated as \\*
$
  \strict[*]{\isAtl[(const,const-near,near)]_\Ar}
    (
      \seqname{A}^1,
      E^1,
      \catname{C}^1,
      \seqname{A}^2,
      E^2,
      \catname{C}^2,
      \funcname{f}_0,
      \funcname{f}_1
    )
$,
iff it is a (constrained) (semi-strict, strict)
$\Triv{E^1}$-$\Triv{E^2}$ m-atlas (near) morphism of $\seqname{A}^1$ to
$\seqname{A}^2$ in the coordinate model spaces $\seqname{C}^1$,
$\seqname{C}^2$.

$\funcseqname{f}$ is a (constrained) (semi-strict, strict) $E^1$-$E^2$
m-atlas (constrained) (near) morphism of $\seqname{A}^1$ to
$\seqname{A}^2$ in the coordinate space $\seqname{C}^1$, abbreviated as
\\*
$
  \strict[*]{\isAtl[(const,const-near,near)]_\Ar}
  (
    \seqname{A}^1,
    E^1,
    \seqname{C}^1,
    \seqname{A}^2,
    E^2,
    \funcname{f}_0,
    \funcname{f}_1
  )
$,
iff it is a (constrained) (semi-strict, strict) $\Triv{E^1}$-$\Triv{E^2}$
m-atlas (near) morphism of $\seqname{A}^1$ to $\seqname{A}^2$ in the
coordinate spaces $\seqname{C}^1$, $\seqname{C}^1$.

$\funcseqname{f}$ is a (constrained) (semi-strict, strict)
\nobreaks
  {
    {$E^1$-$E^2$-},
    {$\catname{C}^1$-$\catname{C}^2$}
  }
m-atlas (near) morphism of $\seqname{A}^1$ to $\seqname{A}^2$ in the
coordinate spaces $\seqname{C}^1$, $\seqname{C}^2$, abbreviated as \\*
$
  \strict[*]{\isAtl[(const,const-near,near)]_\Ar}
  (
    \seqname{A}^1,
    E^1,
    \seqname{C}^1,
    \seqname{A}^2,
    E^2,
    \seqname{C}^2,
    \funcname{f}_0,
    \funcname{f}_1,
    \catname{C}^1,
    \catname{C}^2
  )
$,
iff it is a (constrained) (semi-strict, strict)
\nobreaks
  {
    {$\Triv{E^1}$-$\Triv{E^2}$-},
    {$\catname{C}^1$-$\catname{C}^2$}
  }
m-atlas (near) morphism of $\seqname{A}^1$ to $\seqname{A}^2$ in the
coordinate model spaces $\seqname{C}^1$, $\seqname{C}^2$.

$\funcseqname{f}$ is a (constrained) (semi-strict, strict)
$E^1$-$E^2$-$\catname{C}^1$ m-atlas (near) morphism of $\seqname{A}^1$
to $\seqname{A}^2$ in the coordinate model category $\catname{C}^1$,
abbreviated as \\*
$
  \strict[*]{\isAtl[(const,const-near,near)]_\Ar}
  (
    \seqname{A}^1,
    E^1,
    \seqname{C}^1,
    \seqname{A}^2,
    E^2,
    \funcname{f}_0,
    \funcname{f}_1,
    \catname{C}^1
  )
$,
iff it is a (constrained) (semi-strict, strict)
\nobreaks
  {
    {$\Triv{E^1}$-$\Triv{E^2}$-},
    {$\catname{C}^1$-$\catname{C}^1$}
  }
m-atlas (near) morphism of $\seqname{A}^1$ to $\seqname{A}^2$ in the
coordinate model categories $\catname{C}^1$, $\catname{C}^1$.

Similar definitions apply with restrictions on the admissible
atlases, i.e.,
$\funcseqname{f}$ is a (constrained) (semi-strict, strict)
full (semi-maximal, maximal, full semi-maximal, full maximal)
\nobreaks
  {
    {$\catname{E}^1$-$\catname{E}^2$-},
    {$\seqname{E}^1$-$\seqname{E}^2$-},
    {$\catname{C}^1$-$\catname{C}^2$}
  }
m-atlas (near) morphism of
$\seqname{A}^1$ to $\seqname{A}^2$ in the coordinate spaces
$\seqname{C}^1$, abbreviated as \\*

\begin{equation*}
\begin{array}{lll}
\strict[*]
  {
    \maxfull[**]
      {
        \isAtl[(const,const-near,near)]_\Ar
          (
            \seqname{A}^1,
            \seqname{E}^1,
            \seqname{C}^1,
            \seqname{A}^2,
            \seqname{E}^2,
            \seqname{C}^2,
            \funcname{f}_0,
            \funcname{f}_1,
            \catname{E}^1,
            \catname{E}^2,
            \catname{C}^1,
            \catname{C}^2
          )
      }
  }
  & & \iff
\\*
\strict[**]
  {
    \isAtl[(const,const-near,near)]_\Ar
    (
      \seqname{A}^1,
      \seqname{E}^1,
      \seqname{C}^1,
      \seqname{A}^2,
      \seqname{E}^2,
      \seqname{C}^2,
      \funcname{f}_0,
      \funcname{f}_1,
      \catname{E}^1,
      \catname{E}^2,
      \catname{C}^1,
      \catname{C}^2
    )
  }
  & \land
\\*
\maxfull[**]{\isAtl_\Ob}
  (
    \seqname{A}^1,
    \seqname{E}^1,
    \seqname{C}^1
  )
  & \land
\\*
\maxfull[**]{\isAtl_\Ob}
  (
    \seqname{A}^2,
    \seqname{E}^2,
    \seqname{C}^2
  )
\end{array}
\end{equation*}

\end{definition}

\subsubsection{Proclamations on m-atlas (near) morphisms}
\begin{lemma}[M-atlas (near) morphisms]
\label{lem:M-ATLmorph}
Let
\begin{enumerate}
\item
$\catname{E}^i$, $\catname{C}^i$,
$\catname{E}'^i$, $\catname{C}'^i$, $i=1,2$, be model categories

\item
$\catname{E}^i \subcat[full-] \catname{E}'^i$,
$\catname{C}^i \subcat[full-] \catname{C}'^i$

\item
$\seqname{E}^i \objin \catname{E}^i$,
$\seqname{C}^i \objin \catname{C}^i$,
$\seqname{C}'^i \objin \catname{C}'^i$,
$\seqname{C}^i \submod \seqname{C}'^i$

\item
$\funcname{f}'_1 \maps \seqname{C}'^1 \to \seqname{C}'^2$ be a model
function such that
$\funcname{f}'_1[\seqname{C}^1] \subseteq \seqname{C}^2$ and
$
  \funcname{f}_1 \defeq
  \funcname{f}'_1 \restrictto_{\seqname{C}^1,\seqname{C}^2}
$
is a model function

\item
$\seqname{A}^i$ be an M-atlas
of $\seqname{E}^i$ in the coordinate space $\seqname{C}^i$

\item
$
  \funcseqname{f} \defeq
  (
    \funcname{f}_0 \maps \seqname{E}^1 \to \seqname{E}^2,
    \funcname{f}_1 \maps \seqname{C}^1 \to \seqname{C}^2
  )
$
be a (semi-strict, strict)
\nobreaks
  {
    {($\catname{E}^1$-$\catname{E}^2$-)},
    {$\seqname{E}^1$-$\seqname{E}^2$},
    {(-$\catname{C}^1$-$\catname{C}^2$)}
  }
m-atlas (near) morphism of $\seqname{A}^1$ to $\seqname{A}^2$ in the
coordinate spaces $\seqname{C}^1$, $\seqname{C}^2$

\item
$
  \funcseqname{f}' \defeq
  (
    \funcname{f}_0 \maps \seqname{E}^1 \to \seqname{E}^2,
    \funcname{f}'_1 \maps \seqname{C}^1 \to \seqname{C}'^2
  )
$, \\*
$
  \funcseqname{f}'' \defeq
  (
    \funcname{f}_0 \maps \seqname{E}^1 \to \seqname{E}^2,
    \funcname{f}'_1 \maps \seqname{C}'^1 \to \seqname{C}'^2
  )
$
be pairs of model functions
\end{enumerate}

Then

\begin{enumerate}
\item
\label{M-ATLmorph:morphisnear}
If $\funcseqname{f}$ is an m-atlas morphism then $\funcseqname{f}$ is an
m-atlas near morphism.

\begin{proof}
Let
$(U^i, V^i, \phi^i \maps U^i \toiso V^i) \in \seqname{A}^i$, $i=1,2$,
be charts with
$I \defeq U^1 \cap \funcname{f}^{-1}_0[U^2] \neq \emptyset$,
$u^1 \in I$ and
$u^2 \defeq \funcname{f}_0(u^1)$.
Let
$
  (U'^1, V'^1, \phi^1 \maps U'^1 \toiso V'^1)
  \in \seqname{A}^1
$
and
$
  (U'^2 ,\hat{V}'^2, \phi'^2 \maps U'^2 \toiso \hat{V}'^2)
  \in \seqname{A}^2
$
be charts as in \pagecref{def:M-ATLmorph}. $U'^2 \subseteq I$ and
\crefrange{eq:xin}{eq:fphigeqgphi} of \pagecref{def:M-ATLmorphnear}
hold.
\end{proof}

\item
\label{M-ATLmorph:semimaxnearismorph}
If $\seqname{A}^i$, $i=1,2$, is semi-maximal and $\funcseqname{f}$ is a
(semi-strict, strict)
\nobreaks
  {
    {($\catname{E}^1$-$\catname{E}^2$-)},
    {$\seqname{E}^1$-$\seqname{E}^2$},
    {(-$\catname{C}^1$-$\catname{C}^2$)}
  }
m-atlas near morphism of $\seqname{A}^1$ to $\seqname{A}^2$ in the
coordinate spaces $\seqname{C}^1$, $\seqname{C}^2$ then it is a
(semi-strict, strict)
\nobreaks
  {
    {($\catname{E}^1$-$\catname{E}^2$-)},
    {$\seqname{E}^1$-$\seqname{E}^2$},
    {(-$\catname{C}^1$-$\catname{C}^2$)}
  }
m-atlas morphism of $\seqname{A}^1$ to $\seqname{A}^2$ in the coordinate
spaces $\seqname{C}^1$, $\seqname{C}^2$.

\begin{proof}
\leavevmode
\begin{description}
\item [$\funcseqname{f}$ is an M-atlas morphism.]
Let $(U^i, V^i, \phi^i \maps U^i \toiso V^i) \in \seqname{A}^i$,
$i=1,2$.  Diagram \eqref{eq:AtlN1} in \pagecref{def:M-ATLmorphnear} is
M-locally nearly commutative in $\seqname{C}^2$, i.e., for any
$e^1 \in I \defeq U^1 \cap \funcname{f}^{-1}_0[U^2]$ there are model
neighborhoods $U'^1 \subseteq I$, $V'^1 \subseteq V^1$,
$U'^2 \subseteq U^2$, $V'^2 \subseteq V^2$, $\hat{V}'^2 \subseteq C^2$
and an isomorphism $\hat{\funcname{f}} \maps V'^2 \toiso \hat{V}'^2$
such that \crefrange{eq:xin}{eq:fphigeqgphi} in
{
  \showlabelsinline
  \cref{def:M-ATLmorphnear}
}
hold, as shown in
{
  \showlabelsinline
  \crefrange{fig:AtlN1}{fig:AtlN2}.
}

$\funcname{f}_0[U'^1] \subseteq U'^2$
by \cref{eq:fup1in}
{
  \showlabelsinline
  \cref{def:M-ATLmorphnear}
}.
$\funcname{f}_1[V'^1] \subseteq \hat{V}'^2$
by \cref{eq:fvp1in}
{
  \showlabelsinline
  \cref{def:M-ATLmorphnear}
}.

Since $\seqname{A}^i$ is semi-maximal,
$(U'^i, V'^i, \phi^i \maps U'^i \toiso V'^i) \in \seqname{A}^i$ and
$
  (U'^2, \hat{\funcname{f}}[V'^2], \hat{\funcname{f}} \compose \phi^2)
  \in \seqname{A}^2
$.

Define
$
  \phi'^2 \maps U'^2 \toiso \hat{\funcname{f}}[V'^2] \defeq
  \hat{\funcname{f}} \compose \phi^2
$.
Then Diagram \eqref{eq:AtlM1} is commutative.

\item [$\funcseqname{f}$ is (semi-strict, strict).]
The definitions of strict and semi-strict are the same for m-atlas near
morphisms and for m-atlas morphisms.
\end{description}
\end{proof}

\item
\label{M-ATLmorph:fineistopmorph}
Let $E^i \defeq \Top(\seqname{E}^i)$, $i=1,2$. If each $\seqname{C}^i$
is fine grained then \\*
$
   \hat{\funcname{f}} \defeq
   (
      \funcname{f}_0 \maps E^1 \to E^2,
      \funcname{f}_1 \maps \seqname{C}^1 \to \seqname{C}^2
   )
$
is a (semi-strict, strict) $E^1$-$E^2$
m-atlas (near) morphism of $\seqname{A}^1$ to $\seqname{A}^2$ in the
coordinate spaces $\seqname{C}^1$, $\seqname{C}^2$.

\begin{proof}
\leavevmode
\begin{enumerate}
\item
$\seqname{A}^i$ is an m-atlas of $E^i$ in the coordinate space
$\seqname{C}^i$ by \pagecref{lem:M-atlas}.

\item
$\funcname{f}_0 \maps \seqname{E}^1 \to \seqname{E}^2$ is a model
function, hence continuous, hence by \pagecref{lem:modfunctriv},
$\funcname{f}_0 \maps \Triv{E^1} \to \Triv{E^2}$ is a model function.

\item
$\funcname{f}_1 \maps \seqname{C}^1 \to \seqname{C}^2$ is a model
function by hypothesis.

\item
Any model neighborhood of $\seqname{E}^i$ is open, hence a model
neighborhood of $\Triv{E^i}$.

\item
Let $(U^i,V^i,\phi^i) \in \seqname{A}^i$, $i=1,2$.  All of the model
neighborhoods of $\seqname{E}^i$ mentioned in
\pagecref{def:M-ATLmorphnear} and \pagecref{def:M-ATLmorph} are also
model neighborhoods of $\Triv{E^i}$.
\end{enumerate}
\end{proof}

\item
\label{M-ATLmorph:istosup}
$\funcseqname{f}'$ is an $\seqname{E}^1$-$\seqname{E}^2$ m-atlas (near)
morphism of $\seqname{A}^1$ to $\seqname{A}^2$ in the coordinate spaces
$\seqname{C}^1$, $\seqname{C}'^2$ and $\funcseqname{f}''$ is an
$\seqname{E}^1$-$\seqname{E}^2$ m-atlas (near) morphism of
$\seqname{A}^1$ to $\seqname{A}^2$ in the coordinate spaces
$\seqname{C}'^1$, $\seqname{C}'^2$.

\begin{proof}
$\funcname{f}_0 \maps \seqname{E}^1 \to \seqname{E}^2$ is a model
function by \pagecref{def:M-ATLmorphnear}
(\pagecref{def:M-ATLmorph}).

$\funcname{f}'_1 \maps \seqname{C}'^1 \to \seqname{C}'^2$ and
$\funcname{f}'_1 \maps \seqname{C}^1 \to \seqname{C}'^2$ are model
functions by hypothesis.

$\seqname{C}^i \submod \seqname{C}'^i$ by hypothesis,

$C^i$ is an open subset of $C'^i$ by \pagecref{lem:modsub}.

The expanded model space does not change the atlases or the commutation.
\end{proof}

\item
\label{M-ATLmorph:strictissemi}
If $\funcseqname{f}$ is a strict
\nobreaks
  {
    {($\catname{E}^1$-$\catname{E}^2$-)},
    {$\seqname{E}^1$-$\seqname{E}^2$},
    {(-$\catname{C}^1$-$\catname{C}^2$)}
  }
m-atlas (near) morphism of $\seqname{A}^1$ to $\seqname{A}^2$ in the
coordinate spaces $\seqname{C}^1$, $\seqname{C}^2$ then
$\funcseqname{f}$ is a semi-strict
\nobreaks
  {
    {($\catname{E}^1$-$\catname{E}^2$-)},
    {$\seqname{E}^1$-$\seqname{E}^2$},
    {-($\catname{C}^1$-$\catname{C}^2$)}
  }
m-atlas (near) morphism of $\seqname{A}^1$ to $\seqname{A}^2$ in the
coordinate spaces $\seqname{C}^1$, $\seqname{C}^2$

\begin{proof}
There are four overlapping cases:
\begin{description}
\item[$\seqname{E}^i$ without category:]
$\funcname{f}_0$ is an m-morphism of $\seqname{E}^1$ to $\seqname{E}^2$
by \pagecref{def:M-ATLmorphstrict}. $\funcname{f}_0$ is a local
m-morphism of $\seqname{E}^1$ to $\seqname{E}^2$ by
\pagecref{lem:modMmorph}.

\item[$\seqname{C}^i$ without category:]
$\funcname{f}_1$ is an m-morphism of $\seqname{C}^1$ to $\seqname{C}^2$
by \cref{def:M-ATLmorphstrict}. $\funcname{f}_1$ is a local m-morphism
of $\seqname{C}^1$ to $\seqname{C}^2$ by \cref{lem:modMmorph}.

\item[$\seqname{E}^i$ with category:]
$\funcname{f}_0$ is an $\catname{E}^1$-$\catname{E}^2$ m-morphism of
$\seqname{E}^1$ to $\seqname{E}^2$ by \cref{def:M-ATLmorphstrict}.
$\funcname{f}_0$ is a local $\catname{E}^1$-$\catname{E}^2$ m-morphism
of $\seqname{E}^1$ to $\seqname{E}^2$ by \cref{lem:modMmorph}.

\item[$\seqname{C}^i$ with category:]
$\funcname{f}_1$ is a $\catname{C}^1$-$\catname{C}^2$ m-morphism of
$\seqname{C}^1$ to $\seqname{C}^2$ by \cref{def:M-ATLmorphstrict}.
$\funcname{f}_1$ is a local $\catname{C}^1$-$\catname{C}^2$ m-morphism
of $\seqname{C}^1$ to $\seqname{C}^2$ by \cref{lem:modMmorph}.

\end{description}
\end{proof}

\item
\label{M-ATLmorph:strictSstrictisstrict}
If $\funcseqname{f}$ is a coherent semi-strict
\nobreaks
  {
    {($\catname{E}^1$-$\catname{E}^2$-)},
    {$\seqname{E}^1$-$\seqname{E}^2$},
    {(-$\catname{C}^1$-$\catname{C}^2$)}
  }
m-atlas (near) morphism of $\seqname{A}^1$ to $\seqname{A}^2$ in the
coordinate spaces $\seqname{C}^1$, $\seqname{C}^2$, then
$\funcseqname{f}$ is a strict
\nobreaks
  {
    {($\catname{E}^1$-$\catname{E}^2$-)},
    {$\seqname{E}^1$-$\seqname{E}^2$},
    {(-$\catname{C}^1$-$\catname{C}^2$)}
  }
m-atlas (near) morphism of $\seqname{A}^1$ to $\seqname{A}^2$ in the
coordinate spaces $\seqname{C}^1$, $\seqname{C}^2$.

\begin{proof}
There are four overlapping cases.
\begin{description}
\item[$\seqname{E}^i$ without category:]
\leavevmode
$\funcname{f}_0$ is a local m-morphism of $\seqname{E}^1$ to
$\seqname{E}^2$ by \pagecref{def:M-ATLmorphsemi}.
$\seqname{E}^1 \submod[strict-] \seqname{E}^2$ by \pagecref{def:M-ATLmorphsemi}.
$\Set(\seqname{E}^2)$ is a model neighborhood of $\seqname{E}^2$ by \pagecref{def:modsub}.
The result follows by \cref{mod:sheaf} (Restricted sheaf condition) of
{
  \showlabelsinline
  \cref{def:modsub}
}.

\item[$\seqname{E}^i$ with category:]
\leavevmode
$\funcname{f}_0$ is a local $\catname{E}^1$-$\catname{E}^2$ m-morphism
of $\seqname{E}^1$ to $\seqname{E}^2$ by \pagecref{def:M-ATLmorphsemi}.
$\catname{E}^1 \subcat[strict-] \catname{E}^2$ by
\cref{def:M-ATLmorphsemi}.  $\catname{E}^2$ satisfies the restricted
sheaf condition by \pagecref{def:modcat} and the result follows.

\item[$\seqname{C}^i$ without category:]
\leavevmode
$\funcname{f}_0$ is a local m-morphism of $\seqname{C}^1$ to
$\seqname{C}^2$ by \pagecref{def:M-ATLmorphsemi}.
$\seqname{C}^1 \submod[strict-] \seqname{C}^2$ by \pagecref{def:M-ATLmorphsemi}.
$\Set(\seqname{C}^2)$ is a model neighborhood of $\seqname{C}^2$ by \pagecref{def:modsub}.
The result follows by \cref{mod:sheaf} (Restricted sheaf condition) of
{
  \showlabelsinline
  \cref{def:modsub}
}.

\item[$\seqname{C}^i$ with category:]
\leavevmode
$\funcname{f}_0$ is a local $\catname{C}^1$-$\catname{C}^2$ m-morphism
of $\seqname{C}^1$ to $\seqname{C}^2$ by \pagecref{def:M-ATLmorphsemi}.
$\catname{C}^1 \subcat[strict-] \catname{C}^2$ by
\cref{def:M-ATLmorphsemi}.  $\catname{E}^2$ satisfies the restricted
sheaf condition by \pagecref{def:modcat} and the result follows.
\end{description}
\end{proof}

\item
\label{M-ATLmorph:SstrictisSstricttosup}
If $\funcseqname{f}$ is a semi-strict
\nobreaks
  {
    {($\catname{E}^1$-$\catname{E}^2$-)},
    {$\seqname{E}^1$-$\seqname{E}^2$},
    {(-$\catname{C}^1$-$\catname{C}^2$)}
  }
m-atlas (near) morphism of $\seqname{A}^1$ to $\seqname{A}^2$ in the
coordinate spaces $\seqname{C}^1$, $\seqname{C}^2$ then
$\funcseqname{f}'$ is a semi-strict
\nobreaks
  {
    {($\catname{E}'^1$-$\catname{E}'^2$-)},
    {$\seqname{E}^1$-$\seqname{E}^2$},
    {(-$\catname{C}'^1$-$\catname{C}'^2$)}
  }
m-atlas (near) morphism of $\seqname{A}^1$ to $\seqname{A}^2$ in the
coordinate spaces $\seqname{C}^1$, $\seqname{C}'^2$.

\begin{proof}
There are four overlapping cases.
\begin{description}
\item[$\seqname{E}^i$ without category:]
$\funcname{f}_0$ is a (strict) local m-morphism of $\seqname{E}^1$ to
$\seqname{E}^2$ by \pagecref{def:M-ATLmorphsemi}.

\item[$\seqname{E}^i$ with category:]
$\funcname{f}_0$ is a (strict) local $\catname{E}^1$-$\catname{E}^2$
m-morphism of $\seqname{E}^1$ to $\seqname{E}^2$ by
\pagecref{def:M-ATLmorphsemi}.
$\catname{E}^i \subcat[full-] \catname{E}'^i$ by hypothesis.
$\funcname{f}_0$ is a (strict) local $\catname{E}'^1$-$\catname{E}'^2$
m-morphism of $\seqname{E}^1$ to $\seqname{E}^2$ by
\pagecref{lem:modMmorph}.

\item[$\seqname{C}^i$ without category:]
$\funcname{f}_1$ is a (strict) local m-morphism of $\seqname{C}^1$ to
$\seqname{C}^2$ by \pagecref{def:M-ATLmorphsemi}.
$\seqname{C}^2 \submod \seqname{C}'^2$ by hypothesis.
$\funcname{f}_1$ is a (strict) local m-morphism of $\seqname{C}^1$ to
$\seqname{C}'^2$ by \pagecref{lem:modMmorph}.

\item[$\seqname{C}^i$ with category:]
$\funcname{f}_0$ is a local $\catname{C}^1$-$\catname{C}^2$
m-morphism of $\seqname{C}^1$ to $\seqname{C}^2$ by
\pagecref{def:M-ATLmorphsemi}.
$\catname{C}^i \subcat[full-] \catname{C}'^i$ by hypothesis.
$\funcname{f}_0$ is a (strict) local $\catname{C}'^1$-$\catname{C}'^2$
m-morphism of $\seqname{C}^1$ to $\seqname{C}^2$ by
\pagecref{lem:modMmorph}.
\end{description}
\end{proof}

\item
\label{M-ATLmorph:strictisstricttosup}
If $\funcseqname{f}$ is a strict
\nobreaks
  {
    {($\catname{E}^1$-$\catname{E}^2$-)},
    {$\seqname{E}^1$-$\seqname{E}^2$},
    {(-$\catname{C}^1$-$\catname{C}^2$)}
  }
m-atlas (near) morphism of $\seqname{A}^1$ to $\seqname{A}^2$ in the
coordinate spaces $\seqname{C}^1$, $\seqname{C}^2$ then
$\funcseqname{f}'$ is a strict
\nobreaks
  {
    {($\catname{E}'^1$-$\catname{E}'^2$-)},
    {$\seqname{E}^1$-$\seqname{E}^2$},
    {(-$\catname{C}'^1$-$\catname{C}'^2$)}
  }
m-atlas (near) morphism of $\seqname{A}^1$ to $\seqname{A}^2$ in the
coordinate spaces $\seqname{C}^1$, $\seqname{C}'^2$.

\begin{proof}
There are four overlapping cases.
\begin{description}
\item[$\seqname{E}^i$ without category:]
$\funcname{f}_0$ is an m-morphism of $\seqname{E}^1$ to
$\seqname{E}^2$ by \pagecref{def:M-ATLmorphsemi}.

\item[$\seqname{E}^i$ with category:]
$\funcname{f}_0$ is an $\catname{E}^1$-$\catname{E}^2$
m-morphism of $\seqname{E}^1$ to $\seqname{E}^2$ by
\pagecref{def:M-ATLmorphsemi}.
$\catname{E}^i \subcat[full-] \catname{E}'^i$ by hypothesis.
$\funcname{f}_0$ is an $\catname{E}'^1$-$\catname{E}'^2$
m-morphism of $\seqname{E}^1$ to $\seqname{E}^2$ by
\pagecref{lem:modMmorph}.

\item[$\seqname{C}^i$ without category:]
$\funcname{f}_1$ is an m-morphism of $\seqname{C}^1$ to
$\seqname{C}^2$ by \pagecref{def:M-ATLmorphsemi}.
$\seqname{C}^2 \submod \seqname{C}'^2$ by hypothesis.
$\funcname{f}_1$ is an m-morphism of $\seqname{C}^1$ to
$\seqname{C}'^2$ by \pagecref{lem:modMmorph}.

\item[$\seqname{C}^i$ with category:]
$\funcname{f}_0$ is an $\catname{C}^1$-$\catname{C}^2$
m-morphism of $\seqname{C}^1$ to $\seqname{C}^2$ by
\pagecref{def:M-ATLmorphsemi}.
$\catname{C}^i \subcat[full-] \catname{C}'^i$ by hypothesis.
$\funcname{f}_0$ is a $\catname{C}'^1$-$\catname{C}'^2$
m-morphism of $\seqname{C}^1$ to $\seqname{C}^2$ by
\pagecref{lem:modMmorph}.
\end{description}
\end{proof}

\item
\label{M-ATLmorph:SstrictlocalisSstricttosup}
If $\funcseqname{f}$ is a semi-strict
\nobreaks
  {
    {($\catname{E}^1$-$\catname{E}^2$-)},
    {$\seqname{E}^1$-$\seqname{E}^2$},
    {(-$\catname{C}^1$-$\catname{C}^2$)}
  }
m-atlas (near) morphism of $\seqname{A}^1$ to $\seqname{A}^2$ in the
coordinate spaces $\seqname{C}^1$, $\seqname{C}^2$ and $\funcname{f}'_1$
is a local m-morphism of $\seqname{C}'^1$ to $\seqname{C}'^2$ (a
local $\catname{C}'^1$-$\catname{C}'^2$ m-morphism) of $\seqname{C}'^1$ to
$\seqname{C}'^2$) then $\funcseqname{f}''$ is a semi-strict
\nobreaks
  {
    {($\catname{E}'^1$-$\catname{E}'^2$-)},
    {$\seqname{E}^1$-$\seqname{E}^2$},
    {(-$\catname{C}'^1$-$\catname{C}'^2$)}
  }
m-atlas (near) morphism of $\seqname{A}^1$ to $\seqname{A}^2$ in the
coordinate spaces $\seqname{C}'^1$, $\seqname{C}'^2$.

\begin{proof}
There are four overlapping cases:

\begin{description}
\item[$\seqname{E}^i$ without categories:]
$\funcname{f}_0$ is a local m-morphism of $\seqname{E}^1$ to
$\seqname{E}^2$ by \pagecref{def:M-ATLmorphsemi}.

\item[$\seqname{E}^i$ with categories:]
$\funcname{f}_0$ is a local $\catname{E}^1$-$\catname{E}^2$
m-morphism of $\seqname{E}^1$ to $\seqname{E}^2$ by
\cref{def:M-ATLmorphsemi}.
$\catname{E}^i \subcat[full-] \catname{E}'^i$ by hypothesis.
$\funcname{f}_0$ is a local $\catname{E}'^1$-$\catname{E}'^2$
m-morphism of $\seqname{E}^1$ to $\seqname{E}^2$ by
\pagecref{lem:modMmorph}.

\item[$\seqname{C}'^i$ without categories:]
$\funcname{f}'_1$ is a local m-morphism of $\seqname{C}'^1$ to
$\seqname{C}'^2$ by hypothesis.

\item[$\seqname{C}'^i$ with categories:]
$\funcname{f}'_1$ is a local $\catname{C}'^1$-$\catname{C}'^2$
m-morphism of $\seqname{C}^1$ to $\seqname{C}^2$ by hypothesis.
\end{description}
\end{proof}

\item
\label{M-ATLmorph:strictmorphisstricttosup}
If $\funcseqname{f}$ is a strict
\nobreaks
  {
    {($\catname{E}^1$-$\catname{E}^2$-)},
    {$\seqname{E}^1$-$\seqname{E}^2$},
    {(-$\catname{C}^1$-$\catname{C}^2$)}
  }
m-atlas (near) morphism of $\seqname{A}^1$ to $\seqname{A}^2$ in the
coordinate spaces $\seqname{C}^1$, $\seqname{C}^2$ and $\funcname{f}'_1$
is an m-morphism of $\seqname{C}'^1$ to $\seqname{C}'^2$ (a
$\catname{C}'^1$-$\catname{C}'^2$ m-morphism) of $\seqname{C}'^1$ to
$\seqname{C}'^2$) then $\funcseqname{f}''$ is a strict
\nobreaks
  {
    {($\catname{E}'^1$-$\catname{E}'^2$-)},
    {$\seqname{E}^1$-$\seqname{E}^2$},
    {(-$\catname{C}'^1$-$\catname{C}'^2$)}
  }
m-atlas (near) morphism of $\seqname{A}^1$ to $\seqname{A}^2$ in the
coordinate spaces $\seqname{C}'^1$, $\seqname{C}'^2$.

\begin{proof}
There are four overlapping cases:

\begin{description}
\item[$\seqname{E}^i$ without categories:]
$\funcname{f}_0$ is an m-morphism of $\seqname{E}^1$ to
$\seqname{E}^2$ by \pagecref{def:M-ATLmorph}.

\item[$\seqname{E}^i$ with categories:]
$\funcname{f}_0$ is an $\catname{E}^1$-$\catname{E}^2$
m-morphism of $\seqname{E}^1$ to $\seqname{E}^2$ by
\cref{def:M-ATLmorph}.
$\catname{E}^i \subcat[full-] \catname{E}'^i$ by hypothesis.
$\funcname{f}_0$ is an $\catname{E}'^1$-$\catname{E}'^2$
m-morphism of $\seqname{E}^1$ to $\seqname{E}^2$ by
\pagecref{lem:modMmorph}.

\item[$\seqname{C}'^i$ without categories:]
$\funcname{f}'_1$ is an m-morphism of $\seqname{C}'^1$ to
$\seqname{C}'^2$ by hypothesis.

\item[$\seqname{C}'^i$ with categories:]
$\funcname{f}'_1$ is a $\catname{C}'^1$-$\catname{C}'^2$
m-morphism of $\seqname{C}^1$ to $\seqname{C}^2$ by hypothesis.
\end{description}
\end{proof}
\end{enumerate}
\end{lemma}

\begin{corollary}[M-atlas morphisms]
\label{cor:M-ATLmorph}
Let
\begin{enumerate}
\item
$\catname{C}^i$, $\catname{C}'^i$, $i=1,2$, be model categories

\item
$\catname{C}^i \subcat[full-] \catname{C}'^i$

\item
$\seqname{C}^i \objin \catname{C}^i$

\item
$\seqname{C}'^i \objin \catname{C}'^i$

\item
$\seqname{C}^i \submod \seqname{C}'^i$

\item
$\funcname{f}'_1 \maps \seqname{C}'^1 \to \seqname{C}'^2$ be a model
function such that
$\funcname{f}'_1[\seqname{C}^1] \subseteq \seqname{C}^2$ and
$
  \funcname{f}_1 \defeq
  \funcname{f}'_1 \restrictto_{\seqname{C}^1,\seqname{C}^2}
$
is a model function

\item
$E^i$ be a topological space

\item
$\seqname{A}^i$ be an m-atlas of $E^i$ in the coordinate space
$\seqname{C}^i$

\item
$
  \funcseqname{f} \defeq
  (
    \funcname{f}_0 \maps E^1 \to E^2,
    \funcname{f}_1 \maps \seqname{C}^1 \to \seqname{C}^2
  )
$
be a (semi-strict, strict)
\nobreaks
  {
    {($\catname{E}^1$-$\catname{E}^2$-)},
    {$E^1$-$E^2$},
    {(-$\catname{C}^1$-$\catname{C}^2$)}
  }
m-atlas (near) morphism of $\seqname{A}^1$ to $\seqname{A}^2$ in the
coordinate spaces $\seqname{C}^1$, $\seqname{C}^2$
\end{enumerate}

If
$
  \funcseqname{f} \defeq
  (
    \funcname{f}_0 \maps E^1 \to E^2,
    \funcname{f}_1 \maps \seqname{C}^1 \to \seqname{C}^2
  )
$
is a (semi-strict, strict) $E^1$-$E^2$ m-atlas morphism of
$\seqname{A}^1$ to $\seqname{A}^2$ in the coordinate spaces
$\seqname{C}^1$, $\seqname{C}^2$ then it is a (semi-strict, strict)
$E^1$-$E^2$ m-atlas near morphism of $\seqname{A}^1$ to $\seqname{A}^2$
in the coordinate spaces $\seqname{C}^1$, $\seqname{C}^2$

If $\seqname{A}^i$, $i=1,2$, is semi-maximal and $\funcseqname{f}$ is a
(semi-strict, strict) $E^1$-$E^2$ m-atlas near morphism of
$\seqname{A}^1$ to $\seqname{A}^2$ in the coordinate spaces
$\seqname{C}^1$, $\seqname{C}^2$ then it is a (semi-strict, strict)
$E^1$-$E^2$ m-atlas morphism of $\seqname{A}^1$ to $\seqname{A}^2$ in
the coordinate spaces $\seqname{C}^1$, $\seqname{C}^2$.

If $\funcseqname{f}$ is a semi-strict (strict)
$E^1$-$E^2$-$\catname{C}^1$-$\catname{C}^2$ m-atlas (near) morphism of
$\seqname{A}^1$ to $\seqname{A}^2$ in the coordinate spaces
$\seqname{C}^1$, $\seqname{C}^2$ then $\funcseqname{f}$ is a semi-strict
(strict) $E^1$-$E^2$-$\catname{C}'^1$-$\catname{C}'^2$ m-atlas (near)
morphism of $\seqname{A}^1$ to $\seqname{A}^2$ in the coordinate spaces
$\seqname{C}^1$, $\seqname{C}^2$.

\begin{proof}
The result follows from \pagecref{def:M-ATLmorphabbrev}.
\end{proof}
\end{corollary}

\begin{lemma}[Equivalence of m-atlas (near) morphisms]
\label{lem:M-ATLmorphequiv}
Let
\begin{enumerate}
\item
$\catname{E}^i$, $\catname{C}^i$, $i=1,2$, be model categories
\item
$\seqname{E}^i \objin \catname{E}^i$,
$\seqname{C}^i \objin \catname{C}^i$,
\item
$\seqname{A}^i$ be an m-atlas of
$\seqname{E}^i$ in the coordinate space $\seqname{C}^i$,
\item
$
  \funcseqname{f} \defeq
  (
    \funcname{f}_0 \maps \seqname{E}^1 \to \seqname{E}^2,
    \funcname{f}_1 \maps \seqname{C}^1 \to \seqname{C}^2
  )
$ and
$
  \funcseqname{g} \defeq
  (
    \funcname{g}_0 \maps \seqname{E}^1 \to \seqname{E}^2,
    \funcname{g}_1 \maps \seqname{C}^1 \to \seqname{C}^2
  )
$
be equivalent m-atlas near morphisms of $\seqname{A}^1$ to
$\seqname{A}^2$ in the coordinate spaces $\seqname{C}^1$,
$\seqname{C}^2$.
\end{enumerate}

If $\funcseqname{f}$ is a (coherent) semi-strict
\nobreaks
  {
    {($\catname{E}^1$-$\catname{E}^2$-)},
    {$\seqname{E}^1$-$\seqname{E}^2$},
    {(-$\catname{C}^1$-$\catname{C}^2$)}
  }
m-atlas near morphism of $\seqname{A}^1$ to $\seqname{A}^2$ in the
coordinate spaces $\seqname{C}^1$, $\seqname{C}^2$ then
$\funcseqname{g}$ is a (coherent) semi-strict
\nobreaks
  {
    {($\catname{E}^1$-$\catname{E}^2$-)},
    {$\seqname{E}^1$-$\seqname{E}^2$},
    {(-$\catname{C}^1$-$\catname{C}^2$)}
  }
m-atlas near morphism of $\seqname{A}^1$ to $\seqname{A}^2$ in the
coordinate spaces $\seqname{C}^1$, $\seqname{C}^2$.

\begin{proof}
Let $v^1 \in \seqname{C}^1$ be an arbitrary point. Let
$(U^1,V^1,\phi^1) \in \seqname{A}^1$ be an arbitrary chart such that
$v^1 \in V^1$, $u^1 \defeq \phi^{1-1}(v^1)$,
$(U^2,V^2,\phi^2) \in \seqname{A}^2$ be an arbitrary chart at
$u^2 \defeq \funcname{f}_0(u^1)$, $v^2 \defeq \phi^2(u^2)$ and
$I \defeq U^1 \cap \funcname{f}_0^{-1}[U^2]$, as shown in
\fullcref{fig:AtlE1}.

\begin{figure}[ht]
\[ \bfig
\node v1f(0,0)[v^1]
\node V1f(1000,0)[V^1]
\node C2f(1500,0)[C^2]
\node V2f(2000,0)[V^2]
\node U1f(1000,-500)[U^1]
\node u1(0,-1000)[u^1]
\node Uint(1000,-1000)[{I=U^1 \cap \funcname{f}^{-1}_0[U^2]}]
\node U2(2000,-1000)[U^2]
\node U1g(1000,-1500)[U^1]
\node v1g(0,-2000)[v^1]
\node V1g(1000,-2000)[V^1]
\node C2g(1500,-2000)[C^2]
\node V2g(2000,-2000)[V^2]
%
% y=0
\arrow |a|/^{ (}->/[v1f`V1f;\funcname{i}]
\arrow |a|[V1f`C2f;\funcname{f}_1]
%
% y=-500
\arrow |r|[U1f`V1f;\phi^1]
\arrow |l|//[U1f`V1f;\iso]
\arrow |r|[U1f`U2;\funcname{f}_0]
\arrow |r|[U2`V2f;\phi^2]
\arrow |l|//[U2`V2f;\iso]
%
% y=-1000
\arrow |r|[u1`v1f;\phi^1]
\arrow |a|/^{ (}->/[u1`U1f;\funcname{i}]
\arrow |a|/^{ (}->/[u1`Uint;\funcname{i}]
\arrow |a|/^{ (}->/[u1`U1g;\funcname{i}]
\arrow |r|[u1`v1g;\phi^1]
\arrow |r|/_{ (}->/[Uint`U1f;\funcname{i}]
\arrow |a|[Uint`U2;\funcname{f}_0]
\arrow |r|/^{ (}->/[Uint`U1g;\funcname{i}]
%
% y=-1500
\arrow |a|[U1g`U2;\funcname{f}_0]
\arrow |r|[U1g`V1g;\phi^1]
\arrow |l|//[U1g`V1g;\iso]
\arrow |r|[U2`V2g;\phi^2]
\arrow |l|//[U2`V2g;\iso]
%
% y=-2000
\arrow |a|/^{ (}->/[v1g`V1g;\funcname{i}]
\arrow |a|[V1g`C2g;\funcname{g}_1]
\efig \]
\caption{Uncompleted equivalent m-atlas near morphisms}
\label{fig:AtlE1}
\end{figure}

Let $^{f}U'^1 \subseteq I$, $^{f}V'^1 \subseteq V^1$
$^{f}U'^2 \subseteq U^2$. $^{f}V'^2 \subseteq V^2$ and
$\hat{^{f}V'^2} \subseteq \seqname{C}^2$ be model neighborhoods and
$\hat{\funcname{f}} \maps ^{f}V'^2 \toiso \hat{^{f}V'^2}$ an isomorphism
of $\seqname{C}^2$ ($\catname{C}^2$) as in
\pagecref{def:M-ATLmorphnear}.

Let $^{g}U'^1 \subseteq I$, $^{g}V'^1 \subseteq V^1$
$^{g}U'^2 \subseteq U^2$. $^{g}V'^2 \subseteq V^2$ and
$\hat{^{g}V'^2} \subseteq \seqname{C}^2$ be model neighborhoods and
$\hat{\funcname{g}} \maps ^{g}V'^2 \toiso \hat{^{g}V'^2}$ an isomorphism
of $\seqname{C}^2$ ($\catname{C}^2$) as in \cref{def:M-ATLmorphnear}.

$\hat{^{f}V'^2} \cap \hat{^{g}V'^2}$ is a model neighborhood of
$\seqname{C}^2$ by
{
  \showlabelsinline
  \cref{mod:closed}
}
of \pagecref{def:model}.

There are four overlapping cases.
\begin{description}
\item[$\seqname{E}^i$ without category:]
$\funcname{f}_0 = \funcname{g}_0$ by \pagecref{def:M-ATLmorphequiv}.
$\funcname{f}_0$ is a local m-morphism of $\seqname{E}^1$ to
$\seqname{E}^2$ by \pagecref{def:M-ATLmorphsemi}.

If $\funcseqname{f}$ is coherent then
$\seqname{E}^1 \submod[strict-] \seqname{E}^2$ by
\pagecref{def:M-ATLmorphsemi}.

\item[$\seqname{C}^i$ without category:]
$\funcname{f}_1$ is a local m-morphism of $\seqname{C}^1$ to
$\seqname{C}^2$. \\*
$
  \phi^2 \compose \funcname{f}_0 \compose \phi^{1-1} \maps
  \phi^1[I] \to V^2
$
is a local m-morphism of $\phi^1[I]$ to $\seqname{C}^2$ by
\pagecref{def:M-ATLmorphsemi}.

Let $\hat{V}'^2 \subseteq \hat{^{f}V'^2} \cap \hat{^{g}V'^2}$ be a model
neighborhod of $\seqname{C}^2$ at $v^2$ such that \\*
$
  \funcname{f}_1 \maps V'^1 \defeq \funcname{f}_1^{-1}[\hat{V}'^2)
  \to \hat{V}'^2
$
is a morphism of $\seqname{C}^2$.

The following, shown in
{
  \showlabelsinline
  \crefrange{fig:AtlE2}{fig:AtlE3},
}
are model
neighborhoods of $\seqname{C}^2$ by \pagecref{def:modfunc}:
\begin{align*}
V'^2 & \defeq \hat{f}^{-1}[\hat{V}'^2]
\\*
U'^2 & \defeq \phi^{2-1}[V'^2]
\\*
U'^1 & \defeq \funcname{f}^{-1}_0[U'^2]
\end{align*}

\begin{figure}[ht!]
\[ \bfig
\node v1f(1000,500)[v^1]
%
\node V1f(0,0)[V^1]
\node fVp1(1000,0)[{^f}V'^1]
\node fVhp2(1500,0)[\hat{^{f}V'^2}]
\node fVp2(2000,0)[{^f}V'^2]
\node V2f(2500,0)[V^2]
\node u2(2000,0)[u^2]
%
\node U1f(0,-500)[U^1]
\node fUp1(1000,-500)[{^f}U'^1]
\node fUp2(2000,-500)[{^f}U'^2]
%
\node u1f(600,-1100)[u^1]
%
\node Uint(0,-1500)[{I=U^1 \cap \funcname{f}^{-1}_0[U^2]}]
\node Up1(1000,-1500)[U'^1]
\node Up2(2000,-1500)[U'^2]
\node U2(2500,-1500)[U^2]
%
\node U1g(0,-2000)[U^1]
\node gUp1(1000,-2000)[^{g}U'^1]
\node gUp2(2000,-2000)[^{g}U'^2]
%
\node V1g(0,-2500)[V^1]
\node gVp1(1000,-2500)[^{g}V'^1]
\node gVhp2(1500,-2500)[\hat{^{g}V'^2}]
\node gVp2(2000,-2500)[^{g}V'^2]
\node V2g(2500,-2500)[V^2]
%
\node v1g(1000,-3000)[v^1]
%
% y=500
\arrow |r|/^{ (}->/[v1f`fVp1;\funcname{i}]
%
% y=0
\arrow |a|/_{ (}->/[fVp1`V1f;\funcname{i}]
\arrow |a|[fVp1`fVhp2;\funcname{f}_1]
\arrow |a|[fVp2`fVhp2;\hat{\funcname{f}}]
\arrow |b|[fVp2`fVhp2;\iso]
\arrow |a|/^{ (}->/[fVp2`V2f;\funcname{i}]
%
% y=-500
\arrow |r|[U1f`V1f;\phi^1]
\arrow |l|//[U1f`V1f;\iso]
\arrow |r|[fUp1`fVp1;\phi^1]
\arrow |l|//[fUp1`fVp1;\iso]
\arrow |a|/_{ (}->/[fUp1`U1f;\funcname{i}]
\arrow |a|[fUp1`fUp2;\funcname{f}_0]
\arrow |a|/_{ (}->/[fUp1`Uint;\funcname{i}]
\arrow |r|[fUp2`fVp2;\phi^2]
\arrow |l|//[fUp2`fVp2;\iso]
\arrow |a|/^{ (}->/[fUp2`U2;\funcname{i}]
%
% y=-1000
\arrow |a|/_{ (}->/[u1f`Up1;\funcname{i}]
%
% y=-1500
\arrow |a|/>->/[Uint`U1f;\funcname{i}]
\arrow |a|/>->/[Uint`U1g;\funcname{i}]
\arrow |a|/_{ (}->/[Up1`fUp1;\funcname{i}]
\arrow |a|/_{ (}->/[Up1`Uint;\funcname{i}]
\arrow |r|/^{ (}->/[Up1`gUp1;\funcname{i}]
\arrow |a|[Up1`Up2;\funcname{f}_0]
\arrow |a|/^{ (}->/[Up2`U2;\funcname{i}]
\arrow |a|/^{ (}->/[Up2`fUp2;\funcname{i}]
\arrow |r|/^{ (}->/[gUp2`U2;\funcname{i}]
\arrow |r|/^{ (}->/[Up2`gUp2;\funcname{i}]
%
% y=-2000
\arrow |r|[U1g`V1g;\phi^ 1]
\arrow |a|/_{ (}->/[gUp1`Uint;\funcname{i}]
\arrow |a|/_{ (}->/[gUp1`U1g;\funcname{i}]
\arrow |a|[gUp1`gUp2;\funcname{g}_0]
\arrow |r|[gUp1`gVp1;\phi^1]
\arrow |l|//[gUp1`gVp1;\iso]
\arrow |a|/^{ (}->/[gUp2`U2;\funcname{i}]
\arrow |r|[gUp2`gVp2;\phi^2]
\arrow |l|//[gUp2`gVp2;\iso]
\arrow |r|[U2`V2f;\phi^2]
\arrow |r|[U2`V2g;\phi^2]
\arrow |l|//[U2`V2g;\iso]
%
% y=-2500
\arrow |a|/_{ (}->/[gVp1`V1g;\funcname{i}]
\arrow |a|[gVp1`gVhp2;\funcname{g}_1]
\arrow |a|[gVp2`gVhp2;\hat{\funcname{g}}]
\arrow |b|[gVp2`gVhp2;\iso]
\arrow |a|/^{ (}->/[gVp2`V2g;\funcname{i}]
%
% y=-3000
\arrow |r|/_{ (}->/[v1g`gVp1;\funcname{i}]
\efig \]
\caption{Completed equivalent m-atlas near morphisms}
\label{fig:AtlE2}
\end{figure}

\begin{figure}[ht!]
\[ \bfig
% y=0
\node Up1f(1000,0)[U'^1]
\node Up2f(2000,0)[U'^2]
%
% y=-500
\node u1f(0,-500)[u^1]
\node fUp1(1000,-500)[{^f}U'^1]
\node fUp2(2000,-500)[{^f}U'^2]
%
% y=-1000
\node fVp1(1000,-1000)[{^f}V'^1]
\node fVhp2(1500,-1000)[\hat{^{f}V'^2}]
\node fVp2(2000,-1000)[{^f}V'^2]
%
% y=-1500
\node v1(0,-1500)[v^1]
\node Vp1(1000,-1500)[V'^1]
\node Vhp2(1500,-1500)[\hat{V}'^2]
\node Vp2(2000,-1500)[V'^2]
%
% y=-2000
\node gVp1(1000,-2000)[^{g}V'^1]
\node gVhp2(1500,-2000)[\hat{^{g}V'^2}]
\node gVp2(2000,-2000)[{^g}V'^2]
%
% y=-2500
\node u1g(0,-2500)[u^1]
\node gUp1(1000,-2500)[{^g}U'^1]
\node gUp2(2000,-2500)[{^g}U'^2]
%
% y=-3000
\node Up1g(1000,-3000)[U'^1]
\node Up2g(2000,-3000)[U'^2]
%
% y=0
\arrow |a|[Up1f`Up2f;\funcname{f}_0]
\arrow |a|/^{ (}->/[Up1f`fUp1;\funcname{i}]
\arrow |r|/^{ (}->/[Up2f`fUp2;\funcname{i}]
%
% y=-500
\arrow |a|/^{ (}->/[u1f`Up1f;\funcname{i}]
\arrow |a|/^{ (}->/[u1f`fUp1;\funcname{i}]
\arrow |a|[u1f`fVp1;\phi^1]
\arrow |b|//[u1f`fVp1;\iso]
\arrow |a|[fUp1`fUp2;\funcname{f}_0]
\arrow |r|[fUp1`fVp1;\phi^1]
\arrow |l|//[fUp1`fVp1;\iso]
\arrow |r|[fUp2`fVp2;\phi^2]
\arrow |l|//[fUp2`fVp2;\iso]
%
% y=-1000
\arrow |a|[fVp1`fVhp2;\funcname{f}_1]
\arrow |a|[fVp2`fVhp2;\hat{\funcname{f}}]
\arrow |b|//[fVp2`fVhp2;\iso]
%
% y=-1500
\arrow |a|/^{ (}->/[v1`fVp1;\funcname{i}]
\arrow |a|/^{ (}->/[v1`Vp1;\funcname{i}]
\arrow |a|/^{ (}->/[v1`gVp1;\funcname{i}]
\arrow |r|/_{ (}->/[Vp1`fVp1;\funcname{i}]
\arrow |a|[Vp1`fVhp2;\funcname{f}_1]
\arrow |a|[Vp1`gVhp2;\funcname{g}_1]
\arrow |r|/^{ (}->/[Vp1`gVp1;\funcname{i}]
\arrow |r|/_{ (}->/[Vp2`fVp2;\funcname{i}]
\arrow |r|/^{ (}->/[Vp2`gVp2;\funcname{i}]
%
% y=-2000
\arrow |a|[gVp1`gVhp2;\funcname{g}_1]
\arrow |a|[gVp2`gVhp2;\hat{\funcname{g}}]
\arrow |b|//[gVp2`gVhp2;\iso]
%
% y=-2500
\arrow |a|[u1g`gVp1;\phi^1]
\arrow |b|//[u1g`gVp1;\iso]
\arrow |a|/^{ (}->/[u1g`gUp1;\funcname{i}]
\arrow |a|/^{ (}->/[u1g`Up1g;\funcname{i}]
\arrow |r|[gUp1`gVp1;\phi^1]
\arrow |a|[gUp1`gUp2;\funcname{f}_0]
\arrow |r|[gUp2`gVp2;\phi^2]
\arrow |l|//[gUp2`gVp2;\iso]
%
% y=-3000
\arrow |a|[Up1g`Up2g;\funcname{f}_0]
\arrow |a|/_{ (}->/[Up1g`gUp1;\funcname{i}]
\arrow |a|/_{ (}->/[Up2g`gUp2;\funcname{i}]
\efig \]
\caption{Completed equivalent m-atlas near morphisms}
\label{fig:AtlE3}
\end{figure}

Then
\begin{equation*}
\begin{array} {llll}
  \funcname{g}_1 \maps V'^1 \to \hat{V}'^2
  & =
  & \hat{\funcname{g}} \maps V'^2 \to \hat{V}'^2
  & \compose
\\*
  &
  & \phi^2 \maps U'^2 \to V'^2
  & \compose
\\*
  &
  & \funcname{f}_0 \maps U'^1 \to U'^2
  & \compose
\\*
  &
  & \phi^{1-1} \maps V'^1 \to U'^1
\\*
  & =
  & \hat{\funcname{g}} \maps  V'^2 \to \hat{V}'^2
  & \compose
\\*
  &
  & \hat{\funcname{f}}^{-1} \maps  \hat{V}'^2 \to  V'^2
  & \compose
\\*
  &
  & \hat{\funcname{f}} \maps V'^2 \to \hat{V}'^2
  & \compose
\\*
  &
  & \phi^2 \maps  U'^2 \to '^2
  & \compose
\\*
  &
  & \funcname{f}_0 \maps U'^1 \to U'^2
  & \compose
\\*
  &
  & \phi^{1-1} \maps V'^1 \to U'^1
\\*
  & =
  & \hat{\funcname{g}} \maps  V'^2 \to \hat{V}'^2
  & \compose
\\*
  &
  & \hat{\funcname{f}}^{-1} \maps  \hat{V}'^2 \to  V'^2
  & \compose
\\*
  &
  & \funcname{f}_1 \maps V'^1 \to \hat{V}'^2
\end{array}
\end{equation*}
is a composition of morphisms of $\seqname{C}^2$, hence a morphism.

If $\funcseqname{f}$ is coherent then
$\seqname{C}^1 \submod[strict-] \seqname{C}^2$ by
\pagecref{def:M-ATLmorphsemi}.

\item[$\seqname{E}^i$ with category:]
$\funcname{f}_0 = \funcname{g}_0$ by \cref{def:M-ATLmorphequiv}.
$\funcname{f}_0$ is a local $\catname{E}^1$-$\catname{E}^2$
m-morphism of $\seqname{E}^1$ to $\seqname{E}^2$ by
\cref{def:M-ATLmorphsemi}.

If $\funcseqname{f}$ is coherent then
$\catname{E}^1 \subcat[strict-] \catname{E}^2$ by
\pagecref{def:M-ATLmorphsemi}.

\item[$\seqname{C}^i$ with category:]
$\funcname{f}_1$ is a local $\catname{C}^1$-$\catname{C}^2$ m-morphism
of $\seqname{C}^1$ to $\seqname{C}^2$ and \\*
$
  \phi^2 \compose \funcname{f}_0 \compose \phi^{1-1} \maps
  \phi^1[I] \to V^2
$
is a local $\catname{C}^1$-$\catname{C}^2$ m-morphism of
$\Mod(\phi^1[I],\seqname{C}^2)$ to $\Mod(V^2,\seqname{C}^2)$ by
\pagecref{def:M-ATLmorphsemi}.

Let $U'^2 \subseteq {}^f\!U'^2 \cap ^{g}U'^2$ be a model neighborhod of
$\seqname{C}^2$ at $u^2$ such that \\*
$\funcname{f}_0 \maps U'^1 \defeq \funcname{f}_0^{-1}[U'^2] \to U'^2$ is
a morphism of $\catname{C}^2$.

Let $V'^2 \defeq \phi^2[U'^2]$,
$\hat{^{f}V'^2} \defeq \hat{\funcname{f}}[^{f}V'^2]$,
$\hat{^{g}V'^2} \defeq \hat{\funcname{g}}[^{g}V'^2]$,
$^{f}V'^1 \defeq \funcname{f}_1^{-1}[\hat{^{f}V'^2}]$ and
$^{g}V'^1 \defeq \funcname{g}_1^{-1}[\hat{^{g}V'^2}]$.
All of these are model neighborhoods by \pagecref{def:modfunc} and the
corresponding relative model subspaces of $\seqname{V}^i$ are objects of
$\catname{V}^i$ by
{
  \showlabelsinline
  \cref{modcat:rest}
}
of \pagecref{def:modcat}. The restrictions of morphisms to them are
morphisms by
{
  \showlabelsinline
  \cref{modcat:rest}
}
of \cref{def:modcat}.

Then
\begin{equation*}
\begin{array} {llll}
  \funcname{g}_1 \maps V'^1 \to \hat{V}'^2
  & =
  & \hat{\funcname{g}} \maps V'^2 \to \hat{V}'^2
  & \compose
\\*
  &
  & \phi^2 \maps U'^2 \to V'^2
  & \compose
\\*
  &
  & \funcname{f}_0 \maps U'^1 \to U'^2
  & \compose
\\*
  &
  & \phi^{1-1} \maps V'^1 \to U'^1
\\*
  & =
  & \hat{\funcname{g}} \maps  V'^2 \to \hat{V}'^2
  & \compose
\\*
  &
  & \hat{\funcname{f}}^{-1} \maps  \hat{V}'^2 \to  V'^2
  & \compose
\\*
  &
  & \hat{\funcname{f}} \maps V'^2 \to \hat{V}'^2
  & \compose
\\*
  &
  & \phi^2 \maps  U'^2 \to '^2
  & \compose
\\*
  &
  & \funcname{f}_0 \maps U'^1 \to U'^2
  & \compose
\\*
  &
  & \phi^{1-1} \maps V'^1 \to U'^1
\\*
  & =
  & \hat{\funcname{g}} \maps  V'^2 \to \hat{V}'^2
  & \compose
\\*
  &
  & \hat{\funcname{f}}^{-1} \maps  \hat{V}'^2 \to  V'^2
  & \compose
\\*
  &
  & \funcname{f}_1 \maps V'^1 \to \hat{V}'^2
\end{array}
\end{equation*}
is a composition of morphisms of $\catname{C}^2$, hence a morphism.

If $\funcseqname{f}$ is coherent then
$\catname{C}^1 \subcat[strict-] \catname{C}^2$ by
\pagecref{def:M-ATLmorphsemi}.
\end{description}
\end{proof}
\end{lemma}

\begin{corollary}[Equivalence of m-atlas (near) morphisms]
\label{cor:M-ATLmorphequiv}
Let
\begin{enumerate}
\item
$\catname{E}^i$, $\catname{C}^i$, $i=1,2$, be model categories
\item
$\seqname{E}^i \objin \catname{E}^i$
\item
$\seqname{C}^i \objin \catname{C}^i$
\item
$\seqname{A}^i$ be an m-atlas of $\seqname{E}^i$ in the coordinate space
$\seqname{C}^i$
\item
$\seqname{A}^1$ be full
\item
$
  \funcseqname{f} \defeq
  (
    \funcname{f}_0 \maps \seqname{E}^1 \to \seqname{E}^2,
    \funcname{f}_1 \maps \seqname{C}^1 \to \seqname{C}^2
  )
$
be a coherent semi-strict
\nobreaks
  {
    {($\catname{E}^1$-$\catname{E}^2$-)},
    {$\seqname{E}^1$-$\seqname{E}^2$},
    {(-$\catname{C}^1$-$\catname{C}^2$)}
  }
m-atlas (near) morphism of $\seqname{A}^1$ to $\seqname{A}^2$ in the
coordinate spaces $\seqname{C}^1$, $\seqname{C}^2$.
\item
$
  \funcseqname{g} \defeq
  (
    \funcname{g}_0 \maps \seqname{E}^1 \to \seqname{E}^2,
    \funcname{g}_1 \maps \seqname{C}^1 \to \seqname{C}^2
  )
$
be an m-atlas (near) morphisms of $\seqname{A}^1$ to $\seqname{A}^2$ in
the coordinate spaces $\seqname{C}^1$, $\seqname{C}^2$ equivalent to
$\funcseqname{f}$
\end{enumerate}

Then $\funcseqname{g}$ is a strict
\nobreaks
  {
    {($\catname{E}^1$-$\catname{E}^2$-)},
    {$\seqname{E}^1$-$\seqname{E}^2$},
    {(-$\catname{C}^1$-$\catname{C}^2$)}
  }
m-atlas (near) morphism of $\seqname{A}^1$ to $\seqname{A}^2$ in the
coordinate spaces $\seqname{C}^1$, $\seqname{C}^2$.

\begin{proof}
$\funcseqname{g}$ is a coherent semi-strict
\nobreaks
  {
    {($\catname{E}^1$-$\catname{E}^2$-)},
    {$\seqname{E}^1$-$\seqname{E}^2$},
    {(-$\catname{C}^1$-$\catname{C}^2$)}
  }
m-atlas (near) morphism of $\seqname{A}^1$ to $\seqname{A}^2$ in the
coordinate spaces $\seqname{C}^1$, $\seqname{C}^2$ by
\pagecref{lem:M-ATLmorphequiv}.
The result follows from \cref{M-ATLmorph:strictSstrictisstrict} of
{
  \showlabelsinline
  \pagecref{lem:M-ATLmorph}.
}
\end{proof}
\end{corollary}

\begin{lemma}[Composition of m-atlas (near) morphisms]
\label{lem:M-ATLmorphcomp}
Let
\begin{enumerate}
\item
$\catname{E}^i$, $\catname{C}^i$, $i=1,2,3$, be a model category,
$\seqname{E}^i \objin \catname{E}^i$,
$\seqname{C}^i \objin \catname{C}^i$,
\item
$\seqname{A}^i$ an m-atlas of
$\seqname{E}^i$ in the coordinate space $\seqname{C}^i$
\item
$
  \funcseqname{f}^i \defeq
  (
    \funcname{f}^i_0 \maps \seqname{E}^i \to \seqname{E}^{i+1},
    \funcname{f}^i_1 \maps \seqname{C}^i \to \seqname{C}^{i+1}
  )
$
and
$
  \funcseqname{g}^i \defeq
  (
    \funcname{g}^i_0 \maps \seqname{E}^i \to \seqname{E}^{i+1},
    \funcname{g}^i_1 \maps \seqname{C}^i \to \seqname{C}^{i+1}
  )
$
$i=1,2$,
be equivalent (semi-strict, strict)
\nobreaks
  {
    {($\catname{E}^i$-$\catname{E}^{i+1}$-)},
    {$\seqname{E}^i$-$\seqname{E}^{i+1}$},
    {(-$\catname{C}^i$-$\catname{C}^{i+1}$)}
  }
m-atlas (near) morphisms of $\seqname{A}^i$ to $\seqname{A}^{i+1}$ in
the coordinate spaces $\seqname{C}^i$, $\seqname{C}^{i+1}$.
\end{enumerate}

Then:
\begin{enumerate}
\item
\label{M-morphcomp:nearcompmorph}
If $\funcseqname{f}^1$ is an
\nobreaks
  {
    {($\catname{E}^1$-$\catname{E}^2$-)},
    {$\seqname{E}^1$-$\seqname{E}^2$},
    {(-$\catname{C}^1$-$\catname{C}^2$)}
  }
m-atlas morphism of $\seqname{A}^1$ to $\seqname{A}^2$ in the
coordinate spaces $\seqname{C}^1$, $\seqname{C}^2$ and
$\funcseqname{f}^2$ is an
\nobreaks
  {
    {($\catname{E}^2$-$\catname{E}^3$-)},
    {$\seqname{E}^2$-$\seqname{E}^3$},
    {(-$\catname{C}^2$-$\catname{C}^3$)}
  }
m-atlas near morphism of $\seqname{A}^2$ to $\seqname{A}^3$ in the
coordinate spaces $\seqname{C}^2$, $\seqname{C}^3$ then
$\funcseqname{f}^2 \compose[()] \funcseqname{f}^1$ is an
\nobreaks
  {
    {($\catname{E}^1$-$\catname{E}^3$-)},
    {$\seqname{E}^1$-$\seqname{E}^3$},
    {(-$\catname{C}^1$-$\catname{C}^3$)}
  }
m-atlas near morphism of $\seqname{A}^1$
to $\seqname{A}^3$ in the coordinate spaces $\seqname{C}^1$,
$\seqname{C}^3$.

\begin{proof}
$\funcname{f}^2_0 \compose \funcname{f}^1_0$ and
$\funcname{f}^2_1 \compose \funcname{f}^1_1$ are model functions by
\pagecref{lem:modcomp}.

Let $(U^i, V^i, \phi^i) \in \seqname{A}^i$, $i=1,3$, be a chart such
that
$
  I \defeq
  U^1 \cap (\funcname{f}^2_0 \compose \funcname{f}^1_0)^{-1}[U^3]
  \neq \emptyset
$.
It suffices to show that diagram \eqref{eq:morphcompN1} below is locally
nearly commutative.

\begin{equation}
\label{eq:morphcompN1}
\begin{array} {lll}
D \defeq & \bigl \{
\\
  &&
  \phi^1 \maps I \into V^1,
  \funcname{f}^2_1 \compose \funcname{f}^1_1 \maps  V^1 \to \seqname{C}^3,
\\
  &&
  \funcname{f}^2_0 \compose \funcname{f}^1_0 \maps  I \to U^3,
  \phi^3 \maps U^3 \toiso V^3
\\
& \bigr \}
\end{array}
\end{equation}

Let $u^1 \in I$, $u^2 \defeq \funcname{f}^1_0(u^1)$,
$u^3 \defeq \funcname{f}^2_0(u^2)$.
$(U^2, V^2, \phi^2) \in \seqname{A}^2$ be a chart at $u^2$. Since
$\funcseqname{f}^1$ is a morphism, there exists a subchart
$
  (U'^1, V'^1, \phi^1 \maps U'^1 \toiso V'^1)
  \in \seqname{A}^1
$
at $u^1$, a model neighborhood $U'^2 \subseteq U^2$ of $u^2$ and a chart \\
$
  (U'^2 ,\hat{V}'^2, \phi'^2 \maps U'^2 \toiso \hat{V}'^2)
  \in \seqname{A}^2
$
at $u^2$ such that $\funcname{f}^1_0[U'^1] \subseteq U'^2$ and diagram
\eqref{eq:AtlM1} in \fullcref{def:M-ATLmorph} on \cpageref{eq:AtlM1} is
commutative, as shown in \pagecref{fig:AtlM1}.

Let $I^{2,3} \defeq U'^2 \cap \funcname{f}^{2-1}_0[U^3]$. Since
$\funcseqname{f}^2$ is a near morphism, diagram \eqref{eq:morphcompN2} below
is locally nearly commutative at $u^2$, and thus there exist model
neighborhoods $U''^2 \subseteq I^{2,3}$, $U'^3 \subseteq U^3$,
$V''^2 \subseteq V'^3$,
$V'^3 \subseteq V^3$, $\hat{V}'^3 \subseteq C^3$ and an isomorphism
$\hat{\funcname{f}} \maps V'^3 \toiso \hat{V}'^3$\! such that diagram
\eqref{eq:morphcompN2} below is commutative, i.e.,
\crefrange{eq:xin2}{eq:fphigeqgphi2} below hold, as shown in
{
  \showlabelsinline
  \crefrange{fig:AtlN1}{fig:AtlN2}.
}

\begin{equation}
\label{eq:morphcompN2}
\begin{array} {lll}
D \defeq & \bigl \{
\\
  &&
  \funcname{i} \maps \set{u^2} \into I^{2,3},
  \phi'^2 \maps I^{2,3} \into V'^2,
  \funcname{f}^2_1 \maps  V'^2 \to \seqname{C}^3,
\\
  &&
  \funcname{f}^2_0 \maps I^{2,3} \to U'^3,
  \phi^3 \maps U'^3 \toiso V'^3
\\
& \bigr \}
\end{array}
\end{equation}

\begin{equation}
u^2 \in U''^2
\label{eq:xin2}
\end{equation}
\begin{equation}
\funcname{f}_0[U''^2] \subseteq U'^3
\end{equation}
\begin{equation}
\phi^1[U''^2] = V''^2
\end{equation}
\begin{equation}
\funcname{f}_1[V''^2] \subseteq \hat{V}'^3
\label{eq:fvp1in2}
\end{equation}
\begin{equation}
\phi^3[U'^3] = V'^3
\end{equation}
\begin{equation}
\hat{\funcname{f}} \compose \phi^3 \compose \funcname{f}^2_0 =
\funcname{f}^2_1 \compose \phi'^2
\label{eq:fphigeqgphi2}
\end{equation}

Define $U''^1 \defeq \funcname{f}^1_0[U''^2]$ and
$V''^1 \defeq \funcname{f}^1_1[V''^2]$. Diagram \eqref{eq:morphcompN3}
below is commutative, i.e.,
{
  \showlabelsinline
  \crefrange{eq:xin3}{eq:fphigeqgphi3}
}
below hold, as shown in \cref{fig:morphcompN1,fig:morphcompN2}.

\begin{equation}
\label{eq:morphcompN3}
\begin{array} {lll}
D \defeq & \bigl \{
\\
  &&
  \funcname{i} \maps \set{u^1} \into U''^1,
  \phi^1 \maps U''^1 \into V''^1,
  \funcname{f}^2_1 \compose \funcname{f}^1_1 \maps  V''^1 \to V'^3,
\\
  &&
  \funcname{f}^2_0 \compose \funcname{f}^1_0 \maps U''^1 \to U'^3
  \phi^3 \maps U'^3 \toiso V'^3
\\
& \bigr \}
\end{array}
\end{equation}

\begin{equation}
u^1 \in U''^1 \subseteq U'^ 1
\label{eq:xin3}
\end{equation}
\begin{equation}
\funcname{f}^2_0 \compose \funcname{f}^1_0[U''^1] \subseteq U'^3
\end{equation}
\begin{equation}
\phi^1[U'^3] = V'^3
\end{equation}
\begin{equation}
\funcname{f}^2_1 \compose \funcname{f}^1_1[V''^1] \subseteq \hat{V}'^3
\label{eq:fvp1in3}
\end{equation}
\begin{equation}
\phi^3[U'^3] = V'^3
\end{equation}
\begin{equation}
\begin{array}{lll}
  \hat{\funcname{f}} \compose \phi^3 \compose \funcname{f}^2_0
      \compose \funcname{f}^1_0
  &
    =
    &
      \funcname{f}^2_1 \compose \phi'^2 \compose \funcname{f}^1_0
\\
  &
    =
    &
      \funcname{f}^2_1 \compose \funcname{f}^1_1 \compose \phi^1
\end{array}
\label{eq:fphigeqgphi3}
\end{equation}

\begin{figure}[ht]
\[ \bfig
\node x(-900,0)[u^1]
\node up1(-900,-500)[U'^1]
\node vp1(2000,-500)[V'^1]
\node up2(-900,-1000)[U'^2]
\node uint(0,-1000)[{I^{2,3} = U'^2 \cap \funcname{f}^{-1}_0[U^3]}]
\node upp2(500,-1000)[U''^2]
\node vpp2(1500,-1000)[V''^2]
\node vhp2(2000,-1000)[\hat{V}'^2]
\node E3(-900,-1500)[E^3]
\node u3(0,-1500)[U^3]
\node v3(500,-1500)[V^3]
\node c3(2000,-1500)[C^3]
\arrow |r|/^{ (}->/[x`up1;\funcname{i}]
\arrow |a|/>->>/[up1`vp1;\phi^1]
\arrow |b|//[up1`vp1;\iso]
\arrow |r|[up1`up2;\funcname{f}^1_0]
\arrow |r|[vp1`vhp2;\funcname{f}^1_1]
\arrow |r|[up2`E3;\funcname{f}^2_0]
\arrow |a|/^{ (}->/[uint`up2;\funcname{i}]
\arrow |a|/>->/[uint`vhp2;\phi'^2]
\arrow |r|[uint`u3;\funcname{f}^2_0]
\arrow |r|[vhp2`c3;\funcname{f}^2_1]
\arrow |r|/^{ (}->/[u3`E3;\funcname{i}]
\arrow |a|/>->>/[u3`v3;\phi^3]
\arrow |b|//[u3`v3;\iso]
\efig \]
\caption{Partially completed composition of M-atlas near morphisms}
\label{fig:morphcompN1}
\end{figure}

\begin{figure}[ht]
\[ \bfig
\node x(0,-1000)[u^1]
\node up1(500,0)[U'^1]
\node upp1(500,-1000)[U''^1]
\node vpp1(500,-2000)[V''^1]
\node vp1(500,-2500)[V'^1]
\node up2(1500,0)[U'^2]
\node uint(1500,-500)[{I^{2,3} = U'^2 \cap \funcname{f}^{-1}_0[U^3]}]
\node upp2(1500,-1000)[U''^2]
\node vpp2(1500,-2000)[V''^2]
\node vhp2(1500,-2500)[\hat{V}'^2]
\node E3(2500,0)[E^3]
\node u3(2500,-500)[U^3]
\node up3(2500,-1000)[U'^3]
\node vp3(2500,-1500)[V'^3]
\node vhp3(2500,-2000)[\hat{V}'^3]
\node v3(2500,-1500)[V^3]
\node c3(2500,-2500)[C^3]
\arrow |a|/^{ (}->/[x`upp1;\funcname{i}]
\arrow |a|[up1`up2;\funcname{f}^1_0]
\arrow |r|/^{ (}->/[upp1`up1;\funcname{i}]
\arrow |a|[upp1`upp2;\funcname{f}^1_0]
\arrow |r|/>->>/[upp1`vpp1;\phi^1]
\arrow |l|//[upp1`vpp1;\iso]
\arrow |a|[vpp1`vpp2;\funcname{f}^1_1]
\arrow |r|/^{ (}->/[vpp1`vp1;\funcname{i}]
\arrow |a|[vp1`vhp2;\funcname{f}^1_1]
\arrow |a|[up2`E3;\funcname{f}^2_0]
\arrow |r|/^{ (}->/[uint`up2;\funcname{i}]
\arrow |a|[uint`u3;\funcname{f}^2_0]
\arrow |r|/^{ (}->/[upp2`uint;\funcname{i}]
\arrow |a|[upp2`up3;\funcname{f}^2_0]
\arrow |r|/>->>/[upp2`vpp2;\phi'^2]
\arrow |l|//[upp2`vpp2;\iso]
\arrow |a|[vpp2`vhp3;\funcname{f}^2_1]
\arrow |a|/^{ (}->/[vpp2`vhp2;\funcname{i}]
\arrow |a|[vhp2`c3;\funcname{f}^2_1]
\arrow |r|/^{ (}->/[u3`E3;\funcname{i}]
\arrow |r|/^{ (}->/[up3`u3;\funcname{i}]
\arrow |r|/>->>/[up3`vp3;\phi^3]
\arrow |l|//[up3`vp3;\iso]
\arrow |r|/>->>/[vp3`vhp3;\hat{\funcname{f}}]
\arrow |l|//[vp3`vhp3;\iso]
\arrow |r|/^{ (}->/[vhp3`c3;\funcname{i}]
\efig \]
\caption{Completed composition of M-atlas near morphisms}
\label{fig:morphcompN2}
\end{figure}

\end{proof}

\item
\label{M-morphcomp:morphcompmorph}
If each $\funcseqname{f}^i$ is an $\seqname{E}^1$-$\seqname{E}^{i+1}$
m-atlas morphism of $\seqname{A}^i$ to $\seqname{A}^{i+1}$ in the
coordinate spaces $\seqname{C}^i$, $\seqname{C}^{i+1}$ then
$\funcseqname{f}^2 \compose[()] \funcseqname{f}^1$ is an
$\seqname{E}^1$-$\seqname{E}^3$ m-atlas morphism of $\seqname{A}^1$ to
$\seqname{A}^3$ in the coordinate spaces $\seqname{C}^1$,
$\seqname{C}^3$.

\begin{proof}
$\funcname{f}^2_0 \compose \funcname{f}^1_0$ and
$\funcname{f}^2_1 \compose \funcname{f}^1_1$ are model functions by
\pagecref{lem:modcomp}.

Let $(U^i, V^i, \phi^i) \in \seqname{A}^i$, $i=1,3$, be a chart,
$
  I \defeq
  U^1 \cap (\funcname{f}^2_0 \compose \funcname{f}^1_0)^{-1}[U^3]
  \neq \emptyset
$.

Let $u^1 \in I$, $u^2 \defeq \funcname{f}^1_0(u^1)$,
$u^3 \defeq \funcname{f}^2_0(u^2)$ and
$(U^2, V^2, \phi^2) \in \seqname{A}^2$ be a chart at $u^2$.
Since $\funcseqname{f}^1$ is a morphism by hypothesis, there exists a
subchart \\*
\nobreaks
  {{
    $
      (U'^1, V'^1, \phi^1 \maps U'^1 \toiso V'^1)
      \in \seqname{A}^1
    $
  }}
at $u^1$, a model neighborhood $U'^2 \subseteq U^2$ of $u^2$ and a chart
$
  (U'^2 ,\hat{V}'^2, \phi'^2 \maps U'^2 \toiso \hat{V}'^2)
  \in \seqname{A}^2
$
at $u^2$ such that
\nobreaks
  {{
    $\funcname{f}^1_0[U'^1] \subseteq U'^2$
  }}
and diagram \eqref{eq:AtlM1} in \fullcref{def:M-ATLmorph} on \cpageref{eq:AtlM1} is
commutative, as shown in \pagecref{fig:AtlM1}.

\begin{figure}[ht]

\[ \bfig
\node x(0,0)[u^1]
\node up1(500,0)[U'^1]
\node upp1(500,-1000)[U''^1]
\node vpp1(500,-2000)[V''^1]
\node vp1(500,-2500)[V'^1]
\node up2(1500,0)[U'^2]
\node uint(1500,-500)[{I^{2,3} = U'^2 \cap \funcname{f}^{-1}_0[U^3]}]
\node upp2(1500,-1000)[U''^2]
\node vpp2(1500,-2000)[V''^2]
\node vhp2(1500,-2500)[\hat{V}'^2]
\node E3(2500,0)[E^3]
\node u3(2500,-500)[U^3]
\node up3(2500,-1000)[U'^3]
%\node vp3(2500,-1500)[V'^3]
\node vhp3(2500,-2000)[\hat{V}'^3]
\node v3(2500,-1500)[V^3]
\node c3(2500,-2500)[C^3]
\arrow |a|/^{ (}->/[x`up1;\funcname{i}]
\arrow |a|[up1`up2;\funcname{f}^1_0]
\arrow  |r|/>->>/[up1`vp1;\phi^1]
%\arrow |r|/^{ (}->/[upp1`up1;\funcname{i}]
%\arrow |a|[upp1`upp2;\funcname{f}^1_0]
%\arrow |r|/>->>/[upp1`vpp1;\phi^1]
%\arrow |l|//[upp1`vpp1;\iso]
%\arrow |a|[vpp1`vpp2;\funcname{f}^1_1]
%\arrow |r|/^{ (}->/[vpp1`vp1;\funcname{i}]
\arrow |a|[vp1`vhp2;\funcname{f}^1_1]
\arrow |a|[up2`E3;\funcname{f}^2_0]
\arrow |r|/^{ (}->/[uint`up2;\funcname{i}]
\arrow |a|[uint`u3;\funcname{f}^2_0]
\arrow |a|[uint`vhp2;\phi'^2]
%\arrow |r|/^{ (}->/[upp2`uint;\funcname{i}]
%\arrow |a|[upp2`up3;\funcname{f}^2_0]
%\arrow |r|/>->>/[upp2`vpp2;\phi'^2]
%\arrow |l|//[upp2`vpp2;\iso]
%\arrow |a|[vpp2`vhp3;\funcname{f}^2_1]
%\arrow |a|/^{ (}->/[vpp2`vhp2;\funcname{i}]
\arrow |a|[vhp2`c3;\funcname{f}^2_1]
\arrow |r|/^{ (}->/[u3`E3;\funcname{i}]
\arrow |r|/^{ (}->/[u3`v3;\phi'^3]
%\arrow |r|/^{ (}->/[up3`u3;\funcname{i}]
%\arrow |r|/>->>/[up3`vhp3;\phi'^3]
%\arrow |l|//[up3`vhp3;\iso]
%\arrow |r|/>->>/[vp3`vhp3;\hat{\funcname{f}}]
%\arrow |l|//[vp3`vhp3;\iso]
%\arrow |r|/^{ (}->/[vhp3`c3;\funcname{i}]
\efig \]
\caption{Uncompleted composition of m-atlas morphisms}
\label{fig:morphcompM1}
\end{figure}

Since $\funcseqname{f}^2$ is a morphism by hypothesis, there exists a
subchart \\
$
  (U''^2, V''^2, \phi^2 \maps U''^2 \toiso V''^2)
  \in \seqname{A}^2
$
at $u^2$, a model neighborhood $U'^3 \subseteq U^3$ of $u^3$ and a chart
$
  (U'^3 ,\hat{V}'^3, \phi'^3 \maps U'^3 \toiso \hat{V}'^3)
  \in \seqname{A}^3
$
at $u^2$ such that $\funcname{f}^2_0[U''^2] \subseteq U'^3$ and diagram
\eqref{eq:AtlM2} below is commutative, as shown in
\cref{fig:morphcompM1,fig:morphcompM2}.

\begin{equation}
\begin{array}{lll}
\label{eq:AtlM2}
D \defeq & \bigl \{
\\
  &&
  \funcname{i} \maps \set{u^2} \into U''^2,
  \phi^2 \maps U''^2 \toiso V'^2,
  \funcname{f}_1 \maps  V''^2 \to \hat{V}'^3,
\\
  &&
  \funcname{f}_0 \maps  U'^2 \to U'^3,
  \phi'^3 \maps U'^3 \toiso \hat{V}'^3
\\
& \bigr \}
\end{array}
\end{equation}

\begin{figure}[ht]
\[ \bfig
\node x(0,0)[u^1]
\node up1(500,0)[U'^1]
\node upp1(500,-1000)[U''^1]
\node vpp1(500,-2000)[V''^1]
\node vp1(500,-2500)[V'^1]
\node up2(1500,0)[U'^2]
\node uint(1500,-500)[{I^{2,3} = U'^2 \cap \funcname{f}^{-1}_0[U^3]}]
\node upp2(1500,-1000)[U''^2]
\node vpp2(1500,-2000)[V''^2]
\node vhp2(1500,-2500)[\hat{V}'^2]
\node E3(2500,0)[E^3]
\node u3(2500,-500)[U^3]
\node up3(2500,-1000)[U'^3]
%\node vp3(2500,-1500)[V'^3]
\node vhp3(2500,-2000)[\hat{V}'^3]
\node v3(2500,-1500)[V^3]
\node c3(2500,-2500)[C^3]
\arrow |a|/^{ (}->/[x`up1;\funcname{i}]
\arrow |a|[up1`up2;\funcname{f}^1_0]
\arrow |r|/>->>/[up1`vp1;\phi^1]
\arrow |l|//[up1`vp1;\iso]
% \arrow |r|/^{ (}->/[upp1`up1;\funcname{i}]
% \arrow |a|[upp1`upp2;\funcname{f}^1_0]
% \arrow |r|/>->>/[upp1`vpp1;\phi^1]
% \arrow |l|//[upp1`vpp1;\iso]
% \arrow |a|[vpp1`vpp2;\funcname{f}^1_1]
% \arrow |r|/^{ (}->/[vpp1`vp1;\funcname{i}]
\arrow |a|[vp1`vhp2;\funcname{f}^1_1]
\arrow |a|[up2`E3;\funcname{f}^2_0]
\arrow |r|/^{ (}->/[uint`up2;\funcname{i}]
\arrow |a|[uint`u3;\funcname{f}^2_0]
\arrow |r|/^{ (}->/[upp2`uint;\funcname{i}]
\arrow |a|[upp2`up3;\funcname{f}^2_0]
\arrow |r|/>->>/[upp2`vpp2;\phi'^2]
\arrow |l|//[upp2`vpp2;\iso]
\arrow |a|[vpp2`vhp3;\funcname{f}^2_1]
\arrow |a|/^{ (}->/[vpp2`vhp2;\funcname{i}]
\arrow |a|[vhp2`c3;\funcname{f}^2_1]
\arrow |r|/^{ (}->/[u3`E3;\funcname{i}]
\arrow |r|/^{ (}->/[up3`u3;\funcname{i}]
\arrow |r|/>->>/[up3`vhp3;\phi'^3]
\arrow |l|//[up3`vhp3;\iso]
%\arrow |r|/>->>/[vp3`vhp3;\hat{\funcname{f}}]
%\arrow |l|//[vp3`vhp3;\iso]
\arrow |r|/^{ (}->/[vhp3`c3;\funcname{i}]
\efig \]
\caption{Partially completed composition of m-atlas morphisms}
\label{fig:morphcompM2}
\end{figure}

\begin{figure}[ht]
\[ \bfig
\node x(0,-1000)[u^1]
\node up1(500,0)[U'^1]
\node upp1(500,-1000)[U''^1]
\node vpp1(500,-2000)[V''^1]
\node vp1(500,-2500)[V'^1]
\node up2(1500,0)[U'^2]
\node uint(1500,-500)[{I^{2,3} = U'^2 \cap \funcname{f}^{-1}_0[U^3]}]
\node upp2(1500,-1000)[U''^2]
\node vpp2(1500,-2000)[V''^2]
\node vhp2(1500,-2500)[\hat{V}'^2]
\node E3(2500,0)[E^3]
\node u3(2500,-500)[U^3]
\node up3(2500,-1000)[U'^3]
%\node vp3(2500,-1500)[V'^3]
\node vhp3(2500,-2000)[\hat{V}'^3]
\node v3(2500,-1500)[V^3]
\node c3(2500,-2500)[C^3]
\arrow |a|/^{ (}->/[x`upp1;\funcname{i}]
\arrow |a|[up1`up2;\funcname{f}^1_0]
\arrow |r|/^{ (}->/[upp1`up1;\funcname{i}]
\arrow |a|[upp1`upp2;\funcname{f}^1_0]
\arrow |r|/>->>/[upp1`vpp1;\phi^1]
\arrow |l|//[upp1`vpp1;\iso]
\arrow |a|[vpp1`vpp2;\funcname{f}^1_1]
\arrow |r|/^{ (}->/[vpp1`vp1;\funcname{i}]
\arrow |a|[vp1`vhp2;\funcname{f}^1_1]
\arrow |a|[up2`E3;\funcname{f}^2_0]
\arrow |r|/^{ (}->/[uint`up2;\funcname{i}]
\arrow |a|[uint`u3;\funcname{f}^2_0]
\arrow |r|/^{ (}->/[upp2`uint;\funcname{i}]
\arrow |a|[upp2`up3;\funcname{f}^2_0]
\arrow |r|/>->>/[upp2`vpp2;\phi'^2]
\arrow |l|//[upp2`vpp2;\iso]
\arrow |a|[vpp2`vhp3;\funcname{f}^2_1]
\arrow |a|/^{ (}->/[vpp2`vhp2;\funcname{i}]
\arrow |a|[vhp2`c3;\funcname{f}^2_1]
\arrow |r|/^{ (}->/[u3`E3;\funcname{i}]
\arrow |r|/^{ (}->/[up3`u3;\funcname{i}]
\arrow |r|/>->>/[up3`vhp3;\phi'^3]
\arrow |l|//[up3`vhp3;\iso]
%\arrow |r|/>->>/[vp3`vhp3;\hat{\funcname{f}}]
%\arrow |l|//[vp3`vhp3;\iso]
\arrow |r|/^{ (}->/[vhp3`c3;\funcname{i}]
\efig \]
\caption{Completed composition of M-atlas morphisms}
\label{fig:morphcompM3}
\end{figure}

Define $U''^1 \defeq \funcname{f}^{1-1}_0[U''^2]$,
$
  V''^1 \defeq
  \funcname{f}^{1-1}_1[V''^2] =
  \phi^1[U''^1]
$,
\Cref{fig:morphcompM3} is commutative.
\end{proof}

\item
\label{M-morphcomp:nearcompnear1}
If $\seqname{A}^2$ is semi-maximal then
$
  \funcseqname{f}^2 \compose[()] \funcseqname{f}^1 \defeq
  (
    \funcname{f}^2_0 \compose \funcname{f}^1_0,
    \funcname{f}^2_1 \compose \funcname{f}^1_1
  )
$
is an $\seqname{E}^1$-$\seqname{E}^3$ m-atlas near morphism of
$\seqname{A}^1$ to $\seqname{A}^3$ in the coordinate spaces
$\seqname{C}^1$, $\seqname{C}^3$.

\begin{proof}
Since $\seqname{A}^2$ is semi-maximal and $\funcseqname{f}^1$ is a
(near) morphism, $\funcseqname{f}^1$
is a morphism either by hypothesis
or by \cref{M-ATLmorph:semimaxnearismorph} of
{
  \showlabelsinline
  \fullcref{lem:M-ATLmorph} on
  \cpageref{M-ATLmorph:semimaxnearismorph}.
}

The result follows from \cref{M-morphcomp:morphcompmorph} above,
\end{proof}

\item
\label{M-morphcomp:nearcompnear2}
If $\seqname{A}^2$ and $\seqname{A}^3$ are semi-maximal then
$
  \funcseqname{f}^2 \compose[()] \funcseqname{f}^1 \defeq
  (
    \funcname{f}^2_0 \compose \funcname{f}^1_0,
    \funcname{f}^2_1 \compose \funcname{f}^1_1
  )
$
is an $\seqname{E}^1$-$\seqname{E}^3$ m-atlas morphism of
$\seqname{A}^1$ to $\seqname{A}^3$ in the coordinate spaces
$\seqname{C}^1$, $\seqname{C}^3$.

\begin{proof}
Since $\seqname{A}^2$ is semi-maximal and $\funcseqname{f}^1$ is a
(near) morphism, $\funcseqname{f}^1$
is a morphism either by hypothesis
or by \cref{M-ATLmorph:semimaxnearismorph} of
{
  \showlabelsinline
  \fullcref{lem:M-ATLmorph} on
  \cpageref{M-ATLmorph:semimaxnearismorph}.
}
Similarly, $\funcseqname{f}^2$ is a morphism.

The result follows from \cref{M-morphcomp:morphcompmorph} above,
\end{proof}

\item
\label{M-morphcomp:semicompsemi}
If each $\funcname{f}^i$, $i=1,2$, is a semi-strict
\nobreaks
  {
    {($\catname{E}^1$-$\catname{E}^2$-)},
    {$\seqname{E}^1$-$\seqname{E}^2$},
    {(-$\catname{C}^1$-$\catname{C}^2$)}
  }
m-atlas (near) morphism
then $\funcseqname{f}^2 \compose[()] \funcseqname{f}^1$ is a semi-strict
\nobreaks
  {
    {($\catname{E}^1$-$\catname{E}^2$-)},
    {$\seqname{E}^1$-$\seqname{E}^2$},
    {(-$\catname{C}^1$-$\catname{C}^2$)}
  }
m-atlas (near) morphism.

\begin{proof}
There are four overlapping cases.
\begin{description}
\item[$\seqname{E}^i$ without category:]
If each $\seqname{E}^i \submod \seqname{E}^{i+1}$ and each
$\funcseqname{f}^i_0$ is a local m-morphism of $\seqname{E}^i$ to
$\seqname{E}^{i+1}$, then by \pagecref{lem:modMmorphcomp}
$\funcseqname{f}^2_0 \compose \funcseqname{f}^1_0$ is a local
m-morphism of $\seqname{E}^1$ to $\seqname{E}^3$.

\item[$\seqname{C}^i$ without category:]
If each $\seqname{C}^i \submod \seqname{C}^{i+1}$ and each
$\funcseqname{f}^i_1$ is a local m-morphism of $\seqname{C}^i$ to
$\seqname{C}^{i+1}$, then by \cref{lem:modMmorphcomp}
$\funcseqname{f}^2_1 \compose \funcseqname{f}^1_1$ is a local
morphism of $\seqname{C}^1$ to $\seqname{C}^3$.
Let $(U^i, V^i, \phi^i \maps U^i \toiso V^i) \in \seqname{A}^i$, $i=1,3$,
with
$
  I \defeq
  U^1 \cap (\funcseqname{f}^2_0 \compose \funcseqname{f}^1_0)^{-1}[U^3]
  \neq \emptyset
$.
For every $u^1 \in I$ let $u^2 \defeq \funcseqname{f}^1_0(u^1)$,
$u^3 \defeq \funcseqname{f}^2_0(u^2)$ and
$(U^2, V^2, \phi^2 \maps U^2 \toiso V^2) \in \seqname{A}^2$ be a chart
at $u^2$,
$
  I^1 \defeq
  U^1 \cap \funcseqname{f}^{1-1}_0[U^2]
  \neq \emptyset
$
and
$
  I^2 \defeq
  U^2 \cap \funcseqname{f}^{2-1}_0 [U^3]
  \neq \emptyset
$.
If
$
  \phi^2 \compose \funcname{f}^1_0 \compose \phi^{1-1} \maps
  \phi^1[I^1] \to V^2
$
is a local m-morphism of $\Mod(\phi^1[I^1],\seqname{C}^1)$
to $\Mod(V^2,\seqname{C}^2)$ and
$
  \phi^3 \compose \funcname{f}^2_0 \compose \phi^{2-1} \maps
  \phi^2[I^2] \to V^3
$
is a local m-morphism of $\Mod(\phi^1[I^2],\seqname{C}^1)$
to $\Mod(V^2,\seqname{C}^2)$ then by \cref{def:modlocalMmorph} there
exist model neighborhoods $V'^i \subseteq V^i$, $i \in [1,3]$, such that
$\funcname{f}^i[V'^i] \subseteq V'^{i+1}$ and \\*
$
  \phi^{i+1} \compose \funcname{f}^i_0 \compose \phi^{i-1} \maps
  V'^i \to V'^{i+1}
$
is an m-morphism of $\seqname{C}^{i+1}$, $i=1,2$.
Then
\begin{equation*}
\begin{array}{lll}
  \phi^3 \compose \funcname{f}^2_0 \compose \funcname{f}^1_0 \compose
  \phi^{1-1} \maps V'^1 \to V'^3
  &
  =
  &
  \phi^3 \compose \funcname{f}^2_0 \compose \phi^{2-1} \maps
  V'^2 \to V'^3 \compose
\\
  &&
  \phi^2 \compose \funcname{f}^1_0 \compose \phi^{1-1} \maps
  V'^1 \to V'^2
\end{array}
\end{equation*}
is an m-morphism of $\seqname{C}^3$. Since $u^1 \in I$ is arbitray,
$
  \phi^3 \compose
  \funcname{f}^2_0 \compose \funcname{f}^1_0 \compose
  \phi^{1-1} \maps \phi^1[I] \to V^3
$
is a local m-morphism of $\Mod(\phi^1[I],\seqname{C}^3)$
to $\Mod(V^3,\seqname{C}^3)$.

\item[$\seqname{E}^i$ with category:]
If each $\catname{E}^i \subcat[full-] \catname{E}^{i+1}$ and each
$\funcseqname{f}^i_0$ is a local $\catname{E}^i$-$\catname{E}^{i+1}$
m-morphism of $\seqname{E}^i$ to $\seqname{E}^{i+1}$, then by
\pagecref{lem:modMmorphcomp}
$\funcseqname{f}^2_0 \compose \funcseqname{f}^1_0$ is a local
$\catname{E}^1$-$\catname{E}^3$ m-morphism of $\seqname{E}^1$ to
$\seqname{E}^3$.

\item[$\seqname{C}^i$ with category:]
If each $\catname{C}^i \subcat[full-] \catname{C}^{i+1}$ and each
$\funcseqname{f}^i_1$ is a local $\catname{C}^i$-$\catname{C}^{i+1}$
m-morphism of $\seqname{C}^i$ to $\seqname{C}^{i+1}$, then by
\pagecref{lem:modMmorphcomp}
$\funcseqname{f}^2_1 \compose \funcseqname{f}^1_1$ is a local
$\catname{C}^1$-$\catname{C}^3$ m-morphism of $\seqname{C}^1$ to
$\seqname{C}^3$.
\end{description}
\end{proof}

\item
\label{M-morphcomp:strictcompstrict}
If each $\funcname{f}^i$, $i=1,2$, is strict then
$\funcseqname{f}^2 \compose[()] \funcseqname{f}^1$ is strict.

\begin{proof}
There are four overlapping cases.
\begin{description}
\item[$\seqname{E}^i$ without category:]
If each $\seqname{E}^i \submod \seqname{E}^{i+1}$ and each
$\funcseqname{f}^i_0$ is an m-morphism of $\seqname{E}^i$ to
$\seqname{E}^{i+1}$, then by \pagecref{lem:modMmorphcomp}
$\funcseqname{f}^2_0 \compose \funcseqname{f}^1_0$ is an
m-morphism of $\seqname{E}^1$ to $\seqname{E}^3$.

\item[$\seqname{C}^i$ without category:]
If each $\seqname{C}^i \submod \seqname{C}^{i+1}$ and each
$\funcseqname{f}^i_1$ is an m-morphism of $\seqname{C}^i$ to
$\seqname{C}^{i+1}$, then by \pagecref{lem:modMmorphcomp}
$\funcseqname{f}^2_1 \compose \funcseqname{f}^1_1$ is an
m-morphism of $\seqname{C}^1$ to $\seqname{C}^3$.

\item[$\seqname{E}^i$ with category:]
If each $\catname{E}^i \subcat[full-] \catname{E}^{i+1}$ and each
$\funcseqname{f}^i_0$ is a morphism of $\catname{E}^{i+1}$, then
$\funcseqname{f}^2_0 \compose \funcseqname{f}^1_0$ is a morphism of
$\catname{E}^3$.

\item[$\seqname{C}^i$ with category:]
If each $\catname{C}^i \subcat[full-] \catname{C}^{i+1}$ and each
$\funcseqname{f}^i_1$ is a morphism of $\catname{C}^i$, then
$\funcseqname{f}^2_1 \compose \funcseqname{f}^1_1$ is a
morphism of $\catname{C}^3$.
\end{description}
\end{proof}

\item
\label{M-morphcomp:constcompconst}
If each $\funcseqname{f}^i$ is constrained then
$\funcseqname{f}^2 \compose[()] \funcseqname{f}^1$ is constrained.

\begin{proof}
This is a corollary of \pagecref{lem:modcomp}.
\end{proof}

\item
\label{M-morphcomp:compequivcomp}
$\funcseqname{f}^2 \compose[()] \funcname{f}^1$ is equivalent to
$\funcseqname{g}^2 \compose[()] \funcname{g}^1$.

\begin{proof}
$\funcname{f}^1_0 = \funcname{g}^1_0$ and
$\funcname{f}^2_0 = \funcname{g}^2_0$, hence
$
  \funcname{f}^2_0 \compose \funcname{f}^1_0 =
  \funcname{g}^2_0 \compose \funcname{g}^1_0
$.
\end{proof}
\end{enumerate}
\end{lemma}

\begin{corollary}[Composition of m-atlas (near) morphisms]
\label{cor:M-ATLmorphcomp}
Let $\catname{E}^1$ be a model category,
$\seqname{E}^1 \objin \catname{E}^1$,
$\catname{C}^i$, $i=1,2,3$, be a model category,
$\seqname{C}^i \objin \catname{C}^i$, $\seqname{A}^i$ be an m-atlas of
$\seqname{E}^1$ in the coordinate space $\seqname{C}^i$ and
$
  \funcseqname{f}^i \defeq
    (
      \funcname{f}^i_0 \maps \seqname{E}^1 \to \seqname{E}^1,
      \funcname{f}^i_1 \maps \seqname{C}^i \to \seqname{C}^{i+1}
    )
$, $i=1,2$,
be a (semi-strict, strict)
\nobreaks
  {
    {($\catname{E}^1$-)},
    {$\seqname{E}^1$},
    {(-$\catname{C}^i$-$\catname{C}^{i+1}$)}
  }
m-atlas (near) morphism of
$\seqname{A}^i$ to $\seqname{A}^{i+1}$ in the coordinate spaces
$\seqname{C}^i$, $\seqname{C}^{i+1}$.

Then

\begin{enumerate}
\item
If $\seqname{A}^2$ is semi-maximal then
$
  \funcseqname{f}^2 \compose[()] \funcseqname{f}^1 =
    (
      \funcname{f}^2_0 \compose \funcname{f}^1_0,
      \funcname{f}^2_1 \compose \funcname{f}^1_1
    )
$
is a (semi-strict, strict)
\nobreaks
  {
    {($\catname{E}^1$-)},
    {$\seqname{E}^1$},
    {(-$\catname{C}^1$-$\catname{C}^3$)},
  }
m-atlas near morphism of $\seqname{A}^1$ to $\seqname{A}^3$ in the
coordinate spaces $\seqname{C}^1$, $\seqname{C}^3$.

\begin{proof}
Since each $\funcseqname{f}^i$ is a (semi-strict, strict)
\nobreaks
  {
    {($\catname{E}^1$-$\catname{E}^1$-)},
    {$\seqname{E}^1$-$\seqname{E}^1$},
    {(-$\catname{C}^i$-$\catname{C}^{i+1}$)}
  }
m-atlas (near) morphism of
$\seqname{A}^i$ to $\seqname{A}^{i+1}$ in the coordinate spaces
$\seqname{C}^i$, $\seqname{C}^{i+1}$ by
\pagecref{def:M-ATLmorphabbrev},
and $\seqname{A}^2$ is semi-maximal,
$\funcseqname{f}^2 \compose[()] \funcseqname{f}^1$ is a (semi-strict,
strict)
\nobreaks
  {
    {($\catname{E}^1$-$\catname{E}^1$-)},
    {$\seqname{E}^1$-$\seqname{E}^1$},
    {(-$\catname{C}^i$-$\catname{C}^3$)}
  }
m-atlas near morphism of
$\seqname{A}^1$ to $\seqname{A}^3$ in the coordinate spaces
$\seqname{C}^1$, $\seqname{C}^3$ by \cref{M-morphcomp:nearcompnear1} of
{
  \showlabelsinline
  \pagecref{lem:M-ATLmorphcomp}.
}
Apply \cref{def:M-ATLmorphabbrev}.
\end{proof}

\item
If each $\funcseqname{f}^i$ is a (semi-strict, strict)
\nobreaks
  {
    {($\catname{E}^1$-)$\seqname{E}^1$},
    {$\seqname{E}^1$-$\seqname{E}^1$},
    {(-$\catname{C}^i$-$\catname{C}^{i+1}$)}
  }
m-atlas morphism of $\seqname{A}^i$ to $\seqname{A}^{i+1}$ in the
coordinate spaces $\seqname{C}^i$, $\seqname{C}^{i+1}$ then
$\funcseqname{f}^2 \compose[()] \funcseqname{f}^1$ is a (semi-strict,
strict)
\nobreaks
  {
    {($\catname{E}^1$-$\seqname{E}^1$-)},
    {$\seqname{E}^1$-$\seqname{E}^1$},
    {(-$\catname{C}^1$-$\catname{C}^3$)}
  }
m-atlas morphism of $\seqname{A}^1$ to $\seqname{A}^3$ in the coordinate
spaces $\seqname{C}^1$, $\seqname{C}^3$.

\begin{proof}
Apply \cref{def:M-ATLmorphabbrev} as above.
\end{proof}
\end{enumerate}

Let $\catname{E}^i$ $i=1,2,3$, be a model category,
$\seqname{E}^i \objin \catname{E}^i$,
$\catname{C}^1$ be a model category,
$\seqname{C}^1 \objin \catname{C}^1$, $\seqname{A}^i$ be an m-atlas of
$\seqname{E}^i$ in the coordinate space $\seqname{C}^1$ and \\*
\nobreaks
  {{
    $
      \funcseqname{f}^i \defeq
        (
          \funcname{f}^i_0 \maps \seqname{E}^i \to \seqname{E}^{i+1}
          \funcname{f}^i_1 \maps \seqname{C}^1 \to \seqname{C}^1
        )
    $,
  }}
$i=1,2$,
be a (semi-strict, strict)
\nobreaks
  {
    {($\catname{E}^i$-$\catname{E}^{i+1}$-)},
    {$\seqname{E}^i$-$\seqname{E}^{i+1}$},
    {(-$\catname{C}^1$)}
  }
m-atlas (near) morphism of
$\seqname{A}^i$ to $\seqname{A}^{i+1}$ in the coordinate space
$\seqname{C}^1$.

\begin{enumerate}[resume]
\item
If $\seqname{A}^2$ is semi-maximal then
$
  \funcseqname{f}^2 \compose[()] \funcseqname{f}^1 =
    (
      \funcname{f}^2_0 \compose \funcname{f}^1_0,
      \funcname{f}^2_1 \compose \funcname{f}^1_1
    )
$
is a (semi-strict, strict)
\nobreaks
  {
    {($\catname{E}^1$-$\catname{E}^3$-)},
    {$\seqname{E}^1$-$\seqname{E}^3$},
    {(-$\catname{C}^1$)}
  }
m-atlas near morphism of
$\seqname{A}^1$ to $\seqname{A}^3$ in the coordinate space
$\seqname{C}^1$.

\begin{proof}
Apply \cref{def:M-ATLmorphabbrev} as above.
\end{proof}

\item
If each $\funcseqname{f}^i$ is a (semi-strict, strict)
\nobreaks
  {
    {($\catname{E}^i$-$\catname{E}^{i+1}$-)},
    {$\seqname{E}^i$-$\seqname{E}^{i+1}$},
    {(-$\catname{C}^1$)}
  }
m-atlas morphism of $\seqname{A}^i$ to $\seqname{A}^{i+1}$ in the
coordinate spaces $\seqname{C}^i$, $\seqname{C}^{i+1}$ then
$\funcseqname{f}^2 \compose[()] \funcseqname{f}^1$ is a (semi-strict,
strict) \\*
\nobreaks
  {
    {($\catname{E}^1$-$\catname{E}^3$-)},
    {$\seqname{E}^1$-$\seqname{E}^3$},
    {(-$\catname{C}^1$-$\catname{C}^1$)}
  }
m-atlas morphism of $\seqname{A}^1$ to $\seqname{A}^3$ in the coordinate
spaces $\seqname{C}^1$, $\seqname{C}^1$.

\begin{proof}
Apply \cref{def:M-ATLmorphabbrev} as above.
\end{proof}
\end{enumerate}

Let $\catname{E}^1$, $\catname{C}^1$ be model categories,
$\seqname{E}^1 \objin \catname{E}^1$,
$\seqname{C}^1 \objin \catname{C}^1$, $\seqname{A}^i$, $i=1,2,3$, be an
m-atlas of $\seqname{E}^1$ in the coordinate space $\seqname{C}^1$ and
$
  \funcseqname{f}^i \defeq
    (
      \funcname{f}^i_0 \maps \seqname{E}^1 \to \seqname{E}^1,
      \funcname{f}^i_1 \maps \seqname{C}^1 \to \seqname{C}^1
    )
$, $i=1,2$,
be a (semi-strict, strict) $\seqname{E}^1$ m-atlas (near) morphism of
$\seqname{A}^i$ to $\seqname{A}^{i+1}$ in the coordinate space
$\seqname{C}^1$.

\begin{enumerate}[resume]
\item
If $\seqname{A}^2$ is semi-maximal then
$
  \funcseqname{f}^2 \compose[()] \funcseqname{f}^1 =
    (
      \funcname{f}^2_0 \compose \funcname{f}^1_0,
      \funcname{f}^2_1 \compose \funcname{f}^1_1
    )
$
is a (semi-strict, strict) $\seqname{E}^1$ m-atlas m-atlas near morphism
of $\seqname{A}^1$ to $\seqname{A}^3$ in the coordinate space
$\seqname{C}^1$.

\begin{proof}
Apply \cref{def:M-ATLmorphabbrev} as above.
\end{proof}

\item
If each $\funcseqname{f}^i$ is a (semi-strict, strict) $\seqname{E}^1$
m-atlas morphism of $\seqname{A}^i$ to $\seqname{A}^{i+1}$ in the
coordinate space $\seqname{C}^1$ then
$\funcseqname{f}^2 \compose[()] \funcseqname{f}^1$ is a (semi-strict,
strict) $\seqname{E}^1$ m-atlas morphism of $\seqname{A}^1$ to
$\seqname{A}^3$ in the coordinate space $\seqname{C}^1$.

\begin{proof}
Apply \cref{def:M-ATLmorphabbrev} as above.
\end{proof}
\end{enumerate}

Let $\catname{C}^i$, $i=1,2,3$, be a model category,
$\seqname{C}^i \objin \catname{C}^i$, $E^i$ be a topological space,
$\seqname{A}^i$ be an m-atlas of $E^i$ in the coordinate space
$\seqname{C}^i$,
$\seqname{A}^2$ be semi-maximal and
$
  \funcseqname{f}^i \defeq \\*
    (
      \funcname{f}^i_0 \maps E^i \to E^{i+1},
      \funcname{f}^i_1 \maps \seqname{C}^i \to \seqname{C}^{i+1}
    )
$, $i=1,2$,
be a (semi-strict, strict) $E^i$-$E^{i+1}$ m-atlas near morphism of
$\seqname{A}^i$ to $\seqname{A}^{i+1}$ in the coordinate spaces
$\seqname{C}^i$, $\seqname{C}^{i+1}$. Then

\begin{enumerate}[resume]
\item
$
  \funcseqname{f}^2 \compose[()] \funcseqname{f}^1 =
    (
      \funcname{f}^2_0 \compose \funcname{f}^1_0,
      \funcname{f}^2_1 \compose \funcname{f}^1_1
    )
$
is a (semi-strict, strict) $E^1$-$E^3$ m-atlas near morphisms of
$\seqname{A}^1$ to $\seqname{A}^3$ in the coordinate spaces
$\seqname{C}^1$, $\seqname{C}^3$.

\begin{proof}
Apply \cref{def:M-ATLmorphabbrev} as above.
\end{proof}
\end{enumerate}

Let $\catname{E}^i$, $\catname{C}^i$, $i \in [1,3]$, be model categories,
$\seqname{E}^i \objin \catname{E}^i$,
$\seqname{C}^i \objin \catname{C}^i$ $\seqname{A}^i$ an m-atlas
of $\seqname{E}^i$ in the coordinate space $\seqname{C}^i$,
$\seqname{A}^2$ maximal and
$
  \funcseqname{f}^i \defeq \\
  (
    \funcname{f}^i_0 \maps \seqname{E}^i \to \seqname{E}^{i+1},
    \funcname{f}^i_1 \maps \seqname{C}^i \to \seqname{C}^{i+1}
  )
$
a (semi-strict, strict) $\seqname{E}^i$-$\seqname{E}^{i+1}$ m-atlas near
morphism of $\seqname{A}^i$ to $\seqname{A}^{i+1}$ in the coordinate
spaces $\seqname{C}^i$, $\seqname{C}^{i+1}$. Then

\begin{enumerate}[resume]
\item
$
  \funcseqname{f}^2 \compose[()] \funcseqname{f}^1 =
  (
    \funcname{f}^2_0 \compose \funcname{f}^1_0,
    \funcname{f}^2_1 \compose \funcname{f}^1_1
  )
$
is a (semi-strict, strict) $\seqname{E}^1$-$\seqname{E}^3$ m-atlas near
morphism of $\seqname{A}^1$ to $\seqname{A}^3$ in the coordinate spaces
$\seqname{C}^1$, $\seqname{C}^3$.

\begin{proof}
Apply \cref{def:M-ATLmorphabbrev} as above.
\end{proof}
\end{enumerate}

Let $\catname{C}^i$, $i=1,2,3$, be a model category,
$\seqname{C}^i \objin \catname{C}^i$, $E^i$ be a topological space,
$\seqname{A}^i$ be an m-atlase of $E^i$ in the coordinate space
$\seqname{C}^i$,
$\seqname{A}^2$ maximal and
$
  \funcseqname{f}^i \defeq \\*
  (
    \funcname{f}^i_0 \maps E^i \to E^{i+1},
    \funcname{f}^i_1 \maps \seqname{C}^i \to \seqname{C}^{i+1}
  )
$
be a (semi-strict, strict) $E^i$-$E^{i+1}$ m-atlas near morphism of
$\seqname{A}^i$ to $\seqname{A}^{i+1}$ in the coordinate spaces
$\seqname{C}^i$, $\seqname{C}^{i+1}$. Then

\begin{enumerate}[resume]
\item
$
  \funcseqname{f}^2 \compose[()] \funcseqname{f}^1 =
  (
    \funcname{f}^2_0 \compose \funcname{f}^1_0,
    \funcname{f}^2_1 \compose \funcname{f}^1_1
  )
$
is a (semi-strict, strict) $E^1$-$E^3$ m-atlas near morphism of
$\seqname{A}^1$ to $\seqname{A}^3$ in the coordinate spaces
$\seqname{C}^1$, $\seqname{C}^3$.

\begin{proof}
Apply \cref{def:M-ATLmorphabbrev} as above.
\end{proof}
\end{enumerate}

Let $\catname{C}^i$, $i=1,2,3$, be a model category,
$\seqname{C}^i \objin \catname{C}^i$, $E^i$ be a topological space,
$\seqname{A}^i$ be an m-atlas of $E^i$ in the coordinate space
$\seqname{C}^i$ and
$
  \funcseqname{f}^i \defeq
    (
      \funcname{f}^i_0 \maps E^i \to E^{i+1},
      \funcname{f}^i_1 \maps \seqname{C}^i \to \seqname{C}^{i+1}
    )
$, $i=1,2$,
be a (semi-strict, strict) $E^i$-$E^{i+1}$ m-atlas morphism of
$\seqname{A}^i$ to $\seqname{A}^{i+1}$ in the coordinate spaces
$\seqname{C}^i$, $\seqname{C}^{i+1}$. Then

\begin{enumerate}[resume]
\item
$
  \funcseqname{f}^2 \compose[()] \funcseqname{f}^1 =
    (
      \funcname{f}^2_0 \compose \funcname{f}^1_0,
      \funcname{f}^2_1 \compose \funcname{f}^1_1
    )
$
is a (semi-strict, strict) $E^1$-$E^3$ m-atlas morphism of
$\seqname{A}^1$ to $\seqname{A}^3$ in the coordinate spaces
$\seqname{C}^1$, $\seqname{C}^3$.

\begin{proof}
Apply \cref{def:M-ATLmorphabbrev} as above.
\end{proof}
\end{enumerate}
\end{corollary}

\begin{lemma}[M-atlas identity]
\label{lem:M-ATLid}
Let
\begin{enumerate}
\item
$\catname{E}^i$, $\catname{C}^i$, $i=1,2$, be model categories
\item
$\catname{E}^1 \subcat[full-] \catname{E}^2$
\item
$\catname{C}^1 \subcat[full-] \catname{C}^2$
\item
$\seqname{E}^i \objin \catname{E}^i$
\item
$\seqname{E}^1 \submod \seqname{E}^2$
\item
$\seqname{C}^i \objin \catname{C}^i$
\item
$\seqname{C}^1 \submod \seqname{C}^2$
\item
$\seqname{A}^i$ be an m-atlas of $\seqname{E}^i$ in the coordinate space
$\seqname{C}^i$
\item
$\seqname{A}^1 \subseteq \seqname{A}^2$
\end{enumerate}

The inclusion morphism
$
  \ID%
    [
      {(\seqname{E}^1, \seqname{C}^1)},
      {(\seqname{E}^2, \seqname{C}^2)}
    ]
$
is a semi-strict
\nobreaks
  {
    {($\catname{E}^1$-$\catname{E}^2$-)},
    {$\seqname{E}^1$-$\seqname{E}^2$-},
    {(-$\catname{C}^1$-$\catname{C}^2$)}
  }
m-atlas near morphism of
$\seqname{A}^1$ to $\seqname{A}^2$ in the coordinate spaces
$\seqname{C}^1$, $\seqname{C}^2$.

\begin{proof}
\begin{description}

\item%
  [
    \ensuremath
      {
        \ID%
          [
            {(\seqname{E}^1, \seqname{C}^1)},
            {(\seqname{E}^2, \seqname{C}^2)}
          ]
      }
    is an m-atlas near morphism:
  ]

Let $(U^i,V^i,\phi^i) \in \seqname{A}^i$, $i=1,2$, be charts such that
$I \defeq U^1 \cap U^2 \neq \emptyset$. To show that diagram $D$ below
is M-locally nearly commutative in $\seqname{C}^2$,
let $u^1 \in I$. Since $\seqname{E}^1 \submod \seqname{E}^2$ by
hypothethis, $U^1$ is a model neighborhood of $\seqname{E}^2$ and thus
$I$ is a model neighborhood of $\seqname{E}^2$ and the identity map
$\Id[I,U^2]$ is a morphism of $\seqname{E}^2$. Then define
$U'^1 \defeq I$, $U'^2 \defeq I$, $V'^1 \defeq \phi^1[I]$,
$V'^2 \defeq \phi^2[I]$, $\hat{V}'^2 \defeq V'^1$ and
$
  \hat{\funcname{f}} \defeq
  \phi^1 \compose \phi^{2-1} \maps V'^2 \toiso V'^1
$.
\Crefrange{eq:xin}{eq:fphigeqgphi} in \fullcref{def:M-ATLmorphnear} hold.

\begin{equation*}
\begin{array} {lll}
D \defeq & \bigl \{
\\
  &&
  \phi^1 \maps I \defeq U^1 \cap U^2 \into V^1,
  \Id[{V^1},{\seqname{C}^2}] \maps  V^1 \to \seqname{C}^2,
\\
  &&
  \Id[{U^1},{U^2}] \maps  U^1 \to U^2,
  \phi^2 \maps U^2 \toiso V^2
\\
& \bigr \}
\end{array}
\end{equation*}

\item%
  [
    \ensuremath
      {
        \ID%
          [
            {(\seqname{E}^1, \seqname{C}^1)},
            {(\seqname{E}^2, \seqname{C}^2)}
          ]
      }
    is semi-strict:
  ]

Let $u^1 \in \seqname{E}^1$. Since the model neighborhoods of
$\seqname{E}^1$ form an open cover, there is a model neighborhood $U$ of
$\seqname{E}^1$ at $u^1$. $\Id[{U}]$ is a morphism of $\seqname{E}^1$
and of $\catname{E}^1$; since $\seqname{E}^1 \submod \seqname{E}^2$ and
$\catname{E}^1 \subcat[full-] \catname{E}^2$, $\Id[{U}]$ is also a
morphism of $\seqname{E}^2$ and of $\catname{E}^2$. Thus
$
\Id%
  [
    {\seqname{E}^1},
    {\seqname{E}^2}
  ]
$
is a local m-morphism of $\seqname{E}^1$ to $\seqname{E}^2$ and a local
$\catname{E}^1$-$\catname{E}^2$ m-morphism $\seqname{E}^1$ to $\seqname{E}^2$.

Let $v^1 \in \seqname{C}^1$. Since the model neighborhoods of
$\seqname{C}^1$ form an open cover, there is a model neighborhood $V$ of
$\seqname{C}^1$ at $v^1$. $\Id[{V}]$ is a morphism of $\seqname{C}^1$
and of $\catname{C}^1$; since $\seqname{C}^1 \submod \seqname{C}^2$ and
$\catname{C}^1 \subcat[full-] \catname{C}^2$, $\Id[{V}]$ is also a
morphism of $\seqname{C}^2$ and of $\catname{C}^2$. Thus
$
  \Id%
    [
      {\seqname{C}^1},
      {\seqname{C}^2}
    ]
$
is a local m-morphism of $\seqname{C}^1$ to $\seqname{C}^2$ and a local
$\catname{C}^1$-$\catname{C}^2$ m-morphism $\seqname{C}^1$ to $\seqname{C}^2$.
\end{description}
\end{proof}

If $\seqname{A}^1$ and $\seqname{A}^2$ are semi-maximal then
$
  \ID%
    [
      {(\seqname{E}^1, \seqname{C}^1)},
      {(\seqname{E}^2, \seqname{C}^2)}
    ]
$
is an m-atlas morphism.

\begin{proof}
Since
$
  \ID%
    [
      {(\seqname{E}^1, \seqname{C}^1)},
      {(\seqname{E}^2, \seqname{C}^2)}
    ]
$
is an m-atlas near morphism by the above and the atlasses are
semi-maximal by hypothesis, it is an m-atlas morphism by
\pagecref{lem:M-ATLmorph}.
\end{proof}

\begin{proof}

If $\seqname{E}^1 \objin \catname{E}^2$ and
$\seqname{C}^1 \objin \catname{C}^2$ then
$
\Id%
    [
      {(\seqname{A}^1, \seqname{E}^1, \seqname{C}^1)},
      {(\seqname{A}^2, \seqname{E}^2, \seqname{C}^2)}
    ]
$
is strict:

$\Id[{\seqname{E}^1}] \arin \catname{E}^1
    \subcat[full-] \catname{E}^2
$ and
$
  \Id[{\seqname{C}^1}] \arin \catname{C}^1
    \subcat[full-] \catname{C}^2
$.
\end{proof}

If each $\Top{\seqname{E}^i} \objin \seqname{E}^i$
($\seqname{E}^i \objin \catname{E}^i$)
and each $\Top{\seqname{C}^i} \objin \seqname{C}^i$
($\seqname{C}^i \objin \catname{C}^i$) then
$
  \ID%
    [
      {(\seqname{E}^1, \seqname{C}^1)},
      {(\seqname{E}^2, \seqname{C}^2)}
    ]
$
is a strict
\nobreaks
  {
    {($\catname{E}^1$-$\catname{E}^2$-)},
    {$\seqname{E}^1$-$\seqname{E}^2$-},
    {(-$\catname{C}^1$-$\catname{C}^2$)}
  }
m-atlas morphism of
$\seqname{A}^1$ to $\seqname{A}^2$ in the coordinate spaces
$\seqname{C}^1$, $\seqname{C}^2$.

\begin{proof}
There are four overlapping cases
\begin{description}
\item%
  [
    Each
    \ensuremath
      {
        \Top{\seqname{E}^i} \objin \seqname{E}^i
      }:
  ]
$\Id[{\seqname{E}^1},{\seqname{E}^2}]$ is a morphism of
$\seqname{E}^2$ by \cref{mod:inclusion} of
{
  \showlabelsinline
  \pagecref{def:model}
},
hence an m-morphism of $\seqname{E}^1$ to $\seqname{E}^2$.

\item%
  [
    Each
    \ensuremath
      {
        \seqname{E}^i \objin \catname{E}^2
      }:
  ]
$\Id[{\seqname{E}^1},{\seqname{E}^2}]$ is a morphism of
$\catname{E}^2$ by \cref{modcat:morph} of
{
  \showlabelsinline
  \pagecref{def:modcat}
},
hence an $\catname{E}^1$-$\catname{E}^2$ m-morphism of $\seqname{E}^1$
to $\seqname{E}^2$.

\item%
  [
    Each
    \ensuremath
      {
        \Top{\seqname{C}^i} \objin \seqname{C}^i
      }:
  ]
$\Id[{\seqname{C}^1},{\seqname{C}^2}]$ is a morphism of
$\seqname{C}^2$ by \cref{mod:inclusion} of
{
  \showlabelsinline
  \cref{def:model}
},
hence an m-morphism of $\seqname{C}^1$ to $\seqname{C}^2$.

\item%
  [
    Each
    \ensuremath
      {
        \seqname{C}^i \objin \catname{C}^2
      }:
  ]
$\Id[{\seqname{C}^1},{\seqname{C}^2}]$ is a morphism of
$\catname{C}^2$ by \cref{modcat:morph} of
{
  \showlabelsinline
  \cref{def:modcat}
},
hence a $\catname{C}^1$-$\catname{C^2}$ m-morphism of $\seqname{C}^1$ to
$\seqname{C}^2$.
\end{description}
\end{proof}
\end{lemma}

\subsection{Categories of m-atlases and functors}
\label{sub:M-cat}
This subsection defines categories of m-atlases and functors among them,
and proves some basic results.

\subsubsection{M-atlas categories}
In the following, \textbf{...} will refer to one of
\begin{equation*}
  \seqname{A}^1, \seqname{E}^1, \seqname{C}^1,
  \seqname{A}^2, \seqname{E}^2, \seqname{C}^2,
  \funcname{f}_0, \funcname{f}_1
\end{equation*}
\begin{equation*}
  \seqname{A}^1, E^1, \seqname{C}^1,
  \seqname{A}^2, E^2, \seqname{C}^2,
  \funcname{f}_0, \funcname{f}_1
\end{equation*}
depending on context.

\begin{definition}[Sets of m-atlas morphisms]
\label{def:M-ATLmorphsets}
Let $\seqname{E}^i$ and $\seqname{C}^i$. $i=1,2$, be model spaces. Then
the set of (constrained) (semi-strict, strict)
(full, semi-maximal, maximal, full semi-maximal, full maximal)
$\seqname{E}^1$-$\seqname{E}^2$ m-atlas (near) morphisms in the
coordinate spaces $\seqname{C}^1$,  $\seqname{C}^2$ is

\begin{multline}
  \strict[*]
    {
      \maxfull[*]
        {
          \Atl[(const,const-near,near)]_\Ar
            (\seqname{E}^1, \seqname{C}^1, \seqname{E}^2, \seqname{C}^2)
        }
    }
  \defeq \\*
  \set
  {{
    \bigl (
      (
        \funcname{f}_0,
        \funcname{f}_1
      ),
      (\seqname{A}^1, \seqname{E}^1, \seqname{C}^1),
      (\seqname{A}^2, \seqname{E}^2, \seqname{C}^2)
    \bigr )
  }}%
  [
    {
      \strict[*]
      {
        \maxfull[*]
          {
            \isAtl[(const,const-near,near)]
              (
                \seqname{A}^1, \seqname{E}^1, \seqname{C}^1,
                \seqname{A}^2, \seqname{E}^2, \seqname{C}^2,
                \funcname{f}_0, \funcname{f}_1
              )
          }
      }
    }
  ]*
\end{multline}

Let $E^i$, $i=1,2$, be a topological space and
$\seqname{C}^1=(C^i, \catname{C}^i)$, $i=1,2$, be a model space.
Then the set of (constrained) (semi-strict, strict)
(full, semi-maximal, maximal, full semi-maximal, full maximal)
$E^1$-$E^2$ m-atlas (near) morphisms in the
coordinate spaces $\seqname{C}^1$,  $\seqname{C}^2$ is

\begin{multline}
  \strict[*]
    {
      \maxfull[*]
        {
          \Atl[(const,const-near,near)]_\Ar
            (E^1, \seqname{C}^1, E^2, \seqname{C}^2)
        }
    }
  \defeq \\*
  \set
  {{
    \bigl (
      (
        \funcname{f}_0,
        \funcname{f}_1
      ),
      (\seqname{A}^1, E^1, \seqname{C}^1),
      (\seqname{A}^2, E^2, \seqname{C}^2)
    \bigr )
  }}%
  [
    {
      \strict[*]
      {
        \maxfull[*]
          {
            \isAtl[(const,const-near,near)]
              (
                \seqname{A}^1, E^1, \seqname{C}^1,
                \seqname{A}^2, E^2, \seqname{C}^2,
                \funcname{f}_0, \funcname{f}_1
              )
          }
      }
    }
  ]*
\end{multline}

Let
\begin{enumerate}
\item
$\seqname{E}$ and $\seqname{C}$ be sets of model spaces
\item
$\seqname{P} \defeq \seqname{E} \times \seqname{C}$
\end{enumerate}

\begin{multline}
  \strict[*]
    {
      \maxfull[*]
        {
          \Atl[(const,const-near,near)]_\Ar
            (\seqname{E}, \seqname{C})
        }
    }
  \defeq \\*
  \set
  {{
    \bigl (
      (
        \funcname{f}_0,
        \funcname{f}_1
      ),
      (\seqname{A}^1, \seqname{E}^1, \seqname{C}^1),
      (\seqname{A}^2, \seqname{E}^2, \seqname{C}^2)
    \bigr )
  }}%
  [
    {(\seqname{E}^i, \seqname{C}^i) \in \seqname{P}},
    {
      \strict[*]
      {
        \maxfull[*]
          {
            \isAtl[(const,const-near,near)]
              (
                ...
              )
          }
      }
    }
  ]*
\end{multline}

Let
\begin{enumerate}
\item
$\catname{E}$ be a model category
\item
$\seqname{C}$ be a set of model spaces
\item
$\seqname{P} \defeq \Ob(\catname{E}) \times \seqname{C}$
\end{enumerate}

\begin{multline}
  \strict[*]
    {
      \maxfull[*]
        {
          \Atl[(const,const-near,near)]_\Ar
            (\seqname{E}, \catname{C})
        }
    }
  \defeq \\*
  \set
  {{
    \bigl (
      (
        \funcname{f}_0,
        \funcname{f}_1
      ),
      (\seqname{A}^1, \seqname{E}^1, \seqname{C}^1),
      (\seqname{A}^2, \seqname{E}^2, \seqname{C}^2)
    \bigr )
  }}%
  [
    {(\seqname{E}^i, \seqname{C}^i) \in \seqname{P}},
    {
      \strict[*]
      {
        \maxfull[*]
          {
            \isAtl[(const,const-near,near)]
              (
                ...,
                \catname{E},
                \catname{E}
              )
          }
      }
    }
  ]*
\end{multline}

Let
\begin{enumerate}
\item
$\seqname{E}$ be a set of model spaces
\item
$\seqname{C}$ be a model category
\item
$\seqname{P} \defeq \seqname{E} \times Ob(\catname{C})$
\end{enumerate}

\begin{multline}
  \strict[*]
    {
      \maxfull[*]
        {
          \Atl[(const,const-near,near)]_\Ar
            (\catname{E}, \catname{C})
        }
    }
  \defeq \\*
  \set
  {{
    \bigl (
      (
        \funcname{f}_0,
        \funcname{f}_1
      ),
      (\seqname{A}^1, \seqname{E}^1, \seqname{C}^1),
      (\seqname{A}^2, \seqname{E}^2, \seqname{C}^2)
    \bigr )
  }}%
  [
    {(\seqname{E}^i, \seqname{C}^i) \in \seqname{P}},
    {
      \strict[*]
      {
        \maxfull[*]
          {
            \isAtl[(const,const-near,near)]
              (
                ...,
                \catname{C},
                \catname{C}
              )
          }
      }
    }
  ]*
\end{multline}

Let
\begin{enumerate}
\item
$\catname{E}$ and $\catname{C}$ be model categories
\item
$\seqname{P} \defeq \Ob(\catname{E}) \times Ob(\catname{C})$
\end{enumerate}

\begin{multline}
  \strict[*]
    {
      \maxfull[*]
        {
          \Atl[(const,const-near,near)]_\Ar
            (\catname{E}, \catname{C})
        }
    }
  \defeq \\*
  \set
  {{
    \bigl (
      (
        \funcname{f}_0,
        \funcname{f}_1
      ),
      (\seqname{A}^1, \seqname{E}^1, \seqname{C}^1),
      (\seqname{A}^2, \seqname{E}^2, \seqname{C}^2)
    \bigr )
  }}%
  [
    {(\seqname{E}^i, \seqname{C}^i) \in \seqname{P}},
    {
      \strict[*]
      {
        \maxfull[*]
          {
            \isAtl[(const,const-near,near)]
              (
                ...,
                \catname{E},
                \catname{E},
                \catname{C},
                \catname{C}
              )
          }
      }
    }
  ]*
\end{multline}
\end{definition}

\begin{lemma}[Sets of m-atlas morphisms]
\label{lem:M-ATLmorphsets}
Let $\seqname{E}$ and $\seqname{C}$ be sets of model spaces.

$
  \strict[*]
    {
      \maxfull[*]
        {
          \Atl[(const,const-near,near)]_\Ar
            (\seqname{E}, \seqname{C})
        }
    }
  \subseteq
  \Atl[(near)]_\Ar
    (\seqname{E}, \seqname{C})
$.

\begin{align*}
  \strict[*]
    {
      \maxfull
        {
          \Atl[(const,const-near,near)]_\Ar
            (\seqname{E}, \seqname{C})
        }
    }
  & \subseteq
  \strict[*]
    {
      \maxfull[S-]
        {
          \Atl[(const,const-near,near)]_\Ar
            (\seqname{E}, \seqname{C})
        }
    }
  \\*
  & \subseteq
  \strict[*]
    {
      \full
        {
          \Atl[(const,const-near,near)]_\Ar
            (\seqname{E}, \seqname{C})
        }
    }
\end{align*}

\begin{proof}
The result follows from
{
  \showlabelsinline
  \pagecref{def:M-ATLmorphabbrev}
}
and \cref{def:M-ATLmorphsets} above.
\end{proof}

\begin{equation*}
  \strict[*]
    {
      \maxfull[*]
        {
          \Atl_\Ar
            (\seqname{E}, \seqname{C})
        }
    }
  \subseteq
  \strict[*]
    {
      \maxfull[*]
        {
          \Atl[near]_\Ar
            (\seqname{E}, \seqname{C})
        }
    }
\end{equation*}

\begin{equation*}
  \strict[*]
    {
      \maximal[*]
        {
          \Atl_\Ar
            (\seqname{E}, \seqname{C})
        }
    }
  =
  \strict[*]
    {
      \maximal[*]
        {
          \Atl[near]_\Ar
            (\seqname{E}, \seqname{C})
        }
    }
\end{equation*}

\begin{proof}
The result follows from
\cref{M-ATLmorph:morphisnear} of
{
  \showlabelsinline
  \fullcref{lem:M-ATLmorph} on
  \cpageref{M-ATLmorph:morphisnear},
  \pagecref{def:M-ATLmorphabbrev}
}
and \cref{def:M-ATLmorphsets} above.
\end{proof}

\begin{proof}
The result follows from \pagecref{lem:M-ATLmorph},
{
  \showlabelsinline
  \pagecref{def:M-ATLmorphabbrev}
}
and \cref{def:M-ATLmorphsets} above.
\end{proof}
\end{lemma}

\begin{corollary}[Sets of m-atlas morphisms]
\label{cor:M-ATLmorphsets}
Let $\seqname{E}$ and $\seqname{C}$ be sets of topological spaces.

\begin{equation*}
  \strict[*]
    {
      \maxfull[*]
        {
          \Atl[(const,const-near,near)]_\Ar
            (\seqname{E}, \seqname{C})
        }
    }
  \subseteq
  \Atl[(near)]_\Ar
    (\seqname{E}, \seqname{C})
\end{equation*}

\begin{align*}
  \strict[*]
    {
      \maxfull
        {
          \Atl[(const,const-near,near)]_\Ar
            (\seqname{E}, \seqname{C})
        }
    }
  & \subseteq
  \strict[*]
    {
      \maxfull[S-]
        {
          \Atl[(const,const-near,near)]_\Ar
            (\seqname{E}, \seqname{C})
        }
    }
  \\*
  & \subseteq
  \strict[*]
    {
      \full
        {
          \Atl[(const,const-near,near)]_\Ar
            (\seqname{E}, \seqname{C})
        }
    }
\end{align*}

\begin{equation*}
  \strict[*]
    {
      \maxfull[*]
        {
          \Atl_\Ar
            (\seqname{E}, \seqname{C})
        }
    }
  \subseteq
  \strict[*]
    {
      \maxfull[*]
        {
          \Atl[near]_\Ar
            (\seqname{E}, \seqname{C})
        }
    }
\end{equation*}

\begin{proof}
The result follows from
{
  \showlabelsinline
  \pagecref{def:M-ATLmorphabbrev}
}
and \cref{def:M-ATLmorphsets} above.
\end{proof}
\end{corollary}

\begin{definition}[Categories $\Atl(\seqname{E}, \seqname{C})$]
\label{def:M-cat}
Let $\seqname{E}$ and $\seqname{C}$ be sets of model spaces.
Let $P \defeq \seqname{E} \times \seqname{C}$.

\begin{multline}
  \strict[*]
    {
      \maxfull[*]
        {
          \Atl[(const)]
            (\seqname{E}, \seqname{C})
        }
    }
  \defeq \\*
  \biggl (
    \maxfull[*]
      {
        \Atl_\Ob
          (\seqname{E}, \seqname{C})
      },
    \strict[*]
      {
        \maxfull[*]
          {
            \Atl[(const)]_\Ar
              (\seqname{E}, \seqname{C})
          }
      },
    \compose[A]
  \biggr )
\end{multline}

Let $\seqname{E}^i$ and $\seqname{C}^i$, $i=1,2$, be model spaces,
$\seqname{A}^i$ be an atlas of $\seqname{E}^i$ in the coordinate
space $\seqname{C}^i$ and
$
\funcseqname{f} \defeq
(
  \funcname{f}_0 \maps \seqname{E}^1 \to \seqname{E}^2,
  \funcname{f}_0 \maps \seqname{C}^1 \to \seqname{C}^2
)
$
be a (strict, semi-strict) $\seqname{E}^1$-$\seqname{E}^2$ M-atlas
(near) morphism of $\seqname{A}^1$ to $\seqname{A}^2$ in the coordinate
spaces $\seqname{C}^1$, $\seqname{C}^2$.

The identity functor $\Functor_{\M,\Id}$ is

\begin{equation}
\Functor_{\M,\Id} (\seqname{A}^i, \seqname{E}^i, \seqname{C}^i)
\defeq
(\seqname{A}^i, \seqname{E}^i, \seqname{C}^i)
\end{equation}

This nomenclature will be justified below.

\begin{multline}
\Functor_{\M,\Id}
  \biggl (
    (
      \funcname{f}_0 \maps \seqname{E}^1 \to \seqname{E}^2,
      \funcname{f}_1 \maps \seqname{C}^1 \to \seqname{C}^2,
    ),
    (\seqname{A}^1, \seqname{E}^1, \seqname{C}^1),
    (\seqname{A}^2, \seqname{E}^2, \seqname{C}^2)
  \biggr )
\defeq
\\*
  \bigl (
    (
      \funcname{f}_0 \maps \seqname{E}^1 \to \seqname{E}^2,
      \funcname{f}_1 \maps \seqname{C}^1 \to \seqname{C}^2,
    ),
    (\seqname{A}^1, \seqname{E}^1, \seqname{C}^1),
    (\seqname{A}^2, \seqname{E}^2, \seqname{C}^2)
  \bigr )
\end{multline}

These are fine grained iff $\seqname{C}$ is fine grained.
\end{definition}

\begin{lemma}[$\Atl(\seqname{E}, \seqname{C})$ is a category]
\label{lem:M-ATLiscat}
Let $\seqname{E}$ and $\seqname{C}$ be sets of model spaces. Then
\begin{enumerate}
\item
\label{Mcat:near}
$
  \strict[*]
    {
      \maximal[*]
        {
          \Atl[(const-near,near)]
            (\seqname{E}, \seqname{C})
        }
    }
$
is a category and the identity
morphism for an object $(\seqname{A}^i, \seqname{E}^i, \seqname{C}^i)$
of
$
  \strict[*]
    {
      \maximal[*]
        {
          \Atl[(const-near,near)]_\Ob
            (\seqname{E}, \seqname{C})
        }
    }
$
is
$\Id[{{(\seqname{A}^i, \seqname{E}^i, \seqname{C}^i)}}]$.
\item
\label{Mcat:morph}
$
  \strict[*]
    {
      \maxfull[*]
        {
          \Atl[(const)]
            (\seqname{E}, \seqname{C})
        }
    }
$
is a category and the identity
morphism for an object $(\seqname{A}^i, \seqname{E}^i, \seqname{C}^i)$
of
$
  \strict[*]
    {
      \maxfull[*]
        {
          \Atl[(const)]_\Ob
            (\seqname{E}, \seqname{C})
        }
    }
$
is
$\Id[{{(\seqname{A}^i, \seqname{E}^i, \seqname{C}^i)}}]$.

\begin{proof}
Let $(\seqname{A}^i,\seqname{E}^i,\seqname{C}^i)$, $i \in [1,3]$,
be an object of $\Atl(\seqname{E}, \seqname{C})$ and
let \\
$
  \funcname{m}^i \defeq
  \bigl (
    (\funcname{f}_0^i,\funcname{f}_1^i),
    (\seqname{A}^i,\seqname{E}^i,\seqname{C}^i),
    (\seqname{A}^{i+1},\seqname{E}^{i+1},\seqname{C}^{i+1})
  \bigr )
$
be a morphism of $\Atl(\seqname{E}, \seqname{C})$. Then
\begin{description}
\item [Composition:]
\leavevmode \\*
$
    \funcseqname{m}^2 \compose[A] \funcseqname{m}^1 =
  \bigl (
    (
      \funcname{f}_0^2 \compose \funcname{f}_0^1,
      \funcname{f}_1^2 \compose \funcname{f}_1^1
    ),
    (\seqname{A}^1,\seqname{E}^1,\seqname{C}^1),
    (\seqname{A}^3,\seqname{E}^3,\seqname{C}^3)
  \bigr )
$
is a morphism of $\Atl(\seqname{E}, \seqname{C})$ by
\cref{M-morphcomp:morphcompmorph} of
{
  \showlabelsinline
  \fullcref{lem:M-ATLmorphcomp} on
  \cpageref{M-morphcomp:morphcompmorph}.
}

\item [Associativity:]
\leavevmode \\*
Composition is associative by \pagecref{lem:labcomp}.

\item [Identity:]
\leavevmode \\*
$\Id[{{(\seqname{A}^i, \seqname{E}^i, \seqname{C}^i)}}]$ is an identity
morphism by \cref{lem:labcomp}.
\end{description}
\end{proof}
\end{enumerate}
\end{lemma}

\begin{corollary}[Subcategories of $\Atl(\seqname{E}, \seqname{C})$]
\label{cor:M-ATLiscat}
Let $\seqname{E}$ and $\seqname{C}$ be sets of model spaces. Then
$
  \strict[*]
    {
      \maxfull[*]
        {
          \Atl[(const)]
            (\seqname{E}, \seqname{C})
        }
    }
$
is a subcategory of
$\Atl[(const)] (\seqname{E}, \seqname{C})$.

\begin{remark}
They are not, in general, full subcategories.
\end{remark}

The identity functor $\Functor_{\M,\Id}$ is a functor from each of the
subcategories to each of the containing categories.
\end{corollary}

\begin{definition}[Categories $\Atl(\catname{E}, \catname{C})$]
\label{def:M-cat-cat}
Let $\catname{E}$ and $\catname{C}$ be model categories.
Let $\seqname{P} \defeq Ob(\catname{E}) \times \Ob(\catname{C})$.
Then

\begin{multline}
  \maxfull[*]
    {
      \Atl_\Ob(\catname{E}, \catname{C})
    }
  \defeq \\*
  \union
    [
        {\seqname{E} \objin \catname{E}},
        {\seqname{C} \objin \catname{C}}
    ]
    {{
      \maxfull[*]
        {
          \Atl_\Ob(\seqname{E}, \seqname{C})
        }
    }}
\end{multline}
\begin{multline}
  \strict[*]
    {
      \maxfull[*]
        {
          \Atl[(const,const-near,near)]_\Ar
            (\catname{E}, \catname{C})
        }
    }
  \defeq \\*
    \set
    {{
      \bigl (
        (
          \funcname{f}_0,
          \funcname{f}_1
        ),
        (\seqname{A}^1, \seqname{E}^1, \seqname{C}^1),
        (\seqname{A}^2, \seqname{E}^2, \seqname{C}^2)
      \bigr )
    }}%
    [
      {
        (\seqname{E}^i, \seqname{C}^i) \in \seqname{P}
      },
      {
        \strict[*]
        {
          \maxfull[*]
          {
            \isAtl[(const,const-near,near)]_\Ar
            (
              \seqname{A}^1,
              \seqname{E}^1,
              \seqname{C}^1,
              \seqname{A}^2,
              \seqname{E}^2,
              \seqname{C}^2,
              \funcname{f}_0,
              \funcname{f}_1,
              \catname{E},
              \catname{C}
            )
          }
        }
      }
    ]*
\end{multline}
\begin{multline}
  \strict[*]
  {
    \maxfull[*]
    {
      \Atl[(const,const-near,near)](\catname{E}, \catname{C})
    }
  }
  \defeq \\*
  \bigl (
    \Atl_\Ob(\catname{E}, \catname{C}),
    \Atl[(const,const-near,near)]_\Ar(\catname{E}, \catname{C}),
    \compose[A]
  \bigr )
\end{multline}
\end{definition}

\begin{lemma}[$\Atl(\catname{E}, \catname{C})$ is a category]
\label{lem:M-cat-ATLiscat}
Let $\catname{E}$ and $\catname{C}$ be model categories. Then

\begin{enumerate}
\item
$
  \strict[*]
  {
    \maxfull[*]
    {
      \Atl[(const)](\catname{E}, \catname{C})
    }
  }
$
is a category and the identity morphism for an object
$(\seqname{A}^i, \seqname{E}^i, \seqname{C}^i)$ of
$
  \strict[*]
  {
    \maxfull[*]
    {
      \Atl[(const)](\catname{E}, \catname{C})
    }
  }
$
is $\Id_{(\seqname{A}^i, \seqname{E}^i, \seqname{C}^i}$.

\begin{proof}
Let $(\seqname{A}^i,\seqname{E}^i,\seqname{C}^i)$, $i \in [1,3]$,
be an object of $\Atl(\catname{E}, \catname{C})$ and
let \\
$
  \funcname{m}^i \defeq
  \bigl (
    (\funcname{f}_0^i,\funcname{f}_1^i),
    (\seqname{A}^i,\seqname{E}^i,\seqname{C}^i),
    (\seqname{A}^{i+1},\seqname{E}^{i+1},\seqname{C}^{i+1})
  \bigr )
$
be a morphism of $\Atl(\catname{E}, \catname{C})$. Then
\begin{description}
\item[Composition:]
\leavevmode \\*
$\funcname{f}_0^2 \compose \funcname{f}_0^1$ is a morphism of
$\catname{E}$, $\funcname{f}_1^2 \compose \funcname{f}_1^1$ is a
morphism of $\catname{C}$ and
$
  \funcseqname{f}^2 \compose[()] \funcseqname{f}^1 =
  (
    \funcname{f}_0^2 \compose \funcname{f}_0^1,
    \funcname{f}_1^2 \compose \funcname{f}_1^1
  )
$
is an $E^1$-$E^3$ m-atlas morphism of $\seqname{A}^1$ to $\seqname{A}^3$
in the coordinate spaces $\seqname{C}^1$, $\seqname{C}^2$ by
\cref{M-morphcomp:morphcompmorph} of
{
  \showlabelsinline
  \fullcref{lem:M-ATLmorphcomp} on
  \cpageref{M-morphcomp:morphcompmorph}.
}

\item[Associativity:]
\leavevmode \\*
Composition is associative by \pagecref{lem:labcomp}.

\item[Identity:]
$\Id[{{(\seqname{A}^i, \seqname{E}^i, \seqname{C}^i)}}]$ is an identity
morphism by \cref{lem:labcomp}.
\end{description}
\end{proof}

\item
$
  \strict[*]
  {
    \maximal[*]
    {
      \Atl[(const-near,near)](\catname{E}, \catname{C})
    }
  }
$
is a category and the identity morphism for an object
$(\seqname{A}^i, \seqname{E}^i, \seqname{C}^i)$ of
$
  \strict[*]
  {
    \maximal[*]
    {
      \Atl[(const-near,near)](\catname{E}, \catname{C})
    }
  }
$
is $\Id_{(\seqname{A}^i, \seqname{E}^i, \seqname{C}^i}$.
\end{enumerate}
\end{lemma}

\begin{definition}[Categories $\Atl(\seqname{E}, \seqname{C})$ of topological spaces]
\label{def:T-cat}
Let $\seqname{E}$ be a set of topological spaces and $\seqname{C}$ a
set of model spaces.
Let $\seqname{P} \defeq \seqname{E} \times \seqname{C}$.
Then

\begin{multline}
  \strict[*]
    {
      \maxfull[*]
        {
          \Atl[(const)]
            (\seqname{E}, \seqname{C})
        }
    }
  \defeq \\*
  \biggl (
    \maxfull[*]
      {
        \Atl_\Ob
          (\seqname{E}, \seqname{C})
      },
    \strict[*]
      {
        \maxfull[*]
          {
            \Atl[(const)]_\Ar
              (\seqname{E}, \seqname{C})
          }
      },
    \compose[A]
  \biggr )
\end{multline}

Let
$
  (\seqname{A}^1, \seqname{E}^1, \seqname{C}^1) \in
  \Atl_\Ob(\seqname{E}, \seqname{C})
$.

\begin{equation}
\Id[{{(\seqname{A}^1, \seqname{E}^1, \seqname{C}^1)}}] \defeq
\bigl (
  (\Id[{\seqname{E}^1}], \Id[{\seqname{C}^1}]),
  (\seqname{A}^1, \seqname{E}^1, \seqname{C}^1),
  (\seqname{A}^1, \seqname{E}^1, \seqname{C}^1)
\bigr )
\end{equation}
\end{definition}

\begin{lemma}[$\Atl(\seqname{E}, \seqname{C})$ of topological spaces is a category]
\label{lem:T-ATLiscat}
Let $\seqname{E}$ be a set of topological spaces and $\seqname{C}$ a
set of model spaces. Then $\Atl(\seqname{E}, \seqname{C})$ is a category
and the identity morphism for an object
$(\seqname{A}^1, E^1, \seqname{C}^1)$ is
$\Id[{{(\seqname{A}^1, E^1, \seqname{C}^1)}}]$.

\begin{proof}
Let $(\seqname{A}^i,E^i,\seqname{C}^i)$, $i \in [1,3]$,
be an object of $\Atl(E, \seqname{C})$. and let
$
  \funcname{m}^i \defeq \\
  \bigl (
    (\funcname{f}_0^i,\funcname{f}_1^i),
    (\seqname{A}^i,E^i,\seqname{C}^i),
    (\seqname{A}^{i+1},E^{i+1},\seqname{C}^{i+1})
  \bigr )
$
be a morphism of $\Atl(\seqname{E}, \seqname{C})$. Then
\begin{description}
\item[Composition:]
\leavevmode \\*
$
  \bigl (
    (
      \funcname{f}_0^2 \compose \funcname{f}_0^1,
      \funcname{f}_1^2 \compose \funcname{f}_1^1
    ),
    (\seqname{A}^1,E^1,\seqname{C}^1),
    (\seqname{A}^3,E^3,\seqname{C}^3)
  \bigr )
$
is a morphism of $\Atl(\seqname{E}, \seqname{C})$ by
\pagecref{cor:M-ATLmorphcomp}.

\item[Associativity:]
\leavevmode \\*
Composition is associative by \pagecref{lem:labcomp}.

\item[Identity:]
$\Id[{{(\seqname{A}^i, \seqname{E}^i, \seqname{C}^i)}}]$ is an identity
morphism by \cref{lem:labcomp}.
\end{description}
\end{proof}

Similar results follow with restrictions on the admissible atlases,
the admissible morphisms, or both:
\begin{enumerate}
\item $\full{\Atl}(\catname{E}, \catname{C})$
\item $\maximal[S-]{\Atl}(\catname{E}, \catname{C})$
\item $\maximal{\Atl}(\catname{E}, \catname{C})$
\item $\maxfull[S-]{\Atl}(\catname{E}, \catname{C})$
\item $\maxfull{\Atl}(\catname{E}, \catname{C})$
\item $\strict[semi-]{\full{\Atl}}(\catname{E}, \catname{C})$
\item $\strict[semi-]{\maximal[S-]{\Atl}}(\catname{E}, \catname{C})$
\item $\strict[semi-]{\maximal{\Atl}}(\catname{E}, \catname{C})$
\item $\strict[semi-]{\maxfull[S-]{\Atl}}(\catname{E}, \catname{C})$
\item $\strict[semi-]{\maxfull{\Atl}}(\catname{E}, \catname{C})$
\item $\strict{\full{\Atl}}(\catname{E}, \catname{C})$
\item $\strict{\maximal[S-]{\Atl}}(\catname{E}, \catname{C})$
\item $\strict{\maximal{\Atl}}(\catname{E}, \catname{C})$
\item $\strict{\maxfull[S-]{\Atl}}(\catname{E}, \catname{C})$
\item $\strict{\maxfull{\Atl}}(\catname{E}, \catname{C})$
\end{enumerate}
\end{lemma}

\begin{definition}[$\Triv{\seqname{E}}$ and $\Triv{\Atl}(\seqname{E}, \seqname{C})$]
\label{def:M-triv}
Let $\seqname{E}$ be a set of topological spaces and $\seqname{C}$ a
set of model spaces. Then
\begin{equation}
\Triv{\seqname{E}} \defeq \set {\Triv{E^\mu}}[E^\mu \in \seqname{E}]
\end{equation}
\begin{equation}
\Triv{\Atl}(\seqname{E}, \seqname{C}) \defeq
\Atl \Big ( \Triv{\seqname{E}}, \seqname{C} \Big )
\end{equation}
\end{definition}

\subsubsection{M-atlas functors}
\begin{definition}[$\Functor_{\M,\Top}$]
\label{def:Func-M-Top}
Let $\seqname{E}^i \defeq (E^i,\catname{E}^i)$, $i=1,2$, and
$\seqname{C}^i \defeq (C^i,\catname{C}^i)$
be model spaces and $\seqname{A}^i$ be an m-atlas of
$\seqname{E}^i$ in the coordinate space $\seqname{C}^i$. Then
\begin{equation}
\Functor_{\M,\Top} (\seqname{A}^i, \seqname{E}^i, \seqname{C}^i)
\defeq
(\seqname{A}^i, E^i, \seqname{C}^i)
\end{equation}

Let $\funcname{f}_0 \maps \seqname{E}^1 \to \seqname{E}^2$ and
$\funcname{f}_1 \maps \seqname{C}^1 \to \seqname{C}^2$ be model
functions. Then
\begin{multline}
\Functor_{\M,\Top}
  \bigl (
    (
      \funcname{f}_0 \maps \seqname{E}^1 \to \seqname{E}^2,
      \funcname{f}_1 \maps \seqname{C}^1 \to \seqname{C}^2,
    ),
    (\seqname{A}^1, \seqname{E}^1, \seqname{C}^1),
    (\seqname{A}^2, \seqname{E}^2, \seqname{C}^2)
  \bigr )
\defeq \\
  \bigl (
    (
      \funcname{f}_0 \maps E^1 \to E^2,
      \funcname{f}_1 \maps \seqname{C}^1 \to \seqname{C}^2
    ),
    ( \seqname{A}^1, E^1, \seqname{C}^1 ),
    ( \seqname{A}^2, E^2, \seqname{C}^2 )
  \bigr )
\end{multline}
\end{definition}

\begin{lemma}[$\Functor_{\M,\Top}$ maps M-atlases to M-atlases if the coordinate model space is fine grained.]
\label{lem:Func-M-Top}
Let $\seqname{E} \defeq (E,\catname{E})$ and $\seqname{C}$ be model
spaces, $\seqname{C}$ be fine grained and $\seqname{A}$ be an M-atlas of
$\seqname{E}$ in the model space $\seqname{C}$. Then
$\seqname{A}$ is an M-atlas of $E$ in the model space $\seqname{C}$.

\begin{proof}
Let $(U, V, \phi)$ be an M-chart in $\seqname{A}$, Every open subset of
$V$ is a model neighborhood of $\seqname{C}$ since it is fine grained by
hypothesis. Then $(U, V, \phi)$ is an M-chart of $E$ in the model space
$\seqname{C}$ by \pagecref{lem:M-chart}.

The definition of mutually compatible charts
neighborhood of $\seqname{C}$. $(U, V, \phi)$ is also an M-chart of $E$
in the coordinate space $\seqname{C}$:

Since $U$ is a model neighborhood of $\seqname{E}$, it is open and hence
a model neighborhood of $\Triv{E}$.

$V$ is a model neighborhood of $\seqname{C}$ by hypothesis.

Since $U'$ is a model neighborhood of $\Triv{E}$ and hence open, and
$\phi$ is a homeomorphism, $\phi[U']$ is open.
Since $V$ is a model neighborhood of $\seqname{C}$,
$\seqname{C}$ is fine grained and $\phi[U'] \subseteq V$ is open,
$\phi[U']$ is a model neighborhood of $\seqname{C}$,

Since $V'$ is a model neighborhood of $\seqname{C}$, it is open,
$\phi^{-1}[V']$ is open and hence a model neighborhood of $\Triv{E}$.
\end{proof}
\end{lemma}

\begin{theorem}[$\Functor_{\M,\Top}$ is a functor if $\seqname{C}$ is fine grained.]
\label{the:Func-M-Top}
Let $\seqname{E}$ and $\seqname{C}$ be sets of model spaces and
$\seqname{C}$ be fine grained. Then $\Functor_{\M,\Top}$ is a functor
from $\Atl(\seqname{E}, \seqname{C})$ to
$\Atl \bigl ( \Top(\seqname{E}), \seqname{C} \bigr )$

\begin{proof}
Let
$
  o^i \defeq
  \bigl (
    \seqname{A}^i, \seqname{E}^i \defeq (E^i,\catname{E^i}), \seqname{C}^i
  \bigr )
$,
$i \in [1,3]$, be an object of $\Atl(\seqname{E}, \seqname{C})$,
$
  o'^i \defeq \Functor_{\M,\Top} o^i =
  (\seqname{A}^i, E^i, \seqname{C}^i)
$
the corresponding object of $\Atl \bigl ( \Top(\seqname{E}), \seqname{C} \bigr )$,
$
  m^i \defeq \\*
  \bigl (
    (
      \funcname{f}_0^i \maps E^i \to E^{i+1},
      \funcname{f}_1^i \maps \seqname{C}^i \to \seqname{C}^{i+1},
    ),
    o^i,
    o^{i+1}
  \bigr )
$, $i=1,2$,
be a morphism of $\Atl(\seqname{E}, \seqname{C})$
and
$
  m'^i \defeq \Functor_{\M,\Top} m^i =
  \bigl (
    (
      \funcname{f}_0^i \maps E^i \to E^{i+1},
      \funcname{f}_1^i \maps \seqname{C}^i \to \seqname{C}^{i+1}
    ),
    o'^i,
    o'^{i+1}
  \bigr )
$
the corresponding morphism of $\Atl \bigl ( \Top(\seqname{E}), \seqname{C} \bigr )$.

$\Functor_{\M,\Top}$ satisfies these criteria:
\begin{description}
\item[Preservation of objects:]
\leavevmode \\*
$\seqname{A}^i$ is an M-atlas of $E^i$ in the coordinate space
$\seqname{C}^i$ by \fullcref{lem:Func-M-Top} above.

\item[Preservation of morphisms:]
\leavevmode \\*
$\funcname{f}_0^i \maps E^i \to E^{i+1}$ is continuous and
$\funcname{f}_1^i \maps \seqname{C}^i \to \seqname{C}^{i+1}$ is a model
function.

\item[Preservation of endpoints:]
\begin{align*}
  \Functor_{\M,\Top} o^i
  & =
  o'^i \objin \Atl \bigl ( \Top(\seqname{E}), \seqname{C} \bigr )
\\*
  \Functor_{\M,\Top} m^i
  & =
  \bigl (
    (
      \funcname{f}_0^i \maps E^i \to E^{i+1},
      \funcname{f}_1^i \maps \seqname{C}^i \to \seqname{C}^{i+1}
    ),
    o'^i,
    o'^{i+1}
  \bigr )
  \arin \Atl \bigl ( \Top(\seqname{E}), \seqname{C} \bigr )
\\*
  & \quad
\text{is a morphism from $o'^i$ to $o'^{i+1}$.}
\end{align*}

\item[Identity:]
\leavevmode \\*
Let
$
  \bigl (
    (\Id[{\seqname{E}^i}], \Id[{\seqname{C}^i}]),
   o^i,
   o^i
  \bigr )
$
be an identity morphism of
$o^i$ in $\Atl(\seqname{E}, \seqname{C})$.
Then
$
  \bigl (
    (
      \Id[E^i] \maps \Top(\seqname{E}^i) \to \Top(\seqname{E}^i),
      \Id[{\seqname{C}^i}]
    ),
    o'^i,
    o'^i
  \bigr )
$
is an identity morphism of $o'^i$ in $\Atl \bigl ( \Top(\seqname{E}), \seqname{C} \bigr )$.

\item[Composition:]
\begin{equation*}
\begin{array}{lll}
  \Functor_{\M,\Top} (m^2 \compose[A] m^1)
  & =
  &
  \Functor_{\M,\Top}
    \bigl (
      (
        \funcname{f}_0^2 \compose \funcname{f}_0^1,
        \funcname{f}_1^2 \compose \funcname{f}_1^1
      ),
      o^1,
      o^3
    \bigr )
\\*
  & =
  &
    \bigl (
      (
        \funcname{f}_0^2 \compose \funcname{f}_0^1,
        \funcname{f}_1^2 \compose \funcname{f}_1^1
      ),
      o'^1,
      o'^3
    \bigr )
\\*
  & =
  &
    \bigl (
      (\funcname{f}_0^2, \funcname{f}_1^2),
      o'^1,
      o'^2
    \bigr )
    \compose[A]
    \bigl (
      (\funcname{f}_0^1, \funcname{f}_1^1),
      o'^2,
      o'^3
    \bigr )
\\*
  & =
  &
  \Functor_{\M,\Top}
    \bigl (
      (\funcname{f}_0^2, \funcname{f}_1^2),
      o^2,
      o^3
    \bigr )
  \compose[A]
\\*
  &
  &
    \Functor_{\M,\Top}
    \bigl (
      (\funcname{f}_0^1, \funcname{f}_1^1),
      o^1,
      o^2
    \bigr )
\\*
  & =
  &
  \Functor_{\M,\Top} m^2
  \compose[A]
  \Functor_{\M,\Top} m^1
\end{array}
\end{equation*}
\end{description}
\end{proof}
\end{theorem}

\begin{definition}[$\Functor_{\Top,\M}$]
\label{def:Func-Top-M}
Let $E^i$, $i=1,2$, be a topological space, $\seqname{C}^i$ a model
space and $\seqname{A}^i$ an m-atlas of $E^i$
and $\seqname{A}$ an m-atlas of $E$ in the coordinate space
$\seqname{C}^i$. Then
\begin{equation}
\Functor_{\Top,\M} (\seqname{A}^i, E^i, \seqname{C}^i)
\defeq
\Big ( \seqname{A}, \Triv{E^i}, \seqname{C}^i \Big )
\end{equation}

Let $\funcname{f}_0 \maps E^1 \to E^2$ be continuous and
$\funcname{f}_1 \maps \seqname{C}^1 \to \seqname{C}^2$ be a model
function. Then
\begin{multline}
\Functor_{\Top,\M}
  \bigl (
    (
      \funcname{f}_0 \maps E^1 \to E^2,
      \funcname{f}_1 \maps \seqname{C}^1 \to \seqname{C}^2,
    ),
    (\seqname{A}^1, E^1, \seqname{C}^1),
    (\seqname{A}^2, E^2, \seqname{C}^2)
  \bigr )
\defeq \\
  \biggl (
    \Bigl (
      \funcname{f}_0 \maps \Triv{E^1} \to \Triv{E^2},
      \funcname{f}_1 \maps \seqname{C}^1 \to \seqname{C}^2
     \Bigr ),
    \Bigl ( \seqname{A}^1, \Triv{E^1}, \seqname{C}^1 \Bigr ),
    \Bigl ( \seqname{A}^2, \Triv{E^2}, \seqname{C}^2 \Bigr )
  \biggr )
\end{multline}
\end{definition}

\begin{lemma}[$\Functor_{\Top,\M}$ maps M-atlases to M-atlases]
\label{lem:Func-Top-M}
Let $E$ be a topological space, $C$ be a model space and $\seqname{A}$
be an M-atlas of $E$ in the model space $\seqname{C}$. Then
$\seqname{A}$ is an M-atlas of $\Triv{E}$ in the model space
$\seqname{C}$.

\begin{proof}
Let $(U, V, \phi)$ be a M-chart in $\seqname{A}$, $U' \subseteq U$ a
model neighborhood of $\Triv{E}$ and $V' \subseteq V$ a model
neighborhood of $\seqname{C}$. $(U, V, \phi)$ is also an M-chart of $E$
in the coordinate space $\seqname{C}$:
\end{proof}
\end{lemma}

\begin{theorem}[$\Functor_{\Top,\M}$ is a functor]
\label{the:Func-Top-M}
Let $\seqname{E}$ be a set of topological spaces and $\seqname{C}$ a set
of model spaces, Then $\Functor_{\Top,\M}$ is a functor from
$\Atl(\seqname{E}, \seqname{C})$ to
$\Atl \big ( \Triv{\seqname{E}}, \seqname{C} \big )$.

\begin{proof}
Let $o^i \defeq (\seqname{A}^i, E^i, \seqname{C}^i)$, $i \in [1,3]$,
be an object of $\Atl(\seqname{E}, \seqname{C})$,
$o'^i \defeq (\seqname{A}^i, \Triv{E^i}, \seqname{C}^i)$
the corresponding object of
$\Atl \big ( \Triv{\seqname{E}}, \seqname{C} \big )$,
$
 m^i \defeq
  \bigl (
    (f_0^i \maps  E^i \to E^{i+1}, f_1^i \maps  \seqname{C}^i \to  \seqname{C}^{i+1}),
   o^i ,
   o^{i+1}
 \bigr )
$,
$i=1,2$, be a morphism of $\Atl(\seqname{E}, \seqname{C})$ and
$
  m'^i \defeq
  \bigl (
    (f_0^i, f_1^i),
    o'^i ,
   o'^{i+1}
  \bigr )
$
the corresponding morphism of
$\Atl \big ( \Triv{\seqname{E}}, \seqname{C} \big )$.

$\Functor_{\Top,\M}$ satisfies these criteria:
\begin{description}
\item[Preservation of endpoints:]
$\seqname{A}^i$ is an M-atlas of $\Triv{E^i}$ in the coordinate space
$\seqname{C}^i$ by \fullcref{lem:Func-M-Top} above.
$\funcname{f}_0^i \maps E^i \to E^{i+1}$ is continuous and
$\funcname{f}_1^i \maps \seqname{C}^i \to \seqname{C}^{i+1}$ is a model
function.

\begin{align*}
  \Functor_{\Top,\M} o^i
  & =
  o'^i \objin \Atl \big ( \Triv{\seqname{E}}, \seqname{C} \big )
\\*
  \Functor_{\Top,\M} m^i
  & =
  m'^i \arin \Atl \big ( \Triv{\seqname{E}}, \seqname{C} \big )
\\*
  & \quad
\text{is a morphism from $o'^i$ to $o'^{i+1}$.}
\end{align*}

\item[Identity:]
$
  \Id[o^i] \defeq
  \bigl (
    (\Id[E^i], \Id[\seqname{C}^i]),
    o^i,
    o^i
  \bigr )
$.
Then
$
  \Functor_{\Top,\M} \Id[o^i] =
  \bigl (
    (\Id[{\Triv{E^i}}], \Id[{\seqname{C}^i}]),
    o'^i,
    o'^i
  \bigr )
$
is an identity morphism of
$o'^i$.

\item[Composition:]
\begin{equation*}
\begin{array}{lll}
  \Functor_{\Top,\M} (m^2 \compose[A] m^1)
  & =
  &
    \Functor_{\Top,\M}
      \bigl (
        (f_0^2 \compose f_0^1, f_1^2 \compose f_1^1),
        o^1,
        o^3
      \bigr )
\\*
  & =
  &
      \bigl (
        (f_0^2 \compose f_0^1, f_1^2 \compose f_1^1),
        o'^1,
        o'^3
      \bigr )
\\*
  & =
  &
      \bigl (
        (f_0^2, f_1^2),
        o'^1,
        o'^2
      \bigr )
      \compose[A]
      \bigl (
        (f_0^1, f_1^1),
        o'^2,
        o'^3
      \bigr )
\\*
  & =
  &
    \Functor_{\Top,\M}
      \bigl (
        (f_0^2, f_1^2),
        o^2,
        o^3
      \bigr )
    \compose[A]
\\*
  &
  &
    \Functor_{\Top,\M}
      \bigl (
        (f_0^1, f_1^1),
        o^1,
        o^2
      \bigr )
\\*
  & =
  &
  \Functor_{\Top,\M} m^2
  \compose[A]
  \Functor_{\Top,\M} m^1
\end{array}
\end{equation*}
\end{description}
\end{proof}

$\Functor_{\M,\Top} \compose \Functor_{\Top,\M}$ is the identity functor
on $\Atl(\seqname{E}, \seqname{C})$.

\begin{proof}
\begin{align*}
  \Functor_{\M,\Top} \compose \Functor_{\Top,\M} o^i
  & =
  \Functor_{\M,\Top} o'^i
\\
  & =
   o^i
\end{align*}

\begin{align*}
  \Functor_{\M,\Top} \compose \Functor_{\Top,\M} m^i
  & =
  \Functor_{\M,\Top} m'^i
\\
  & =
  m^i
\end{align*}

\end{proof}
\end{theorem}

\section{Associated model spaces and functors}
{
  \showlabelsinline
  \label{sec:assmod}
}

\begin{definition}[Coordinate model spaces associated with m-atlases]
\label{def:M-modC}
Let $\seqname{A}^i$, $i=1,2$, be an m-atlas of $\seqname{E}^i$ in
the coordinate space $\seqname{C}^i$,
$\funcseqname{f} \defeq (\funcname{f}_0,\funcname{f}_1)$ be an
$\seqname{\seqname{E}}^1$-$\seqname{\seqname{E}}^2$ m-atlas (near)
morphism $(\funcname{f}_0,\funcname{f}_1)$ of $\seqname{A}^1$ to
$\seqname{A}^2$ in the coordinate spaces $\seqname{\seqname{C}}^1$, Then:

The minimal coordinate model space with neighborhoods in the m-atlas
$\seqname{A}^i$ of $\seqname{E}^i$ in the coordinate space
$\seqname{C}^i$ is

\begin{multline}
\Functor^{\minimal{\M}}_2 (\seqname{A}^i, \seqname{E}^i, \seqname{C}^i)
\defeq \\
\minimal{\Mod}
  \left (
    \seqname{C}^i,
    \pi_2[\seqname{A}^i],
    \set
      {\phi' \compose \phi^{-1}}%
      [
        \equant
          {
            {(U,V,\phi) \in \seqname{A}^i},
            {(U',V',\phi') \in \seqname{A}^i}
          }
          {
            U \cap U' \ne \emptyset
          }
      ]
  \right )
\end{multline}

The coordinate mapping associated with the
$\seqname{\seqname{E}}^1$-$\seqname{\seqname{E}}^2$ m-atlas (near)
morphism $(\funcname{f}_0,\funcname{f}_1)$ of $\seqname{A}^1$ to
$\seqname{A}^2$ in the coordinate spaces $\seqname{\seqname{C}}^1$,
$\seqname{\seqname{C}}^2$ is

\begin{multline}
\Functor^{\minimal{\M}}_2
  \bigl (
    (
      \funcname{f}_0,
      \funcname{f}_1
    ),
    (\seqname{A}^1, \seqname{E}^1, \seqname{C}^1),
    (\seqname{A}^2, \seqname{E}^2, \seqname{C}^2)
  \bigr )
\defeq \\
  \funcname{f}_1
  \maps
  \Functor^{\minimal{\M}}_2 (\seqname{A}^1, \seqname{E}^1, \seqname{C}^1)
  \to
  \Functor^{\minimal{\M}}_2 (\seqname{A}^2, \seqname{E}^2, \seqname{C}^2)
\end{multline}

If it is a model function then it is also the coordinate m-morphism
associated with the $\seqname{\seqname{E}}^1$-$\seqname{\seqname{E}}^2$
m-atlas (near) morphism $(\funcname{f}_0,\funcname{f}_1)$ of
$\seqname{A}^1$ to $\seqname{A}^2$ in the coordinate spaces
$\seqname{\seqname{C}}^1$, $\seqname{C}^2$.

\end{definition}

\begin{lemma}[Coordinate model spaces associated with m-atlases]
\label{lem:M-ismodC}
Let $\seqname{A}^i$ be an m-atlas of $\seqname{E}^i$ in the coordinate space
$\seqname{C}^i$, $i=1,2$.

$
  \Functor^{\minimal{\M}}_2
    (\seqname{A}^i, \seqname{E}^i, \seqname{C}^i)
$
is a model space.

\begin{proof}
$
  \Functor^{\minimal{\M}}_2
    (\seqname{A}^i, \seqname{E}^i, \seqname{C}^i)
$
satisfies the conditions for a model space.

\begin{enumerate}
\item Since $\pi_2[\seqname{A}^i]$ is an open cover of
$\union{\pi_2[\seqname{A}^i]}$, the set of finite intersections is also
an open cover.

\item Finite intersections of finite intersections are finite intersections

\item Restrictions of continuous functions are continuous

\item If $\funcname{f} \maps A \to B$ is a morphism of
$
  \Functor^{\minimal{\M}}_2
    (\seqname{A}^i, \seqname{E}^i, \seqname{C}^i)
$,
$A, A', B, B'$ model meighborhoods of
$
  \Functor^{\minimal{\M}}_2
    (\seqname{A}^i, \seqname{E}^i, \seqname{C}^i)
$,
$A' \subseteq A$, $B' \subseteq B$ and $\funcname{f}[A'] \subseteq B'$
then since $\funcname{f} \maps A \to B$ is a morphism it is a
restriction of a transition function between sets in
$\pi_2[\seqname{A}^i]$ and its restrictions are also, hence morphisms,
and thus $\funcname{f} \restrictto_{A'} A' \to B'$ is a morphism.

\item
If $(U,V,\phi) \in \seqname{A}^i$ then
$\Id[{V}] = \phi \compose \phi^{-1}$ is a transition function and hence a
morphism of
$
  \Functor^{\minimal{\M}}_2
    (\seqname{A}^i, \seqname{E}^i, \seqname{C}^i)
$.
If $A, A'$ are objects of
$
  \pi_2
    \bigl (
      \Functor^{\minimal{\M}}_2
        (\seqname{A}^i, \seqname{E}^i, \seqname{C}^i)
    \bigr )
$
and $A' \subseteq A$ then the inclusion map
$\funcname{i} \maps A' \hookrightarrow A$ is a restriction of an
identity morphism of
$
  \Functor^{\minimal{\M}}_2
    (\seqname{A}^i, \seqname{E}^i, \seqname{C}^i)
$
and hence a morphism.

\item Restricted sheaf condition: let
\begin{enumerate}
\item $U_\alpha,V_\alpha$, $\alpha \prec \Alpha$, be model
neighborhoods of
$
    \Functor^{\minimal{\M}}_2
      (\seqname{A}^i, \seqname{E}^i, \seqname{C}^i)
$
\item $\funcname{f}_\alpha \maps U_\alpha \to V_\alpha$ be a morphism of
$
      \Functor^{\minimal{\M}}_2
        (\seqname{A}^i, \seqname{E}^i, \seqname{C}^i)
$
\item $U \defeq \union[\alpha \prec \Alpha]{U_\alpha}$
\item $V \defeq \union[\alpha \prec \Alpha]{V_\alpha}$
\item $\funcname{f} \maps U \to V$ be continuous and
$
  \uquant
    {{\alpha \prec \Alpha},{x \in U_\alpha}}
    {\funcname{f}(x) = \funcname{f}_\alpha(x)}
$
\end{enumerate}
Then $\funcname{f}$ is a morphism of $\seqname{C}^i$ and hence a
morphism of
$
      \Functor^{\minimal{\M}}_2
        (\seqname{A}^i, \seqname{E}^i, \seqname{C}^i)
$.
\end{enumerate}
\end{proof}

Let $\funcseqname{f} \defeq (\funcname{f}_0, \funcname{f}_1)$ be an
$\seqname{E}^1$-$\seqname{E}^2$ m-atlas (near) morphism from
$\seqname{A}^1$ to $\seqname{A}^2$ in the coordinate spaces
$\seqname{C}^1$, $\seqname{C}^2$.  If $\funcseqname{f}$ is a morphism or
each $\seqname{A}^i$ is semi-maximal then \\*
$
  \funcname{f}_1
  \maps
  \Functor^{\minimal{\M}}_2 (\seqname{A}^1, \seqname{E}^1, \seqname{C}^1)
  \to
  \Functor^{\minimal{\M}}_2 (\seqname{A}^2, \seqname{E}^2, \seqname{C}^2)
$
is well defined, i.e.,
$
  \funcname{f}_1
    \biggl [
      \Functor^{\minimal{\M}}_2
        (\seqname{A}^1, \seqname{E}^1, \seqname{C}^1)
    \biggr ]
  \subseteq
      \Functor^{\minimal{\M}}_2
        (\seqname{A}^2, \seqname{E}^2, \seqname{C}^2)
$.

\begin{proof}
$\funcseqname{f}$ is a morphism either by hypothesis or by
\pagecref{lem:M-ATLmorph}.
Let
$
  v^1 \in
  \Functor^{\minimal{\M}}_2
    (\seqname{A}^1, \seqname{E}^1, \seqname{C}^1)
$,
$(U^1,V^1,\phi^1) \in \seqname{A}^1$ be a chart with $v^1 \in V^1$,
$u^1 \defeq \phi^{1-1}(v^1)$, $u^2 \defeq \funcname{f}_0(u^1)$ and
$(U^2,V^2,\phi^2) \in \seqname{A}^2$ be a chart at $u^2$.
Let $(U'^1, V'^1, \phi^1 \maps U'^1 \toiso V'^1) \in \seqname{A}^1$,
$
  (U'^2 ,\hat{V}'^2, \phi'^2 \maps U'^2 \toiso \hat{V}'^2)
  \in \seqname{A}^2
$
be as in \pagecref{def:M-ATLmorph}. Then
$
  \funcname{f}_1(v^1) \in
  \hat{V}'^2 \subseteq
  \Functor^{\minimal{\M}}_2
    (\seqname{A}^1, \seqname{E}^1, \seqname{C}^1)
$.
\end{proof}
\end{lemma}

\begin{definition}[Model spaces associated with m-atlases]
\label{def:M-modE}
Let $\seqname{A}^i$, $i=1,2$, be an m-atlas of $\seqname{E}^i$ in the
coordinate space $\seqname{C}^i$ and
$\funcseqname{f} \defeq (\funcname{f}_0,\funcname{f}_1)$ be an
$\seqname{\seqname{E}}^1$-$\seqname{\seqname{E}}^2$ m-atlas (near)
morphism $(\funcname{f}_0,\funcname{f}_1)$ of $\seqname{A}^1$ to
$\seqname{A}^2$ in the coordinate spaces $\seqname{\seqname{C}}^1$,
Then:

The minimal model space with neighborhoods in the m-atlas
$\seqname{A}^i$ of $\seqname{E}^i$ in the coordinate space
$\seqname{C}^i$ is

\begin{multline}
\Functor^{\minimal{\M}}_1 (\seqname{A}^i, \seqname{E}^i, \seqname{C}^i)
\defeq \\
\minimal{\Mod}
  \left (
    \seqname{E}^i,
    \pi_1[\seqname{A}^i],
    \set
      {\phi'^{-1} \compose \phi}%
      [
        \equant
          {
            {(U,V,\phi) \in \seqname{A}^i},
            {(U',V',\phi') \in \seqname{A}^i}
          }
          {
            U \cap U' \ne \emptyset
          }
      ]
  \right )
\end{multline}

The mapping associated with the $\seqname{E}^1$-$\seqname{E}^2$
m-atlas morphism $(\funcname{f}_0,\funcname{f}_1)$ of
$\seqname{A}^1$ to $\seqname{A}^2$ in the coordinate spaces
$\seqname{C}^1$, $\seqname{C}^2$ is

\begin{multline}
\Functor^{\minimal{\M}}_1
  \bigl (
    (
      \funcname{f}_0,
      \funcname{f}_1
    ),
    (\seqname{A}^1, \seqname{E}^1, \seqname{C}^1),
    (\seqname{A}^2, \seqname{E}^2, \seqname{C}^2)
  \bigr )
\defeq \\
  \funcname{f}_0
  \maps
  \Functor^{\minimal{\M}}_1 (\seqname{A}^1, \seqname{E}^1, \seqname{C}^1)
  \to
  \Functor^{\minimal{\M}}_1 (\seqname{A}^2, \seqname{E}^2, \seqname{C}^2)
\end{multline}

If it is a model function then the it is also the m-atlas morphism
associated with the $\seqname{E}^1$-$\seqname{E}^2$ m-atlas morphism
$(\funcname{f}_0,\funcname{f}_1)$ of $\seqname{A}^1$ to $\seqname{A}^2$
in the coordinate spaces $\seqname{C}^1$, $\seqname{C}^2$.
\end{definition}

\begin{lemma}[Model spaces associated with m-atlases]
\label{lem:M-ismodE}
Let $\seqname{A}$ be an m-atlas of $\seqname{E}$ in the coordinate space
$\seqname{C}$. Then
$\Functor^{\minimal{\M}}_1 (\seqname{A}, \seqname{E}, \seqname{C})$ is a
model space.

\begin{proof}
The result follows from \Pagecref{lem:minmod}.
\end{proof}
\end{lemma}

\begin{theorem}[Functors from m-atlases to model spaces]
\label{the:M-isfunc}
Let $\seqname{E}$ and and $\seqname{C}$ be sets of model spaces.

$\Functor^{\minimal{\M}}_1$ is a functor from
$\Atl(\seqname{E}, \seqname{C})$ to
$
  \Trivcat[full-submod-]{\seqname{E}}  \subcat[full-]
  \Trivcat[submod-]{\seqname{E}}
$

$\Functor^{\minimal{\M}}_1$ is a functor from
$\full{\Atl}(\seqname{E}, \seqname{C})$ to
$
  \Trivcat[full-submod-]{\seqname{E}} \subcat[full-]
  \Trivcat[submod-]{\seqname{E}}
$

$\Functor^{\minimal{\M}}_2$ is a functor from
$\maximal[S-]{\Atl}(\seqname{E}, \seqname{C})$ to
$\Trivcat[submod-]{\seqname{C}}$.

$\Functor^{\minimal{\M}}_2$ is a functor from
$\full{\Atl}(\seqname{E}, \seqname{C})$ to
$
  \Trivcat[full-submod-]{\seqname{C}} \subcat[full-]
  \Trivcat[submod-]{\seqname{C}}
$

\begin{proof}
$
  \Trivcat[full-submod-]{\seqname{E}}  \subcat[full-]
  \Trivcat[submod-]{\seqname{E}}
$
and
$
  \Trivcat[full-submod-]{\seqname{C}} \subcat[full-]
  \Trivcat[submod-]{\seqname{C}}
$
by \pagecref{lem:trivmod}.

Let $o^i \defeq (\seqname{A}^i, E^i, C^i)$, $i \in [1,3]$, be
objects in $\Atl(\seqname{E}, \seqname{C})$
and let \\
$
m^i \defeq
  \bigl (
    (\funcname{f}^i_0, \funcname{f}^i_1),
    o^i,
    o^{i+1}
  \bigr )
$
be a morphism in $\Atl(\seqname{E}, \seqname{C})$.

$
  \Functor^{\minimal{\M}}_1 \maps
  \Atl(\seqname{E}, \seqname{C}) \to
  \Trivcat[submod-]{\seqname{E}}
$:
\begin{description}
\item[Preservation of endpoints:]
$
\Functor^{\minimal{\M}}_1(\funcname{m}^i) =
\funcname{f}^i_0 \maps \Functor^{\minimal{\M}}_1 o^i \to
\Functor^{\minimal{\M}}_1 o^{i+1}
$
\item[Composition:]
\begin{align*}
  \Functor^{\minimal{\M}}_1(\funcname{m}^2 \compose[A] \funcname{m}^1)
  & =
  \Functor^{\minimal{\M}}_1
    \bigl (
      (
        \funcname{f}^2_0 \compose \funcname{f}^1_0,
        \funcname{f}^2_1 \compose \funcname{f}^1_1
      )
      (\seqname{A}^1, E^1, C^1),
      (\seqname{A}^3, E^3, C^3)
    \bigr )
\\*
  & =
  \funcname{f}^2_0 \compose \funcname{f}^1_0 \maps
   \Functor^{\minimal{\M}}_1 o^1 \to
    \Functor^{\minimal{\M}}_1 o^3
\\*
  & =
  \bigl (
    \funcname{f}^2_0 \maps
    \Functor^{\minimal{\M}}_1 o^2 \to
    \Functor^{\minimal{\M}}_1 o^3
   \bigr )
  \compose
  \bigl (
    \funcname{f}^1_0 \maps
    \Functor^{\minimal{\M}}_1 o^1 \to
  \Functor^{\minimal{\M}}_1 o^2
 \bigr )
\\*
  & =
  \Functor^{\minimal{\M}}_1
    \bigl (
      (
        \funcname{f}^2_0,
        \funcname{f}^2_1
      )
      (\seqname{A}^2, E^2, C^2),
      (\seqname{A}^3, E^3, C^3)
    \bigr )
  \compose
\\*
  & \quad
  \Functor^{\minimal{\M}}_1
    \bigl (
      (
        \funcname{f}^1_0,
       \funcname{f}^1_1
      )
     (\seqname{A}^1, E^1, C^1),
      (\seqname{A}^2, E^2, C^2)
    \bigr )
\\*
  & =
  \Functor^{\minimal{\M}}_1(\funcname{m}^2)
  \compose
  \Functor^{\minimal{\M}}_1(\funcname{m}^1)
\end{align*}

\item[Identity:]
\leavevmode
\begin{enumerate}
\item
$
\Functor^{\minimal{\M}}_1(\Id[{o^i}]) =
\Functor^{\minimal{\M}}_1
  \bigl (
    (\Id[{E^i}], \Id[{C^i}]),
    (\seqname{A}^i, E^i, C^i),
    (\seqname{A}^i, E^i, C^i)
  \bigr )
=
\\
\Id[{E^i}] \maps \Functor^{\minimal{\M}}_1 o^i \to
\Functor^{\minimal{\M}}_1 o^i
$
\item
$
\Id_{\Functor^{\minimal{\M}}_1} o^i =
\Id[{E^i}] \maps \Functor^{\minimal{\M}}_1 o^i \to
\Functor^{\minimal{\M}}_1 o^i
$
\end{enumerate}
\end{description}

The proof for
$
  \Functor^{\minimal{\M}}_1 \maps
  \full{\Atl}(\seqname{E}, \seqname{C}) \to
  \Trivcat[full-submod-]{\seqname{E}} \subcat[full-]
  $
is identical.

$
  \Functor^{\minimal{\M}}_2 \maps
  \Atl(\seqname{E}, \seqname{C}) \to
  \optriv{\seqname{C}}
$:

\begin{description}
\item[Preservation of endpoints:]\quad\\*
$
\Functor^{\minimal{\M}}_2(\funcname{m}^i) =
\funcname{f}^i_1 \maps \Functor^{\minimal{\M}}_2 o^i \to
\Functor^{\minimal{\M}}_2 o^{i+1}
$
\item $\Functor (g \compose f) = \Functor (g) \compose \Functor (f)$:
$
\Functor^{\minimal{\M}}_2(\funcname{m}^2 \compose[A] \funcname{m}^1)
=
\\
\Functor^{\minimal{\M}}_2
  \bigl (
    (
      \funcname{f}^2_0 \compose \funcname{f}^1_0,
      \funcname{f}^2_1 \compose \funcname{f}^1_1
    )
    (\seqname{A}^1, E^1, C^1),
    (\seqname{A}^3, E^3, C^3)
  \bigr )
=
\\
\funcname{f}^2_0 \compose \funcname{f}^1_0 \maps
  \Functor^{\minimal{\M}}_2 o^1 \to
  \Functor^{\minimal{\M}}_2 o^3
=
\\
\bigl (
  \funcname{f}^2_1 \maps
  \Functor^{\minimal{\M}}_2 o^2 \to
  \Functor^{\minimal{\M}}_2 o^3
  \bigr )
\compose
\bigl (
  \funcname{f}^1_1 \maps
  \Functor^{\minimal{\M}}_2 o^1 \to
  \Functor^{\minimal{\M}}_2 o^2
\bigr )
=
\\
\Functor^{\minimal{\M}}_2
  \bigl (
    (
      \funcname{f}^2_0,
      \funcname{f}^2_1
    )
    (\seqname{A}^2, E^2, C^2),
    (\seqname{A}^3, E^3, C^3)
  \bigr )
\compose \\
\Functor^{\minimal{\M}}_2
  \bigl (
    (
      \funcname{f}^1_0,
      \funcname{f}^1_1
    )
    (\seqname{A}^1, E^1, C^1),
    (\seqname{A}^2, E^2, C^2)
  \bigr )
=
\\
\Functor^{\minimal{\M}}_2(\funcname{m}^2)
\compose
\Functor^{\minimal{\M}}_2(\funcname{m}^1)
$
\item $\Functor (\Id[{A}]) = \Id_{\Functor (A)}$:
\begin{enumerate}
\item
$
\Functor^{\minimal{\M}}_2(\Id[{o^i}]) =
\Functor^{\minimal{\M}}_2
  \bigl (
    (\Id[{E^i}], \Id[{C^i}]),
    (\seqname{A}^i, E^i, C^i),
    (\seqname{A}^i, E^i, C^i)
  \bigr )
=
\\
\Id[{C^i}] \maps \Functor^{\minimal{\M}}_2 o^i \to
\Functor^{\minimal{\M}}_2 o^i
$
\item
$
\Id_{\Functor^{\minimal{\M}}_2} o^i =
\Id[{C^i}] \maps \Functor^{\minimal{\M}}_2 o^i \to
\Functor^{\minimal{\M}}_2 o^i
$
\end{enumerate}
\end{description}

The proof for
$
  \Functor^{\minimal{\M}}_2 \maps
  \full{\Atl}(\seqname{E}, \seqname{C}) \to
  \Trivcat[full-submod-]{\seqname{C}} \subcat[full-]
$
is identical.
\end{proof}
\end{theorem}

\section{Classic m-atlas morphisms and functors}
\label{sec:M-ATLcmorph}
This subsection introduces an alternate definition of and taxonomy for
morphisms between m-atlases, defines categories of m-atlases, defines
functors amomg them and proves some basic reasults.  It introduces the
notion of classic m-atlas morphisms, which is very similar to
the conventional definitions for a manifold, although it does not
require the space to be locally Euclidean. The reader can safely skip
it, but it does help to establish a context for other notions that are
used in the exposition of local coordinate spaces.

While it is easy to define satisfactory notions of classic m-morphisms,
there are technical difficulties in defining a satisfactory notion of
classic near m-morphisms.

\subsection{Classic m-atlas morphisms}
\label{sub:M-ATLcmorph}
This subsection introduces the notion of classic m-atlas morphisms, and
proves some basic results.

\subsubsection{Definitions of classic m-atlas morphisms}
\begin{definition}[Classic m-atlas morphisms for model spaces]
\label{def:M-ATLcmorph}
Let $\catname{E}^i$, $\catname{C}^i$, $i=1,2$, be model categories,
$\seqname{E}^i \objin \catname{E}^i$,
$\seqname{C}^i \objin \catname{C}^i$
and $\seqname{A}^i$ be an m-atlas
of $\seqname{E}^i$ in the coordinate space $\seqname{C}^i$.

$\funcname{f} \maps \seqname{E}^1 \to \seqname{E}^2$ is a classic
\nobreaks
  {
    {$\seqname{E}^1$-$\seqname{E}^2$},
    {(-$\catname{C}^1$-$\catname{C}^2$)}
  }
m-atlas morphism of $\seqname{A}^1$ to $\seqname{A}^2$ in the coordinate
spaces $\seqname{C}^1$, $\seqname{C}^2$, abbreviated as
$
  \isAtl[classic]_\Ar
  (
    \seqname{A}^1,
    \seqname{E}^1,
    \seqname{C}^1,
    \seqname{A}^2,
    \seqname{E}^2,
    \seqname{C}^2,
    \funcname{f}
  )
$.
and a classic
\nobreaks
 {
   {$\catname{E}^1$-$\catname{E}^2$},
   {(-$\catname{C}^1$-$\catname{C}^2$)}
  }
m-atlas morphism of $\seqname{A}^1$ to $\seqname{A}^2$ in the coordinate
model categories $\catname{C}^1$, $\catname{C}^2$, abbreviated as
$
  \isAtl[classic]_\Ar
  (
    \seqname{A}^1,
    \catname{E}^1,
    \catname{C}^1,
    \seqname{A}^2,
    \catname{E}^2,
    \catname{C}^2,
    \funcname{f}
  )
$,
iff
\nobreaks
  {{
    $\seqname{C}^1 \submod \seqname{C}^2$
  }}
($\catname{C}^1 \subcat \catname{C}^2)$
and for any $(U^i, V^i, \phi^i \maps U^i \toiso V^i) \in \seqname{A}^i$,
$i=1,2$, with $I \defeq U^1 \cap \funcname{f}^{-1}[U^2] \neq \emptyset$,
$
  \phi^2 \compose \funcname{f} \compose \phi^{1-1} \maps
    \phi^1[I] \to V^2
$
is a morphism of $\seqname{C}^2$ ($\catname{C}^2$).

$\funcname{f}$ is also a constrained classic
\nobreaks
  {
    {$\seqname{E}^1$-$\seqname{E}^2$},
    {(-$\catname{C}^1$-$\catname{C}^2$)}
  }
m-atlas
morphism of $\seqname{A}^1$ to $\seqname{A}^2$ in the coordinate spaces
$\seqname{C}^1$, $\seqname{C}^2$, abbreviated as \\*
$
  \isAtl[constrained,classic]_\Ar
    (
      \seqname{A}^1,
      E^1,
      \seqname{C}^1,
      \seqname{A}^2,
      E^2,
      \seqname{C}^2,
      \funcname{f}_0,
      \funcname{f}_1
    )
$,
and a constrained classic
\nobreaks
  {
    {$\catname{E}^1$-$\catname{E}^2$},
    {(-$\catname{C}^1$-$\catname{C}^2$)}
  }
m-atlas morphism of
$\seqname{A}^1$ to $\seqname{A}^2$ in the coordinate model categories
$\catname{C}^1$, $\catname{C}^2$, abbreviated as
$
  \isAtl[constrained,classic]_\Ar
    (
      \seqname{A}^1,
      E^1,
      \catname{C}^1,
      \seqname{A}^2,
      E^2,
      \catname{C}^2,
      \funcname{f}_0,
      \funcname{f}_1
    )
$,
iff $\funcname{f}$ is constrained, i.e.,
$\funcname{f}[\seqname{E}]$ is contained in a model neighborhood of
$\seqname{E}^2$.

Let $\funcname{f}$ be a classic
\nobreaks
  {
    {$\seqname{E}^1$-$\seqname{E}^2$},
    {(-$\catname{C}^1$-$\catname{C}^2$)}
  }
m-atlas morphism of $\seqname{A}^1$ to $\seqname{A}^2$ in the
coordinate spaces $\seqname{C}^1$, $\seqname{C}^2$.

The triple
$
  \Bigl (
    (\funcname{f}),
    (\seqname{A}^1, \seqname{E}^1, \seqname{C}^1),
    (\seqname{A}^2, \seqname{E}^2, \seqname{C}^2)
  \Bigr )
$
will refer to $\funcname{f}$ considered as a classic
\nobreaks
  {
    {$\seqname{E}^1$-$\seqname{E}^2$},
    {(-$\catname{C}^1$-$\catname{C}^2$)}
  }
and
\nobreaks
  {
    {$\catname{E}^1$-$\catname{E}^2$},
    {(-$\catname{C}^1$-$\catname{C}^2$)}
  }
m-atlas morphism of $\seqname{A}^1$ to $\seqname{A}^2$ in the coordinate
spaces $\seqname{C}^1$, $\seqname{C}^2$.
\end{definition}

\begin{definition}[Classic m-atlas morphisms for topological spaces]
\label{def:M-ATLcmorphtop}
Let $\catname{C}^i$, $i=1,2$, be a model category, $\catname{E}^i$ be a
topological category, $E^i \objin \catname{E}^i$,
$\seqname{C}^i \objin \catname{C}^i$ and $\seqname{A}^i$ be an m-atlas
of $E^i$ in the coordinate space $\seqname{C}^i$.

$\funcname{f} \maps E^1 \to E^2$ is a classic
\nobreaks
  {
    {$E^1$-$E^2$},
    {(-$\catname{C}^1$-$\catname{C}^2$)}
  }
m-atlas morphism of $\seqname{A}^1$ to $\seqname{A}^2$ in the coordinate
spaces $\seqname{C}^1$, $\seqname{C}^2$, abbreviated as
$
  \isAtl[classic]_\Ar
  (
    \seqname{A}^1,
    E^1,
    \seqname{C}^1,
    \seqname{A}^2,
    E^2,
    \seqname{C}^2,
    \funcname{f}
  )
$.
and a classic
\nobreaks
  {
    {$\catname{E}^1$-$\catname{E}^2$},
    {(-$\catname{C}^1$-$\catname{C}^2$)}
  }
m-atlas morphism of $\seqname{A}^1$ to $\seqname{A}^2$ in the coordinate
model categories $\catname{C}^1$, $\catname{C}^2$, abbreviated as
$
  \isAtl[classic]_\Ar
  (
    \seqname{A}^1,
    \catname{E}^1,
    \catname{C}^1,
    \seqname{A}^2,
    \catname{E}^2,
    \catname{C}^2,
    \funcname{f}
  )
$,
iff $\funcname{f}$ is a classic
\nobreaks
  {
    {$\Triv{E^1}$-$\Triv{E^2}$},
    {(-$\catname{C}^1$-$\catname{C}^2$)}
  }
m-atlas
morphism of $\seqname{A}^1$ to $\seqname{A}^2$ in the coordinate spaces
$\seqname{C}^1$, $\seqname{C}^2$.

The triple
$
  \Bigl (
    (\funcname{f}),
    (\seqname{A}^1, E^1, \seqname{C}^1),
    (\seqname{A}^2, E^2, \seqname{C}^2)
  \Bigr )
$
will refer to $\funcname{f}$ considered as a classic
\nobreaks
  {
    {$E^1$-$E^2$}%
    {(-$\catname{C}^1$-$\catname{C}^2$)}
  }
and
\nobreaks
  {
    {$\catname{E}^1$-$\catname{E}^2$}%
    {(-$\catname{C}^1$-$\catname{C}^2$)}
  }
m-atlas morphism of $\seqname{A}^1$ to $\seqname{A}^2$ in the coordinate
spaces $\seqname{C}^1$, $\seqname{C}^2$.
\end{definition}

\begin{definition}[Classic m-atlas inclusion and identity morphisms]
\label{def:M-ATLcid}
Let
\begin{enumerate}
\item
$\catname{E}^i$, $\catname{C}^i$, $i=1,2$, be model categories
\item
$\seqname{E}^i \objin \catname{E}^i$
\item
$\seqname{E}^1 \submod \seqname{E}^2$
\item
$\seqname{C}^i \objin \catname{C}^i$
\item
\nobreaks
  {
    {$\seqname{C}^1 \submod \seqname{C}^2$},
    {($\catname{C}^1 \subcat \catname{C}^2$)}
  }
\item
$\seqname{A}^i$ be an m-atlas of $\seqname{E}^i$ in the coordinate
space $\seqname{C}^i$
\item
$\seqname{A}^1 \subseteq \seqname{A}^2$
\end{enumerate}

The classic
\nobreaks
  {
    {$\seqname{E}^1$-$\seqname{E}^2$},
    {(-$\catname{C}^1$-$\catname{C}^2$)}
  }
m-atlas inclusion morphism of $\seqname{A}^1$ to $\seqname{A}^2$ in the
coordinate spaces $\seqname{C}^1$, $\seqname{C}^2$ is
$
  \Id%
    [
      {(\seqname{E}^1, \seqname{C}^1)},
      {(\seqname{E}^2, \seqname{C}^2)}
    ]
  \defeq
  \Id%
    [
      {\seqname{E}^1},
      {\seqname{E}^2}
    ]
$.

The classic m-atlas inclusion morphism of
$(\seqname{A}^1, \seqname{E}^1, \seqname{C}^1)$ to
$(\seqname{A}^2, \seqname{E}^2, \seqname{C}^2)$ is \\*
$
  \Id%
    [
      {(\seqname{A}^1, \seqname{E}^1, \seqname{C}^1)},
      {(\seqname{A}^2, \seqname{E}^2, \seqname{C}^2)}
    ]
  \defeq
  \biggl (
    \ID%
    [
      {(\seqname{E}^1)},
      {(\seqname{E}^2)}
    ],
    (\seqname{A}^1, \seqname{E}^1, \seqname{C}^1),
    (\seqname{A}^2, \seqname{E}^2, \seqname{C}^2)
  \biggr )
$,
$
 \Id%
   [
     {\seqname{E}^1},
     {\seqname{E}^2}
   ]
$
considered as a classic
\nobreaks
  {
    {$\seqname{E}^1$-$\seqname{E}^2$},
    {(-$\catname{C}^1$-$\catname{C}^2$)}
  }
m-atlas morphism of $\seqname{A}^1$ to $\seqname{A}^2$ in the coordinate
spaces $\seqname{C}^1$, $\seqname{C}^2$.

The classic
\nobreaks
  {
    {$\seqname{E}^i$},
    {(-$\catname{C}^i$)}
  }
m-atlas identity morphism of $\seqname{A}^i$
in the coordinate space $\seqname{C}^i$ is
$
  \Id[{{(\seqname{E}^i, \seqname{C}^i)}}] \defeq
  \Id%
    [
      {\seqname{E}^i}
    ]
$.

The classic m-atlas identity morphism of
$(\seqname{A}^i, \seqname{E}^i, \seqname{C}^i)$ is \\*
$
  \Id[{{(\seqname{A}^i, \seqname{E}^i, \seqname{C}^i)}}] \defeq
  \bigl (
    \ID[{{(\seqname{E}^i)}}],
    (\seqname{A}^i, \seqname{E}^i, \seqname{C}^i),
    (\seqname{A}^i, \seqname{E}^i, \seqname{C}^i)
  \bigr )
$,\hspace{1pt}
$
 \Id%
   [
     {\seqname{E}^i}
   ]
$
considered as a classic
\nobreaks
  {
    {$\seqname{E}^i$},
    {(-$\catname{C}^i$)}
  }
m-atlas morphism of
$\seqname{A}^i$ in the coordinate space $\seqname{C}^i$.

Let
\begin{enumerate}
\item
$E^i$ be a topological space, $i=1,2$,
\item
$E^1 \subsp E^2$
\item
$\catname{C}^i$ be a model category
\item
$\seqname{C}^i \objin \catname{C}^i$
\item
$\catname{C}^1 \subcat \catname{C}^2$
\item
$\seqname{A}^i$ be an m-atlas of $E^i$ in the coordinate space
$\seqname{C}^i$
\item
$\seqname{A}^1 \subseteq \seqname{A}^2$
\end{enumerate}

The classic
\nobreaks
  {
    {$E^1$-$E^2$},
    {(-$\catname{C}^1$-$\catname{C}^2$)}
  }
m-atlas inclusion morphism of
$\seqname{A}^1$ to $\seqname{A}^2$ in the coordinate spaces
$\seqname{C}^1$, $\seqname{C}^2$ is
$
 \Id%
   [
     {E^1},
     {E^2}
   ]
$.

The classic m-atlas inclusion morphism of
$(\seqname{A}^1, E^1, \seqname{C}^1)$ to
$(\seqname{A}^2, E^2, \seqname{C}^2)$ is \\*
$
  \Id%
    [
        {(\seqname{A}^1, E^1, \seqname{C}^1)},
        {(\seqname{A}^2, E^2, \seqname{C}^2)}
    ]
  \defeq
  \biggl (
    \ID%
    [
      {(E^1)},
      {(E^2)}
    ],
    (\seqname{A}^1, E^1, \seqname{C}^1),
    (\seqname{A}^2, E^2, \seqname{C}^2)
  \biggr )
$,
$
 \Id%
   [
     {E^1},
     {E^2}
   ]
$
considered as a classic $E^1$-$E^2$ m-atlas morphism of
$\seqname{A}^1$ to $\seqname{A}^2$ in the coordinate spaces
$\seqname{C}^1$, $\seqname{C}^2$.

The classic $E^i$ m-atlas identity morphism of $\seqname{A}^i$ in the
coordinate space $\seqname{C}^i$ is
$
 \Id%
   [
     {(E^i, \seqname{C}^i)}
   ]
$.

The classic m-atlas identity morphism of $(\seqname{A}^i, E^i, \seqname{C}^i)$ is \\*
$
  \Id[{{(\seqname{A}^i, E^i, \seqname{C}^i)}}] \defeq
  \bigl (
    \ID[{{(E^i)}}],
    (\seqname{A}^i, E^i, \seqname{C}^i),
    (\seqname{A}^i, E^i, \seqname{C}^i)
  \bigr )
$,
$
 \Id%
   [
     {E^i}
   ]
$
considered as an $E^i$ m-atlas morphism of $\seqname{A}^i$ in
the coordinate space $\seqname{C}^i$.
\end{definition}

\subsubsection{Definitions of semi-strict and strict}
\begin{definition}[Semi-strict classic m-atlas morphisms]
\label{def:M-ATLcmorphsemi}
Let $\catname{E}^i$, $\catname{C}^i$, $i=1,2$, be model categories,
$\seqname{E}^i \objin \catname{E}^i$,
$\seqname{C}^i \objin \catname{C}^i$ and $\seqname{A}^i$ be an m-atlas
of $\seqname{E}^i$ in the coordinate space $\seqname{C}^i$.

Let $\funcname{f} \maps \seqname{E}^1 \to \seqname{E}^2$ be a classic
$\seqname{E}^1$-$\seqname{E}^2$ m-atlas morphism of
$\seqname{A}^1$ to $\seqname{A}^2$ in the coordinate spaces
$\seqname{C}^1$, $\seqname{C}^2$.

$\funcname{f}$ is a semi-strict classic $\seqname{E}^1$-$\seqname{E}^2$
m-atlas morphism of $\seqname{A}^1$ to $\seqname{A}^2$ in the
coordinate spaces $\seqname{C}^1$, $\seqname{C}^2$, abbreviated as
$
  \strict[semi-]{\isAtl[classic]_\Ar}
  (
    \seqname{A}^1,
    \seqname{E}^1,
    \seqname{C}^1,
    \seqname{A}^2,
    \seqname{E}^2,
    \seqname{C}^2,
    \funcname{f}
  )
$,
iff $\funcname{f}$ is a local m-morphism of $\seqname{E}^1$ to
$\seqname{E}^2$.

\begin{remark}
Definitions of semi-strict classic m-atlas morphisms for
topological spaces would be pointless, as any semi-strict classic
m-atlas morphism would be strict.
\end{remark}

It is a semi-strict classic
$\catname{E}^1$-$\catname{E}^2$-$\seqname{E}^1$-$\seqname{E}^2$ m-atlas
morphism of $\seqname{A}^1$ to $\seqname{A}^2$ in the coordinate
spaces $\seqname{C}^1$, $\seqname{C}^2$, abbreviated as \\*
$
  \strict[semi-]{\isAtl[classic]_\Ar}
  (
    \seqname{A}^1,
    \seqname{E}^1,
    \seqname{C}^1,
    \seqname{A}^2,
    \seqname{E}^2,
    \seqname{C}^2,
    \funcname{f},
    \catname{E}^1,
    \catname{E}^2
  )
$,
iff $\funcname{f}$ is a local $\catname{E}^1$-$\catname{E}^2$
m-morphism of $\seqname{E}^1$ to $\seqname{E}^2$.

Let $\funcname{f} \maps \seqname{E}^1 \to \seqname{E}^2$ be a classic
$\seqname{E}^1$-$\seqname{E}^2$-$\catname{C}^1$-$\catname{C}^2$ m-atlas
morphism of $\seqname{A}^1$ to $\seqname{A}^2$ in the coordinate
spaces $\seqname{C}^1$, $\seqname{C}^2$.

$\funcname{f}$ is a semi-strict classic
$\seqname{E}^1$-$\seqname{E}^2$-$\catname{C}^1$-$\catname{C}^2$
m-atlas morphism of $\seqname{A}^1$ to $\seqname{A}^2$ in the
coordinate spaces $\seqname{C}^1$, $\seqname{C}^2$, abbreviated as \\*
$
  \strict[semi-]{\isAtl[classic]_\Ar}
  (
    \seqname{A}^1,
    \seqname{E}^1,
    \seqname{C}^1,
    \seqname{A}^2,
    \seqname{E}^2,
    \seqname{C}^2,
    \funcname{f},
    \catname{C}^1,
    \catname{C}^2
  )
$,
iff $\funcname{f}$ is a local m-morphism of $\seqname{E}^1$ to
$\seqname{E}^2$.

It is a semi-strict classic
$\catname{E}^1$-$\catname{E}^2$-%
$\seqname{E}^1$-$\seqname{E}^2$-%
$\catname{C}^1$-$\catname{C}^2$
m-atlas morphism of $\seqname{A}^1$ to $\seqname{A}^2$ in the coordinate
spaces $\seqname{C}^1$, $\seqname{C}^2$, abbreviated as \\*
$
  \strict[semi-]{\isAtl[classic]_\Ar}
  (
    \seqname{A}^1,
    \seqname{E}^1,
    \seqname{C}^1,
    \seqname{A}^2,
    \seqname{E}^2,
    \seqname{C}^2,
    \funcname{f},
    \catname{E}^1,
    \catname{E}^2,
    \catname{C}^1,
    \catname{C}^2
  )
$,
iff $\funcname{f}$ is a local $\catname{E}^1$-$\catname{E}^2$
m-morphism of $\seqname{E}^1$ to $\seqname{E}^2$.
\end{definition}

\begin{definition}[Strict classic m-atlas morphisms]
\label{def:M-ATLcmorphstrict}
Let $\catname{E}^i$, $\catname{C}^i$, $i=1,2$, be model categories,
$\seqname{E}^i \objin \catname{E}^i$,
$\seqname{C}^i \objin \catname{C}^i$, $\seqname{A}^i$ be an m-atlas of
$\seqname{E}^i$ in the coordinate space $\seqname{C}^i$.

Let $\funcname{f} \maps \seqname{E}^1 \to \seqname{E}^2$ be a classic
$\seqname{E}^1$-$\seqname{E}^2$ m-atlas morphism of
$\seqname{A}^1$ to $\seqname{A}^2$ in the coordinate spaces
$\seqname{C}^1$, $\seqname{C}^2$.

$\funcname{f}$ is a strict classic $\seqname{E}^1$-$\seqname{E}^2$
m-atlas morphism of $\seqname{A}^1$ to $\seqname{A}^2$ in the
coordinate spaces $\seqname{C}^1$, $\seqname{C}^2$, abbreviated as
$
  \strict{\isAtl[classic]_\Ar}
  (
    \seqname{A}^1,
    \seqname{E}^1,
    \seqname{C}^1,
    \seqname{A}^2,
    \seqname{E}^2,
    \seqname{C}^2,
    \funcname{f}
  )
$,
iff $\funcname{f}$ is an m-morphism of $\seqname{E}^1$ to $\seqname{E}^2$.

It is a strict classic
$\catname{E}^1$-$\catname{E}^2$-$\seqname{E}^1$-$\seqname{E}^2$ m-atlas
morphism of $\seqname{A}^1$ to $\seqname{A}^2$ in the coordinate
spaces $\seqname{C}^1$, $\seqname{C}^2$, abbreviated as \\*
$
  \strict{\isAtl[classic]_\Ar}
  (
    \seqname{A}^1,
    \seqname{E}^1,
    \seqname{C}^1,
    \seqname{A}^2,
    \seqname{E}^2,
    \seqname{C}^2,
    \funcname{f},
    \catname{E}^1,
    \catname{E}^2
  )
$,
iff $\funcname{f}$ is an $\catname{E}^1$-$\catname{E}^2$ m-morphism of
$\seqname{E}^1$ to $\seqname{E}^2$.

Let $\funcname{f} \maps \seqname{E}^1 \to \seqname{E}^2$ be a classic
$\seqname{E}^1$-$\seqname{E}^2$-$\catname{C}^1$-$\catname{C}^2$ m-atlas
morphism of $\seqname{A}^1$ to $\seqname{A}^2$ in the coordinate
spaces $\seqname{C}^1$, $\seqname{C}^2$.

$\funcname{f}$ is a strict classic
$\seqname{E}^1$-$\seqname{E}^2$-$\catname{C}^1$-$\catname{C}^2$
m-atlas morphism of $\seqname{A}^1$ to $\seqname{A}^2$ in the coordinate
spaces $\seqname{C}^1$, $\seqname{C}^2$, abbreviated as
$
  \strict{\isAtl[classic]_\Ar}
  (
    \seqname{A}^1,
    \seqname{E}^1,
    \seqname{C}^1,
    \seqname{A}^2,
    \seqname{E}^2,
    \seqname{C}^2,
    \funcname{f},
    \catname{C}^1,
    \catname{C}^2
  )
$,
iff $\funcname{f}$ is an m-morphism of $\seqname{E}^1$ to $\seqname{E}^2$.

It is a strict classic
$\catname{E}^1$-$\catname{E}^2$-%
$\seqname{E}^1$-$\seqname{E}^2$-%
$\catname{C}^1$-$\catname{C}^2$
m-atlas morphism of $\seqname{A}^1$ to $\seqname{A}^2$ in the coordinate
spaces $\seqname{C}^1$, $\seqname{C}^2$, abbreviated as \\*
$
  \strict{\isAtl[classic]_\Ar}
  (
    \seqname{A}^1,
    \seqname{E}^1,
    \seqname{C}^1,
    \seqname{A}^2,
    \seqname{E}^2,
    \seqname{C}^2,
    \funcname{f},
    \catname{E}^1,
    \catname{E}^2,
    \catname{C}^1,
    \catname{C}^2
  )
$,
iff $\funcname{f}$ is an $\catname{E}^1$-$\catname{E}^2$
m-morphism of $\seqname{E}^1$ to $\seqname{E}^2$.

\end{definition}

\subsubsection{Abbreviated nomenclature for classic m-atlas morphisms}
\begin{definition}[Abbreviated nomenclature for classic m-atlas morphisms]
\label{def:M-ATLcmorphabbrev}
Let $\catname{E}^i$, $\catname{C}^i$, $i=1,2$, be model categories,
$\seqname{E}^i \objin \catname{E}^i$,
$\seqname{C}^i \objin \catname{C}^i$, $\seqname{A}^i$ be an m-atlas of
$\seqname{E}^i$ in the coordinate space $\seqname{C}^i$ and
$\funcname{f} \maps \seqname{E}^1 \to \seqname{E}^2$ be a classic
$\seqname{E}^1$-$\seqname{E}^2$ m-atlas morphism of
$\seqname{A}^1$ to $\seqname{A}^2$ in the coordinate spaces
$\seqname{C}^1$,
$\seqname{C}^2$.

$\funcname{f}$ is a (full) (maximal) (constrained) (semi-strict, strict)
classic $\seqname{E}^1$-$\seqname{E}^2$ m-atlas morphism of
$\seqname{A}^1$ to $\seqname{A}^2$ in the coordinate space
$\seqname{C}^1$, abbreviated as \\*
$
  \strict[*]{\maxfull[*]{\isAtl[(constrained,)classic]_\Ar}}
  (
    \seqname{A}^1,
    \seqname{E}^1,
    \seqname{C}^1,
    \seqname{A}^2,
    \seqname{E}^2,
    \funcname{f}
  )
$,
iff it is a (full) (maximal) (constrained) (semi-strict, strict) classic
$\seqname{E}^1$-$\seqname{E}^2$ m-atlas morphism of
$\seqname{A}^1$ to $\seqname{A}^2$ in the coordinate spaces
$\seqname{C}^1$, $\seqname{C}^1$.

$\funcname{f}$ is a (full) (maximal) (constrained) (semi-strict, strict)
classic
$\catname{E}^1$-$\catname{E}^2$-$\seqname{E}^1$-$\seqname{E}^2$
m-atlas morphism of $\seqname{A}^1$ to $\seqname{A}^2$ in the
coordinate space $\seqname{C}^1$, abbreviated as \\*
$
  \strict[*]{\maxfull[*]{\isAtl[(constrained,)classic]_\Ar}}
  (
    \seqname{A}^1,
    \seqname{E}^1,
    \seqname{C}^1,
    \seqname{A}^2,
    \seqname{E}^2,
    \funcname{f},
    \catname{E}^1,
    \catname{E}^2
  )
$,
iff it is a (full) (maximal) (constrained) (semi-strict, strict) classic
$\catname{E}^1$-$\catname{E}^2$-$\seqname{E}^1$-$\seqname{E}^2$
m-atlas morphism of $\seqname{A}^1$ to $\seqname{A}^2$ in the
coordinate spaces $\seqname{C}^1$, $\seqname{C}^1$.

Let $\catname{C}^i$, $i=1,2$, be a model category,
$\seqname{C}^i \objin \catname{C}^i$,
$E^i$ be a topological space,
$\seqname{A}^i$ be an m-atlas of $E^i$ in the coordinate space
$\seqname{C}^i$ and $\funcname{f} \maps E^1 \to E^2$ be a classic
$E^1$-$E^2$ m-atlas morphism of $\seqname{A}^1$ to
$\seqname{A}^2$ in the coordinate spaces $\seqname{C}^1$, $\seqname{C}^2$.

$\funcname{f}$ is a (full) (maximal) (constrained) (semi-strict, strict)
classic $E^1$-$E^2$ m-atlas morphism
of $\seqname{A}^1$ to $\seqname{A}^2$ in the coordinate spaces
$\seqname{C}^1$, $\seqname{C}^2$, abbreviated as \\*
$
  \strict[*]{\maxfull[*]{\isAtl[(constrained,)classic]_\Ar}}
  (
    \seqname{A}^1,
    E^1,
    \seqname{C}^1,
    \seqname{A}^2,
    E^2,
    \seqname{C}^2,
    \funcname{f}
  )
$,
iff it is a (full) (maximal) (constrained) (semi-strict, strict) classic
$\Triv{E}^1$-$\Triv{E}^2$ m-atlas
morphism of $\seqname{A}^1$ to $\seqname{A}^2$ in the coordinate spaces
$\seqname{C}^1$, $\seqname{C}^2$.

$\funcname{f}$ is a (full) (maximal) (constrained) (semi-strict, strict)
classic $E^1$-$E^2$ m-atlas morphism of $\seqname{A}^1$ to
$\seqname{A}^2$ in the coordinate space $\seqname{C}^1$, abbreviated
as \\*
$
  \strict[*]{\maxfull[*]{\isAtl[(constrained,)classic]_\Ar}}
  (
    \seqname{A}^1,
    E^1,
    \seqname{C}^1,
    \seqname{A}^2,
    E^2,
    \funcname{f}
  )
$,
iff it is a (full) (maximal) (constrained) (semi-strict, strict) classic
$E^1$-$E^2$ m-atlas morphism of $\seqname{A}^1$ to $\seqname{A}^2$ in
the coordinate spaces $\seqname{C}^1$, $\seqname{C}^1$.
\end{definition}

\subsubsection{Proclamations on classic m-atlas morphisms}
\begin{lemma}[Classic m-atlas morphisms]
\label{lem:M-ATLcmorph}
Let $\seqname{E}$ be a set of topological spaces,
$E^i \in \seqname{E}$, $i=1,2$,
$\seqname{C}^i$ be a model space,
$\seqname{A}^i$ be an m-atlas of $\Triv{E^i}$ in the coordinate space
$\seqname{C}^i$.

\begin{enumerate}
\item
Let $\funcname{f} \maps \Triv{E^1} \to \Triv{E^2}$ be a classic
$\Triv{E^1}$-$\Triv{E^2}$ m-atlas morphism of $\seqname{A}^1$ to
$\seqname{A}^2$ in the coordinate spaces $\seqname{C}^1$,
$\seqname{C}^2$ and $\seqname{C}^1 \submod \seqname{C}^2$. Then
$\funcname{f}$ is a strict constrained classic
\nobreaks
  {
    {$\Trivcat{\seqname{E}}$-$\Trivcat{\seqname{E}}$-},
    {$\Triv{E^1}$-$\Triv{E^2}$}
  }
m-atlas morphism of $\seqname{A}^1$ to $\seqname{A}^2$ in the coordinate
spaces $\seqname{C}^1$, $\seqname{C}^2$.

\begin{proof}
\leavevmode
\begin{description}
\item[$\funcname{f}$ is constrained:]
\leavevmode
$\funcname{f}[E^1] \subseteq E^2$ by hypothesis.
$E^2$ is a model neighborhood of $\Triv{E^2}$ by \pagecref{def:trivmod}.
\item[$\funcname{f}$ is strict:]
\leavevmode
$\Trivcat{\seqname{E}} \subcat[full-] \Trivcat{\seqname{E}}$.
$\funcname{f}$ is continuous by \pagecref{def:modfunc}.
$E^i \in \seqname{E}$ by hypothesis,
$E^i \objin  \Trivcat{\seqname{E}}$ and
$\funcname{f} \arin \Trivcat{\seqname{E}}$ by \cref{def:trivmod}.
$\seqname{C}^1 \submod \seqname{C}^2$
by \pagecref{def:M-ATLcmorph}.
\end{description}
\end{proof}

\item
If $E^1$ is an open subspace of $E^2$ then
$\funcname{f} \maps \Triv{E^1} \to \Triv{E^2}$ is also a strict
constrained classic $\Triv{E^1}$-$\Triv{E^2}$ m-atlas morphism of
$\seqname{A}^1$ to $\seqname{A}^2$ in the coordinate spaces
$\seqname{C}^1$, $\seqname{C}^2$.

\begin{proof}
Each $E^i \in \seqname{E}$ by hypothesis and each
$\Triv{E^i} \objin \Trivcat{\seqname{E}}$ by \cref{def:trivmod}.

$E^1 \subsp[open-] E^2$ by hypothesis.

$\Triv{E^1} \submod \Triv{E^2}$ by \pagecref{lem:trivmod}.

$\funcname{f} \arin \Trivcat{\seqname{E}}$ by \cref{def:trivmod}.
\end{proof}
\end{enumerate}

Let
\begin{enumerate}
\item
$\catname{E}^i$, $\catname{C}^i$,
$\catname{E}'^i$, $\catname{C}'^i$, $i=1,2$, be model categories
\item
$\catname{E}^i \subcat[full-] \catname{E}'^i$
\item
$\catname{C}^i \subcat[full-] \catname{C}'^i$
\item
$\catname{C}'^1 \subcat[full-] \catname{C}'^2$
\item
$\seqname{E}^i \objin \catname{E}^i$
\item
$\seqname{C}^i \objin \catname{C}^i$
\item
$\seqname{C}'^i \objin \catname{C}'^i$
\item
$\seqname{C}^i \submod \seqname{C}'^i$
\item
$\seqname{A}^i$ be an m-atlas of $\seqname{E}^i$ in the coordinate space
$\seqname{C}^i$
\item
$\funcname{f} \maps \seqname{E}^1 \to \seqname{E}^2$ be a
(semi-strict, strict)
\nobreaks
  {
    {($\catname{E}^1$-$\catname{E}^2$-)},
    {$\seqname{E}^1$-$\seqname{E}^2$},
    {(-$\catname{C}^1$-$\catname{C}^2$)}
  }
m-atlas morphism of $\seqname{A}^1$ to $\seqname{A}^2$ in the
coordinate spaces $\seqname{C}^1$, $\seqname{C}^2$
(model function?)
\end{enumerate}

If $\funcname{f}$ is a (semi-strict, strict)
$\catname{E}^1$-%
$\catname{E}^2$-%
$\seqname{E}^1$-%
$\seqname{E}^2$
(%
$\seqname{E}^1$-%
$\seqname{E}^2$%
)
classic m-atlas morphism of
$\seqname{A}^1$ to $\seqname{A}^2$ in the coordinate spaces
$\seqname{C}^1$, $\seqname{C}^2$, then $\funcname{f}$ is a
(semi-strict, strict)
$\catname{E}'^1$-%
$\catname{E}'^2$-%
$\seqname{E}^1$-%
$\seqname{E}^2$
(
$\seqname{E}^1$-%
$\seqname{E}^2$%
)
classic m-atlas morphism of $\seqname{A}^1$ to $\seqname{A}^2$ in the
coordinate spaces $\seqname{C}'^1$, $\seqname{C}'^2$.

\begin{proof}
\leavevmode
\begin{description}
\item[$\funcname{f}$ remains a classic m-atlas morphism with the expanded categories.]
$\catname{C}'^1 \subcat[full-] \catname{C}'^2$ by hypothesis.

Let $(U^i, V^i, \phi^i \maps U^i \toiso V^i) \in \seqname{A}^i$, $i=1,2$,
with $I \defeq U^1 \cap \funcname{f}^{-1}[U^2] \neq \emptyset$.
$\phi^2 \compose \funcname{f} \compose \phi^{-1}$ is a morphism of
$\catname{C}^2 \subcat[full-] \catname{C}'^2$ by \cref{def:M-ATLcmorph}.
\item[$\funcname{f}$ remains semi-strict (strict).]
If $\funcname{f}$ is a local morphism of $\catname{E}^2$ then
$\funcname{f}$ is a local morphism of $\catname{E}'^2$.    If
$\funcname{f}$ is a local morphism of $\seqname{E}^2$ then
$\funcname{f}$ is a local morphism of $\seqname{E}'^2$.    If
$\funcname{f} \arin \catname{E}^2$ then
$\funcname{f} \arin \catname{E}'^2$.  If $\funcname{f}$ is a
morphism of $\seqname{E}^2$ then $\funcname{f}$ is a morphism of
$\seqname{E}'^2$.
\end{description}
\end{proof}
\end{lemma}

\begin{lemma}[Composition of classic m-atlas morphisms]
\label{lem:M-ATLcmorphcomp}
Let $\catname{E}^i$, $\catname{C}^i$, $i \in [1,3]$, be a model
category,
$\seqname{E}^i \objin \catname{E}^i$,
$\seqname{C}^i \objin \catname{C}^i$, $\seqname{A}^i$ be an m-atlas
of $\seqname{E}^i$ in the coordinate space $\seqname{C}^i$ and
$\funcname{f}^i \maps \seqname{E}^i \to \seqname{E}^{i+1}$, $i=1,2$, be
a (semi-strict, strict) classic $\seqname{E}^i$-$\seqname{E}^{i+1}$
m-atlas morphism of $\seqname{A}^i$ to $\seqname{A}^{i+1}$ in the
coordinate spaces $\seqname{C}^i$, $\seqname{C}^{i+1}$.

$
  \funcname{f}^2 \compose \funcname{f}^1
  \maps \seqname{E}^1 \to \seqname{E}^3
$
is a (semi-strict, strict) classic $\seqname{E}^1$-$\seqname{E}^3$
m-atlas morphism of $\seqname{A}^1$ to $\seqname{A}^3$ in the coordinate
spaces $\seqname{C}^1$, $\seqname{C}^3$.

\begin{proof}
\leavevmode
\begin{description}
\item[\ensuremath{\funcname{f}^2 \compose \funcname{f}^1} is a classic $\seqname{E}^1$-$\seqname{E}^3$ m-atlas morphism.]
If $\catname{C}^1 \subcat \catname{C}^2$ and
$\catname{C}^2 \subcat \catname{C}^3$ then
$\catname{C}^1 \subcat \catname{C}^3$.
Let
$(U^1, V^1, \phi^1 \maps U^1 \toiso V^1) \in \seqname{A}^1$,
$(U^3 ,V^3, \phi^3 \maps U^3 \toiso V^3) \in \seqname{A}^3$
with
$
  I \defeq
  U^1 \cap (\funcname{f}^2 \compose \funcname{f}^1)^{-1}[U^3]
  \neq \emptyset
$,
$(\funcname{f}^2 \compose \funcname{f}^1)^{-1}[U^3]$ is a model
neighborhood by \pagecref{def:modfunc} and
$I$ is a model neighborhood by \cref{mod:closed} of
{
  \showlabelsinline
  \pagecref{def:model}
}.

For any $u^1 \in I$ and any chart
$(U^2 ,V^2, \phi^2 \maps U^2 \toiso V^2) \in \seqname{A}^2$ at
$\funcname{f}^1(u^1)$, define
$I^{1,2} \defeq U^1 \cap \funcname{f}^{1-1}[U^2]$,
$I^{2,3} \defeq U^2 \cap \funcname{f}^{2-1}[U^3]$ and
$I^{1,3}_{u^1} \defeq U^1 \cap \funcname{f}^{1-1}[I^{2,3}] \subseteq I$.
$I^{1,3}_{u^1}$ is a model neighborhood by
{
  \showlabelsinline
  \cref{def:modfunc}
}
 and
\cref{mod:closed} of
{
  \showlabelsinline
  \cref{def:model}
}.

$
  \phi^2 \compose \funcname{f}^1 \compose \phi^{1-1} \maps
    \phi^1[I^{1,2}] \to V^2
$
is a morphism of $\catname{C}^2 \subcat \catname{C}^3$ by hypothesis and
$
  \phi^2 \compose \funcname{f}^1 \compose \phi^{1-1} \maps
    \phi^1[I^{1,3}_{u^1}] \to V^2
$
is a morphism of $\catname{C}^2 \subcat \catname{C}^3$ by
{
  \showlabelsinline
  \cref{mod:restriction}.
}
of \cref{def:model}.
$
  \phi^3 \compose \funcname{f}^2 \compose \phi^{2-1} \maps
    \phi^2[I^{2,3}] \to V^3
$
is a morphism of $\catname{C}^3$ by hypothesis and
and thus
$
  (
    \phi^3 \compose
    \funcname{f}^2 \compose
    \funcname{f}^1 \compose
    \phi^{1-1}
    \maps \phi^1[I^{1,3}_{u^1}] \to V^3
  )
  =
  (
    \phi^3 \compose \funcname{f}^2 \compose \phi^{2-1} \maps
      \phi^2[I^{2,3}] \to V^3
  )
  \compose
  (
    \phi^2 \compose \funcname{f}^1 \compose \phi^{1-1} \maps
      \phi^1[I^{1,3}_{u^1}] \to V^2
  )
$
is a morphism of $\catname{C}^3$.

$
  (
    \phi^3 \compose
    \funcname{f}^2 \compose
    \funcname{f}^1 \compose
    \phi^{1-1}
    \maps \phi^1[I^{1,3}_{u^1}] \to V^3
  )
$
agrees with
$
  (
    \phi^3 \compose
    \funcname{f}^2 \compose
    \funcname{f}^1 \compose
    \phi^{1-1}
    \maps \phi^1[I] \to V^3
  )
$
on $\phi^1[I^{1,3}_{u^1}]$ and
$\phi^1[I] = \union[u^1 \in I]{\phi^1[I^{1,3}_{u^1}]}$,
thus
$
  (
    \phi^3 \compose
    \funcname{f}^2 \compose
    \funcname{f}^1 \compose
    \phi^{1-1}
    \maps \phi^1[I] \to V^3
  )
$
is a morphism of $\catname{C}^3$ by
{
  \showlabelsinline
  \cref{mod:sheaf}
}
of \cref{def:model}.

\item [\ensuremath{\funcname{f}^2 \compose \funcname{f}^1} is (semi-strict, strict).]
If $\seqname{E}^1 \submod \seqname{E}^2$ and
$\seqname{E}^2 \submod \seqname{E}^3$ then
$\seqname{E}^1 \submod \seqname{E}^3$. If
$\seqname{C}^1 \submod \seqname{C}^2$ and
$\seqname{C}^2 \submod \seqname{C}^3$ then
$\seqname{C}^1 \submod \seqname{C}^3$.
If $\catname{C}^1 \subcat \catname{C}^2$ and
$\catname{C}^2 \subcat \catname{C}^3$ then
$\catname{C}^1 \subcat \catname{C}^3$.

If $\funcname{f}^1$ is a local morphism of $\seqname{E}^2$ and
$\funcname{f}^2$ is a local morphism of $\seqname{E}^3$ then
$\funcname{f}^2 \compose \funcname{f}^1$ is a local morphism of
$\seqname{E}^3$ by \pagecref{lem:modlocalMmorphcomp}.

If $\funcname{f}^1$ is a local $\catname{E}^1$-$\catname{E}^2$
morphism of $\seqname{E}^2$ and $\funcname{f}^2$ is a local
$\catname{E}^2$-$\catname{E}^3$ morphism of $\seqname{E}^3$ then
$\funcname{f}^2 \compose \funcname{f}^1$ is a local
$\catname{E}^1$-$\catname{E}^3$ morphism of $\seqname{E}^3$ by
\cref{lem:modlocalMmorphcomp}.

If $\funcname{f}^1$ is a morphism of $\seqname{E}^2 \seqname{E}^3$ and
$\funcname{f}^2$ is a morphism of $\seqname{E}^3$ then
$\funcname{f}^2 \compose \funcname{f}^1$ is a morphism of
$\seqname{E}^3$.
\end{description}
\end{proof}
Similar results apply with restrictions on the admissible atlases.
\end{lemma}

\begin{lemma}[Classic m-atlas inclusion and identity]
\label{lem:M-ATLcid}
Let
\begin{enumerate}
\item
$\catname{E}^i$, $\catname{C}^i$, $i=1,2$, be model categories
\item
$\catname{E}^1 \subcat[full-] \catname{E}^2$
\item
$\seqname{E}^i \objin \catname{E}^i$
\item
$\seqname{E}^1 \submod \seqname{E}^2$
\item
$\seqname{C}^i \objin \catname{C}^i$
\item
$\seqname{C}^1 \submod \seqname{C}^2$
($\catname{C}^1 \subcat[full-] \catname{C}^2$)
\item
$\seqname{A}^i$ be a maximal m-atlas of $\seqname{E}^i$ in the
coordinate space $\seqname{C}^i$
\item
$\seqname{A}^1 \subseteq \seqname{A}^2$
\end{enumerate}

The classic m-atlas inclusion morphism
$
 \Id%
   [
     {(\seqname{E}^1,\seqname{C}^1)},
     {(\seqname{E}^2,\seqname{C}^2)}
   ]
  =
  \ID%
    [
      {(\seqname{E}^1)},
      {(\seqname{E}^2)}
    ]
$
is a semi-strict classic
\nobreaks
  {
    {$\seqname{E}^1$-$\seqname{E}^2$},
    {(-$\catname{C}^1$-$\catname{C}^2$)}
  }
m-atlas morphism of $\seqname{A}^1$ to $\seqname{A}^2$ in the coordinate
spaces $\seqname{C}^1$, $\seqname{C}^2$ and a classic
\nobreaks
  {
    {$\catname{E}^1$-$\catname{E}^2$},
    {(-$\catname{C}^1$-$\catname{C}^2$)}
  }
m-atlas morphism of $\seqname{A}^1$ to $\seqname{A}^2$ in the coordinate
spaces $\seqname{C}^1$, $\seqname{C}^2$.

\begin{proof}
\mbox{}
\begin{description}
\item%
[
  \ensuremath
  {
    \Id%
      [
        {\seqname{E}^1},
        {\seqname{E}^2}
      ]
  }
  is a classic m-atlas morphism:
]
\mbox{}

Let $(U^i, V^i, \phi^i) \in \seqname{A}^i$, $i=1,2$, with
$I \defeq U^1 \cap U^2 \neq \emptyset$.

Since $\seqname{A}^1 \subseteq \seqname{A}^2$, $(U^1, V^1, \phi^1)$ is
compatible with $(U^2, V^2, \phi^2)$, \\*
$
  \phi^2 \compose
  \Id%
    [
      {\seqname{E}^1},
      {\seqname{E}^2}
    ]
  \compose
  \phi^{1-1}
  \maps \phi^1[I] \to V^2
  =
  \phi^2 \compose
  \phi^{1-1}
  \maps \phi^1[I] \to V^2
$
is a morphism of $\seqname{C}^2$ and
a morphism of $\catname{C}^2$.

\item%
[
  \ensuremath
  {
    \Id%
      [
        {\seqname{E}^1},
        {\seqname{E}^2}
      ]
  }
  is semi-striict:
]
\end{description}
\end{proof}

If each $\Top(\seqname{E}^i) \objin \seqname{E}^i$
($\seqname{E}^i \objin \catname{E}^i$)
and
\nobreaks{{$\seqname{C}^1 \submod \seqname{C}^2$}}
\nobreaks{{($\catname{C}^1 \subcat \catname{C}^2$)}}
then the classic m-atlas inclusion morphism
$
  \Id%
    [
     {(\seqname{E}^1,\seqname{C}^1)},
     {(\seqname{E}^2,\seqname{C}^2)}
    ]
$
is a strict classic
\nobreaks
  {
    {($\catname{E}^1$-$\catname{E}^2$-)},
    {$\seqname{E}^1$-$\seqname{E}^2$},
    {(-$\catname{C}^1$-$\catname{C}^2$)}
  }
m-atlas morphism of
$\seqname{A}^1$ to $\seqname{A}^2$ in the coordinate spaces
$\seqname{C}^1$, $\seqname{C}^2$.

\begin{proof}
\mbox{}
\begin{description}
\item%
[
  \ensuremath
  {
    \Id%
      [
        {(\seqname{E}^1,\seqname{C}^1)},
        {(\seqname{E}^2,\seqname{C}^2)}
      ]
  }
  is a classic m-atlas morphism
]

$
 \Id%
   [
     {(\seqname{E}^1,\seqname{C}^1)},
     {(\seqname{E}^2,\seqname{C}^2)}
   ]
$
is a classic
\nobreaks{$\seqname{E}^1$-$\seqname{E}^2$}%
\nobreaks{(-$\catname{C}^1$-$\catname{C}^2$)}
m-atlas morphism of $\seqname{A}^1$ to $\seqname{A}^2$ in the coordinate
spaces $\seqname{C}^1$, $\seqname{C}^2$ and a classic
\nobreaks{$\catname{E}^1$-$\catname{E}^2$}%
\nobreaks{(-$\catname{C}^1$-$\catname{C}^2$)}
m-atlas morphism of $\seqname{A}^1$ to $\seqname{A}^2$ in the coordinate
spaces $\seqname{C}^1$, $\seqname{C}^2$.
by the above proof.

\item%
  [
    \ensuremath
      {
        \Id%
          [
            {(\seqname{E}^1,\seqname{C}^1)},
            {(\seqname{E}^2,\seqname{C}^2)}
          ]
      }
    is strict:
  ]


If each $\Top(\seqname{E}^i) \objin \seqname{E}^i$ then
$\Id[{\seqname{E}^1},{\seqname{E}^2}]$ is a morphism of
$\seqname{E}^2$ by \cref{mod:inclusion} of
{
  \showlabelsinline
  \pagecref{def:model}
},
hence an m-morphism of $\seqname{E}^1$ to $\seqname{E}^2$.

If each
$\seqname{E}^i \objin \catname{E}^i$
$\Id[{\seqname{E}^1},{\seqname{E}^2}]$ is a morphism of
$\catname{E}^2$ by \cref{modcat:morph} of
{
  \showlabelsinline
  \pagecref{def:modcat}
},
hence an $\catname{E}^1$-$\catname{E}^2$ m-morphism of $\seqname{E}^1$
to $\seqname{E}^2$.
\end{description}
\end{proof}
\end{lemma}

\begin{corollary}[Classic m-atlas identity]
\label{cor:M-ATLcid}
Let $\catname{E}$, $\catname{C}$, be model categories,
$\seqname{E} \objin \catname{E}$,
$\seqname{C} \objin \catname{C}$,
$\seqname{A}$ an m-atlas of
$\seqname{E}$ in the coordinate spaces $\seqname{C}$.

The identity morphism
$\seqname{f} \defeq \Id[{{(\seqname{A}, \seqname{E}, \seqname{C})}}]$ is a
strict m-atlas morphism of $\seqname{A}$ to $\seqname{A}$ in the
coordinate spaces $\seqname{C}$, $\seqname{C}$.
\end{corollary}

\subsection{Categories of classic m-atlases and functors}
\label{sub:C-cat}
This subsection defines categories of m-atlases with classic m-atlas
morphisms and functors among them, and proves some basic results.

\subsubsection{Classic m-atlas categories}
\begin{definition}[Sets of classic m-atlas morphisms]
Let $\seqname{E}^i$ and $\seqname{C}^i$. $i=1,2$, be model spaces.
Then
\begin{multline}
  \strict[*]
    {
      \maxfull[*]
        {
          \Atl[classic]_\Ar
            (\seqname{E}^1, \seqname{C}^1, \seqname{E}^2, \seqname{C}^2)
        }
    }
  \defeq \\*
  \set
  {{
    \bigl (
      (\funcname{f}),
      (\seqname{A}^1, \seqname{E}^1, \seqname{C}^1),
      (\seqname{A}^2, \seqname{E}^2, \seqname{C}^2)
    \bigr )
  }}%
  [
    {
      \strict[*]
      {
        \maxfull[*]
          {
            \isAtl[classic]_\Ar
            (
              \seqname{A}^1, \seqname{E}^1, \seqname{C}^1,
              \seqname{A}^2, \seqname{E}^2, \seqname{C}^2,
              \funcname{f}
            )
          }
      }
    }
  ]
\end{multline}
\end{definition}

\begin{definition}[Categories $\Atl^\Classic(\seqname{E}, \seqname{C})$]
\label{def:C-cat}
Let $\seqname{E}$ and $\seqname{C}$ be sets of model spaces.
Let $P \defeq \seqname{E} \times \seqname{C}$.
Then
\begin{equation}
  \strict[*]
    {
      \maxfull[*]
        {
          \Atl[classic]_\Ob(\seqname{E}, \seqname{C})
        }
    }
  \defeq
  \union
    [
      {(\seqname{E}^\mu, \seqname{C}^\mu) \in \seqname{P}},
    ]
    {{
      \strict[*]
        {
          \maxfull[*]
            {
              \Atl[classic]_\Ob
              (
                \seqname{E}^\mu,
                \seqname{C}^\mu
              )
            }
        }
    }}
\end{equation}
\begin{equation}
  \strict[*]
    {
      \maxfull[*]
        {
          \Atl[classic]_\Ar(\seqname{E}, \seqname{C})
        }
    }
  \defeq
  \union
    [
      {(\seqname{E}^\mu, \seqname{C}^\mu) \in \seqname{P}},
      {(\seqname{E}^\nu, \seqname{C}^\nu) \in \seqname{P}}
    ]
    {{
      \strict[*]
        {
          \maxfull[*]
            {
              \Atl[classic]_\Ar
              (
                \seqname{E}^\mu,
                \seqname{C}^\mu,
                \seqname{E}^\nu,
                \seqname{C}^\nu
              )
            }
        }
    }}
\end{equation}
\begin{equation}
  \strict[*]
    {
      \maxfull[*]
        {
          \Atl[classic](\seqname{E}, \seqname{C})
        }
    }
  \defeq
  \bigl (
    \maxfull[*]
      {
        \Atl[classic]_\Ob(\seqname{E}, \seqname{C})
      },
    \maxfull[*]
      {
        \Atl[classic]_\Ar(\seqname{E}, \seqname{C})
      },
    \compose[A]
  \bigr )
\end{equation}

Let $\seqname{E}^i$ and $\seqname{C}^i$, $i=1,2$, be model spaces,
$\seqname{A}^i$ be an atlas of $\seqname{E}^i$ in the coordinate
space $\seqname{C}^i$ and
$\funcname{f} \maps \seqname{E}^1 \to \seqname{E}^2$ be a
(strict, semi-strict) $\seqname{E}^1$-$\seqname{E}^2$ classic M-atlas
 morphism of $\seqname{A}^1$ to $\seqname{A}^2$ in the coordinate
spaces $\seqname{C}^1$, $\seqname{C}^2$.

The identity functor $\Functor_{\Classic,\Id}$ is

\begin{equation}
\Functor_{\Classic,\Id} (\seqname{A}^i, \seqname{E}^i, \seqname{C}^i)
\defeq
(\seqname{A}^i, \seqname{E}^i, \seqname{C}^i)
\end{equation}
\begin{multline}
\Functor_{\Classic,\Id}
  \bigl (
    (\funcname{f} \maps \seqname{E}^1 \to \seqname{E}^),
    (\seqname{A}^1, \seqname{E}^1, \seqname{C}^1),
    (\seqname{A}^2, \seqname{E}^2, \seqname{C}^2)
  \bigr )
\defeq
\\*
  \bigl (
    (\funcname{f} \maps \seqname{E}^1 \to \seqname{E}^2 ),
    (\seqname{A}^1, \seqname{E}^1, \seqname{C}^1),
    (\seqname{A}^2, \seqname{E}^2, \seqname{C}^2)
  \bigr )
\end{multline}

This nomenclature will be justified below.
\end{definition}

\begin{lemma}[$\Atl^\Classic(\seqname{E}, \seqname{C})$ is a category]
\label{lem:M-ATLCiscat}
Let $\seqname{E}$ and $\seqname{C}$ be sets of model spaces. Then
$\Atl[Classic](\seqname{E}, \seqname{C})$ is a category and the identity
morphism for an object $(\seqname{A}^i, \seqname{E}^i, \seqname{C}^i)$
of $\Atl[Classic]_\Ob(\seqname{E}, \seqname{C})$ is
$\Id[{{(\seqname{A}^i, \seqname{E}^i, \seqname{C}^i)}}]$.

\begin{proof}
Let $(\seqname{A}^i,\seqname{E}^i,\seqname{C}^i)$, $i \in [1,3]$,
be an object of $\Atl[Classic](\seqname{E}, \seqname{C})$ and
let \\
$
  \funcname{m}^i \defeq
  \bigl (
    (\funcname{f}^i),
    (\seqname{A}^i,\seqname{E}^i,\seqname{C}^i),
    (\seqname{A}^{i+1},\seqname{E}^{i+1},\seqname{C}^{i+1})
  \bigr )
$
be a morphism of $\Atl[Classic](\seqname{E}, \seqname{C})$. Then
\begin{description}
\item [Composition:]
$\funcname{f}^2 \compose \funcname{f}^1$ is a classic
$\seqname{E}^1$-$\seqname{E}^3$ m-atlas morphism of $\seqname{A}^1$ to
$\seqname{A}^3$ in the coordinate spaces $\seqname{C}^1$,
$\seqname{C}^3$ by \pagecref{lem:M-ATLcmorphcomp} and
$
  (\funcname{f}^2) \compose[()] (\funcname{f}^1) =
  (\funcname{f}^2 \compose \funcname{f}^1)
$.
\item [Associativity:]
Composition is associative by \pagecref{lem:labcomp}.
\item [Identity:]
$\Id[{{(\seqname{A}^i, \seqname{E}^i, \seqname{C}^i)}}]$ is an identity
morphism by \cref{lem:labcomp}.
\end{description}
\end{proof}
\end{lemma}


\begin{corollary}[Subcategories of $\Atl^\Classic(\seqname{E}, \seqname{C})$]
\label{cor:ATLCiscat}
Let $\seqname{E}$ and $\seqname{C}$ be sets of model spaces. Then
all of the following are subcategories of
$\Atl[Classic](\seqname{E}, \seqname{C})$.

\begin{enumerate}
\item $\full{\Atl[Classic]}(\seqname{E}, \seqname{C})$
\item $\maximal[S-]{\Atl[Classic]}(\seqname{E}, \seqname{C})$
\item $\maximal{\Atl[Classic]}(\seqname{E}, \seqname{C})$
\item $\maxfull[S-]{\Atl[Classic]}(\seqname{E}, \seqname{C})$
\item $\maxfull{\Atl[Classic]}(\seqname{E}, \seqname{C})$
\item $\strict[semi-]{\Atl[Classic]}(\seqname{E}, \seqname{C})$
\item $\strict[semi-]{\full{\Atl[Classic]}}(\seqname{E}, \seqname{C})$
\item $\strict[semi-]{\maximal[S-]{\Atl[Classic]}}(\seqname{E}, \seqname{C})$
\item $\strict[semi-]{\maximal{\Atl[Classic]}}(\seqname{E}, \seqname{C})$
\item $\strict[semi-]{\maxfull[S-]{\Atl[Classic]}}](\seqname{E}, \seqname{C})$
\item $\strict[semi-]{\maxfull{\Atl[Classic]}}(\seqname{E}, \seqname{C})$
\item $\strict{\Atl[Classic]}(\seqname{E}, \seqname{C})$
\item $\strict{\full{\Atl[Classic]}}(\seqname{E}, \seqname{C})$
\item $\strict{\maximal[S-]{\Atl[Classic]}}(\seqname{E}, \seqname{C})$
\item $\strict{\maximal{\Atl[Classic]}}(\seqname{E}, \seqname{C})$
\item $\strict{\maxfull[S-]{\Atl[Classic]}}(\seqname{E}, \seqname{C})$
\item $\strict{\maxfull{\Atl[Classic]}}(\seqname{E}, \seqname{C})$
\end{enumerate}

\begin{remark}
They are not, in general, full subcategories.
\end{remark}

The identity functor $\Functor_{\Classic,\Id}$ is a functor from each of the
subcategories to each of the containing categories.
\end{corollary}

\begin{definition}[Categories $\Atl^\Classic(\catname{E}, \catname{C})$]
\label{def:C-cat-cat}
Let $\catname{E}$ and $\catname{C}$ be model categories.
Let $\seqname{P} \defeq Ob(\catname{E}) \times \Ob(\catname{C})$.
Then
\begin{equation}
  \Atl[\Classic]_\Ob(\catname{E}, \catname{C})
  \defeq
  \union
    [
        {\seqname{E} \objin \catname{E}},
        {\seqname{C} \objin \catname{C}}
    ]
    {{
      \Atl[\Classic]_\Ob(\seqname{E}, \seqname{C})
    }}
\end{equation}
\begin{multline}
\Atl[Classic]_\Ar(\catname{E}, \catname{C})
\defeq
  \set
  {{
    \bigl (
      (\funcname{f}),
      (\seqname{A}^1, \seqname{E}^1, \seqname{C}^1),
      (\seqname{A}^2, \seqname{E}^2, \seqname{C}^2)
    \bigr )
  }}%
  [
    {
      (\seqname{E}^i, \seqname{C}^i) \in \seqname{P}
    },
    {
      \strict{\isAtl_\Ar}
      (
        \seqname{A}^1,
        \seqname{E}^1,
        \seqname{C}^1,
        \seqname{A}^2,
        \seqname{E}^2,
        \seqname{C}^2,
        \funcname{f},
        \catname{E},
        \catname{C}
      )
    }
  ]*
\end{multline}
\begin{equation}
  \Atl[\Classic](\catname{E}, \catname{C})
  \defeq
  \bigl (
    \Atl[\Classic]_\Ob(\catname{E}, \catname{C}),
    \Atl[\Classic]_\Ar(\catname{E}, \catname{C}),
    \compose[A]
  \bigr )
\end{equation}

Similar definitions apply with restrictions on the admissible atlases,
the admissible morphisms, or both:
\begin{enumerate}
\item $\full{\Atl[\Classic]}$
\item $\maximal[S-]{\Atl[\Classic]}$
\item $\maximal{\Atl[\Classic]}$
\item $\maxfull[S-]{\Atl[\Classic]}$
\item $\maxfull{\Atl[\Classic]}$
\item $\strict[semi-]{\full{\Atl[\Classic]}}$
\item $\strict[semi-]{\maximal[S-]{\Atl[\Classic]}}$
\item $\strict[semi-]{\maximal{\Atl[\Classic]}}$
\item $\strict[semi-]{\maxfull[S-]{\Atl[\Classic]}}$
\item $\strict[semi-]{\maxfull{\Atl[\Classic]}}$
\item $\strict{\full{\Atl[\Classic]}}$
\item $\strict{\maximal[S-]{\Atl[\Classic]}}$
\item $\strict{\maximal{\Atl[\Classic]}}$
\item $\strict{\maxfull[S-]{\Atl[\Classic]}}$
\item $\strict{\maxfull{\Atl[\Classic]}}$
\end{enumerate}
e.g.,

\begin{equation}
  \maxfull[S-]{\Atl[\Classic]_\Ob}l(\catname{E}, \catname{C})
  \defeq
  \union
    [
      {(\seqname{E}^\mu, \seqname{C}^\mu) \in \seqname{P}},
    ]
    {{
      \maxfull[S-]{\Atl[\Classic]_\Ob}(\seqname{E}^\mu, \seqname{C}^\mu)
    }}
\end{equation}
\begin{equation}
  \strict[semi-]{\maxfull[S-]{\Atl[\Classic]_\Ar}}(\catname{E}, \catname{C})
  \defeq
  \union
    [
      {(\seqname{E}^\mu, \seqname{C}^\mu) \in \seqname{P}},
      {(\seqname{E}^\nu, \seqname{C}^\nu) \in \seqname{P}}
    ]
    {{
      \strict[semi-]{\maxfull[S-]{\Atl[\Classic]_\Ar}}
      (
        \seqname{E}^\mu,
        \seqname{C}^\mu,
        \seqname{E}^\nu,
        \seqname{C}^\nu
      )
    }}
\end{equation}
\begin{equation}
  \strict[semi-]{\maxfull[S-]{\Atl[\Classic]}}(\catname{E}, \catname{C})
  \defeq
  \biggl (
    \maxfull[S-]{\Atl[\Classic]_\Ob}l(\catname{E}, \catname{C}),
    \strict[semi-]{\maxfull[S-]{\Atl[\Classic]_\Ar}}(\seqname{E}, \seqname{C}),
    \compose[A]
  \biggr )
\end{equation}

Let $(\seqname{A}^1, \seqname{E}^1, \seqname{C}^1) \in \Atl[\Classic]_\Ob(\catname{E}, \catname{C})$.
\begin{equation}
\Id[{{(\seqname{A}^1, \seqname{E}^1, \seqname{C}^1)}}] \defeq
\bigl (
  (\Id[{\seqname{E}^1}], \Id[{\seqname{C}^1}]),
  (\seqname{A}^1, \seqname{E}^1, \seqname{C}^1),
  (\seqname{A}^1, \seqname{E}^1, \seqname{C}^1)
\bigr )
\end{equation}
\end{definition}

\begin{lemma}[$\Atl^\Classic(\catname{E}, \catname{C})$ is a category]
\label{lem:C-cat-ATLCiscat}
Let $\catname{E}$ and $\catname{C}$ be model categories. Then
$\Atl[\Classic](\catname{E}, \catname{C})$ is a category and the identity morphism
for an object $(\seqname{A}^i, \seqname{E}^i, \seqname{C}^i)$ of
$\Atl[\Classic]_\Ob(\catname{E}, \catname{C})$ is
$\Id[{{(\seqname{A}^i, \seqname{E}^i, \seqname{C}^i)}}]$.

\begin{proof}
Let $(\seqname{A}^i,\seqname{E}^i,\seqname{C}^i)$, $i \in [1,3]$,
be an object of $\Atl[\Classic](\catname{E}, \catname{C})$ and
let \\
$
  \funcname{m}^i \defeq
  \bigl (
    (\funcname{f}^i),
    (\seqname{A}^i,\seqname{E}^i,\seqname{C}^i),
    (\seqname{A}^{i+1},\seqname{E}^{i+1},\seqname{C}^{i+1})
  \bigr )
$, $i=1,2$,
be a morphism of $\Atl[\Classic](\catname{E}, \catname{C})$. Then
\begin{description}
\item[Composition:]
$\funcname{f}^2 \compose \funcname{f}^1$ is a morphism of $\catname{E}$ and
is an $E^1$-$E^3$ m-atlas morphism of $\seqname{A}^1$ to $\seqname{A}^3$
in the coordinate spaces $\seqname{C}^1$, $\seqname{C}^3$ by
\pagecref{lem:M-ATLcmorphcomp}.

\item[Associativity:]
Composition is associative by \pagecref{lem:labcomp}.

\item[Identity:]
$\Id[{{(\seqname{A}^i, \seqname{E}^i, \seqname{C}^i)}}]$ is an identity
morphism by \cref{lem:labcomp}.
\end{description}
\end{proof}

Similar results follow with restrictions on the admissible atlases,
the admissible morphisms, or both.
\end{lemma}

\subsubsection{Classic m-atlas functors}
\begin{definition}[$\Functor_{\M,\Classic}$]
\label{def:Func-M-C}
Let $\seqname{E}^i$, $i=1,2$, $\seqname{C}^i$ be model spaces and
$\seqname{A}^i$ an m-atlas of $\seqname{E}^i$ in the coordinate model
space $\seqname{C}^i$. Then
\begin{equation}
\Functor_{\M,\Classic} (\seqname{A}^i, \seqname{E}^i, \seqname{C}^i)
\defeq (\seqname{A}^i, \seqname{E}^i, \seqname{C}^i)
\end{equation}

Let
$
  \funcseqname{f} \defeq
    (
    \funcname{f}_0 \maps \seqname{E}^1 \to \seqname{E}^2,
    \funcname{f}_1 \maps \seqname{C}^1 \to \seqname{C}^2,
    )
$
be an M-atlas morphism of $\seqname{A}^1$ to $\seqname{A}^2$ in the
coordinate spaces $\seqname{C}^1$, $\seqname{C}^2$. Then
\begin{multline}
\Functor_{\M,\Classic}
  \bigl (
    (
      \funcname{f}_0 \maps \seqname{E}^1 \to \seqname{E}^2,
      \funcname{f}_1 \maps \seqname{C}^1 \to \seqname{C}^2,
    ),
    (\seqname{A}^1, \seqname{E}^1, \seqname{C}^1),
    (\seqname{A}^2, \seqname{E}^2, \seqname{C}^2)
  \bigr )
\defeq \\
  \bigl (
    (
      \funcname{f}_0 \maps \seqname{E}^1 \to \seqname{E}^2
    ),
    (\seqname{A}^1, \seqname{E}^1, \seqname{C}^1),
    (\seqname{A}^2, \seqname{E}^2, \seqname{C}^2)
  \bigr )
\end{multline}
\end{definition}

\begin{theorem}[$\Functor_{\M,\Classic}$ is a functor]
\label{the:Func-M-C}
Let $\seqname{E}$ and $\seqname{C}$ be sets of model spaces, Then
$\Functor_{\M,\Classic}$ is a functor from
$\Atl(\seqname{E}, \seqname{C})$ to
$\Atl[classic](\seqname{E}, \seqname{C})$.

\begin{proof}
Let $o^i \defeq  (\seqname{A}^i, \seqname{E}^i, \seqname{C}^i)$,
$i \in [1,3]$, be an object of $\Atl(\seqname{E}, \seqname{C})$ and \\*
$
  m^i \defeq
  \bigl (
    (f_0^i, f_1^i),
    o^i,
    o^{i+1}
  \bigr )
$, $i=1,2$,
be a morphism of $\Atl(\seqname{E}, \seqname{C})$.

$\Functor_{\M,\Classic}$ satisfies these criteria:
\begin{description}
\item[Preservation of endpoints:]
\begin{align*}
  \Functor_{\M,\Classic} o^i
  & =
  o^i \objin \Atl[classic](\seqname{E}, \seqname{C})
\\*
  \Functor_{\Top,\M} m^i
  & =
  \bigl (
    (
      \funcname{f}_0 \maps \seqname{E}^i \to \seqname{E}^{i+1}
    ),
    (\seqname{A}^i, \seqname{E}^i, \seqname{C}^i),
    (\seqname{A}^{i+1}, \seqname{E}^{i+1}, \seqname{C}^{i+1})
  \bigr )
\\*
  & \quad
  \text{is a morphism from $o^i$ to $o^{i+1}$.}
\end{align*}

\item[Identity:]
Let
$
  \bigl (
    (\Id[{\seqname{E}^i}], \Id[{\seqname{C}^i}]),
   o^i,
   o^i
  \bigr )
$
be an identity morphism of
$o^i$ in $\Atl(\seqname{E}, \seqname{C})$.
Then
$
  \bigl (
    (
     \Id[\seqname{E}^i] \maps \seqname{E}^i \to \seqname{E}^i
    ),
    o^i,
    o^i
  \bigr )
$
is an identity morphism of $o^i$ in $\Atl[classic](\seqname{E}, \seqname{C})$.

\item[Composition:]
\begin{equation*}
\begin{array}{lll}
  \Functor_{\M,\Classic} (m^2 \compose[A] m^i)
  & =
  &
  \Functor_{\M,\Classic}
    \bigl (
      (f_0^2 \compose f_0^1, f_1^2 \compose f_1^1),
      o^1,
      o^3
    \bigr )
\\*
  & =
  &
    \bigl (
      (f_0^2 \compose f_0^1),
      o^1,
      o^3
    \bigr )
\\*
  & =
  &
    \bigl (
      (f_0^2),
      o^1,
      o^2
    \bigr )
    \compose[A]
    \bigl (
      (f_0^1),
      o^2,
      o^3
    \bigr )
\\*
  & =
  &
  \Functor_{\M,\Classic}
    \bigl (
      (f_0^2, f_1^2),
      o^2,
      o^3
    \bigr )
  \compose[A]
\\*
  &
  &
    \Functor_{\M,\Classic}
    \bigl (
      (f_0^1, f_1^1),
      o^1,
      o^2
    \bigr )
\\*
  & =
  &
  \Functor_{\M,\Classic} m^2
  \compose[A]
  \Functor_{\M,\Classic} m^1
\end{array}
\end{equation*}
\end{description}
\end{proof}
\end{theorem}

\part{Equivalence of manifolds}
\label{part:man}
For a manifold\footnote{
  The literature has several different definitions of a manifold.
  This paper uses one chosen for ease of exposition.
}
the coordinate category is open subsets of a Banach space or more
generally a Fr\'echet space, with an appropriate choice of morphisms.
Choosing a separating hyperplane and half space, with open sets in the
chosen half space, allows manifolds with boundary. Similarly, choosing a
ball with tails allows a manifold with tails.

For differentiable manifolds the coordinate category is similar, but the
morphisms are limited to those sufficiently differentiable, in order to
impose a diferentiability constraint on the transition functions
$\funcname{t}^\alpha_\beta=\phi_\beta \compose \phi_\alpha^{-1}$.

This part of the paper defines $\Ck$-atlases, $\Ck$-manifolds, local
coordinate spaces equivalent to $\Ck$-manifolds, categories of them and
functors, and gives some basic results.

\section {Linear spaces and linear model spaces}
{
  \showlabelsinline
  \label{sec:lin}
}

This subsection defines spaces and related model spaces suitable for use
as the coordinate spaces of generalized $\Ck$ manifolds.

\begin{definition}[Linear spaces]
{
  \showlabelsinline
  \label{def:lin}
}
Let $C$ be a locally arcwise connected topological subspace of a real
(complex) Banach or Fr\'echet space. Then $C$ is a real (complex) linear
space.  If $C$ has a non-void interior, i.e., contains a ball, then it
is non-degenerate.

\begin{remark}
This paper uses the term \textit{ball} to refer to balls of the
underlying space but uses the terms \textit{open set} and
\textit{neighborhoods} to refer to the relative topology.
\end{remark}

Let $\catname{C}$ be a small category whose objects are real (complex)
linear spaces and whose morphisms are $\Ck$ functions. Then
$\catname{C}$ is a $\Ck$ linear category.  If each object of
$\catname{C}$ is non-degenerate then $\catname{C}$ is non-degenerate.

Let $C$ be a real (complex) linear space. Then $\trivcat[\Ck-]{C}$, the
category of all $\Ck$ functions between nonvoid open subspaces of $C$,
is the trivial $\Ck$ linear category of $C$.
\end{definition}

\begin{lemma}[Open subsets of linear spaces are locally arcwise connected]
Let $U$ be an open subset of the real (complex) linear space $S$. Then
$U$ is locally arcwise connected.

\begin{proof}
An open subset of a locally arcwise connected space is locally arcwise
connected.
\end{proof}
\end{lemma}

\begin{definition}[Linear model spaces]
Let $S$ be a real (complex) linear space and
$\seqname{S} \defeq (S, \catname{S})$ a model space for $S$.
Then $\seqname{S}$ is a real (complex) linear model space.

Let $\seqname{S} \defeq (S, \catname{S})$ be a real (complex) linear
model space such that every morphism of $\catname{S}$ is a $\Ck$
function. Then $\seqname{S}$ is a real (complex) $\Ck$ linear model
space.

Let $\catname{S}$ be a small category whose objects are real (complex)
$\Ck$ linear model spaces and whose morphisms are $\Ck$ model functions.
Then $\catname{S}$ is a $\Ck$ linear model category.
\end{definition}

\begin{definition}[Trivial \ensuremath{\Ck} linear model spaces]
\label{def:trivck}
Let $C$ be a real (complex) linear space and $\catname{C}$ the category
of all $\Ck$ functions between open sets of $C$. Then
$\Triv[\Ck-]{C} \defeq (C, \catname{C})$ is the trivial $\Ck$ linear
model space of $C$ and $\Triv[\Ck-]{C}$ is a real (complex) trivial
$\Ck$ linear model space.

Let $\seqname{C}$ be a set of real (complex) linear spaces.

$
  \Triv[\Ck-]{\seqname{C}}
  \defeq
  \set
    {\Triv[\Ck-]{C'}}%
    [C' \in \seqname{C}]
$
is the set of trivial $\Ck$ linear model spaces of $\seqname{C}$.

The category of trivial $\Ck$ linear model spaces of $\seqname{C}$,
abbreviated $\Trivcat[\Ck-]{\seqname{C}}$, is the category whose
objects are $\Triv[\Ck-]{\seqname{C}}$ and whose morphisms are all the
$\Ck$ functions among them.

$\optriv[\Ck-]{\seqname{C}}$, the category of open trivial $\Ck$ model
spaces in $\seqname{C}$, is the category whose objects are
$\set {\Triv[\Ck-]{U}}[U \in \op{\seqname{C}}]$, the trivial $\Ck$
linear model spaces of non-null open sets of spaces in $\seqname{C}$,
and whose morphisms are all the $\Ck$ functions among them.
\end{definition}

\begin{lemma}[Trivial \ensuremath{\Ck} linear model spaces]
\label{lem:trivck}
Let $C^i$, $i=1,2$, be a real (complex) linear space and
$\catname{C}^i \defeq \Cat \biggl ( \Triv[\Ck-]{C^i} \biggr )$.

\begin{enumerate}
\item
$\Triv[\Ck-]{C^i}$ is a $\Ck$ linear model space.

\begin{proof}
\leavevmode
\begin{enumerate}
\item $\Ob(\catname{C}^i)$ is an open cover for $C^i$.

\item $\Ob(\catname{C}^i)$ is closed under finite intersections.

\item The morphisms of $\catname{C}^i$ are $\Ck$, hence continuous.

\item If $\funcname{f} \maps A \to B$ is a morphism,
$A' \objin \catname{C}^i \subseteq A \objin \catname{C}^i$,
$B' \objin \catname{C}^i \subseteq B \objin \catname{C}^i$
and $\funcname{f}[A'] \subseteq B'$ then
$\funcname{f} \restrictto_{A'} \maps A' \to B'$ is $\Ck$ and thus a
morphism.

\item If $A' \objin \catname{C}^i \subseteq A \objin \catname{C}^i$
then the inclusion map $\funcname{i} \maps A' \hookrightarrow A$ is
$\Ck$ and thus a morphism.

\item Restricted sheaf condition: Whenever
\begin{enumerate}
\item $U_\alpha$ and $V_\alpha$, $\alpha \prec \Alpha$,
are objects of $\catname{C}^i$.

\item $\funcname{f}_\alpha \maps U_\alpha \to V_\alpha$ are morphisms of
$\catname{C}^i$.

\item
$U \defeq \union[\alpha \prec \Alpha]{U_\alpha} \objin \catname{C}^i$,

\item
$V \defeq \union[\alpha \prec \Alpha]{V_\alpha} \objin \catname{C}^i$

\item
$\funcname{f} \maps U \to V$ is a continuous function and
for every $\alpha \prec \Alpha$,
$\funcname{f}$ agrees with $\funcname{f}_\alpha$ on $U_\alpha$
\end{enumerate}
then $\funcname{f}$ is $\Ck$ and thus a morphism of $\catname{C}^i$.
\end{enumerate}
\end{proof}

\item
Let $\funcname{f} \maps C^1 \to C^2$ be $\Ck$. Then $\funcname{f}$
is a model function from $\Triv[\Ck-]{C^1}$ to $\Triv[\Ck-]{C^2}$.

\begin{proof}
Let $U^i$ be a model neighborhood of $\Triv[\Ck-]{C^i}$, $i=1,2$.
\begin{enumerate}
\item
Since $U^2$ is a model neighborhood of $\Triv[\Ck-]{C^2}$, $U^2$ is open,
$\funcname{f}^{-1}[U^2]$ is open and thus a model neighborhood of
$\Triv[\Ck-]{C^1}$.
\item
$\funcname{f}[U^1] \subseteq C^2$. Since $C^2$ is open, it is a model
neighborhood of $\Triv[\Ck-]{C^2}$.
\end{enumerate}
\end{proof}

\item
Let $C^1$ be an open subspace of $C^2$. Then
$\Triv[\Ck-]{C^1} \submod \Triv[\Ck-]{C^2}$.

\begin{proof}
By \cref{def:trivck} above, $\Triv[\Ck-]{C^i} = (C^i, \catname{C}^i)$,
where $\catname{C}^i)$ is the category of all $\Ck$ functions between
open sets of $C^i$.

\begin{description}
\item[\ensuremath{\Triv[\Ck-]{C^1} \submod \Triv[\Ck-]{C^2}}:]
\mbox{}
\begin{enumerate}
\item $C^1$ is an open subspace of $C^2$ by hypothesis.

\item Every object of $\catname{C}^1$ is open in $C^1$, thus open in
$C^2$ and thus an object of $\catname{C}^2$. If
$\funcname{f} \maps U \to V \arin \catname{C}^1$ then $U$
and $V$ are open in $C^1$ and $\funcname{f}$ is $\Ck$ in $C^1$,
thus $U$ and $V$ are open in $C^2$ and $\funcname{f}$ is $\Ck$ in
$C^2$ and $\funcname{f} \maps U \to V \arin \catname{C}^2$.
If $\funcname{f} \maps U \to V$ is a morphism of
$\catname{C}^2$, $U \objin \catname{C}^1$ and
$V \objin \catname{C}^1$, then $U$ and $V$ are open in $C^1$ and
$\funcname{f}$ is $\Ck$ in $C^2$, thus $\Ck$ in $C^1$, and
$\funcname{f} \maps U \to V \arin \catname{C}^1$.

\item If $U \objin \catname{C}^2$ then $U$ is open in $C^2$, hence
$U \cap C^1$ is open in $C^2$, hence $U \cap C^1$ is open in $C^1$ and
$U \objin \catname{C}^1$.

\item If $U \objin \catname{C}^2$ and $U \subseteq C^1$ then U is open
is $C^2$, hence open is $C^1$, hence $U \objin \catname{C}^1$.
\end{enumerate}
\end{description}
\end{proof}
\end{enumerate}
\end{lemma}

\begin{definition}[ \ensuremath{\Ck} singleton categories]
\label {def:singCk}
Let $\seqname{C}$ be a $\Ck$ linear model space. Then the $\Ck$
singleton category of $\seqname{C}$, abbreviated
$\singcat[\Ck-]{\seqname{C}}$, is the category whose sole object is
$\seqname{C}$ and whose morphisms are all of the $\Ck$ model functions
from $\seqname{C}$ to itself.
\end{definition}

\section{$\Ck$-nearly commutative diagrams}
\label{sec:Ck-NCD}
Let $C$ be a linear space, $\catname{C} \defeq \trivcat[\Ck-]{S}$,
$\seqname{C} \defeq \Triv[\Ck-]{C}$ and $D$ a tree with two branches,
whose nodes are topological spaces $U_i$ and $V^j$ and whose links are
continuous functions $\funcname{f}_i \maps U_i \to U_{i+1}$ and
$\funcname{f}'_j \maps U_j \to U_{j+1}$ between the spaces:

\begin{align*}
D = \{
  &
    \funcname{f}_0 \maps U_0 = V_0 \to U_1,
    \dotsc,
    \funcname{f}_{m - 1} \maps U_{m - 1} \to U_m,
\\
  &
    \funcname{f}'_0 \maps  U_0 = V_0 \to V_1,
    \dotsc,
    \funcname{f}'_{m-1} \maps V_{m-1} \to V_n
\}
\end{align*}
with $U_0 = V_0$, $U_m \subseteq C$ and $V_n \subseteq C$ open, as shown
in \pagecref{fig:NCDb}.

\begin{definition}[\ensuremath{\Ck}-nearly commutative diagrams]
\label{def:Ck-NCD}
$D$ is (left,right,strongly) $\Ck$-nearly commutative in linear space
$C$ iff $D$ is (left,right,strongly) M-nearly commutative in model space
$\seqname{C}$.

\end{definition}

\begin{definition}[\ensuremath{\Ck}-nearly commutative diagrams at a point]
\label{def:Ck-NCDp}
Let $C$, $\catname{C}$, $\seqname{C}$ and $D$ be as above and $x$ be an
element of the initial node. $D$ is (left,right,strongly) $\Ck$-nearly
commutative in $C$ at $x$ iff $D$ is (left,right,strongly) M-nearly
commutative in model space $\seqname{C}$ at $x$.

\end{definition}

\begin{definition}[\ensuremath{\Ck}-locally nearly commutative diagrams]
\label{def:Ck-NCDl}
Let $C$, $\catname{C}$, $\seqname{C}$ and $D$ be as above. $D$ is
(left,right,strongly) $\Ck$-locally nearly commutative in $C$ iff $D$ is
(left,right,strongly) M-locally nearly commutative in model space
$\seqname{C}$.

\end{definition}

\begin{lemma}
\label{lem:Ck-NCD}
Let $C$, $\catname{C}$, $\seqname{C}$ and $D$ be as above.  $D$ is
$\Ck$-nearly commutative in linear space $C$ iff There is a $\Ck$
diffeomorphism $\hat{\funcname{f}} \maps U_m \toiso V_n$ making the
graph a commutative diagram, as shown in \cref{fig:NCDa}.

\begin{proof}
A set is a model neighborhood of $\seqname{C}$ (an object of
$\trivcat[\Ck-]{S}$) iff it is an open subset of $A$. A function is a
morphism of $\seqname{C}$ (a morphism of $\trivcat[\Ck-]{S}$) iff it is
a $\Ck$ function between open subsets of $S$.
\end{proof}
\end{lemma}

\section{$\Ck$ charts}
{
  \showlabelsinline
  \label{sec:ckcharts}
}
\begin{definition}[\ensuremath{\Ck} charts]
\label{def:Ck-chart}
Let $C$ be a linear space and $E$ a topological space. A $\Ck$\footnote{
With $k \in \mathbb{N} \cup \{ \infty, \omega \}$.}
chart $(U, V, \phi)$ of $E$ in the coordinate space $C$ consists of

\begin{enumerate}
\item A nonvoid open subset $U \subseteq E$, known as a coordinate patch
\item An open subset $V \subseteq C$
\item A homeomorphism $\phi \maps U \toiso V$, known as a
coordinate function
\end{enumerate}
\begin{remark}
I consider it clearer to explicate the range, rather than the
conventional usage of specifying only the domain and function or the
minimalist usage of specifying only the function.
\end{remark}

A chart $(U, V, \phi)$ is non-degenerate iff $V$ contains a ball of the
underlying Banach or Fr\'echet space.
\end{definition}

\begin{lemma}[\ensuremath{\Ck} charts]
\label{lem:Ck-chart}
Let $C$ be a linear space and $E$ a topological space. The triple
$(U, V, \phi)$ is a $\Ck$ chart of $E$ in the coordinate space $C$ iff
it is an m-chart of $\Triv{E}$ in the coordinate space $\Triv[\Ck-]{C}$.

\begin{proof}
$U$ is a model neighborhood of $\Triv{E}$ iff it is a nonvoid open set
of $E$.  $V$ is a model neighborhood of $\Triv[\Ck-]{C}$ iff it is a
nonvoid open set of $C$.  $\phi$ is required to be a homeomorphis from
$U$ to $V$ in either case.
\end{proof}
\end{lemma}

\begin{definition}[\ensuremath{\Ck} subcharts]
\label{def:Ck-subchart}
Let $(U, V, \phi)$ be a $\Ck$ chart of $E$ in the coordinate space $C$
and $U'$ be a nonvoid open subset of $U$. Then
$(U', V', \phi') \defeq (U', \phi[U'], \phi \restrictto_{U',V'})$
is a subchart of $(U, V, \phi)$.

By abuse of language we will write $(U', V', \phi)$ for
$(U', V', \phi')$.
\end{definition}

\begin{lemma}[\ensuremath{\Ck} subcharts]
\label{lem:Ck-subchart}
Let $(U, V, \phi)$ be a $\Ck$ chart of $E$ in the coordinate space $C$
and $(U', V', \phi')$ be a triple. Then $(U', V', \phi')$ is a subchart
of $(U, V, \phi)$ iff it is a subchart of $(U, V, \phi)$ considered as a
chart of $\Triv{E}$ in the coordinate space $\Triv[\Ck-]{C}$.

\begin{proof}
$(U',V',\phi')$ satisfies the conditions of \cref{def:Ck-chart};
\begin{enumerate}
\item
$U$ is open by \cref{def:Ck-chart}
\item
$U'$ is required to be a nonvoid open subset of $U$ in either case,
and thus a model neighborhood of $\Triv{E}$.
\item
$\phi$ is a homeomorphism, so $V'=\phi[U']$ is open,
and thus a model neighborhood of $\Triv[\Ck-]{C}$.
\item
$U'$ is a model neighborhood of $\Triv{E}$ iff it is an open set of $E$.
\item
$\phi$ is a homeomorphism, so
$\phi \restrictto_{U'} \maps U' \toiso \phi[U']$ is also.
\end{enumerate}
\end{proof}
\end{lemma}

\begin{corollary}[\ensuremath{\Ck} subcharts]
\label{cor:Ck-subchart}
Let $(U, V, \phi)$ be a $\Ck$ chart of $E$ in the coordinate space $C$
and $(U', V', \phi')$ be a subchart of $(U, V, \phi)$. Then
$(U', V', \phi')$ is a $\Ck$ chart of $E$ in the coordinate space $C$.

\begin{proof}
The result follows from \fullcref{lem:Ck-chart} above,
{
  \showlabelsinline
  \pagecref{lem:M-chart}
}
and \pagecref{def:M-subchart}.
\end{proof}
\end{corollary}

\begin{definition}[\ensuremath{\Ck} compatibility]
\label{Ck-compat}
Let $(U, V, \phi)$ and $(U',V', \phi')$ be $\Ck$ charts of $E$ in the
coordinate space $C$. Then $(U, V, \phi)$ is $\Ck$ compatible with
$(U', V', \phi')$ iff either
\begin{enumerate}
\item $U$ and $U'$ are disjoint
\item The transition function $\funcname{t}=\phi' \compose \phi^{-1} \restrictto_{\phi[U \cap U']}$
is a $\Ck$ diffeomorphism.
\end{enumerate}
\end{definition}

\begin{lemma}[Symmetry of \ensuremath{\Ck} compatibility]
\label{lem:Ck-compatsym}
Let $(U, V, \phi)$ and $(U',V', \phi')$ be $\Ck$ charts of $E$ in the
coordinate space $C$.
Then $(U, V, \phi)$ is $\Ck$ compatible with $(U', V', \phi')$ iff
$(U', V', \phi')$ is $\Ck$ compatible with $(U,V,\phi)$.

\begin{proof}
It suffices to prove the implication in only one direction.
\begin{enumerate}
\item $U \cap U' = U' \cap U$.
\item Since the transition function
$\funcname{t}=\phi' \compose \phi^{-1} \restrictto_{\phi[U \cap U']}$
is a  $\Ck$ diffeomorphism of $C$, so is
$\funcname{t}^{-1} = \phi \compose \phi'^{-1} \restrictto_{\phi'[U \cap U']}$.
\end{enumerate}
\end{proof}
\end{lemma}

\begin{lemma}[\ensuremath{\Ck} compatibility of subcharts]
Let $(U_i, V_i, \phi_i)$, $i=1,2$, be $\Ck$ charts of $E$ in the
coordinate space $C$, $(U'_i, V'_i, \phi'_i)$ be subcharts and
$(U_1, V_1, \phi_1)$ be $\Ck$ compatible with $(U_2, V_2, \phi_2)$. Then
$(U'_1, V'_1, \phi'_1)$ is $\Ck$ compatible with
$(U'_2, V'_2, \phi'_2)$.

\begin{proof}
If $U_1 \cap U_2 = \emptyset$ then $U'_1 \cap U'_2 = \emptyset$. If
$U'_1 \cap U'_2 = \emptyset$ then $(U'_1, V'_1, \phi'_1)$ is $\Ck$
compatible with $(U'_2, V'_2, \phi'_2)$. Otherwise, the transition
function
$
  t^1_2 \defeq
  \phi_2 \compose \phi^{-1}_1 \restrictto_{\phi_1[U_1 \cap U_2]}
$
is a $\Ck$ diffeomorphism and hence
$
  t^1_2 \restrictto_{\phi_1[U'_1 \cap U'_2]} \maps
  \phi_1[U'_1 \cap U'_2] \toiso
  \phi_2[U'_1 \cap U'_2]
$
is a $\Ck$ diffeomorphism.
\end{proof}
\end{lemma}

\begin{corollary}[\ensuremath{\Ck} compatibility with subcharts]
Let $(U, V, \phi)$ be a $\Ck$ charts of $E$ in
the coordinate space $C$ and $(U', V', \phi')$ a
subchart. Then $(U', V', \phi')$ is $\Ck$ compatible with
$(U,V,\phi)$.

\begin{proof}
$(U, V, \phi)$ is $\Ck$ compatible with itself and is a subchart of
itself,
\end{proof}
\end{corollary}

\begin{definition}[Covering by \ensuremath{\Ck} charts]
Let $\seqname{A}$ be a set of charts of $E$ in the coordinate space $C$.
$\seqname{A}$ covers $E$ iff $\pi_1[\seqname{A}]$ covers $E$, i.e.,
$E = \union{{\pi_1[\seqname{A}]}}$.
\end{definition}

\section{$\Ck$-atlases}
\label{sec:ckatlases}
A set of charts can be atlases for different coordinate spaces even if
it is for the same total space.  In order to aggregate them into
categories, there must be a way to distinguish them. Including the
two\footnote{The total space is redundant, but convenient.}
spaces in the definitions of the categories serves the purpose.

\begin{definition}[\ensuremath{\Ck}-atlases]
\label{def:Ck-atlas}
Let $\seqname{A}$ be a set of mutually $\Ck$ compatible charts of $E$ in
the coordinate space $C$. $\seqname{A}$ is a $\Ck$-atlas of $E$ in
the coordinate space $C$, abbreviated
$\isAtl[\Ck]_\Ob(\seqname{A}, E, C)$, iff $\seqname{A}$ covers $E$.

$\seqname{A}$ is a full $\Ck$-atlas of $E$ in
the coordinate space $C$, abbreviated
$\full{\isAtl[\Ck]_\Ob}(\seqname{A}, E, C)$, iff
\begin{enumerate}
\item $\pi_1[\seqname{A}]$ covers $E$
\item $\pi_2[\seqname{A}]$ covers $C$.
\end{enumerate}

$\seqname{A}$ is non-degenerate iff it contains a non-degenerate chart.

By abuse of language we write $U \in \seqname{A}$ for
$U \in \pi_1[\seqname{A}]$.

Let $E$ be a topological space and $C$ a linear space. Then
\begin{equation}
\Atl[\Ck]_\Ob(E,C) \defeq
\set
{{
  (\seqname{A}, E, C)
}}%
[
  {{
    \isAtl[\Ck]_\Ob(\seqname{A}, E, C}
  )}
]
\end{equation}
\begin{equation}
\full{\Atl[\Ck]_\Ob}(E,C) \defeq
\set
{{
  (\seqname{A}, E, C)
}}%
[
  {{
    \full{\isAtl[\Ck]_\Ob}(\seqname{A}, E, C}
  )}
]
\end{equation}

Let $\seqname{E}$ be a set of topological spaces and $\seqname{C}$ a
set of linear spaces. Then
\begin{equation}
\Atl[\Ck]_\Ob(\seqname{E}, \seqname{C}) \defeq
  \union
    [
        {E_\mu \in \seqname{E}},
        {C_\mu \in \seqname{C}}
    ]
    {{
      \Atl_\Ob(E_\mu, C_\mu)
    }}
\end{equation}
\begin{equation}
\full{\Atl[\Ck]_\Ob}(\seqname{E}, C) \defeq
\set
{{
  (\seqname{A}, E, C)
}}%
[
  {{
    \full{\isAtl[\Ck]_\Ob}(\seqname{A}, E, C}
  )}
]
\end{equation}
\end{definition}

\begin{lemma}[\ensuremath{\Ck}-atlases]
\label{lem:Ck-atlas}
Let $E$ be a topological space, $C$ be a linear space and $\seqname{A}$
be a $\Ck$-atlas of $E$ in the coordinate space $C$.

If $\seqname{A}$ is non-degenerate then $C$ is non-degenerate.

\begin{proof}
Let $(U,V,\phi)$ be a non-degenerate chart of $\seqname{A}$. $V$
contains a ball and is contained in $C$.
\end{proof}

If $C$ is non-degenerate and $\seqname{A}$ is full then $\seqname{A}$ is
non-degenerate.

\begin{proof}
Let $B$ be a ball in $C$ with center $v$. Since $\seqname{A}$ is full,
it contains a chart $(U,V,\phi)$ with $v \in V$. Then $V$ contains a
ball with center $v$.
\end{proof}
\end{lemma}

\begin{definition}[Compatibility of charts with \ensuremath{\Ck}-atlases]
\label{def:CkCompat}
A chart $(U, V, \phi)$ of $E$ in the coordinate space $C$ is $\Ck$
compatible with a $\Ck$-atlas $\seqname{A}$ iff it is $\Ck$ compatible with
every chart in the atlas.
\end{definition}

\begin{lemma}[Compatibility of subcharts with \ensuremath{\Ck}-atlases]
\label{lem:CkCompat}
Let $\seqname{A}$ be a $\Ck$-atlas of $E$ in the coordinate space $C$
and $\seqname{C}_1 = (U_1, V_1, \phi_1)$ a $\Ck$ chart in
$\seqname{A}$. Then any subchart of $\seqname{C}_1$ is $\Ck$ compatible
with $\seqname{A}$.

\begin{proof}
Let $\seqname{C}' = (U', V', \phi')$ be a subchart of $\seqname{C}_1$ and
$\seqname{C}_2 = (U_2, V_2, \phi_2)$ another chart in $\seqname{A}$.
\begin{enumerate}
\item If $U_1 \cap U_2 = \emptyset$, then $U' \cap U_2 = \emptyset$.
\item If $U' \cap U_2 = \emptyset$ then $\seqname{C}'$ is $\Ck$
compatible with $\seqname{C}_2$.
\item Otherwise the transition function
$t^1_2 \defeq \phi_2 \compose \phi^{-1}_1 \restrictto_{\phi_1[U_1 \cap U_2]}$
is a $\Ck$ diffeomorphism and thus
$t^1_2 \restrictto_{\phi_1[U' \cap U_2]}$ is a $\Ck$ diffeomorphism.
\end{enumerate}
\end{proof}
\end{lemma}

\begin{lemma}[Extensions of \ensuremath{\Ck}-atlases]
\label{Ck-atl:extensions}
Let $\seqname{A}$ be a $\Ck$ atlas of $\seqname{E}$ in the coordinate
space $C$ and $(U_i,V_i,\phi_i)$, $i=1,2$, be a $\Ck$ chart of $E$ in
the coordinate space $C$ $\Ck$ compatible with $\seqname{A}$ in the
coordinate space $C$. Then $(U_1,V_1,\phi_1)$ is $\Ck$ compatible with
$(U_2,V_2,\phi_2)$ in the coordinate space $C$.

\begin{proof}
If $U_1 \cap U_2 = \emptyset$ then $(U_1,V_1,\phi_1)$ is
$\Ck$ compatible with $(U_2,V_2,\phi_2)$. Otherwise,
$\phi_2 \compose \phi^{-1}_1 \restrictto_{\phi_1[U_1 \cap U_2]} \maps
\phi_1[U_1 \cap U_2] \toiso \phi_2[U_1 \cap U_2]$
is a homeomorphism.  It remains to show that
$\phi_2 \compose \phi^{-1}_1 \restrictto_{\phi_1[U_1 \cap U_2]}$
is a $\Ck$ diffeomorphism.
Let $(U'_\alpha,V'_\alpha,\phi'_\alpha)$, $\alpha \prec \Alpha$, be
charts in $\seqname{A}$ such that
$U_1 \cap U_2 \subseteq \union[\alpha \prec \Alpha]{U'_\alpha}$ and
$U_1 \cap U_2 \cap U'_\alpha \neq \emptyset$, $\alpha \prec \Alpha$.
Since the charts are $\Ck$ compatible with $(U'_\alpha,V'_\alpha,\phi'_\alpha)$,
$\phi_2 \compose \phi'^{-1}_\alpha \restrictto_{U_1 \cap U_2 \cap U'_\alpha}$
and
$\phi'_\alpha \compose \phi^{-1}_1 \restrictto_{U_1 \cap U_2 \cap U'_\alpha}$
are $\Ck$ diffeomorphisms and thus
$\phi_2 \compose \phi^{-1}_1 = \phi_2 \compose \phi'^{-1}_\alpha \compose \phi'_\alpha \compose \phi^{-1}_1$
is a $\Ck$ diffeomorphism.
\end{proof}
\end{lemma}

\begin{definition}[Maximal \ensuremath{\Ck}-atlases]
\label{def:Ck-ATLmax}
Let $E$ be a topological space, $C$ be a linear space and $\seqname{A}$
be a non-degenerate $\Ck$-atlas of $E$ in the coordinate space $C$.

$\seqname{A}$ is a maximal $\Ck$-atlas of $E$ in the coordinate space
$C$, abbreviated \\*
$\maximal{\isAtl[\Ck]_\Ob}(\seqname{A}, E,C)$, iff $\seqname{A}$ is a
$\Ck$-atlas that cannot be extended by adding an additional $\Ck$
compatible chart.

\begin{equation}
  \maxfull{\isAtl[\Ck]_\Ob}(\seqname{A}, E, C) \defeq
    \full{\isAtl[\Ck]_\Ob}(\seqname{A}, E, C) \land
    \maximal{\isAtl[\Ck]_\Ob}(\seqname{A}, E, C)
\end{equation}

$\seqname{A}$ is a semi-maximal $\Ck$-atlas of $E$ in
the coordinate space $\seqname{C}$, abbreviated
$\maximal[S-]{\isAtl[\Ck]_\Ob}(\seqname{A}, E, \seqname{C})$, iff whenever
$(U,V,\phi) \in \seqname{A}$, $U' \subseteq U, V' \subseteq V$ and
$V'' \subseteq C$ are open, $\phi[U'] = V'$ and
$\phi' \maps V' \toiso V''$ is a $\Ck$ diffeomorphism then
$(U', V'', \phi' \compose \phi) \in \seqname{A}$.

\begin{equation}
  \maxfull[S-]{\isAtl[\Ck]_\Ob}(\seqname{A}, E, C) \defeq
    \full{\isAtl[\Ck]_\Ob}(\seqname{A}, E, C) \land
    \maximal[S-]{\isAtl[\Ck]_\Ob}(\seqname{A}, E, C)
\end{equation}
\end{definition}

\begin{lemma}[Maximal $\Ck$-atlases are semi-maximal $\Ck$-atlases]
Let $E$ be a topological space, $C$ a $\Ck$-linear space and
$\seqname{A}$ a maximal $\Ck$-atlas of $E$ in the coordinate space $C$.
Then $\seqname{A}$ is a semi-maximal $\Ck$-atlas of $E$ in the
coordinate space $C$.

\begin{proof}
Let $(U,V,\phi) \in \seqname{A}$, $U' \subseteq U, V' \subseteq V$ and
$V'' \subseteq C$ be open, $\phi[U'] = V'$ and
$\phi' \maps V' \toiso V''$ be a $\Ck$ diffeomorphism.  $(U', V', \phi)$
is a subchart of $(U,V,\phi)$ and by \pagecref{lem:CkCompat} is
$\Ck$ compatible with the charts of $\seqname{A}$. Since $\phi'$ is a
$\Ck$ diffeomorphism, $(U', V'', \phi' \compose \phi)$ is $\Ck$
compatible with the charts of $\seqname{A}$.  Since $\seqname{A}$ is
maximal, $(U', V'', \phi' \compose \phi)$ is a chart of $\seqname{A}$.
\end{proof}
\end{lemma}

\begin{theorem}[Existence and uniqueness of maximal \ensuremath{\Ck}-atlases]
Let $\seqname{A}$ be a $\Ck$-atlas of $E$ in the coordinate space $C$.
Then there exists a unique maximal $\Ck$-atlas
$\maximal{\Atlas^\Ck}(\seqname{A}, E, C)$ of $E$ in the coordinate space $C$
compatible with $\seqname{A}$.

\begin{proof}
Let $\seqname{P}$ be the set of all $\Ck$-atlases of $\seqname{E}$ in
the coordinate space $C$ containing $\seqname{A}$ and $\Ck$ compatible
in the coordinate space $\seqname{C}$ with all of the $\Ck$ charts in
$\seqname{A}$. Let $\maximal{\seqname{P}}$ be a maximal chain of
$\seqname{A}$. Then $A'=\union{\maximal{\seqname{P}}}$ is a maximal
$\Ck$ atlas of $\seqname{E}$ in the coordinate space $\seqname{C}$
$\Ck$ compatible with $\seqname{A}$.
Uniqueness follows from \pagecref{Ck-atl:extensions}.
\end{proof}
\end{theorem}

\begin{definition}[Sets of $\Ck$-atlases]
\label{def:Ck-ATLsets}

Let $E$ be a topological space and $C$ be a linear space.

\begin{multline}
\maxfull[*]{\Atl[\Ck]_\Ob(E, C)} \defeq
\set
{{
   (\seqname{A}, E, C)
}}%
[
  {
    \maxfull[*]
    {
      \isAtl[\Ck]_\Ob
      (\seqname{A}, E, C)
    }
  }
]
\end{multline}

Let $\seqname{E}$ be a set of topological spaces and $\seqname{C}$ a
set of linear spaces. Then

\begin{multline}
\maxfull[*]{\Atl[\Ck]_\Ob (\seqname{E}, \seqname{C})} \defeq
  \union
    [
        {E \in \seqname{E}},
        {C \in \seqname{C}}
    ]
    {{
      \maxfull[*]{\Atl[\Ck]_\Ob(E, C)}
    }}
\end{multline}

\end{definition}

\section{$\Ck$-atlas (near) morphisms and functors}
\label{sec:ckmorph}
This section introduces a taxonomy for morphisms between $\Ck$-atlases,
defines categories of $\Ck$-atlases, defines functors amomg them,
defines functors between them and categories of m-atlases, defines
inverse functors and proves some basic reasults.

\subsection{$\Ck$-atlas (near) morphisms}
This subsection introduces the notions of $\Ck$-atlas near
morphisms and of $\Ck$-atlas morphisms, and proves some basic
results. $\Ck$-atlas near morphisms between semi-maximal atlases
will be proven to be $\Ck$-atlas morphisms.

\subsubsection{Definitions of \ensuremath{\Ck}-atlas (near) morphisms}
\begin{definition}[\ensuremath{\Ck}-atlas near morphisms]
\label{def:Ck-ATLmorphnear}
Let $E^i$, $i=1,2$, be a topological space, $C^i$ a linear space,
$\seqname{A}^i$ a $\Ck$-atlases of $E^i$ in the coordinate space $C^i$ and
$\funcseqname{f} \defeq (\funcname{f}_0, \funcname{f}_1)$ a
pair\footnote{
  The conventional definition uses only the first of the two
  functions and a slightly different compatibility condition.
}
of functions.

$\funcseqname{f}$ is an $E^1$-$E^2$ $\Ck$ near morphism of
$\seqname{A}^1$ to $\seqname{A}^2$ in the coordinate spaces $C^1$,
$C^2$, abbreviated as
$
\isAtl[\Ck,near]_\Ar
  (
    \seqname{A}^1,
    E ^1,
    C^1,
    \seqname{A}^2,
    E^2,
    C^2,
    \funcname{f}_0,
    \funcname{f}_1
  )
$,
iff
\begin{enumerate}
\item $\funcname{f}_0 \maps E^1 \to E^2$ is a continuous function.
\item $\funcname{f}_1 \maps C^1 \to C^2$ is a $\Ck$ function.
\item for any
$(U^i, V^i, \phi^i \maps U^i \toiso V^i) \in \seqname{A}^i$, $i=1,2$,
diagram  \eqref{eq:AtlN1} in
{
  \showlabelsinline
  \pagecref{def:M-ATLmorphnear}
}
is $\Ck$-locally nearly commutative in $\seqname{C}^2$, i.e., for any
$u^1 \in I$ there are open sets $U'^1 \subseteq I$, $V'^1 \subseteq V^1$,
$U'^2 \subseteq U^2$, $V'^2 \subseteq V^2$,
$\hat{V}'^2 \subseteq C^2$ and a $\Ck$ diffeomorphism
$\hat{\funcname{f}} \maps \hat{V}'^2 \toiso V'^2$ such that
\crefrange{eq:xin}{eq:fphigeqgphi} on
\cpagerefrange{eq:xin}{eq:fphigeqgphi} in
{
  \showlabelsinline
  \pagecref{def:M-ATLmorph} hold.
}
\end{enumerate}

$\funcseqname{f}$ is also a
full (semi-maximal, maximal, full semi-maximal, full maximal)
$E^1$-$E^2$ $\Ck$ near morphism of
$\seqname{A}^1$ to $\seqname{A}^2$ in the coordinate spaces $C^1$,
$C^2$, abbreviated as
$
  \maxfull[**]
    {
      \isAtl[\Ck,near]_\Ar
        (
          \seqname{A}^1, E^1, C^1,
          \seqname{A}^2, E^2, C^2,
          \funcname{f}_0, \funcname{f}_1
        )
    }
$,
iff \\*
$
  \maxfull[**]
    {
      \isAtl[\Ck,near]_\Ob
        (
          \seqname{A}^1, E^1, C^1
        )
    }
  \land
  \maxfull[**]
    {
      \isAtl[\Ck,near]_\Ob
        (
          \seqname{A}^2, E^2, C^2
        )
    }
$

The identity morphism of $(\seqname{A}^i, E^i, C^i)$ is
\begin{equation}
\Id[{{(\seqname{A}^i, E^i, C^i)}}] \defeq
\bigl (
  (\Id[{E^i}], \Id[{C^i}]),
  (\seqname{A}^i, E^i, C^i),
  (\seqname{A}^i, E^i, C^i)
\bigr )
\end{equation}
This nomenclature will be justified later.

Let $E^i$, $i=1,2$, be topological spaces and $C^i$ be linear
spaces. Then

\begin{multline}
\maxfull[*]{\Atl[\Ck,near]_\Ar}(E^1, C^1, E^2, C^2) \defeq
  \set
  {{
    \bigl (
      (
        \funcname{f}_0,
        \funcname{f}_1
      ),
      (\seqname{A}^1, E^1, C^1),
      (\seqname{A}^2, E^2, C^2)
    \bigr )
  }}%
  [{{
    \maxfull[*]{\isAtl[\Ck,near]_\Ar}
    (
      \seqname{A}^1,
      E^1,
      C^1,
      \seqname{A}^2,
      E^2,
      C^2,
      \funcname{f}_0,
      \funcname{f}_1
    )
  }}]*
\end{multline}

The identity morphism of $(\seqname{A}^i, E^i, C^i)$ is
\begin{equation}
\Id[{{(\seqname{A}^i, E^i, C^i)}}] \defeq
\bigl (
  (\Id[{E^i}], \Id[{C^i}]),
  (\seqname{A}^i, E^i, C^i),
  (\seqname{A}^i, E^i, C^i)
\bigr )
\end{equation}
This nomenclature will be justified later.
\end{definition}

\begin{definition}[\ensuremath{\Ck}-atlas morphisms]
\label{def:Ck-ATLmorph}
Let $E^i$, $i=1,2$, be a topological space, $C^i$ a linear space,
$\seqname{A}^i$ a $\Ck$-atlases of $E^i$ in the coordinate space $C^i$ and
$\funcseqname{f} \defeq (\funcname{f}_0, \funcname{f}_1)$ a
pair\footnote{
  The conventional definition uses only the first of the two
  functions and a slightly different compatibility condition.
}
of functions.

$\funcseqname{f}$ is an $E^1$-$E^2$ $\Ck$-morphism of $\seqname{A}^1$ to
$\seqname{A}^2$ in the coordinate spaces $C^1$, $C^2$, abbreviated as
$\isAtl[\Ck]_\Ar(\seqname{A}^1, E^1, C^1$,
                  $\seqname{A}^2$, $E^2$, $C^2$,
                  $\funcname{f}_0$, $\funcname{f}_1)$,
iff
\begin{enumerate}
\item $\funcname{f}_0 \maps E^1 \to E^2$ is a continuous function.
\item $\funcname{f}_1 \maps C^1 \to C^2$ is a $\Ck$ function.
\item for any $u^1 \in E^1$, any chart
$(U^1, V^1, \phi^1 \maps U^1 \toiso V^1) \in \seqname{A}^1$ at $u^1$
and any chart
$(U^2, V^2, \phi^2 \maps U^2 \toiso V^2) \in \seqname{A}^2$ at
$\funcname{f}_0(u^1)$ there exists a subchart \\*
$
  (U'^1, V'^1, \phi^1 \maps U'^1 \toiso V'^1)
  \in \seqname{A}^1
$
at $u^1$, an open set $U'^2 \subseteq U^2$ and a chart
$
  (U'^2 ,\hat{V}'^2, \phi'^2 \maps U'^2 \toiso \hat{V}'^2)
  \in \seqname{A}^2
$
at $\funcname{f}_0(u^1)$ such that
$\funcname{f}_0[U'^1] \subseteq U'^2$ and diagram \eqref{eq:AtlM1} in
\pagecref{def:M-ATLmorph} is commutative, as shown in \pagecref{fig:AtlM1}.
\end{enumerate}

$\funcseqname{f}$ is also a
full (semi-maximal, maximal, full semi-maximal, full maximal)
$E^1$-$E^2$ $\Ck$ morphism of $\seqname{A}^1$ to $\seqname{A}^2$ in the
coordinate spaces $C^1$, $C^2$, abbreviated as \\*
$
  \maxfull[**]
    {
      \isAtl[\Ck]_\Ar
        (
          \seqname{A}^1, E^1, C^1,
          \seqname{A}^2, E^2, C^2,
          \funcname{f}_0, \funcname{f}_1
        )
    }
$,
iff \\*
$
  \maxfull[**]
    {
      \isAtl[\Ck]_\Ob
        (
          \seqname{A}^1, E^1, C^1
        )
    }
  \land
  \maxfull[**]
    {
      \isAtl[\Ck]_\Ob
        (
          \seqname{A}^2, E^2, C^2
        )
    }
$

The identity morphism of $(\seqname{A}^i, E^i, C^i)$ is
\begin{equation}
\Id[{{(\seqname{A}^i, E^i, C^i)}}] \defeq
\bigl (
  (\Id[{E^i}], \Id[{C^i}]),
  (\seqname{A}^i, E^i, C^i),
  (\seqname{A}^i, E^i, C^i)
\bigr )
\end{equation}
This nomenclature will be justified later.

Let $E^i$, $i=1,2$, be topological spaces and $C^i$ be linear
spaces. Then
\begin{multline}
\Atl[\Ck]_\Ar(E^1, C^1, E^2, C^2) \defeq
  \set
  {{
    \bigl (
      (
        \funcname{f}_0,
        \funcname{f}_1
      ),
      (\seqname{A}^1, E^1, C^1),
      (\seqname{A}^2, E^2, C^2)
    \bigr )
  }}%
  [{{
    \isAtl[\Ck]_\Ar
    (
      \seqname{A}^1,
      E^1,
      C^1,
      \seqname{A}^2,
      E^2,
      C^2,
      \funcname{f}_0,
      \funcname{f}_1
    )
  }}]*
\end{multline}
\begin{multline}
\full{\Atl[\Ck]_\Ar}(E^1, C^1, E^2, C^2) \defeq
  \set
  {{
    \bigl (
      (
        \funcname{f}_0,
        \funcname{f}_1
      ),
      (\seqname{A}^1, E^1, C^1),
      (\seqname{A}^2, E^2, C^2)
    \bigr )
  }}%
  [{{
    \full{\isAtl[\Ck]_\Ar}
    (
      \seqname{A}^1,
      E^1,
      C^1,
      \seqname{A}^2,
      E^2,
      C^2,
      \funcname{f}_0,
      \funcname{f}_1
    )
  }}]*
\end{multline}
\begin{multline}
\maximal{\Atl[\Ck]_\Ar}(E^1, C^1, E^2, C^2) \defeq
  \set
  {{
    \bigl (
      (
        \funcname{f}_0,
        \funcname{f}_1
      ),
      (\seqname{A}^1, E^1, C^1),
      (\seqname{A}^2, E^2, C^2)
    \bigr )
  }}%
  [{{
    \maximal{\isAtl[\Ck]_\Ar}
    (
      \seqname{A}^1,
      E^1,
      C^1,
      \seqname{A}^2,
      E^2,
      C^2,
      \funcname{f}_0,
      \funcname{f}_1
    )
  }}]*
\end{multline}
\begin{multline}
\maximal[S-]{\Atl[\Ck]_\Ar}(E^1, C^1, E^2, C^2) \defeq
  \set
  {{
    \bigl (
      (
        \funcname{f}_0,
        \funcname{f}_1
      ),
      (\seqname{A}^1, E^1, C^1),
      (\seqname{A}^2, E^2, C^2)
    \bigr )
  }}%
  [{{
    \maximal[S-]{\isAtl[\Ck]_\Ar}
    (
      \seqname{A}^1,
      E^1,
      C^1,
      \seqname{A}^2,
      E^2,
      C^2,
      \funcname{f}_0,
      \funcname{f}_1
    )
  }}]*
\end{multline}
\begin{multline}
\maxfull{\Atl[\Ck]_\Ar}(E^1, C^1, E^2, C^2) \defeq
  \set
  {{
    \bigl (
      (
        \funcname{f}_0,
        \funcname{f}_1
      ),
      (\seqname{A}^1, E^1, C^1),
      (\seqname{A}^2, E^2, C^2)
    \bigr )
  }}%
  [{{
    \maxfull{\isAtl[\Ck]_\Ar}
    (
      \seqname{A}^1,
      E^1,
      C^1,
      \seqname{A}^2,
      E^2,
      C^2,
      \funcname{f}_0,
      \funcname{f}_1
    )
  }}]*
\end{multline}
\begin{multline}
\maxfull[S-]{\Atl[\Ck]_\Ar}(E^1, C^1, E^2, C^2) \defeq
  \set
  {{
    \bigl (
      (
        \funcname{f}_0,
        \funcname{f}_1
      ),
      (\seqname{A}^1, E^1, C^1),
      (\seqname{A}^2, E^2, C^2)
    \bigr )
  }}%
  [{{
    \maxfull[S-]{\isAtl[\Ck]_\Ar}
    (
      \seqname{A}^1,
      E^1,
      C^1,
      \seqname{A}^2,
      E^2,
      C^2,
      \funcname{f}_0,
      \funcname{f}_1
    )
  }}]*
\end{multline}
\end{definition}

\begin{definition}[Equivalence of \ensuremath{\Ck}-atlas (near) morphisms]
\label{def:Ck-ATLmorphequiv}
Let $E^i$, $i=1,2$, be a topological space, $C^i$ a linear space,
$\seqname{A}^i$ a $\Ck$-atlases of $E^i$ in the coordinate space $C^i$ and
$\funcseqname{f} \defeq (\funcname{f}_0, \funcname{f}_1)$ and
$\funcseqname{g} \defeq (\funcname{g}_0, \funcname{g}_1)$ $E^1$-$E^2$
$\Ck$ (near) morphisms of $\seqname{A}^1$ to $\seqname{A}^2$ in the
coordinate spaces $C^1$, $C^2$.

$\funcseqname{f}$ is $\Ck$-equivalent to $\funcseqname{g}$ iff
$\funcname{f}_0 = \funcname{g}_0$.
\end{definition}

\subsubsection{Proclamations on \ensuremath{\Ck}-atlas (near) morphisms}
\begin{lemma}[\ensuremath{\Ck}-atlas (near) morphisms]
\label{lem:Ck-ATLmorph}
Let $\seqname{E}$ be a set of topological spaces, $\seqname{C}$ a set of
$\Ck$ linear spaces, $E^i \in \seqname{E}$, $i=1,2$,
$C^i \in \seqname{C}$, $\seqname{A}^i$ a $\Ck$-atlas of $E^i$ in the
coordinate space $C^i$ and
$
  \funcseqname{f} \defeq
  (
    \funcname{f}_0 \maps E^1 \to E^2,
    \funcname{f}_1 \maps C^1 \to C^2
  )
$
a pair of functions.

$\funcseqname{f}$ is an $E^1$-$E^2$ $\Ck$ (near) morphism of
$\seqname{A}^1$ to $\seqname{A}^2$ in the coordinate spaces $C^1$, $C^2$
iff $\funcseqname{f}$ is a strict
$\Trivcat{\seqname{E}}$-%
$\Trivcat{\seqname{E}}$-%
$\Triv{E^1}$-$\Triv{E^2}$-%
$\optriv[\Ck-]{\seqname{C}}$-%
$\optriv[\Ck-]{\seqname{C}}$
m-atlas (near) morphism of $\seqname{A}^1$ to $\seqname{A}^2$ in the
coordinate spaces $\Triv[\Ck-]{\seqname{C}^1}$,
$\Triv[\Ck-]{\seqname{C}^2}$.

\begin{proof}
The model neighborhoods of $\Triv[\Ck-]{\seqname{C}^i}$ are the open sets
of $\seqname{C}^i$,
$\Trivcat{\seqname{E}} \subcat[full-] \Trivcat{\seqname{E}}$,
the morphisms of $\Trivcat{\seqname{E}}$ are the continuous
functions between spaces in $\seqname{E}$,
$\Trivcat[\Ck-]{\seqname{C}} \subcat[full-] \Trivcat[\Ck-]{\seqname{C}}$,
and the morphisms of $\Trivcat[\Ck-]{\seqname{C}}$ are the $\Ck$
functions between spaces in $\seqname{C}$.
\end{proof}

$\funcseqname{f}$ is an $E^1$-$E^2$ $\Ck$-morphism of $\seqname{A}^1$ to
$\seqname{A}^2$ in the coordinate spaces $C^1$, $C^2$ iff
$\funcseqname{f}$ is a strict
$E^1$-$E^2$-%
$\Trivcat[\Ck-]{\seqname{C}}$-%
$\Trivcat[\Ck-]{\seqname{C}}$
m-atlas morphism of $\seqname{A}^1$ to $\seqname{A}^2$ in the coordinate
spaces $\Triv[\Ck-]{\seqname{C}^1}$, $\Triv[\Ck-]{\seqname{C}^2}$.

\begin{proof}
The model neighborhoods of $\Triv[\Ck-]{\seqname{C}^i}$ are the open sets
of $\seqname{C}^i$,
$\Trivcat{\seqname{E}} \subcat[full-] \Trivcat{\seqname{E}}$,
$\Trivcat[\Ck-]{\seqname{C}} \subcat[full-] \Trivcat[\Ck-]{\seqname{C}}$
and the morphisms of $\Trivcat[\Ck-]{\seqname{C}}$ are the $\Ck$
functions between spaces in $\seqname{C}$.
\end{proof}
\end{lemma}

\begin{corollary}[\ensuremath{\Ck}-atlas (near) morphisms]
\label{cor:Ck-ATLmorph}
Let $E^i$, $i=1,2$, be a topological space, $C^i$ a linear space,
$\seqname{A}^i$ a semi-maximal $\Ck$-atlas of $E^i$ in the coordinate
space $C^i$ and
$\funcseqname{f} \defeq (\funcname{f}_0, \funcname{f}_1)$ an $E^1$-$E^2$
$\Ck$ near morphism of $\seqname{A}^1$ to $\seqname{A}^2$ in the
coordinate spaces $C^1$, $C^2$. Then $\funcseqname{f}$ is an $E^1$-$E^2$
$\Ck$ morphism of $\seqname{A}^1$ to $\seqname{A}^2$ in the coordinate
spaces $C^1$, $C^2$.

\begin{proof}
The result follows from \pagecref{lem:M-ATLmorph}.
\end{proof}

Let $E^i$, $i=1,2,3$, be a topological space, $C^i$ a linear space,
$\seqname{A}^i$ a $\Ck$-atlas of $E^i$ in the coordinate space $C^i$ and
$(\funcname{f}^i_0, \funcname{f}^i_1)$ an $E^i$-$E^{i+1}$
$\Ck$-morphism of $\seqname{A}^i$ to $\seqname{A}^{i+1}$ in the
coordinate spaces $C^i$, $C^{i+1}$. Then
$
  (
    \funcname{f}^2_0 \compose \funcname{f}^1_0,
    \funcname{f}^2_1 \compose \funcname{f}^1_1
  )
$
is a $E^1$-$E^3$ $\Ck$-morphism of $\seqname{A}^1$ to $\seqname{A}^3$ in
the coordinate spaces $C^1$, $C^3$.

\begin{proof}
The result follows from
\cref{M-morphcomp:morphcompmorph} of
{
  \showlabelsinline
  \fullcref{lem:M-ATLmorphcomp} on
  \cpageref{M-morphcomp:morphcompmorph}.
}
\end{proof}
\end{corollary}

\begin{lemma}[Composition of equivalent \ensuremath{\Ck}-atlas (near) morphisms]
\label{lem:Ck-ATLmorphequiv}
Let $E^i$, $i=1,2,3$, be a topological space,
$C^i$ be a $\Ck$ linear space,
$\seqname{A}^i$ be a $\Ck$-atlas of $E^i$ in the coordinate space $C^i$ and
$
  \funcseqname{f}^i \defeq
  (
    \funcname{f}^i_0 \maps \seqname{E}^i \to \seqname{E}^{i+1},
    \funcname{f}^i_1 \maps \seqname{C}^i \to \seqname{C}^{i+1}
  )
$
and \\*
$
  \funcseqname{g}^i \defeq
  (
    \funcname{g}^i_0 \maps \seqname{E}^i \to \seqname{E}^{i+1},
    \funcname{g}^i_1 \maps \seqname{C}^i \to \seqname{C}^{i+1}
  )
$
be $\Ck$-equivalent $\seqname{E}^i$-$\seqname{E}^{i+1}$ m-atlas (near)
morphisms of $\seqname{A}^i$ to $\seqname{A}^{i+1}$ in the coordinate
spaces $\seqname{C}^i$, $\seqname{C}^{i+1}$.

Then $\funcseqname{f}^2 \compose[()] \funcname{f}^1$ is $\Ck$-equivalent
to $\funcseqname{g}^2 \compose[()] \funcname{g}^1$.

\begin{proof}
$\funcname{f}^1_0 = \funcname{g}^1_0$ and
$\funcname{f}^2_0 = \funcname{g}^2_0$, hence
$
  \funcname{f}^2_0 \compose \funcname{f}^1_0 =
  \funcname{g}^2_0 \compose \funcname{g}^1_0
$.
\end{proof}
\end{lemma}

\subsection{Categories of $\Ck$ atlases and functors}

\subsubsection{Categories of $\Ck$ atlases}
This subsubsection defines categories of $\Ck$-atlases with $\Ck$-atlas
morphisms as morphisms.  It does not define categories of $\Ck$-atlases
with $\Ck$-atlas near morphisms as morphisms becaue the composition of
two $\Ck$-atlas near morphisms is not in general a $\Ck$-atlas near
morphism, and requiring the atlasses to be semi-maximal would cause all
$\Ck$-atlas near morphisms to be $\Ck$-atlas morphisms.

\begin{definition}[Categories of $\Ck$ atlases]
\label{def:Ck-ATLcat}
Let $\seqname{E}$ be a set of topological spaces and $\seqname{C}$ a
set of linear spaces.
Let $P \defeq \seqname{E} \times \seqname{C}$. Then
\begin{equation}
  \maxfull[*]
    {
      \Atl[\Ck]_\Ar(\seqname{E}, \seqname{C})
    }
  \defeq
    \union
      [
          {(E^\mu, C^\mu) \in \seqname{P}},
          {(E^\nu, C^\nu) \in \seqname{P}}
      ]
      {{
        \maxfull[*]
        {
          \Atl[\Ck]_\Ar
          (
            E^\mu,
            C^\mu,
            E^\nu,
            C^\nu
          )
        }
      }}
\end{equation}
\begin{equation}
  \maxfull[*]
    {
      \Atl[\Ck](\seqname{E}, \seqname{C})
    }
  \defeq
  \Bigl (
    \maxfull[*]
      {
        \Atl[\Ck]_\Ob(\seqname{E}, \seqname{C})
      },
    \maxfull[*]
      {
        \Atl[\Ck]_\Ar(\seqname{E}, \seqname{C})
      },
    \compose[A]
  \Bigr )
\end{equation}

\begin{remark}
It is pointless to define categories of near morphisms, since \\*
$
  \maximal[**]
  {
    is\Atl[\Ck,near]_\Ar
    (
      E^\mu,
      C^\mu,
      E^\nu,
      C^\nu
    )
  }
\iff
  \maximal[**]
  {
    is\Atl[\Ck]_\Ar
    (
      E^\mu,
      C^\mu,
      E^\nu,
      C^\nu
    )
  }
$.
\end{remark}
\end{definition}

\begin{lemma}[\ensuremath{\Atl[\Ck](\seqname{E}, \seqname{C}} is a category]
\label{lem:Ck-ATLiscat}
Let $\seqname{E}$ be a set of topological spaces and $\seqname{C}$ a set
of linear spaces. Then each of
$\maxfull[*]{\Atl[\Ck]}(\seqname{E}, \seqname{C})$ is a category.

Let
$
  (\seqname{A}^i, E^i, C^i)
  \in
  \Atl[\Ck]_\Ob(\seqname{E}, \seqname{C})
$.
Then $\Id[{{(\seqname{A}^i, E^i, C^i)}}]$ is the
identity morphism for $(\seqname{A}^i, E^i, C^i)$.

\begin{proof}
Let $(\seqname{A}^i,E^i,C^i)$, $i \in [1,3]$,
be an object of $\Atl[\Ck](\seqname{E}, \seqname{C})$ and
let
$
  \funcname{m}^i \defeq \\
  \bigl (
    (\funcname{f}_0^i,\funcname{f}_1^i),
    (\seqname{A}^i,E^i,C^i),
    (\seqname{A}^{i+1},E^{i+1},C^{i+1})
  \bigr )
$, $i=1,2$,
be a morphism of $\Atl[\Ck](\seqname{E}, \seqname{C})$. Then
\begin{description}
\item[Composition:]
$
  \bigl (
    (
      \funcname{f}_0^2 \compose \funcname{f}_0^1,
      \funcname{f}_1^2 \compose \funcname{f}_1^1
    ),
    (\seqname{A}^1,E^1,C^1),
    (\seqname{A}^3,E^3,C^3)
  \bigr )
$
is a morphism of \\*
$\Atl[\Ck](\seqname{E}, \seqname{C})$ by
\pagecref{cor:Ck-ATLmorph}.

\item[Associativity:]
Composition is associative by \pagecref{lem:labcomp}.

\item[Identity:]
$\Id[{{(\seqname{A}^i, E^i, C^i)}}]$ is an identity
morphism by \cref{lem:labcomp}.
\end{description}
\end{proof}
\end{lemma}

\subsubsection{\ensuremath{\Ck} atlas functors}
\begin{definition}[Functors from \ensuremath{\Ck} atlases to m-atlases]
\label{def:Func-Ck-M}
Let $E^i$, $i=1,2$, be a topological space, $C^i$ a linear space,
$\seqname{A}^i$ $\Ck$-atlases of $E^i$ in the coordinate space $C^i$,
$\funcname{f}_0 \maps E^1 \to E^2$ continuous and
$\funcname{f}_1 \maps C^1 \to C^2$ $\Ck$. Then
\begin{equation}
\Functor^\Ck_{\Ck,\M} \bigl (\seqname{A}^i, E^i, C^i \bigl )
\defeq
  \Bigl (
    \seqname{A}^i,
    \Triv{E^i},
    \Triv[\Ck-]{C^i}
  \Bigr )
\end{equation}
\begin{multline}
\Functor^\Ck_{\Ck,\M}
  \bigl (
    (
      \funcname{f}_0,
      \funcname{f}_1
    ),
    (\seqname{A}^1, E^1, C^1),
    (\seqname{A}^2, E^2, C^2)
  \bigr )
\defeq \\
  \biggl (
    \Bigl (
      \funcname{f}_0,
      \funcname{f}_1
     \Bigr ) ,
    \Bigl ( \seqname{A}^1, \Triv{E^1}, \Triv[\Ck-]{C^1} \Bigr ) ,
    \Bigl ( \seqname{A}^2, \Triv{E^2}, \Triv[\Ck-]{C^2} \Bigr )
  \biggr )
\end{multline}
\end{definition}

\begin{theorem}[Functors from \ensuremath{\Ck} atlases to m-atlases]
\label{the:Func-Ck-M}
Let $\seqname{E}$ be a set of topological spaces and $\seqname{C}$ a
set of linear spaces. Then
$\Functor^\Ck_{\Ck,\M}$ is a functor from
$\Atl[\Ck] \bigl ( \seqname{E}, \seqname{C} \bigr )$
to
$\Atl \bigl ( \Triv{\seqname{E}}, \Triv[\Ck-]{\seqname{C}} \bigr )$

\begin{proof}
Let $o^i \defeq (\seqname{A}^i, E^i, C^i)$,
$i \in [1,3]$, be an object of $\Atl[\Ck](\seqname{E}, \seqname{C})$, \\*
$
  o'^i \defeq
  \Functor^\Ck_{\Ck,\M} o^i
  =
  \bigl (
     \seqname{A}^i,
     \Triv{E^i},
     \Triv[\Ck-]{C^i}
   \bigr )
$,
$
  \funcname{m}^i \defeq
    \bigl (
      (\funcseqname{f}^i_0, \funcseqname{f}^i_1),
      o^i,
      o^{i+1}
    \bigr )
$, $i=1,2$,
be a morphism from $o^i$ to $o^{i+1}$ and
$
  \funcname{m}'^i \defeq
  \Functor^\Ck_{\Ck,\M} \funcname{m}^i
  =
    \bigl (
      (\funcseqname{f}^i_0, \funcseqname{f}^i_1).
      o'^i,
      o'^{i+1}
    \bigr )
$,
be the corresponding morphism in
$\Atl(\Triv{\seqname{E}}, \Triv[\Ck-]{\seqname{C}})$.

\begin{description}
\item[Preservation of endpoints:]
\begin{align*}
  \Functor^\Ck_{\Ck,\M} o^i
  & =
  \bigl (
     \seqname{A}^i,
     \Triv{E^i},
     \Triv[\Ck-]{C^i}
   \bigr )
\\*
  & =
  o'^i
\\*
  \Functor^\Ck_{\Ck,\M} \funcname{m}^i
  & =
  \funcname{m}'^i
\\*
  & =
    \bigl (
      (
        \funcseqname{f}^i_0,
        \funcseqname{f}^i_1
      ),
     o'^i,
     o'^{i+1}
    \bigr )
\end{align*}
$\Functor^\Ck_{\Ck,\M} \funcname{m}^i$ is a morphism from
$\Functor^\Ck_{\Ck,\M} o^i$ to $\Functor^\Ck_{\Ck,\M} o^{i+1}$:

\item[Identity:]
\begin{align*}
  \Functor^\Ck_{\Ck,\M} \Id[{o^i}]
  & =
  \Functor^\Ck_{\Ck,\M}
    \bigl (
      (\Id[{E^1}], \Id[{C^1}]),
      o^i,
      o^i
    \bigr )
\\*
  & =
  \bigl (
    (\Id[{\Triv{E^i}}], \Id[{C^i}]),
    o'^i,
    o'^i
  \bigr )
\\*
  & =
  \bigl (
    (\Id[{\Triv{E^i}}], \Id[{C^i}]),
    \Functor^\Ck_{\Ck,\M} o^i,
    \Functor^\Ck_{\Ck,\M} o^i
  \bigr )
\\*
  & =
  \Id[{\Functor^\Ck_{\Ck,\M} o^i}]
\end{align*}

\item[Composition:]
\begin{align*}
  m^2 \compose[A] m^1
  &=
  \bigl (
    (f_0^2 \compose f_0^1, f_1^2 \compose f_1^1),
    o^1,
    o^3
  \bigr )
\\*
  \Functor^\Ck_{\Ck,\M} (m^2 \compose[A] m^1)
  &=
  \bigl (
    (f_0^2 \compose f_0^1, f_1^2 \compose f_1^1),
    o'^1,
    o'^3
  \bigr )
\\*
  &=
  \bigl (
    (f_0^2, f_1^2),
    o'^2,
    o'^3
  \bigr ) \compose[A]
  \bigl (
    (f_0^1, f_1^1),
    o'^1,
    o'^2
  \bigr )
\\*
  &=
    \Functor^\Ck_{\Ck,\M}
      \bigl (
        (f_0^2, f_1^2),
        o^2,
        o^3
      \bigr ) \compose[A]
    \Functor^\Ck_{\Ck,\M}
      \bigl (
        (f_0^1, f_1^1),
        o^1,
        o^2
      \bigr )
\\*
  &=
  \Functor^\Ck_{\Ck,\M} m^2 \compose[A] \Functor^\Ck_{\Ck,\M} m^1
\end{align*}
\end{description}
\end{proof}
\end{theorem}

\begin{definition}[Functors from m-atlases to \ensuremath{\Ck} atlases]
\label{def:Func-M-Ck}
Let $\seqname{E}^i$, $i=1,2$, be model spaces, $\seqname{C}^i$ be linear
model spaces, $\seqname{A}^i$ a maximal m-atlas of $\seqname{E}^i$ in
the coordinate space $\seqname{C}^i$ and $(\funcname{f}_0,
\seqname{f}_1)$ an $\seqname{E}^1$-$\seqname{E}^2$ m-atlas morphism of
$\seqname{A}^1$ to $\seqname{A}^2$ in the coordinate spaces
$\seqname{C}^1$, $\seqname{C}^2$, Then
\begin{equation}
\Functor^\Ck_{\M,\Ck} (\seqname{A}^i, \seqname{E}^i, \seqname{C}^i)
\defeq
  \bigl (
     \seqname{A}^i,
     \pi_1(\seqname{E}^i),
     \pi_1(\seqname{C}^i)
   \bigr )
\end{equation}
\begin{multline}
\Functor^\Ck_{\M,\Ck}
  \bigl (
    (
      \funcname{f}_0,
      \funcname{f}_1
    ),
    (\seqname{A}^1, \seqname{E}^1, \seqname{C}^1),
    (\seqname{A}^2, \seqname{E}^2, \seqname{C}^2)
  \bigr )
\defeq \\
  \Bigl (
    \bigl (
      \funcname{f}_0,
      \funcname{f}_1
    \bigr ),
    \bigl ( \seqname{A}^1, \pi_1(\seqname{E}^1), \pi_1(\seqname{C}^1) \bigr ),
    \bigl ( \seqname{A}^2, \pi_1(\seqname{E}^2), \pi_1(\seqname{C}^2) \bigr )
  \Bigr )
\end{multline}
\end{definition}

\begin{theorem}[Functors from m-atlases to \ensuremath{\Ck} atlases]
\label{the:Func-M-Ck}
Let $\seqname{E}$ be a set of model spaces and $\seqname{C}$ a
set of $\Ck$ linear model spaces. Then
$\Functor^\Ck_{\M,\Ck}$ is a functor from
$\Atl(\seqname{E}, \seqname{C})$
to
$\Atl[\Ck] \bigl (\pi_1[\seqname{E}], \pi_1[\seqname{C}] \bigr )$.

\begin{proof}
Let
$
o^i \defeq
\bigl (
  \seqname{A}^i, (E^i,\catname{E}^i), (C^i,\catname{C}^i)
\bigr )
$,
$i \in [1,3]$, be an object of $\Atl[\Ck](\seqname{E}, \seqname{C})$ and \\*
$
\funcname{m}^i \defeq
  \bigl (
    (\funcseqname{f}^i_0, \funcseqname{f}^i_1).
    o^i,
    o^{i+1}
  \bigr )
$, $i=1,2$,
be a morphism from $o^i$ to $o^{i+1}$.

$\Functor^\Ck_{\M,\Ck} (\funcname{m}^1)$ is a morphism from
$\Functor^\Ck_{\M,\Ck} o^1$ to $\Functor^\Ck_{\M,\Ck} o^2$:

\begin{equation*}
\begin{split}
  &
\Functor^\Ck_{\M,\Ck} (\funcname{m}^1) =
\\*
  &
\Functor^\Ck_{\M,\Ck}
  \bigl (
    (
      \funcseqname{f}^1_0,
      \funcseqname{f}^1_1
    ),
    (\seqname{A}^1, \seqname{E}^1, \seqname{C}^1),
    (\seqname{A}^2, \seqname{E}^2, \seqname{C}^2)
  \bigr )
=
\\*
  &
  \Bigl (
    (
      \funcseqname{f}^1_0,
      \funcseqname{f}^1_1
    ),
    \bigl (
      \seqname{A}^1, \pi_1(\seqname{E}^1), \pi_1(\seqname{C}^1)
    \bigr ),
    \bigl (
      \seqname{A}^2, \pi_1(\seqname{E}^2), \pi_1(\seqname{C}^2)
    \bigr )
  \Bigr )
\end{split}
\end{equation*}

$\Functor^\Ck_{\M,\Ck}$ maps identity functions to identity functions:

\begin{equation*}
\begin{split}
  &
\Functor^\Ck_{\M,\Ck} \Id[{{(\seqname{A}^i, \seqname{E}^i, \seqname{C}^i)}}]
=
\\*
  &
\Functor^\Ck_{\M,\Ck}
  \bigl (
    (\Id[{\seqname{E}^i}], \Id[{\seqname{C}^i}]),
    (\seqname{A}^i, \seqname{E}^i, \seqname{C}^i),
    (\seqname{A}^i, \seqname{E}^i, \seqname{C}^i)
  \bigr )
=
\\*
  &
  \bigl (
    (\Id[{\pi_1(\seqname{E}^i)}], \Id[{\pi_1(\seqname{C}^i)}]),
    (\seqname{A}^i, \pi_1(\seqname{E}^i), \pi_1(\seqname{C}^i)
    ),
    (\seqname{A}^i, \pi_1(\seqname{E}^i), \pi_1(\seqname{C}^i)
    )
  \bigr )
=
\\*
  &
  \bigl (
    (\Id[{E^i}], \Id[{C^i}]),
    \Functor^\Ck_{\M,\Ck} (\seqname{A}^i, \seqname{E}^i, \seqname{C}^i),
    \Functor^\Ck_{\M,\Ck} (\seqname{A}^i, \seqname{E}^i, \seqname{C}^i)
  \bigr )
=
\\*
  &
\Id[{{\Functor^\Ck_{\M,\Ck} (\seqname{A}^i, \seqname{E}^i, \seqname{C}^i)}}]
\end{split}
\end{equation*}

$
  \Functor^\Ck_{\M,\Ck} m^2 \compose[A] \Functor^\Ck_{\M,\Ck} m^1
  =
  \Functor^\Ck_{\M,\Ck} (m^2 \compose[A] m^1)
$:

\begin{align*}
  m^2 \compose[A] m^1
  & =
  \bigl (
    (f_0^2 \compose f_0^1, f_1^2 \compose f_1^1),
    (\seqname{A}^1, \seqname{E}^1, \seqname{C}^1),
    (\seqname{A}^3, \seqname{E}^3, \seqname{C}^3)
  \bigr )
\\*
  \Functor^\Ck_{\M,\Ck}
    \bigl (
       (\seqname{A}^i, \seqname{E}^i, \seqname{C}^i)
    \bigr )
  & =
  \bigl (
    \seqname{A}^i,
    \pi_1(\seqname{E}^i),
    \pi_1(\seqname{C}^i)
  \bigr )
\\*
  \Functor^\Ck_{\M,\Ck} (m^i)
  & =
    \Bigl (
      (f_0^i, f_1^i),
      \bigl (
        \seqname{A}^i,
        \pi_1(\seqname{E}^i),
        \pi_1(\seqname{C}^i)
      \bigr ),
      \bigl (
        \seqname{A}^{i+1},
        \pi_1(\seqname{E}^{i+1}),
        \pi_1(\seqname{C}^{i+1})
      \bigr )
    \Bigr )
\\*
  \Functor^\Ck_{\M,\Ck} m^2 \compose[A] \Functor^\Ck_{\M,\Ck} m^1
  & =
    \Bigl (
      (f_0^2 \compose f_0^1, f_1^2 \compose f_1^2),
      \bigl (
        \seqname{A}^1,
        \pi_1(\seqname{E}^1),
        \pi_1(\seqname{C}^1)
      \bigr ),
      \bigl (
        \seqname{A}^3,
        \pi_1(\seqname{E}^3),
        \pi_1(\seqname{C}^3)
      \bigr )
    \Bigr )
\\*
  \Functor^\Ck_{\M,\Ck} (m^2 \compose m^1)
  & =
    \Bigl (
      (f_0^2 \compose f_0^1, f_1^2 \compose f_1^2),
      \bigl (
        \seqname{A}^1,
        \pi_1(\seqname{E}^1),
        \pi_1(\seqname{C}^1)
      \bigr ),
      \bigl (
        \seqname{A}^3,
        \pi_1(\seqname{E}^3),
        \pi_1(\seqname{C}^3)
      \bigr )
    \Bigr )
\end{align*}

\end{proof}
\end{theorem}

\section{Associated model spaces and functors}
{
  \showlabelsinline
  \label{sec:assmodC}
}
$\Functor^\Ck_{\Ck,\M}$ is an obvious functor of $\Ck$-atlases to
m-atlases, but $\Triv{E}$ and $\Triv[\Ck-]{C}$ may have more model
neighborhhoods or more morphisms than needed for consistency with the
atlas. There exist, however, model spaces with the minimum model
neighborhoods and morphisms needed.

\begin{definition}[Coordinate model spaces associated with \ensuremath{\Ck}-atlases]
\label{def:Ck-modC}
Let $\seqname{A}^i$, $i=1,2$, be a $\Ck$-atlas of $E^i$ in
the coordinate space $C^i$, $\funcname{f}_0 \maps E^1 \to E^2$ a
continuous function and $\funcname{f}_1 \maps C^1 \to C^2$ a $\Ck$
function. Then

\begin{multline}
\Functor^{\minimal{\Ck}}_2 (\seqname{A}^i, E^i, C^i)
\defeq \\
\minimal{\Mod}
  \left (
    C^i,
    \pi_2[\seqname{A}^i],
    \set
      {\phi' \compose \phi^{-1}}%
      [
        \equant
          {
            {(U,V,\phi) \in \seqname{A}^i},
            {(U',V',\phi') \in \seqname{A}^i}
          }
          {
            U \cap U' \ne \emptyset
          }
      ]
  \right )
\end{multline}
\begin{multline}
\Functor^{\minimal{\Ck}}_2
  \bigl (
    (
      \funcname{f}_0,
      \funcname{f}_1
    ),
    (\seqname{A}^1, E^1, C^1),
    (\seqname{A}^2, E^2, C^2)
  \bigr )
\defeq \\
  \funcname{f}_1
  \maps
  \Functor^{\minimal{\Ck}}_2 (\seqname{A}^1, E^1, C^1)
  \to
  \Functor^{\minimal{\Ck}}_2 (\seqname{A}^2, E^2, C^2)
\end{multline}

The minimal coordinate $\Ck$ model space with neighborhoods in the
$\Ck$-atlas $\seqname{A}^i$ of $E^i$ in the coordinate space $C^i$ is
$\Functor^{\minimal{\Ck}}_2 (\seqname{A}^1, E^i, C^i)$.

The coordinate mapping associated with the
$\seqname{E}^1$-$\seqname{E}^2$ $\Ck$-atlas morphism
$(\funcname{f}_0,\funcname{f}_1)$ of $\seqname{A}^1$ to $\seqname{A}^2$
in the coordinate spaces $\seqname{C}^1$, $\seqname{C}^2$ is
$
  \funcname{f}_1
  \maps
  \Functor^{\minimal{\Ck}}_2 (\seqname{A}^1, E^1, C^1)
  \to
  \Functor^{\minimal{\Ck}}_2 (\seqname{A}^2, E^2, C^2)
$.
If it is a model function then it is also the
coordinate m-atlas morphism associated with the
$\seqname{E}^1$-$\seqname{E}^2$ $\Ck$-atlas morphism
$(\funcname{f}_0,\funcname{f}_1)$ of $\seqname{A}^1$ to $\seqname{A}^2$
in the coordinate spaces $\seqname{C}^1$, $\seqname{C}^2$.
\end{definition}

\begin{lemma}[Coordinate model spaces associated with \ensuremath{\Ck}-atlases]
\label{lem:Ck-ismodC}
Let $\seqname{A}^i$ be a $\Ck$-atlas of $E^i$ in the coordinate space
$C^i$, $i=1,2$. Then
$\Functor^{\minimal{\Ck}}_2 (\seqname{A}^i, E^i, C^i)$ is a $\Ck$ linear
model space.

\begin{proof}
$\Functor^{\minimal{\Ck}}_2 (\seqname{A}^i, E^i, C^i)$ satisfies the
conditions for a model space.
\begin{enumerate}
\item Since $\pi_2[\seqname{A}^i]$ is an open cover of
$\union{\pi_2[\seqname{A}^i]}$, the set of finite intersections is also an
open cover.

\item Finite intersections of finite intersections are finite intersections

\item Restrictions of continuous functions are continuous

\item If $\funcname{f} \maps A \to B$ is a morphism of
$\Functor^{\minimal{\Ck}}_2 (\seqname{A}^i, E^i, C^i)$,
$A, A', B, B'$ model meighborhoods of
$\Functor^{\minimal{\Ck}}_2 (\seqname{A}^i, E^i, C^i)$,
$A' \subseteq A$, $B' \subseteq B$ and $\funcname{f}[A'] \subseteq B'$ then since
$\funcname{f} \maps A \to B$ is a morphism it is a restriction of a transition
function between its restrictions to sets in $\pi_2[\seqname{A}^i]$ and
its restrictions are also, hence morphisms, and thus
$\funcname{f} \restrictto_{A'} A' \to B'$ is a morphism.

\item
If $(U,V,\phi) \in \seqname{A}^i$ then $\Id[{V}] = \phi \compose \phi^{-1}A$
is a transition function and hence a morphism of
$\Functor^{\minimal{\Ck}}_2 (\seqname{A}^i, E^i, C^i)$. If $A, A'$ are
objects of
$
  \pi_2
    \bigl (
      \Functor^{\minimal{\Ck}}_2
        (\seqname{A}^i, E^i, C^i)
    \bigr )
$
and $A' \subseteq A$ then the inclusion map
$\funcname{i} \maps A' \hookrightarrow A$ is a restriction of an
identity morphism of
$
  \Functor^{\minimal{\Ck}}_2
    (\seqname{A}^i, E^i, C^i)
$
and hence a morphism.

\item Restricted sheaf condition: let
\begin{enumerate}
\item $U_\alpha,V_\alpha$, $\alpha \prec \Alpha$, be an object of
$
\pi_2
  \biggl (
    \Functor^{\minimal{\Ck}}_2
      (\seqname{A}^i, E^i, C^i)
  \biggr )
$
\item $\funcname{f}_\alpha \maps U_\alpha \to V_\alpha$ be a morphism of
$
  \pi_2
    \biggl (
      \Functor^{\minimal{\Ck}}_2
        (\seqname{A}^i, E^i, C^i)
    \biggr )
$
\item $U \defeq \union[\alpha \prec \Alpha]{U_\alpha}$
\item $V \defeq \union[\alpha \prec \Alpha]{V_\alpha}$
\item $\funcname{f} \maps U \to V$ be continuous and
$
  \uquant
    {{\alpha \prec \Alpha},{x \in U_\alpha}}
    {\funcname{f}(x) = \funcname{f}_\alpha(x)}
$
\end{enumerate}
Then $\funcname{f}$ is $\Ck$ and hence a morphism of
$
  \pi_2
    \biggl (
      \Functor^{\minimal{\Ck}}_2
        (\seqname{A}^i, E^i, C^i)
    \biggr )
$.
\end{enumerate}
\end{proof}

$\Functor^{\minimal{\Ck}}_2 (\seqname{A}^i, E^i, C^i)$ is
non-degenerate iff $\seqname{A}^i$ is non-degenerate.

\begin{proof}
If $\seqname{A}^i$ is non-degenerate then by \pagecref{def:Ck-atlas}
there is at least one chart $(U,V,\phi) \in \seqname{A}$ such that $V$
contains a ball of the underlying Banach or Fr\'echet space.
$U \subseteq \union{\pi_2[\seqname{A}^i]}$, so by \pagecref{def:Ck-modC}
that ball is contained in
$\Functor^{\minimal{\Ck}}_2 (\seqname{A}^i, E^i, C^i)$.

Conversly, if $\Functor^{\minimal{\Ck}}_2 (\seqname{A}^i, E^i, C^i)$ is
non-degenerate then by \pagecref{def:lin} it contains a ball $B$ at
$v \in \Functor^{\minimal{\Ck}}_2 (\seqname{A}^i, E^i, C^i)$
with radius $r$. By \cref{def:Ck-modC} there is a chart
$(U,V,\phi) \in \seqname{A}^i$ at $v$.  Since $V$ is open in $C^i$,
there is a ball $B'$ at $v$ with rafius $r' > 0$ such that
$B' \cap C^i \subseteq V$. Then the ball at $v$ with radius $\min(r,r')$
is contained in $V$.
\end{proof}

Let $\seqname{A}^i$, $i=1,2$, be a $\Ck$-atlas of $E^i$ in the
coordinate space $C^i$, $\funcname{f}_0 \maps E^1 \to E^2$ a continuous
function, $\funcname{f}_1 \maps C^1 \to C^2$ a $\Ck$ function and
$\funcseqname{f} \defeq (\funcname{f}_0, \funcname{f}_1)$ either a
semi-maximal $\Ck$-atlas near morphism or a $\Ck$-atlas morphism from
$\seqname{A}^1$ to $\seqname{A}^2$. Then
$
  \funcname{f}_1
  \maps
  \Functor^{\minimal{\Ck}}_2 (\seqname{A}^1, E^1, C^1)
  \to
  \Functor^{\minimal{\Ck}}_2 (\seqname{A}^2, E^2, C^2)
$
is well defined, i.e.,
$
  \funcname{f}_1
    \biggl [
      \Functor^{\minimal{\Ck}}_2
        (\seqname{A}^1, E^1, C^1)
    \biggr ]
  \subseteq
      \Functor^{\minimal{\Ck}}_2
        (\seqname{A}^2, E^2, C^2)
$.

\begin{proof}
$\funcseqname{f}$ is a morphism either by \pagecref{cor:Ck-ATLmorph} or
by hypothesis. By \pagecref{def:Ck-ATLmorph} there exists a subchart
$
  (U'^1, V'^1, \phi^1 \maps U'^1 \toiso V'^1)
  \in \seqname{A}^1
$
at $u^1$, an open set $U'^2 \subseteq U^2$ and a chart
$
  (U'^2 ,\hat{V}'^2, \phi'^2 \maps U'^2 \toiso \hat{V}'^2)
  \in \seqname{A}^2
$
at $u^2$ such that
$\funcname{f}_0[U'^1] \subseteq U'^2$ and diagram \eqref{eq:AtlM1} in
\pagecref{def:M-ATLmorph} is commutative, as shown in
\pagecref{fig:AtlM1}. Then $\funcname{f}_1 (v^1) \in \hat{V}'^2$.
\end{proof}
\end{lemma}

\begin{definition}[Model spaces associated with \ensuremath{\Ck}-atlases]
\label{def:Ck-modE}
Let $\seqname{A}^i$, $i=1,2$, be a $\Ck$-atlas of $E^i$ in the coordinate
space $C^i$ and
$\funcseqname{f} \defeq (\funcname{f}_0, \funcname{f}_1)$ an
$\seqname{E}^1$-$\seqname{E}^2$ $\Ck$-atlas morphism of $\seqname{A}^1$
to $\seqname{A}^2$ in the coordinate spaces $\seqname{C}^1$,
$\seqname{C}^2$. Then:

The minimal $\Ck$ model space with neighborhoods
in the $\Ck$-atlas $\seqname{A}^i$ of $E^i$ in the coordinate space
$C^i$ is

\begin{multline}
\Functor^{\minimal{\Ck}}_1 (\seqname{A}^i, E^i, C^i)
\defeq \\
\minimal{\Mod}
  \left (
    E^i,
    \pi_1[\seqname{A}^i],
    \set
      {\phi'^{-1} \compose \phi}%
      [
        \equant
          {
            {(U,V,\phi) \in \seqname{A}^i},
            {(U',V',\phi') \in \seqname{A}^i}
          }
          {
            U \cap U' \ne \emptyset
          }
      ]
  \right )
\end{multline}

The mapping associated with the $\seqname{E}^1$-$\seqname{E}^2$
$\Ck$-atlas morphism $(\funcname{f}_0,\funcname{f}_1)$ of
$\seqname{A}^1$ to $\seqname{A}^2$ in the coordinate spaces
$\seqname{C}^1$, $\seqname{C}^2$ is

\begin{multline}
\Functor^{\minimal{\Ck}}_1
  \bigl (
    (
      \funcname{f}_0,
      \funcname{f}_1
    ),
    (\seqname{A}^1, E^1, C^1),
    (\seqname{A}^2, E^2, C^2)
  \bigr )
\defeq \\
  \funcname{f}_0
  \maps
  \Functor^{\minimal{\Ck}}_1 (\seqname{A}^1, E^1, C^1)
  \to
  \Functor^{\minimal{\Ck}}_1 (\seqname{A}^2, E^2, C^2)
\end{multline}

If it is a model function then the it is also the m-atlas morphism
associated with the $\seqname{E}^1$-$\seqname{E}^2$ $\Ck$-atlas morphism
$(\funcname{f}_0,\funcname{f}_1)$ of $\seqname{A}^1$ to $\seqname{A}^2$
in the coordinate spaces $\seqname{C}^1$, $\seqname{C}^2$.
\end{definition}

\begin{lemma}[Model spaces associated with \ensuremath{\Ck}-atlases]
\label{lem:Ck-ismodE}
Let $\seqname{A}$ be a $\Ck$-atlas of $E$ in the coordinate
space $C$. Then $\Functor^{\minimal{\Ck}}_1 (\seqname{A}, E, C)$ is a
model space.

\begin{proof}
The result follows from \Pagecref{lem:minmod}.
\end{proof}
\end{lemma}

\begin{theorem}[Functors from \ensuremath{\Ck} atlases to model spaces]
Let $\seqname{E}$ be a set of topological spaces and $\seqname{C}$ a set
of linear spaces. Then $\Functor^{\minimal{\Ck}}_1$ is a functor from
$\Atl[\Ck](\seqname{E}, \seqname{C})$ to $\Trivcat{\seqname{E}}$,
$\Functor^{\minimal{\Ck}}_1$ is a functor from
$\full{\Atl[\Ck]}(\seqname{E}, \seqname{C})$ to $\Trivcat{\seqname{E}}$,
$\Functor^{\minimal{\Ck}}_2$ is a functor from
$\maximal[S-]{\Atl[\Ck]}(\seqname{E}, \seqname{C})$ to
$\optriv[\Ck-]{\seqname{C}}$ and $\Functor^{\minimal{\Ck}}_2$ is a
functor from $\full{\Atl[\Ck]}(\seqname{E}, \seqname{C})$ to
$\Triv[\Ck-]{\seqname{C}}$.

\begin{proof}
Let $o^i \defeq (\seqname{A}^i, E^i, C^i)$, $i \in [1,3]$, be
objects in $\Atl[\Ck](\seqname{E}, \seqname{C})$
and let \\
$
m^i \defeq
  \bigl (
    (\funcname{f}^i_0, \funcname{f}^i_1),
    o^i,
    o^{i+1}
  \bigr )
$
be a morphism in $\Atl[\Ck](\seqname{E}, \seqname{C})$.

$
  \Functor^{\minimal{\Ck}}_1 \maps
  \Atl[\Ck](\seqname{E}, \seqname{C}) \to
  \Triv{\seqname{E}}
$:
\begin{description}
\item[Preservation of endpoints:]
$
\Functor^{\minimal{\Ck}}_1(\funcname{m}^i) =
\funcname{f}^i_0 \maps \Functor^{\minimal{\Ck}}_1 o^i \to
\Functor^{\minimal{\Ck}}_1 o^{i+1}
$
\item[Composition:]
\begin{align*}
  \Functor^{\minimal{\Ck}}_1(\funcname{m}^2 \compose[A] \funcname{m}^1)
  & =
  \Functor^{\minimal{\Ck}}_1
    \bigl (
      (
        \funcname{f}^2_0 \compose \funcname{f}^1_0,
        \funcname{f}^2_1 \compose \funcname{f}^1_1
      )
      (\seqname{A}^1, E^1, C^1),
      (\seqname{A}^3, E^3, C^3)
    \bigr )
\\*
  & =
  \funcname{f}^2_0 \compose \funcname{f}^1_0 \maps
   \Functor^{\minimal{\Ck}}_1 o^1 \to
    \Functor^{\minimal{\Ck}}_1 o^3
\\*
  & =
  \bigl (
    \funcname{f}^2_0 \maps
    \Functor^{\minimal{\Ck}}_1 o^2 \to
    \Functor^{\minimal{\Ck}}_1 o^3
   \bigr )
  \compose
  \bigl (
    \funcname{f}^1_0 \maps
    \Functor^{\minimal{\Ck}}_1 o^1 \to
  \Functor^{\minimal{\Ck}}_1 o^2
 \bigr )
\\*
  & =
  \Functor^{\minimal{\Ck}}_1
    \bigl (
      (
        \funcname{f}^2_0,
        \funcname{f}^2_1
      )
      (\seqname{A}^2, E^2, C^2),
      (\seqname{A}^3, E^3, C^3)
    \bigr )
  \compose
\\*
  & \quad
  \Functor^{\minimal{\Ck}}_1
    \bigl (
      (
        \funcname{f}^1_0,
       \funcname{f}^1_1
      )
     (\seqname{A}^1, E^1, C^1),
      (\seqname{A}^2, E^2, C^2)
    \bigr )
\\*
  & =
  \Functor^{\minimal{\Ck}}_1(\funcname{m}^2)
  \compose
  \Functor^{\minimal{\Ck}}_1(\funcname{m}^1)
\end{align*}

\item[Identity:]
\leavevmode
\begin{enumerate}
\item
$
\Functor^{\minimal{\Ck}}_1(\Id[{o^i}]) =
\Functor^{\minimal{\Ck}}_1
  \bigl (
    (\Id[{E^i}], \Id[{C^i}]),
    (\seqname{A}^i, E^i, C^i),
    (\seqname{A}^i, E^i, C^i)
  \bigr )
=
\\
\Id[{E^i}] \maps \Functor^{\minimal{\Ck}}_1 o^i \to
\Functor^{\minimal{\Ck}}_1 o^i
$
\item
$
\Id_{\Functor^{\minimal{\Ck}}_1} o^i =
\Id[{E^i}] \maps \Functor^{\minimal{\Ck}}_1 o^i \to
\Functor^{\minimal{\Ck}}_1 o^i
$
\end{enumerate}
\end{description}

The proof for
$
\Functor^{\minimal{\Ck}}_1 \maps
\full{\Atl[\Ck]}(\seqname{E}, \seqname{C}) \to \Triv{\seqname{E}}
$
is identical.

$
  \Functor^{\minimal{\Ck}}_2 \maps
  \Atl[\Ck](\seqname{E}, \seqname{C}) \to
  \optriv[\Ck-]{\seqname{C}}
$:

\begin{description}
\item[Preservation of endpoints:]\quad\\*
$
\Functor^{\minimal{\Ck}}_2(\funcname{m}^i) =
\funcname{f}^i_1 \maps \Functor^{\minimal{\Ck}}_2 o^i \to
\Functor^{\minimal{\Ck}}_2 o^{i+1}
$
\item $\Functor (g \compose f) = \Functor (g) \compose \Functor (f)$:
$
\Functor^{\minimal{\Ck}}_2(\funcname{m}^2 \compose[A] \funcname{m}^1)
=
\\
\Functor^{\minimal{\Ck}}_2
  \bigl (
    (
      \funcname{f}^2_0 \compose \funcname{f}^1_0,
      \funcname{f}^2_1 \compose \funcname{f}^1_1
    )
    (\seqname{A}^1, E^1, C^1),
    (\seqname{A}^3, E^3, C^3)
  \bigr )
=
\\
\funcname{f}^2_0 \compose \funcname{f}^1_0 \maps
  \Functor^{\minimal{\Ck}}_2 o^1 \to
  \Functor^{\minimal{\Ck}}_2 o^3
=
\\
\bigl (
  \funcname{f}^2_1 \maps
  \Functor^{\minimal{\Ck}}_2 o^2 \to
  \Functor^{\minimal{\Ck}}_2 o^3
  \bigr )
\compose
\bigl (
  \funcname{f}^1_1 \maps
  \Functor^{\minimal{\Ck}}_2 o^1 \to
  \Functor^{\minimal{\Ck}}_2 o^2
\bigr )
=
\\
\Functor^{\minimal{\Ck}}_2
  \bigl (
    (
      \funcname{f}^2_0,
      \funcname{f}^2_1
    )
    (\seqname{A}^2, E^2, C^2),
    (\seqname{A}^3, E^3, C^3)
  \bigr )
\compose \\
\Functor^{\minimal{\Ck}}_2
  \bigl (
    (
      \funcname{f}^1_0,
      \funcname{f}^1_1
    )
    (\seqname{A}^1, E^1, C^1),
    (\seqname{A}^2, E^2, C^2)
  \bigr )
=
\\
\Functor^{\minimal{\Ck}}_2(\funcname{m}^2)
\compose
\Functor^{\minimal{\Ck}}_2(\funcname{m}^1)
$
\item $\Functor (\Id[{A}]) = \Id_{\Functor (A)}$:
\begin{enumerate}
\item
$
\Functor^{\minimal{\Ck}}_2(\Id[{o^i}]) =
\Functor^{\minimal{\Ck}}_2
  \bigl (
    (\Id[{E^i}], \Id[{C^i}]),
    (\seqname{A}^i, E^i, C^i),
    (\seqname{A}^i, E^i, C^i)
  \bigr )
=
\\
\Id[{C^i}] \maps \Functor^{\minimal{\Ck}}_2 o^i \to
\Functor^{\minimal{\Ck}}_2 o^i
$
\item
$
\Id_{\Functor^{\minimal{\Ck}}_2} o^i =
\Id[{C^i}] \maps \Functor^{\minimal{\Ck}}_2 o^i \to
\Functor^{\minimal{\Ck}}_2 o^i
$
\end{enumerate}
\end{description}

The proof for
$
  \Functor^{\minimal{\Ck}}_2 \maps
  \full{\Atl[\Ck]}(\seqname{E}, \seqname{C}) \to
  \Triv[\Ck-]{\seqname{C}}
$
is identical.
\end{proof}
\end{theorem}

\section{Classic \ensuremath{\Ck}-atlas morphisms and functors}
\label{sec:Ck-ATLcmorph}
This section introduces an alternate definition of and taxonomy for
morphisms between $\Ck$-atlases, defines categories of $\Ck$-atlases,
defines functors amomg them and proves some basic reasults.  It
introduces the notion of classic $\Ck$-atlas morphisms, which is
equivalent to the conventional definitions for a manifold.

\subsection{Classic \ensuremath{\Ck}-atlas  morphisms}
\label{sub:Ck-ATLcmorph}
This subsection introduces the notion of classic $\Ck$-atlas morphisms,
and proves some basic results.

\subsubsection{Definitions of classic \ensuremath{\Ck}-atlas morphisms}

\begin{definition}[Classic \ensuremath{\Ck}-atlas morphisms]
\label{def:Ck-ATLcmorph}
Let $E^i$, $i=1,2$, be a topological space, $C^i$ be a linear space,
$\seqname{A}^i$ be a $\Ck$-atlases of $E^i$ in the coordinate space
$C^i$ and $\funcname{f}$ be a continuous function.

$\funcname{f}$ is a $E^1$-$E^2$ classic $\Ck$ morphism of
$\seqname{A}^1$ to $\seqname{A}^2$ in the coordinate spaces $C^1$,
$C^2$, abbreviated as
$
\isAtl[classic,\Ck]_\Ar
  (
    \seqname{A}^1,
    E ^1,
    C^1,
    \seqname{A}^2,
    E^2,
    C^2,
    \funcname{f}
  )
$,
iff for any
$(U^i, V^i, \phi^i \maps U^i \toiso V^i) \in \seqname{A}^i$,
$i=1,2$, with $I \defeq U^1 \cap \funcname{f}^{-1}[U^2] \neq \emptyset$,
$
  \phi^2 \compose \funcname{f} \compose \phi^{1-1} \maps
    \phi^1[I] \to V^2
$
is a $\Ck$ function.
\end{definition}

\begin{remark}
Definitions of constrained, semistrict and strict classic $\Ck$-atlas
morphisms would be pointless, as any classic $\Ck$-atlas
morphism would be constrained and strict.
\end{remark}

\subsubsection{Proclamations on classic \ensuremath{\Ck}-atlas morphisms}
\begin{lemma}[Classic \ensuremath{\Ck}-atlas morphisms]
\label{lem:Ck-ATLcmorph}
Let $\seqname{E}$ be a set of topological spaces, $\seqname{C}$ a set of
$\Ck$ linear spaces, $E^i \in \seqname{E}$, $i=1,2$,
$C^i \in \seqname{C}$, $\seqname{A}^i$ a $\Ck$-atlas of $E^i$ in the
coordinate space $C^i$ and $\funcname{f} \maps E^1 \to E^2$ a function.

$\funcname{f}$ is an
$E^1$-$E^2$ classic $\Ck$-atlas near morphism of $\seqname{A}^1$ to
$\seqname{A}^2$ in the coordinate spaces $C^1$, $C^2$ iff $\funcname{f}$
is an
$E^1$-$E^2$-%
$\optriv[\Ck-]{\seqname{C}}$-%
$\optriv[\Ck-]{\seqname{C}}$
classic m-atlas near morphism of $\seqname{A}^1$ to $\seqname{A}^2$ in
the coordinate spaces $\Triv[\Ck-]{\seqname{C}^1}$,
$\Triv[\Ck-]{\seqname{C}^2}$.

Let $E^1$ be an open subspace of $E^2$. Then $\funcname{f}$ is an
$E^1$-$E^2$ classic $\Ck$-atlas morphism of $\seqname{A}^1$ to
$\seqname{A}^2$ in the coordinate spaces $C^1$, $C^2$ iff $\funcname{f}$
is a strict $E^1$-$E^2$ classic m-atlas morphism of $\seqname{A}^1$ to
$\seqname{A}^2$ in the coordinate spaces $\Triv[\Ck-]{\seqname{C}^1}$,
$\Triv[\Ck-]{\seqname{C}^2}$.

$\funcname{f}$ is an $E^1$-$E^2$ classic $\Ck$-atlas morphism of
$\seqname{A}^1$ to $\seqname{A}^2$ in the coordinate spaces $C^1$, $C^2$
iff $\funcname{f}$ is a strict
$E^1$-$E^2$-%
$\optriv[\Ck-]{\seqname{C}}$-%
$\optriv[\Ck-]{\seqname{C}}$
classic m-atlas morphism of $\seqname{A}^1$ to $\seqname{A}^2$ in the
coordinate spaces $\Triv[\Ck-]{\seqname{C}^1}$,
$\Triv[\Ck-]{\seqname{C}^2}$.


\end{lemma}

\section {$\Ck$ manifolds}
\label{sec:ckman}
Conventionally a manifold is different from its atlases, but
$\maximal{\Atl[\Ck]}(\seqname{E}, \seqname{C})$ in
\pagecref{def:Ck-ATLcat} encourages treating them on an equal footing.
All of the results for maximal atlases carry directly over to results
for manifolds.

\begin{definition}[\ensuremath{\Ck} manifolds]
\label{def:man}
Let $E$ be a topological space, $C$ a linear space and $\seqname{A}$ a
maximal\footnote{
Requiring that the atlas be full would eliminate some pathologies.}
$\Ck$-atlas of $E$ in the coordinate space $C$. Then $(E, C,
\seqname{A})$ is a $\Ck$ manifold.

Let $\seqname{E}$ be a set of topological spaces and $\seqname{C}$ be a
set of linear spaces. Then
\begin{equation}
\Man^\Ck_\Ob(\seqname{E}, \seqname{C}) \defeq
\maximal{\Atl[\Ck]}(\seqname{E}, \seqname{C})
\end{equation}
\begin{remark}
The manifold $(E, C, \seqname{A})$ corresponds to the object
$(\seqname{A}, E, C)$.
\end{remark}
\end{definition}

\begin{definition}[\ensuremath{\Ck} manifold morphisms]
\label{def:manmorph}
Let $S^i \defeq (E^i,C^i,\seqname{A}^i)$, $i=1,2$, be $\Ck$ manifolds
and $(\funcname{f}_0 \maps E^1 \to E^2, \funcname{f}_1 \maps C^1 \to C^2)$
be an $E^1$-$E^2$ $\Ck$-morphism of $\seqname{A}^1$ to $\seqname{A}^2$
in the coordinate spaces $C^1$, $C^2$. Then
$(\funcname{f}_0, \funcname{f}_1)$ is a $\Ck$ morphism of $S^1$ to
$S^2$.

Let $E^i$, $i=1,2$, be a topological space and $C^i$ be a linear space.
Then
\begin{multline}
\Man^\Ck_\Ar(E^1, C^1, E^2, C^2) \defeq
\Atl[\Ck]_\Ar(E^1, C^1, E^2, C^2)
\end{multline}

Let $\seqname{E}$ be a set of topological spaces and $\seqname{C}$ a
set of linear spaces. Then
\begin{equation}
\Man^\Ck_\Ar(\seqname{E}, \seqname{C}) \defeq
\maximal{\Atl[\Ck]_\Ar}(\seqname{E}, \seqname{C})
\end{equation}
\begin{equation}
\Man^\Ck(\seqname{E}, \seqname{C}) \defeq
\maximal{\Atl[\Ck]}(\seqname{E}, \seqname{C})
\end{equation}
\end{definition}

\begin{theorem}[Categories of \ensuremath{\Ck} manifolds]
Let $\seqname{E}$ be a set of topological spaces and $\seqname{C}$ a set
of linear spaces. Then $\Man^\Ck(\seqname{E}, \seqname{C})$ is a
category and the identity morphism of $(\seqname{A}, E, C)$ is an
identity morphism.

\begin{proof}
{
  \showlabelsinline
  The result follows directly from \cref{def:man,def:manmorph} above and
}
\pagecref{lem:Ck-ATLiscat}.
\end{proof}
\end{theorem}

\part{Equivalence of fiber bundles}
\label{part:bun}
For fiber bundles\footnote{
  The literature has several definitions of fiber bundle. This paper
  uses one chosen for clarity of exposition. It differs from
  \cite[p.~8]{TopFib} in that, e.g., it uses the machinery of maximal
  atlases rather than equivalence classes of coordinate bundles, the
  nomenclature differs in several minor regards.
},
the adjunct spaces are the base space $\seqname{X} = (X,\catname{X})$,
the fiber $\seqname{Y} = (Y, \catname{Y})$
and the group $G$;
the category of the coordinate space is the category of Cartesian
products $\{U \times Y \mid U \objin \catname{X}\}$ of model
neighborhoods in the base space with the entire fiber, with morphisms
$\funcname{t} \maps U \times Y \to U \times Y$ that preserve the
fibers, i.e., $\pi_1 \compose \funcname{t} = \pi_1$, and are
generated by the group action on the fiber
(Equation~\eqref{eq:transitiongroup}).

The sole adjunct functions are the projection $\pi \maps E \onto X$, the
group operation and the group action on the fiber.

This part of the paper defines bundle atlases, fiber bundles, local
coordinate spaces equivalent to fiber bundles, categories of them and
functors, and gives basic results.

\begin{definition}[Trivial group category of groups]
Let $\seqname{G}$ be a set of topological groups. The trivial group
category of $\seqname{G}$, abbreviated
$\Trivcat[\mathrm{group}-]{\seqname{G}}$, is the category of all
continuous homomorphisms between groups of $\seqname{G}$. By abuse of
language it will be shortened to $\Trivcat{\seqname{G}}$ when the
meaning is clear from context.
\end{definition}

\begin{definition}[Group actions]
Let $Y$ be a topological space, $G$ a topological group, $\rho$ an
effective group action of $G$ on $Y$, $y \in Y$ and $g \in G$. Then
$y \star g \defeq \rho(y,g)$.

Let $X$ be a topological space and $x \in X$. Then
$(x, y) \star g \defeq (x, y \star g)$.

A $\star$ with a subscript, superset, underset or overset refers to the
group action $\rho$ with the same subscript, superset, underset or
overset.

\begin{remark}
This notation is only used when it is clear from context what the group
action is.
\end{remark}
\end{definition}

\begin{definition}[Protobundles]
\label{def:ProtoBun}
Let $\seqname{B} \defeq (E, X, Y, \pi, G, \rho)$, where $E$, $X$ and $Y$
are topological spaces, $G$ a topological group, $\pi \maps E \onto X$ a
continuous surjection and $\rho$ an efective group action of $G$ on $Y$.
Then $\seqname{B}$ is a protobundle.
\begin{remark}
While this definition does not itself require $E$ to have a local product
structure nor require $\pi$
to have the Covering Homotopy Property, only those protobundles having
an atlas are of interest, and for them \pagecref{def:BunAtl} imposes
additional constraints.
\end{remark}
\end{definition}

\section{$G$-$rho$ model spaces}
\label{sec:grhomod}
\begin{definition}[$G$-$\rho$-model spaces]
{
  \showlabelsinline
  \label{def:Grhomod}
}
Let $\seqname{XY} \defeq (X \times Y, \catname{XY})$ be a model space,
$G$ a topological group and $\rho$ an effective group action of $G$ on
$Y$ such that the objects of $\catname{XY}$ are products of open sets
with $Y$ and the morphisms are fiber-preserving automorphisms generated
by the group action, i.e.,
\begin{equation}
\uquant
  {
    U \objin \catname{XY}
  }
  {
    \equant
      {
        V \in \op{X}
      }
      {
        U = V \times Y
      }
  }
\end{equation}
\begin{equation}
\label{eq:fGrho}
\uquant
  {
   \funcname{f} \arin \catname{XY} \maps V \times Y \toiso V \times Y
  }
  {
    \equant
    {
      \funcname{g} \in G^V
    }
    {
    \uquant
      {
        {(x,y) \in V \times Y}
      }
      {
        \funcname{f}(x,y) =
          \bigl ( x, y \star \funcname{g}(x) \bigr )
      }
    }
  }
\end{equation}

Then $\seqname{XY}$ is a $G$-$\rho$ model space of $X \times Y$,
abbreviated
$
\mathrm{isG{\rho}}
(
  \seqname{XY},
  Y,
  G,
  \rho
)
$,
$G_{\seqname{XY},\funcname{f}} \defeq \range(\funcname{g})$
is the set of group elements associated with $\funcname{f}$
and
$
  G_\seqname{XY} \defeq
  \union%
    [{\funcname{f} \arin \catname{XY}}]
    {G_{\seqname{XY},\funcname{f}}}
$
is the set of group elements associated with $\seqname{XY}$ morphisms.
\end{definition}

\begin{lemma}[$G$-$\rho$-model spaces]
\label{lem:Grhomod}
Let $\seqname{XY}$ be a $G$-$\rho$ model space of $X \times Y$ and
\nobreaks
{{
  $\funcname{f} \arin \catname{XY} \maps V \times Y \toiso V \times Y$.
}}

The function $\funcname{g}$ in \cref{eq:fGrho} is unique.

\begin{proof}
The group action is effective.
\end{proof}

$G_\seqname{XY}$ is unique.

\begin{proof}
The function $\funcname{g}$ is unique.
\end{proof}
\end{lemma}

\begin{definition}[Morphisms of $G$-$\rho$-model spaces]
\label{def:Grhomorph}
Let
$X^i.Y^i$, $i=1,2$, be topological spaces,
$G^i$ a topological group,
$\rho^i$ an effective group action on $Y^i$ and
\nobreaks
{{
  $\seqname{XY}^i \defeq (X^i \times Y^i, \catname{XY}^i)$
}}
a $G^i$-$\rho^i$ model space of $X^i \times Y^i$.
Then a model function
$\funcname{f}_C \maps \seqname{XY}^1 \to \seqname{XY}^2$ is a
$G^1$-$G^2$-$\rho^1$-$\rho^2$ morphism of $X^1 \times Y^1$ to $X^2 \times Y^2$,
abbreviated
$
\mathrm{isG{\rho}morph}
(
  \seqname{XY}^1,
  G^1,
  \rho^1,
  \seqname{XY}^2,
  G^2,
  \rho^2,
  \funcname{f}_C
)
$
iff it preserves the group action, i.e., there is a continuous
homomorphism $\funcname{f}_G \maps G^1 \to G^2$ such that
\fullcref{fig:Grho} is commutative, i.e.,
{
  \showlabelsinline
  \cref{eq:Grho}
}
below holds.

\begin{figure}
\[ \bfig
\node xyg1(0,1000)[\bigl ( (x^1,y^1), g^1 \bigr )]
\node xys1(0,0)[(x^1, y^1 \star g^1)]
\node xyg2(1000,1000)[\bigl ( (x^2,y^2), g^2 \bigr )]
\node xys2(1000,0)[(x^2, y^2 \star g^2)]
\arrow |r|[xyg1`xyg2;\funcname{f}_C \times \funcname{f}_G]
\arrow |a|[xyg1`xys1;\star^1]
\arrow |r|[xys1`xys2;\funcname{f}_C]
\arrow |a|[xyg2`xys2;\star^2]
\efig \]
\caption{Preserving group action}
\label{fig:Grho}
\end{figure}
\begin{equation}
\showlabelsinline
\label{eq:Grho}
\equant
  {
    \funcname{f}_G \maps G_1 \to G_2
  }
  {
    \uquant
    {
      {
        (x,y) \in X^1 \times Y^1
      },
      {
        g \in G^1
      }
    }
    {
      \funcname{f}_C \bigl ( (x,y) \star^1 g \bigr )
      =
      \funcname{f}_C \bigl ( (x,y) \bigr ) \star^2 \funcname{f}_G (g)
    }
  }
\end{equation}
\end{definition}

\begin{lemma}[Morphisms of $G$-$\rho$-model spaces]
\label{lem:Grhomorph}
Let $X^i.Y^i$, $i=1,2,3$, be topological spaces, $G^i$ a topological
group, $\rho^i$ an effective group action on $Y^i$,
\nobreaks
{{
  $\seqname{XY}^i \defeq (X^i \times Y^i, \catname{XY}^i)$
}}
a $G^i$-$\rho^i$ model space of $X^i \times Y^i$ and
$\funcname{f}^i_C \maps X^i \times Y^i \to X^{i+1} \times Y^{i+1}$ a
$G^i$-$G^{i+1}$-$\rho^i$-$\rho^{i+1}$ morphism of $X^i \times Y^i$ to
$X^{i+1} \times Y^{i+1}$.

\begin{enumerate}
\item
\label{grhomorph:fgunique}
The function $\funcname{f}_G$ in \cref{eq:Grho} is unique.

\begin{proof}
The group action is effective.
\end{proof}

\item
\label{grhomorph:fiber}
$\funcname{f}^i_C$ preserves fibers, i.e.,
$
  \pi_1 \bigl ( \funcname{f}^i_C(x,y) \bigr ) =
  \pi_1 \bigl ( \funcname{f}^i_C(x,y') \bigr )
$
for $x$ in $X^i$ and $y,y'$ in $ Y^i$,
\begin{proof}
Since $\rho^i$ is effective, there is a $g \in G^i$ such that
$y' = y \star^i g$ and thus by \cref{def:Grhomorph}
\begin{equation}
\begin{split}
  \pi_1 \bigl ( \funcname{f}^i_C(x,y') \bigr )
  & =
\\
  \pi_1 \bigl ( \funcname{f}^i_C(x,y \star^i g) \bigr )
  & =
\\
  \pi_1 \bigl ( \funcname{f}^i_C(x,y) \star^{i+1} \funcname{f}^i_G(g) \bigr )
  & =
\\
  \pi_1 \bigl ( \funcname{f}^i_C(x,y) \bigr )
\end{split}
\end{equation}
\end{proof}

\item
\label{grhomorph:fxunique}
There exists a unique function $\funcname{f}^i_X \maps X^i \to X^{i+1}$
such that
$\funcname{f}^i_X \compose \pi_1 = \pi_1 \compose \funcname{f}^i_C$.

\begin{proof}
Define
$\funcname{f}^i_X(x) = \pi_1 \bigl ( \funcname{f}^i_C(x,y) \bigr )$,
where $y$ is an arbitrary point of $Y^i$. It does not depend on the
choice of $y$ because $\funcname{f}^i_C$ preserves fibers.
\end{proof}

\begin{remark}
$\funcname{f}^i_C$ may be a twisted product:  there need not exist
$\funcname{f}^i_Y \maps Y^i \to Y^{i+1}$ such that
$\funcname{f}^i_C = \funcname{f}^i_X \times \funcname{f}^i_Y$.
\end{remark}

\item
\label{grhomorph:comp}
$\funcname{f}^2_C \compose \funcname{f}^1_C$ is a
$G^1$-$G^3$-$\rho^1$-$\rho^3$ morphism of $X^1 \times Y^1$ to $X^3 \times Y^3$.

\begin{proof}

\begin{figure}
\[ \bfig
\node xyg1(0,1500)[\bigl ( (x^1,y^1), g^1 \bigr )]
\node xys1(0,0)[(x^1, y^1 \star g^1)]
\node xyg2(1000,1000)[\bigl ( (x^2,y^2), g^2 \bigr )]
\node xys2(1000,0)[(x^2, y^2 \star g^2)]
\node xyg3(2000,1500)[\bigl ( (x^3,y^3), g^3 \bigr )]
\node xys3(2000,0)[(x^3, y^3 \star g^3)]
\arrow |r|[xyg1`xyg3;\funcname{f}^2_C \compose \funcname{f}^1_C \times%
  \funcname{f}^2_G \compose \funcname{f}^1_G]
\arrow |r|[xyg1`xyg2;\funcname{f}^1_C \times \funcname{f}^1_G]
\arrow |a|[xyg1`xys1;\star^1]
\arrow |r|[xys1`xys2;\funcname{f}^1_C]
\arrow |a|[xyg2`xys2;\star^2]
\arrow |r|[xyg2`xyg3;\funcname{f}^2_C \times \funcname{f}^2_G]
\arrow |r|[xys2`xys3;\funcname{f}^1_C]
\arrow |a|[xyg3`xys3;\star^3]
\efig \]
\caption{Preserving group actions}
\label{fig:Grho2}
\end{figure}

Let $\funcname{f}^i_G \maps G^i \to G^{i+1}$ be a continuous
homomorphism such that
\begin{equation}
\label{eq:Grho2}
    \uquant
    {
      {
        (x^i,y^i) \in X^i \times Y^i
      },
      {
        g^i \in G^i
      }
    }
    {
      \funcname{f}^i_C \bigl ( (x^i,y^i) \star g^i \bigr )
      =
      \funcname{f}^i_C \bigl ( (x^i,y^i) \bigr ) \star \funcname{f}^i_G (g^i)
    }
\end{equation}
Let $(x^1,y^1) \in X^1 \times Y^1$ and $g^1 \in G^1$. Then \fullcref{fig:Grho2} is
commutative and
\begin{align}
  \funcname{f}^2_C \compose \funcname{f}^1_C
  \bigl ( (x^1,y^1) \star g^1 \bigr ) = \\*
  & =
  \funcname{f}^2_C
  \Bigl (
    \funcname{f}^1_C \bigl ( (x^1,y^1) \bigr ) \star
    \funcname{f}^1_G (g^1)
  \Bigr )
\\*
  & =
  \funcname{f}^2_C
  \Bigl (
    \funcname{f}^1_C \bigl ( (x^1,y^1) \bigr )
  \Bigr )
  \star\funcname{f}^2_G \compose \funcname{f}^1_G (g^1) = \\*
\\*
  & =
  \funcname{f}^2_C \compose \funcname{f}^1_C \bigl ( (x^1,y^1) \bigr )
  \star\funcname{f}^2_G \compose \funcname{f}^1_G (g^1)
\end{align}
\end{proof}
\end{enumerate}
\end{lemma}

\begin{definition}[Categories of $G$-$\rho$-model spaces]
\label{def:Grhomodcat}
Let $X^\alpha.Y^\alpha$, $\alpha \prec \Alpha$, be topological spaces,
$G^\alpha$ a topological group, $\rho^\alpha$ an effective group action
on $Y^\alpha$, \\
$
  \seqname{XY}^\alpha \defeq
  (X^\alpha \times Y^\alpha, \catname{XY}^\alpha)
$
a $G^\alpha$-$\rho^\alpha$ model space of $X^\alpha \times Y^\alpha$,
$
  \catseqname{XYG\rho}_\Ob \defeq \\
  \set
    {\seqname{XY}^\alpha}%
    [\alpha \prec \Alpha]
$,
$
  \catseqname{XYG\rho}_\Ar \defeq
  \set
    {\funcname{f}_C \maps X^\alpha \times Y^\alpha \to X^\beta \times Y^\beta}%
    [
      {
        \mathrm{isG{\rho}morph}
        (
          \seqname{XY}^\alpha,
          G^\alpha,
          \rho^\alpha,
          \seqname{XY}^\beta,
          G^\beta,
          \rho^\beta,
          \funcname{f}_C
        )
      },
      \alpha \prec \Alpha,
      \beta \prec \Alpha
    ]*
$
and
$
  \catseqname{XYG\rho} \defeq \\
  (
    \catseqname{XYG\rho}_\Ob,
    \catseqname{XYG\rho}_\Ar
  )
$.
Then any subcategory of $\catseqname{XYG\rho}$ is a $G$-$\rho$ model category.
\end{definition}

\begin{definition}[Trivial $G$-$\rho$-model spaces]
\label{def:Grhomodtriv}
Let $X$ and $Y$ be topological spaces, $G$ a topological group, $\rho$
an effective group action of $G$ on $Y$ and $\catname{XY}$ the category
of all products of open subsets of $X$ with $Y$ and all homeomorphisms
induced by the group action, i.e,

\begin{equation}
\Ob(\catname{XY})
\defeq
\set
  {V \times Y}%
  [V \in \op{X}]
\end{equation}
\begin{multline}
\Ar(\catname{XY})
\defeq
\set
  {
   \funcname{f} \maps V \times Y \toiso V \times Y
  }%
  [
    {V \in \op{X}},
    {\pi_1 \compose \funcname{f} = \pi_1},
    {
      \equant
        {
         \funcname{g} \in G^V,
        }
        {
          \uquant
            {
              {x \in V},
              {y \in Y}
            }
            {
               f(x, y) = \bigl ( x, y \star \funcname{g}(x) \bigr )
            }
        }
    }
  ]*
\end{multline}
Then the trivial $G$-$\rho$ model space of $X,Y$, abbreviated
$\Triv[G-\rho-]{X,Y}$, is $(X \times Y, \catname{XY})$ and $\Triv[G-\rho-]{X,Y}$ is
a trivial $G$-$\rho$ model space of $X,Y$.

\begin{remark}
Let $G'$ be a topological group and $\rho'$ an effective group action of
$G'$ on $Y$ such that $\Triv[G-\rho-]{X,Y} = \Triv[G'-\rho'-]{X,Y}$.
Although $G'$ must be isomorphic to $G$. it need not have the same
topology.
\end{remark}

The identity morphism of $\Triv[G-\rho-]{X,Y}$ is
$\Id_{\Triv[G-\rho-]{X,Y}} \defeq \Id_{X \times Y}$.
\end{definition}

\begin{lemma}[The trivial $G$-$\rho$ model space of $X,Y$ is a $G$-$\rho$ model space of $X,Y$]
\label{lem:Grhomodtriv}
Let $X$ and $Y$ be topological spaces, $G$ a topological group and
$\rho$ an effective group action on $Y$. Then $\Triv[G-\rho-]{X,Y}$ is
a $G$-$\rho$ model space of $X \times Y$.

\begin{proof}
$\Triv[G-\rho-]{X,Y}$ satisfies the conditions of
\pagecref{def:Grhomod}:
\begin{enumerate}
\item Product with fiber:
\[
\uquant
  {
    U \objin \pi_2 \biggl ( \Triv[G-\rho-]{X,Y} \biggr )
  }
  {
    \equant
      {
        V \in \op{X}
      }
      {
        U = V \times Y
      }
  }
\]
by \cref{def:Grhomodtriv}.
\item Generated by $\rho$:
Let
$
  \funcname{f} \arin
  \pi_2 \biggl ( \Triv[G-\rho-]{X,Y} \biggr )
  \maps V \times Y \toiso V \times Y
$.
Then
\[
  \equant
    {
      \funcname{g} \in G^V
    }
    {
    \uquant
      {
        {(x,y) \in V \times Y}
      }
      {
        \funcname{f}(x,y) =
          \bigl ( x, y \star \funcname{g}(x) \bigr )
      }
    }
\]
by \cref{def:Grhomodtriv}.
\end{enumerate}
\end{proof}
\end{lemma}

\begin{lemma}[Morphisms of trivial $G$-$\rho$-model spaces]
\label{lem:Grhomodtrivmorph}
Let $X^i$ and $Y^i$, $i=1,2$, be topological spaces, $G^i$ a topological
group, $\rho^i$ an effective group action of $G^i$ on $Y^i$ and
$
  \funcseqname{f} \defeq f_X \times f_Y \maps
  \Triv[G^1-\rho^1-]{X^1,Y^1} \to \Triv[G^2-\rho^2-]{X^2,Y^2}
$
a pair of continuous functions. Then $\funcseqname{f}$ is a model
function from $\Triv[G^1-\rho^1-]{X^1,Y^1}$ to
$\Triv[G^2-\rho^3-]{X^2,Y^3}$.
\begin{proof}
\leavevmode
\begin{enumerate}
\item
$\funcseqname{f}$ is a Cartesian product of continuous functions, hence
continuous.

\item
Let $m^2 \defeq U^2 \times Y^2$ be a model neighborhood of
$\Triv[G^2-\rho^2-]{X^2,Y^2}$. Then
$U^1 \defeq \funcname{f}_X^{-1}(U^2)$ is open and
$\funcseqname{f}^{-1}(m^2) = U^1 \times Y^1$ is a model function.

\item
Let $m^1 \defeq U^1 \times Y^1$ be a model neighborhood of
$\Triv[G^1-\rho^1-]{X^1,Y^1}$. $X^2 \times Y^2$ is a model neighborhood
of $\Triv[G^2-\rho^2-]{X^2,Y^2}$ and
$\funcseqname{f}(m^1) \subseteq X^2 \times Y^2$.
\end{enumerate}
\end{proof}
\end{lemma}

\begin{definition}[Categories of trivial $G$-$\rho$-model spaces]
\label{def:Grhomodtrivcat}
Let
$
  \seqname{B}^\alpha \defeq
  (E^\alpha, X^\alpha, Y^\alpha, \pi^\alpha, G^\alpha, \rho^\alpha)
$,
$\alpha \prec \Alpha$, be a protobundle and
$\seqname{B} \defeq \set{{\seqname{B}^\alpha}}[\alpha \prec \Alpha]$
be a set of protobundles.

The trivial coordinate model category of $\seqname{B}$,
$\Trivcat[\Bun-]{\seqname{B}}$, is the category with objects all trivial
$G_\alpha$-$\rho_\alpha$ model spaces of $X_\alpha, Y_\alpha$,
$\alpha \prec \Alpha$ and morphisms all continuous functions compatible
with the group action:

\begin{equation}
\Trivcat[\Bun-]{\seqname{B}}_\Ob \defeq
\set
  {
    \Triv[G-\rho-]{X,Y}
  }%
  [
    {
      (
        E,
        X,
        Y,
        \pi,
        G,
        \rho
      )
      \in \seqname{B}
    }
  ]
\end{equation}
\begin{multline}
\Trivcat[\Bun-]{\seqname{B}}_\Ar \defeq
\set
  {
    \funcname{f}_C \maps
    X^1 \times Y^1 \to
    X^2 \times Y^2
  }%
  [
    {
      \equant
        {
          {
            (
              E^i,
              X^i,
              Y^i,
              \pi^i,
              G^i,
              \rho^i
            )
            \in \seqname{B}
          },
          \funcname{f}_G \maps G^1 \to G^2
        }
        {
          \mathrm{isG{\rho}morph}
          (
            \Triv[G-\rho-]{X^1,Y^1},
            G^1,
            \rho^1,
            \Triv[G-\rho-]{X^2,Y^2},
            G^2,
            \rho^2,
            \funcname{f}^C
          )
        }
    }
  ]*
\end{multline}
\begin{equation}
\Trivcat[\Bun-]{\seqname{B}} \defeq
\bigl (
  \Trivcat[\Bun-]{\seqname{B}}_\Ob,
  \Trivcat[\Bun-]{\seqname{B}}_\Ar
\bigr )
\end{equation}

\begin{remark}
The morphisms $\funcname{f}_C$ may be twisted products:  there need not
exist $\funcname{f}_Y \maps Y^1 \to Y^2$ such that
$\funcname{f}_C = \funcname{f}_X \times \funcname{f}_Y$.
\end{remark}

The trivial product coordinate model category of $\seqname{B}$,
$\Trivcat[\BunProd-]{\seqname{B}}$, is the category with objects all
trivial $G^\alpha$-$\rho^\alpha$ model spaces of $X^\alpha, Y^\alpha$,
$\alpha \prec \Alpha$ and morphisms all products of continuous functions
compatible with the group action:

\begin{equation}
\Trivcat[\BunProd-]{\seqname{B}}_\Ob \defeq \Trivcat[\Bun-]{\seqname{B}}_\Ob
\end{equation}
\begin{multline}
\Trivcat[\BunProd-]{\seqname{B}}_\Ar \defeq
\set
  {
    \funcname{f}_C \defeq  \funcname{f}_X \times \funcname{f}_Y \maps
    X^1 \times Y^1 \to
    X^2 \times Y^2
  }%
  [
    {
      \equant
        {
          {
            (
              E^i,
              X^i,
              Y^i,
              \pi^i,
              G^i,
              \rho^i
            )
            \in \seqname{B}
          },
          \funcname{f}_G \maps G^1 \to G^2
        }
        {
          \mathrm{isG{\rho}morph}
          (
            \Triv[G-\rho-]{X^1,Y^1},
            G^1,
            \rho^1,
            \Triv[G-\rho-]{X^2,Y^2},
            G^2,
            \rho^2,
            \funcname{f}_C
          )
        }
    }
  ]*
\end{multline}
\begin{equation}
\Trivcat[\BunProd-]{\seqname{B}} \defeq
\Bigl (
  \Trivcat[\BunProd-]{\seqname{B}}_\Ob,
  \Trivcat[\BunProd-]{\seqname{B}}_\Ar
\Bigr )
\end{equation}

Any subcategory of $\Trivcat[\Bun-]{\seqname{B}}$ is a trivial coordinate
model category and any subcategory of $\Trivcat[\BunProd-]{\seqname{B}}$ is
a trivial product coordinate model category.
\end{definition}

\begin{lemma}[The trivial coordinate model category of $\seqname{B}$ is a category]
\label{lem:Grhomodtrivcat}
Let \\
$
  \seqname{B}^\alpha \defeq
  (E^\alpha, X^\alpha, Y^\alpha, \pi^\alpha, G^\alpha, \rho^\alpha)
$,
$\alpha \prec \Alpha$, be a protobundle and
$\seqname{B} \defeq \set{{\seqname{B}^\alpha}}[\alpha \prec \Alpha]$
be a set of protobundles.

Then $\Trivcat[\Bun-]{\seqname{B}}$ is a category and the identity
morphism for object $\seqname{B}^\alpha$ is
$\Id_{X^\alpha \times Y^\alpha}$.

\begin{proof}
\leavevmode
\begin{description}
\item[Composition:]
The composition of $G$-$\rho$ morphisms is a $G$-$\rho$ morphism by
\pagecref{lem:Grhomorph}.
\item[Associativity:]
Morphisms are simply functions and composition of morphisms is simply
composition of functions.
\item[Unit:]
The identity morphisms are simply identity functions and composition of
morphisms is simply composition of functions.
\end{description}
\end{proof}

$\Trivcat[\BunProd-]{\seqname{B}}$ is a category and the identity
morphism for object $\seqname{B}^\alpha$ is
$\Id_{X^\alpha \times Y^\alpha}$.

\begin{proof}
\leavevmode
\begin{description}
\item[Composition:]
The composition of $G$-$\rho$ morphisms is a $G$-$\rho$ morphism by
\pagecref{lem:Grhomorph}. Let
$\Triv[G^i-\rho^i-]{X^i,Y^i} \objin\Trivcat[\BunProd-]{\seqname{B}},$
$i=1,2,3$,
$
  \funcname{f}^i_C =
  \funcname{f}^i_X \times \funcname{f}^i_Y \maps
  X^i \times Y^i \to X^{i+1} \times Y^{i+1}
  \arin  \Trivcat[\BunProd-]{\seqname{B}}
$.
Then
$
  \funcname{f}^{i+1}_C \compose \funcname{f}^i_C =
  (\funcname{f}^{i+1}_X \compose \funcname{f}^i_X)
  \times
  (\funcname{f}^{i+1}_Y \compose \funcname{f}^i_Y)
$.

\item[Associativity:]
Morphisms are simply functions and composition of morphisms is simply
composition of functions.

\item[Unit:]
The identity morphisms are simply identity functions and composition of
morphisms is simply composition of functions.
\end{description}
\end{proof}
\end{lemma}

\section{$G$-$rho$-nearly commutative diagrams}
\label{sec:grhoncd}
Let
\begin{enumerate}
\item
$X$ and $Y$ be topological spaces
\item
$G$ be a topological group,
\item
$\rho$ an effective group action on $Y$
\item
$\seqname{C}=(C,\catname{C}) \defeq \Triv[G-\rho-]{X,Y}$
\item
$D$ be, as in diagram \eqref{eq:NCD} of \pageref{part:ncd}, a tree with
two branches, whose nodes are topological spaces $U_i$ and $V^j$ and
whose links are continuous functions $\funcname{f}_i \maps U_i \to
U_{i+1}$ and $\funcname{f}'_j \maps U_j \to U_{j+1}$ between the spaces,
with $U_0 = V_0$, $U_m \subseteq C$ open and $V_n \subseteq C$ open, as
shown in \pagecref{fig:NCDb}.
\end{enumerate}

\begin{definition}[$G$-$\rho$-nearly commutative diagrams]
$D$ is nearly commutative in $X,Y,\pi,\rho$ iff $D$ is nearly
commutative in category $\catname{C}$.
\end{definition}

\begin{definition}[$G$-$\rho$-nearly commutative diagrams at a point]
\label{def:grhoNCD}
Let $\catname{C}$ and $D$ be as above and $x$ be an element of the
initial node. $D$ is nearly commutative in $X,Y,\pi,\rho$ at $x$ iff
$D$ is nearly commutative in $\catname{C}$ at $x$.
\end{definition}

\begin{definition}[$G$-$\rho$-locally nearly commutative diagrams]
Let $\catname{C}$ and $D$ be as above. $D$ is locally nearly commutative
in $X,Y,\pi,\rho$ iff $D$ is locally nearly commutative in
$\catname{C}$.
\end{definition}

\section{Bundle charts}
{
  \showlabelsinline
  \label{sec:buncharts}
}
\begin{definition}[$Y$-$\pi$-bundle charts]
\label{def:Ypichart}
Let $E,X,Y$ be topological spaces and \linebreak
\nobreaks
{{
  $\pi \maps E \to X$.
}}
A $Y$-$\pi$-bundle chart $C \defeq (U, V \times Y, \phi)$ of $E$ in the
coordinate space $X \times Y$ consists of

\begin{figure}[ht]
\[ \bfig
\node u(0,1000)[U]
\node vy(0,0)[V \times Y]
\node x(1000,0)[X]
\arrow |l|[u`vy;\phi]
\arrow |r|[u`vy;\iso]
\arrow |a|[u`x;\pi]
\arrow |r|[vy`x;\pi_1]
\efig \]
\caption{$\phi$ preserves fibers}
\label{fig:Ypichart}
\end{figure}

\begin{enumerate}
\item a nonvoid open set $U \subseteq E$, known as a coordinate patch
\item An open set $V \times Y \subseteq X \times Y$
\item a homeomorphism $\phi \maps U \toiso V \times Y$, known as a
coordinate function, that preserves fibers. i.e.,
$\pi_1 \compose \phi = \pi$.
\end{enumerate}

$C$ covers $u$ iff $u \in U$.
\end{definition}

\begin{lemma}[$Y$-$\pi$-bundle charts]
\label{lem:Ypichart}
Let $E,X,Y$ be topological spaces,
$G$ a topological group, $\rho$ an effective group action of $G$ on
$Y$
and $\seqname{XY} \defeq (X \times Y, \catname{XY})$
a $G$-$\rho$ model space of $X \times Y$ such that
for every
$e \in E$ there is a $Y$-$\pi$-bundle chart of $E$ in the coordinate
space $X \times Y$ in $\seqname{A}$ that covers $e$.

Then every chart in $\seqname{A}$ is a $Y$-$\pi$-bundle chart of $E$ in
the coordinate space $X \times Y$.

\begin{proof}
Let $e \in E$, $(U,V,\phi)$ be an arbitrary $Y$-$\pi$-bundle chart of
$E$ in the coordinate space $X \times Y$ in $\seqname{A}$ that covers
$e$ and $(U',V',\phi')$ be an arbitrary m-chart in $\seqname{A}$ that
covers $e$.

By \pagecref{def:Grhomodtriv} the transition function is of the form
$\funcname{t}(x, y) = \bigl ( x, y \star \funcname{g}(x) \bigr )$.
Then $\pi_1 \compose \funcname{t} = \pi_1$,
$\pi_1 \compose \phi' = \pi_1 \compose \phi =\pi$.
\end{proof}
\end{lemma}

\begin{lemma}[$Y$-$\pi$-bundle charts in \ensuremath{\Triv[G-\rho-]{X,Y}}]
\label{lem:GrhomodtrivYpichart}
Let $X$ and $Y$ be topological spaces, $G$ a topological group, $\rho$
an effective group action of $G$ on $Y$ and
$\pi \defeq \pi_1 \maps X \times Y \to Y$.
$C \defeq (U, V \times Y, \phi)$ is a $Y$-$\pi$-bundle chart of $E$ in
the coordinate space $X \times Y$ iff $C$ is an m-chart of $E$ in the
coordinate space $\Triv[G-\rho-]{X,Y}$.

\end{lemma}

\begin{lemma}[Properties of projection]
\label{def:piepi}
Let $(U, V \times Y, \phi)$ be a $Y$-$\pi$-bundle chart of $E$ in the
coordinate space $X \times Y$ and $v \in V$. Then

$\pi \restrictto_U$ is a surjection.
\begin{proof}
Let $v \in V$, $y \in Y$ and
$u \defeq \phi^{-1} \bigl ( (v,y) \bigr ) \in U$. Then $\pi(u) = v$.
\end{proof}
$\pi^{-1}[\set{v}]$ is homeomorphic to $Y$.
\begin{proof}
$\phi$ and $\phi^{-1}$ are homeomorphisms, so their restrictions are
homeomorphisms and thus $\phi^{-1}[\set{v}]$ is homeomorphic to
$\set{v} \times Y$, which is homeomorphic to $Y$.
\end{proof}
\end{lemma}

\begin{definition}[$Y$-$\pi$ subcharts]
\label{def:Ypisub}
Let $(U, V \times Y, \phi)$ be a $Y$-$\pi$-bundle chart of $E$ in the
coordinate space $X \times Y$ and $U'$ be a nonvoid open subset of $U$.
Then
$(U', V' \times Y, \phi') \defeq (U', \phi[U'], \phi \restrictto_{U'})$
is a subchart of $(U, V \times Y, \phi)$.
\end{definition}

\begin{lemma}[$Y$-$\pi$ subcharts]
Let $(U, V \times Y, \phi)$ be a $Y$-$\pi$-bundle chart of $E$ in the
coordinate space $X \times Y$ and $(U', V' \times Y, \phi')$ a subchart
of $(U, V \times Y, \phi)$. Then $(U', V' \times Y, \phi')$ is a
$Y$-$\pi$-bundle chart of $E$ in the coordinate space $X \times Y$.
\begin{proof}
$(U', V' \times Y, \phi')$ satisfies the conditions of \cref{def:Ypichart}:
\begin{enumerate}
\item $U' \subseteq E$ is open.
\item $\phi[U']$ os open since $\phi$ is a homeomorphism.
\item $\phi \restrictto_{U'} \maps U' \to V' \times Y$ is the
restriction of a homeomorphism and thus a homeomorphism.
$\phi \restrictto_{U'}$ preserves fibers because $\phi$ does.
\end{enumerate}
\end{proof}
\end{lemma}

\begin{definition}[$G$-$\rho$-compatibility]
\label{def:Grho-compat}
Let $E,X,Y$ be topological spaces, $G$ a topological group,
$\rho \maps Y \times G \to Y$ an effective right action of
$G$ on $Y$, $\pi \maps E \onto X$ surjective
and $(U^\mu, V^\mu \times Y, \phi^\mu)$,
$(U^\nu, V^\nu \times Y, \phi^\nu)$ $Y$-$\pi$-bundle charts. Then
$(U^\mu, V^\mu \times Y, \phi^\mu)$ and
$(U^\nu, V^\nu \times Y, \phi^\nu)$ are $G$-$\rho$-compatible iff either
\begin{enumerate}
\item $U^\mu$ and $U^\nu$ are disjoint
\item
The transition function
$
  \funcname{t}^\mu_\nu \defeq
    \phi^\mu \compose
    {\phi^\nu}^{-1} \restrictto_{\phi^\nu[U^\mu \cap U^\nu]}
$
is generated by the group action, i.e., there is a continuous function
$\funcname{g}^\mu_\nu \maps \pi_1[\phi^\nu[U^\mu \cap U^\nu]] \to G$
such that
\begin{equation}
\showlabelsinline
\label{eq:transitiongroup}
\uquant
  {{(x,y) \in \phi^\nu[U^\mu \cap U^\nu]}}
  {
    \funcname{t}^\mu_\nu(x,y) =
    \bigl ( x, y \star \funcname{g}^\mu_\nu(x) \bigr )
  }
\end{equation}
\end{enumerate}
\end{definition}

\begin{lemma}[$G$-$\rho$-compatibility]
\label{lem:Grho-compat}
Let $E,X,Y$ be topological spaces, $G$ a topological group,
$\rho \maps Y \times G \to Y$ an effective right action of
$G$ on $Y$, $\pi \maps E \onto X$ surjective
and $(U^i, V^i \times Y, \phi^i)$, $i =1,2$ $Y$-$\pi$-bundle charts.
Then $(U^1, V^1 \times Y, \phi^1)$ and
$(U^2, V^2 \times Y, \phi^2)$ are $G$-$\rho$-compatible iff
$(U^1, V^1 \times Y, \phi^1)$ and
$(U^2, V^2 \times Y, \phi^2)$ are compatible m-charts of
$\Triv{E}$ in the coordinate space
$\Triv[G-\rho-]{X,Y}$.

\begin{proof}
\begin{enumerate}
\leavevmode
\item
Each $U^i$ is a model neighborhood of $\Triv{E}$ and each $V^i \times Y$
is a model neighborhood of $\Triv[G-\rho-]{X,Y}$.

\item
If $U^1 \cap U^2 = \emptyset$ then the lemma is trivially true.

\item
If the charts are $G$-$\rho$-compatible then the transition function
$
  \funcname{t}^1_2 \defeq
    \phi^1 \compose
    {\phi^2}^{-1} \restrictto_{\phi^2[U^1 \cap U^2]}
$
is generated by the group action, i.e., there is a continuous function
$\funcname{g}^1_2 \maps \pi_1[\phi^2[U^1 \cap U^2]] \to G$
such that
$
  \uquant
    {{(x,y) \in \phi^2[U^1 \cap U^2]}}
    {
      \funcname{t}^1_2(x,y) =
      \bigl ( x, y \star \funcname{g}^1_2(x) \bigr )
    }
$.
Then $\pi_1 \compose \funcname{t}^1_2 = \pi_1$, $\funcname{t}^1_2$ is a
morphism of $\Triv[G-\rho-]{X,Y}$ and the charts are compatible m-charts
of $\Triv[G-\rho-]{X,Y}$.

\item
If the charts are compatible m-charts of $\Triv[G-\rho-]{X,Y}$ then the
transition function is an isomorphism of $\Triv[G-\rho-]{X,Y}$, hence
generated by the group action, and the charts are $G$-$\rho$-compatible.
\end{enumerate}
\end{proof}
\end{lemma}

\begin{corollary}[Symmetry of $G$-$\rho$ compatibility]
\label{cor:Grho-compatsym}
Let $(U^\mu, V^\mu \times Y, \phi^\mu)$ and \linebreak
$(U^\nu, V^\nu \times Y, \phi^\nu)$ be $Y$-$\pi$-bundle charts of $E$ in the
coordinate space $X \times Y$. Then
$(U^\mu, V^\mu \times Y, \phi^\mu)$ is $G$-$\rho$-compatible with
$(U^\nu, V^\nu \times Y, \phi^\nu)$ iff
$(U^\nu, V^\nu \times Y, \phi^\nu)$ is $G$-$\rho$-compatible with
$(U^\mu, V^\mu \times Y, \phi^\mu)$.
\begin{proof}
The result follows from \pagecref{lem:M-compatsym}.
\end{proof}
\end{corollary}

\begin{lemma}[$G$-$\rho$-compatibility of subcharts]
Let $(U^i, V^i \times Y, \phi^i)$ be a $Y$-$\pi$-bundle chart of $E$ in
the coordinate space $X \times Y$, $(U'^i, V'^i \times Y, \phi'^i)$ a
subchart and $(U^1, V^1, \phi^1)$ be $G$-$\rho$-compatible with
$(U^2, V^2, \phi^2)$.  Then $(U'^1, V'^1 \times Y, \phi'^1)$ is
$G$-$\rho$-compatible with $(U'^2, V'^2 \times Y, \phi'^2)$.

\begin{proof}
If $U^1 \cap U^2 = \emptyset$ then $U'^1 \cap U'^2 = \emptyset$. If
\nobreaks
{{
  $U'^1 \cap U'^2 = \emptyset$
}}
then
\nobreaks
{{
  $(U'^1, V'^1 \times Y, \phi'^1)$
}}
is $G$-$\rho$-compatible with $(U'^2, V'^2 \times Y, \phi'^2)$. Otherwise,
the transition function
\nobreaks
{{
  $
    t^1_2 \defeq
    \phi^1 \compose \phi^{2-1} \restrictto_{\phi^2[U^1 \cap U^2]}
  $
}}
is generated by the group action and hence
\nobreaks
{{
  $
    t^1_2 \restrictto_{\phi^2[U'^1 \cap U'^2]}
    \maps
    \phi^2[U'^1 \cap U'^2]
    \toiso
    \phi^1[U'^1 \cap U'^2]
  $
}}
is generated by the group action.
\end{proof}
\end{lemma}

\begin{corollary}[$G$-$\rho$-compatibility with subcharts]
Let $(U, V \times Y, \phi)$, be a $Y$-$\pi$-bundle chart of $E$ in the
coordinate space $X \times Y$ and $(U', V' \times Y, \phi')$ a subchart.
Then $(U', V' \times Y, \phi')$ is $G$-$\rho$-compatible with
$(U, V \times Y, \phi)$.

\begin{proof}
$(U, V \times Y, \phi)$ is $G$-$\rho$-compatible with itself and is a
subchart of itself,
\end{proof}
\end{corollary}

\begin{definition}[Covering by $Y$-$\pi$-bundle charts]
Let $\seqname{A}$ be a set of $Y$-$\pi$-bundle
charts of $E$ in the coordinate space $X \times Y$.
$\seqname{A}$ covers $E$ iff $E = \union{{\pi_1[\seqname{A}]}}$.
\end{definition}

\begin{lemma}
Let $\seqname{A}$ be a set of $Y$-$\pi$-bundle charts of $E$ in the
coordinate space $X \times Y$ that covers $E$ and $x \in X$. Then
$\pi^{-1}[\set{x}]$ is homeomorphic to $Y$.
\begin{proof}
Since $\seqname{A}$ covers $E$, there is a chart $(U, V, \phi)$ in
$\seqname{A}$ containing $x$.
Then $\pi^{-1}[\set{x}]$ is homeomorphic to $Y$ by
\pagecref{def:piepi}.
\end{proof}
\end{lemma}

\section{Bundle atlases}
\label{sec:bunatlases}
A set of charts can be atlases for different fiber bundles even if it is
for the same total model space, base space and fiber. In order to
aggregate atlases into categories, there must be a way to distinguish
them. Including the spaces\footnote{
The spaces are redundant, but convenient.},
group and group action in the definitions of the categories serves the
purpose.

\begin{definition}[Bundle atlases]
\label{def:BunAtl}
Let $\seqname{B} \defeq (E, X, Y, \pi, G, \rho)$, be a protobundle. Then
$(\seqname{A},\seqname{B})$ is a bundle atlas, $\seqname{A}$ is a bundle
atlas of $B$, abbreviated $\isAtl[\Bun]_\Ob(\seqname{A}, \seqname{B})$
and $\seqname{A}$ is a $\pi$-$G$-$\rho$-bundle atlas of $E$ in the
coordinate space $X \times Y$, abbreviated
$\isAtl[\Bun]_\Ob(\seqname{A}, E, X, Y, \pi, G, \rho)$, iff
$\seqname{A}$ consists of a set of mutually $G$-$\rho$-compatible
$Y$-$\pi$-bundle charts of $E$ in the coordinate space $X \times Y$ that
covers $E$\footnote{
  There is no need to introduce the concept of a full
  $\pi$-$G$-$\rho$-bundle atlas because a $\pi$-$G$-$\rho$-bundle
  atlas is automatically full.
}.

By abuse of language we write $U \in \seqname{A}$ for
$U \in \pi_1[\seqname{A}]$.

\begin{remark}
The definition of a $\pi$-$G$-$\rho$-bundle atlas is by design similar
to the definition of a coordinate bundle in \cite[p.~7]{TopFib}, but
there are significant differences. This paper will use the term bundle
atlas to avoid confusion.
\end{remark}

Let
$
  \seqname{B}^\alpha \defeq
  (E^\alpha, X^\alpha, Y^\alpha, \pi^\alpha, G^\alpha, \rho^\alpha)
$,
$\alpha \prec \Alpha$, be a protobundle and
$\seqname{B} \defeq \set{{\seqname{B}^\alpha}}[\alpha \prec \Alpha]$
be a set of protobundles.
Then
\begin{multline}
\Atl[\Bun]_\Ob(\seqname{B}^\alpha) \defeq
  \set
  {{
    (
      \seqname{A},
      \seqname{B}^\alpha
    )
  }}%
  [
    {{
      \isAtl[\Bun]_\Ob
        (
          \seqname{A},
          E^\alpha,
          X^\alpha,
          Y^\alpha,
          G^\alpha,
          \pi^\alpha,
          \rho^\alpha
        )
    }}
  ]
\end{multline}

\begin{equation}
\Atl[\Bun]_\Ob \seqname{B} \defeq
  \union[\alpha \prec \Alpha]{\Atl[\Bun]_\Ob(\seqname{B}^\alpha)}
\end{equation}
\end{definition}

\begin{lemma}[Bundle atlases]
\label{lem:BunAtl}
Let $E,X,Y$ be topological spaces, $G$ a topological group,
$\pi \maps E \onto X$ surjective and $\rho \maps Y \times G \to Y$ an
effective right action of $G$ on $Y$. Then $\seqname{A}$ is a
$\pi$-$G$-$\rho$-bundle atlas of $E$ in the coordinate space
$X \times Y$ iff it is an m-atlas of $\Triv{E}$ in the coordinate space
$\Triv[G-\rho-]{X,Y}$ and every coordinate function preserves fibers.

\begin{proof}
If $\seqname{A}$ is a $\pi$-$G$-$\rho$-bundle atlas of $E$ in the
coordinate space $X \times Y$ then
\begin{enumerate}
\item Every chart in $\seqname{A}$ is a $Y$-$\pi$-bundle chart of $E$ in
the coordinate space $X \times Y$, and hence its coordinate function
preserves fibers.
\item The charts in $\seqname{A}$ are mutually $G$-$\rho$-compatible;
hence the transition functions are generated by the group action and are
morphisms of $\Triv[G-\rho-]{X,Y}$.
\end{enumerate}

If $\seqname{A}$ is an m-atlas of $\Triv{E}$ in the coordinate space
$\Triv[G-\rho-]{X,Y}$ then the transition functions are generated by the
group action and and thus the charts are mutually $G$-$\rho$-compatible.

If every coordinate function preserves fibers then the m-charts of
$\seqname{A}$ are $Y$-$\pi$-bundle charts.
\end{proof}
\end{lemma}

\begin{definition}[Compatibility of charts with bundle atlases]
A $Y$-$\pi$-bundle chart $(U, V \times Y, \phi)$ of $E$ in the
coordinate space $X \times Y$ is $G$-$\rho$-compatible with a
$\pi$-$G$-$\rho$-bundle atlas $\seqname{A}$ iff it is
$G$-$\rho$-compatible with every chart in the atlas.
\end{definition}

\begin{lemma}[Compatibility of subcharts with bundle atlases]
\label{lem:BunCompat}
Let $\seqname{A}$ be a $\pi$-$G$-$\rho$-bundle atlas of $E$ in the
coordinate space $X \times Y$ and
$\seqname{C}^1 = (U, V \times Y, \phi)$ a $Y$-$\pi$-bundle chart in
$\seqname{A}$. Then any subchart of $\seqname{C}^1$ is
$G$-$\rho$-compatible with $\seqname{A}$.

\begin{proof}
Let $\seqname{C}'^1 = (U'^1, V'^1 \times Y, \phi'^1)$ be a subchart of
$\seqname{C}^1$ and $\seqname{C}^2 = (U^2, V^2 \times Y, \phi^2)$
another chart in $\seqname{A}$.

\begin{enumerate}
\item If $U^1 \cap U^2 = \emptyset$, then $U' \cap U^2 = \emptyset$.
\item If $U' \cap U^2 = \emptyset$ then $\seqname{C}'^1$ is
$G$-$\rho$-compatible with $\seqname{C}^2$.
\item Otherwise the transition function
$
  t^1_2 \defeq
  \phi^1 \compose \phi^{2-1} \restrictto_{\phi^1[U^1 \cap U^2]}
$
is generated by the group action and thus
$t^1_2 \restrictto_{\phi^2[U'^1 \cap U^2]}$ is generated by the group
action.
\end{enumerate}
\end{proof}
\end{lemma}

\begin{lemma}[Extensions of bundle atlases]
\label{def:piGrho-atl:extensions}
Let $\seqname{A}$ be a $\pi$-$G$-$\rho$-atlas of $\seqname{E}$ in the
coordinate space $X \times Y$ and $(U_i,V_i,\phi_i)$, $i=1,2$, be a
$\pi$-$G$-$\rho$ chart of $E$ in the coordinate space $X \times Y$
$G$-$\rho$-compatible with $\seqname{A}$ in the coordinate space $X
\times Y$.  Then $(U_1,V_1,\phi_1)$ is $G$-$\rho$-compatible with
$(U_2,V_2,\phi_2)$ in the coordinate space $X \times Y$.

\begin{proof}
If $U_1 \cap U_2 = \emptyset$ then $(U_1,V_1,\phi_1)$ is
$G$-$\rho$-compatible with $(U_2,V_2,\phi_2)$. Otherwise,
$
  \phi_2 \compose \phi^{-1}_1 \restrictto_{\phi_1[U_1 \cap U_2]}
  \maps
  \phi_1[U_1 \cap U_2] \toiso
  \phi_2[U_1 \cap U_2]
$
is a homeomorphism.  It remains to show that
$\phi_2 \compose \phi^{-1}_1 \restrictto_{\phi_1[U_1 \cap U_2]}$ is
generated by the group action, Let $(U'_\alpha,V'_\alpha,\phi'_\alpha)$,
$\alpha \prec \Alpha$, be charts in $\seqname{A}$ such that
$U_1 \cap U_2 \subseteq \union[\alpha \prec \Alpha]{U'_\alpha}$ and
$U_1 \cap U_2 \cap U'_\alpha \neq \emptyset$, $\alpha \prec \Alpha$.
Since the charts are $G$-$\rho$-compatible with
$(U'_\alpha,V'_\alpha,\phi'_\alpha)$,
$
\phi_2 \compose \phi'^{-1}_\alpha \restrictto_{U_1 \cap U_2 \cap
U'_\alpha}
$
and
$\phi'_\alpha \compose \phi^{-1}_1 \restrictto_{U^1 \cap U^2 \cap
U'_\alpha}
$
are generated by the group action and thus
$\phi^2 \compose \phi^{-1}_1 = \phi^2 \compose \phi'^{-1}_\alpha \compose
\phi'_\alpha \compose \phi^{-1}_1$
is generated by the group action.
\end{proof}
\end{lemma}

\begin{definition}[Maximal bundle atlases]
\label{def:Bun-ATLmax}
Let $\seqname{A}$ be a $\pi$-$G$-$\rho$-bundle atlas of $E$ in the
coordinate space $X \times Y$. $\seqname{A}$ is a maximal
$\pi$-$G$-$\rho$-bundle atlas, abbreviated \\*
$
  \maximal{\isAtl[\Bun]_\Ob}
    (
      \seqname{A},
      E,
      X,
      Y,
      G,
      \pi,
      \rho
    )
$,
iff it cannot be extended by adding an additional $G$-$\rho$-compatible
$Y$-$\pi$-bundle chart.

$\seqname{A}$ is a semi-maximal $\pi$-$G$-$\rho$-bundle atlas
of $\seqname{E}$ in the coordinate space $\seqname{C}$, abbreviated
$
  \maximal[S-]{\isAtl[\Bun]_\Ob}
    (
      \seqname{A},
      E,
      X,
      Y,
      G,
      \pi,
      \rho
    )
$,
iff whenever $(U, V \times Y, \phi) \in \seqname{A}$,
$U' \subseteq U, V' \times Y \subseteq V \times Y$ and
$V'' \times Y \subseteq X \times Y$ are open sets,
$\phi[U'] = V' \times Y$,
$\phi' \maps V' \times Y \toiso V'' \times Y$ is a fiber
preserving homeomorphism generated by the group action then
$(U', V'' \times Y, \phi' \compose \phi) \in \seqname{A}$.

Let
$
  \seqname{B}^\alpha \defeq
  (E^\alpha, X^\alpha, Y^\alpha, \pi^\alpha, G^\alpha, \rho^\alpha)
$,
$\alpha \prec \Alpha$, be a protobundle and
$\seqname{B} \defeq \set{{\seqname{B}^\alpha}}[\alpha \prec \Alpha]$
be a set of protobundles.
Then
\begin{multline}
\maximal[**]{\Atl[\Bun]_\Ob(\seqname{B}^\alpha)} \defeq
  \set
  {{
    (
      \seqname{A},
      \seqname{B}^\alpha
    )
  }}%
  [
    {{
      \maximal[**]
        {
          \isAtl[\Bun]_\Ob
            \biggl (
              \seqname{A},
              E^\alpha,
              X^\alpha,
              Y^\alpha,
              G^\alpha,
              \pi^\alpha,
              \rho^\alpha
            \biggr )
        }
    }}
  ]
\end{multline}

\begin{equation}
\maximal[**]{\Atl[\Bun]_\Ob} \seqname{B} \defeq
  \union
    [\alpha \prec \Alpha]
    {\maximal[**]{\Atl[\Bun]_\Ob}(\seqname{B}^\alpha)}
\end{equation}
\end{definition}

\begin{lemma}[Maximal $\pi$-$G$-$\rho$-bundle atlases are semi-maximal $\pi$-$G$-$\rho$-bundle atlases]
Let $E,X,Y$ be topological spaces, $G$ a topological group,
$\pi \maps E \onto X$ surjective, $\rho \maps Y \times G \to Y$ an
effective right action of $G$ on $Y$ and $\seqname{A}$ a maximal
$\pi$-$G$-$\rho$-bundle atlas of $E$ in the coordinate space $X \times Y$.
Then $\seqname{A}$ is a semi-maximal $\pi$-$G$-$\rho$-bundle atlas of $E$ in the
coordinate space $X \times Y$.

\begin{proof}
Let $(U, V \times Y, \phi) \in \seqname{A}$,
$U' \subseteq U, V' \times Y \subseteq V \times Y$ and
$V'' \times Y \subseteq X \times Y$ be open sets and
$\phi[U'] = V' \times Y$, $\phi' \maps V' \times Y \toiso V'' \times Y$
be a fiber preserving homeomorphism generated by the group action.
$(U', V', \phi)$ is a subchart of $(U,V,\phi)$ and by
\pagecref{lem:BunCompat} is $G$-$\rho$-compatible with the charts of
$\seqname{A}$.  Since $\phi'$ is a fiber preserving homeomorphism
generated by the group action, $(U', V'' \times Y, \phi' \compose \phi)$
is $G$-$\rho$-compatible with the charts of $\seqname{A}$. Since
$\seqname{A}$ is maximal, $(U', V'', \phi' \compose \phi)$ is a chart of
$\seqname{A}$.
\end{proof}
\end{lemma}

\begin{theorem}[Existence and uniqueness of maximal $\pi$-$G$-$\rho$-bundle atlases]
Let $\seqname{B} \defeq (E, X, Y, \pi, G, \rho)$ be a protobundle and
$\seqname{A}$ $\pi$-$G$-$\rho$-bundle atlas of $E$ in the coordinate
space $X \times Y$. Then there exists a unique maximal
$\pi$-$G$-$\rho$-bundle atlas
$\maximal{\Atlas^\Bun}(\seqname{A},\seqname{B})$ of $E$ in the
coordinate space $X \times Y$ $G$-$\rho$-compatible with $\seqname{A}$.

\begin{proof}
Let $\seqname{P}$ be the set of all $\pi$-$G$-$\rho$-bundle atlases $E$
in the coordinate space $X \times Y$ containing $\seqname{A}$ and
$G$-$\rho$ compatible in the coordinate space $X \times Y$ with
$\seqname{A}$. Let $\maximal{\seqname{P}}$ be a maximal chain of
$\seqname{P}$. Then $A'=\union{\maximal{\seqname{P}}}$ is a maximal
$\pi$-$G$-$\rho$-bundle atlas of $\seqname{E}$ in the coordinate space
$X \times Y$ $G$-$\rho$ compatible with $\seqname{A}$. Uniqueness
follows from \pagecref{def:piGrho-atl:extensions}.
\end{proof}
\end{theorem}

\begin{lemma}[Existence and uniqueness of projection for atlases in $G$-$\rho$-model spaces]
\label{lem:mGrho}
Let $E$, $X$ and $Y$ be topological spaces, $G$ a topological group,
$\rho$ an effective action of $G$ on $Y$,
$\seqname{C} \defeq (X \times Y, \catname{XY})$ a $G$-$\rho$ model space
of $X \times Y$ and $\seqname{A}$ an m-atlas of $E$ in the coordinate
model space $\seqname{C}$.  Then there exists a unique function
$\pi \maps \seqname{E} \to X$ such that for any chart $(U, V, \phi)$ in
$\seqname{A}$, $\pi \restrictto_U = \pi_1 \compose \phi$.   If
$\seqname{A}$ is full then $\pi$ is surjective.

\begin{proof}
Let $(U, V, \phi)$ be an arbitrary chart in $\seqname{A}$ and define
$\pi(e \in U) = \pi_1 \compose \phi(e)$. $\pi(e)$ does not depend on the
choice of chart because the morphisms of a $G$-$\rho$ model space
preserve fibers. $\pi$ is continuous because it is continuous on each
coordinate patch.

Let $x \in X$. If $\seqname{A}$ is full then there exists a chart
$(U, V, \phi)$ in $\seqname{A}$ such that $x \in \pi_1[V]$. Let $u$ be an
arbitrary point in $\phi^{-1}[\set{x} \times Y]$. Then
$\pi(u) = \pi_1(\phi(u)) = x$.
\end{proof}
\end{lemma}
\section{Bundle-atlas (near) morphisms and functors}
\label{sec:bunatlmorph}
This section introduces a taxonomy for morphisms between bundle atlases,
defines categories of bundle atlases and constructs functors from them
to categories of m-atlases.  It only constructs reverse functors for
subcategories of $\Trivcat[\BunProd-]{\seqname{B}}$.

\subsection{Bundle-atlas (near) morphisms}
\subsubsection{Definitions of bundle-atlas (near) morphisms}
\begin{definition}[Bundle-atlas near morphisms]
\label{def:Bun-ATLmorphnear}
Let
$
  \seqname{B}^i
  \defeq
  (
    E^i,
    X^i,
    Y^i,
    \pi^i,
    G^i,
    \rho^i
  )
$, \\*
$i=1,2$,
be a protobundle and let $\seqname{A}^i$ be a
$\pi^i$-$G^i$-$\rho^i$-atlas of $E^i$ in the coordinate space
\nobreaks
{{
  $C^i = X^i \times Y^i$.
}}
Then
$
  \funcseqname{f}
  \defeq
  (
    \funcname{f}_E \maps E^1 \to E^2,
    \funcname{f}_X \maps X^1 \to X^2,
    \funcname{f}_Y \maps Y^1 \to Y^2,
    \funcname{f}_G \maps G^1 \to G^2
  )
$
is a $\seqname{B^1}$-$\seqname{B^2}$ bundle-atlas morphism from
$\seqname{A}^1$ to $\seqname{A}^2$, abbreviated \\
\nobreaks
{{
  $
    \isAtl[\Bun,near]_\Ar
      (
        \seqname{A^1}, \seqname{B^1},
        \seqname{A^2}, \seqname{B^2},
        \funcseqname{f}
      )
  $,
}}
iff
\begin{enumerate}
\item all four functions are continuous
\item $\funcname{f}_G$ is a homomorphism
\item $\funcseqname{f}$ commutes with $\pi^i$ and $\rho^i$, i.e.,
\begin{enumerate}
\item $\pi^2 \compose \funcname{f}_E = \funcname{f}_X \compose \pi^1$
\item
$
  \uquant
    {{y \in Y^1},{g \in G^1}}
    {{
      \funcname{f}_Y(y) \star^2 \funcname{f}_G(g) =
      \funcname{f}_Y(y \star^1 g)
    }}
$
\end{enumerate}
\item For any charts
$(U^i, V^i, \phi^i \maps U^i \toiso V^i) \in \seqname{A}^i$. $i=1,2$,
with $I \defeq U^1 \cap \funcname{f}^{-1}_E[U^2] \neq \emptyset$,
diagram \eqref{eq:Bun-AtlN1} below
is locally nearly commutative in $X,Y,\pi,\rho$, i.e., for any
$u^1 \in I$ there are open sets $U'^1 \subseteq I$,
$V'^1 \subseteq V^1$, $U'^2 \subseteq U^2$, $V'^2 \subseteq V^2$,
$\hat{V}'^2 \subseteq X^2 \times Y^2$ and a homeomorphism
$\hat{\funcname{f}} \maps V'^2 \toiso \hat{V}'^2$%
\footnote
{
  This reverses the direction of the arrow
  $\hat{\funcname{f}} \maps U_m \toiso V_n$ from
  \pagecref{sec:grhoncd}.
  This is permissible since $\hat{\funcname{f}}$ is a homeomorphism
}
such that \crefrange{eq:Bun-xin}{eq:Bun-fphigeqgphi} below hold, as
shown in
{
  \showlabelsinline
  \cref{fig:Bun-AtlN1,fig:Bun-AtlN2}
}.

\begin{equation}
\label{eq:Bun-AtlN1}
\begin{array} {lll}
D \defeq & \bigl \{
\\
  &&
  \funcname{f}_E \maps I \defeq U^1 \cap \funcname{f}^{-1}_E[U^2] \to U^2,
  \phi^2 \maps U^2 \toiso V^2,
\\
  &&
  \phi^1 \maps I \into V^1,
  \funcname{f}_X \times \funcname{f}_Y \maps V^1 \to X^2 \times Y^2
\\
& \bigr \}
\end{array}
\end{equation}

\begin{equation}
u^1 \in U'^1
\label{eq:Bun-xin}
\end{equation}

\begin{equation}
\label{eq:Bun-fup1in}
\funcname{f}_E[U'^1] \subseteq U'^2
\end{equation}

\begin{equation}
\phi^1[U'^1] = V'^1
\end{equation}

\begin{equation}
\funcname{f}_X \times \funcname{f}_Y[V'^1] \subseteq \hat{V}'^2
\label{eq:Bun-fvp1in}
\end{equation}

\begin{equation}
\phi^2[U'^2] = V'^2
\end{equation}

\begin{equation}
\hat{\funcname{f}} \compose \phi^2 \compose \funcname{f}_E =
\funcname{f}_X \times \funcname{f}_Y \compose \phi^1
\label{eq:Bun-fphigeqgphi}
\end{equation}

\begin{remark}
$(U'^i, V'^i, \phi^i \restrictto_{U'^i, V'^i})$ need not be a chart of
$\seqname{A}^i$.
\end{remark}

\begin{figure}[ht]
\[ \bfig
\node uint(0,0)[{I = U^1 \cap \funcname{f}^{-1}_0[U^2]}]
\node u2(2000,0)[U^2]
\node v1(0,-1000)[V^1]
\node c21(1000,-1000)[C^2]
\node v2(2000,-1000)[V^2]
\arrow |a|[uint`u2;\funcname{f}_E]
\arrow |r|/>->/[uint`v1;\phi^1]
\arrow |r|[u2`v2;\phi^2]
\arrow |l|//[u2`v2;\iso]
\arrow |a|[v1`c21;\funcname{f}_X \times \funcname{f}_Y]
\efig \]
\caption{Uncompleted m-atlas near morphism}
\label{fig:Bun-AtlN1}
\end{figure}

\begin{figure}[ht]
\[ \bfig
\node x(500,0)[u^1]
\node uint(-300,-500)[{I=U^1 \cap \funcname{f}^{-1}_0[U^2]}]
\node up1(500,-500)[U'^1]
\node up2(1500,-500)[U'^2]
\node u2(2000,-500)[U^2]
\node v1(-300,-1000)[V^1]
\node vp1(500,-1000)[V'^1]
\node vhp2(1000,-1000)[\hat{V}'^2]
\node vp2(1500,-1000)[V'^2]
\node v2(2000,-1000)[V^2]
\arrow |r|/^{ (}->/[x`up1;\funcname{i}]
\arrow |r|/>->/[uint`v1;\phi^1]
\arrow |a|/_{ (}->/[up1`uint;\funcname{i}]
\arrow |a|[up1`up2;\funcname{f}_E]
\arrow |r|[up1`vp1;{\phi^1}]
\arrow |l|//[up1`vp1;\iso]
\arrow |a|/^{ (}->/[up2`u2;i]
\arrow |r|[up2`vp2;\phi^2]
\arrow |l|//[up2`vp2;\iso]
\arrow |r|[u2`v2;\phi^2]
\arrow |l|//[u2`v2;\iso]
\arrow |a|/_{ (}->/[vp1`v1;\funcname{i}]
\arrow |a|[vp1`vhp2;\funcname{f}_X \times \funcname{f}_Y]
\arrow |a|/ >.>>/[vp2`vhp2;\hat{\funcname{f}}]
\arrow |b|//[vp2`vhp2;\iso]
\arrow |a|/^{ (}->/[vp2`v2;\funcname{i}]
%\arrow |a|/^{ (}->/[v2`c2;\funcname{i}]
\efig \]
\caption{Completed m-atlas near morphism}
\label{fig:Bun-AtlN2}
\end{figure}
\end{enumerate}

If $\seqname{A}^1$ and $\seqname{A}^2$ are semi-maximal (maximal)
atlases then $\funcseqname{f}$ is also a semi-maximal (maximal)
$\seqname{B^1}$-$\seqname{B^2}$ bundle-atlas near morphism from
$\seqname{A}^1$ to $\seqname{A}^2$, abbreviated \\*
$
  \maximal[**] {\isAtl[\Bun,near]_\Ar}
    (
      \seqname{A^1}, \seqname{B^1},
      \seqname{A^2}, \seqname{B^2},
      \funcseqname{f}
    )
$

The identity morphism of $(\seqname{A}^i, \seqname{B}^i)$ is
\begin{equation}
\Id[{{(\seqname{A}^i, \seqname{B}^i)}}] \defeq
\Bigl (
  \bigl (
    \Id[{E^i}],
    \Id[{X^i}],
    \Id[{Y^i}],
    \Id[{G^i}]
  \bigr ),
  \bigl ( \seqname{A}^i, \seqname{B}^i\bigr ),
  \bigl ( \seqname{A}^i, \seqname{B}^i\bigr )
\Bigr )
\end{equation}
This nomenclature will be justified later.
\end{definition}

\begin{definition}[Bundle-atlas morphisms]
\label{def:Bun-ATLmorph}
Let
$
  \seqname{B}^i
  \defeq
  (
    E^i,
    X^i,
    Y^i,
    \pi^i,
    G^i,
    \rho^i
  )
$,
$i=1,2$, be a protobundle and let $\seqname{A}^i$ be a
$\pi^i$-$G^i$-$\rho^i$-atlas of $E^i$ in the coordinate space
\nobreaks
  {{
    $C^i \defeq X^i \times Y^i$.
  }}
Then
$
  \funcseqname{f}
  \defeq
  (
    \funcname{f}_E \maps E^1 \to E^2,
    \funcname{f}_X \maps X^1 \to X^2,
    \funcname{f}_Y \maps Y^1 \to Y^2,
    \funcname{f}_G \maps G^1 \to G^2
  )
$
is a $\seqname{B^1}$-$\seqname{B^2}$ bundle-atlas morphism from
$\seqname{A}^1$ to $\seqname{A}^2$, abbreviated \\
$
  \isAtl[\Bun]_\Ar
    (
      \seqname{A^1}, \seqname{B^1},
      \seqname{A^2}, \seqname{B^2},
      \funcseqname{f}
    )
$,
iff
\begin{enumerate}
\item all four functions are continuous
\item $\funcname{f}_G$ is a homomorphism
\item $\funcseqname{f}$ commutes with $\pi^i$ and $\rho^i$, i.e.,
\begin{enumerate}
\item $\pi^2 \compose \funcname{f}_E = \funcname{f}_X \compose \pi^1$
\item
$
  \uquant
    {{y \in Y^1},{g \in G^1}}
    {{
      \funcname{f}_Y(y) \star^2 \funcname{f}_G(g) =
      \funcname{f}_Y(y \star^1 g)
    }}
$
\end{enumerate}
\item For any charts
$(U^i, V^i, \phi^i \maps U^i \toiso V^i) \in \seqname{A}^i$. $i=1,2$,
and any $u^1 \in I \defeq U^1 \cap \funcname{f}^{-1}_E[U^2]$,
there exists a subchart
$
  (U'^1, V'^1, \phi^1 \maps U'^1 \toiso V'^1)
  \in \seqname{A}^1
$
at $u^1$, an open set $U'^2 \subseteq U^2$ and a chart
$
  (U'^2 ,\hat{V}'^2, \phi'^2 \maps U'^2 \toiso \hat{V}'^2)
  \in \seqname{A}^2
$
at $\funcname{f}_E(u^1)$ such that
$\funcname{f}_E[U'^1] \subseteq U'^2$ and diagram \eqref{eq:Bun-AtlM1}
below is commutative, as shown in \cref{fig:Bun-AtlM1}.
\end{enumerate}

\begin{equation}
\begin{array}{lll}
\label{eq:Bun-AtlM1}
D \defeq & \Biggl \{
\\
  &&
  \funcname{i} \maps \set{u^1} \into U'^1,
  \phi^1 \maps U'^1 \toiso V'^1,
  \funcname{f}_X \times \funcname{f}_Y \maps V'^1 \to \hat{V}'^2,
\\
  &&
  \funcname{f}_E \maps  U'^1 \to U'^2,
  \phi'^2 \maps U'^2 \toiso \hat{V}'^2
\\
& \Biggr \}
\end{array}
\end{equation}

\begin{figure}[ht]
\[ \bfig
\node e1(500,0)[u^1]
\node e2(1500,0)[u^2]
\node u1(-300,-500)[U^1]
\node up1(500,-500)[U'^1]
\node up2(1500,-500)[U'^2]
\node u2(2000,-500)[U^2]
\node v1(-300,-1000)[V^1]
\node vp1(500,-1000)[V'^1]
\node vhp2(1500,-1000)[\hat{V}'^2]
\node v2(2000,-1000)[V^2]
\arrow |a|[e1`e2;\funcname{f}_E]
\arrow |r|/^{ (}->/[e1`up1;\funcname{i}]
\arrow |r|/^{ (}->/[e2`up2;\funcname{i}]
\arrow |r|/>->/[u1`v1;\phi^1]
\arrow |l|//[u1`v1;\iso]
\arrow |a|/_{ (}->/[up1`u1;\funcname{i}]
\arrow |a|[up1`up2;\funcname{f}_E]
\arrow |r|[up1`vp1;{\phi^1}]
\arrow |l|//[up1`vp1;\iso]
\arrow |a|/^{ (}->/[up2`u2;i]
\arrow |r|[up2`vhp2;\phi'^2]
\arrow |l|//[up2`vhp2;\iso]
\arrow |r|[u2`v2;\phi^2]
\arrow |l|//[u2`v2;\iso]
\arrow |a|/_{ (}->/[vp1`v1;\funcname{i}]
\arrow |a|[vp1`vhp2;\funcname{f}_X \times \funcname{f}_Y]
\efig \]
\caption{Completed Bundle-atlas morphism}
\label{fig:Bun-AtlM1}
\end{figure}

If $\seqname{A}^1$ and $\seqname{A}^2$ are semi-maximal (maximal)
atlases then $\funcseqname{f}$ is also a semi-maximal (maximal)
$\seqname{B^1}$-$\seqname{B^2}$ bundle-atlas morphism from
$\seqname{A}^1$ to $\seqname{A}^2$, abbreviated \\*
$
  \maximal[**] {\isAtl[\Bun]_\Ar}
    (
      \seqname{A^1}, \seqname{B^1},
      \seqname{A^2}, \seqname{B^2},
      \funcseqname{f}
    )
$

The identity morphism of $(\seqname{A}^i, \seqname{B}^i)$ is
\begin{equation}
\Id[{{(\seqname{A}^i, \seqname{B}^i)}}] \defeq
\Bigl (
\bigl (
  \Id[{E^i}],
  \Id[{X^i}],
  \Id[{Y^i}],
  \Id[{G^i}]
\bigr ),
\bigl ( \seqname{A}^i, \seqname{B}^i\bigr ),
\bigl ( \seqname{A}^i, \seqname{B}^i\bigr )
\Bigr )
\end{equation}
This nomenclature will be justified later.
\end{definition}

\begin{definition}[Equivalence of bundle-atlas (near) morphisms]
\label{def:Bun-ATLmorphequiv}
\leavevmode \\*
$
  \funcseqname{f}
  \defeq
  (
    \funcname{f}_E \maps E^1 \to E^2,
    \funcname{f}_X \maps X^1 \to X^2,
    \funcname{f}_Y \maps Y^1 \to Y^2,
    \funcname{f}_G \maps G^1 \to G^2
  )
$
is bundle equivalent to
$
  \funcseqname{g}
  \defeq
  (
    \funcname{g}_E \maps E^1 \to E^2,
    \funcname{g}_X \maps X^1 \to X^2,
    \funcname{g}_Y \maps Y^1 \to Y^2,
    \funcname{g}_G \maps G^1 \to G^2
  )
$
as a $\seqname{B^1}$-$\seqname{B^2}$ bundle-atlas (near) morphism from
$\seqname{A}^1$ to $\seqname{A}^2$, abbreviated $\funcseqname{f}$ and
$\funcseqname{g}$ are bundle equivalent $\seqname{B^1}$-$\seqname{B^2}$
bundle-atlas (near) morphisms from $\seqname{A}^1$ to $\seqname{A}^2$, iff
$
  \seqname{B}^i
  \defeq
  (
    E^i,
    X^i,
    Y^i,
    \pi^i,
    G^i,
    \rho^i
  )
$,
$i=1,2$, is a protobundle, $\seqname{A}^i$ is a
$\pi^i$-$G^i$-$\rho^i$-atlas of $E^i$ in the coordinate space \\*
$C^i \defeq X^i \times Y^i$, $\funcseqname{f}$ and $\funcseqname{g}$ are
$\seqname{B^1}$-$\seqname{B^2}$ bundle-atlas (near) morphisms from
$\seqname{A}^1$ to $\seqname{A}^2$, $\funcname{f}_E = \funcname{g}_E$,
$\funcname{f}_X = \funcname{g}_X$ and $\funcname{f}_G = \funcname{g}_G$,

By abuse of language we shall write $\funcseqname{f}$ is bundle
equivalent to $\funcseqname{g}$ when the protobundles and bundle
atlases are understood by context.
\end{definition}

\subsubsection{Proclamations on bundle-atlas (near) morphisms}
\begin{lemma}[Bundle-atlas (near) morphisms]
\label{lem:Bun-ATLmorph}
Let
$
  \seqname{B}^i
  \defeq
  (
    E^i,
    X^i,
    Y^i,
    \pi^i,
    G^i,
    \rho^i
  )
$, $i=1,2$, be a protobundle and let
$\seqname{A}^i$ be a $\pi^i$-$G^i$-$\rho^i$-atlas of $E^i$ in the coordinate space
\nobreaks
{{
  $C^i = X^i \times Y^i$.
}}
Then
$
  \funcseqname{f}
  \defeq
  (
    \funcname{f}_E \maps E^1 \to E^2,
    \funcname{f}_X \maps X^1 \to X^2,
    \funcname{f}_Y \maps Y^1 \to Y^2,
    \funcname{f}_G \maps G^1 \to G^2
  )
$
is a $\seqname{B^1}$-$\seqname{B^2}$ bundle-atlas (near) morphism from
$\seqname{A}^1$ to $\seqname{A}^2$ iff
$\funcname{f}_X \times \funcname{f}_Y$ is a
$G^1$-$G^2$-$\rho^1$-$\rho^2$ morphism from
$\Triv[G^1-\rho^1-]{X^1,Y^1}$ to $\Triv[G^2-\rho^2-]{X^2,Y^2}$ and
$(\funcname{f}_E, \funcname{f}_X \times \funcname{f}_Y)$ is a
$\Triv{E^1}$-$\Triv{E^2}$ m-atlas (near) morphism from $\seqname{A}^1$
to $\seqname{A}^2$ in the coordinate spaces
$\Triv[G^1-\rho^1-]{X^1,Y^1}$, $\Triv[G^2-\rho^2-]{X^2,Y^2}$.

\begin{proof}
\leavevmode
\begin{enumerate}
\item
If $\funcseqname{f}$ is a $\seqname{B^1}$-$\seqname{B^2}$ bundle-atlas
morphism then $\funcname{f}_X \times \funcname{f}_Y$ is a model function
from $\Triv[G^1-\rho^1-]{X^1,Y^1}$ to $\Triv[G^2-\rho^2-]{X^2,Y^2}$
by \pagecref{lem:Grhomodtrivmorph}
and $\funcname{f}_G$ is the function asserted to exist in
{
  \showlabelsinline
  \cref{eq:Grho}
}
(preservation of group action) of \pagecref{def:Grhomorph},
so $\funcname{f}_X \times \funcname{f}_Y$ is a
$G^1$-$G^2$-$\rho^1$-$\rho^2$ morphism from
$\Triv[G^1-\rho^1-]{X^1,Y^1}$ to $\Triv[G^2-\rho^2-]{X^2,Y^2}$.

A diagram is locally nearly commutative in $X^2,Y^2,\pi^2,\rho^2$ iff it
is m-locally nearly commutative in
$\pi_2 \biggl ( \Triv[G^2-\rho^2-]{X^2,Y^2} \biggr )$,
thus $(\funcname{f}_E, \funcname{f}_X \times \funcname{f}_Y)$ is an
m-atlas morphism in the coordinate spaces $\Triv[G^1-\rho^1-]{X^1,Y^1}$,
$\Triv[G^2-\rho^2-]{X^2,Y^2}$.

If $\funcname{f}_X \times \funcname{f}_Y$ is a
$G^1$-$G^2$-$\pi^1$-$\pi^2$ morphism from $\Triv[G^1-\rho^1-]{X^1,Y^1}$
to $\Triv[G^2-\rho^2-]{X^2,Y^2}$ then $\funcseqname{f}$ commutes with
$\rho^i$.

\item
If $(\funcname{f}_E, \funcname{f}_X \times \funcname{f}_Y)$ is a
$\Triv{E^1}$-$\Triv{E^2}$ m-atlas morphism from $\seqname{A}^1$ to
$\seqname{A}^2$ in the coordinate spaces $\Triv[G^1-\rho^1-]{X^1,Y^1}$,
$\Triv[G^2-\rho^2-]{X^2,Y^2}$ then $\funcseqname{f}$ commutes with
$\pi^i$.

A diagram is locally nearly commutative in $X^2,Y^2,\pi^2,\rho^2$ iff it
is m-locally nearly commutative in
$\pi_2 \biggl ( \Triv[G^2-\rho^2-]{X^2,Y^2} \biggr )$
thus $(\funcname{f}_E, \funcname{f}_X \times \funcname{f}_Y)$ is a
m-atlas morphism in the coordinate space $\Triv[G^2-\rho^2-]{X^2,Y^2}$,
so $\funcseqname{f}$ is a $\seqname{B^1}$-$\seqname{B^2}$ bundle-atlas
morphism from $\seqname{A}^1$ to $\seqname{A}^2$.
\end{enumerate}
\end{proof}
\end{lemma}

\begin{corollary}[Bundle-atlas morphisms]
\label{cor:Bun-ATLmorph}
Let
$
  \seqname{B}^i
  \defeq
  (
    E^i,
    X^i,
    Y^i,
    \pi^i,
    G^i,
    \rho^i
  )
$,
$i=1,2,3$, be a protobundle,
$\seqname{A}^i$ be a $\pi^i$-$G^i$-$\rho^i$-atlas of $E^i$ in the
coordinate space $C^i = X^i \times Y^i$ and
$
  \funcseqname{f}^i
  \defeq
  (
    \funcname{f}^i_E \maps E^i \to E^{i+1},
    \funcname{f}^i_X \maps X^i \to X^{i+1},
    \funcname{f}^i_Y \maps Y^i \to Y^{i+1},
    \funcname{f}^i_G \maps G^i \to G^{i+1}
  )
$
a $\seqname{B^i}$-$\seqname{B^{i+1}}$ bundle-atlas (near) morphism from
$\seqname{A}^i$ to $\seqname{A}^{i+i}$.

If each $\seqname{A}^i$ is semi-maximal or if each $\funcseqname{f}^i$
is a morphism then $\funcseqname{f}^2 \compose[()] \funcseqname{f}^1$ is
a bundle-atlas (near) morphism from $\seqname{A}^1$ to $\seqname{A}^3$.

\begin{proof}
$
  \funcname{f}^1_X \times \funcname{f} 1_Y \compose
  \funcname{f}^2_X \times \funcname{f}^2_Y
$
is a $G^1$-$G^3$-$\pi^1$-$\pi^3$ morphism from
$\Triv[G^1-\rho^1-]{X^1,Y^1}$ to $\Triv[G^3-\rho^3-]{X^3,Y^3}$ by
\cref{grhomorph:comp} of
{
  \showlabelsinline
  \fullcref{lem:Grhomorph} on
  \cpageref{grhomorph:comp}
}
and
$
  (\funcname{f}^1_E, \funcname{f}^1_X \times \funcname{f}^1_Y)
  \compose[()]
  (\funcname{f}^2_E, \funcname{f}^2_X \times \funcname{f}^2_Y)
$
is a $\Triv{E^1}$-$\Triv{E^3}$ m-atlas morphism from $\seqname{A}^1$ to
$\seqname{A}^3$ in the coordinate space $\Triv[G^3-\rho^3-]{X^3,Y^3}$ by
\cref{M-morphcomp:morphcompmorph} of
{
  \showlabelsinline
  \fullcref{lem:M-ATLmorphcomp} on
  \cpageref{M-morphcomp:morphcompmorph}.
}
\end{proof}
\end{corollary}

\begin{lemma}[Composition of equivalent bundle-atlas (near) morphisms]
\label{lem:Bun-ATLmorphequiv}
Let
\begin{enumerate}
\item
$
  \seqname{B}^i
  \defeq
  (
    E^i,
    X^i,
    Y^i,
    \pi^i,
    G^i,
    \rho^i
  )
$, $i=1,2,$, be a protobundle
\item
$\seqname{A}^i$ be a $\pi^i$-$G^i$-$\rho^i$-atlas of $E^i$ in the coordinate space
$C^i = X^i \times Y^i$
\item
$
  \funcseqname{f}^i \defeq
  (
    \funcname{f}^i_E \maps E^i \to E^{i+1},
    \funcname{f}^i_X \maps X^i \to X^{i+1},
    \funcname{f}^i_Y \maps Y^i \to Y^{i+1},
    \funcname{f}^i_G \maps G^i \to G^{i+1}
  )
$
and
$
  \funcseqname{g}^i \defeq
  (
    \funcname{g}^i_E \maps E^i \to E^{i+1},
    \funcname{g}^i_X \maps X^i \to X^{i+1},
    \funcname{g}^i_Y \maps Y^i \to Y^{i+1},
    \funcname{g}^i_G \maps G^i \to G^{i+1}
  )
$
be bundle equivalent $\seqname{B^i}$-$\seqname{B^{i+1}}$ bundle-atlas
(near) morphisms from $\seqname{A}^1$ to $\seqname{A}^2$
\end{enumerate}

Then $\funcseqname{f}^{i+1} \compose[()] \funcname{f}^1$ is $\Ck$-equivalent
to $\funcseqname{g}^2 \compose[()] \funcname{g}^1$.

\begin{proof}
$\funcname{f}^1_E = \funcname{g}^1_E$, and
$\funcname{f}^2_E = \funcname{g}^2_E$, hence
$
  \funcname{f}^2_E \compose \funcname{f}^1_E =
  \funcname{g}^2_E \compose \funcname{g}^1_E
$.
$\funcname{f}^1_X = \funcname{g}^1_X$, and
$\funcname{f}^2_X = \funcname{g}^2_X$, hence
$
  \funcname{f}^2_X \compose \funcname{f}^1_X =
  \funcname{g}^2_X \compose \funcname{g}^1_X
$.
$\funcname{f}^1_G = \funcname{g}^1_G$, and
$\funcname{f}^2_G = \funcname{g}^2_G$, hence
$
  \funcname{f}^2_G \compose \funcname{f}^1_G =
  \funcname{g}^2_G \compose \funcname{g}^1_G
$.
\end{proof}
\end{lemma}

\subsection{Categories of bundle atlases and functors}
This subsection only defines categories whose morphisms are bundle atlas
morphisms; bundle atlas near morphisms do not satisfy the requiremens for
forming a category.

\begin{definition}[Sets of bundle atlas morphisms]
\label{def:BunATLsets}
Let \\*
$
  \seqname{B}^\alpha \defeq
  (E^\alpha, X^\alpha, Y^\alpha, \pi^\alpha, G^\alpha, \rho^\alpha)
$,
$\alpha \prec \Alpha$, be a protobundle and
$\seqname{B} \defeq \set{{\seqname{B}^\alpha}}[\alpha \prec \Alpha]$
be a set of protobundles.

\begin{multline}
  \maximal[*]
    {
      \Atl[\Bun]_\Ar
        (
          \seqname{B}^\alpha,
          \seqname{B}^\beta
        )
    }
  \defeq
  \set
    {
      \bigl (
        \funcseqname{f},
        (\seqname{A}^\alpha, \seqname{B}^\alpha),
        (\seqname{A}^\beta, \seqname{B}^\beta)
      \bigr )
    }%
    [
      {
        \maximal[*]
          {
            \isAtl[\Bun]_\Ar
              (
                \seqname{A}^\alpha, \seqname{B}^\alpha,
                \seqname{A}^\beta, \seqname{B}^\beta,
                \funcseqname{f}
              )
          }
      }
    ]*
\end{multline}

\begin{multline}
  \maximal[*]
    {
      \Atl[\Bun]_\Ar
        (
          \seqname{B}^\alpha
        )
    }
  \defeq
  \set
    {
      \bigl (
        \funcseqname{f},
        (\seqname{A}^\alpha, \seqname{B}^\alpha),
        (\seqname{A}^\alpha, \seqname{B}^\alpha)
      \bigr )
    }%
    [
      {
        \maximal[*]
          {
            \isAtl[\Bun]_\Ar
              (
                \seqname{A}^\alpha, \seqname{B}^\alpha,
                \seqname{A}^\alpha, \seqname{B}^\alpha,
                \funcseqname{f}
              )
          }
      }
    ]*
\end{multline}

\begin{equation}
\maximal[*]{\Atl[\Bun]_\Ar} \seqname{B}
\defeq
\union
  [
    \seqname{B^\mu \in \seqname{B}},
    \seqname{B}^\nu \in \seqname{B}
  ]
  {
    \maximal[*]{\Atl[\Bun]_\Ar}
      (
        \seqname{B}_\mu,
        \seqname{B}_\nu
      )
  }
\end{equation}
\end{definition}

\begin{definition}[Categories of bundle atlases]
\label{def:Bun-Atlcat}
Let \\*
$
  \seqname{B}^\alpha \defeq
  (E^\alpha, X^\alpha, Y^\alpha, \pi^\alpha, G^\alpha, \rho^\alpha)
$,
$\alpha \prec \Alpha$, be a protobundle and
$\seqname{B} \defeq \set{{\seqname{B}^\alpha}}[\alpha \prec \Alpha]$
be a set of protobundles. Then

\begin{equation}
\maximal[*]
  {
    \Atl[\Bun]
      (
        \seqname{B}^\alpha
      )
  }
\defeq
  \biggl (
    \maximal[*]
      {
        \Atl[\Bun]_\Ob
          (
            \seqname{B}^\alpha
          )
      },
    \maximal[*]
      {
      \Atl[\Bun]_\Ar
        (
          \seqname{B}^\alpha,
          \seqname{B}^\alpha
        )
      },
    \compose[A]
  \biggr )
\end{equation}

\begin{equation}
\maximal[*]{\Atl[\Bun]} \seqname{B}
\defeq
  \biggl (
    \maximal[*]{\Atl[\Bun]_\Ob}
      (
        \seqname{B}
      ),
    \maximal[*]{\Atl[\Bun]_\Ar}
      (
        \seqname{B}
      ),
    \compose[A]
  \biggr )
\end{equation}
\end{definition}

\begin{lemma}[\ensuremath{\Atl[\Bun] \seqname{B}} is a category]
\label{lem:BunATLiscat}
Let
$
  \seqname{B}^\alpha \defeq
  (E^\alpha, X^\alpha, Y^\alpha, \pi^\alpha, G^\alpha, \rho^\alpha)
$,
$\alpha \prec \Alpha$, be a protobundle and
$\seqname{B} \defeq \set{{\seqname{B}^\alpha}}[\alpha \prec \Alpha]$
be a set of protobundles.
Then $\maximal[*]{\Atl[\Bun] \seqname{B}}$ is a category.

Let
$
  (\seqname{A}^\alpha, \seqname{B}^\alpha)
  \in
  \Atl[\Bun]_\Ob \seqname{B}
$.
Then $\Id[{{(\seqname{A}^\alpha, \seqname{B}^\alpha)}}]$ is the
identity morphism for $(\seqname{A}^\alpha, \seqname{B}^\alpha)$.

\begin{proof}
Let $(\seqname{A}^i,\seqname{B}^i)$, $i \in [1,3]$,
be an object of $\Atl[\Bun] \seqname{B}$ and
let \\
$
  \funcname{m}^i \defeq
  \bigl (
    \funcseqname{f}^i,
    (\seqname{A}^i,\seqname{B}^i),
    (\seqname{A}^{i+1},\seqname{B}^{i+1})
  \bigr )
$
be a morphism of $\Atl[\Bun] \seqname{B}$. Then
\begin{description}
\item[Composition:]
$
  \bigl (
    \funcseqname{f}^2 \compose[()] \funcseqname{f}^1,
    (\seqname{A}^1,\seqname{E}^1,\seqname{C}^1),
    (\seqname{A}^3,\seqname{E}^3,\seqname{C}^3)
  \bigr )
$
is a morphism of $\Atl[\Bun] \seqname{B}$ by
\pagecref{cor:Bun-ATLmorph}.

\item[Associativity:]
Composition is associative by \pagecref{lem:labcomp}.

\item[Identity:]
$\Id_{(\seqname{A}^i, \seqname{E}^i, \seqname{C}^i}$ is an identity
morphism by \cref{lem:labcomp}.
\end{description}
\end{proof}
\end{lemma}

\begin{definition}[Functor from Bundle atlases to m-atlases]
\label{def:FuncBuntoM}
Let \\
$
  \seqname{B}^i
  \defeq
  (
    E^i,
    X^i,
    Y^i,
    \pi^i,
    G^i,
    \rho^i
  )
$,
$i=1,2$, be a protobundle, let $\seqname{A}^i$ be a
$\pi^i$-$G^i$-$\rho^i$-atlas of $E^i$ in the coordinate space
$C^i = X^i \times Y^i$ and let \\*
$
  \funcseqname{f}
  \defeq
  (
    \funcname{f}_E \maps E^1 \to E^2,
    \funcname{f}_X \maps X^1 \to X^2,
    \funcname{f}_Y \maps Y^1 \to Y^2,
    \funcname{f}_G \maps G^1 \to G^2
  )
$
be a $\seqname{B^1}$-$\seqname{B^2}$ bundle-atlas morphism from
$\seqname{A}^1$ to $\seqname{A}^2$. Then
\begin{equation}
\Functor^\Bun_{\Bun,\M} (\seqname{A}^i, \seqname{B}^i)
\defeq
\biggl ( \seqname{A}^i, \Triv{E^i}, \Triv[G^i-\rho^i-]{X^i,Y^i} \biggr )
\end{equation}
\begin{multline}
\Functor^\Bun_{\Bun,\M}
  \bigl (
    \funcseqname{f},
    (\seqname{A}^1, \seqname{B}^1),
    (\seqname{A}^2, \seqname{B}^2)
  \bigr )
\defeq \\
  \Biggl (
    (\funcseqname{f}_E, \funcseqname{f}_X \times \funcseqname{f}_Y),
    \biggl ( \seqname{A}^1, \Triv{E^1}, \Triv[G^1-\rho^1-]{C^1} \biggr ),
    \biggl ( \seqname{A}^2, \Triv{E^2}, \Triv[G^2-\rho^2-]{C^2} \biggr )
  \Biggr )
\end{multline}
\end{definition}

\begin{theorem}[Functor from Bundle atlases to m-atlases]
\label{the:FuncBuntoM}
Let $\seqname{E}$ be a set of topological spaces,
$
  \seqname{B}^\alpha \defeq
  (E^\alpha \in \seqname{E}, X^\alpha, Y^\alpha, \pi^\alpha, G^\alpha, \rho^\alpha)
$,
$\alpha \prec \Alpha$, be a protobundle, \\
$\seqname{B} \defeq \set{{\seqname{B}^\alpha}}[\alpha \prec \Alpha]$
be a set of protobundles and
$
  \seqname{C}^\alpha \defeq
  \Triv[G^\alpha-\rho^\alpha-]{(X^\alpha, Y^\alpha)}
$.
Then $\Functor^\Bun_{\Bun,\M}$ is a functor from
$\maximal[*]{\Atl[\Bun]} \seqname{B}$ to
$
  \maximal[*]
    {
      \Atl
        \Bigl (
          \Trivcat{\seqname{E}}, \Trivcat[\Bun-]{\seqname{B}}
        \Bigr )
    }
$
and a functor from \\*
$\maximal[*]{\Atl[\Bun]} \seqname{B}$ to
$
  \maximal[*]
    {
      \Atl
        \Bigl (
          \Trivcat{\seqname{E}}, \Trivcat[\BunProd-]{\seqname{B}}
        \Bigr )
    }
$.

\begin{proof}
Let $o^i \defeq (\seqname{A}^i, B^i)$, $i \in [1,3]$, be an object of
$\Atl[\Bun](\seqname{B})$,
$
  o'^i \defeq \Functor^\Bun_{\Bun,\M} o^i = \\*
  \Bigl (
    \seqname{A}^i,
    \Triv{E^i},
    \Triv[G^i-\rho^i-]{C^i}
  \Bigr )
$
be the corresponding object of
$
  \Atl
  \Bigl (
    \Trivcat{\seqname{E}},
    \Trivcat[\Bun-]{\seqname{B}}
  \Bigr )
$, \\
$
\funcname{m}^i \defeq
  \bigl (
    \funcseqname{f}^i,
    o^i,
    o^{i+1}
  \bigr )
$, $i=1,2$,
be a morphism of $\Atl[\Bun](\seqname{B})$ from $o^i$ to $o^{i+1}$ and
\begin{align*}
  \funcname{m}'^i
  & \defeq
  \Functor^\Bun_{\Bun,\M} \funcname{m}^i
\\*
  & =
  \Biggl (
    (\funcname{f}^i_E, \funcname{f}^i_X \times \funcname{f}^i_Y),
    \biggl ( \seqname{A}^i, \Triv{E^i}, \Triv[G^i-\rho^i-]{C^i} \biggl ),
    \biggl (
      \seqname{A}^{i+1}, \Triv{E^{i+1}}, \Triv[G^{i+1}-\rho^{i+1}-]{C^2}
    \biggr )
  \Biggr )
\\*
  & =
  \bigl (
    (\funcname{f}^i_E, \funcname{f}^i_X \times \funcname{f}^i_Y),
    o'^i,
    o'^{i+1}
  \bigr )
\end{align*}
be the corresponding morphism of
$\Atl \biggl ( \Trivcat{\seqname{E}}, \Trivcat[\Bun-]{\seqname{B}} \biggr )$.

\begin{description}
\item[Preservation of endpoints:]
$\Triv{E^i} \objin \Trivcat{\seqname{E}}$.
$\Triv[G^i-\rho^i-]{C^i} \arin \Trivcat[\BunProd-]{\seqname{B}}$.

By \pagecref{the:FuncBuntoM},
$(\funcname{f}_E, \funcname{f}_X \times \funcname{f}_Y)$ is a
$\Triv{E^i}$-$\Triv{E^{i+1}}$ m-atlas morphism from $\seqname{A}^i$ to
$\seqname{A}^{i+1}$ in the coordinate spaces
$\Triv[G^i-\rho^i-]{X^i,Y^i}$,
$\Triv[G^{i+1}-\rho^{i+1}-]{X^{i+1},Y^{i+1}}$.

\begin{align*}
  \Functor^\Bun_{\Bun,\M} \funcname{m}^i
  & =
  \Functor^\Bun_{\Bun,\M}
    \bigl (
      \funcseqname{f}^i,
      o^i,
      o^{i+1}
    \bigr )
\\*
  & =
  \bigl (
    (
      \funcname{f}^i_E,
      \funcname{f}^i_X \times \funcname{f}^i_Y
    ),
    o'^i,
    o'^{i+1}
  \bigr )
\end{align*}

\item[Identity:]
\begin{align*}
  \Functor^\Bun_{\Bun,\M} \Id[{{(\seqname{A}^i, \seqname{B}^i)}}]
  & =
  \Functor^\Bun_{\Bun,\M}
    \bigl (
      (\Id[{E^i}], \Id[{X^i}], \Id[{Y^i}], \Id[{G^i}]),
      o^i,
      o^i
    \bigr )
\\*
  & =
  \bigl (
    (\Id[{\Triv{E^i}}], \Id[{C^i}]),
    o'^i,
    o'^i
  \bigr )
\\*
  & =
  \Bigl (
    (\Id[{\Triv{E^i}}], \Id[{C^i}]),
    \Functor^\Bun_{\Bun,\M} o^i,
    \Functor^\Bun_{\Bun,\M} o^i
  \Bigr )
\\*
  & =
  \Id[{\Functor^\Bun_{\Bun,\M} o^i}]
\end{align*}

\item[Composition:]
\begin{align*}
  m^2 \compose[A] m^1
  & =
  \bigl (
    (
      \funcname{f}^2_0 \compose \funcname{f}^1_0,
      \funcname{f}^2_X \compose \funcname{f}^1_X,
      \funcname{f}^2_Y \compose \funcname{f}^1_Y,
      \funcname{f}^2_G \compose \funcname{f}^1_G
    ),
    o^1,
    o^3
  \bigr )
\\*
  \Functor^\Bun_{\Bun,\M} (m^2 \compose m^1)
  & =
    \bigl (
      (
        \funcname{f}_E^2 \compose \funcname{f}_E^1,
        (\funcname{f}_X^2 \compose \funcname{f}_X^1) \times
        (\funcname{f}_Y^2 \compose \funcname{f}_Y^1)
      ),
      o'^1,
      o'^3
    \bigr )
\\*
  & =
    \bigl (
      (
        \funcname{f}_E^2 \compose \funcname{f}_E^1,
        (\funcname{f}^2_X \times \funcname{f}^2_Y) \compose
        (\funcname{f}^1_X \times \funcname{f}^1_Y)
      ),
      o'^1,
      o'^3
    \bigr )
\\*
  & =
  \Functor^\Bun_{\Bun,\M} m^2 \compose[A] \Functor^\Bun_{\Bun,\M} m^1
\end{align*}
\end{description}
\end{proof}
\end{theorem}

\begin{lemma}[Base space functions derived from bundle-atlas morphisms]
\label{lem:BunAtlbasemorph}
Let $\seqname{E}^i$, $i=1,2$, be a model space, $X^i$, $Y^i$ topological
spaces, $G^i$ a topological group, $\rho^i$ an effective action of $G^i$
on $Y^i$, $\seqname{C}^i \defeq (X^i \times Y^i, \catname{XY}^i)$ a
$G^i$-$\rho^i$ model space of $X^i \times Y^i$,
$\funcname{f}_C \maps \seqname{C}^1 \to \seqname{C}^2$ a
$G^1$-$G^2$-$\rho^1$-$\rho^2$ morphism of
$X^1 \times Y^1$ to $X^2 \times Y^2$, i.e., a model function that
preserves group action, $\seqname{A}^i$ an m-atlas of $E^i$ in the
coordinate model space $\seqname{C}^i$ and
$
  \funcseqname{f} \defeq
  (
    \funcname{f}_E \maps \seqname{E^1} \to \seqname{E}^2,
    \funcname{f}_C \maps \seqname{C}^1 \to \seqname{C}^2
  )
$
an $\seqname{E}^1$-$\seqname{E}^2$ m-atlas morphism of $\seqname{A}^1$
to $\seqname{A}^2$ in the coordinate spaces $\seqname{C}^1$,
$\seqname{C}^2$.

Then there exists a unique function $\funcname{f}_X \maps X^1 \to X^2$
such that $\funcname{f}_X \compose \pi_1 = \pi_1 \compose \funcname{f}_C$.

\begin{proof}
Let $x$ in $X^1$, $y,y'$ in $Y^1$, Then
$\funcname{f}_C$ preserves fibers, i.e., \\*
$
\pi_1 \bigl ( \funcname{f}_C(x,y) \bigr ) =
\pi_1 \bigl ( \funcname{f}_C(x,y') \bigr )
$.
by \pagecref{lem:Grhomorph}. Define
$\funcname{f}_X(x) \defeq \pi_1 \bigl ( \funcname{f}_C(x,y) \bigr )$
\end{proof}

\begin{remark}
There need not exist $\funcname{f}_Y \maps Y^1 \to Y^2$ such that
$\funcname{f}_C = \funcname{f}_X \times \funcname{f}_Y$.
\end{remark}
\end{lemma}

\begin{definition}[Functor from m-atlases to Bundle atlases]
\label{def:FuncMtoBun}
Let $\catseqname{E}$ be a trivial model category, $\catname{C}$ be a
trivial product coordinate model category with objects \\*
$
\set
  {
    \seqname{C}^\alpha \defeq
    (X^\alpha \times Y^\alpha, \catname{XY}^\alpha)
  }%
  [\alpha \prec \Alpha]
$,
$G$ a group valued function on $\Ob(\seqname{C})$ and $\rho$ a function
valued function on $\Ob(\seqname{C})$ such that for every
$
  \seqname{C}^\alpha \defeq
  (X^\alpha \times Y^\alpha, \catname{XY}^\alpha) \objin \catname{C}
$,
$\rho(\seqname{C}^\alpha)$ is an effective action of
$G(\seqname{C}^\alpha)$ on $Y^\alpha$ and
$
\Triv%
  [
    G(\seqname{C}^\alpha)-\rho(\seqname{C}^\alpha)-
  ]
  {X^\alpha, Y^\alpha} =
  \seqname{C}^\alpha
$.

Let $\seqname{E}^i \defeq (E^i, \catname{E}^i) \objin \catseqname{E}$,
$i=1,2$, be a trivial model space,
$
  \seqname{C}^i \defeq (X^i \times Y^i, \catname{XY}^i)
  \objin \catname{C}
$
$G^i \defeq G(\seqname{C}^i)$, $\rho^i \defeq \rho(\seqname{C}^i)$,
$\seqname{A}^i$ a full m-atlas of $\seqname{E}^i$ in the coordinate
space $\seqname{C}^i$,
$\pi^i \maps E^i \onto X^i$ the unique function asserted in
\cref{lem:mGrho} and
$
  \seqname{B}^i \defeq
    (
      E^i,
      X^i,
      Y^i,
      \pi^i,
      G^i,
      \rho^i
    )
$.
Then define
\begin{equation}
\Functor^\Bun_{\M-G-\rho,\Bun}
(\seqname{A}^i, \seqname{E}^i, \seqname{C}^i)
\defeq
(\seqname{A}^i, \seqname{B}^i)
\end{equation}

Let
$
  \funcseqname{f} \defeq
  (
    \funcname{f}_E \maps \seqname{E^1} \to \seqname{E}^2,
    \funcname{f}_C \defeq \funcname{f}_X \times \funcname{f}_Y \maps
    \seqname{C}^1 \to \seqname{C}^2
  )
$
be an $\seqname{E}^1$-$\seqname{E}^2$ m-atlas morphism of
$\seqname{A}^1$ to $\seqname{A}^2$ in the coordinate spaces
$\seqname{C}^1, \seqname{C}^2$ that preserves the group action,
$\funcname{f}_G \maps G^1 \to G^2$ the unique function asserted to exist
in \cref{eq:Grho} (preservation of group action) of
\pagecref{def:Grhomorph}.

Then
\begin{multline}
\Functor^\Bun_{\M-G-\rho,\Bun}
  \bigl (
    \funcseqname{f},
    (\seqname{A}^1, \seqname{E^1}, \seqname{C}^1),
    (\seqname{A}^2, \seqname{E^2}, \seqname{C}^2)
  \bigr )
\defeq \\
  \bigl (
    (\funcname{f}_E, \funcname{f}_X, \funcname{f}_Y, \funcname{f}_G),
    (\seqname{A}^1, \seqname{B}^1),
    (\seqname{A}^2, \seqname{B}^2)
  \bigr )
\end{multline}
\end{definition}

\begin{theorem}[Functor from m-atlases to bundle atlases]
\label{the:FuncMtoBun}
Let
\begin{enumerate}
\item
$\catname{C}$ be a trivial product coordinate model category,

\item
$\catname{E}$ be a model category,

\item
$G$ be a group valued function on $\Ob(\catname{C})$

\item
$\rho$ be a function valued function on
$\Ob(\seqname{C})$ such that for every
$
  \seqname{C}^\alpha \defeq
  (X^\alpha \times Y^\alpha, \catname{XY}^\alpha) \objin \catname{C}
$,
$\rho(\seqname{C}^\alpha)$ is an effective action of
$G(\seqname{C}^\alpha)$ on $Y^\alpha$ and
$
  \seqname{C}^\alpha =
  \Triv%
    [
      G(\seqname{C}^\alpha)-\rho(\seqname{C}^\alpha)-
    ]
    {X^\alpha, Y^\alpha}
$.

\item
$\pi$ be the unique function valued function on
$\full{\Atl_\Ob}(\catname{E}, \catname{C})$ such that for every
$
\bigl (
  \seqname{A}^\alpha,
  (E^\alpha,\catname{E}^\alpha),
  (C^\alpha, \catname{C}^\alpha)
  \bigr )
  \in \full{\Atl_\Ob}(\catname{E}, \catname{C})
$ and every $(U, V, \phi) \in \seqname{A}^\alpha$,
$\pi_1 \compose \phi = \pi(\seqname{A}^\alpha) \restrictto_U$

\item
$
  \seqname{B} \defeq
  \set
  {
    \seqname{B}^\alpha \defeq
      \bigl (
        E^\alpha,
        X^\alpha,
        Y^\alpha,
        G(\seqname{C}^\alpha),
        \pi(\seqname{A}^\alpha),
        \rho(\seqname{C}^\alpha)
      \bigr )
  }%
  [
    {
      \Bigl (
        \seqname{A}^\alpha,
        (E^\alpha,\catname{E}^\alpha),
        \bigl (
          C^\alpha \defeq X^\alpha \times Y^\alpha,
          \catname{C}^\alpha
        \bigr )
      \Bigr ) \in \full{\Atl_\Ob}(\catname{E}, \catname{C})
    }
  ]*
$
\end{enumerate}

Then $\Functor^\Bun_{\M-G-\rho,\Bun}$ is a functor from
$
\Atl
  \Bigl (
    \Trivcat{\seqname{E}},
    \Trivcat[\BunProd-]{\seqname{B}}
  \Bigr )
$
to $\Atl[\Bun](\seqname{B})$.

\begin{proof}
Let
$
  o^i \defeq
  \bigl (
    \seqname{A}^i,
    \seqname{E}^i,
    \seqname{C}^i \defeq (X^i \times Y^i, \catname{XY}^i)
  \bigr )
  \objin \catname{M})
$,
$i \in [1,3]$,
$G^i \defeq G(o^i)$, \\
$\pi^i \defeq \pi(o^i)$,
$\rho^i \defeq \rho(o^i)$,
$\seqname{A}^i$ an m-atlas of $E^i$ in the coordinate model space $\seqname{C}^i$, \\
$
  \seqname{B}^i \defeq
    (
      E^i,
      X^i,
      Y^i,
      \pi^i,
      G^i,
      \rho^i
    )
$,
$
  o'^i \defeq
  (
    \seqname{A}^i,
    \seqname{B}^i
  )
$,
$
  \funcname{m}^i \defeq
    \bigl (
      \funcseqname{f}^i,
      o^i,
      o^{i+1}
    \bigr )
    \arin \catname{M}
$, $i=1,2$,
an $E^i$-$E^{i+1}$ m-atlas morphism of $\seqname{A}^i$ to
$\seqname{A}^{i+1}$ in the coordinate spaces
$\seqname{C}^i, \seqname{C}^{i+1}$ that preserves the group action,
$\funcname{f}_G \maps G^i \to G^{i+1}$ the unique function asserted to
exist in \cref{eq:Grho} (preservation of group action) of
\pagecref{def:Grhomorph} and $\pi^i \maps E^i \onto X^i$ the unique
function asserted in \cref{lem:mGrho}.

Let $\funcname{f}^i_G \maps G^i \to G^{i+1}$ be the function asserted by
\pagecref{eq:Grho}; let
$\funcname{f}^i_1 = \funcname{f}^i_X \times \funcname{f}^i_Y$ be the
decomposition given by \pagecref{def:Grhomodtriv}, and let
$
  \funcname{m}'^i \defeq
  \Functor^\Bun_{\M-G-\rho,\Bun} \funcname{m}^i = \\
    \bigl (
      \funcseqname{f}^i,
      (
        \funcname{f}^i_0
        \funcname{f}^i_X
        \funcname{f}^i_Y
        \funcname{f}^i_G
      ),
      o'^i,
      o'^{i+1}
    \bigr )
$.

\begin{description}
\item[Preservation of endpoints:]
$E^i$, $X^i$ and $Y^i$ are topological spaces.

$ G^i$ is a topological group,

$\pi^i \maps E^i \onto X^i$ is a continuous surjection.

$\rho^i \defeq \rho(\seqname{C}^i)$ is an effective action of
$G^i \defeq G(\seqname{C}^i)$ on $Y^\alpha$.

Hence $B^i$ is a protobundle.

\item[Identity:]
\begin{align*}
  \Functor^\Bun_{\M-G-\rho,\Bun}  \Id[{{(\seqname{A}^i, \seqname{B}^i)}}]
  & =
   \Id[{{(\seqname{A}^i, \seqname{B}^i)}}]
    \bigl (
      (\Id[{E^i}], \Id[{X^i}], \Id[{Y^i}], \Id[{G^i}]),
      o^i,
      o^i
    \bigr )
\\*
  & =
  \bigl (
    (\Id[{\Triv{E^i}}], \Id[{C^i}]),
    o'^i,
    o'^i
  \bigr )
\\*
  & =
  \Bigl (
    (\Id[{\Triv{E^i}}], \Id[{C^i}]),
    \Functor^\Bun_{\M-G-\rho,\Bun} o^i,
    \Functor^\Bun_{\M-G-\rho,\Bun} o^i
  \Bigr )
\\*
  & =
  \Id_{\Functor^\Bun_{\M-G-\rho,\Bun} o^i}
\end{align*}

\item[Composition:]
\end{description}
$\Functor^\Bun_{\M-G-\rho,\Bun} \funcname{m}^i$ is a morphism from
$\Functor^\Bun_{\M-G-\rho,\Bun} o^i$ to
$\Functor^\Bun_{\M-G-\rho,\Bun} o^{i+1}$:

\begin{align*}
\Functor^\Bun_{\M-G-\rho,\Bun} (\funcname{m}^i)
  & =
\Functor^\Bun_{\M-G-\rho,\Bun}
  (
    \funcseqname{f}^i,
    o^i,
    o^{i+1}
  )
\\
  & =
  \bigl (
    (
      \funcname{f}^i_0,
      \funcname{f}^i_X,
      \funcname{f}^i_Y,
      \funcname{f}^i_G
    ),
    (\seqname{A}^i, \seqname{B}^i),
    (\seqname{A}^{i+1}, \seqname{B}^{i+1})
  \bigr )
\\
  & =
  \Bigl (
    (
      \funcname{f}^i_0,
      \funcname{f}^i_X,
      \funcname{f}^i_Y,
      \funcname{f}^i_G
    ),
    \Functor^\Bun_{\M-G-\rho,\Bun} o^i,
    \Functor^\Bun_{\M-G-\rho,\Bun} o^{i+1}
  \Bigr )
\end{align*}

$\Functor^\Bun_{\M-G-\rho,\Bun}$ maps identity functions to identity
functions:

\begin{equation}
\begin{split}
  &
\Functor^\Bun_{\M-G-\rho,\Bun} \Id[{o^i}] =
\\
  &
\Functor^\Bun_{\M-G-\rho,\Bun}
  \bigl (
    (\Id[{E^i}], \Id[{C^i}]),
    o^i,
    o^i
  \bigr ) =
\\
  &
  \bigl (
    (\Id[{E^i}], \Id[{X^i}], \Id[{Y^i}], \Id[{G^i}]),
    (\seqname{A}^i, \seqname{B}^i),
    (\seqname{A}^i, \seqname{B}^i)
  \bigr ) =
\\
  &
  \bigl (
    (\Id[{E^i}], \Id[{X^i}], \Id[{Y^i}], \Id[{G^i}]),
    \Functor^\Bun_{\M-G-\rho,\Bun} o^i,
    \Functor^\Bun_{\M-G-\rho,\Bun} o^i
  \bigr ) =
\\
  &
\Id_{\Functor^\Bun_{\M-G-\rho,\Bun} o^i}
\end{split}
\end{equation}

$
  \Functor^\Bun_{\M-G-\rho,\Bun} m^2 \compose[A] \Functor^\Bun_{\M-G-\rho,\Bun} m^1
  =
  \Functor^\Bun_{\M-G-\rho,\Bun} (m^2 \compose[A] m^1)
$:

\begin{enumerate}
\item
$
  m^2 \compose[A] m^1 =
  \bigl (
    (
      \funcname{f}_0^2 \compose \funcname{f}_0^1,
      \funcname{f}_1^2 \compose \funcname{f}_1^1
    ),
    o^1,
    o^3
  \bigr )
  $
\item
$
  \Functor^\Bun_{\M-G-\rho,\Bun}
    \bigl (
       o^i
    \bigr )
  =
  (\seqname{A}^i, \seqname{B}^i)
$
\item
$
  \Functor^\Bun_{\M-G-\rho,\Bun} (m^i)
    = \\
    \bigl (
      (
        \funcname{f}^i_0,
        \funcname{f}^i_X,
        \funcname{f}^i_Y,
        \funcname{f}^i_G
      ),
      (\seqname{A}^i, \seqname{B}^i),
      (\seqname{A}^{i+1}, \seqname{B}^{i+1})
    \bigr )
$
\item
$
  \Functor^\Bun_{\M-G-\rho,\Bun} m^2 \compose[A] \Functor^\Bun_{\M-G-\rho,\Bun} m^1
  = \\
    \bigl (
      (
        \funcname{f}^2_0 \compose \funcname{f}^1_0,
        \funcname{f}^2_X \compose \funcname{f}^1_X,
        \funcname{f}^2_Y \compose \funcname{f}^1_Y,
        \funcname{f}^2_G \compose \funcname{f}^1_G
      ),
      (\seqname{A}^1, \seqname{B}^1),
      (\seqname{A}^3, \seqname{B}^3)
    \bigr )
$
\item
$
  \Functor^\Bun_{\M-G-\rho,\Bun} (m^2 \compose m^1)
  = \\
    \bigl (
      (
        \funcname{f}^2_0 \compose \funcname{f}^1_0,
        \funcname{f}^2_X \compose \funcname{f}^1_X,
        \funcname{f}^2_Y \compose \funcname{f}^1_Y,
        \funcname{f}^2_G \compose \funcname{f}^1_G
      ),
      (\seqname{A}^1, \seqname{B}^1),
      (\seqname{A}^3, \seqname{B}^3)
    \bigr )
$
\end{enumerate}
\end{proof}

$\Functor^\Bun_{\Bun,\M} \compose \Functor^\Bun_{\M-G-\rho,\Bun} = \Id$.
\begin{proof}
Expanding the definitions, we have
\begin{enumerate}
\item
$
  \Functor^\Bun_{\M-G-\rho,\Bun}
  (\seqname{A}^i, \seqname{E}^i, \seqname{C}^i) =
  (\seqname{A}^i, \seqname{B}^i)
$
\item
$
  \Functor^\Bun_{\Bun,\M} (\seqname{A}^i, \seqname{B}^i) =
  (\seqname{A}^i, \Triv{E^i}, \Triv[G^i-\rho^i-]{X^i,Y^i})
  (\seqname{A}^i, \seqname{E}^i, \seqname{C}^i)
$,
since by \cref{def:FuncMtoBun}, $\seqname{E}^i$ and
$\seqname{C}^i$ are trivial.
\item
$
  \Functor^\Bun_{\M-G-\rho,\Bun}
    \bigl (
      (
        \funcseqname{f}^1_E,
        \funcseqname{f}^1_X \times \funcseqname{f}^1_Y
      ),
      (\seqname{A}^1, \seqname{E^1}, \seqname{C}^1),
      (\seqname{A}^2, \seqname{E^2}, \seqname{C}^2)
    \bigr ) = \\
    \bigl (
      (
        \funcname{f}^1_E,
        \funcname{f}^1_X,
        \funcname{f}^1_Y,
        \funcname{f}^1_G
      ),
      (\seqname{A}^1, \seqname{B}^1),
      (\seqname{A}^2, \seqname{B}^2)
    \bigr )
$
\item
$
  \Functor^\Bun_{\Bun,\M}
    \bigl (
      (
        \funcname{f}^1_E,
        \funcname{f}^1_X,
        \funcname{f}^1_Y,
        \funcname{f}^1_G
      ),
      (\seqname{A}^1, \seqname{B}^1),
      (\seqname{A}^2, \seqname{B}^2)
    \bigr ) = \\
    \bigl (
      (
        \funcseqname{f}^1_E,
        \funcseqname{f}^1_X \times \funcseqname{f}^1_Y
      ),
      (\seqname{A}^1, \Triv{E^1}, \Triv[G^1-\rho^1-]{C^1}),
      (\seqname{A}^2, \Triv{E^2}, \Triv[G^2-\rho^2-]{C^2})
    \bigr ) = \\
    \bigl (
      (
        \funcseqname{f}^1_E,
        \funcseqname{f}^1_X \times \funcseqname{f}^1_Y
      ),
      (\seqname{A}^1, \seqname{E}^1, \seqname{C}^1),
      (\seqname{A}^2, \seqname{E}^2, \seqname{C}^2)
    \bigr )
$,
since by \cref{def:FuncMtoBun}, $\seqname{E}^i$ and $\seqname{C}^i$ are trivial.
\end{enumerate}
\end{proof}
\end{theorem}

\section{Associated model spaces and functors}
{
  \showlabelsinline
  \label{sec:assmodB}
}
\begin{definition}[Coordinate model spaces associated with bundle atlases]
\label{def:ismodBF}
Let \\
$
  \seqname{B}^i
  \defeq
  (
    E^i,
    X^i,
    Y^i,
    \pi^i,
    G^i,
    \rho^i
  )
$
be a protobundle
$\seqname{A}^i$ be a $\pi^i$-$G^i$-$\rho^i$-bundle
atlas of $E^i$ in the coordinate space $C^i = X^i \times Y^i$ and
$
  \funcseqname{f} \\
  \defeq
  (
    \funcname{f}_E \maps E^1 \to E^2,
    \funcname{f}_X \maps X^1 \to X^2,
    \funcname{f}_Y \maps Y^1 \to Y^2,
    \funcname{f}_G \maps G^1 \to G^2
  )
$
be a $\seqname{B}^1$-$\seqname{B}^2$ bundle-atlas (near) morphism from
$\seqname{A}^1$ to $\seqname{A}^2$. Then:

The minimal $G$-$\rho$ coordinate model space with neighborhoods in
$\seqname{A}^i$ is $\Functor^\Bun_2 (\seqname{A}^i, \seqname{B}^i)$.

\begin{multline}
\Functor^\Bun_2 (\seqname{A}^i, \seqname{B}^i)
\defeq \\
\minimal{\Mod}
  \Biggl (
    C^i,
    \pi_2[\seqname{A}^i],
    \set
      {\phi' \compose \phi^{-1}}%
      [
        \equant
          {
            {(U,V,\phi) \in \seqname{A}^i},
            {(U',V',\phi') \in \seqname{A}^i}
          }
          {
            U \cap U' \ne \emptyset
          }
      ]
  \Biggr )
\end{multline}

The coordinate mapping associated with the
$\seqname{B^1}$-$\seqname{B^2}$ bundle-atlas (near) morphism
$\funcseqname{f}$ from $\seqname{A^1}$ to $\seqname{A^2}$ is

\begin{equation}
\Functor^\Bun_2
  \bigl (
    \funcseqname{f},
    (\seqname{A}^1, \seqname{B}^1),
    (\seqname{A}^2, \seqname{B}^2)
  \bigr )
\defeq
  \funcname{f}_X \times \funcname{f}_Y
  \maps
  \Functor^\Bun_2 (\seqname{A}^1, \seqname{B}^1)
  \to
  \Functor^\Bun_2 (\seqname{A}^2, \seqname{B}^2)
\end{equation}

If it is a model function then it is also the coordinate
$G^1$-$G^2$-$\rho^1$-$\rho^2$ (near) morphism associated with the
$\seqname{B^1}$-$\seqname{B^2}$ bundle-atlas (near) morphism
$\funcseqname{f}$ from $\seqname{A^1}$ to $\seqname{A^2}$.
\end{definition}

\begin{lemma}[Coordinate model spaces associated with bundle atlases]
{
  \showlabelsinline
  \label{lem:ismodBF}
}
Let \\*
$
  \seqname{B}
  \defeq
  (
    E,
    X,
    Y,
    G,
    \pi,
    \rho
  )
$
and let $\seqname{A}$ be a $\pi$-$G$-$\rho$-bundle atlas of $E$ in the
coordinate space $C = X \times Y$. Then
$\Functor^\Bun_2 (\seqname{A}, \seqname{B})$ is a model space.

\begin{proof}
$\Functor^\Bun_2 (\seqname{A}, \seqname{B})$ satisfies the
conditions for a model space.
for a model space. Let
$
  \catname{C} \defeq
  \Cat \bigl ( \Functor^\Bun_2 (\seqname{A}, \seqname{B}) \bigr )
$.
\begin{enumerate}
\item Since $\pi_2[\seqname{A}]$ is an open cover of
$\union{\pi_2[\seqname{A}]}$, the set of finite intersections is also an
open cover.

\item Finite intersections of finite intersections are finite intersections

\item Restrictions of continuous functions are continuous

\item If $\funcname{f} \maps A \to B$ is a morphism of
$\Functor^\Bun_2 (\seqname{A}, \seqname{B})$ $A, A', B, B'$ objects of
$\catname{C}$ $A' \subseteq A$, $B' \subseteq B$ and
$\funcname{f}[A'] \subseteq B'$ then since $\funcname{f} \maps A \to B$
is a morphism it is a restriction of a transition function between its
restrictions to sets in $\pi_2[\seqname{A}]$ and its restrictions are
also, hence morphisms, and thus
$\funcname{f} \restrictto_{A'} A' \to B'$ is a morphism.

\item
If $(U,V,\phi) \in \seqname{A}$ then $\Id[{V}] = \phi \compose \phi^{-1}A$
is a transition function and hence a morphism of
$\Functor^\Bun_2 (\seqname{A}, \seqname{B})$. If $A, A'$ objects of
$\catname{C}$ and
$A' \subseteq A$ then the inclusion map
$\funcname{i} \maps A' \hookrightarrow A$ is a restriction of an
identity morphism of $\Functor^\Bun_2 (\seqname{A}, \seqname{B})$ and
hence a morphism.

\item Restricted sheaf condition: let
\begin{enumerate}
\item $U_\alpha,V_\alpha$, $\alpha \prec \Alpha$, be an object of
$\catname{C}$

\item $\funcname{f}_\alpha \maps U_\alpha \to V_\alpha$ be a morphism of
$\catname{C}$

\item $U \defeq \union[\alpha \prec \Alpha]{U_\alpha}$

\item $V \defeq \union[\alpha \prec \Alpha]{V_\alpha}$

\item $\funcname{f} \maps U \to V$ be continuous and
$
  \uquant
    {{\alpha \prec \Alpha},{x \in U_\alpha}}
    {\funcname{f}(x) = \funcname{f}_\alpha(x)}
$
\end{enumerate}
Then $\funcname{f}$ is generated by the group action and hence a
morphism of $\catname{C}$
\end{enumerate}
\end{proof}
\end{lemma}

\begin{definition}[Model spaces associated with bundle atlases]
\label{def:ismodT}
Let \\
$
  \seqname{B}^i
  \defeq
  (
    E^i,
    X^i,
    Y^i,
    \pi^i,
    G^i,
    \rho^i
  )
$,
$i=1,2$, and $\seqname{A}^i$ be a $\pi^i$-$G^i$-$\rho^i$-bundle atlas of
$E^i$ in the coordinate space $C^i = X^i \times Y^i$.   Then

The minimal model space with neighborhoods in the
atlas $\seqname{A}^i$ is

\begin{multline}
\Functor^\Bun_1 (\seqname{B}^i, \seqname{A}^i)
\defeq \\
\minimal{\Mod}
  \bigl (
    E^i,
    \pi_1[\seqname{A}^i],
    \set
      {\phi'^{-1} \compose \phi}%
      [
        \equant
          {
            {(U,V,\phi) \in \seqname{A}^i},
            {(U',V',\phi') \in \seqname{A}^i}
          }
          {
            U \cap U' \ne \emptyset
          }
      ]
  \bigr )
\end{multline}

The mapping associated with the $\seqname{B^1}$-$\seqname{B^2}$
bundle-atlas (near) morphism $\funcseqname{f}$ from $\seqname{A^1}$ to
$\seqname{A^2}$ is

\begin{equation}
\Functor^\Bun_1
  \bigl (
    \funcseqname{f},
    (\seqname{A}^1, \seqname{B}^1),
    (\seqname{A}^2, \seqname{B}^2)
  \bigr )
\defeq
  \funcname{f}_E
  \maps
  \Functor^\Bun_1 (\seqname{A}^1, \seqname{B}^1)
  \to
  \Functor^\Bun_1 (\seqname{A}^2, \seqname{B}^2)
\end{equation}

If it is a model function then it is also the m-atlas (near) morphism
associated with the $\seqname{B^1}$-$\seqname{B^2}$ bundle-atlas
(near) morphism $\funcseqname{f}$ from $\seqname{A^1}$ to
$\seqname{A^2}$.
\end{definition}

\begin{lemma}[Model spaces associated with bundle atlases]
\label{lem:ismodT}
Let \\
$
  \seqname{B}
  \defeq
  (
    E,
    X,
    Y,
    G,
    \pi,
    \rho
  )
$
and let $\seqname{A}$ be a $\pi$-$G$-$\rho$-bundle atlas of $E$ in the
coordinate space $C = X \times Y$. Then
$\Functor^\Bun_1 (\seqname{A}, \seqname{B})$ is a model space.

\begin{proof}
\Pagecref{lem:minmod}
\end{proof}
\end{lemma}

\begin{theorem}[Functors from bundle atlases to model spaces]
\label{the:FuncBuntoMod}
Let \\*
$
  \seqname{B}^\alpha \defeq
  (E^\alpha, X^\alpha, Y^\alpha, \pi^\alpha, G^\alpha, \rho^\alpha)
$,
$\alpha \prec \Alpha$, be a protobundle, and
$\seqname{B} \defeq \set{{\seqname{B}^\alpha}}[\alpha \prec \Alpha]$
be a set of protobundles. Then
$\Functor^\Bun_1$ is a functor from
$\Atl[\Bun] \seqname{B}$ to $\Triv{\seqname{E}}$ and
$\Functor^\Bun_2$ is a functor from
$\Atl[\Bun] \seqname{B}$ to $\Trivcat[\Bun-]{\seqname{B}}$.

\begin{proof}
Let $o^i \defeq (\seqname{A}^i, E^i, C^i)$, $i \in [1,3]$, be
objects in
$
\Bun
  (
    \seqname{E},
    \seqname{X},
    \seqname{Y},
    \funcseqname{\pi},
    \funcseqname{G},
    \funcseqname{\rho}
  )
$
and let
$
m^i \defeq
  \bigl (
    (\funcname{f}^i_0, \funcname{f}^i_1),
    o^i,
    o^{i+1}
  \bigr )
$
be a morphism in $\Atl[\Bun](\seqname{E}, \seqname{C})$.

$\Functor^\Bun_1$ is a functor from $\Atl[\Bun] \seqname{B}$ to
$\Triv{\seqname{E}}$.
\begin{description}
\item[Preservation of endpoints:]
\begin{align*}
  \Functor^\Bun_2 \funcname{m}^i
  & =
  \Functor^\Bun_1
    \bigl (
      \funcseqname{f}^i,
      o^i,
      o^{i+1}
    \bigr )
\\*
  & =
  \funcname{f}^i_E \maps \Functor^\Bun_1 o^i \to o^{i+1}
\end{align*}
by \pagecref{def:ismodT}.

\item[Composition:]
\begin{align*}
  \Functor^\Bun_1(\funcname{m}^2 \compose[A] \funcname{m}^1)
  & =
  \Functor^\Bun_1
    \bigl (
      \funcseqname{f}^2 \compose \funcseqname{f}^1
      (\seqname{A}^1, E^1, C^1),
      (\seqname{A}^3, E^3, C^3)
    \bigr )
\\*
  & =
  \funcname{f}^2_E \compose \funcname{f}^1_E \maps
    \Functor^\Bun_1 o^1 \to
    \Functor^\Bun_1 o^3
\\*
  & =
  \bigl (
    \funcname{f}^2_E \maps
    \Functor^\Bun_1 o^2 \to
    \Functor^\Bun_1 o^3
    \bigr )
  \compose
\\*
  & \quad
  \bigl (
    \funcname{f}^1_E \maps
    \Functor^\Bun_1 o^1 \to
    \Functor^\Bun_1 o^2
  \bigr )
\\*
  & =
  \Functor^\Bun_1
    \bigl (
      \funcseqname{f}^2,
      (\seqname{A}^2, E^2, C^2),
      (\seqname{A}^3, E^3, C^3)
    \bigr )
  \compose
\\*
  & \quad
\Functor^\Bun_1
  \bigl (
    \funcseqname{f}^1,
    (\seqname{A}^1, E^1, C^1),
    (\seqname{A}^2, E^2, C^2)
  \bigr )
\\*
  & =
  \Functor^\Bun_1 \funcname{m}^2
  \compose
  \Functor^\Bun_1 \funcname{m}^1
\end{align*}

\item[Identity:]\quad
\begin{enumerate}
\item
$
\Functor^\Bun_1(\Id[{o^i}]) =
\Functor^\Bun_1
  \bigl (
    (\Id[{E^i}], \Id[{C^i}]),
    (\seqname{A}^i, E^i, C^i),
    (\seqname{A}^i, E^i, C^i)
  \bigr )
=
\\
\Id[{E^i}] \maps \Functor^\Bun_1 o^i \to
\Functor^\Bun_1 o^i
$
\item
$
\Id[{\Functor^\Bun_1}] o^i =
\Id[{E^i}] \maps \Functor^\Bun_1 o^i \to
\Functor^\Bun_1 o^i
$
\end{enumerate}
\end{description}

$\Functor^\Bun_2$ is a functor from $\Atl[\Bun] \seqname{B}$ to
$\Trivcat[\Bun-]{\seqname{B}}$.

\begin{description}
\item[Preservation of endpoints:]
$
\Functor^\Bun_2(\funcname{m}^i) =
\funcname{f}^i_1 \maps \Functor^\Bun_2 o^i \to
\Functor^\Bun_2 o^{i+1}
$
\item[Composition:]
\begin{align*}
  \Functor^\Bun_2(\funcname{m}^2 \compose[A] \funcname{m}^1)
  & =
\Functor^\Bun_2
  \bigl (
    (
      \funcname{f}^2_0 \compose \funcname{f}^1_0,
      \funcname{f}^2_1 \compose \funcname{f}^1_1
    )
    (\seqname{A}^1, E^1, C^1),
    (\seqname{A}^3, E^3, C^3)
  \bigr )
\\*
  & =
\funcname{f}^2_0 \compose \funcname{f}^1_0 \maps
  \Functor^\Bun_2 o^1 \to
  \Functor^\Bun_2 o^3
\\*
  & =
\bigl (
  \funcname{f}^2_1 \maps
  \Functor^\Bun_2 o^2 \to
  \Functor^\Bun_2 o^3
  \bigr )
\compose
\bigl (
  \funcname{f}^1_1 \maps
  \Functor^\Bun_2 o^1 \to
  \Functor^\Bun_2 o^2
\bigr )
\\*
  & =
\Functor^\Bun_2
  \bigl (
    (
      \funcname{f}^2_0,
      \funcname{f}^2_1
    )
    (\seqname{A}^2, E^2, C^2),
    (\seqname{A}^3, E^3, C^3)
  \bigr )
\compose
\\*
  & \quad
\Functor^\Bun_2
  \bigl (
    (
      \funcname{f}^1_0,
      \funcname{f}^1_1
    )
    (\seqname{A}^1, E^1, C^1),
    (\seqname{A}^2, E^2, C^2)
  \bigr )
\\*
  & =
\Functor^\Bun_2(\funcname{m}^2)
\compose
\Functor^\Bun_2(\funcname{m}^1)
\end{align*}

\item[Identity:]\quad
\begin{enumerate}
\item
$
\Functor^\Bun_2 \Id[{o^i}] =
\Functor^\Bun_2
  \bigl (
    (\Id[{E^i}], \Id[{C^i}]),
    (\seqname{A}^i, E^i, C^i),
    (\seqname{A}^i, E^i, C^i)
  \bigr )
=
\\
\Id[{C^i}] \maps \Functor^\Bun_2 o^i \to
\Functor^\Bun_2 o^i
$
\item
$
\Id[{\Functor^\Bun_2}] o^i =
\Id[{C^i}] \maps \Functor^\Bun_2 o^i \to
\Functor^\Bun_2 o^i
$
\end{enumerate}
\end{description}
\end{proof}
\end{theorem}

\section{Fiber bundles}
\label{sec:fib}
Conventionally a fiber bundle is different from its atlases, but
\pagecref{def:Bun-Atlcat} encourages treating them on an equal footing. All
of the results for maximal bundle atlases carry directly over to results
for fiber bundles.

\subsection{Definitions}
\begin{definition}[fiber bundles]
\label{def:Bun}
Let $E$, $X$ and $Y$ be topological spaces, $\pi \maps E \onto X$ surjective,
$G$ a topological group,
$\rho \maps Y \times G \to Y$ an effective right action of
$G$ on $Y$ and $\seqname{A}$
a maximal $\pi$-$G$-$\rho$-bundle atlas of $E$ in the coordinate space $X \times Y$.
Then $(E, X, Y, \pi, G, \rho, \seqname{A})$ is a fiber bundle.

Let
$
  \seqname{B}
  \defeq
  \set
    {{
      \seqname{B}^\alpha \defeq
      (
        E^\alpha,
        X^\alpha,
        Y^\alpha,
        \pi^\alpha,
        G^\alpha,
        \rho^\alpha
      )
    }}%
    [
      \alpha \prec \Alpha
    ]
$,
where $E^\alpha, X^\alpha, Y^\alpha$ are topological spaces,
$G^\alpha$ a topological group,
$\pi^\alpha \maps E^\alpha \onto X^\alpha$ surjective and
$\rho^\alpha \maps Y^\alpha \times G^\alpha \to Y^\alpha$
an effective right action of $G^\alpha$ on $Y^\alpha$. Then
\begin{multline}
\Bun_\Ob \seqname{B} \defeq \maximal{\Atl[\Bun]_\Ob} \seqname{B}
\end{multline}
\end{definition}

\begin{definition}[Bundle maps]
\label{def:Bunmorph}
Let \\
$
  \seqname{B}
  \defeq
  \set
    {{
      \seqname{B}^\alpha \defeq
      (
        E^\alpha,
        X^\alpha,
        Y^\alpha,
        \pi^\alpha,
        G^\alpha,
        \rho^\alpha
      )
    }}%
    [
      \alpha \prec \Alpha
    ]
$,
where $E^\alpha, X^\alpha, Y^\alpha$ are topological spaces,
$G^\alpha$ a topological group,
$\pi^\alpha \maps E^\alpha \onto X^\alpha$ surjective and
$\rho^\alpha \maps Y^\alpha \times G^\alpha \to Y^\alpha$
an effective right action of $G^\alpha$ on $Y^\alpha$. Then
\begin{multline}
\Bun_\Ar \seqname{B} \defeq \maximal{\Atl[\Bun]_\Ar} \seqname{B}
\end{multline}
\begin{multline}
\Bun \seqname{B}
\defeq
  \bigl (
    \Bun_\Ob \seqname{B},
    \Bun_\Ar \seqname{B},
    \compose[A]
  \bigr )
\end{multline}

Let $(\seqname{A}^i,\seqname{B}^i) \in \Bun_\Ob \seqname{B}$, $i=1,2$.
Then \\*
$
  \funcseqname{f}
  \defeq
  (
    \funcname{f}_E \maps E^1 \to E^2,
    \funcname{f}_X \maps X^1 \to X^2,
    \funcname{f}_Y \maps Y^1 \to Y^2,
    \funcname{f}_G \maps G^1 \to G^2
  )
$
is a bundle map from $(\seqname{A}^1,\seqname{B}^1)$ to
$(\seqname{A}^2,\seqname{B}^2)$ iff it is a
$\seqname{B^1}$-$\seqname{B^2}$ bundle-atlas morphism from $A^1$ to
$A^2$. The identity morphism for $(\seqname{A}^i,\seqname{B}^i)$ is
\begin{equation}
\Id[{{(\seqname{A}^i,\seqname{B}^i)}}] \defeq
\bigl (
  (\Id[{E^i}], \Id[{X^i}], \Id[{Y^i}], \Id[{G^i}]),
  (\seqname{A}^i,\seqname{B}^i),
  (\seqname{A}^i,\seqname{B}^i)
\bigr )
\end{equation}
\end{definition}

\subsection{Categories of fiber bundles}
\begin{theorem}[Categories of fiber bundles]
\label{the:Bun}
Let \\
$
  \seqname{B}
  \defeq
  \set
    {{
      \seqname{B}^\alpha \defeq
      (
        E^\alpha,
        X^\alpha,
        Y^\alpha,
        G^\alpha,
        \pi^\alpha,
        \rho^\alpha
      )
    }}%
    [
      \alpha \prec \Alpha
    ]
$,
where $E^\alpha, X^\alpha, Y^\alpha$ are topological spaces,
$G^\alpha$ a topological group,
$\pi^\alpha \maps E^\alpha \onto X^\alpha$ surjective and
$\rho^\alpha \maps Y^\alpha \times G^\alpha \to Y^\alpha$
an effective right action of $G^\alpha$ on $Y^\alpha$. Then
$\Bun \seqname{B}$ is a category and $Id_{B^\alpha}$ is the identity
morphism for $B^\alpha$.

\begin{proof}
{
  \showlabelsinline
  The result follows directly from \fullcref{def:Bun},
}
\fullcref{def:Bunmorph} and \pagecref{lem:BunATLiscat}.
\end{proof}
\end{theorem}

\part{Future directions}
\label{part:fut}
If this paradigm proves useful, it can be extended to include a set of
admissible functions on the model neighborhoods of the charts,
possibly using the language of sheaves. That might be desirable for
coordinate spaces more general than Fr\'echet spaces.

Further work is needed to determine whether there is a productive way to
define a notion of classic m-atlas near morphisms.

Further work is needed to determine whether it is productive to define
(semi-)strict m-atlas morphisms in terms of concrete categories over a
category of model spaces rather than directly in terms of a model
category.

The extension of paracompactness to model spaces is intended to be
useful for partitions of unity on fiber bundles. Further work is
needed to determine whether that is actually the case.

The definitions given here include some fairly strong conditions,
e.g., AOC. Further work is needed to determine whether they should be
relaxed for applications beyond manifolds and fiber bundles.

Further work is needed to determine whether the concept of
category-based atlases\footnote{As opposed to pseudogroup based}
of model spaces has general utility.

If the concept of nearly commutative diagrams proves useful, further
work is needed to determine whether a more general definition has
utility.

Further work is needed to determine where it is most fruitful to require
full subcategories and where just subcategories.

Further work is needed to determine conditions for mappings associated
with atlas morphisms to be model functions.

This paper uses the language of category theory as an organizing
principle, but defines various notions concretely with sets. It may be
desirable to abstract away some of the details, in the spirit of, e.g.,
topoi.

\begin{thebibliography}{9}

\bibitem[Ad\'amek, Herrlich, Strecker, 1990] {JoyCat} Ji\v{r}\'i Ad\'amek,
Horst Herrlich,
George E. Strecker.
\textit{Abstract and Concrete Categories The Joy of Cats},
John Wiley and Sons, Inc., 1990.

\bibitem[Kelley, 1955] {GenTop} John L. Kelley,
\textit{General Topology},
D. Van Nostrand Company (first edition), 1955.

\bibitem[Kobayashi, 1996] {FoundDiffGeo1} Shoshichi Kobayashi,
Katsumi Nomizu,
\textit{Foundations of Differential Geometry, Volume I},
ISBN 0-471-15733-3, John Wiley and Sons, Inc., 1996.

%\bibitem[Lee, 2000] {DiffPhys} Jeffrey Marc lee,
%\textit{Differential Geometry},
%Analysis and Physics, 2000.

\bibitem[Mac Lane, 1998] {CftWM} Saunders Mac Lane,
\textit{Categories for the Working Mathemation, 2nd edition},
ISBN 0-387-98403-8, Springer-Verlag, 1998.

\bibitem[Metz, 2018] {LCS} Shmuel (Seymour J.) Metz,
\textit%
{
  Local Coordinate Spaces: a proposed unification of manifolds with
  fiber bundles, and associated machinery
},
\href{https://arxiv.org/pdf/1801.05775}{{\tt arXiv:1801.05775}, 2018}

\bibitem[Steenrod, 1999] {TopFib} Norman Steenrod,
\textit{The Topology Of Fibre Bundles}, ISBN 0-691-00548-6,
Prineton University Press (seventh printing), 1999.

\end{thebibliography}

\end{document}
